% =========================================================================
% Figura (clon de «Tema 2. Figuras», diapositivas «Término cuadrupolar»):
% (a) equipotenciales, (b) equipotenciales y líneas de campo y (c) campo
% eléctrico del término cuadrupolar con simetría de revolución,
% phi_2 ~ Q_zz P2(cos t)/r^3, Q_zz > 0. Datos: figuras/datos/cua_puntual_*.
% =========================================================================
\begin{figure}[t]
\centering
\newcommand{\icono}{% dos dipolos opuestos unidos por su carga negativa
  \draw[-{Stealth[length=2.8pt,width=2.6pt]}, mmarojo, line width=0.9pt] (0,0) -- (0,0.95);
  \draw[-{Stealth[length=2.8pt,width=2.6pt]}, mmarojo, line width=0.9pt] (0,0) -- (0,-0.95);
  \fill[mmarojo] (0,0) circle (0.2);}
\begin{minipage}[t]{0.325\textwidth}\centering
{\footnotesize Equipotenciales}\\[3pt]
\begin{tikzpicture}[x=0.198cm, y=0.198cm]
  \renewcommand{\EQ}[1]{\draw[mmaEQsolo] plot coordinates {#1};}%
  \colorlet{marcoc}{grispizarra!60}
  \panelMMA{10}{white}{0}{% Cuadrupolo puntual - equipotenciales
% Generado por scripts/generar_datos.py (no editar a mano).
\EQ{(-0.220,-0.155) (-0.275,-0.195) (-0.319,-0.225) (-0.375,-0.266) (-0.425,-0.300) (-0.475,-0.336) (-0.525,-0.371) (-0.568,-0.400) (-0.625,-0.441) (-0.675,-0.475) (-0.725,-0.510) (-0.775,-0.544) (-0.821,-0.575) (-0.875,-0.612) (-0.925,-0.646) (-0.969,-0.675) (-1.025,-0.713) (-1.075,-0.745) (-1.121,-0.775) (-1.175,-0.810) (-1.225,-0.841) (-1.275,-0.872) (-1.321,-0.900) (-1.375,-0.933) (-1.425,-0.963) (-1.475,-0.991) (-1.525,-1.020) (-1.575,-1.047) (-1.625,-1.074) (-1.674,-1.100) (-1.723,-1.125) (-1.774,-1.150) (-1.825,-1.175) (-1.875,-1.198) (-1.925,-1.220) (-1.975,-1.241) (-2.025,-1.262) (-2.075,-1.281) (-2.125,-1.300) (-2.175,-1.317) (-2.225,-1.334) (-2.275,-1.349) (-2.325,-1.364) (-2.375,-1.377) (-2.425,-1.389) (-2.475,-1.400) (-2.525,-1.410) (-2.575,-1.419) (-2.625,-1.426) (-2.675,-1.433) (-2.725,-1.438) (-2.775,-1.441) (-2.825,-1.444) (-2.875,-1.445) (-2.925,-1.445) (-2.975,-1.443) (-3.025,-1.440) (-3.075,-1.435) (-3.125,-1.429) (-3.175,-1.421) (-3.225,-1.412) (-3.275,-1.401) (-3.325,-1.388) (-3.375,-1.373) (-3.425,-1.357) (-3.475,-1.338) (-3.525,-1.318) (-3.575,-1.295) (-3.625,-1.270) (-3.675,-1.242) (-3.725,-1.212) (-3.775,-1.179) (-3.825,-1.142) (-3.875,-1.102) (-3.925,-1.058) (-3.975,-1.010) (-4.025,-0.956) (-4.072,-0.900) (-4.110,-0.850) (-4.145,-0.800) (-4.175,-0.753) (-4.206,-0.700) (-4.232,-0.650) (-4.256,-0.600) (-4.277,-0.550) (-4.297,-0.500) (-4.314,-0.450) (-4.329,-0.400) (-4.343,-0.350) (-4.354,-0.300) (-4.364,-0.250) (-4.371,-0.200) (-4.377,-0.150) (-4.382,-0.100) (-4.384,-0.050) (-4.385,0.000) (-4.384,0.050) (-4.382,0.100) (-4.377,0.150) (-4.371,0.200) (-4.364,0.250) (-4.354,0.300) (-4.343,0.350) (-4.329,0.400) (-4.314,0.450) (-4.297,0.500) (-4.277,0.550) (-4.256,0.600) (-4.232,0.650) (-4.206,0.700) (-4.177,0.750) (-4.150,0.792) (-4.110,0.850) (-4.075,0.896) (-4.030,0.950) (-3.985,1.000) (-3.934,1.050) (-3.878,1.100) (-3.825,1.142) (-3.780,1.175) (-3.725,1.212) (-3.675,1.242) (-3.625,1.270) (-3.575,1.295) (-3.525,1.318) (-3.475,1.338) (-3.425,1.357) (-3.375,1.373) (-3.325,1.388) (-3.275,1.401) (-3.225,1.412) (-3.175,1.421) (-3.125,1.429) (-3.075,1.435) (-3.025,1.440) (-2.975,1.443) (-2.925,1.445) (-2.875,1.445) (-2.825,1.444) (-2.775,1.441) (-2.725,1.438) (-2.675,1.433) (-2.625,1.426) (-2.575,1.419) (-2.525,1.410) (-2.475,1.400) (-2.425,1.389) (-2.375,1.377) (-2.325,1.364) (-2.275,1.349) (-2.225,1.334) (-2.175,1.317) (-2.126,1.300) (-2.075,1.281) (-2.025,1.262) (-1.975,1.241) (-1.925,1.220) (-1.875,1.198) (-1.826,1.175) (-1.775,1.151) (-1.725,1.126) (-1.675,1.100) (-1.627,1.075) (-1.580,1.050) (-1.535,1.025) (-1.491,1.000) (-1.447,0.975) (-1.400,0.948) (-1.350,0.918) (-1.300,0.888) (-1.250,0.857) (-1.200,0.825) (-1.150,0.794) (-1.100,0.762) (-1.050,0.729) (-1.007,0.700) (-0.950,0.663) (-0.900,0.629) (-0.858,0.600) (-0.800,0.561) (-0.750,0.527) (-0.700,0.493) (-0.650,0.458) (-0.604,0.425) (-0.550,0.388) (-0.500,0.353) (-0.450,0.318) (-0.400,0.283) (-0.354,0.250) (-0.300,0.213) (-0.250,0.177) (-0.220,0.155)}
\EQ{(0.206,0.144) (0.248,0.175) (0.300,0.213) (0.350,0.248) (0.400,0.283) (0.450,0.318) (0.496,0.350) (0.550,0.388) (0.600,0.423) (0.650,0.458) (0.700,0.493) (0.748,0.525) (0.800,0.561) (0.850,0.595) (0.895,0.625) (0.950,0.663) (1.000,0.696) (1.045,0.725) (1.100,0.762) (1.150,0.794) (1.200,0.825) (1.250,0.857) (1.300,0.888) (1.350,0.918) (1.400,0.948) (1.447,0.975) (1.491,1.000) (1.535,1.025) (1.580,1.050) (1.625,1.074) (1.674,1.100) (1.723,1.125) (1.774,1.150) (1.825,1.175) (1.875,1.198) (1.925,1.220) (1.975,1.241) (2.025,1.262) (2.075,1.281) (2.125,1.300) (2.175,1.317) (2.225,1.334) (2.275,1.349) (2.325,1.364) (2.375,1.377) (2.425,1.389) (2.475,1.400) (2.525,1.410) (2.575,1.419) (2.625,1.426) (2.675,1.433) (2.725,1.438) (2.775,1.441) (2.825,1.444) (2.875,1.445) (2.925,1.445) (2.975,1.443) (3.025,1.440) (3.075,1.435) (3.125,1.429) (3.175,1.421) (3.225,1.412) (3.275,1.401) (3.325,1.388) (3.375,1.373) (3.425,1.357) (3.475,1.338) (3.525,1.318) (3.575,1.295) (3.625,1.270) (3.675,1.242) (3.725,1.212) (3.775,1.179) (3.825,1.142) (3.875,1.102) (3.925,1.058) (3.975,1.010) (4.025,0.956) (4.072,0.900) (4.110,0.850) (4.145,0.800) (4.175,0.753) (4.206,0.700) (4.232,0.650) (4.256,0.600) (4.277,0.550) (4.297,0.500) (4.314,0.450) (4.329,0.400) (4.343,0.350) (4.354,0.300) (4.364,0.250) (4.371,0.200) (4.377,0.150) (4.382,0.100) (4.384,0.050) (4.385,0.000) (4.384,-0.050) (4.382,-0.100) (4.377,-0.150) (4.371,-0.200) (4.364,-0.250) (4.354,-0.300) (4.343,-0.350) (4.329,-0.400) (4.314,-0.450) (4.297,-0.500) (4.277,-0.550) (4.256,-0.600) (4.232,-0.650) (4.206,-0.700) (4.177,-0.750) (4.150,-0.792) (4.110,-0.850) (4.075,-0.896) (4.030,-0.950) (3.985,-1.000) (3.934,-1.050) (3.878,-1.100) (3.825,-1.142) (3.780,-1.175) (3.725,-1.212) (3.675,-1.242) (3.625,-1.270) (3.575,-1.295) (3.525,-1.318) (3.475,-1.338) (3.425,-1.357) (3.375,-1.373) (3.325,-1.388) (3.275,-1.401) (3.225,-1.412) (3.175,-1.421) (3.125,-1.429) (3.075,-1.435) (3.025,-1.440) (2.975,-1.443) (2.925,-1.445) (2.875,-1.445) (2.825,-1.444) (2.775,-1.441) (2.725,-1.438) (2.675,-1.433) (2.625,-1.426) (2.575,-1.419) (2.525,-1.410) (2.475,-1.400) (2.425,-1.389) (2.375,-1.377) (2.325,-1.364) (2.275,-1.349) (2.225,-1.334) (2.175,-1.317) (2.126,-1.300) (2.075,-1.281) (2.025,-1.262) (1.975,-1.241) (1.925,-1.220) (1.875,-1.198) (1.826,-1.175) (1.775,-1.151) (1.725,-1.126) (1.675,-1.100) (1.627,-1.075) (1.580,-1.050) (1.535,-1.025) (1.491,-1.000) (1.447,-0.975) (1.400,-0.948) (1.350,-0.918) (1.300,-0.888) (1.250,-0.857) (1.200,-0.825) (1.150,-0.794) (1.100,-0.762) (1.050,-0.729) (1.007,-0.700) (0.950,-0.663) (0.900,-0.629) (0.858,-0.600) (0.800,-0.561) (0.750,-0.527) (0.700,-0.493) (0.650,-0.458) (0.604,-0.425) (0.550,-0.388) (0.500,-0.353) (0.450,-0.318) (0.400,-0.283) (0.354,-0.250) (0.300,-0.213) (0.250,-0.177) (0.206,-0.144)}
\EQ{(-0.220,-0.155) (-0.275,-0.195) (-0.319,-0.225) (-0.375,-0.266) (-0.425,-0.300) (-0.475,-0.336) (-0.525,-0.371) (-0.568,-0.400) (-0.625,-0.441) (-0.675,-0.475) (-0.725,-0.510) (-0.775,-0.545) (-0.820,-0.575) (-0.875,-0.613) (-0.925,-0.647) (-0.968,-0.675) (-1.025,-0.714) (-1.075,-0.747) (-1.119,-0.775) (-1.175,-0.812) (-1.225,-0.844) (-1.275,-0.875) (-1.325,-0.906) (-1.375,-0.937) (-1.425,-0.967) (-1.475,-0.996) (-1.525,-1.025) (-1.569,-1.050) (-1.614,-1.075) (-1.660,-1.100) (-1.707,-1.125) (-1.755,-1.150) (-1.800,-1.173) (-1.850,-1.197) (-1.900,-1.221) (-1.950,-1.244) (-2.000,-1.266) (-2.050,-1.288) (-2.100,-1.308) (-2.150,-1.327) (-2.200,-1.346) (-2.250,-1.363) (-2.300,-1.380) (-2.350,-1.396) (-2.400,-1.410) (-2.450,-1.424) (-2.500,-1.436) (-2.550,-1.447) (-2.600,-1.458) (-2.650,-1.467) (-2.700,-1.475) (-2.750,-1.481) (-2.800,-1.487) (-2.850,-1.491) (-2.900,-1.494) (-2.950,-1.496) (-3.000,-1.497) (-3.050,-1.496) (-3.100,-1.493) (-3.150,-1.490) (-3.200,-1.485) (-3.250,-1.478) (-3.300,-1.470) (-3.350,-1.460) (-3.400,-1.449) (-3.450,-1.436) (-3.500,-1.421) (-3.550,-1.405) (-3.600,-1.386) (-3.650,-1.365) (-3.700,-1.343) (-3.750,-1.318) (-3.800,-1.290) (-3.850,-1.260) (-3.900,-1.228) (-3.950,-1.192) (-4.000,-1.153) (-4.050,-1.110) (-4.088,-1.075) (-4.125,-1.038) (-4.160,-1.000) (-4.200,-0.954) (-4.243,-0.900) (-4.275,-0.855) (-4.312,-0.800) (-4.342,-0.750) (-4.369,-0.700) (-4.395,-0.650) (-4.417,-0.600) (-4.438,-0.550) (-4.457,-0.500) (-4.473,-0.450) (-4.488,-0.400) (-4.501,-0.350) (-4.512,-0.300) (-4.521,-0.250) (-4.529,-0.200) (-4.534,-0.150) (-4.539,-0.100) (-4.541,-0.050) (-4.542,0.000) (-4.541,0.050) (-4.539,0.100) (-4.534,0.150) (-4.529,0.200) (-4.521,0.250) (-4.512,0.300) (-4.501,0.350) (-4.488,0.400) (-4.473,0.450) (-4.457,0.500) (-4.438,0.550) (-4.417,0.600) (-4.395,0.650) (-4.369,0.700) (-4.342,0.750) (-4.312,0.800) (-4.279,0.850) (-4.243,0.900) (-4.203,0.950) (-4.160,1.000) (-4.125,1.038) (-4.088,1.075) (-4.050,1.110) (-4.003,1.150) (-3.950,1.192) (-3.904,1.225) (-3.850,1.260) (-3.800,1.290) (-3.750,1.318) (-3.700,1.343) (-3.650,1.365) (-3.600,1.386) (-3.550,1.405) (-3.500,1.421) (-3.450,1.436) (-3.400,1.449) (-3.350,1.460) (-3.300,1.470) (-3.250,1.478) (-3.200,1.485) (-3.150,1.490) (-3.100,1.493) (-3.050,1.496) (-3.000,1.497) (-2.950,1.496) (-2.900,1.494) (-2.850,1.491) (-2.800,1.487) (-2.750,1.481) (-2.700,1.475) (-2.650,1.467) (-2.600,1.458) (-2.550,1.447) (-2.500,1.436) (-2.450,1.424) (-2.400,1.410) (-2.350,1.396) (-2.300,1.380) (-2.250,1.363) (-2.200,1.346) (-2.150,1.327) (-2.100,1.308) (-2.050,1.288) (-2.000,1.266) (-1.950,1.244) (-1.900,1.221) (-1.850,1.197) (-1.804,1.175) (-1.755,1.150) (-1.707,1.125) (-1.660,1.100) (-1.614,1.075) (-1.569,1.050) (-1.525,1.025) (-1.482,1.000) (-1.425,0.967) (-1.375,0.937) (-1.325,0.906) (-1.275,0.875) (-1.225,0.844) (-1.175,0.812) (-1.125,0.779) (-1.081,0.750) (-1.025,0.714) (-0.975,0.680) (-0.930,0.650) (-0.875,0.613) (-0.825,0.579) (-0.784,0.550) (-0.725,0.510) (-0.675,0.475) (-0.625,0.441) (-0.575,0.406) (-0.532,0.375) (-0.475,0.336) (-0.425,0.300) (-0.375,0.266) (-0.325,0.230) (-0.284,0.200) (-0.225,0.160) (-0.220,0.155)}
\EQ{(0.206,0.144) (0.248,0.175) (0.300,0.213) (0.350,0.248) (0.400,0.283) (0.450,0.318) (0.496,0.350) (0.550,0.388) (0.600,0.423) (0.650,0.458) (0.700,0.493) (0.747,0.525) (0.800,0.562) (0.850,0.596) (0.894,0.625) (0.950,0.664) (1.000,0.697) (1.043,0.725) (1.100,0.763) (1.150,0.796) (1.196,0.825) (1.250,0.859) (1.300,0.891) (1.350,0.921) (1.397,0.950) (1.450,0.982) (1.500,1.011) (1.550,1.040) (1.600,1.068) (1.650,1.095) (1.700,1.122) (1.750,1.148) (1.800,1.173) (1.850,1.197) (1.900,1.221) (1.950,1.244) (2.000,1.266) (2.050,1.288) (2.100,1.308) (2.150,1.327) (2.200,1.346) (2.250,1.363) (2.300,1.380) (2.350,1.396) (2.400,1.410) (2.450,1.424) (2.500,1.436) (2.550,1.447) (2.600,1.458) (2.650,1.467) (2.700,1.475) (2.750,1.481) (2.800,1.487) (2.850,1.491) (2.900,1.494) (2.950,1.496) (3.000,1.497) (3.050,1.496) (3.100,1.493) (3.150,1.490) (3.200,1.485) (3.250,1.478) (3.300,1.470) (3.350,1.460) (3.400,1.449) (3.450,1.436) (3.500,1.421) (3.550,1.405) (3.600,1.386) (3.650,1.365) (3.700,1.343) (3.750,1.318) (3.800,1.290) (3.850,1.260) (3.900,1.228) (3.950,1.192) (4.000,1.153) (4.050,1.110) (4.088,1.075) (4.125,1.038) (4.160,1.000) (4.200,0.954) (4.243,0.900) (4.275,0.855) (4.312,0.800) (4.342,0.750) (4.369,0.700) (4.395,0.650) (4.417,0.600) (4.438,0.550) (4.457,0.500) (4.473,0.450) (4.488,0.400) (4.501,0.350) (4.512,0.300) (4.521,0.250) (4.529,0.200) (4.534,0.150) (4.539,0.100) (4.541,0.050) (4.542,0.000) (4.541,-0.050) (4.539,-0.100) (4.534,-0.150) (4.529,-0.200) (4.521,-0.250) (4.512,-0.300) (4.501,-0.350) (4.488,-0.400) (4.473,-0.450) (4.457,-0.500) (4.438,-0.550) (4.417,-0.600) (4.395,-0.650) (4.369,-0.700) (4.342,-0.750) (4.312,-0.800) (4.279,-0.850) (4.243,-0.900) (4.203,-0.950) (4.160,-1.000) (4.125,-1.038) (4.088,-1.075) (4.050,-1.110) (4.003,-1.150) (3.950,-1.192) (3.904,-1.225) (3.850,-1.260) (3.800,-1.290) (3.750,-1.318) (3.700,-1.343) (3.650,-1.365) (3.600,-1.386) (3.550,-1.405) (3.500,-1.421) (3.450,-1.436) (3.400,-1.449) (3.350,-1.460) (3.300,-1.470) (3.250,-1.478) (3.200,-1.485) (3.150,-1.490) (3.100,-1.493) (3.050,-1.496) (3.000,-1.497) (2.950,-1.496) (2.900,-1.494) (2.850,-1.491) (2.800,-1.487) (2.750,-1.481) (2.700,-1.475) (2.650,-1.467) (2.600,-1.458) (2.550,-1.447) (2.500,-1.436) (2.450,-1.424) (2.400,-1.410) (2.350,-1.396) (2.300,-1.380) (2.250,-1.363) (2.200,-1.346) (2.150,-1.327) (2.100,-1.308) (2.050,-1.288) (2.000,-1.266) (1.950,-1.244) (1.900,-1.221) (1.850,-1.197) (1.804,-1.175) (1.755,-1.150) (1.707,-1.125) (1.660,-1.100) (1.614,-1.075) (1.569,-1.050) (1.525,-1.025) (1.482,-1.000) (1.425,-0.967) (1.375,-0.937) (1.325,-0.906) (1.275,-0.875) (1.225,-0.844) (1.175,-0.812) (1.125,-0.779) (1.081,-0.750) (1.025,-0.714) (0.975,-0.680) (0.930,-0.650) (0.875,-0.613) (0.825,-0.579) (0.784,-0.550) (0.725,-0.510) (0.675,-0.475) (0.625,-0.441) (0.575,-0.406) (0.532,-0.375) (0.475,-0.336) (0.425,-0.300) (0.375,-0.266) (0.325,-0.230) (0.284,-0.200) (0.225,-0.160) (0.206,-0.144)}
\EQ{(-0.220,-0.155) (-0.275,-0.195) (-0.319,-0.225) (-0.375,-0.266) (-0.425,-0.300) (-0.475,-0.336) (-0.525,-0.371) (-0.568,-0.400) (-0.625,-0.441) (-0.674,-0.475) (-0.725,-0.510) (-0.775,-0.545) (-0.819,-0.575) (-0.875,-0.614) (-0.925,-0.647) (-0.975,-0.681) (-1.025,-0.715) (-1.075,-0.748) (-1.125,-0.781) (-1.175,-0.814) (-1.225,-0.846) (-1.271,-0.875) (-1.325,-0.910) (-1.375,-0.941) (-1.425,-0.971) (-1.473,-1.000) (-1.525,-1.031) (-1.575,-1.060) (-1.625,-1.089) (-1.675,-1.116) (-1.725,-1.144) (-1.775,-1.170) (-1.825,-1.196) (-1.875,-1.221) (-1.925,-1.246) (-1.975,-1.270) (-2.025,-1.293) (-2.075,-1.315) (-2.125,-1.336) (-2.175,-1.357) (-2.222,-1.375) (-2.275,-1.395) (-2.325,-1.413) (-2.375,-1.430) (-2.425,-1.446) (-2.475,-1.461) (-2.525,-1.474) (-2.575,-1.487) (-2.625,-1.499) (-2.675,-1.510) (-2.725,-1.520) (-2.775,-1.529) (-2.825,-1.536) (-2.875,-1.542) (-2.925,-1.548) (-2.975,-1.552) (-3.025,-1.555) (-3.075,-1.556) (-3.125,-1.556) (-3.175,-1.556) (-3.225,-1.553) (-3.275,-1.550) (-3.325,-1.545) (-3.375,-1.538) (-3.425,-1.530) (-3.475,-1.521) (-3.525,-1.510) (-3.575,-1.497) (-3.625,-1.483) (-3.675,-1.467) (-3.725,-1.449) (-3.775,-1.429) (-3.825,-1.407) (-3.875,-1.383) (-3.925,-1.357) (-3.975,-1.329) (-4.021,-1.300) (-4.075,-1.264) (-4.125,-1.227) (-4.175,-1.187) (-4.217,-1.150) (-4.270,-1.100) (-4.318,-1.050) (-4.362,-1.000) (-4.400,-0.952) (-4.439,-0.900) (-4.473,-0.850) (-4.500,-0.807) (-4.533,-0.750) (-4.559,-0.700) (-4.583,-0.650) (-4.605,-0.600) (-4.624,-0.550) (-4.642,-0.500) (-4.658,-0.450) (-4.672,-0.400) (-4.684,-0.350) (-4.695,-0.300) (-4.704,-0.250) (-4.711,-0.200) (-4.717,-0.150) (-4.720,-0.100) (-4.723,-0.050) (-4.724,0.000) (-4.723,0.050) (-4.720,0.100) (-4.717,0.150) (-4.711,0.200) (-4.704,0.250) (-4.695,0.300) (-4.684,0.350) (-4.672,0.400) (-4.658,0.450) (-4.642,0.500) (-4.625,0.548) (-4.605,0.600) (-4.583,0.650) (-4.559,0.700) (-4.533,0.750) (-4.504,0.800) (-4.475,0.847) (-4.439,0.900) (-4.402,0.950) (-4.362,1.000) (-4.325,1.042) (-4.275,1.095) (-4.225,1.143) (-4.175,1.187) (-4.127,1.225) (-4.075,1.264) (-4.025,1.298) (-3.981,1.325) (-3.925,1.357) (-3.875,1.383) (-3.825,1.407) (-3.775,1.429) (-3.725,1.449) (-3.675,1.467) (-3.625,1.483) (-3.575,1.497) (-3.525,1.510) (-3.475,1.521) (-3.425,1.530) (-3.375,1.538) (-3.325,1.545) (-3.275,1.550) (-3.225,1.553) (-3.175,1.556) (-3.125,1.556) (-3.075,1.556) (-3.025,1.555) (-2.975,1.552) (-2.925,1.548) (-2.875,1.542) (-2.825,1.536) (-2.775,1.529) (-2.725,1.520) (-2.675,1.510) (-2.625,1.499) (-2.575,1.487) (-2.525,1.474) (-2.475,1.461) (-2.425,1.446) (-2.375,1.430) (-2.325,1.413) (-2.275,1.395) (-2.225,1.376) (-2.175,1.357) (-2.125,1.336) (-2.075,1.315) (-2.025,1.293) (-1.975,1.270) (-1.925,1.246) (-1.875,1.221) (-1.825,1.196) (-1.775,1.170) (-1.725,1.144) (-1.675,1.116) (-1.625,1.089) (-1.575,1.060) (-1.525,1.031) (-1.475,1.001) (-1.432,0.975) (-1.375,0.941) (-1.325,0.910) (-1.275,0.878) (-1.232,0.850) (-1.175,0.814) (-1.125,0.781) (-1.078,0.750) (-1.025,0.715) (-0.975,0.681) (-0.929,0.650) (-0.875,0.614) (-0.825,0.579) (-0.783,0.550) (-0.725,0.510) (-0.675,0.475) (-0.625,0.441) (-0.575,0.406) (-0.532,0.375) (-0.475,0.336) (-0.425,0.300) (-0.375,0.266) (-0.325,0.230) (-0.284,0.200) (-0.225,0.160) (-0.220,0.155)}
\EQ{(0.206,0.144) (0.248,0.175) (0.300,0.213) (0.350,0.248) (0.400,0.283) (0.450,0.318) (0.496,0.350) (0.550,0.389) (0.600,0.423) (0.650,0.458) (0.700,0.493) (0.747,0.525) (0.800,0.562) (0.850,0.596) (0.893,0.625) (0.950,0.665) (1.000,0.698) (1.050,0.732) (1.100,0.765) (1.150,0.797) (1.193,0.825) (1.250,0.862) (1.300,0.894) (1.350,0.925) (1.400,0.956) (1.450,0.986) (1.500,1.016) (1.550,1.046) (1.600,1.074) (1.646,1.100) (1.691,1.125) (1.737,1.150) (1.785,1.175) (1.833,1.200) (1.883,1.225) (1.934,1.250) (1.987,1.275) (2.042,1.300) (2.099,1.325) (2.150,1.346) (2.200,1.367) (2.250,1.386) (2.300,1.404) (2.350,1.421) (2.400,1.438) (2.450,1.453) (2.500,1.468) (2.550,1.481) (2.600,1.493) (2.650,1.505) (2.700,1.515) (2.750,1.524) (2.800,1.532) (2.850,1.539) (2.900,1.545) (2.950,1.550) (3.000,1.553) (3.050,1.556) (3.100,1.556) (3.150,1.556) (3.200,1.555) (3.250,1.552) (3.300,1.547) (3.350,1.542) (3.400,1.534) (3.450,1.526) (3.500,1.515) (3.550,1.504) (3.600,1.490) (3.650,1.475) (3.700,1.458) (3.750,1.439) (3.800,1.418) (3.850,1.396) (3.900,1.371) (3.950,1.343) (4.000,1.313) (4.050,1.281) (4.094,1.250) (4.150,1.207) (4.189,1.175) (4.244,1.125) (4.294,1.075) (4.340,1.025) (4.375,0.984) (4.421,0.925) (4.450,0.884) (4.489,0.825) (4.519,0.775) (4.546,0.725) (4.571,0.675) (4.594,0.625) (4.615,0.575) (4.633,0.525) (4.650,0.475) (4.665,0.425) (4.678,0.375) (4.690,0.325) (4.699,0.275) (4.708,0.225) (4.714,0.175) (4.719,0.125) (4.722,0.075) (4.723,0.025) (4.723,-0.050) (4.720,-0.100) (4.717,-0.150) (4.711,-0.200) (4.704,-0.250) (4.695,-0.300) (4.684,-0.350) (4.672,-0.400) (4.658,-0.450) (4.642,-0.500) (4.625,-0.548) (4.605,-0.600) (4.583,-0.650) (4.559,-0.700) (4.533,-0.750) (4.504,-0.800) (4.475,-0.847) (4.439,-0.900) (4.402,-0.950) (4.362,-1.000) (4.325,-1.042) (4.275,-1.095) (4.225,-1.143) (4.175,-1.187) (4.127,-1.225) (4.075,-1.264) (4.025,-1.298) (3.981,-1.325) (3.925,-1.357) (3.875,-1.383) (3.825,-1.407) (3.775,-1.429) (3.725,-1.449) (3.675,-1.467) (3.625,-1.483) (3.575,-1.497) (3.525,-1.510) (3.475,-1.521) (3.425,-1.530) (3.375,-1.538) (3.325,-1.545) (3.275,-1.550) (3.225,-1.553) (3.175,-1.556) (3.125,-1.556) (3.075,-1.556) (3.025,-1.555) (2.975,-1.552) (2.925,-1.548) (2.875,-1.542) (2.825,-1.536) (2.775,-1.529) (2.725,-1.520) (2.675,-1.510) (2.625,-1.499) (2.575,-1.487) (2.525,-1.474) (2.475,-1.461) (2.425,-1.446) (2.375,-1.430) (2.325,-1.413) (2.275,-1.395) (2.225,-1.376) (2.175,-1.357) (2.125,-1.336) (2.075,-1.315) (2.025,-1.293) (1.975,-1.270) (1.925,-1.246) (1.875,-1.221) (1.825,-1.196) (1.775,-1.170) (1.725,-1.144) (1.675,-1.116) (1.625,-1.089) (1.575,-1.060) (1.525,-1.031) (1.475,-1.001) (1.432,-0.975) (1.375,-0.941) (1.325,-0.910) (1.275,-0.878) (1.232,-0.850) (1.175,-0.814) (1.125,-0.781) (1.078,-0.750) (1.025,-0.715) (0.975,-0.681) (0.929,-0.650) (0.875,-0.614) (0.825,-0.579) (0.783,-0.550) (0.725,-0.510) (0.675,-0.475) (0.625,-0.441) (0.575,-0.406) (0.532,-0.375) (0.475,-0.336) (0.425,-0.300) (0.375,-0.266) (0.325,-0.230) (0.284,-0.200) (0.225,-0.160) (0.206,-0.144)}
\EQ{(-0.220,-0.155) (-0.275,-0.195) (-0.319,-0.225) (-0.375,-0.266) (-0.425,-0.300) (-0.475,-0.336) (-0.525,-0.371) (-0.568,-0.400) (-0.625,-0.441) (-0.674,-0.475) (-0.725,-0.511) (-0.775,-0.545) (-0.819,-0.575) (-0.875,-0.614) (-0.925,-0.648) (-0.975,-0.682) (-1.025,-0.716) (-1.075,-0.750) (-1.125,-0.783) (-1.175,-0.816) (-1.225,-0.848) (-1.275,-0.881) (-1.325,-0.913) (-1.375,-0.944) (-1.424,-0.975) (-1.475,-1.006) (-1.525,-1.037) (-1.575,-1.066) (-1.625,-1.096) (-1.675,-1.124) (-1.721,-1.150) (-1.766,-1.175) (-1.812,-1.200) (-1.859,-1.225) (-1.907,-1.250) (-1.957,-1.275) (-2.008,-1.300) (-2.060,-1.325) (-2.114,-1.350) (-2.171,-1.375) (-2.225,-1.398) (-2.275,-1.419) (-2.325,-1.438) (-2.375,-1.457) (-2.425,-1.475) (-2.475,-1.492) (-2.525,-1.508) (-2.575,-1.523) (-2.625,-1.538) (-2.675,-1.551) (-2.725,-1.563) (-2.775,-1.574) (-2.825,-1.585) (-2.875,-1.594) (-2.925,-1.602) (-2.975,-1.609) (-3.025,-1.615) (-3.075,-1.620) (-3.125,-1.623) (-3.175,-1.626) (-3.225,-1.627) (-3.275,-1.627) (-3.325,-1.626) (-3.375,-1.624) (-3.425,-1.620) (-3.475,-1.615) (-3.525,-1.609) (-3.575,-1.601) (-3.625,-1.592) (-3.675,-1.581) (-3.725,-1.569) (-3.775,-1.555) (-3.825,-1.539) (-3.875,-1.522) (-3.925,-1.503) (-3.975,-1.482) (-4.025,-1.459) (-4.075,-1.434) (-4.125,-1.407) (-4.175,-1.378) (-4.218,-1.350) (-4.275,-1.311) (-4.323,-1.275) (-4.375,-1.232) (-4.425,-1.188) (-4.464,-1.150) (-4.500,-1.113) (-4.535,-1.075) (-4.575,-1.028) (-4.616,-0.975) (-4.650,-0.928) (-4.685,-0.875) (-4.716,-0.825) (-4.744,-0.775) (-4.770,-0.725) (-4.793,-0.675) (-4.815,-0.625) (-4.835,-0.575) (-4.853,-0.525) (-4.869,-0.475) (-4.883,-0.425) (-4.895,-0.375) (-4.906,-0.325) (-4.916,-0.275) (-4.923,-0.225) (-4.929,-0.175) (-4.934,-0.125) (-4.937,-0.075) (-4.938,-0.025) (-4.938,0.050) (-4.936,0.100) (-4.932,0.150) (-4.926,0.200) (-4.920,0.250) (-4.911,0.300) (-4.901,0.350) (-4.889,0.400) (-4.876,0.450) (-4.861,0.500) (-4.844,0.550) (-4.825,0.600) (-4.804,0.650) (-4.782,0.700) (-4.757,0.750) (-4.730,0.800) (-4.701,0.850) (-4.669,0.900) (-4.635,0.950) (-4.600,0.996) (-4.556,1.050) (-4.512,1.100) (-4.475,1.139) (-4.438,1.175) (-4.400,1.210) (-4.354,1.250) (-4.300,1.292) (-4.255,1.325) (-4.200,1.362) (-4.150,1.393) (-4.100,1.421) (-4.050,1.447) (-4.000,1.471) (-3.950,1.493) (-3.900,1.513) (-3.850,1.531) (-3.800,1.547) (-3.750,1.562) (-3.700,1.575) (-3.650,1.586) (-3.600,1.596) (-3.550,1.605) (-3.500,1.612) (-3.450,1.618) (-3.400,1.622) (-3.350,1.625) (-3.300,1.627) (-3.250,1.627) (-3.200,1.627) (-3.150,1.625) (-3.100,1.622) (-3.050,1.617) (-3.000,1.612) (-2.950,1.605) (-2.900,1.598) (-2.850,1.589) (-2.800,1.580) (-2.750,1.569) (-2.700,1.557) (-2.650,1.544) (-2.600,1.531) (-2.550,1.516) (-2.500,1.500) (-2.450,1.484) (-2.400,1.466) (-2.350,1.448) (-2.300,1.429) (-2.250,1.409) (-2.200,1.388) (-2.150,1.366) (-2.100,1.344) (-2.050,1.320) (-2.000,1.296) (-1.950,1.272) (-1.900,1.246) (-1.850,1.220) (-1.800,1.194) (-1.750,1.167) (-1.700,1.139) (-1.650,1.110) (-1.600,1.081) (-1.550,1.051) (-1.506,1.025) (-1.450,0.991) (-1.400,0.960) (-1.350,0.929) (-1.305,0.900) (-1.250,0.865) (-1.200,0.832) (-1.151,0.800) (-1.100,0.766) (-1.050,0.733) (-1.001,0.700) (-0.950,0.665) (-0.900,0.631) (-0.855,0.600) (-0.800,0.563) (-0.750,0.528) (-0.700,0.493) (-0.650,0.459) (-0.603,0.425) (-0.550,0.389) (-0.500,0.353) (-0.450,0.318) (-0.400,0.283) (-0.354,0.250) (-0.300,0.213) (-0.250,0.177) (-0.220,0.155)}
\EQ{(0.206,0.144) (0.248,0.175) (0.300,0.213) (0.350,0.248) (0.400,0.283) (0.450,0.318) (0.496,0.350) (0.550,0.389) (0.600,0.423) (0.650,0.459) (0.700,0.493) (0.746,0.525) (0.800,0.563) (0.850,0.597) (0.891,0.625) (0.950,0.665) (1.000,0.699) (1.050,0.733) (1.100,0.766) (1.150,0.799) (1.200,0.832) (1.250,0.865) (1.300,0.897) (1.345,0.925) (1.400,0.960) (1.450,0.991) (1.500,1.021) (1.548,1.050) (1.600,1.081) (1.650,1.110) (1.700,1.139) (1.750,1.167) (1.800,1.194) (1.850,1.220) (1.900,1.246) (1.950,1.272) (2.000,1.296) (2.050,1.320) (2.100,1.344) (2.150,1.366) (2.200,1.388) (2.250,1.409) (2.300,1.429) (2.350,1.448) (2.400,1.466) (2.450,1.484) (2.500,1.500) (2.550,1.516) (2.600,1.531) (2.650,1.544) (2.700,1.557) (2.750,1.569) (2.800,1.580) (2.850,1.589) (2.900,1.598) (2.950,1.605) (3.000,1.612) (3.050,1.617) (3.100,1.622) (3.150,1.625) (3.200,1.627) (3.250,1.627) (3.300,1.627) (3.350,1.625) (3.400,1.622) (3.450,1.618) (3.500,1.612) (3.550,1.605) (3.600,1.596) (3.650,1.586) (3.700,1.575) (3.750,1.562) (3.800,1.547) (3.850,1.531) (3.900,1.513) (3.950,1.493) (4.000,1.471) (4.050,1.447) (4.100,1.421) (4.150,1.393) (4.200,1.362) (4.250,1.328) (4.300,1.292) (4.350,1.253) (4.400,1.210) (4.438,1.175) (4.475,1.139) (4.512,1.100) (4.550,1.058) (4.597,1.000) (4.635,0.950) (4.669,0.900) (4.700,0.851) (4.730,0.800) (4.757,0.750) (4.782,0.700) (4.804,0.650) (4.825,0.600) (4.844,0.550) (4.861,0.500) (4.876,0.450) (4.889,0.400) (4.901,0.350) (4.911,0.300) (4.920,0.250) (4.926,0.200) (4.932,0.150) (4.936,0.100) (4.938,0.050) (4.939,0.000) (4.938,-0.050) (4.936,-0.100) (4.932,-0.150) (4.926,-0.200) (4.920,-0.250) (4.911,-0.300) (4.901,-0.350) (4.889,-0.400) (4.876,-0.450) (4.861,-0.500) (4.844,-0.550) (4.825,-0.600) (4.804,-0.650) (4.782,-0.700) (4.757,-0.750) (4.730,-0.800) (4.701,-0.850) (4.669,-0.900) (4.635,-0.950) (4.600,-0.996) (4.556,-1.050) (4.512,-1.100) (4.475,-1.139) (4.438,-1.175) (4.400,-1.210) (4.354,-1.250) (4.300,-1.292) (4.255,-1.325) (4.200,-1.362) (4.150,-1.393) (4.100,-1.421) (4.050,-1.447) (4.000,-1.471) (3.950,-1.493) (3.900,-1.513) (3.850,-1.531) (3.800,-1.547) (3.750,-1.562) (3.700,-1.575) (3.650,-1.586) (3.600,-1.596) (3.550,-1.605) (3.500,-1.612) (3.450,-1.618) (3.400,-1.622) (3.350,-1.625) (3.300,-1.627) (3.250,-1.627) (3.200,-1.627) (3.150,-1.625) (3.100,-1.622) (3.050,-1.617) (3.000,-1.612) (2.950,-1.605) (2.900,-1.598) (2.850,-1.589) (2.800,-1.580) (2.750,-1.569) (2.700,-1.557) (2.650,-1.544) (2.600,-1.531) (2.550,-1.516) (2.500,-1.500) (2.450,-1.484) (2.400,-1.466) (2.350,-1.448) (2.300,-1.429) (2.250,-1.409) (2.200,-1.388) (2.150,-1.366) (2.100,-1.344) (2.050,-1.320) (2.000,-1.296) (1.950,-1.272) (1.900,-1.246) (1.850,-1.220) (1.800,-1.194) (1.750,-1.167) (1.700,-1.139) (1.650,-1.110) (1.600,-1.081) (1.550,-1.051) (1.506,-1.025) (1.450,-0.991) (1.400,-0.960) (1.350,-0.929) (1.305,-0.900) (1.250,-0.865) (1.200,-0.832) (1.151,-0.800) (1.100,-0.766) (1.050,-0.733) (1.001,-0.700) (0.950,-0.665) (0.900,-0.631) (0.855,-0.600) (0.800,-0.563) (0.750,-0.528) (0.700,-0.493) (0.650,-0.459) (0.603,-0.425) (0.550,-0.389) (0.500,-0.353) (0.450,-0.318) (0.400,-0.283) (0.354,-0.250) (0.300,-0.213) (0.250,-0.177) (0.206,-0.144)}
\EQ{(-0.220,-0.155) (-0.275,-0.195) (-0.319,-0.225) (-0.375,-0.266) (-0.425,-0.300) (-0.475,-0.336) (-0.525,-0.371) (-0.567,-0.400) (-0.625,-0.441) (-0.674,-0.475) (-0.725,-0.511) (-0.775,-0.546) (-0.818,-0.575) (-0.875,-0.615) (-0.925,-0.649) (-0.975,-0.683) (-1.025,-0.717) (-1.074,-0.750) (-1.125,-0.785) (-1.175,-0.818) (-1.224,-0.850) (-1.275,-0.884) (-1.325,-0.916) (-1.375,-0.948) (-1.417,-0.975) (-1.475,-1.011) (-1.525,-1.042) (-1.575,-1.073) (-1.620,-1.100) (-1.675,-1.133) (-1.725,-1.162) (-1.775,-1.190) (-1.825,-1.218) (-1.875,-1.246) (-1.925,-1.273) (-1.975,-1.299) (-2.025,-1.325) (-2.075,-1.350) (-2.125,-1.374) (-2.175,-1.398) (-2.225,-1.421) (-2.275,-1.443) (-2.325,-1.465) (-2.375,-1.486) (-2.425,-1.505) (-2.475,-1.525) (-2.525,-1.543) (-2.575,-1.560) (-2.625,-1.577) (-2.675,-1.593) (-2.725,-1.608) (-2.775,-1.622) (-2.825,-1.635) (-2.875,-1.647) (-2.925,-1.658) (-2.975,-1.668) (-3.025,-1.677) (-3.075,-1.686) (-3.125,-1.693) (-3.175,-1.699) (-3.225,-1.704) (-3.275,-1.708) (-3.325,-1.711) (-3.375,-1.713) (-3.425,-1.713) (-3.475,-1.713) (-3.525,-1.711) (-3.575,-1.708) (-3.625,-1.704) (-3.675,-1.698) (-3.725,-1.691) (-3.775,-1.683) (-3.825,-1.674) (-3.875,-1.663) (-3.925,-1.650) (-3.975,-1.636) (-4.025,-1.621) (-4.075,-1.604) (-4.125,-1.585) (-4.175,-1.564) (-4.225,-1.542) (-4.275,-1.518) (-4.325,-1.491) (-4.375,-1.462) (-4.425,-1.431) (-4.472,-1.400) (-4.525,-1.362) (-4.572,-1.325) (-4.625,-1.280) (-4.675,-1.234) (-4.725,-1.185) (-4.775,-1.130) (-4.822,-1.075) (-4.861,-1.025) (-4.897,-0.975) (-4.925,-0.932) (-4.961,-0.875) (-4.989,-0.825) (-5.015,-0.775) (-5.040,-0.725) (-5.062,-0.675) (-5.082,-0.625) (-5.100,-0.577) (-5.117,-0.525) (-5.133,-0.475) (-5.146,-0.425) (-5.158,-0.375) (-5.168,-0.325) (-5.177,-0.275) (-5.184,-0.225) (-5.190,-0.175) (-5.195,-0.125) (-5.197,-0.075) (-5.199,-0.025) (-5.198,0.050) (-5.196,0.100) (-5.193,0.150) (-5.187,0.200) (-5.181,0.250) (-5.173,0.300) (-5.163,0.350) (-5.152,0.400) (-5.140,0.450) (-5.125,0.500) (-5.109,0.550) (-5.092,0.600) (-5.072,0.650) (-5.051,0.700) (-5.028,0.750) (-5.002,0.800) (-4.975,0.850) (-4.946,0.900) (-4.914,0.950) (-4.879,1.000) (-4.842,1.050) (-4.801,1.100) (-4.758,1.150) (-4.710,1.200) (-4.659,1.250) (-4.602,1.300) (-4.550,1.343) (-4.507,1.375) (-4.450,1.415) (-4.400,1.447) (-4.353,1.475) (-4.309,1.500) (-4.260,1.525) (-4.208,1.550) (-4.150,1.575) (-4.100,1.595) (-4.050,1.613) (-4.000,1.629) (-3.950,1.644) (-3.900,1.657) (-3.850,1.668) (-3.800,1.679) (-3.750,1.687) (-3.700,1.695) (-3.650,1.701) (-3.600,1.706) (-3.550,1.710) (-3.500,1.712) (-3.450,1.713) (-3.400,1.713) (-3.350,1.712) (-3.300,1.710) (-3.250,1.706) (-3.200,1.702) (-3.150,1.696) (-3.100,1.689) (-3.050,1.682) (-3.000,1.673) (-2.950,1.663) (-2.900,1.653) (-2.850,1.641) (-2.800,1.628) (-2.750,1.615) (-2.700,1.600) (-2.650,1.585) (-2.600,1.569) (-2.550,1.552) (-2.500,1.534) (-2.450,1.515) (-2.400,1.496) (-2.350,1.475) (-2.300,1.454) (-2.250,1.432) (-2.200,1.410) (-2.150,1.386) (-2.100,1.362) (-2.050,1.338) (-2.000,1.312) (-1.950,1.286) (-1.900,1.259) (-1.850,1.232) (-1.800,1.204) (-1.750,1.176) (-1.705,1.150) (-1.650,1.118) (-1.600,1.088) (-1.550,1.058) (-1.500,1.027) (-1.457,1.000) (-1.400,0.964) (-1.350,0.932) (-1.300,0.900) (-1.250,0.867) (-1.200,0.835) (-1.150,0.801) (-1.100,0.768) (-1.050,0.734) (-1.000,0.700) (-0.950,0.666) (-0.900,0.632) (-0.854,0.600) (-0.800,0.563) (-0.750,0.528) (-0.700,0.494) (-0.650,0.459) (-0.603,0.425) (-0.550,0.389) (-0.500,0.353) (-0.450,0.318) (-0.400,0.283) (-0.354,0.250) (-0.300,0.213) (-0.250,0.177) (-0.220,0.155)}
\EQ{(0.206,0.144) (0.248,0.175) (0.300,0.213) (0.350,0.248) (0.400,0.283) (0.450,0.318) (0.496,0.350) (0.550,0.389) (0.600,0.423) (0.650,0.459) (0.700,0.494) (0.746,0.525) (0.800,0.563) (0.850,0.598) (0.900,0.632) (0.950,0.666) (1.000,0.700) (1.050,0.734) (1.100,0.768) (1.148,0.800) (1.200,0.835) (1.250,0.867) (1.300,0.900) (1.350,0.932) (1.400,0.964) (1.450,0.996) (1.497,1.025) (1.550,1.058) (1.600,1.088) (1.650,1.118) (1.700,1.147) (1.748,1.175) (1.792,1.200) (1.837,1.225) (1.883,1.250) (1.929,1.275) (1.975,1.299) (2.025,1.325) (2.075,1.350) (2.125,1.374) (2.175,1.398) (2.225,1.421) (2.275,1.443) (2.325,1.465) (2.375,1.486) (2.425,1.505) (2.475,1.525) (2.525,1.543) (2.575,1.560) (2.625,1.577) (2.675,1.593) (2.725,1.608) (2.775,1.622) (2.825,1.635) (2.875,1.647) (2.925,1.658) (2.975,1.668) (3.025,1.677) (3.075,1.686) (3.125,1.693) (3.175,1.699) (3.225,1.704) (3.275,1.708) (3.325,1.711) (3.375,1.713) (3.425,1.713) (3.475,1.713) (3.525,1.711) (3.575,1.708) (3.625,1.704) (3.675,1.698) (3.725,1.691) (3.775,1.683) (3.825,1.674) (3.875,1.663) (3.925,1.650) (3.975,1.636) (4.025,1.621) (4.075,1.604) (4.125,1.585) (4.175,1.564) (4.225,1.542) (4.275,1.518) (4.325,1.491) (4.375,1.462) (4.425,1.431) (4.472,1.400) (4.525,1.362) (4.572,1.325) (4.625,1.280) (4.675,1.234) (4.725,1.185) (4.775,1.130) (4.822,1.075) (4.861,1.025) (4.897,0.975) (4.925,0.932) (4.961,0.875) (4.989,0.825) (5.015,0.775) (5.040,0.725) (5.062,0.675) (5.082,0.625) (5.100,0.577) (5.117,0.525) (5.133,0.475) (5.146,0.425) (5.158,0.375) (5.168,0.325) (5.177,0.275) (5.184,0.225) (5.190,0.175) (5.195,0.125) (5.197,0.075) (5.199,0.025) (5.198,-0.050) (5.196,-0.100) (5.193,-0.150) (5.187,-0.200) (5.181,-0.250) (5.173,-0.300) (5.163,-0.350) (5.152,-0.400) (5.140,-0.450) (5.125,-0.500) (5.109,-0.550) (5.092,-0.600) (5.072,-0.650) (5.051,-0.700) (5.028,-0.750) (5.002,-0.800) (4.975,-0.850) (4.946,-0.900) (4.914,-0.950) (4.879,-1.000) (4.842,-1.050) (4.801,-1.100) (4.758,-1.150) (4.710,-1.200) (4.659,-1.250) (4.602,-1.300) (4.550,-1.343) (4.507,-1.375) (4.450,-1.415) (4.400,-1.447) (4.353,-1.475) (4.309,-1.500) (4.260,-1.525) (4.208,-1.550) (4.150,-1.575) (4.100,-1.595) (4.050,-1.613) (4.000,-1.629) (3.950,-1.644) (3.900,-1.657) (3.850,-1.668) (3.800,-1.679) (3.750,-1.687) (3.700,-1.695) (3.650,-1.701) (3.600,-1.706) (3.550,-1.710) (3.500,-1.712) (3.450,-1.713) (3.400,-1.713) (3.350,-1.712) (3.300,-1.710) (3.250,-1.706) (3.200,-1.702) (3.150,-1.696) (3.100,-1.689) (3.050,-1.682) (3.000,-1.673) (2.950,-1.663) (2.900,-1.653) (2.850,-1.641) (2.800,-1.628) (2.750,-1.615) (2.700,-1.600) (2.650,-1.585) (2.600,-1.569) (2.550,-1.552) (2.500,-1.534) (2.450,-1.515) (2.400,-1.496) (2.350,-1.475) (2.300,-1.454) (2.250,-1.432) (2.200,-1.410) (2.150,-1.386) (2.100,-1.362) (2.050,-1.338) (2.000,-1.312) (1.950,-1.286) (1.900,-1.259) (1.850,-1.232) (1.800,-1.204) (1.750,-1.176) (1.705,-1.150) (1.650,-1.118) (1.600,-1.088) (1.550,-1.058) (1.500,-1.027) (1.457,-1.000) (1.400,-0.964) (1.350,-0.932) (1.300,-0.900) (1.250,-0.867) (1.200,-0.835) (1.150,-0.801) (1.100,-0.768) (1.050,-0.734) (1.000,-0.700) (0.950,-0.666) (0.900,-0.632) (0.854,-0.600) (0.800,-0.563) (0.750,-0.528) (0.700,-0.494) (0.650,-0.459) (0.603,-0.425) (0.550,-0.389) (0.500,-0.353) (0.450,-0.318) (0.400,-0.283) (0.354,-0.250) (0.300,-0.213) (0.250,-0.177) (0.206,-0.144)}
\EQ{(-0.220,-0.155) (-0.275,-0.195) (-0.319,-0.225) (-0.375,-0.266) (-0.425,-0.300) (-0.475,-0.336) (-0.525,-0.371) (-0.567,-0.400) (-0.625,-0.441) (-0.674,-0.475) (-0.725,-0.511) (-0.775,-0.546) (-0.817,-0.575) (-0.875,-0.616) (-0.925,-0.650) (-0.975,-0.684) (-1.025,-0.719) (-1.072,-0.750) (-1.125,-0.787) (-1.175,-0.820) (-1.220,-0.850) (-1.275,-0.887) (-1.325,-0.920) (-1.372,-0.950) (-1.425,-0.985) (-1.475,-1.017) (-1.525,-1.048) (-1.568,-1.075) (-1.625,-1.110) (-1.675,-1.141) (-1.725,-1.171) (-1.774,-1.200) (-1.825,-1.230) (-1.875,-1.259) (-1.925,-1.287) (-1.975,-1.314) (-2.025,-1.342) (-2.075,-1.368) (-2.125,-1.394) (-2.175,-1.420) (-2.225,-1.444) (-2.275,-1.468) (-2.325,-1.492) (-2.375,-1.515) (-2.425,-1.537) (-2.475,-1.558) (-2.525,-1.579) (-2.575,-1.599) (-2.625,-1.618) (-2.675,-1.637) (-2.725,-1.654) (-2.775,-1.671) (-2.825,-1.687) (-2.875,-1.702) (-2.925,-1.717) (-2.975,-1.730) (-3.025,-1.743) (-3.075,-1.755) (-3.125,-1.765) (-3.175,-1.775) (-3.225,-1.784) (-3.275,-1.792) (-3.325,-1.799) (-3.375,-1.805) (-3.425,-1.811) (-3.475,-1.815) (-3.525,-1.818) (-3.575,-1.820) (-3.625,-1.820) (-3.675,-1.820) (-3.725,-1.819) (-3.775,-1.816) (-3.825,-1.813) (-3.875,-1.808) (-3.925,-1.802) (-3.975,-1.795) (-4.025,-1.786) (-4.075,-1.776) (-4.125,-1.765) (-4.175,-1.753) (-4.225,-1.739) (-4.275,-1.723) (-4.325,-1.706) (-4.375,-1.688) (-4.425,-1.667) (-4.475,-1.645) (-4.525,-1.622) (-4.575,-1.596) (-4.625,-1.568) (-4.675,-1.538) (-4.725,-1.506) (-4.770,-1.475) (-4.825,-1.435) (-4.868,-1.400) (-4.925,-1.352) (-4.975,-1.305) (-5.025,-1.255) (-5.075,-1.200) (-5.118,-1.150) (-5.150,-1.109) (-5.194,-1.050) (-5.225,-1.004) (-5.259,-0.950) (-5.289,-0.900) (-5.316,-0.850) (-5.341,-0.800) (-5.365,-0.750) (-5.386,-0.700) (-5.406,-0.650) (-5.424,-0.600) (-5.441,-0.550) (-5.456,-0.500) (-5.469,-0.450) (-5.481,-0.400) (-5.491,-0.350) (-5.500,-0.300) (-5.508,-0.250) (-5.514,-0.200) (-5.519,-0.150) (-5.522,-0.100) (-5.524,-0.050) (-5.525,0.000) (-5.524,0.050) (-5.522,0.100) (-5.519,0.150) (-5.514,0.200) (-5.508,0.250) (-5.500,0.300) (-5.491,0.350) (-5.481,0.400) (-5.469,0.450) (-5.456,0.500) (-5.441,0.550) (-5.424,0.600) (-5.406,0.650) (-5.386,0.700) (-5.365,0.750) (-5.341,0.800) (-5.316,0.850) (-5.289,0.900) (-5.259,0.950) (-5.228,1.000) (-5.194,1.050) (-5.157,1.100) (-5.125,1.141) (-5.075,1.200) (-5.030,1.250) (-4.980,1.300) (-4.927,1.350) (-4.875,1.395) (-4.825,1.435) (-4.775,1.472) (-4.725,1.506) (-4.675,1.538) (-4.625,1.568) (-4.575,1.596) (-4.525,1.622) (-4.475,1.645) (-4.425,1.667) (-4.375,1.688) (-4.325,1.706) (-4.275,1.723) (-4.225,1.739) (-4.175,1.753) (-4.125,1.765) (-4.075,1.776) (-4.025,1.786) (-3.975,1.795) (-3.925,1.802) (-3.875,1.808) (-3.825,1.813) (-3.775,1.816) (-3.725,1.819) (-3.675,1.820) (-3.625,1.820) (-3.575,1.820) (-3.525,1.818) (-3.475,1.815) (-3.425,1.811) (-3.375,1.805) (-3.325,1.799) (-3.275,1.792) (-3.225,1.784) (-3.175,1.775) (-3.125,1.765) (-3.075,1.755) (-3.025,1.743) (-2.975,1.730) (-2.925,1.717) (-2.875,1.702) (-2.825,1.687) (-2.775,1.671) (-2.725,1.654) (-2.675,1.637) (-2.625,1.618) (-2.578,1.600) (-2.525,1.579) (-2.475,1.558) (-2.425,1.537) (-2.375,1.515) (-2.325,1.492) (-2.275,1.468) (-2.225,1.444) (-2.175,1.420) (-2.125,1.394) (-2.075,1.368) (-2.025,1.342) (-1.975,1.314) (-1.925,1.287) (-1.875,1.259) (-1.825,1.230) (-1.775,1.201) (-1.725,1.171) (-1.675,1.141) (-1.625,1.110) (-1.575,1.079) (-1.528,1.050) (-1.475,1.017) (-1.425,0.985) (-1.375,0.952) (-1.325,0.920) (-1.275,0.887) (-1.225,0.854) (-1.183,0.825) (-1.125,0.787) (-1.075,0.753) (-1.025,0.719) (-0.975,0.684) (-0.925,0.650) (-0.875,0.616) (-0.825,0.581) (-0.781,0.550) (-0.725,0.511) (-0.675,0.476) (-0.625,0.441) (-0.575,0.406) (-0.532,0.375) (-0.475,0.336) (-0.425,0.300) (-0.375,0.266) (-0.325,0.230) (-0.284,0.200) (-0.225,0.160) (-0.220,0.155)}
\EQ{(0.206,0.144) (0.248,0.175) (0.300,0.213) (0.350,0.248) (0.400,0.283) (0.450,0.318) (0.496,0.350) (0.550,0.389) (0.600,0.424) (0.650,0.459) (0.700,0.494) (0.745,0.525) (0.800,0.564) (0.850,0.598) (0.900,0.633) (0.950,0.667) (0.998,0.700) (1.050,0.736) (1.100,0.770) (1.145,0.800) (1.200,0.837) (1.250,0.870) (1.296,0.900) (1.350,0.936) (1.400,0.968) (1.449,1.000) (1.500,1.032) (1.550,1.064) (1.600,1.095) (1.649,1.125) (1.700,1.156) (1.750,1.186) (1.800,1.215) (1.850,1.244) (1.900,1.273) (1.949,1.300) (1.995,1.325) (2.041,1.350) (2.088,1.375) (2.137,1.400) (2.186,1.425) (2.237,1.450) (2.289,1.475) (2.343,1.500) (2.398,1.525) (2.450,1.548) (2.500,1.569) (2.550,1.589) (2.600,1.609) (2.650,1.628) (2.700,1.646) (2.750,1.663) (2.800,1.679) (2.850,1.695) (2.900,1.710) (2.950,1.724) (3.000,1.737) (3.050,1.749) (3.100,1.760) (3.150,1.771) (3.200,1.780) (3.250,1.788) (3.300,1.796) (3.350,1.803) (3.400,1.808) (3.450,1.813) (3.500,1.816) (3.550,1.819) (3.600,1.820) (3.650,1.820) (3.700,1.820) (3.750,1.818) (3.800,1.815) (3.850,1.811) (3.900,1.805) (3.950,1.799) (4.000,1.791) (4.050,1.782) (4.100,1.771) (4.150,1.759) (4.200,1.746) (4.250,1.731) (4.300,1.715) (4.350,1.697) (4.400,1.678) (4.450,1.657) (4.500,1.634) (4.550,1.609) (4.600,1.582) (4.650,1.553) (4.696,1.525) (4.750,1.489) (4.800,1.453) (4.850,1.415) (4.898,1.375) (4.950,1.329) (5.000,1.280) (5.050,1.228) (5.097,1.175) (5.138,1.125) (5.175,1.076) (5.211,1.025) (5.244,0.975) (5.274,0.925) (5.300,0.880) (5.325,0.832) (5.350,0.782) (5.375,0.726) (5.396,0.675) (5.415,0.625) (5.433,0.575) (5.448,0.525) (5.462,0.475) (5.475,0.425) (5.486,0.375) (5.496,0.325) (5.504,0.275) (5.511,0.225) (5.516,0.175) (5.521,0.125) (5.523,0.075) (5.525,0.025) (5.524,-0.050) (5.522,-0.100) (5.519,-0.150) (5.514,-0.200) (5.508,-0.250) (5.500,-0.300) (5.491,-0.350) (5.481,-0.400) (5.469,-0.450) (5.456,-0.500) (5.441,-0.550) (5.424,-0.600) (5.406,-0.650) (5.386,-0.700) (5.365,-0.750) (5.341,-0.800) (5.316,-0.850) (5.289,-0.900) (5.259,-0.950) (5.228,-1.000) (5.194,-1.050) (5.157,-1.100) (5.125,-1.141) (5.075,-1.200) (5.030,-1.250) (4.980,-1.300) (4.927,-1.350) (4.875,-1.395) (4.825,-1.435) (4.775,-1.472) (4.725,-1.506) (4.675,-1.538) (4.625,-1.568) (4.575,-1.596) (4.525,-1.622) (4.475,-1.645) (4.425,-1.667) (4.375,-1.688) (4.325,-1.706) (4.275,-1.723) (4.225,-1.739) (4.175,-1.753) (4.125,-1.765) (4.075,-1.776) (4.025,-1.786) (3.975,-1.795) (3.925,-1.802) (3.875,-1.808) (3.825,-1.813) (3.775,-1.816) (3.725,-1.819) (3.675,-1.820) (3.625,-1.820) (3.575,-1.820) (3.525,-1.818) (3.475,-1.815) (3.425,-1.811) (3.375,-1.805) (3.325,-1.799) (3.275,-1.792) (3.225,-1.784) (3.175,-1.775) (3.125,-1.765) (3.075,-1.755) (3.025,-1.743) (2.975,-1.730) (2.925,-1.717) (2.875,-1.702) (2.825,-1.687) (2.775,-1.671) (2.725,-1.654) (2.675,-1.637) (2.625,-1.618) (2.578,-1.600) (2.525,-1.579) (2.475,-1.558) (2.425,-1.537) (2.375,-1.515) (2.325,-1.492) (2.275,-1.468) (2.225,-1.444) (2.175,-1.420) (2.125,-1.394) (2.075,-1.368) (2.025,-1.342) (1.975,-1.314) (1.925,-1.287) (1.875,-1.259) (1.825,-1.230) (1.775,-1.201) (1.725,-1.171) (1.675,-1.141) (1.625,-1.110) (1.575,-1.079) (1.528,-1.050) (1.475,-1.017) (1.425,-0.985) (1.375,-0.952) (1.325,-0.920) (1.275,-0.887) (1.225,-0.854) (1.183,-0.825) (1.125,-0.787) (1.075,-0.753) (1.025,-0.719) (0.975,-0.684) (0.925,-0.650) (0.875,-0.616) (0.825,-0.581) (0.781,-0.550) (0.725,-0.511) (0.675,-0.476) (0.625,-0.441) (0.575,-0.406) (0.532,-0.375) (0.475,-0.336) (0.425,-0.300) (0.375,-0.266) (0.325,-0.230) (0.284,-0.200) (0.225,-0.160) (0.206,-0.144)}
\EQ{(-0.220,-0.155) (-0.275,-0.195) (-0.319,-0.225) (-0.375,-0.266) (-0.425,-0.300) (-0.475,-0.336) (-0.525,-0.371) (-0.567,-0.400) (-0.625,-0.442) (-0.673,-0.475) (-0.725,-0.512) (-0.775,-0.547) (-0.817,-0.575) (-0.875,-0.616) (-0.924,-0.650) (-0.975,-0.686) (-1.025,-0.720) (-1.069,-0.750) (-1.125,-0.788) (-1.175,-0.822) (-1.225,-0.856) (-1.275,-0.890) (-1.325,-0.923) (-1.375,-0.956) (-1.425,-0.989) (-1.475,-1.022) (-1.519,-1.050) (-1.575,-1.086) (-1.625,-1.118) (-1.675,-1.149) (-1.725,-1.180) (-1.775,-1.211) (-1.825,-1.242) (-1.875,-1.271) (-1.924,-1.300) (-1.975,-1.330) (-2.025,-1.359) (-2.075,-1.387) (-2.125,-1.415) (-2.175,-1.442) (-2.225,-1.468) (-2.275,-1.495) (-2.325,-1.520) (-2.375,-1.545) (-2.425,-1.570) (-2.475,-1.593) (-2.525,-1.617) (-2.575,-1.639) (-2.625,-1.661) (-2.675,-1.682) (-2.725,-1.703) (-2.775,-1.723) (-2.825,-1.742) (-2.875,-1.761) (-2.925,-1.778) (-2.975,-1.796) (-3.025,-1.812) (-3.075,-1.827) (-3.125,-1.842) (-3.175,-1.856) (-3.225,-1.869) (-3.275,-1.881) (-3.325,-1.893) (-3.375,-1.904) (-3.425,-1.913) (-3.475,-1.922) (-3.525,-1.930) (-3.575,-1.937) (-3.625,-1.944) (-3.675,-1.949) (-3.725,-1.953) (-3.775,-1.957) (-3.825,-1.959) (-3.875,-1.961) (-3.925,-1.961) (-3.975,-1.960) (-4.025,-1.959) (-4.075,-1.956) (-4.125,-1.952) (-4.175,-1.948) (-4.225,-1.942) (-4.275,-1.934) (-4.325,-1.926) (-4.375,-1.917) (-4.425,-1.906) (-4.475,-1.894) (-4.525,-1.880) (-4.575,-1.866) (-4.624,-1.850) (-4.675,-1.832) (-4.725,-1.813) (-4.775,-1.793) (-4.825,-1.770) (-4.875,-1.746) (-4.925,-1.721) (-4.975,-1.693) (-5.025,-1.664) (-5.075,-1.632) (-5.123,-1.600) (-5.175,-1.562) (-5.223,-1.525) (-5.275,-1.482) (-5.325,-1.438) (-5.365,-1.400) (-5.415,-1.350) (-5.450,-1.312) (-5.500,-1.254) (-5.544,-1.200) (-5.575,-1.158) (-5.616,-1.100) (-5.649,-1.050) (-5.675,-1.007) (-5.708,-0.950) (-5.735,-0.900) (-5.759,-0.850) (-5.782,-0.800) (-5.804,-0.750) (-5.823,-0.700) (-5.842,-0.650) (-5.858,-0.600) (-5.874,-0.550) (-5.887,-0.500) (-5.900,-0.450) (-5.911,-0.400) (-5.920,-0.350) (-5.929,-0.300) (-5.936,-0.250) (-5.941,-0.200) (-5.946,-0.150) (-5.949,-0.100) (-5.951,-0.050) (-5.951,0.000) (-5.951,0.050) (-5.949,0.100) (-5.946,0.150) (-5.941,0.200) (-5.936,0.250) (-5.929,0.300) (-5.920,0.350) (-5.911,0.400) (-5.900,0.450) (-5.887,0.500) (-5.874,0.550) (-5.858,0.600) (-5.842,0.650) (-5.823,0.700) (-5.804,0.750) (-5.782,0.800) (-5.759,0.850) (-5.735,0.900) (-5.708,0.950) (-5.679,1.000) (-5.650,1.048) (-5.616,1.100) (-5.581,1.150) (-5.550,1.192) (-5.504,1.250) (-5.461,1.300) (-5.425,1.339) (-5.375,1.390) (-5.325,1.438) (-5.284,1.475) (-5.225,1.524) (-5.175,1.562) (-5.125,1.599) (-5.075,1.632) (-5.025,1.664) (-4.975,1.693) (-4.925,1.721) (-4.875,1.746) (-4.825,1.770) (-4.775,1.793) (-4.725,1.813) (-4.675,1.832) (-4.625,1.850) (-4.575,1.866) (-4.525,1.880) (-4.475,1.894) (-4.425,1.906) (-4.375,1.917) (-4.325,1.926) (-4.275,1.934) (-4.225,1.942) (-4.175,1.948) (-4.125,1.952) (-4.075,1.956) (-4.025,1.959) (-3.975,1.960) (-3.925,1.961) (-3.875,1.961) (-3.825,1.959) (-3.775,1.957) (-3.725,1.953) (-3.675,1.949) (-3.625,1.944) (-3.575,1.937) (-3.525,1.930) (-3.475,1.922) (-3.425,1.913) (-3.375,1.904) (-3.325,1.893) (-3.275,1.881) (-3.225,1.869) (-3.175,1.856) (-3.125,1.842) (-3.075,1.827) (-3.025,1.812) (-2.975,1.796) (-2.925,1.778) (-2.875,1.761) (-2.825,1.742) (-2.775,1.723) (-2.725,1.703) (-2.675,1.682) (-2.625,1.661) (-2.575,1.639) (-2.525,1.617) (-2.475,1.593) (-2.425,1.570) (-2.375,1.545) (-2.325,1.520) (-2.275,1.495) (-2.225,1.468) (-2.175,1.442) (-2.125,1.415) (-2.075,1.387) (-2.025,1.359) (-1.975,1.330) (-1.925,1.301) (-1.881,1.275) (-1.825,1.242) (-1.775,1.211) (-1.725,1.180) (-1.676,1.150) (-1.625,1.118) (-1.575,1.086) (-1.525,1.054) (-1.480,1.025) (-1.425,0.989) (-1.375,0.956) (-1.328,0.925) (-1.275,0.890) (-1.225,0.856) (-1.179,0.825) (-1.125,0.788) (-1.075,0.754) (-1.033,0.725) (-0.975,0.686) (-0.925,0.651) (-0.875,0.616) (-0.825,0.581) (-0.781,0.550) (-0.725,0.512) (-0.675,0.476) (-0.625,0.442) (-0.575,0.406) (-0.531,0.375) (-0.475,0.336) (-0.425,0.300) (-0.375,0.266) (-0.325,0.230) (-0.284,0.200) (-0.225,0.160) (-0.220,0.155)}
\EQ{(0.206,0.144) (0.248,0.175) (0.300,0.213) (0.350,0.248) (0.400,0.283) (0.450,0.318) (0.496,0.350) (0.550,0.389) (0.600,0.424) (0.650,0.459) (0.700,0.494) (0.745,0.525) (0.800,0.564) (0.850,0.599) (0.900,0.634) (0.950,0.668) (0.997,0.700) (1.050,0.737) (1.100,0.771) (1.143,0.800) (1.200,0.839) (1.250,0.873) (1.300,0.906) (1.350,0.940) (1.400,0.973) (1.442,1.000) (1.500,1.038) (1.550,1.070) (1.597,1.100) (1.650,1.134) (1.700,1.165) (1.750,1.196) (1.798,1.225) (1.850,1.257) (1.900,1.286) (1.950,1.316) (2.000,1.344) (2.050,1.373) (2.099,1.400) (2.144,1.425) (2.191,1.450) (2.238,1.475) (2.286,1.500) (2.335,1.525) (2.385,1.550) (2.437,1.575) (2.489,1.600) (2.544,1.625) (2.600,1.650) (2.650,1.672) (2.700,1.693) (2.750,1.713) (2.800,1.733) (2.850,1.751) (2.900,1.770) (2.950,1.787) (3.000,1.804) (3.050,1.820) (3.100,1.835) (3.150,1.849) (3.200,1.863) (3.249,1.875) (3.300,1.887) (3.350,1.898) (3.400,1.909) (3.450,1.918) (3.500,1.926) (3.550,1.934) (3.600,1.941) (3.650,1.946) (3.700,1.951) (3.750,1.955) (3.800,1.958) (3.850,1.960) (3.900,1.961) (3.950,1.961) (4.000,1.960) (4.050,1.958) (4.100,1.954) (4.150,1.950) (4.200,1.945) (4.250,1.938) (4.300,1.930) (4.350,1.922) (4.400,1.911) (4.450,1.900) (4.500,1.887) (4.550,1.873) (4.600,1.858) (4.650,1.841) (4.700,1.823) (4.750,1.803) (4.800,1.782) (4.850,1.759) (4.900,1.734) (4.950,1.707) (5.000,1.679) (5.047,1.650) (5.100,1.616) (5.150,1.581) (5.191,1.550) (5.250,1.503) (5.300,1.461) (5.339,1.425) (5.390,1.375) (5.438,1.325) (5.475,1.284) (5.524,1.225) (5.563,1.175) (5.599,1.125) (5.633,1.075) (5.664,1.025) (5.694,0.975) (5.721,0.925) (5.747,0.875) (5.771,0.825) (5.793,0.775) (5.814,0.725) (5.833,0.675) (5.850,0.625) (5.866,0.575) (5.881,0.525) (5.894,0.475) (5.905,0.425) (5.916,0.375) (5.925,0.325) (5.932,0.275) (5.939,0.225) (5.944,0.175) (5.947,0.125) (5.950,0.075) (5.951,0.025) (5.951,-0.050) (5.949,-0.100) (5.946,-0.150) (5.941,-0.200) (5.936,-0.250) (5.929,-0.300) (5.920,-0.350) (5.911,-0.400) (5.900,-0.450) (5.887,-0.500) (5.874,-0.550) (5.858,-0.600) (5.842,-0.650) (5.823,-0.700) (5.804,-0.750) (5.782,-0.800) (5.759,-0.850) (5.735,-0.900) (5.708,-0.950) (5.679,-1.000) (5.650,-1.048) (5.616,-1.100) (5.581,-1.150) (5.550,-1.192) (5.504,-1.250) (5.461,-1.300) (5.425,-1.339) (5.375,-1.390) (5.325,-1.438) (5.284,-1.475) (5.225,-1.524) (5.175,-1.562) (5.125,-1.599) (5.075,-1.632) (5.025,-1.664) (4.975,-1.693) (4.925,-1.721) (4.875,-1.746) (4.825,-1.770) (4.775,-1.793) (4.725,-1.813) (4.675,-1.832) (4.625,-1.850) (4.575,-1.866) (4.525,-1.880) (4.475,-1.894) (4.425,-1.906) (4.375,-1.917) (4.325,-1.926) (4.275,-1.934) (4.225,-1.942) (4.175,-1.948) (4.125,-1.952) (4.075,-1.956) (4.025,-1.959) (3.975,-1.960) (3.925,-1.961) (3.875,-1.961) (3.825,-1.959) (3.775,-1.957) (3.725,-1.953) (3.675,-1.949) (3.625,-1.944) (3.575,-1.937) (3.525,-1.930) (3.475,-1.922) (3.425,-1.913) (3.375,-1.904) (3.325,-1.893) (3.275,-1.881) (3.225,-1.869) (3.175,-1.856) (3.125,-1.842) (3.075,-1.827) (3.025,-1.812) (2.975,-1.796) (2.925,-1.778) (2.875,-1.761) (2.825,-1.742) (2.775,-1.723) (2.725,-1.703) (2.675,-1.682) (2.625,-1.661) (2.575,-1.639) (2.525,-1.617) (2.475,-1.593) (2.425,-1.570) (2.375,-1.545) (2.325,-1.520) (2.275,-1.495) (2.225,-1.468) (2.175,-1.442) (2.125,-1.415) (2.075,-1.387) (2.025,-1.359) (1.975,-1.330) (1.925,-1.301) (1.881,-1.275) (1.825,-1.242) (1.775,-1.211) (1.725,-1.180) (1.676,-1.150) (1.625,-1.118) (1.575,-1.086) (1.525,-1.054) (1.480,-1.025) (1.425,-0.989) (1.375,-0.956) (1.328,-0.925) (1.275,-0.890) (1.225,-0.856) (1.179,-0.825) (1.125,-0.788) (1.075,-0.754) (1.033,-0.725) (0.975,-0.686) (0.925,-0.651) (0.875,-0.616) (0.825,-0.581) (0.781,-0.550) (0.725,-0.512) (0.675,-0.476) (0.625,-0.442) (0.575,-0.406) (0.531,-0.375) (0.475,-0.336) (0.425,-0.300) (0.375,-0.266) (0.325,-0.230) (0.284,-0.200) (0.225,-0.160) (0.206,-0.144)}
\EQ{(-0.220,-0.155) (-0.275,-0.195) (-0.319,-0.225) (-0.375,-0.266) (-0.425,-0.300) (-0.475,-0.336) (-0.525,-0.371) (-0.567,-0.400) (-0.625,-0.442) (-0.673,-0.475) (-0.725,-0.512) (-0.775,-0.547) (-0.825,-0.582) (-0.875,-0.617) (-0.923,-0.650) (-0.975,-0.687) (-1.025,-0.721) (-1.067,-0.750) (-1.125,-0.790) (-1.175,-0.824) (-1.225,-0.859) (-1.275,-0.893) (-1.323,-0.925) (-1.375,-0.960) (-1.425,-0.994) (-1.472,-1.025) (-1.525,-1.060) (-1.575,-1.093) (-1.624,-1.125) (-1.675,-1.158) (-1.725,-1.190) (-1.775,-1.222) (-1.820,-1.250) (-1.875,-1.285) (-1.925,-1.316) (-1.975,-1.346) (-2.023,-1.375) (-2.075,-1.406) (-2.125,-1.436) (-2.175,-1.465) (-2.225,-1.493) (-2.275,-1.521) (-2.325,-1.549) (-2.372,-1.575) (-2.419,-1.600) (-2.466,-1.625) (-2.514,-1.650) (-2.563,-1.675) (-2.613,-1.700) (-2.664,-1.725) (-2.716,-1.750) (-2.770,-1.775) (-2.825,-1.800) (-2.875,-1.822) (-2.925,-1.844) (-2.975,-1.865) (-3.025,-1.885) (-3.075,-1.904) (-3.125,-1.923) (-3.175,-1.942) (-3.225,-1.959) (-3.275,-1.976) (-3.325,-1.993) (-3.375,-2.008) (-3.425,-2.023) (-3.475,-2.037) (-3.525,-2.051) (-3.575,-2.064) (-3.625,-2.076) (-3.675,-2.087) (-3.725,-2.097) (-3.775,-2.107) (-3.825,-2.116) (-3.875,-2.124) (-3.925,-2.131) (-3.975,-2.138) (-4.025,-2.143) (-4.075,-2.148) (-4.125,-2.152) (-4.175,-2.155) (-4.225,-2.157) (-4.275,-2.158) (-4.325,-2.158) (-4.375,-2.158) (-4.425,-2.156) (-4.475,-2.154) (-4.525,-2.150) (-4.575,-2.146) (-4.625,-2.140) (-4.675,-2.134) (-4.725,-2.126) (-4.775,-2.117) (-4.825,-2.108) (-4.875,-2.097) (-4.925,-2.085) (-4.975,-2.071) (-5.025,-2.057) (-5.075,-2.041) (-5.125,-2.024) (-5.175,-2.006) (-5.225,-1.986) (-5.275,-1.965) (-5.325,-1.942) (-5.375,-1.918) (-5.425,-1.892) (-5.475,-1.864) (-5.525,-1.835) (-5.575,-1.804) (-5.618,-1.775) (-5.675,-1.735) (-5.722,-1.700) (-5.775,-1.658) (-5.814,-1.625) (-5.869,-1.575) (-5.921,-1.525) (-5.970,-1.475) (-6.016,-1.425) (-6.050,-1.385) (-6.098,-1.325) (-6.136,-1.275) (-6.172,-1.225) (-6.200,-1.183) (-6.236,-1.125) (-6.266,-1.075) (-6.294,-1.025) (-6.320,-0.975) (-6.344,-0.925) (-6.367,-0.875) (-6.388,-0.825) (-6.408,-0.775) (-6.426,-0.725) (-6.443,-0.675) (-6.459,-0.625) (-6.473,-0.575) (-6.486,-0.525) (-6.498,-0.475) (-6.509,-0.425) (-6.518,-0.375) (-6.526,-0.325) (-6.533,-0.275) (-6.539,-0.225) (-6.543,-0.175) (-6.547,-0.125) (-6.549,-0.075) (-6.550,-0.025) (-6.550,0.037) (-6.548,0.100) (-6.545,0.150) (-6.541,0.200) (-6.536,0.250) (-6.530,0.300) (-6.522,0.350) (-6.513,0.400) (-6.503,0.450) (-6.492,0.500) (-6.480,0.550) (-6.466,0.600) (-6.451,0.650) (-6.435,0.700) (-6.417,0.750) (-6.398,0.800) (-6.378,0.850) (-6.356,0.900) (-6.332,0.950) (-6.307,1.000) (-6.280,1.050) (-6.251,1.100) (-6.225,1.143) (-6.189,1.200) (-6.154,1.250) (-6.118,1.300) (-6.079,1.350) (-6.037,1.400) (-6.000,1.442) (-5.950,1.496) (-5.900,1.546) (-5.850,1.593) (-5.800,1.637) (-5.754,1.675) (-5.700,1.717) (-5.654,1.750) (-5.600,1.787) (-5.550,1.819) (-5.500,1.850) (-5.456,1.875) (-5.409,1.900) (-5.360,1.925) (-5.307,1.950) (-5.251,1.975) (-5.200,1.996) (-5.150,2.015) (-5.100,2.033) (-5.050,2.049) (-5.000,2.064) (-4.950,2.078) (-4.900,2.091) (-4.850,2.102) (-4.800,2.113) (-4.750,2.122) (-4.700,2.130) (-4.650,2.137) (-4.600,2.143) (-4.550,2.148) (-4.500,2.152) (-4.450,2.155) (-4.400,2.157) (-4.350,2.158) (-4.300,2.158) (-4.250,2.158) (-4.200,2.156) (-4.150,2.153) (-4.100,2.150) (-4.050,2.146) (-4.000,2.140) (-3.950,2.134) (-3.900,2.128) (-3.850,2.120) (-3.800,2.111) (-3.750,2.102) (-3.700,2.092) (-3.650,2.081) (-3.600,2.070) (-3.550,2.057) (-3.500,2.044) (-3.450,2.030) (-3.400,2.016) (-3.350,2.001) (-3.300,1.985) (-3.250,1.968) (-3.200,1.951) (-3.150,1.933) (-3.100,1.914) (-3.050,1.895) (-3.001,1.875) (-2.950,1.854) (-2.900,1.833) (-2.850,1.811) (-2.800,1.789) (-2.750,1.766) (-2.700,1.742) (-2.650,1.718) (-2.600,1.694) (-2.550,1.669) (-2.500,1.643) (-2.450,1.617) (-2.400,1.590) (-2.350,1.563) (-2.300,1.535) (-2.250,1.507) (-2.200,1.479) (-2.150,1.450) (-2.100,1.421) (-2.050,1.391) (-2.000,1.361) (-1.950,1.331) (-1.900,1.300) (-1.850,1.269) (-1.800,1.238) (-1.750,1.206) (-1.702,1.175) (-1.650,1.142) (-1.600,1.109) (-1.550,1.076) (-1.500,1.044) (-1.450,1.010) (-1.400,0.977) (-1.350,0.943) (-1.300,0.910) (-1.250,0.876) (-1.200,0.842) (-1.150,0.807) (-1.103,0.775) (-1.050,0.739) (-1.000,0.704) (-0.950,0.669) (-0.900,0.634) (-0.851,0.600) (-0.800,0.565) (-0.750,0.529) (-0.709,0.500) (-0.650,0.459) (-0.602,0.425) (-0.550,0.389) (-0.500,0.354) (-0.450,0.319) (-0.400,0.283) (-0.354,0.250) (-0.300,0.213) (-0.250,0.177) (-0.220,0.155)}
\EQ{(0.206,0.144) (0.248,0.175) (0.300,0.213) (0.350,0.248) (0.400,0.283) (0.450,0.319) (0.496,0.350) (0.550,0.389) (0.600,0.424) (0.650,0.459) (0.700,0.494) (0.744,0.525) (0.800,0.565) (0.850,0.599) (0.900,0.634) (0.950,0.669) (0.995,0.700) (1.050,0.739) (1.100,0.773) (1.150,0.807) (1.200,0.842) (1.249,0.875) (1.300,0.910) (1.350,0.943) (1.397,0.975) (1.450,1.010) (1.500,1.044) (1.548,1.075) (1.600,1.109) (1.650,1.142) (1.700,1.174) (1.750,1.206) (1.800,1.238) (1.850,1.269) (1.900,1.300) (1.950,1.331) (2.000,1.361) (2.050,1.391) (2.100,1.421) (2.150,1.450) (2.193,1.475) (2.237,1.500) (2.282,1.525) (2.327,1.550) (2.372,1.575) (2.419,1.600) (2.466,1.625) (2.514,1.650) (2.563,1.675) (2.613,1.700) (2.664,1.725) (2.716,1.750) (2.770,1.775) (2.825,1.800) (2.875,1.822) (2.925,1.844) (2.975,1.865) (3.025,1.885) (3.075,1.904) (3.125,1.923) (3.175,1.942) (3.225,1.959) (3.275,1.976) (3.325,1.993) (3.375,2.008) (3.425,2.023) (3.475,2.037) (3.525,2.051) (3.575,2.064) (3.625,2.076) (3.675,2.087) (3.725,2.097) (3.775,2.107) (3.825,2.116) (3.875,2.124) (3.925,2.131) (3.975,2.138) (4.025,2.143) (4.075,2.148) (4.125,2.152) (4.175,2.155) (4.225,2.157) (4.275,2.158) (4.325,2.158) (4.375,2.158) (4.425,2.156) (4.475,2.154) (4.525,2.150) (4.575,2.146) (4.625,2.140) (4.675,2.134) (4.725,2.126) (4.775,2.117) (4.825,2.108) (4.875,2.097) (4.925,2.085) (4.975,2.071) (5.025,2.057) (5.075,2.041) (5.125,2.024) (5.175,2.006) (5.225,1.986) (5.275,1.965) (5.325,1.942) (5.375,1.918) (5.425,1.892) (5.475,1.864) (5.525,1.835) (5.575,1.804) (5.618,1.775) (5.675,1.735) (5.722,1.700) (5.775,1.658) (5.814,1.625) (5.869,1.575) (5.921,1.525) (5.970,1.475) (6.016,1.425) (6.050,1.385) (6.098,1.325) (6.136,1.275) (6.172,1.225) (6.200,1.183) (6.236,1.125) (6.266,1.075) (6.294,1.025) (6.320,0.975) (6.344,0.925) (6.367,0.875) (6.388,0.825) (6.408,0.775) (6.426,0.725) (6.443,0.675) (6.459,0.625) (6.473,0.575) (6.486,0.525) (6.498,0.475) (6.509,0.425) (6.518,0.375) (6.526,0.325) (6.533,0.275) (6.539,0.225) (6.543,0.175) (6.547,0.125) (6.549,0.075) (6.550,0.025) (6.550,-0.037) (6.548,-0.100) (6.545,-0.150) (6.541,-0.200) (6.536,-0.250) (6.530,-0.300) (6.522,-0.350) (6.513,-0.400) (6.503,-0.450) (6.492,-0.500) (6.480,-0.550) (6.466,-0.600) (6.451,-0.650) (6.435,-0.700) (6.417,-0.750) (6.398,-0.800) (6.378,-0.850) (6.356,-0.900) (6.332,-0.950) (6.307,-1.000) (6.280,-1.050) (6.251,-1.100) (6.225,-1.143) (6.189,-1.200) (6.154,-1.250) (6.118,-1.300) (6.079,-1.350) (6.037,-1.400) (6.000,-1.442) (5.950,-1.496) (5.900,-1.546) (5.850,-1.593) (5.800,-1.637) (5.754,-1.675) (5.700,-1.717) (5.654,-1.750) (5.600,-1.787) (5.550,-1.819) (5.500,-1.850) (5.456,-1.875) (5.409,-1.900) (5.360,-1.925) (5.307,-1.950) (5.251,-1.975) (5.200,-1.996) (5.150,-2.015) (5.100,-2.033) (5.050,-2.049) (5.000,-2.064) (4.950,-2.078) (4.900,-2.091) (4.850,-2.102) (4.800,-2.113) (4.750,-2.122) (4.700,-2.130) (4.650,-2.137) (4.600,-2.143) (4.550,-2.148) (4.500,-2.152) (4.450,-2.155) (4.400,-2.157) (4.350,-2.158) (4.300,-2.158) (4.250,-2.158) (4.200,-2.156) (4.150,-2.153) (4.100,-2.150) (4.050,-2.146) (4.000,-2.140) (3.950,-2.134) (3.900,-2.128) (3.850,-2.120) (3.800,-2.111) (3.750,-2.102) (3.700,-2.092) (3.650,-2.081) (3.600,-2.070) (3.550,-2.057) (3.500,-2.044) (3.450,-2.030) (3.400,-2.016) (3.350,-2.001) (3.300,-1.985) (3.250,-1.968) (3.200,-1.951) (3.150,-1.933) (3.100,-1.914) (3.050,-1.895) (3.001,-1.875) (2.950,-1.854) (2.900,-1.833) (2.850,-1.811) (2.800,-1.789) (2.750,-1.766) (2.700,-1.742) (2.650,-1.718) (2.600,-1.694) (2.550,-1.669) (2.500,-1.643) (2.450,-1.617) (2.400,-1.590) (2.350,-1.563) (2.300,-1.535) (2.250,-1.507) (2.200,-1.479) (2.150,-1.450) (2.100,-1.421) (2.050,-1.391) (2.000,-1.361) (1.950,-1.331) (1.900,-1.300) (1.850,-1.269) (1.800,-1.238) (1.750,-1.206) (1.702,-1.175) (1.650,-1.142) (1.600,-1.109) (1.550,-1.076) (1.500,-1.044) (1.450,-1.010) (1.400,-0.977) (1.350,-0.943) (1.300,-0.910) (1.250,-0.876) (1.200,-0.842) (1.150,-0.807) (1.103,-0.775) (1.050,-0.739) (1.000,-0.704) (0.950,-0.669) (0.900,-0.634) (0.851,-0.600) (0.800,-0.565) (0.750,-0.529) (0.709,-0.500) (0.650,-0.459) (0.602,-0.425) (0.550,-0.389) (0.500,-0.354) (0.450,-0.319) (0.400,-0.283) (0.354,-0.250) (0.300,-0.213) (0.250,-0.177) (0.206,-0.144)}
\EQ{(-0.220,-0.155) (-0.275,-0.195) (-0.319,-0.225) (-0.375,-0.266) (-0.424,-0.300) (-0.475,-0.336) (-0.525,-0.371) (-0.567,-0.400) (-0.625,-0.442) (-0.673,-0.475) (-0.725,-0.512) (-0.775,-0.547) (-0.825,-0.583) (-0.875,-0.618) (-0.922,-0.650) (-0.975,-0.688) (-1.025,-0.722) (-1.075,-0.757) (-1.125,-0.792) (-1.173,-0.825) (-1.225,-0.861) (-1.275,-0.896) (-1.318,-0.925) (-1.375,-0.964) (-1.425,-0.998) (-1.475,-1.032) (-1.525,-1.066) (-1.575,-1.100) (-1.625,-1.133) (-1.675,-1.167) (-1.725,-1.200) (-1.775,-1.233) (-1.825,-1.266) (-1.875,-1.298) (-1.925,-1.330) (-1.975,-1.362) (-2.025,-1.394) (-2.074,-1.425) (-2.125,-1.457) (-2.175,-1.488) (-2.225,-1.519) (-2.275,-1.549) (-2.325,-1.580) (-2.375,-1.609) (-2.425,-1.639) (-2.475,-1.668) (-2.525,-1.697) (-2.575,-1.725) (-2.619,-1.750) (-2.665,-1.775) (-2.710,-1.800) (-2.757,-1.825) (-2.803,-1.850) (-2.850,-1.874) (-2.900,-1.900) (-2.949,-1.925) (-2.999,-1.950) (-3.050,-1.975) (-3.100,-1.999) (-3.150,-2.022) (-3.200,-2.045) (-3.250,-2.068) (-3.300,-2.090) (-3.350,-2.111) (-3.400,-2.132) (-3.450,-2.152) (-3.500,-2.172) (-3.550,-2.191) (-3.600,-2.210) (-3.650,-2.228) (-3.700,-2.246) (-3.750,-2.263) (-3.800,-2.279) (-3.850,-2.295) (-3.900,-2.310) (-3.950,-2.324) (-4.000,-2.338) (-4.050,-2.352) (-4.100,-2.364) (-4.150,-2.376) (-4.200,-2.387) (-4.250,-2.398) (-4.300,-2.408) (-4.350,-2.417) (-4.400,-2.425) (-4.450,-2.433) (-4.500,-2.440) (-4.550,-2.447) (-4.600,-2.452) (-4.650,-2.457) (-4.700,-2.462) (-4.750,-2.465) (-4.800,-2.468) (-4.850,-2.469) (-4.900,-2.470) (-4.950,-2.471) (-5.000,-2.470) (-5.050,-2.469) (-5.100,-2.467) (-5.150,-2.464) (-5.200,-2.460) (-5.250,-2.455) (-5.300,-2.449) (-5.350,-2.443) (-5.400,-2.435) (-5.450,-2.427) (-5.500,-2.417) (-5.550,-2.407) (-5.600,-2.395) (-5.650,-2.383) (-5.700,-2.370) (-5.750,-2.355) (-5.800,-2.339) (-5.850,-2.323) (-5.900,-2.305) (-5.950,-2.286) (-6.000,-2.265) (-6.050,-2.244) (-6.100,-2.221) (-6.150,-2.196) (-6.200,-2.171) (-6.250,-2.144) (-6.300,-2.115) (-6.350,-2.085) (-6.400,-2.053) (-6.450,-2.019) (-6.500,-1.983) (-6.545,-1.950) (-6.600,-1.906) (-6.650,-1.864) (-6.695,-1.825) (-6.748,-1.775) (-6.799,-1.725) (-6.846,-1.675) (-6.892,-1.625) (-6.925,-1.586) (-6.974,-1.525) (-7.013,-1.475) (-7.049,-1.425) (-7.083,-1.375) (-7.116,-1.325) (-7.147,-1.275) (-7.175,-1.226) (-7.200,-1.181) (-7.225,-1.134) (-7.250,-1.084) (-7.275,-1.031) (-7.300,-0.975) (-7.320,-0.925) (-7.340,-0.875) (-7.358,-0.825) (-7.375,-0.775) (-7.391,-0.725) (-7.405,-0.675) (-7.419,-0.625) (-7.431,-0.575) (-7.443,-0.525) (-7.453,-0.475) (-7.462,-0.425) (-7.470,-0.375) (-7.477,-0.325) (-7.483,-0.275) (-7.488,-0.225) (-7.492,-0.175) (-7.495,-0.125) (-7.497,-0.075) (-7.498,-0.025) (-7.498,0.050) (-7.496,0.100) (-7.494,0.150) (-7.490,0.200) (-7.486,0.250) (-7.480,0.300) (-7.474,0.350) (-7.466,0.400) (-7.457,0.450) (-7.448,0.500) (-7.437,0.550) (-7.425,0.600) (-7.412,0.650) (-7.398,0.700) (-7.383,0.750) (-7.367,0.800) (-7.349,0.850) (-7.330,0.900) (-7.310,0.950) (-7.289,1.000) (-7.266,1.050) (-7.242,1.100) (-7.217,1.150) (-7.190,1.200) (-7.161,1.250) (-7.131,1.300) (-7.100,1.350) (-7.066,1.400) (-7.031,1.450) (-7.000,1.492) (-6.955,1.550) (-6.913,1.600) (-6.875,1.644) (-6.825,1.698) (-6.775,1.749) (-6.725,1.797) (-6.675,1.843) (-6.625,1.886) (-6.577,1.925) (-6.525,1.965) (-6.477,2.000) (-6.425,2.036) (-6.375,2.069) (-6.325,2.100) (-6.275,2.129) (-6.225,2.157) (-6.175,2.184) (-6.125,2.209) (-6.075,2.232) (-6.025,2.255) (-5.977,2.275) (-5.925,2.295) (-5.875,2.314) (-5.825,2.331) (-5.775,2.347) (-5.725,2.362) (-5.675,2.376) (-5.625,2.389) (-5.575,2.401) (-5.525,2.412) (-5.475,2.422) (-5.425,2.431) (-5.375,2.439) (-5.325,2.446) (-5.275,2.452) (-5.225,2.457) (-5.175,2.462) (-5.125,2.465) (-5.075,2.468) (-5.025,2.470) (-4.975,2.471) (-4.925,2.471) (-4.875,2.470) (-4.825,2.469) (-4.775,2.466) (-4.725,2.463) (-4.675,2.460) (-4.625,2.455) (-4.575,2.450) (-4.525,2.444) (-4.475,2.437) (-4.425,2.430) (-4.375,2.421) (-4.325,2.412) (-4.275,2.403) (-4.225,2.393) (-4.175,2.382) (-4.125,2.370) (-4.075,2.358) (-4.025,2.345) (-3.975,2.331) (-3.925,2.317) (-3.875,2.302) (-3.825,2.287) (-3.775,2.271) (-3.725,2.254) (-3.675,2.237) (-3.625,2.219) (-3.575,2.201) (-3.525,2.182) (-3.475,2.162) (-3.425,2.142) (-3.375,2.122) (-3.325,2.100) (-3.275,2.079) (-3.225,2.057) (-3.175,2.034) (-3.125,2.011) (-3.075,1.987) (-3.025,1.963) (-2.975,1.938) (-2.925,1.913) (-2.875,1.887) (-2.825,1.861) (-2.775,1.835) (-2.725,1.808) (-2.675,1.781) (-2.625,1.753) (-2.575,1.725) (-2.531,1.700) (-2.487,1.675) (-2.425,1.639) (-2.375,1.609) (-2.325,1.580) (-2.276,1.550) (-2.225,1.519) (-2.175,1.488) (-2.125,1.457) (-2.075,1.426) (-2.025,1.394) (-1.975,1.362) (-1.925,1.330) (-1.878,1.300) (-1.825,1.266) (-1.775,1.233) (-1.725,1.200) (-1.675,1.167) (-1.625,1.133) (-1.576,1.100) (-1.525,1.066) (-1.475,1.032) (-1.428,1.000) (-1.375,0.964) (-1.325,0.930) (-1.282,0.900) (-1.225,0.861) (-1.175,0.827) (-1.125,0.792) (-1.075,0.757) (-1.029,0.725) (-0.975,0.688) (-0.925,0.652) (-0.875,0.618) (-0.825,0.583) (-0.779,0.550) (-0.725,0.512) (-0.675,0.477) (-0.625,0.442) (-0.575,0.407) (-0.531,0.375) (-0.475,0.336) (-0.425,0.300) (-0.375,0.266) (-0.325,0.230) (-0.284,0.200) (-0.225,0.160) (-0.220,0.155)}
\EQ{(0.206,0.144) (0.248,0.175) (0.300,0.213) (0.350,0.248) (0.400,0.283) (0.450,0.319) (0.496,0.350) (0.550,0.389) (0.600,0.424) (0.650,0.460) (0.700,0.495) (0.744,0.525) (0.800,0.565) (0.850,0.600) (0.900,0.635) (0.950,0.670) (0.993,0.700) (1.050,0.740) (1.100,0.774) (1.150,0.809) (1.200,0.844) (1.245,0.875) (1.300,0.913) (1.350,0.947) (1.400,0.981) (1.450,1.015) (1.500,1.049) (1.550,1.083) (1.600,1.117) (1.650,1.150) (1.700,1.183) (1.750,1.216) (1.800,1.249) (1.850,1.282) (1.900,1.314) (1.950,1.346) (1.995,1.375) (2.050,1.410) (2.100,1.442) (2.150,1.473) (2.194,1.500) (2.250,1.534) (2.300,1.565) (2.350,1.594) (2.400,1.624) (2.444,1.650) (2.500,1.682) (2.550,1.711) (2.600,1.739) (2.650,1.767) (2.700,1.795) (2.750,1.822) (2.800,1.848) (2.850,1.874) (2.900,1.900) (2.949,1.925) (2.999,1.950) (3.050,1.975) (3.100,1.999) (3.150,2.022) (3.200,2.045) (3.250,2.068) (3.300,2.090) (3.350,2.111) (3.400,2.132) (3.450,2.152) (3.500,2.172) (3.550,2.191) (3.600,2.210) (3.650,2.228) (3.700,2.246) (3.750,2.263) (3.800,2.279) (3.850,2.295) (3.900,2.310) (3.950,2.324) (4.000,2.338) (4.050,2.352) (4.100,2.364) (4.150,2.376) (4.200,2.387) (4.250,2.398) (4.300,2.408) (4.350,2.417) (4.400,2.425) (4.450,2.433) (4.500,2.440) (4.550,2.447) (4.600,2.452) (4.650,2.457) (4.700,2.462) (4.750,2.465) (4.800,2.468) (4.850,2.469) (4.900,2.470) (4.950,2.471) (5.000,2.470) (5.050,2.469) (5.100,2.467) (5.150,2.464) (5.200,2.460) (5.250,2.455) (5.300,2.449) (5.350,2.443) (5.400,2.435) (5.450,2.427) (5.500,2.417) (5.550,2.407) (5.600,2.395) (5.650,2.383) (5.700,2.370) (5.750,2.355) (5.800,2.339) (5.850,2.323) (5.900,2.305) (5.950,2.286) (6.000,2.265) (6.050,2.244) (6.100,2.221) (6.150,2.196) (6.200,2.171) (6.250,2.144) (6.300,2.115) (6.350,2.085) (6.400,2.053) (6.450,2.019) (6.500,1.983) (6.545,1.950) (6.600,1.906) (6.650,1.864) (6.695,1.825) (6.748,1.775) (6.799,1.725) (6.846,1.675) (6.892,1.625) (6.925,1.586) (6.974,1.525) (7.013,1.475) (7.049,1.425) (7.083,1.375) (7.116,1.325) (7.147,1.275) (7.175,1.226) (7.200,1.181) (7.225,1.134) (7.250,1.084) (7.275,1.031) (7.300,0.975) (7.320,0.925) (7.340,0.875) (7.358,0.825) (7.375,0.775) (7.391,0.725) (7.405,0.675) (7.419,0.625) (7.431,0.575) (7.443,0.525) (7.453,0.475) (7.462,0.425) (7.470,0.375) (7.477,0.325) (7.483,0.275) (7.488,0.225) (7.492,0.175) (7.495,0.125) (7.497,0.075) (7.498,0.025) (7.498,-0.050) (7.496,-0.100) (7.494,-0.150) (7.490,-0.200) (7.486,-0.250) (7.480,-0.300) (7.474,-0.350) (7.466,-0.400) (7.457,-0.450) (7.448,-0.500) (7.437,-0.550) (7.425,-0.600) (7.412,-0.650) (7.398,-0.700) (7.383,-0.750) (7.367,-0.800) (7.349,-0.850) (7.330,-0.900) (7.310,-0.950) (7.289,-1.000) (7.266,-1.050) (7.242,-1.100) (7.217,-1.150) (7.190,-1.200) (7.161,-1.250) (7.131,-1.300) (7.100,-1.350) (7.066,-1.400) (7.031,-1.450) (7.000,-1.492) (6.955,-1.550) (6.913,-1.600) (6.875,-1.644) (6.825,-1.698) (6.775,-1.749) (6.725,-1.797) (6.675,-1.843) (6.625,-1.886) (6.577,-1.925) (6.525,-1.965) (6.477,-2.000) (6.425,-2.036) (6.375,-2.069) (6.325,-2.100) (6.275,-2.129) (6.225,-2.157) (6.175,-2.184) (6.125,-2.209) (6.075,-2.232) (6.025,-2.255) (5.977,-2.275) (5.925,-2.295) (5.875,-2.314) (5.825,-2.331) (5.775,-2.347) (5.725,-2.362) (5.675,-2.376) (5.625,-2.389) (5.575,-2.401) (5.525,-2.412) (5.475,-2.422) (5.425,-2.431) (5.375,-2.439) (5.325,-2.446) (5.275,-2.452) (5.225,-2.457) (5.175,-2.462) (5.125,-2.465) (5.075,-2.468) (5.025,-2.470) (4.975,-2.471) (4.925,-2.471) (4.875,-2.470) (4.825,-2.469) (4.775,-2.466) (4.725,-2.463) (4.675,-2.460) (4.625,-2.455) (4.575,-2.450) (4.525,-2.444) (4.475,-2.437) (4.425,-2.430) (4.375,-2.421) (4.325,-2.412) (4.275,-2.403) (4.225,-2.393) (4.175,-2.382) (4.125,-2.370) (4.075,-2.358) (4.025,-2.345) (3.975,-2.331) (3.925,-2.317) (3.875,-2.302) (3.825,-2.287) (3.775,-2.271) (3.725,-2.254) (3.675,-2.237) (3.625,-2.219) (3.575,-2.201) (3.525,-2.182) (3.475,-2.162) (3.425,-2.142) (3.375,-2.122) (3.325,-2.100) (3.275,-2.079) (3.225,-2.057) (3.175,-2.034) (3.125,-2.011) (3.075,-1.987) (3.025,-1.963) (2.975,-1.938) (2.925,-1.913) (2.875,-1.887) (2.825,-1.861) (2.775,-1.835) (2.725,-1.808) (2.675,-1.781) (2.625,-1.753) (2.575,-1.725) (2.531,-1.700) (2.487,-1.675) (2.425,-1.639) (2.375,-1.609) (2.325,-1.580) (2.276,-1.550) (2.225,-1.519) (2.175,-1.488) (2.125,-1.457) (2.075,-1.426) (2.025,-1.394) (1.975,-1.362) (1.925,-1.330) (1.878,-1.300) (1.825,-1.266) (1.775,-1.233) (1.725,-1.200) (1.675,-1.167) (1.625,-1.133) (1.576,-1.100) (1.525,-1.066) (1.475,-1.032) (1.428,-1.000) (1.375,-0.964) (1.325,-0.930) (1.282,-0.900) (1.225,-0.861) (1.175,-0.827) (1.125,-0.792) (1.075,-0.757) (1.029,-0.725) (0.975,-0.688) (0.925,-0.652) (0.875,-0.618) (0.825,-0.583) (0.779,-0.550) (0.725,-0.512) (0.675,-0.477) (0.625,-0.442) (0.575,-0.407) (0.531,-0.375) (0.475,-0.336) (0.425,-0.300) (0.375,-0.266) (0.325,-0.230) (0.284,-0.200) (0.225,-0.160) (0.206,-0.144)}
\EQ{(-0.220,-0.155) (-0.275,-0.195) (-0.319,-0.225) (-0.375,-0.266) (-0.424,-0.300) (-0.475,-0.336) (-0.525,-0.371) (-0.567,-0.400) (-0.625,-0.442) (-0.672,-0.475) (-0.725,-0.513) (-0.775,-0.548) (-0.825,-0.583) (-0.875,-0.618) (-0.921,-0.650) (-0.975,-0.689) (-1.025,-0.724) (-1.075,-0.759) (-1.125,-0.794) (-1.170,-0.825) (-1.225,-0.864) (-1.275,-0.899) (-1.325,-0.934) (-1.375,-0.968) (-1.421,-1.000) (-1.475,-1.038) (-1.525,-1.072) (-1.575,-1.107) (-1.625,-1.141) (-1.674,-1.175) (-1.725,-1.210) (-1.775,-1.244) (-1.821,-1.275) (-1.875,-1.312) (-1.925,-1.346) (-1.969,-1.375) (-2.025,-1.413) (-2.075,-1.446) (-2.118,-1.475) (-2.175,-1.513) (-2.225,-1.546) (-2.270,-1.575) (-2.325,-1.611) (-2.375,-1.644) (-2.424,-1.675) (-2.475,-1.708) (-2.525,-1.740) (-2.575,-1.772) (-2.620,-1.800) (-2.675,-1.834) (-2.725,-1.866) (-2.775,-1.896) (-2.822,-1.925) (-2.875,-1.957) (-2.925,-1.988) (-2.975,-2.017) (-3.025,-2.047) (-3.073,-2.075) (-3.125,-2.105) (-3.175,-2.134) (-3.225,-2.163) (-3.275,-2.191) (-3.325,-2.219) (-3.375,-2.246) (-3.425,-2.274) (-3.474,-2.300) (-3.521,-2.325) (-3.569,-2.350) (-3.617,-2.375) (-3.665,-2.400) (-3.715,-2.425) (-3.765,-2.450) (-3.816,-2.475) (-3.867,-2.500) (-3.920,-2.525) (-3.973,-2.550) (-4.025,-2.574) (-4.075,-2.596) (-4.125,-2.619) (-4.175,-2.640) (-4.225,-2.662) (-4.275,-2.683) (-4.325,-2.703) (-4.375,-2.723) (-4.425,-2.743) (-4.475,-2.762) (-4.525,-2.781) (-4.575,-2.799) (-4.625,-2.817) (-4.675,-2.834) (-4.725,-2.851) (-4.775,-2.867) (-4.825,-2.883) (-4.875,-2.899) (-4.925,-2.914) (-4.975,-2.928) (-5.025,-2.942) (-5.075,-2.956) (-5.125,-2.968) (-5.175,-2.981) (-5.225,-2.993) (-5.275,-3.004) (-5.325,-3.015) (-5.375,-3.025) (-5.425,-3.035) (-5.475,-3.044) (-5.525,-3.053) (-5.575,-3.061) (-5.625,-3.069) (-5.675,-3.075) (-5.725,-3.082) (-5.775,-3.088) (-5.825,-3.093) (-5.875,-3.098) (-5.925,-3.102) (-5.975,-3.105) (-6.025,-3.108) (-6.075,-3.110) (-6.125,-3.112) (-6.175,-3.113) (-6.225,-3.113) (-6.275,-3.113) (-6.325,-3.112) (-6.375,-3.110) (-6.425,-3.108) (-6.475,-3.105) (-6.525,-3.101) (-6.575,-3.097) (-6.625,-3.092) (-6.675,-3.086) (-6.725,-3.080) (-6.775,-3.072) (-6.825,-3.065) (-6.875,-3.056) (-6.925,-3.046) (-6.975,-3.036) (-7.025,-3.025) (-7.075,-3.013) (-7.125,-3.001) (-7.175,-2.987) (-7.225,-2.973) (-7.275,-2.958) (-7.325,-2.942) (-7.375,-2.925) (-7.425,-2.907) (-7.475,-2.888) (-7.525,-2.868) (-7.575,-2.848) (-7.625,-2.826) (-7.675,-2.803) (-7.725,-2.779) (-7.775,-2.754) (-7.825,-2.728) (-7.875,-2.701) (-7.920,-2.675) (-7.975,-2.642) (-8.025,-2.611) (-8.075,-2.579) (-8.117,-2.550) (-8.175,-2.509) (-8.222,-2.475) (-8.275,-2.434) (-8.318,-2.400) (-8.375,-2.352) (-8.425,-2.308) (-8.475,-2.262) (-8.514,-2.225) (-8.550,-2.189) (-8.589,-2.150) (-8.625,-2.111) (-8.675,-2.056) (-8.723,-2.000) (-8.765,-1.950) (-8.800,-1.905) (-8.842,-1.850) (-8.875,-1.804) (-8.912,-1.750) (-8.946,-1.700) (-8.975,-1.654) (-9.008,-1.600) (-9.037,-1.550) (-9.065,-1.500) (-9.092,-1.450) (-9.118,-1.400) (-9.142,-1.350) (-9.165,-1.300) (-9.188,-1.250) (-9.209,-1.200) (-9.229,-1.150) (-9.248,-1.100) (-9.267,-1.050) (-9.284,-1.000) (-9.300,-0.951) (-9.316,-0.900) (-9.330,-0.850) (-9.344,-0.800) (-9.357,-0.750) (-9.368,-0.700) (-9.379,-0.650) (-9.390,-0.600) (-9.399,-0.550) (-9.407,-0.500) (-9.415,-0.450) (-9.422,-0.400) (-9.428,-0.350) (-9.433,-0.300) (-9.437,-0.250) (-9.441,-0.200) (-9.444,-0.150) (-9.446,-0.100) (-9.447,-0.050) (-9.447,0.000) (-9.447,0.050) (-9.446,0.100) (-9.444,0.150) (-9.441,0.200) (-9.437,0.250) (-9.433,0.300) (-9.428,0.350) (-9.422,0.400) (-9.415,0.450) (-9.407,0.500) (-9.399,0.550) (-9.390,0.600) (-9.379,0.650) (-9.368,0.700) (-9.357,0.750) (-9.344,0.800) (-9.330,0.850) (-9.316,0.900) (-9.300,0.950) (-9.284,1.000) (-9.267,1.050) (-9.248,1.100) (-9.229,1.150) (-9.209,1.200) (-9.188,1.250) (-9.165,1.300) (-9.142,1.350) (-9.118,1.400) (-9.092,1.450) (-9.065,1.500) (-9.037,1.550) (-9.008,1.600) (-8.978,1.650) (-8.950,1.693) (-8.912,1.750) (-8.878,1.800) (-8.842,1.850) (-8.804,1.900) (-8.765,1.950) (-8.725,1.998) (-8.680,2.050) (-8.636,2.100) (-8.600,2.138) (-8.564,2.175) (-8.525,2.214) (-8.488,2.250) (-8.450,2.285) (-8.406,2.325) (-8.350,2.373) (-8.300,2.414) (-8.255,2.450) (-8.200,2.491) (-8.153,2.525) (-8.100,2.562) (-8.050,2.595) (-8.003,2.625) (-7.950,2.657) (-7.900,2.686) (-7.850,2.714) (-7.800,2.741) (-7.750,2.767) (-7.700,2.791) (-7.650,2.814) (-7.600,2.837) (-7.550,2.858) (-7.500,2.878) (-7.450,2.898) (-7.400,2.916) (-7.350,2.933) (-7.300,2.950) (-7.250,2.966) (-7.200,2.980) (-7.150,2.994) (-7.100,3.007) (-7.050,3.019) (-7.000,3.031) (-6.950,3.041) (-6.900,3.051) (-6.850,3.060) (-6.800,3.069) (-6.750,3.076) (-6.700,3.083) (-6.650,3.089) (-6.600,3.094) (-6.550,3.099) (-6.500,3.103) (-6.450,3.106) (-6.400,3.109) (-6.350,3.111) (-6.300,3.112) (-6.250,3.113) (-6.200,3.113) (-6.150,3.112) (-6.100,3.111) (-6.050,3.109) (-6.000,3.107) (-5.950,3.103) (-5.900,3.100) (-5.850,3.095) (-5.800,3.090) (-5.750,3.085) (-5.700,3.079) (-5.650,3.072) (-5.600,3.065) (-5.550,3.057) (-5.500,3.049) (-5.450,3.040) (-5.400,3.030) (-5.350,3.020) (-5.300,3.010) (-5.250,2.998) (-5.200,2.987) (-5.151,2.975) (-5.100,2.962) (-5.050,2.949) (-5.000,2.935) (-4.950,2.921) (-4.900,2.906) (-4.850,2.891) (-4.800,2.875) (-4.750,2.859) (-4.700,2.843) (-4.650,2.825) (-4.600,2.808) (-4.550,2.790) (-4.500,2.771) (-4.450,2.752) (-4.400,2.733) (-4.350,2.713) (-4.300,2.693) (-4.250,2.672) (-4.200,2.651) (-4.150,2.629) (-4.100,2.608) (-4.050,2.585) (-4.000,2.562) (-3.950,2.539) (-3.900,2.516) (-3.850,2.492) (-3.800,2.467) (-3.750,2.443) (-3.700,2.418) (-3.650,2.392) (-3.600,2.366) (-3.550,2.340) (-3.500,2.314) (-3.450,2.287) (-3.400,2.260) (-3.350,2.233) (-3.300,2.205) (-3.250,2.177) (-3.203,2.150) (-3.159,2.125) (-3.116,2.100) (-3.050,2.062) (-3.000,2.032) (-2.950,2.002) (-2.904,1.975) (-2.850,1.942) (-2.800,1.912) (-2.750,1.881) (-2.700,1.850) (-2.650,1.819) (-2.600,1.787) (-2.550,1.756) (-2.502,1.725) (-2.450,1.692) (-2.400,1.660) (-2.350,1.627) (-2.300,1.595) (-2.250,1.562) (-2.200,1.529) (-2.156,1.500) (-2.100,1.463) (-2.050,1.430) (-2.006,1.400) (-1.950,1.363) (-1.900,1.329) (-1.858,1.300) (-1.800,1.261) (-1.750,1.227) (-1.700,1.193) (-1.650,1.158) (-1.602,1.125) (-1.550,1.090) (-1.500,1.055) (-1.457,1.025) (-1.400,0.986) (-1.350,0.951) (-1.300,0.916) (-1.250,0.881) (-1.206,0.850) (-1.150,0.811) (-1.100,0.776) (-1.050,0.741) (-1.000,0.706) (-0.956,0.675) (-0.900,0.636) (-0.850,0.600) (-0.800,0.566) (-0.750,0.530) (-0.708,0.500) (-0.650,0.460) (-0.601,0.425) (-0.550,0.389) (-0.500,0.354) (-0.450,0.319) (-0.400,0.283) (-0.354,0.250) (-0.300,0.213) (-0.250,0.177) (-0.220,0.155)}
\EQ{(0.206,0.144) (0.248,0.175) (0.300,0.213) (0.350,0.248) (0.400,0.283) (0.450,0.319) (0.496,0.350) (0.550,0.389) (0.600,0.424) (0.650,0.460) (0.700,0.495) (0.743,0.525) (0.800,0.566) (0.849,0.600) (0.900,0.636) (0.950,0.671) (0.992,0.700) (1.050,0.741) (1.098,0.775) (1.150,0.811) (1.200,0.846) (1.242,0.875) (1.300,0.916) (1.349,0.950) (1.400,0.986) (1.450,1.020) (1.493,1.050) (1.550,1.090) (1.600,1.124) (1.650,1.158) (1.700,1.193) (1.748,1.225) (1.800,1.261) (1.850,1.295) (1.895,1.325) (1.950,1.363) (2.000,1.396) (2.043,1.425) (2.100,1.463) (2.150,1.496) (2.194,1.525) (2.250,1.562) (2.300,1.595) (2.347,1.625) (2.400,1.660) (2.450,1.692) (2.500,1.724) (2.550,1.756) (2.600,1.787) (2.650,1.819) (2.700,1.850) (2.750,1.881) (2.800,1.912) (2.850,1.942) (2.900,1.972) (2.946,2.000) (3.000,2.032) (3.050,2.062) (3.100,2.091) (3.150,2.120) (3.200,2.148) (3.247,2.175) (3.292,2.200) (3.337,2.225) (3.382,2.250) (3.428,2.275) (3.474,2.300) (3.521,2.325) (3.569,2.350) (3.617,2.375) (3.665,2.400) (3.715,2.425) (3.765,2.450) (3.816,2.475) (3.867,2.500) (3.920,2.525) (3.973,2.550) (4.025,2.574) (4.075,2.596) (4.125,2.619) (4.175,2.640) (4.225,2.662) (4.275,2.683) (4.325,2.703) (4.375,2.723) (4.425,2.743) (4.475,2.762) (4.525,2.781) (4.575,2.799) (4.625,2.817) (4.675,2.834) (4.725,2.851) (4.775,2.867) (4.825,2.883) (4.875,2.899) (4.925,2.914) (4.975,2.928) (5.025,2.942) (5.075,2.956) (5.125,2.968) (5.175,2.981) (5.225,2.993) (5.275,3.004) (5.325,3.015) (5.375,3.025) (5.425,3.035) (5.475,3.044) (5.525,3.053) (5.575,3.061) (5.625,3.069) (5.675,3.075) (5.725,3.082) (5.775,3.088) (5.825,3.093) (5.875,3.098) (5.925,3.102) (5.975,3.105) (6.025,3.108) (6.075,3.110) (6.125,3.112) (6.175,3.113) (6.225,3.113) (6.275,3.113) (6.325,3.112) (6.375,3.110) (6.425,3.108) (6.475,3.105) (6.525,3.101) (6.575,3.097) (6.625,3.092) (6.675,3.086) (6.725,3.080) (6.775,3.072) (6.825,3.065) (6.875,3.056) (6.925,3.046) (6.975,3.036) (7.025,3.025) (7.075,3.013) (7.125,3.001) (7.175,2.987) (7.225,2.973) (7.275,2.958) (7.325,2.942) (7.375,2.925) (7.425,2.907) (7.475,2.888) (7.525,2.868) (7.575,2.848) (7.625,2.826) (7.675,2.803) (7.725,2.779) (7.775,2.754) (7.825,2.728) (7.875,2.701) (7.920,2.675) (7.975,2.642) (8.025,2.611) (8.075,2.579) (8.117,2.550) (8.175,2.509) (8.222,2.475) (8.275,2.434) (8.318,2.400) (8.375,2.352) (8.425,2.308) (8.475,2.262) (8.514,2.225) (8.550,2.189) (8.589,2.150) (8.625,2.111) (8.675,2.056) (8.723,2.000) (8.765,1.950) (8.800,1.905) (8.842,1.850) (8.875,1.804) (8.912,1.750) (8.946,1.700) (8.975,1.654) (9.008,1.600) (9.037,1.550) (9.065,1.500) (9.092,1.450) (9.118,1.400) (9.142,1.350) (9.165,1.300) (9.188,1.250) (9.209,1.200) (9.229,1.150) (9.248,1.100) (9.267,1.050) (9.284,1.000) (9.300,0.951) (9.316,0.900) (9.330,0.850) (9.344,0.800) (9.357,0.750) (9.368,0.700) (9.379,0.650) (9.390,0.600) (9.399,0.550) (9.407,0.500) (9.415,0.450) (9.422,0.400) (9.428,0.350) (9.433,0.300) (9.437,0.250) (9.441,0.200) (9.444,0.150) (9.446,0.100) (9.447,0.050) (9.447,0.000) (9.447,-0.050) (9.446,-0.100) (9.444,-0.150) (9.441,-0.200) (9.437,-0.250) (9.433,-0.300) (9.428,-0.350) (9.422,-0.400) (9.415,-0.450) (9.407,-0.500) (9.399,-0.550) (9.390,-0.600) (9.379,-0.650) (9.368,-0.700) (9.357,-0.750) (9.344,-0.800) (9.330,-0.850) (9.316,-0.900) (9.300,-0.950) (9.284,-1.000) (9.267,-1.050) (9.248,-1.100) (9.229,-1.150) (9.209,-1.200) (9.188,-1.250) (9.165,-1.300) (9.142,-1.350) (9.118,-1.400) (9.092,-1.450) (9.065,-1.500) (9.037,-1.550) (9.008,-1.600) (8.978,-1.650) (8.950,-1.693) (8.912,-1.750) (8.878,-1.800) (8.842,-1.850) (8.804,-1.900) (8.765,-1.950) (8.725,-1.998) (8.680,-2.050) (8.636,-2.100) (8.600,-2.138) (8.564,-2.175) (8.525,-2.214) (8.488,-2.250) (8.450,-2.285) (8.406,-2.325) (8.350,-2.373) (8.300,-2.414) (8.255,-2.450) (8.200,-2.491) (8.153,-2.525) (8.100,-2.562) (8.050,-2.595) (8.003,-2.625) (7.950,-2.657) (7.900,-2.686) (7.850,-2.714) (7.800,-2.741) (7.750,-2.767) (7.700,-2.791) (7.650,-2.814) (7.600,-2.837) (7.550,-2.858) (7.500,-2.878) (7.450,-2.898) (7.400,-2.916) (7.350,-2.933) (7.300,-2.950) (7.250,-2.966) (7.200,-2.980) (7.150,-2.994) (7.100,-3.007) (7.050,-3.019) (7.000,-3.031) (6.950,-3.041) (6.900,-3.051) (6.850,-3.060) (6.800,-3.069) (6.750,-3.076) (6.700,-3.083) (6.650,-3.089) (6.600,-3.094) (6.550,-3.099) (6.500,-3.103) (6.450,-3.106) (6.400,-3.109) (6.350,-3.111) (6.300,-3.112) (6.250,-3.113) (6.200,-3.113) (6.150,-3.112) (6.100,-3.111) (6.050,-3.109) (6.000,-3.107) (5.950,-3.103) (5.900,-3.100) (5.850,-3.095) (5.800,-3.090) (5.750,-3.085) (5.700,-3.079) (5.650,-3.072) (5.600,-3.065) (5.550,-3.057) (5.500,-3.049) (5.450,-3.040) (5.400,-3.030) (5.350,-3.020) (5.300,-3.010) (5.250,-2.998) (5.200,-2.987) (5.151,-2.975) (5.100,-2.962) (5.050,-2.949) (5.000,-2.935) (4.950,-2.921) (4.900,-2.906) (4.850,-2.891) (4.800,-2.875) (4.750,-2.859) (4.700,-2.843) (4.650,-2.825) (4.600,-2.808) (4.550,-2.790) (4.500,-2.771) (4.450,-2.752) (4.400,-2.733) (4.350,-2.713) (4.300,-2.693) (4.250,-2.672) (4.200,-2.651) (4.150,-2.629) (4.100,-2.608) (4.050,-2.585) (4.000,-2.562) (3.950,-2.539) (3.900,-2.516) (3.850,-2.492) (3.800,-2.467) (3.750,-2.443) (3.700,-2.418) (3.650,-2.392) (3.600,-2.366) (3.550,-2.340) (3.500,-2.314) (3.450,-2.287) (3.400,-2.260) (3.350,-2.233) (3.300,-2.205) (3.250,-2.177) (3.203,-2.150) (3.159,-2.125) (3.116,-2.100) (3.050,-2.062) (3.000,-2.032) (2.950,-2.002) (2.904,-1.975) (2.850,-1.942) (2.800,-1.912) (2.750,-1.881) (2.700,-1.850) (2.650,-1.819) (2.600,-1.787) (2.550,-1.756) (2.502,-1.725) (2.450,-1.692) (2.400,-1.660) (2.350,-1.627) (2.300,-1.595) (2.250,-1.562) (2.200,-1.529) (2.156,-1.500) (2.100,-1.463) (2.050,-1.430) (2.006,-1.400) (1.950,-1.363) (1.900,-1.329) (1.858,-1.300) (1.800,-1.261) (1.750,-1.227) (1.700,-1.193) (1.650,-1.158) (1.602,-1.125) (1.550,-1.090) (1.500,-1.055) (1.457,-1.025) (1.400,-0.986) (1.350,-0.951) (1.300,-0.916) (1.250,-0.881) (1.206,-0.850) (1.150,-0.811) (1.100,-0.776) (1.050,-0.741) (1.000,-0.706) (0.956,-0.675) (0.900,-0.636) (0.850,-0.600) (0.800,-0.566) (0.750,-0.530) (0.708,-0.500) (0.650,-0.460) (0.601,-0.425) (0.550,-0.389) (0.500,-0.354) (0.450,-0.319) (0.400,-0.283) (0.354,-0.250) (0.300,-0.213) (0.250,-0.177) (0.206,-0.144)}
\EQ{(10.000,-7.071) (9.950,-7.036) (9.900,-7.000) (9.850,-6.965) (9.800,-6.930) (9.758,-6.900) (9.700,-6.859) (9.652,-6.825) (9.600,-6.788) (9.550,-6.753) (9.500,-6.718) (9.450,-6.682) (9.405,-6.650) (9.350,-6.611) (9.300,-6.576) (9.250,-6.541) (9.200,-6.505) (9.157,-6.475) (9.100,-6.435) (9.051,-6.400) (9.000,-6.364) (8.950,-6.329) (8.900,-6.293) (8.850,-6.258) (8.804,-6.225) (8.750,-6.187) (8.700,-6.152) (8.650,-6.117) (8.600,-6.081) (8.556,-6.050) (8.500,-6.010) (8.450,-5.975) (8.400,-5.940) (8.350,-5.904) (8.309,-5.875) (8.250,-5.834) (8.202,-5.800) (8.150,-5.763) (8.100,-5.728) (8.050,-5.692) (8.000,-5.657) (7.955,-5.625) (7.900,-5.586) (7.850,-5.551) (7.800,-5.515) (7.750,-5.480) (7.708,-5.450) (7.650,-5.409) (7.601,-5.375) (7.550,-5.339) (7.500,-5.303) (7.450,-5.268) (7.400,-5.233) (7.354,-5.200) (7.300,-5.162) (7.250,-5.127) (7.200,-5.091) (7.150,-5.056) (7.106,-5.025) (7.050,-4.985) (7.000,-4.950) (6.950,-4.914) (6.900,-4.879) (6.859,-4.850) (6.800,-4.808) (6.753,-4.775) (6.700,-4.738) (6.650,-4.702) (6.600,-4.667) (6.550,-4.632) (6.505,-4.600) (6.450,-4.561) (6.400,-4.525) (6.350,-4.490) (6.300,-4.455) (6.258,-4.425) (6.200,-4.384) (6.152,-4.350) (6.100,-4.313) (6.050,-4.278) (6.000,-4.243) (5.950,-4.207) (5.904,-4.175) (5.850,-4.137) (5.800,-4.101) (5.750,-4.066) (5.700,-4.031) (5.657,-4.000) (5.600,-3.960) (5.551,-3.925) (5.500,-3.889) (5.450,-3.854) (5.400,-3.818) (5.350,-3.783) (5.303,-3.750) (5.250,-3.712) (5.200,-3.677) (5.150,-3.642) (5.100,-3.606) (5.056,-3.575) (5.000,-3.536) (4.950,-3.500) (4.900,-3.465) (4.850,-3.429) (4.808,-3.400) (4.750,-3.359) (4.702,-3.325) (4.650,-3.288) (4.600,-3.253) (4.550,-3.217) (4.500,-3.182) (4.455,-3.150) (4.400,-3.111) (4.350,-3.076) (4.300,-3.041) (4.250,-3.005) (4.207,-2.975) (4.150,-2.935) (4.101,-2.900) (4.050,-2.864) (4.000,-2.828) (3.950,-2.793) (3.900,-2.758) (3.854,-2.725) (3.800,-2.687) (3.750,-2.652) (3.700,-2.616) (3.650,-2.581) (3.606,-2.550) (3.550,-2.510) (3.500,-2.475) (3.450,-2.440) (3.400,-2.404) (3.359,-2.375) (3.300,-2.334) (3.253,-2.300) (3.200,-2.263) (3.150,-2.227) (3.100,-2.192) (3.050,-2.157) (3.005,-2.125) (2.950,-2.086) (2.900,-2.051) (2.850,-2.015) (2.800,-1.980) (2.758,-1.950) (2.700,-1.909) (2.652,-1.875) (2.600,-1.839) (2.550,-1.803) (2.500,-1.768) (2.450,-1.732) (2.404,-1.700) (2.350,-1.662) (2.300,-1.626) (2.250,-1.591) (2.200,-1.556) (2.157,-1.525) (2.100,-1.485) (2.051,-1.450) (2.000,-1.414) (1.950,-1.379) (1.900,-1.344) (1.850,-1.308) (1.803,-1.275) (1.750,-1.238) (1.700,-1.202) (1.650,-1.167) (1.600,-1.131) (1.556,-1.100) (1.500,-1.061) (1.450,-1.025) (1.400,-0.990) (1.350,-0.955) (1.308,-0.925) (1.250,-0.884) (1.202,-0.850) (1.150,-0.813) (1.100,-0.778) (1.050,-0.743) (1.000,-0.707) (0.955,-0.675) (0.900,-0.637) (0.850,-0.601) (0.800,-0.566) (0.750,-0.531) (0.708,-0.500) (0.650,-0.460) (0.601,-0.425) (0.550,-0.389) (0.500,-0.354) (0.450,-0.319) (0.400,-0.283) (0.354,-0.250) (0.300,-0.213) (0.250,-0.177) (0.206,-0.144)}
\EQ{(-0.220,-0.155) (-0.275,-0.195) (-0.319,-0.225) (-0.375,-0.266) (-0.424,-0.300) (-0.475,-0.336) (-0.525,-0.371) (-0.566,-0.400) (-0.625,-0.442) (-0.672,-0.475) (-0.725,-0.513) (-0.775,-0.548) (-0.825,-0.584) (-0.875,-0.619) (-0.920,-0.650) (-0.975,-0.690) (-1.025,-0.725) (-1.075,-0.760) (-1.125,-0.796) (-1.167,-0.825) (-1.225,-0.866) (-1.273,-0.900) (-1.325,-0.937) (-1.375,-0.972) (-1.425,-1.008) (-1.475,-1.043) (-1.520,-1.075) (-1.575,-1.114) (-1.625,-1.149) (-1.675,-1.185) (-1.725,-1.220) (-1.768,-1.250) (-1.825,-1.291) (-1.874,-1.325) (-1.925,-1.361) (-1.975,-1.397) (-2.025,-1.432) (-2.075,-1.467) (-2.121,-1.500) (-2.175,-1.538) (-2.225,-1.573) (-2.275,-1.609) (-2.325,-1.644) (-2.369,-1.675) (-2.425,-1.715) (-2.475,-1.750) (-2.525,-1.786) (-2.575,-1.821) (-2.616,-1.850) (-2.675,-1.892) (-2.722,-1.925) (-2.775,-1.962) (-2.825,-1.998) (-2.875,-2.033) (-2.925,-2.068) (-2.970,-2.100) (-3.025,-2.139) (-3.075,-2.174) (-3.125,-2.210) (-3.175,-2.245) (-3.217,-2.275) (-3.275,-2.316) (-3.323,-2.350) (-3.375,-2.387) (-3.425,-2.422) (-3.475,-2.457) (-3.525,-2.493) (-3.571,-2.525) (-3.625,-2.563) (-3.675,-2.599) (-3.725,-2.634) (-3.775,-2.669) (-3.818,-2.700) (-3.875,-2.740) (-3.924,-2.775) (-3.975,-2.811) (-4.025,-2.846) (-4.066,-2.875) (-4.125,-2.917) (-4.172,-2.950) (-4.225,-2.988) (-4.275,-3.023) (-4.325,-3.058) (-4.375,-3.094) (-4.419,-3.125) (-4.475,-3.164) (-4.525,-3.200) (-4.575,-3.235) (-4.625,-3.270) (-4.667,-3.300) (-4.725,-3.341) (-4.773,-3.375) (-4.825,-3.412) (-4.875,-3.447) (-4.925,-3.483) (-4.975,-3.518) (-5.021,-3.550) (-5.075,-3.589) (-5.125,-3.624) (-5.175,-3.659) (-5.225,-3.695) (-5.268,-3.725) (-5.325,-3.765) (-5.374,-3.800) (-5.425,-3.836) (-5.475,-3.871) (-5.525,-3.907) (-5.575,-3.942) (-5.622,-3.975) (-5.675,-4.013) (-5.725,-4.048) (-5.775,-4.084) (-5.825,-4.119) (-5.869,-4.150) (-5.925,-4.190) (-5.975,-4.225) (-6.025,-4.260) (-6.075,-4.296) (-6.117,-4.325) (-6.175,-4.366) (-6.223,-4.400) (-6.275,-4.437) (-6.325,-4.472) (-6.375,-4.508) (-6.425,-4.543) (-6.470,-4.575) (-6.525,-4.614) (-6.575,-4.649) (-6.625,-4.685) (-6.675,-4.720) (-6.718,-4.750) (-6.775,-4.791) (-6.824,-4.825) (-6.875,-4.861) (-6.925,-4.897) (-6.975,-4.932) (-7.025,-4.967) (-7.071,-5.000) (-7.125,-5.038) (-7.175,-5.073) (-7.225,-5.109) (-7.275,-5.144) (-7.319,-5.175) (-7.375,-5.215) (-7.425,-5.250) (-7.475,-5.286) (-7.525,-5.321) (-7.566,-5.350) (-7.625,-5.392) (-7.672,-5.425) (-7.725,-5.462) (-7.775,-5.498) (-7.825,-5.533) (-7.875,-5.568) (-7.920,-5.600) (-7.975,-5.639) (-8.025,-5.675) (-8.075,-5.710) (-8.125,-5.745) (-8.167,-5.775) (-8.225,-5.816) (-8.273,-5.850) (-8.325,-5.887) (-8.375,-5.922) (-8.425,-5.957) (-8.475,-5.993) (-8.521,-6.025) (-8.575,-6.063) (-8.625,-6.099) (-8.675,-6.134) (-8.725,-6.170) (-8.768,-6.200) (-8.825,-6.240) (-8.874,-6.275) (-8.925,-6.311) (-8.975,-6.346) (-9.025,-6.382) (-9.075,-6.417) (-9.122,-6.450) (-9.175,-6.488) (-9.225,-6.523) (-9.275,-6.558) (-9.325,-6.594) (-9.369,-6.625) (-9.425,-6.665) (-9.475,-6.700) (-9.525,-6.735) (-9.575,-6.771) (-9.617,-6.800) (-9.675,-6.841) (-9.723,-6.875) (-9.775,-6.912) (-9.825,-6.947) (-9.875,-6.983) (-9.925,-7.018) (-9.970,-7.050) (-10.000,-7.071)}
\EQ{(0.206,0.144) (0.248,0.175) (0.300,0.213) (0.350,0.248) (0.400,0.283) (0.450,0.319) (0.496,0.350) (0.550,0.389) (0.600,0.424) (0.650,0.460) (0.700,0.495) (0.743,0.525) (0.800,0.566) (0.849,0.600) (0.900,0.637) (0.950,0.672) (1.000,0.707) (1.050,0.743) (1.096,0.775) (1.150,0.813) (1.200,0.849) (1.250,0.884) (1.300,0.919) (1.344,0.950) (1.400,0.990) (1.450,1.025) (1.500,1.061) (1.550,1.096) (1.591,1.125) (1.650,1.167) (1.697,1.200) (1.750,1.238) (1.800,1.273) (1.850,1.308) (1.900,1.344) (1.945,1.375) (2.000,1.414) (2.050,1.450) (2.100,1.485) (2.150,1.520) (2.192,1.550) (2.250,1.591) (2.298,1.625) (2.350,1.662) (2.400,1.697) (2.450,1.732) (2.500,1.768) (2.546,1.800) (2.600,1.839) (2.650,1.874) (2.700,1.909) (2.750,1.945) (2.793,1.975) (2.850,2.015) (2.899,2.050) (2.950,2.086) (3.000,2.121) (3.050,2.157) (3.100,2.192) (3.147,2.225) (3.200,2.263) (3.250,2.298) (3.300,2.334) (3.350,2.369) (3.394,2.400) (3.450,2.440) (3.500,2.475) (3.550,2.510) (3.600,2.546) (3.642,2.575) (3.700,2.616) (3.748,2.650) (3.800,2.687) (3.850,2.722) (3.900,2.758) (3.950,2.793) (3.995,2.825) (4.050,2.864) (4.100,2.899) (4.150,2.935) (4.200,2.970) (4.243,3.000) (4.300,3.041) (4.349,3.075) (4.400,3.111) (4.450,3.147) (4.500,3.182) (4.550,3.217) (4.596,3.250) (4.650,3.288) (4.700,3.323) (4.750,3.359) (4.800,3.394) (4.844,3.425) (4.900,3.465) (4.950,3.500) (5.000,3.536) (5.050,3.571) (5.091,3.600) (5.150,3.642) (5.197,3.675) (5.250,3.712) (5.300,3.748) (5.350,3.783) (5.400,3.818) (5.445,3.850) (5.500,3.889) (5.550,3.924) (5.600,3.960) (5.650,3.995) (5.692,4.025) (5.750,4.066) (5.798,4.100) (5.850,4.137) (5.900,4.172) (5.950,4.207) (6.000,4.243) (6.046,4.275) (6.100,4.313) (6.150,4.349) (6.200,4.384) (6.250,4.419) (6.293,4.450) (6.350,4.490) (6.399,4.525) (6.450,4.561) (6.500,4.596) (6.550,4.632) (6.600,4.667) (6.647,4.700) (6.700,4.738) (6.750,4.773) (6.800,4.808) (6.850,4.844) (6.894,4.875) (6.950,4.914) (7.000,4.950) (7.050,4.985) (7.100,5.020) (7.142,5.050) (7.200,5.091) (7.248,5.125) (7.300,5.162) (7.350,5.197) (7.400,5.233) (7.450,5.268) (7.495,5.300) (7.550,5.339) (7.600,5.374) (7.650,5.409) (7.700,5.445) (7.743,5.475) (7.800,5.515) (7.849,5.550) (7.900,5.586) (7.950,5.622) (8.000,5.657) (8.050,5.692) (8.096,5.725) (8.150,5.763) (8.200,5.798) (8.250,5.834) (8.300,5.869) (8.344,5.900) (8.400,5.940) (8.450,5.975) (8.500,6.010) (8.550,6.046) (8.591,6.075) (8.650,6.117) (8.697,6.150) (8.750,6.187) (8.800,6.223) (8.850,6.258) (8.900,6.293) (8.945,6.325) (9.000,6.364) (9.050,6.399) (9.100,6.435) (9.150,6.470) (9.192,6.500) (9.250,6.541) (9.298,6.575) (9.350,6.611) (9.400,6.647) (9.450,6.682) (9.500,6.718) (9.546,6.750) (9.600,6.788) (9.650,6.824) (9.700,6.859) (9.750,6.894) (9.793,6.925) (9.850,6.965) (9.899,7.000) (9.950,7.036) (10.000,7.071)}
\EQ{(-10.000,7.071) (-9.950,7.036) (-9.900,7.000) (-9.850,6.965) (-9.800,6.930) (-9.758,6.900) (-9.700,6.859) (-9.652,6.825) (-9.600,6.788) (-9.550,6.753) (-9.500,6.718) (-9.450,6.682) (-9.405,6.650) (-9.350,6.611) (-9.300,6.576) (-9.250,6.541) (-9.200,6.505) (-9.157,6.475) (-9.100,6.435) (-9.051,6.400) (-9.000,6.364) (-8.950,6.329) (-8.900,6.293) (-8.850,6.258) (-8.804,6.225) (-8.750,6.187) (-8.700,6.152) (-8.650,6.117) (-8.600,6.081) (-8.556,6.050) (-8.500,6.010) (-8.450,5.975) (-8.400,5.940) (-8.350,5.904) (-8.309,5.875) (-8.250,5.834) (-8.202,5.800) (-8.150,5.763) (-8.100,5.728) (-8.050,5.692) (-8.000,5.657) (-7.955,5.625) (-7.900,5.586) (-7.850,5.551) (-7.800,5.515) (-7.750,5.480) (-7.708,5.450) (-7.650,5.409) (-7.601,5.375) (-7.550,5.339) (-7.500,5.303) (-7.450,5.268) (-7.400,5.233) (-7.354,5.200) (-7.300,5.162) (-7.250,5.127) (-7.200,5.091) (-7.150,5.056) (-7.106,5.025) (-7.050,4.985) (-7.000,4.950) (-6.950,4.914) (-6.900,4.879) (-6.859,4.850) (-6.800,4.808) (-6.753,4.775) (-6.700,4.738) (-6.650,4.702) (-6.600,4.667) (-6.550,4.632) (-6.505,4.600) (-6.450,4.561) (-6.400,4.525) (-6.350,4.490) (-6.300,4.455) (-6.258,4.425) (-6.200,4.384) (-6.152,4.350) (-6.100,4.313) (-6.050,4.278) (-6.000,4.243) (-5.950,4.207) (-5.904,4.175) (-5.850,4.137) (-5.800,4.101) (-5.750,4.066) (-5.700,4.031) (-5.657,4.000) (-5.600,3.960) (-5.551,3.925) (-5.500,3.889) (-5.450,3.854) (-5.400,3.818) (-5.350,3.783) (-5.303,3.750) (-5.250,3.712) (-5.200,3.677) (-5.150,3.642) (-5.100,3.606) (-5.056,3.575) (-5.000,3.536) (-4.950,3.500) (-4.900,3.465) (-4.850,3.429) (-4.808,3.400) (-4.750,3.359) (-4.702,3.325) (-4.650,3.288) (-4.600,3.253) (-4.550,3.217) (-4.500,3.182) (-4.455,3.150) (-4.400,3.111) (-4.350,3.076) (-4.300,3.041) (-4.250,3.005) (-4.207,2.975) (-4.150,2.935) (-4.101,2.900) (-4.050,2.864) (-4.000,2.828) (-3.950,2.793) (-3.900,2.758) (-3.854,2.725) (-3.800,2.687) (-3.750,2.652) (-3.700,2.616) (-3.650,2.581) (-3.606,2.550) (-3.550,2.510) (-3.500,2.475) (-3.450,2.440) (-3.400,2.404) (-3.359,2.375) (-3.300,2.334) (-3.253,2.300) (-3.200,2.263) (-3.150,2.227) (-3.100,2.192) (-3.050,2.157) (-3.005,2.125) (-2.950,2.086) (-2.900,2.051) (-2.850,2.015) (-2.800,1.980) (-2.758,1.950) (-2.700,1.909) (-2.652,1.875) (-2.600,1.839) (-2.550,1.803) (-2.500,1.768) (-2.450,1.732) (-2.404,1.700) (-2.350,1.662) (-2.300,1.626) (-2.250,1.591) (-2.200,1.556) (-2.157,1.525) (-2.100,1.485) (-2.051,1.450) (-2.000,1.414) (-1.950,1.379) (-1.900,1.344) (-1.850,1.308) (-1.803,1.275) (-1.750,1.238) (-1.700,1.202) (-1.650,1.167) (-1.600,1.131) (-1.556,1.100) (-1.500,1.061) (-1.450,1.025) (-1.400,0.990) (-1.350,0.955) (-1.308,0.925) (-1.250,0.884) (-1.202,0.850) (-1.150,0.813) (-1.100,0.778) (-1.050,0.743) (-1.000,0.707) (-0.955,0.675) (-0.900,0.637) (-0.850,0.601) (-0.800,0.566) (-0.750,0.531) (-0.708,0.500) (-0.650,0.460) (-0.601,0.425) (-0.550,0.389) (-0.500,0.354) (-0.450,0.319) (-0.400,0.283) (-0.354,0.250) (-0.300,0.213) (-0.250,0.177) (-0.220,0.155)}
\EQ{(4.306,-10.000) (4.350,-9.950) (4.392,-9.900) (4.425,-9.860) (4.473,-9.800) (4.513,-9.750) (4.550,-9.701) (4.588,-9.650) (4.624,-9.600) (4.660,-9.550) (4.695,-9.500) (4.725,-9.455) (4.761,-9.400) (4.793,-9.350) (4.824,-9.300) (4.855,-9.250) (4.884,-9.200) (4.913,-9.150) (4.941,-9.100) (4.969,-9.050) (4.995,-9.000) (5.021,-8.950) (5.046,-8.900) (5.071,-8.850) (5.095,-8.800) (5.118,-8.750) (5.140,-8.700) (5.162,-8.650) (5.183,-8.600) (5.203,-8.550) (5.223,-8.500) (5.242,-8.450) (5.261,-8.400) (5.279,-8.350) (5.296,-8.300) (5.313,-8.250) (5.329,-8.200) (5.344,-8.150) (5.359,-8.100) (5.373,-8.050) (5.387,-8.000) (5.400,-7.951) (5.413,-7.900) (5.425,-7.850) (5.436,-7.800) (5.447,-7.750) (5.457,-7.700) (5.467,-7.650) (5.476,-7.600) (5.484,-7.550) (5.492,-7.500) (5.500,-7.450) (5.507,-7.400) (5.513,-7.350) (5.519,-7.300) (5.524,-7.250) (5.529,-7.200) (5.533,-7.150) (5.537,-7.100) (5.540,-7.050) (5.542,-7.000) (5.544,-6.950) (5.545,-6.900) (5.546,-6.850) (5.547,-6.800) (5.546,-6.750) (5.546,-6.700) (5.544,-6.650) (5.543,-6.600) (5.540,-6.550) (5.537,-6.500) (5.534,-6.450) (5.530,-6.400) (5.525,-6.350) (5.520,-6.300) (5.515,-6.250) (5.509,-6.200) (5.502,-6.150) (5.495,-6.100) (5.487,-6.050) (5.478,-6.000) (5.469,-5.950) (5.460,-5.900) (5.450,-5.850) (5.439,-5.800) (5.428,-5.750) (5.416,-5.700) (5.404,-5.650) (5.391,-5.600) (5.378,-5.550) (5.363,-5.500) (5.349,-5.450) (5.334,-5.400) (5.318,-5.350) (5.301,-5.300) (5.284,-5.250) (5.267,-5.200) (5.249,-5.150) (5.230,-5.100) (5.210,-5.050) (5.190,-5.000) (5.170,-4.950) (5.150,-4.903) (5.127,-4.850) (5.104,-4.800) (5.081,-4.750) (5.057,-4.700) (5.033,-4.650) (5.008,-4.600) (4.982,-4.550) (4.956,-4.500) (4.929,-4.450) (4.901,-4.400) (4.875,-4.354) (4.844,-4.300) (4.814,-4.250) (4.784,-4.200) (4.753,-4.150) (4.725,-4.107) (4.688,-4.050) (4.655,-4.000) (4.625,-3.955) (4.587,-3.900) (4.551,-3.850) (4.516,-3.800) (4.479,-3.750) (4.441,-3.700) (4.403,-3.650) (4.364,-3.600) (4.325,-3.551) (4.284,-3.500) (4.250,-3.458) (4.201,-3.400) (4.159,-3.350) (4.125,-3.311) (4.075,-3.254) (4.026,-3.200) (3.981,-3.150) (3.935,-3.100) (3.900,-3.063) (3.864,-3.025) (3.825,-2.985) (3.775,-2.934) (3.725,-2.883) (3.675,-2.834) (3.625,-2.785) (3.575,-2.736) (3.537,-2.700) (3.500,-2.665) (3.457,-2.625) (3.403,-2.575) (3.350,-2.526) (3.300,-2.481) (3.250,-2.437) (3.209,-2.400) (3.151,-2.350) (3.100,-2.306) (3.050,-2.263) (3.006,-2.225) (2.950,-2.178) (2.900,-2.137) (2.856,-2.100) (2.800,-2.054) (2.750,-2.014) (2.702,-1.975) (2.650,-1.933) (2.609,-1.900) (2.550,-1.854) (2.500,-1.814) (2.450,-1.775) (2.400,-1.737) (2.353,-1.700) (2.300,-1.660) (2.255,-1.625) (2.200,-1.583) (2.156,-1.550) (2.100,-1.508) (2.056,-1.475) (2.000,-1.433) (1.956,-1.400) (1.900,-1.359) (1.854,-1.325) (1.800,-1.285) (1.752,-1.250) (1.700,-1.212) (1.650,-1.175) (1.600,-1.139) (1.550,-1.103) (1.500,-1.067) (1.450,-1.031) (1.408,-1.000) (1.350,-0.959) (1.304,-0.925) (1.250,-0.887) (1.200,-0.851) (1.150,-0.815) (1.100,-0.780) (1.059,-0.750) (1.000,-0.708) (0.953,-0.675) (0.900,-0.637) (0.850,-0.602) (0.800,-0.566) (0.750,-0.531) (0.707,-0.500) (0.650,-0.460) (0.601,-0.425) (0.550,-0.389) (0.500,-0.354) (0.450,-0.319) (0.400,-0.283) (0.354,-0.250) (0.300,-0.213) (0.250,-0.177) (0.206,-0.144)}
\EQ{(-0.220,-0.155) (-0.275,-0.195) (-0.319,-0.225) (-0.375,-0.266) (-0.424,-0.300) (-0.475,-0.336) (-0.525,-0.372) (-0.566,-0.400) (-0.625,-0.442) (-0.672,-0.475) (-0.725,-0.513) (-0.775,-0.549) (-0.825,-0.584) (-0.875,-0.620) (-0.918,-0.650) (-0.975,-0.691) (-1.024,-0.725) (-1.075,-0.762) (-1.125,-0.797) (-1.175,-0.833) (-1.225,-0.869) (-1.269,-0.900) (-1.325,-0.941) (-1.373,-0.975) (-1.425,-1.013) (-1.475,-1.049) (-1.525,-1.085) (-1.575,-1.121) (-1.625,-1.157) (-1.675,-1.194) (-1.718,-1.225) (-1.775,-1.267) (-1.820,-1.300) (-1.875,-1.340) (-1.922,-1.375) (-1.975,-1.415) (-2.023,-1.450) (-2.075,-1.489) (-2.123,-1.525) (-2.175,-1.564) (-2.222,-1.600) (-2.275,-1.640) (-2.320,-1.675) (-2.375,-1.717) (-2.418,-1.750) (-2.475,-1.795) (-2.525,-1.834) (-2.575,-1.873) (-2.625,-1.913) (-2.671,-1.950) (-2.725,-1.993) (-2.764,-2.025) (-2.825,-2.075) (-2.875,-2.116) (-2.916,-2.150) (-2.975,-2.199) (-3.025,-2.241) (-3.065,-2.275) (-3.123,-2.325) (-3.175,-2.371) (-3.225,-2.415) (-3.265,-2.450) (-3.321,-2.500) (-3.375,-2.549) (-3.425,-2.595) (-3.475,-2.642) (-3.525,-2.689) (-3.563,-2.725) (-3.600,-2.760) (-3.641,-2.800) (-3.692,-2.850) (-3.742,-2.900) (-3.791,-2.950) (-3.840,-3.000) (-3.875,-3.037) (-3.911,-3.075) (-3.950,-3.117) (-4.000,-3.171) (-4.049,-3.225) (-4.093,-3.275) (-4.137,-3.325) (-4.175,-3.369) (-4.222,-3.425) (-4.264,-3.475) (-4.300,-3.519) (-4.345,-3.575) (-4.384,-3.625) (-4.422,-3.675) (-4.460,-3.725) (-4.497,-3.775) (-4.534,-3.825) (-4.569,-3.875) (-4.600,-3.919) (-4.638,-3.975) (-4.672,-4.025) (-4.700,-4.068) (-4.737,-4.125) (-4.768,-4.175) (-4.799,-4.225) (-4.825,-4.269) (-4.858,-4.325) (-4.887,-4.375) (-4.915,-4.425) (-4.942,-4.475) (-4.969,-4.525) (-4.995,-4.575) (-5.020,-4.625) (-5.045,-4.675) (-5.069,-4.725) (-5.093,-4.775) (-5.116,-4.825) (-5.138,-4.875) (-5.159,-4.925) (-5.180,-4.975) (-5.200,-5.024) (-5.220,-5.075) (-5.239,-5.125) (-5.258,-5.175) (-5.275,-5.223) (-5.293,-5.275) (-5.310,-5.325) (-5.325,-5.373) (-5.341,-5.425) (-5.356,-5.475) (-5.371,-5.525) (-5.384,-5.575) (-5.398,-5.625) (-5.410,-5.675) (-5.422,-5.725) (-5.434,-5.775) (-5.445,-5.825) (-5.455,-5.875) (-5.465,-5.925) (-5.474,-5.975) (-5.483,-6.025) (-5.491,-6.075) (-5.498,-6.125) (-5.505,-6.175) (-5.512,-6.225) (-5.518,-6.275) (-5.523,-6.325) (-5.528,-6.375) (-5.532,-6.425) (-5.536,-6.475) (-5.539,-6.525) (-5.542,-6.575) (-5.544,-6.625) (-5.545,-6.675) (-5.546,-6.725) (-5.547,-6.775) (-5.547,-6.825) (-5.546,-6.875) (-5.545,-6.925) (-5.543,-6.975) (-5.541,-7.025) (-5.538,-7.075) (-5.535,-7.125) (-5.531,-7.175) (-5.526,-7.225) (-5.522,-7.275) (-5.516,-7.325) (-5.510,-7.375) (-5.503,-7.425) (-5.496,-7.475) (-5.489,-7.525) (-5.480,-7.575) (-5.471,-7.625) (-5.462,-7.675) (-5.452,-7.725) (-5.442,-7.775) (-5.430,-7.825) (-5.419,-7.875) (-5.407,-7.925) (-5.394,-7.975) (-5.380,-8.025) (-5.366,-8.075) (-5.352,-8.125) (-5.337,-8.175) (-5.321,-8.225) (-5.305,-8.275) (-5.288,-8.325) (-5.270,-8.375) (-5.252,-8.425) (-5.233,-8.475) (-5.213,-8.525) (-5.193,-8.575) (-5.173,-8.625) (-5.151,-8.675) (-5.129,-8.725) (-5.106,-8.775) (-5.083,-8.825) (-5.059,-8.875) (-5.034,-8.925) (-5.008,-8.975) (-4.982,-9.025) (-4.955,-9.075) (-4.927,-9.125) (-4.900,-9.173) (-4.870,-9.225) (-4.840,-9.275) (-4.809,-9.325) (-4.777,-9.375) (-4.750,-9.417) (-4.711,-9.475) (-4.677,-9.525) (-4.642,-9.575) (-4.606,-9.625) (-4.575,-9.668) (-4.532,-9.725) (-4.500,-9.766) (-4.453,-9.825) (-4.413,-9.875) (-4.375,-9.920) (-4.328,-9.975) (-4.306,-10.000)}
\EQ{(0.206,0.144) (0.248,0.175) (0.300,0.213) (0.350,0.248) (0.400,0.283) (0.450,0.319) (0.495,0.350) (0.550,0.389) (0.600,0.424) (0.650,0.460) (0.700,0.495) (0.742,0.525) (0.800,0.566) (0.848,0.600) (0.900,0.637) (0.950,0.673) (1.000,0.708) (1.050,0.744) (1.094,0.775) (1.150,0.815) (1.199,0.850) (1.250,0.887) (1.300,0.923) (1.350,0.959) (1.400,0.995) (1.443,1.025) (1.500,1.067) (1.546,1.100) (1.600,1.139) (1.650,1.175) (1.700,1.212) (1.750,1.248) (1.800,1.285) (1.850,1.322) (1.900,1.359) (1.950,1.396) (2.000,1.433) (2.050,1.470) (2.100,1.508) (2.150,1.546) (2.200,1.583) (2.250,1.621) (2.300,1.660) (2.350,1.698) (2.400,1.737) (2.450,1.775) (2.500,1.814) (2.546,1.850) (2.600,1.893) (2.640,1.925) (2.700,1.973) (2.750,2.014) (2.795,2.050) (2.850,2.095) (2.900,2.137) (2.946,2.175) (3.000,2.220) (3.050,2.263) (3.094,2.300) (3.150,2.349) (3.200,2.393) (3.250,2.437) (3.293,2.475) (3.349,2.525) (3.400,2.572) (3.450,2.618) (3.500,2.665) (3.537,2.700) (3.575,2.736) (3.615,2.775) (3.666,2.825) (3.717,2.875) (3.767,2.925) (3.815,2.975) (3.864,3.025) (3.900,3.063) (3.935,3.100) (3.975,3.144) (4.025,3.198) (4.071,3.250) (4.115,3.300) (4.150,3.340) (4.200,3.399) (4.243,3.450) (4.275,3.489) (4.325,3.550) (4.364,3.600) (4.400,3.646) (4.441,3.700) (4.475,3.745) (4.516,3.800) (4.550,3.848) (4.587,3.900) (4.621,3.950) (4.650,3.992) (4.688,4.050) (4.721,4.100) (4.750,4.146) (4.784,4.200) (4.814,4.250) (4.844,4.300) (4.873,4.350) (4.900,4.398) (4.925,4.443) (4.950,4.489) (4.975,4.536) (5.000,4.585) (5.025,4.634) (5.050,4.685) (5.075,4.737) (5.100,4.791) (5.125,4.846) (5.149,4.900) (5.170,4.950) (5.190,5.000) (5.210,5.050) (5.230,5.100) (5.249,5.150) (5.267,5.200) (5.284,5.250) (5.301,5.300) (5.318,5.350) (5.334,5.400) (5.349,5.450) (5.363,5.500) (5.378,5.550) (5.391,5.600) (5.404,5.650) (5.416,5.700) (5.428,5.750) (5.439,5.800) (5.450,5.850) (5.460,5.900) (5.469,5.950) (5.478,6.000) (5.487,6.050) (5.495,6.100) (5.502,6.150) (5.509,6.200) (5.515,6.250) (5.520,6.300) (5.525,6.350) (5.530,6.400) (5.534,6.450) (5.537,6.500) (5.540,6.550) (5.543,6.600) (5.544,6.650) (5.546,6.700) (5.546,6.750) (5.547,6.800) (5.546,6.850) (5.545,6.900) (5.544,6.950) (5.542,7.000) (5.540,7.050) (5.537,7.100) (5.533,7.150) (5.529,7.200) (5.524,7.250) (5.519,7.300) (5.513,7.350) (5.507,7.400) (5.500,7.450) (5.492,7.500) (5.484,7.550) (5.476,7.600) (5.467,7.650) (5.457,7.700) (5.447,7.750) (5.436,7.800) (5.425,7.850) (5.413,7.900) (5.400,7.950) (5.387,8.000) (5.373,8.050) (5.359,8.100) (5.344,8.150) (5.329,8.200) (5.313,8.250) (5.296,8.300) (5.279,8.350) (5.261,8.400) (5.242,8.450) (5.223,8.500) (5.203,8.550) (5.183,8.600) (5.162,8.650) (5.140,8.700) (5.118,8.750) (5.095,8.800) (5.071,8.850) (5.046,8.900) (5.021,8.950) (4.995,9.000) (4.969,9.050) (4.941,9.100) (4.913,9.150) (4.884,9.200) (4.855,9.250) (4.825,9.299) (4.793,9.350) (4.761,9.400) (4.728,9.450) (4.700,9.492) (4.660,9.550) (4.625,9.599) (4.588,9.650) (4.551,9.700) (4.513,9.750) (4.475,9.798) (4.433,9.850) (4.400,9.890) (4.350,9.949) (4.306,10.000)}
\EQ{(-4.306,10.000) (-4.350,9.950) (-4.392,9.900) (-4.425,9.860) (-4.473,9.800) (-4.513,9.750) (-4.550,9.701) (-4.588,9.650) (-4.624,9.600) (-4.660,9.550) (-4.695,9.500) (-4.725,9.455) (-4.761,9.400) (-4.793,9.350) (-4.824,9.300) (-4.855,9.250) (-4.884,9.200) (-4.913,9.150) (-4.941,9.100) (-4.969,9.050) (-4.995,9.000) (-5.021,8.950) (-5.046,8.900) (-5.071,8.850) (-5.095,8.800) (-5.118,8.750) (-5.140,8.700) (-5.162,8.650) (-5.183,8.600) (-5.203,8.550) (-5.223,8.500) (-5.242,8.450) (-5.261,8.400) (-5.279,8.350) (-5.296,8.300) (-5.313,8.250) (-5.329,8.200) (-5.344,8.150) (-5.359,8.100) (-5.373,8.050) (-5.387,8.000) (-5.400,7.951) (-5.413,7.900) (-5.425,7.850) (-5.436,7.800) (-5.447,7.750) (-5.457,7.700) (-5.467,7.650) (-5.476,7.600) (-5.484,7.550) (-5.492,7.500) (-5.500,7.450) (-5.507,7.400) (-5.513,7.350) (-5.519,7.300) (-5.524,7.250) (-5.529,7.200) (-5.533,7.150) (-5.537,7.100) (-5.540,7.050) (-5.542,7.000) (-5.544,6.950) (-5.545,6.900) (-5.546,6.850) (-5.547,6.800) (-5.546,6.750) (-5.546,6.700) (-5.544,6.650) (-5.543,6.600) (-5.540,6.550) (-5.537,6.500) (-5.534,6.450) (-5.530,6.400) (-5.525,6.350) (-5.520,6.300) (-5.515,6.250) (-5.509,6.200) (-5.502,6.150) (-5.495,6.100) (-5.487,6.050) (-5.478,6.000) (-5.469,5.950) (-5.460,5.900) (-5.450,5.850) (-5.439,5.800) (-5.428,5.750) (-5.416,5.700) (-5.404,5.650) (-5.391,5.600) (-5.378,5.550) (-5.363,5.500) (-5.349,5.450) (-5.334,5.400) (-5.318,5.350) (-5.301,5.300) (-5.284,5.250) (-5.267,5.200) (-5.249,5.150) (-5.230,5.100) (-5.210,5.050) (-5.190,5.000) (-5.170,4.950) (-5.150,4.903) (-5.127,4.850) (-5.104,4.800) (-5.081,4.750) (-5.057,4.700) (-5.033,4.650) (-5.008,4.600) (-4.982,4.550) (-4.956,4.500) (-4.929,4.450) (-4.901,4.400) (-4.875,4.354) (-4.844,4.300) (-4.814,4.250) (-4.784,4.200) (-4.753,4.150) (-4.725,4.107) (-4.688,4.050) (-4.655,4.000) (-4.625,3.955) (-4.587,3.900) (-4.551,3.850) (-4.516,3.800) (-4.479,3.750) (-4.441,3.700) (-4.403,3.650) (-4.364,3.600) (-4.325,3.551) (-4.284,3.500) (-4.250,3.458) (-4.201,3.400) (-4.159,3.350) (-4.125,3.311) (-4.075,3.254) (-4.026,3.200) (-3.981,3.150) (-3.935,3.100) (-3.900,3.063) (-3.864,3.025) (-3.825,2.985) (-3.775,2.934) (-3.725,2.883) (-3.675,2.834) (-3.625,2.785) (-3.575,2.736) (-3.537,2.700) (-3.500,2.665) (-3.457,2.625) (-3.403,2.575) (-3.350,2.526) (-3.300,2.481) (-3.250,2.437) (-3.209,2.400) (-3.151,2.350) (-3.100,2.306) (-3.050,2.263) (-3.006,2.225) (-2.950,2.178) (-2.900,2.137) (-2.856,2.100) (-2.800,2.054) (-2.750,2.014) (-2.702,1.975) (-2.650,1.933) (-2.609,1.900) (-2.550,1.854) (-2.500,1.814) (-2.450,1.775) (-2.400,1.737) (-2.353,1.700) (-2.300,1.660) (-2.255,1.625) (-2.200,1.583) (-2.156,1.550) (-2.100,1.508) (-2.056,1.475) (-2.000,1.433) (-1.956,1.400) (-1.900,1.359) (-1.854,1.325) (-1.800,1.285) (-1.752,1.250) (-1.700,1.212) (-1.650,1.175) (-1.600,1.139) (-1.550,1.103) (-1.500,1.067) (-1.450,1.031) (-1.408,1.000) (-1.350,0.959) (-1.304,0.925) (-1.250,0.887) (-1.200,0.851) (-1.150,0.815) (-1.100,0.780) (-1.059,0.750) (-1.000,0.708) (-0.953,0.675) (-0.900,0.637) (-0.850,0.602) (-0.800,0.566) (-0.750,0.531) (-0.707,0.500) (-0.650,0.460) (-0.601,0.425) (-0.550,0.389) (-0.500,0.354) (-0.450,0.319) (-0.400,0.283) (-0.354,0.250) (-0.300,0.213) (-0.250,0.177) (-0.220,0.155)}
\EQ{(-0.220,-0.155) (-0.275,-0.195) (-0.319,-0.225) (-0.375,-0.266) (-0.424,-0.300) (-0.475,-0.337) (-0.525,-0.372) (-0.575,-0.407) (-0.625,-0.443) (-0.671,-0.475) (-0.725,-0.514) (-0.775,-0.549) (-0.825,-0.585) (-0.875,-0.620) (-0.917,-0.650) (-0.975,-0.692) (-1.022,-0.725) (-1.075,-0.763) (-1.125,-0.799) (-1.175,-0.835) (-1.225,-0.872) (-1.275,-0.908) (-1.325,-0.944) (-1.368,-0.975) (-1.425,-1.017) (-1.470,-1.050) (-1.525,-1.091) (-1.571,-1.125) (-1.625,-1.166) (-1.671,-1.200) (-1.725,-1.241) (-1.770,-1.275) (-1.825,-1.317) (-1.869,-1.350) (-1.925,-1.394) (-1.965,-1.425) (-2.025,-1.472) (-2.075,-1.512) (-2.123,-1.550) (-2.175,-1.592) (-2.216,-1.625) (-2.275,-1.674) (-2.325,-1.715) (-2.367,-1.750) (-2.425,-1.799) (-2.475,-1.843) (-2.525,-1.886) (-2.570,-1.925) (-2.625,-1.975) (-2.675,-2.020) (-2.725,-2.066) (-2.762,-2.100) (-2.800,-2.136) (-2.841,-2.175) (-2.893,-2.225) (-2.943,-2.275) (-2.993,-2.325) (-3.042,-2.375) (-3.090,-2.425) (-3.125,-2.463) (-3.160,-2.500) (-3.200,-2.545) (-3.249,-2.600) (-3.293,-2.650) (-3.336,-2.700) (-3.375,-2.747) (-3.418,-2.800) (-3.450,-2.840) (-3.497,-2.900) (-3.535,-2.950) (-3.572,-3.000) (-3.608,-3.050) (-3.644,-3.100) (-3.675,-3.146) (-3.712,-3.200) (-3.744,-3.250) (-3.775,-3.299) (-3.807,-3.350) (-3.837,-3.400) (-3.866,-3.450) (-3.894,-3.500) (-3.922,-3.550) (-3.948,-3.600) (-3.974,-3.650) (-3.999,-3.700) (-4.023,-3.750) (-4.047,-3.800) (-4.069,-3.850) (-4.091,-3.900) (-4.112,-3.950) (-4.132,-4.000) (-4.152,-4.050) (-4.170,-4.100) (-4.188,-4.150) (-4.206,-4.200) (-4.222,-4.250) (-4.238,-4.300) (-4.253,-4.350) (-4.267,-4.400) (-4.280,-4.450) (-4.293,-4.500) (-4.305,-4.550) (-4.316,-4.600) (-4.327,-4.650) (-4.337,-4.700) (-4.346,-4.750) (-4.354,-4.800) (-4.362,-4.850) (-4.369,-4.900) (-4.376,-4.950) (-4.381,-5.000) (-4.386,-5.050) (-4.391,-5.100) (-4.394,-5.150) (-4.397,-5.200) (-4.400,-5.250) (-4.401,-5.300) (-4.402,-5.350) (-4.402,-5.400) (-4.402,-5.450) (-4.401,-5.500) (-4.399,-5.550) (-4.396,-5.600) (-4.393,-5.650) (-4.389,-5.700) (-4.385,-5.750) (-4.380,-5.800) (-4.374,-5.850) (-4.367,-5.900) (-4.360,-5.950) (-4.352,-6.000) (-4.343,-6.050) (-4.334,-6.100) (-4.323,-6.150) (-4.313,-6.200) (-4.301,-6.250) (-4.289,-6.300) (-4.276,-6.350) (-4.262,-6.400) (-4.247,-6.450) (-4.232,-6.500) (-4.216,-6.550) (-4.200,-6.598) (-4.182,-6.650) (-4.163,-6.700) (-4.144,-6.750) (-4.125,-6.798) (-4.104,-6.850) (-4.082,-6.900) (-4.060,-6.950) (-4.036,-7.000) (-4.012,-7.050) (-3.987,-7.100) (-3.961,-7.150) (-3.934,-7.200) (-3.907,-7.250) (-3.878,-7.300) (-3.850,-7.347) (-3.818,-7.400) (-3.786,-7.450) (-3.753,-7.500) (-3.725,-7.542) (-3.684,-7.600) (-3.650,-7.648) (-3.611,-7.700) (-3.575,-7.747) (-3.533,-7.800) (-3.500,-7.841) (-3.450,-7.900) (-3.406,-7.950) (-3.361,-8.000) (-3.325,-8.039) (-3.275,-8.092) (-3.225,-8.142) (-3.175,-8.191) (-3.125,-8.238) (-3.085,-8.275) (-3.029,-8.325) (-2.975,-8.372) (-2.925,-8.414) (-2.880,-8.450) (-2.825,-8.494) (-2.784,-8.525) (-2.725,-8.569) (-2.682,-8.600) (-2.625,-8.640) (-2.575,-8.674) (-2.525,-8.707) (-2.475,-8.739) (-2.425,-8.770) (-2.375,-8.800) (-2.325,-8.829) (-2.275,-8.858) (-2.225,-8.886) (-2.175,-8.912) (-2.125,-8.938) (-2.075,-8.964) (-2.025,-8.988) (-1.975,-9.012) (-1.925,-9.035) (-1.875,-9.057) (-1.825,-9.078) (-1.775,-9.099) (-1.725,-9.119) (-1.675,-9.139) (-1.625,-9.158) (-1.577,-9.175) (-1.525,-9.193) (-1.475,-9.210) (-1.425,-9.226) (-1.375,-9.242) (-1.325,-9.257) (-1.275,-9.272) (-1.225,-9.285) (-1.175,-9.298) (-1.125,-9.311) (-1.075,-9.323) (-1.025,-9.335) (-0.975,-9.345) (-0.925,-9.356) (-0.875,-9.365) (-0.825,-9.375) (-0.775,-9.383) (-0.725,-9.391) (-0.675,-9.399) (-0.625,-9.406) (-0.575,-9.412) (-0.525,-9.418) (-0.475,-9.423) (-0.425,-9.428) (-0.375,-9.432) (-0.325,-9.436) (-0.275,-9.439) (-0.225,-9.442) (-0.175,-9.444) (-0.125,-9.446) (-0.075,-9.447) (-0.025,-9.447) (0.050,-9.447) (0.100,-9.446) (0.150,-9.445) (0.200,-9.443) (0.250,-9.441) (0.300,-9.438) (0.350,-9.434) (0.400,-9.430) (0.450,-9.426) (0.500,-9.421) (0.550,-9.415) (0.600,-9.409) (0.650,-9.402) (0.700,-9.395) (0.750,-9.387) (0.800,-9.379) (0.850,-9.370) (0.900,-9.361) (0.950,-9.351) (1.000,-9.340) (1.050,-9.329) (1.100,-9.317) (1.150,-9.305) (1.200,-9.292) (1.250,-9.278) (1.300,-9.264) (1.349,-9.250) (1.400,-9.234) (1.450,-9.218) (1.500,-9.202) (1.550,-9.185) (1.600,-9.167) (1.650,-9.148) (1.700,-9.129) (1.750,-9.109) (1.800,-9.089) (1.850,-9.068) (1.900,-9.046) (1.946,-9.025) (2.000,-9.000) (2.050,-8.976) (2.100,-8.951) (2.150,-8.925) (2.198,-8.900) (2.244,-8.875) (2.289,-8.850) (2.333,-8.825) (2.400,-8.785) (2.450,-8.755) (2.497,-8.725) (2.550,-8.691) (2.600,-8.657) (2.647,-8.625) (2.700,-8.587) (2.750,-8.551) (2.800,-8.513) (2.849,-8.475) (2.900,-8.434) (2.942,-8.400) (3.000,-8.351) (3.050,-8.307) (3.100,-8.262) (3.139,-8.225) (3.192,-8.175) (3.242,-8.125) (3.291,-8.075) (3.338,-8.025) (3.375,-7.985) (3.425,-7.929) (3.471,-7.875) (3.513,-7.825) (3.550,-7.779) (3.592,-7.725) (3.625,-7.682) (3.667,-7.625) (3.700,-7.578) (3.736,-7.525) (3.770,-7.475) (3.800,-7.428) (3.833,-7.375) (3.863,-7.325) (3.892,-7.275) (3.921,-7.225) (3.948,-7.175) (3.974,-7.125) (4.000,-7.075) (4.024,-7.025) (4.048,-6.975) (4.071,-6.925) (4.093,-6.875) (4.114,-6.825) (4.134,-6.775) (4.154,-6.725) (4.173,-6.675) (4.191,-6.625) (4.208,-6.575) (4.224,-6.525) (4.240,-6.475) (4.255,-6.425) (4.269,-6.375) (4.282,-6.325) (4.295,-6.275) (4.307,-6.225) (4.318,-6.175) (4.329,-6.125) (4.338,-6.075) (4.347,-6.025) (4.356,-5.975) (4.364,-5.925) (4.370,-5.875) (4.377,-5.825) (4.382,-5.775) (4.387,-5.725) (4.391,-5.675) (4.395,-5.625) (4.398,-5.575) (4.400,-5.525) (4.401,-5.475) (4.402,-5.425) (4.402,-5.375) (4.402,-5.325) (4.401,-5.275) (4.399,-5.225) (4.396,-5.175) (4.393,-5.125) (4.389,-5.075) (4.384,-5.025) (4.379,-4.975) (4.373,-4.925) (4.366,-4.875) (4.358,-4.825) (4.350,-4.775) (4.342,-4.725) (4.332,-4.675) (4.322,-4.625) (4.311,-4.575) (4.299,-4.525) (4.287,-4.475) (4.274,-4.425) (4.260,-4.375) (4.245,-4.325) (4.230,-4.275) (4.214,-4.225) (4.197,-4.175) (4.180,-4.125) (4.161,-4.075) (4.142,-4.025) (4.122,-3.975) (4.102,-3.925) (4.080,-3.875) (4.058,-3.825) (4.035,-3.775) (4.011,-3.725) (3.987,-3.675) (3.961,-3.625) (3.935,-3.575) (3.908,-3.525) (3.880,-3.475) (3.852,-3.425) (3.825,-3.380) (3.792,-3.325) (3.760,-3.275) (3.728,-3.225) (3.700,-3.183) (3.661,-3.125) (3.626,-3.075) (3.590,-3.025) (3.554,-2.975) (3.516,-2.925) (3.478,-2.875) (3.438,-2.825) (3.400,-2.778) (3.357,-2.725) (3.315,-2.675) (3.275,-2.629) (3.227,-2.575) (3.182,-2.525) (3.136,-2.475) (3.100,-2.436) (3.050,-2.384) (3.000,-2.332) (2.950,-2.282) (2.900,-2.232) (2.850,-2.184) (2.800,-2.136) (2.762,-2.100) (2.725,-2.066) (2.681,-2.025) (2.626,-1.975) (2.575,-1.930) (2.525,-1.886) (2.484,-1.850) (2.426,-1.800) (2.375,-1.757) (2.325,-1.715) (2.277,-1.675) (2.225,-1.633) (2.185,-1.600) (2.125,-1.552) (2.075,-1.512) (2.029,-1.475) (1.975,-1.433) (1.933,-1.400) (1.875,-1.355) (1.825,-1.317) (1.775,-1.279) (1.725,-1.241) (1.675,-1.203) (1.625,-1.166) (1.575,-1.128) (1.525,-1.091) (1.475,-1.054) (1.425,-1.017) (1.375,-0.981) (1.333,-0.950) (1.275,-0.908) (1.230,-0.875) (1.175,-0.835) (1.126,-0.800) (1.075,-0.763) (1.025,-0.727) (0.975,-0.692) (0.925,-0.656) (0.882,-0.625) (0.825,-0.585) (0.777,-0.550) (0.725,-0.514) (0.675,-0.478) (0.625,-0.443) (0.575,-0.407) (0.531,-0.375) (0.475,-0.337) (0.425,-0.301) (0.375,-0.266) (0.325,-0.230) (0.284,-0.200) (0.225,-0.160) (0.206,-0.144)}
\EQ{(0.206,0.144) (0.248,0.175) (0.300,0.213) (0.350,0.248) (0.400,0.283) (0.450,0.319) (0.495,0.350) (0.550,0.390) (0.600,0.425) (0.650,0.460) (0.700,0.496) (0.742,0.525) (0.800,0.567) (0.847,0.600) (0.900,0.638) (0.950,0.674) (1.000,0.710) (1.050,0.745) (1.092,0.775) (1.150,0.817) (1.196,0.850) (1.250,0.890) (1.299,0.925) (1.350,0.963) (1.400,0.999) (1.450,1.036) (1.500,1.073) (1.550,1.110) (1.600,1.147) (1.650,1.184) (1.700,1.222) (1.750,1.260) (1.800,1.298) (1.850,1.336) (1.900,1.374) (1.950,1.413) (1.997,1.450) (2.050,1.492) (2.092,1.525) (2.150,1.572) (2.200,1.612) (2.246,1.650) (2.300,1.694) (2.350,1.736) (2.396,1.775) (2.450,1.821) (2.500,1.864) (2.541,1.900) (2.598,1.950) (2.650,1.997) (2.700,2.043) (2.750,2.089) (2.789,2.125) (2.825,2.160) (2.867,2.200) (2.918,2.250) (2.968,2.300) (3.017,2.350) (3.066,2.400) (3.113,2.450) (3.150,2.490) (3.200,2.545) (3.249,2.600) (3.293,2.650) (3.336,2.700) (3.375,2.747) (3.418,2.800) (3.450,2.840) (3.497,2.900) (3.535,2.950) (3.572,3.000) (3.608,3.050) (3.644,3.100) (3.675,3.146) (3.712,3.200) (3.744,3.250) (3.775,3.299) (3.807,3.350) (3.837,3.400) (3.866,3.450) (3.894,3.500) (3.922,3.550) (3.948,3.600) (3.974,3.650) (3.999,3.700) (4.023,3.750) (4.047,3.800) (4.069,3.850) (4.091,3.900) (4.112,3.950) (4.132,4.000) (4.152,4.050) (4.170,4.100) (4.188,4.150) (4.206,4.200) (4.222,4.250) (4.238,4.300) (4.253,4.350) (4.267,4.400) (4.280,4.450) (4.293,4.500) (4.305,4.550) (4.316,4.600) (4.327,4.650) (4.337,4.700) (4.346,4.750) (4.354,4.800) (4.362,4.850) (4.369,4.900) (4.376,4.950) (4.381,5.000) (4.386,5.050) (4.391,5.100) (4.394,5.150) (4.397,5.200) (4.400,5.250) (4.401,5.300) (4.402,5.350) (4.402,5.400) (4.402,5.450) (4.401,5.500) (4.399,5.550) (4.396,5.600) (4.393,5.650) (4.389,5.700) (4.385,5.750) (4.380,5.800) (4.374,5.850) (4.367,5.900) (4.360,5.950) (4.352,6.000) (4.343,6.050) (4.334,6.100) (4.323,6.150) (4.313,6.200) (4.301,6.250) (4.289,6.300) (4.276,6.350) (4.262,6.400) (4.247,6.450) (4.232,6.500) (4.216,6.550) (4.200,6.598) (4.182,6.650) (4.163,6.700) (4.144,6.750) (4.125,6.798) (4.104,6.850) (4.082,6.900) (4.060,6.950) (4.036,7.000) (4.012,7.050) (3.987,7.100) (3.961,7.150) (3.934,7.200) (3.907,7.250) (3.878,7.300) (3.850,7.347) (3.818,7.400) (3.786,7.450) (3.753,7.500) (3.725,7.542) (3.684,7.600) (3.650,7.648) (3.611,7.700) (3.575,7.747) (3.533,7.800) (3.500,7.841) (3.450,7.900) (3.406,7.950) (3.361,8.000) (3.325,8.039) (3.275,8.092) (3.225,8.142) (3.175,8.191) (3.125,8.238) (3.085,8.275) (3.029,8.325) (2.975,8.372) (2.925,8.414) (2.880,8.450) (2.825,8.494) (2.784,8.525) (2.725,8.569) (2.682,8.600) (2.625,8.640) (2.575,8.674) (2.525,8.707) (2.475,8.739) (2.425,8.770) (2.375,8.800) (2.325,8.829) (2.275,8.858) (2.225,8.886) (2.175,8.912) (2.125,8.938) (2.075,8.964) (2.025,8.988) (1.975,9.012) (1.925,9.035) (1.875,9.057) (1.825,9.078) (1.775,9.099) (1.725,9.119) (1.675,9.139) (1.625,9.158) (1.577,9.175) (1.525,9.193) (1.475,9.210) (1.425,9.226) (1.375,9.242) (1.325,9.257) (1.275,9.272) (1.225,9.285) (1.175,9.298) (1.125,9.311) (1.075,9.323) (1.025,9.335) (0.975,9.345) (0.925,9.356) (0.875,9.365) (0.825,9.375) (0.775,9.383) (0.725,9.391) (0.675,9.399) (0.625,9.406) (0.575,9.412) (0.525,9.418) (0.475,9.423) (0.425,9.428) (0.375,9.432) (0.325,9.436) (0.275,9.439) (0.225,9.442) (0.175,9.444) (0.125,9.446) (0.075,9.447) (0.025,9.447) (-0.050,9.447) (-0.100,9.446) (-0.150,9.445) (-0.200,9.443) (-0.250,9.441) (-0.300,9.438) (-0.350,9.434) (-0.400,9.430) (-0.450,9.426) (-0.500,9.421) (-0.550,9.415) (-0.600,9.409) (-0.650,9.402) (-0.700,9.395) (-0.750,9.387) (-0.800,9.379) (-0.850,9.370) (-0.900,9.361) (-0.950,9.351) (-1.000,9.340) (-1.050,9.329) (-1.100,9.317) (-1.150,9.305) (-1.200,9.292) (-1.250,9.278) (-1.300,9.264) (-1.349,9.250) (-1.400,9.234) (-1.450,9.218) (-1.500,9.202) (-1.550,9.185) (-1.600,9.167) (-1.650,9.148) (-1.700,9.129) (-1.750,9.109) (-1.800,9.089) (-1.850,9.068) (-1.900,9.046) (-1.946,9.025) (-2.000,9.000) (-2.050,8.976) (-2.100,8.951) (-2.150,8.925) (-2.198,8.900) (-2.244,8.875) (-2.289,8.850) (-2.333,8.825) (-2.400,8.785) (-2.450,8.755) (-2.497,8.725) (-2.550,8.691) (-2.600,8.657) (-2.647,8.625) (-2.700,8.587) (-2.750,8.551) (-2.800,8.513) (-2.849,8.475) (-2.900,8.434) (-2.942,8.400) (-3.000,8.351) (-3.050,8.307) (-3.100,8.262) (-3.139,8.225) (-3.192,8.175) (-3.242,8.125) (-3.291,8.075) (-3.338,8.025) (-3.375,7.985) (-3.425,7.929) (-3.471,7.875) (-3.513,7.825) (-3.550,7.779) (-3.592,7.725) (-3.625,7.682) (-3.667,7.625) (-3.700,7.578) (-3.736,7.525) (-3.770,7.475) (-3.800,7.428) (-3.833,7.375) (-3.863,7.325) (-3.892,7.275) (-3.921,7.225) (-3.948,7.175) (-3.974,7.125) (-4.000,7.075) (-4.024,7.025) (-4.048,6.975) (-4.071,6.925) (-4.093,6.875) (-4.114,6.825) (-4.134,6.775) (-4.154,6.725) (-4.173,6.675) (-4.191,6.625) (-4.208,6.575) (-4.224,6.525) (-4.240,6.475) (-4.255,6.425) (-4.269,6.375) (-4.282,6.325) (-4.295,6.275) (-4.307,6.225) (-4.318,6.175) (-4.329,6.125) (-4.338,6.075) (-4.347,6.025) (-4.356,5.975) (-4.364,5.925) (-4.370,5.875) (-4.377,5.825) (-4.382,5.775) (-4.387,5.725) (-4.391,5.675) (-4.395,5.625) (-4.398,5.575) (-4.400,5.525) (-4.401,5.475) (-4.402,5.425) (-4.402,5.375) (-4.402,5.325) (-4.401,5.275) (-4.399,5.225) (-4.396,5.175) (-4.393,5.125) (-4.389,5.075) (-4.384,5.025) (-4.379,4.975) (-4.373,4.925) (-4.366,4.875) (-4.358,4.825) (-4.350,4.775) (-4.342,4.725) (-4.332,4.675) (-4.322,4.625) (-4.311,4.575) (-4.299,4.525) (-4.287,4.475) (-4.274,4.425) (-4.260,4.375) (-4.245,4.325) (-4.230,4.275) (-4.214,4.225) (-4.197,4.175) (-4.180,4.125) (-4.161,4.075) (-4.142,4.025) (-4.122,3.975) (-4.102,3.925) (-4.080,3.875) (-4.058,3.825) (-4.035,3.775) (-4.011,3.725) (-3.987,3.675) (-3.961,3.625) (-3.935,3.575) (-3.908,3.525) (-3.880,3.475) (-3.852,3.425) (-3.825,3.380) (-3.792,3.325) (-3.760,3.275) (-3.728,3.225) (-3.700,3.183) (-3.661,3.125) (-3.626,3.075) (-3.590,3.025) (-3.554,2.975) (-3.516,2.925) (-3.478,2.875) (-3.438,2.825) (-3.400,2.778) (-3.357,2.725) (-3.315,2.675) (-3.275,2.629) (-3.227,2.575) (-3.182,2.525) (-3.136,2.475) (-3.100,2.436) (-3.050,2.384) (-3.000,2.332) (-2.950,2.282) (-2.900,2.232) (-2.850,2.184) (-2.800,2.136) (-2.762,2.100) (-2.725,2.066) (-2.681,2.025) (-2.626,1.975) (-2.575,1.930) (-2.525,1.886) (-2.484,1.850) (-2.426,1.800) (-2.375,1.757) (-2.325,1.715) (-2.277,1.675) (-2.225,1.633) (-2.185,1.600) (-2.125,1.552) (-2.075,1.512) (-2.029,1.475) (-1.975,1.433) (-1.933,1.400) (-1.875,1.355) (-1.825,1.317) (-1.775,1.279) (-1.725,1.241) (-1.675,1.203) (-1.625,1.166) (-1.575,1.128) (-1.525,1.091) (-1.475,1.054) (-1.425,1.017) (-1.375,0.981) (-1.333,0.950) (-1.275,0.908) (-1.230,0.875) (-1.175,0.835) (-1.126,0.800) (-1.075,0.763) (-1.025,0.727) (-0.975,0.692) (-0.925,0.656) (-0.882,0.625) (-0.825,0.585) (-0.777,0.550) (-0.725,0.514) (-0.675,0.478) (-0.625,0.443) (-0.575,0.407) (-0.531,0.375) (-0.475,0.337) (-0.425,0.301) (-0.375,0.266) (-0.325,0.230) (-0.284,0.200) (-0.225,0.160) (-0.220,0.155)}
\EQ{(-0.220,-0.155) (-0.275,-0.195) (-0.319,-0.225) (-0.375,-0.266) (-0.424,-0.300) (-0.475,-0.337) (-0.525,-0.372) (-0.575,-0.407) (-0.625,-0.443) (-0.671,-0.475) (-0.725,-0.514) (-0.775,-0.549) (-0.825,-0.585) (-0.875,-0.621) (-0.916,-0.650) (-0.975,-0.693) (-1.020,-0.725) (-1.075,-0.765) (-1.124,-0.800) (-1.175,-0.838) (-1.225,-0.874) (-1.275,-0.911) (-1.325,-0.948) (-1.375,-0.985) (-1.425,-1.022) (-1.475,-1.060) (-1.525,-1.098) (-1.575,-1.136) (-1.625,-1.174) (-1.675,-1.213) (-1.723,-1.250) (-1.775,-1.291) (-1.818,-1.325) (-1.875,-1.370) (-1.925,-1.411) (-1.973,-1.450) (-2.025,-1.493) (-2.075,-1.535) (-2.122,-1.575) (-2.175,-1.621) (-2.225,-1.664) (-2.265,-1.700) (-2.321,-1.750) (-2.375,-1.800) (-2.425,-1.846) (-2.475,-1.894) (-2.525,-1.942) (-2.575,-1.991) (-2.625,-2.042) (-2.675,-2.093) (-2.725,-2.146) (-2.775,-2.200) (-2.820,-2.250) (-2.864,-2.300) (-2.900,-2.342) (-2.949,-2.400) (-2.989,-2.450) (-3.025,-2.495) (-3.068,-2.550) (-3.100,-2.593) (-3.142,-2.650) (-3.175,-2.697) (-3.212,-2.750) (-3.245,-2.800) (-3.275,-2.846) (-3.309,-2.900) (-3.340,-2.950) (-3.369,-3.000) (-3.398,-3.050) (-3.425,-3.099) (-3.450,-3.146) (-3.475,-3.195) (-3.500,-3.245) (-3.525,-3.298) (-3.549,-3.350) (-3.571,-3.400) (-3.592,-3.450) (-3.612,-3.500) (-3.631,-3.550) (-3.650,-3.600) (-3.667,-3.650) (-3.684,-3.700) (-3.700,-3.750) (-3.715,-3.800) (-3.729,-3.850) (-3.742,-3.900) (-3.755,-3.950) (-3.767,-4.000) (-3.777,-4.050) (-3.787,-4.100) (-3.797,-4.150) (-3.805,-4.200) (-3.813,-4.250) (-3.820,-4.300) (-3.826,-4.350) (-3.831,-4.400) (-3.835,-4.450) (-3.839,-4.500) (-3.842,-4.550) (-3.844,-4.600) (-3.845,-4.650) (-3.846,-4.700) (-3.846,-4.750) (-3.845,-4.800) (-3.843,-4.850) (-3.840,-4.900) (-3.837,-4.950) (-3.833,-5.000) (-3.828,-5.050) (-3.822,-5.100) (-3.816,-5.150) (-3.808,-5.200) (-3.800,-5.250) (-3.791,-5.300) (-3.782,-5.350) (-3.771,-5.400) (-3.760,-5.450) (-3.747,-5.500) (-3.734,-5.550) (-3.721,-5.600) (-3.706,-5.650) (-3.690,-5.700) (-3.674,-5.750) (-3.656,-5.800) (-3.638,-5.850) (-3.619,-5.900) (-3.600,-5.948) (-3.578,-6.000) (-3.556,-6.050) (-3.533,-6.100) (-3.509,-6.150) (-3.484,-6.200) (-3.458,-6.250) (-3.431,-6.300) (-3.403,-6.350) (-3.375,-6.399) (-3.344,-6.450) (-3.313,-6.500) (-3.280,-6.550) (-3.250,-6.594) (-3.211,-6.650) (-3.175,-6.699) (-3.137,-6.750) (-3.100,-6.797) (-3.057,-6.850) (-3.015,-6.900) (-2.975,-6.946) (-2.926,-7.000) (-2.879,-7.050) (-2.830,-7.100) (-2.779,-7.150) (-2.727,-7.200) (-2.675,-7.247) (-2.625,-7.291) (-2.585,-7.325) (-2.525,-7.375) (-2.475,-7.414) (-2.429,-7.450) (-2.375,-7.490) (-2.326,-7.525) (-2.275,-7.560) (-2.225,-7.594) (-2.176,-7.625) (-2.125,-7.657) (-2.075,-7.688) (-2.025,-7.717) (-1.975,-7.745) (-1.925,-7.772) (-1.875,-7.799) (-1.825,-7.824) (-1.775,-7.849) (-1.725,-7.873) (-1.675,-7.896) (-1.625,-7.918) (-1.575,-7.939) (-1.525,-7.959) (-1.475,-7.979) (-1.425,-7.998) (-1.375,-8.016) (-1.325,-8.034) (-1.276,-8.050) (-1.225,-8.066) (-1.175,-8.082) (-1.125,-8.096) (-1.075,-8.110) (-1.025,-8.123) (-0.975,-8.136) (-0.925,-8.148) (-0.875,-8.159) (-0.825,-8.170) (-0.775,-8.179) (-0.725,-8.189) (-0.675,-8.197) (-0.625,-8.205) (-0.575,-8.213) (-0.525,-8.219) (-0.475,-8.225) (-0.425,-8.231) (-0.375,-8.236) (-0.325,-8.240) (-0.275,-8.244) (-0.225,-8.247) (-0.175,-8.249) (-0.125,-8.251) (-0.075,-8.252) (-0.025,-8.253) (0.050,-8.253) (0.100,-8.252) (0.150,-8.250) (0.200,-8.248) (0.250,-8.245) (0.300,-8.242) (0.350,-8.238) (0.400,-8.234) (0.450,-8.228) (0.500,-8.223) (0.550,-8.216) (0.600,-8.209) (0.650,-8.201) (0.700,-8.193) (0.750,-8.184) (0.800,-8.175) (0.850,-8.164) (0.900,-8.153) (0.950,-8.142) (1.000,-8.130) (1.050,-8.117) (1.100,-8.103) (1.150,-8.089) (1.200,-8.074) (1.250,-8.058) (1.300,-8.042) (1.350,-8.025) (1.400,-8.007) (1.450,-7.989) (1.500,-7.969) (1.548,-7.950) (1.600,-7.928) (1.650,-7.907) (1.700,-7.884) (1.750,-7.861) (1.800,-7.837) (1.850,-7.812) (1.900,-7.786) (1.950,-7.759) (2.000,-7.731) (2.050,-7.702) (2.096,-7.675) (2.150,-7.642) (2.200,-7.610) (2.250,-7.577) (2.300,-7.543) (2.350,-7.508) (2.395,-7.475) (2.450,-7.434) (2.493,-7.400) (2.550,-7.354) (2.600,-7.313) (2.644,-7.275) (2.700,-7.225) (2.750,-7.178) (2.800,-7.130) (2.850,-7.080) (2.900,-7.028) (2.949,-6.975) (2.993,-6.925) (3.036,-6.875) (3.075,-6.828) (3.117,-6.775) (3.150,-6.733) (3.193,-6.675) (3.225,-6.630) (3.263,-6.575) (3.296,-6.525) (3.325,-6.480) (3.359,-6.425) (3.389,-6.375) (3.418,-6.325) (3.445,-6.275) (3.472,-6.225) (3.497,-6.175) (3.521,-6.125) (3.545,-6.075) (3.567,-6.025) (3.589,-5.975) (3.609,-5.925) (3.629,-5.875) (3.647,-5.825) (3.665,-5.775) (3.682,-5.725) (3.698,-5.675) (3.713,-5.625) (3.728,-5.575) (3.741,-5.525) (3.754,-5.475) (3.765,-5.425) (3.776,-5.375) (3.787,-5.325) (3.796,-5.275) (3.804,-5.225) (3.812,-5.175) (3.819,-5.125) (3.825,-5.075) (3.830,-5.025) (3.835,-4.975) (3.839,-4.925) (3.842,-4.875) (3.844,-4.825) (3.845,-4.775) (3.846,-4.725) (3.846,-4.675) (3.845,-4.625) (3.843,-4.575) (3.841,-4.525) (3.837,-4.475) (3.833,-4.425) (3.828,-4.375) (3.823,-4.325) (3.816,-4.275) (3.809,-4.225) (3.801,-4.175) (3.792,-4.125) (3.783,-4.075) (3.772,-4.025) (3.761,-3.975) (3.749,-3.925) (3.736,-3.875) (3.722,-3.825) (3.708,-3.775) (3.692,-3.725) (3.676,-3.675) (3.659,-3.625) (3.641,-3.575) (3.622,-3.525) (3.602,-3.475) (3.582,-3.425) (3.560,-3.375) (3.538,-3.325) (3.514,-3.275) (3.490,-3.225) (3.465,-3.175) (3.439,-3.125) (3.412,-3.075) (3.384,-3.025) (3.355,-2.975) (3.325,-2.926) (3.294,-2.875) (3.262,-2.825) (3.229,-2.775) (3.200,-2.733) (3.160,-2.675) (3.125,-2.627) (3.087,-2.575) (3.050,-2.527) (3.009,-2.475) (2.975,-2.432) (2.928,-2.375) (2.886,-2.325) (2.850,-2.284) (2.800,-2.228) (2.752,-2.175) (2.705,-2.125) (2.657,-2.075) (2.609,-2.025) (2.559,-1.975) (2.508,-1.925) (2.456,-1.875) (2.403,-1.825) (2.350,-1.777) (2.300,-1.731) (2.250,-1.687) (2.209,-1.650) (2.151,-1.600) (2.100,-1.556) (2.050,-1.514) (2.003,-1.475) (1.950,-1.431) (1.900,-1.391) (1.850,-1.350) (1.800,-1.311) (1.755,-1.275) (1.700,-1.232) (1.659,-1.200) (1.600,-1.155) (1.550,-1.117) (1.500,-1.079) (1.450,-1.041) (1.400,-1.004) (1.350,-0.966) (1.300,-0.929) (1.250,-0.893) (1.200,-0.856) (1.158,-0.825) (1.100,-0.783) (1.055,-0.750) (1.000,-0.711) (0.951,-0.675) (0.900,-0.639) (0.850,-0.603) (0.800,-0.567) (0.750,-0.532) (0.706,-0.500) (0.650,-0.461) (0.600,-0.425) (0.550,-0.390) (0.500,-0.354) (0.450,-0.319) (0.400,-0.284) (0.354,-0.250) (0.300,-0.213) (0.250,-0.177) (0.206,-0.144)}
\EQ{(0.206,0.144) (0.248,0.175) (0.300,0.213) (0.350,0.248) (0.400,0.284) (0.450,0.319) (0.495,0.350) (0.550,0.390) (0.600,0.425) (0.650,0.461) (0.700,0.496) (0.742,0.525) (0.800,0.567) (0.846,0.600) (0.900,0.639) (0.950,0.675) (1.000,0.711) (1.050,0.747) (1.100,0.783) (1.150,0.819) (1.192,0.850) (1.250,0.893) (1.294,0.925) (1.350,0.966) (1.396,1.000) (1.450,1.041) (1.495,1.075) (1.550,1.117) (1.594,1.150) (1.650,1.193) (1.691,1.225) (1.750,1.271) (1.800,1.311) (1.850,1.350) (1.900,1.391) (1.942,1.425) (2.000,1.472) (2.050,1.514) (2.093,1.550) (2.150,1.599) (2.200,1.643) (2.250,1.687) (2.293,1.725) (2.348,1.775) (2.400,1.823) (2.450,1.870) (2.500,1.918) (2.550,1.967) (2.600,2.017) (2.650,2.068) (2.700,2.120) (2.750,2.173) (2.797,2.225) (2.842,2.275) (2.875,2.313) (2.925,2.372) (2.969,2.425) (3.009,2.475) (3.049,2.525) (3.087,2.575) (3.124,2.625) (3.160,2.675) (3.195,2.725) (3.225,2.770) (3.262,2.825) (3.294,2.875) (3.325,2.925) (3.355,2.975) (3.384,3.025) (3.412,3.075) (3.439,3.125) (3.465,3.175) (3.490,3.225) (3.514,3.275) (3.538,3.325) (3.560,3.375) (3.582,3.425) (3.602,3.475) (3.622,3.525) (3.641,3.575) (3.659,3.625) (3.675,3.673) (3.692,3.725) (3.708,3.775) (3.722,3.825) (3.736,3.875) (3.749,3.925) (3.761,3.975) (3.772,4.025) (3.783,4.075) (3.792,4.125) (3.801,4.175) (3.809,4.225) (3.816,4.275) (3.823,4.325) (3.828,4.375) (3.833,4.425) (3.837,4.475) (3.841,4.525) (3.843,4.575) (3.845,4.625) (3.846,4.675) (3.846,4.725) (3.845,4.775) (3.844,4.825) (3.842,4.875) (3.839,4.925) (3.835,4.975) (3.830,5.025) (3.825,5.075) (3.819,5.125) (3.812,5.175) (3.804,5.225) (3.796,5.275) (3.787,5.325) (3.776,5.375) (3.765,5.425) (3.754,5.475) (3.741,5.525) (3.728,5.575) (3.713,5.625) (3.698,5.675) (3.682,5.725) (3.665,5.775) (3.647,5.825) (3.629,5.875) (3.609,5.925) (3.589,5.975) (3.567,6.025) (3.545,6.075) (3.521,6.125) (3.497,6.175) (3.472,6.225) (3.445,6.275) (3.418,6.325) (3.389,6.375) (3.359,6.425) (3.328,6.475) (3.300,6.519) (3.263,6.575) (3.229,6.625) (3.193,6.675) (3.156,6.725) (3.125,6.765) (3.077,6.825) (3.036,6.875) (3.000,6.917) (2.950,6.974) (2.903,7.025) (2.855,7.075) (2.805,7.125) (2.753,7.175) (2.700,7.225) (2.650,7.269) (2.600,7.313) (2.555,7.350) (2.500,7.395) (2.450,7.434) (2.400,7.471) (2.350,7.508) (2.300,7.543) (2.253,7.575) (2.200,7.610) (2.150,7.642) (2.100,7.673) (2.054,7.700) (2.000,7.731) (1.950,7.759) (1.900,7.786) (1.850,7.812) (1.800,7.837) (1.750,7.861) (1.700,7.884) (1.650,7.907) (1.600,7.928) (1.550,7.949) (1.500,7.969) (1.450,7.989) (1.400,8.007) (1.350,8.025) (1.300,8.042) (1.250,8.058) (1.200,8.074) (1.150,8.089) (1.100,8.103) (1.050,8.117) (1.000,8.130) (0.950,8.142) (0.900,8.153) (0.850,8.164) (0.800,8.175) (0.750,8.184) (0.700,8.193) (0.650,8.201) (0.600,8.209) (0.550,8.216) (0.500,8.223) (0.450,8.228) (0.400,8.234) (0.350,8.238) (0.300,8.242) (0.250,8.245) (0.200,8.248) (0.150,8.250) (0.100,8.252) (0.050,8.253) (0.000,8.253) (-0.050,8.253) (-0.100,8.252) (-0.150,8.250) (-0.200,8.248) (-0.250,8.245) (-0.300,8.242) (-0.350,8.238) (-0.400,8.234) (-0.450,8.228) (-0.500,8.223) (-0.550,8.216) (-0.600,8.209) (-0.650,8.201) (-0.700,8.193) (-0.750,8.184) (-0.800,8.175) (-0.850,8.164) (-0.900,8.153) (-0.950,8.142) (-1.000,8.130) (-1.050,8.117) (-1.100,8.103) (-1.150,8.089) (-1.200,8.074) (-1.250,8.058) (-1.300,8.042) (-1.350,8.025) (-1.400,8.007) (-1.450,7.989) (-1.500,7.969) (-1.548,7.950) (-1.600,7.928) (-1.650,7.907) (-1.700,7.884) (-1.750,7.861) (-1.800,7.837) (-1.850,7.812) (-1.900,7.786) (-1.950,7.759) (-2.000,7.731) (-2.050,7.702) (-2.096,7.675) (-2.150,7.642) (-2.200,7.610) (-2.250,7.577) (-2.300,7.543) (-2.350,7.508) (-2.395,7.475) (-2.450,7.434) (-2.493,7.400) (-2.550,7.354) (-2.600,7.313) (-2.644,7.275) (-2.700,7.225) (-2.750,7.178) (-2.800,7.130) (-2.850,7.080) (-2.900,7.028) (-2.949,6.975) (-2.993,6.925) (-3.036,6.875) (-3.075,6.828) (-3.117,6.775) (-3.150,6.733) (-3.193,6.675) (-3.225,6.630) (-3.263,6.575) (-3.296,6.525) (-3.325,6.480) (-3.359,6.425) (-3.389,6.375) (-3.418,6.325) (-3.445,6.275) (-3.472,6.225) (-3.497,6.175) (-3.521,6.125) (-3.545,6.075) (-3.567,6.025) (-3.589,5.975) (-3.609,5.925) (-3.629,5.875) (-3.647,5.825) (-3.665,5.775) (-3.682,5.725) (-3.698,5.675) (-3.713,5.625) (-3.728,5.575) (-3.741,5.525) (-3.754,5.475) (-3.765,5.425) (-3.776,5.375) (-3.787,5.325) (-3.796,5.275) (-3.804,5.225) (-3.812,5.175) (-3.819,5.125) (-3.825,5.075) (-3.830,5.025) (-3.835,4.975) (-3.839,4.925) (-3.842,4.875) (-3.844,4.825) (-3.845,4.775) (-3.846,4.725) (-3.846,4.675) (-3.845,4.625) (-3.843,4.575) (-3.841,4.525) (-3.837,4.475) (-3.833,4.425) (-3.828,4.375) (-3.823,4.325) (-3.816,4.275) (-3.809,4.225) (-3.801,4.175) (-3.792,4.125) (-3.783,4.075) (-3.772,4.025) (-3.761,3.975) (-3.749,3.925) (-3.736,3.875) (-3.722,3.825) (-3.708,3.775) (-3.692,3.725) (-3.676,3.675) (-3.659,3.625) (-3.641,3.575) (-3.622,3.525) (-3.602,3.475) (-3.582,3.425) (-3.560,3.375) (-3.538,3.325) (-3.514,3.275) (-3.490,3.225) (-3.465,3.175) (-3.439,3.125) (-3.412,3.075) (-3.384,3.025) (-3.355,2.975) (-3.325,2.926) (-3.294,2.875) (-3.262,2.825) (-3.229,2.775) (-3.200,2.733) (-3.160,2.675) (-3.125,2.627) (-3.087,2.575) (-3.050,2.527) (-3.009,2.475) (-2.975,2.432) (-2.928,2.375) (-2.886,2.325) (-2.850,2.284) (-2.800,2.228) (-2.752,2.175) (-2.705,2.125) (-2.657,2.075) (-2.609,2.025) (-2.559,1.975) (-2.508,1.925) (-2.456,1.875) (-2.403,1.825) (-2.350,1.777) (-2.300,1.731) (-2.250,1.687) (-2.209,1.650) (-2.151,1.600) (-2.100,1.556) (-2.050,1.514) (-2.003,1.475) (-1.950,1.431) (-1.900,1.391) (-1.850,1.350) (-1.800,1.311) (-1.755,1.275) (-1.700,1.232) (-1.659,1.200) (-1.600,1.155) (-1.550,1.117) (-1.500,1.079) (-1.450,1.041) (-1.400,1.004) (-1.350,0.966) (-1.300,0.929) (-1.250,0.893) (-1.200,0.856) (-1.158,0.825) (-1.100,0.783) (-1.055,0.750) (-1.000,0.711) (-0.951,0.675) (-0.900,0.639) (-0.850,0.603) (-0.800,0.567) (-0.750,0.532) (-0.706,0.500) (-0.650,0.461) (-0.600,0.425) (-0.550,0.390) (-0.500,0.354) (-0.450,0.319) (-0.400,0.284) (-0.354,0.250) (-0.300,0.213) (-0.250,0.177) (-0.220,0.155)}
\EQ{(-0.220,-0.155) (-0.275,-0.195) (-0.319,-0.225) (-0.375,-0.266) (-0.424,-0.300) (-0.475,-0.337) (-0.525,-0.372) (-0.575,-0.407) (-0.625,-0.443) (-0.671,-0.475) (-0.725,-0.514) (-0.775,-0.550) (-0.825,-0.586) (-0.875,-0.622) (-0.925,-0.658) (-0.975,-0.694) (-1.019,-0.725) (-1.075,-0.767) (-1.121,-0.800) (-1.175,-0.840) (-1.223,-0.875) (-1.275,-0.914) (-1.323,-0.950) (-1.375,-0.989) (-1.422,-1.025) (-1.475,-1.066) (-1.520,-1.100) (-1.575,-1.143) (-1.616,-1.175) (-1.675,-1.222) (-1.725,-1.263) (-1.771,-1.300) (-1.825,-1.344) (-1.875,-1.386) (-1.921,-1.425) (-1.975,-1.472) (-2.025,-1.515) (-2.064,-1.550) (-2.120,-1.600) (-2.174,-1.650) (-2.225,-1.698) (-2.275,-1.746) (-2.325,-1.795) (-2.375,-1.845) (-2.425,-1.896) (-2.475,-1.949) (-2.522,-2.000) (-2.567,-2.050) (-2.611,-2.100) (-2.650,-2.146) (-2.695,-2.200) (-2.735,-2.250) (-2.774,-2.300) (-2.812,-2.350) (-2.849,-2.400) (-2.885,-2.450) (-2.919,-2.500) (-2.950,-2.546) (-2.985,-2.600) (-3.016,-2.650) (-3.046,-2.700) (-3.075,-2.750) (-3.100,-2.794) (-3.125,-2.841) (-3.150,-2.889) (-3.175,-2.939) (-3.200,-2.991) (-3.225,-3.045) (-3.249,-3.100) (-3.270,-3.150) (-3.290,-3.200) (-3.309,-3.250) (-3.326,-3.300) (-3.343,-3.350) (-3.359,-3.400) (-3.374,-3.450) (-3.389,-3.500) (-3.402,-3.550) (-3.414,-3.600) (-3.426,-3.650) (-3.436,-3.700) (-3.446,-3.750) (-3.455,-3.800) (-3.462,-3.850) (-3.469,-3.900) (-3.476,-3.950) (-3.481,-4.000) (-3.485,-4.050) (-3.489,-4.100) (-3.491,-4.150) (-3.493,-4.200) (-3.494,-4.250) (-3.494,-4.300) (-3.493,-4.350) (-3.492,-4.400) (-3.489,-4.450) (-3.486,-4.500) (-3.482,-4.550) (-3.477,-4.600) (-3.471,-4.650) (-3.464,-4.700) (-3.456,-4.750) (-3.447,-4.800) (-3.438,-4.850) (-3.428,-4.900) (-3.416,-4.950) (-3.404,-5.000) (-3.391,-5.050) (-3.377,-5.100) (-3.362,-5.150) (-3.346,-5.200) (-3.329,-5.250) (-3.311,-5.300) (-3.292,-5.350) (-3.272,-5.400) (-3.251,-5.450) (-3.229,-5.500) (-3.206,-5.550) (-3.182,-5.600) (-3.157,-5.650) (-3.131,-5.700) (-3.103,-5.750) (-3.075,-5.799) (-3.044,-5.850) (-3.013,-5.900) (-2.981,-5.950) (-2.950,-5.995) (-2.912,-6.050) (-2.875,-6.100) (-2.837,-6.150) (-2.800,-6.196) (-2.755,-6.250) (-2.712,-6.300) (-2.675,-6.342) (-2.625,-6.396) (-2.575,-6.447) (-2.525,-6.496) (-2.475,-6.544) (-2.425,-6.589) (-2.384,-6.625) (-2.325,-6.675) (-2.275,-6.716) (-2.231,-6.750) (-2.175,-6.792) (-2.130,-6.825) (-2.075,-6.863) (-2.025,-6.897) (-1.982,-6.925) (-1.925,-6.961) (-1.875,-6.991) (-1.825,-7.020) (-1.775,-7.048) (-1.725,-7.074) (-1.676,-7.100) (-1.625,-7.125) (-1.575,-7.149) (-1.525,-7.172) (-1.475,-7.194) (-1.425,-7.215) (-1.375,-7.236) (-1.325,-7.255) (-1.275,-7.274) (-1.225,-7.292) (-1.175,-7.309) (-1.125,-7.325) (-1.075,-7.340) (-1.025,-7.355) (-0.975,-7.369) (-0.925,-7.382) (-0.875,-7.395) (-0.825,-7.406) (-0.775,-7.417) (-0.725,-7.427) (-0.675,-7.437) (-0.625,-7.446) (-0.575,-7.454) (-0.525,-7.461) (-0.475,-7.468) (-0.425,-7.474) (-0.375,-7.479) (-0.325,-7.484) (-0.275,-7.488) (-0.225,-7.492) (-0.175,-7.494) (-0.125,-7.496) (-0.075,-7.498) (-0.025,-7.498) (0.050,-7.498) (0.100,-7.497) (0.150,-7.495) (0.200,-7.493) (0.250,-7.490) (0.300,-7.486) (0.350,-7.482) (0.400,-7.477) (0.450,-7.471) (0.500,-7.465) (0.550,-7.458) (0.599,-7.450) (0.650,-7.441) (0.700,-7.432) (0.750,-7.422) (0.800,-7.412) (0.850,-7.400) (0.900,-7.388) (0.950,-7.376) (1.000,-7.362) (1.050,-7.348) (1.100,-7.333) (1.150,-7.317) (1.200,-7.300) (1.250,-7.283) (1.300,-7.265) (1.350,-7.246) (1.400,-7.226) (1.450,-7.205) (1.500,-7.183) (1.550,-7.161) (1.600,-7.137) (1.650,-7.113) (1.700,-7.088) (1.750,-7.061) (1.800,-7.034) (1.850,-7.005) (1.900,-6.976) (1.950,-6.945) (2.000,-6.913) (2.050,-6.880) (2.094,-6.850) (2.150,-6.811) (2.198,-6.775) (2.250,-6.735) (2.295,-6.700) (2.350,-6.654) (2.400,-6.611) (2.441,-6.575) (2.495,-6.525) (2.547,-6.475) (2.597,-6.425) (2.644,-6.375) (2.690,-6.325) (2.725,-6.286) (2.775,-6.227) (2.817,-6.175) (2.850,-6.133) (2.893,-6.075) (2.925,-6.031) (2.964,-5.975) (2.997,-5.925) (3.025,-5.881) (3.060,-5.825) (3.089,-5.775) (3.117,-5.725) (3.144,-5.675) (3.170,-5.625) (3.195,-5.575) (3.218,-5.525) (3.241,-5.475) (3.262,-5.425) (3.282,-5.375) (3.302,-5.325) (3.320,-5.275) (3.338,-5.225) (3.354,-5.175) (3.370,-5.125) (3.384,-5.075) (3.398,-5.025) (3.410,-4.975) (3.422,-4.925) (3.433,-4.875) (3.443,-4.825) (3.452,-4.775) (3.460,-4.725) (3.467,-4.675) (3.474,-4.625) (3.479,-4.575) (3.484,-4.525) (3.488,-4.475) (3.491,-4.425) (3.493,-4.375) (3.494,-4.325) (3.494,-4.275) (3.494,-4.225) (3.492,-4.175) (3.490,-4.125) (3.487,-4.075) (3.483,-4.025) (3.478,-3.975) (3.473,-3.925) (3.466,-3.875) (3.459,-3.825) (3.450,-3.775) (3.441,-3.725) (3.431,-3.675) (3.420,-3.625) (3.408,-3.575) (3.395,-3.525) (3.382,-3.475) (3.367,-3.425) (3.352,-3.375) (3.335,-3.325) (3.318,-3.275) (3.300,-3.227) (3.280,-3.175) (3.260,-3.125) (3.238,-3.075) (3.216,-3.025) (3.193,-2.975) (3.168,-2.925) (3.143,-2.875) (3.117,-2.825) (3.089,-2.775) (3.061,-2.725) (3.031,-2.675) (3.001,-2.625) (2.969,-2.575) (2.936,-2.525) (2.902,-2.475) (2.867,-2.425) (2.831,-2.375) (2.800,-2.334) (2.755,-2.275) (2.715,-2.225) (2.675,-2.176) (2.632,-2.125) (2.589,-2.075) (2.550,-2.031) (2.500,-1.976) (2.452,-1.925) (2.404,-1.875) (2.355,-1.825) (2.305,-1.775) (2.254,-1.725) (2.201,-1.675) (2.150,-1.628) (2.100,-1.582) (2.050,-1.537) (2.008,-1.500) (1.950,-1.450) (1.900,-1.407) (1.850,-1.365) (1.802,-1.325) (1.750,-1.283) (1.710,-1.250) (1.650,-1.202) (1.600,-1.163) (1.552,-1.125) (1.500,-1.085) (1.455,-1.050) (1.400,-1.008) (1.357,-0.975) (1.300,-0.933) (1.257,-0.900) (1.200,-0.858) (1.155,-0.825) (1.100,-0.785) (1.053,-0.750) (1.000,-0.712) (0.950,-0.676) (0.900,-0.640) (0.850,-0.604) (0.800,-0.568) (0.750,-0.532) (0.706,-0.500) (0.650,-0.461) (0.600,-0.425) (0.550,-0.390) (0.500,-0.354) (0.450,-0.319) (0.400,-0.284) (0.354,-0.250) (0.300,-0.213) (0.250,-0.177) (0.206,-0.144)}
\EQ{(0.206,0.144) (0.248,0.175) (0.300,0.213) (0.350,0.248) (0.400,0.284) (0.450,0.319) (0.495,0.350) (0.550,0.390) (0.600,0.425) (0.650,0.461) (0.700,0.496) (0.741,0.525) (0.800,0.568) (0.845,0.600) (0.900,0.640) (0.949,0.675) (1.000,0.712) (1.050,0.748) (1.100,0.785) (1.150,0.821) (1.200,0.858) (1.250,0.895) (1.300,0.933) (1.350,0.970) (1.400,1.008) (1.450,1.046) (1.500,1.085) (1.550,1.124) (1.600,1.163) (1.647,1.200) (1.700,1.242) (1.741,1.275) (1.800,1.324) (1.850,1.365) (1.892,1.400) (1.950,1.450) (2.000,1.493) (2.050,1.537) (2.092,1.575) (2.147,1.625) (2.200,1.674) (2.250,1.722) (2.300,1.770) (2.350,1.820) (2.400,1.871) (2.450,1.923) (2.499,1.975) (2.545,2.025) (2.589,2.075) (2.625,2.117) (2.674,2.175) (2.715,2.225) (2.750,2.269) (2.793,2.325) (2.825,2.367) (2.867,2.425) (2.900,2.472) (2.936,2.525) (2.969,2.575) (3.000,2.624) (3.031,2.675) (3.061,2.725) (3.089,2.775) (3.117,2.825) (3.143,2.875) (3.168,2.925) (3.193,2.975) (3.216,3.025) (3.238,3.075) (3.260,3.125) (3.280,3.175) (3.299,3.225) (3.318,3.275) (3.335,3.325) (3.352,3.375) (3.367,3.425) (3.382,3.475) (3.395,3.525) (3.408,3.575) (3.420,3.625) (3.431,3.675) (3.441,3.725) (3.450,3.775) (3.459,3.825) (3.466,3.875) (3.473,3.925) (3.478,3.975) (3.483,4.025) (3.487,4.075) (3.490,4.125) (3.492,4.175) (3.494,4.225) (3.494,4.275) (3.494,4.325) (3.493,4.375) (3.491,4.425) (3.488,4.475) (3.484,4.525) (3.479,4.575) (3.474,4.625) (3.467,4.675) (3.460,4.725) (3.452,4.775) (3.443,4.825) (3.433,4.875) (3.422,4.925) (3.410,4.975) (3.398,5.025) (3.384,5.075) (3.370,5.125) (3.354,5.175) (3.338,5.225) (3.320,5.275) (3.302,5.325) (3.282,5.375) (3.262,5.425) (3.241,5.475) (3.218,5.525) (3.195,5.575) (3.170,5.625) (3.144,5.675) (3.117,5.725) (3.089,5.775) (3.060,5.825) (3.029,5.875) (3.000,5.921) (2.964,5.975) (2.929,6.025) (2.900,6.066) (2.856,6.125) (2.825,6.165) (2.776,6.225) (2.734,6.275) (2.700,6.314) (2.650,6.369) (2.600,6.422) (2.550,6.472) (2.500,6.520) (2.450,6.567) (2.400,6.611) (2.355,6.650) (2.300,6.696) (2.250,6.735) (2.200,6.774) (2.150,6.811) (2.100,6.846) (2.058,6.875) (2.000,6.913) (1.950,6.945) (1.901,6.975) (1.850,7.005) (1.800,7.034) (1.750,7.061) (1.700,7.088) (1.650,7.113) (1.600,7.137) (1.550,7.161) (1.500,7.183) (1.450,7.205) (1.402,7.225) (1.350,7.246) (1.300,7.265) (1.250,7.283) (1.201,7.300) (1.150,7.317) (1.100,7.333) (1.050,7.348) (1.000,7.362) (0.950,7.376) (0.900,7.388) (0.850,7.400) (0.800,7.412) (0.750,7.422) (0.700,7.432) (0.650,7.441) (0.600,7.450) (0.550,7.458) (0.500,7.465) (0.450,7.471) (0.400,7.477) (0.350,7.482) (0.300,7.486) (0.250,7.490) (0.200,7.493) (0.150,7.495) (0.100,7.497) (0.050,7.498) (0.000,7.498) (-0.050,7.498) (-0.100,7.497) (-0.150,7.495) (-0.200,7.493) (-0.250,7.490) (-0.300,7.486) (-0.350,7.482) (-0.400,7.477) (-0.450,7.471) (-0.500,7.465) (-0.550,7.458) (-0.599,7.450) (-0.650,7.441) (-0.700,7.432) (-0.750,7.422) (-0.800,7.412) (-0.850,7.400) (-0.900,7.388) (-0.950,7.376) (-1.000,7.362) (-1.050,7.348) (-1.100,7.333) (-1.150,7.317) (-1.200,7.300) (-1.250,7.283) (-1.300,7.265) (-1.350,7.246) (-1.400,7.226) (-1.450,7.205) (-1.500,7.183) (-1.550,7.161) (-1.600,7.137) (-1.650,7.113) (-1.700,7.088) (-1.750,7.061) (-1.800,7.034) (-1.850,7.005) (-1.900,6.976) (-1.950,6.945) (-2.000,6.913) (-2.050,6.880) (-2.094,6.850) (-2.150,6.811) (-2.198,6.775) (-2.250,6.735) (-2.295,6.700) (-2.350,6.654) (-2.400,6.611) (-2.441,6.575) (-2.495,6.525) (-2.547,6.475) (-2.597,6.425) (-2.644,6.375) (-2.690,6.325) (-2.725,6.286) (-2.775,6.227) (-2.817,6.175) (-2.850,6.133) (-2.893,6.075) (-2.925,6.031) (-2.964,5.975) (-2.997,5.925) (-3.025,5.881) (-3.060,5.825) (-3.089,5.775) (-3.117,5.725) (-3.144,5.675) (-3.170,5.625) (-3.195,5.575) (-3.218,5.525) (-3.241,5.475) (-3.262,5.425) (-3.282,5.375) (-3.302,5.325) (-3.320,5.275) (-3.338,5.225) (-3.354,5.175) (-3.370,5.125) (-3.384,5.075) (-3.398,5.025) (-3.410,4.975) (-3.422,4.925) (-3.433,4.875) (-3.443,4.825) (-3.452,4.775) (-3.460,4.725) (-3.467,4.675) (-3.474,4.625) (-3.479,4.575) (-3.484,4.525) (-3.488,4.475) (-3.491,4.425) (-3.493,4.375) (-3.494,4.325) (-3.494,4.275) (-3.494,4.225) (-3.492,4.175) (-3.490,4.125) (-3.487,4.075) (-3.483,4.025) (-3.478,3.975) (-3.473,3.925) (-3.466,3.875) (-3.459,3.825) (-3.450,3.775) (-3.441,3.725) (-3.431,3.675) (-3.420,3.625) (-3.408,3.575) (-3.395,3.525) (-3.382,3.475) (-3.367,3.425) (-3.352,3.375) (-3.335,3.325) (-3.318,3.275) (-3.300,3.227) (-3.280,3.175) (-3.260,3.125) (-3.238,3.075) (-3.216,3.025) (-3.193,2.975) (-3.168,2.925) (-3.143,2.875) (-3.117,2.825) (-3.089,2.775) (-3.061,2.725) (-3.031,2.675) (-3.001,2.625) (-2.969,2.575) (-2.936,2.525) (-2.902,2.475) (-2.867,2.425) (-2.831,2.375) (-2.800,2.334) (-2.755,2.275) (-2.715,2.225) (-2.675,2.176) (-2.632,2.125) (-2.589,2.075) (-2.550,2.031) (-2.500,1.976) (-2.452,1.925) (-2.404,1.875) (-2.355,1.825) (-2.305,1.775) (-2.254,1.725) (-2.201,1.675) (-2.150,1.628) (-2.100,1.582) (-2.050,1.537) (-2.008,1.500) (-1.950,1.450) (-1.900,1.407) (-1.850,1.365) (-1.802,1.325) (-1.750,1.283) (-1.710,1.250) (-1.650,1.202) (-1.600,1.163) (-1.552,1.125) (-1.500,1.085) (-1.455,1.050) (-1.400,1.008) (-1.357,0.975) (-1.300,0.933) (-1.257,0.900) (-1.200,0.858) (-1.155,0.825) (-1.100,0.785) (-1.053,0.750) (-1.000,0.712) (-0.950,0.676) (-0.900,0.640) (-0.850,0.604) (-0.800,0.568) (-0.750,0.532) (-0.706,0.500) (-0.650,0.461) (-0.600,0.425) (-0.550,0.390) (-0.500,0.354) (-0.450,0.319) (-0.400,0.284) (-0.354,0.250) (-0.300,0.213) (-0.250,0.177) (-0.220,0.155)}
\EQ{(-0.220,-0.155) (-0.275,-0.195) (-0.319,-0.225) (-0.375,-0.266) (-0.424,-0.300) (-0.475,-0.337) (-0.525,-0.372) (-0.575,-0.408) (-0.625,-0.443) (-0.671,-0.475) (-0.725,-0.515) (-0.775,-0.550) (-0.825,-0.586) (-0.875,-0.622) (-0.925,-0.659) (-0.975,-0.695) (-1.017,-0.725) (-1.075,-0.768) (-1.119,-0.800) (-1.175,-0.842) (-1.219,-0.875) (-1.275,-0.917) (-1.319,-0.950) (-1.375,-0.994) (-1.416,-1.025) (-1.475,-1.071) (-1.525,-1.111) (-1.574,-1.150) (-1.625,-1.191) (-1.667,-1.225) (-1.725,-1.274) (-1.775,-1.316) (-1.815,-1.350) (-1.872,-1.400) (-1.925,-1.447) (-1.975,-1.492) (-2.025,-1.538) (-2.064,-1.575) (-2.117,-1.625) (-2.167,-1.675) (-2.217,-1.725) (-2.265,-1.775) (-2.300,-1.812) (-2.335,-1.850) (-2.375,-1.894) (-2.424,-1.950) (-2.467,-2.000) (-2.500,-2.040) (-2.548,-2.100) (-2.587,-2.150) (-2.624,-2.200) (-2.661,-2.250) (-2.696,-2.300) (-2.725,-2.343) (-2.762,-2.400) (-2.794,-2.450) (-2.824,-2.500) (-2.850,-2.545) (-2.875,-2.589) (-2.900,-2.635) (-2.925,-2.683) (-2.950,-2.733) (-2.975,-2.786) (-3.000,-2.841) (-3.025,-2.899) (-3.046,-2.950) (-3.065,-3.000) (-3.083,-3.050) (-3.100,-3.099) (-3.117,-3.150) (-3.132,-3.200) (-3.146,-3.250) (-3.159,-3.300) (-3.171,-3.350) (-3.183,-3.400) (-3.193,-3.450) (-3.202,-3.500) (-3.211,-3.550) (-3.218,-3.600) (-3.224,-3.650) (-3.230,-3.700) (-3.235,-3.750) (-3.238,-3.800) (-3.241,-3.850) (-3.243,-3.900) (-3.244,-3.950) (-3.244,-4.000) (-3.243,-4.050) (-3.241,-4.100) (-3.238,-4.150) (-3.234,-4.200) (-3.229,-4.250) (-3.224,-4.300) (-3.217,-4.350) (-3.210,-4.400) (-3.201,-4.450) (-3.192,-4.500) (-3.182,-4.550) (-3.170,-4.600) (-3.158,-4.650) (-3.145,-4.700) (-3.130,-4.750) (-3.115,-4.800) (-3.099,-4.850) (-3.081,-4.900) (-3.063,-4.950) (-3.043,-5.000) (-3.022,-5.050) (-3.001,-5.100) (-2.978,-5.150) (-2.954,-5.200) (-2.928,-5.250) (-2.902,-5.300) (-2.875,-5.348) (-2.845,-5.400) (-2.814,-5.450) (-2.782,-5.500) (-2.750,-5.549) (-2.715,-5.600) (-2.678,-5.650) (-2.640,-5.700) (-2.601,-5.750) (-2.560,-5.800) (-2.525,-5.840) (-2.475,-5.896) (-2.425,-5.949) (-2.375,-6.000) (-2.325,-6.049) (-2.275,-6.096) (-2.225,-6.140) (-2.184,-6.175) (-2.125,-6.224) (-2.075,-6.263) (-2.026,-6.300) (-1.975,-6.337) (-1.925,-6.372) (-1.875,-6.406) (-1.825,-6.438) (-1.775,-6.469) (-1.725,-6.499) (-1.679,-6.525) (-1.634,-6.550) (-1.587,-6.575) (-1.538,-6.600) (-1.487,-6.625) (-1.433,-6.650) (-1.377,-6.675) (-1.325,-6.697) (-1.275,-6.718) (-1.225,-6.737) (-1.175,-6.756) (-1.125,-6.773) (-1.075,-6.790) (-1.025,-6.806) (-0.975,-6.821) (-0.925,-6.835) (-0.875,-6.849) (-0.825,-6.861) (-0.775,-6.873) (-0.725,-6.884) (-0.675,-6.895) (-0.625,-6.904) (-0.575,-6.913) (-0.525,-6.921) (-0.475,-6.928) (-0.425,-6.935) (-0.375,-6.941) (-0.325,-6.946) (-0.275,-6.950) (-0.225,-6.954) (-0.175,-6.956) (-0.125,-6.959) (-0.075,-6.960) (-0.025,-6.961) (0.050,-6.960) (0.100,-6.959) (0.150,-6.958) (0.200,-6.955) (0.250,-6.952) (0.300,-6.948) (0.350,-6.943) (0.400,-6.938) (0.450,-6.932) (0.500,-6.925) (0.550,-6.917) (0.600,-6.909) (0.650,-6.899) (0.700,-6.890) (0.750,-6.879) (0.800,-6.867) (0.850,-6.855) (0.900,-6.842) (0.950,-6.828) (1.000,-6.814) (1.050,-6.798) (1.100,-6.782) (1.150,-6.765) (1.200,-6.747) (1.250,-6.728) (1.300,-6.708) (1.350,-6.687) (1.400,-6.665) (1.450,-6.642) (1.500,-6.619) (1.550,-6.594) (1.600,-6.568) (1.650,-6.541) (1.700,-6.513) (1.750,-6.484) (1.800,-6.454) (1.845,-6.425) (1.900,-6.389) (1.950,-6.355) (1.992,-6.325) (2.050,-6.282) (2.092,-6.250) (2.150,-6.204) (2.200,-6.162) (2.242,-6.125) (2.297,-6.075) (2.350,-6.025) (2.400,-5.975) (2.448,-5.925) (2.494,-5.875) (2.538,-5.825) (2.575,-5.782) (2.621,-5.725) (2.660,-5.675) (2.697,-5.625) (2.732,-5.575) (2.766,-5.525) (2.799,-5.475) (2.830,-5.425) (2.859,-5.375) (2.888,-5.325) (2.915,-5.275) (2.941,-5.225) (2.966,-5.175) (2.989,-5.125) (3.012,-5.075) (3.033,-5.025) (3.053,-4.975) (3.072,-4.925) (3.090,-4.875) (3.107,-4.825) (3.123,-4.775) (3.138,-4.725) (3.151,-4.675) (3.164,-4.625) (3.176,-4.575) (3.187,-4.525) (3.197,-4.475) (3.206,-4.425) (3.214,-4.375) (3.221,-4.325) (3.227,-4.275) (3.232,-4.225) (3.236,-4.175) (3.239,-4.125) (3.242,-4.075) (3.243,-4.025) (3.244,-3.975) (3.243,-3.925) (3.242,-3.875) (3.240,-3.825) (3.237,-3.775) (3.232,-3.725) (3.227,-3.675) (3.221,-3.625) (3.214,-3.575) (3.207,-3.525) (3.198,-3.475) (3.188,-3.425) (3.177,-3.375) (3.165,-3.325) (3.153,-3.275) (3.139,-3.225) (3.124,-3.175) (3.109,-3.125) (3.092,-3.075) (3.075,-3.027) (3.056,-2.975) (3.036,-2.925) (3.015,-2.875) (2.993,-2.825) (2.970,-2.775) (2.946,-2.725) (2.921,-2.675) (2.895,-2.625) (2.867,-2.575) (2.839,-2.525) (2.809,-2.475) (2.778,-2.425) (2.750,-2.381) (2.713,-2.325) (2.678,-2.275) (2.643,-2.225) (2.606,-2.175) (2.575,-2.135) (2.528,-2.075) (2.488,-2.025) (2.450,-1.980) (2.402,-1.925) (2.358,-1.875) (2.312,-1.825) (2.275,-1.785) (2.225,-1.733) (2.175,-1.683) (2.125,-1.633) (2.075,-1.585) (2.038,-1.550) (2.000,-1.515) (1.957,-1.475) (1.901,-1.425) (1.850,-1.381) (1.800,-1.337) (1.756,-1.300) (1.700,-1.253) (1.650,-1.212) (1.605,-1.175) (1.550,-1.131) (1.500,-1.091) (1.450,-1.052) (1.400,-1.013) (1.351,-0.975) (1.300,-0.936) (1.252,-0.900) (1.200,-0.861) (1.152,-0.825) (1.100,-0.787) (1.051,-0.750) (1.000,-0.713) (0.950,-0.677) (0.900,-0.640) (0.850,-0.604) (0.800,-0.568) (0.750,-0.532) (0.705,-0.500) (0.650,-0.461) (0.600,-0.425) (0.550,-0.390) (0.500,-0.354) (0.450,-0.319) (0.400,-0.284) (0.354,-0.250) (0.300,-0.213) (0.250,-0.177) (0.206,-0.144)}
\EQ{(0.206,0.144) (0.248,0.175) (0.300,0.213) (0.350,0.248) (0.400,0.284) (0.450,0.319) (0.495,0.350) (0.550,0.390) (0.600,0.425) (0.650,0.461) (0.700,0.497) (0.750,0.532) (0.800,0.568) (0.845,0.600) (0.900,0.640) (0.948,0.675) (1.000,0.713) (1.050,0.750) (1.100,0.787) (1.150,0.823) (1.200,0.861) (1.250,0.898) (1.300,0.936) (1.350,0.974) (1.400,1.013) (1.448,1.050) (1.500,1.091) (1.543,1.125) (1.600,1.171) (1.650,1.212) (1.697,1.250) (1.750,1.295) (1.800,1.337) (1.844,1.375) (1.900,1.424) (1.950,1.469) (2.000,1.515) (2.038,1.550) (2.075,1.585) (2.117,1.625) (2.167,1.675) (2.217,1.725) (2.265,1.775) (2.300,1.812) (2.335,1.850) (2.375,1.894) (2.424,1.950) (2.467,2.000) (2.500,2.040) (2.548,2.100) (2.587,2.150) (2.624,2.200) (2.661,2.250) (2.696,2.300) (2.725,2.343) (2.762,2.400) (2.794,2.450) (2.824,2.500) (2.850,2.545) (2.875,2.589) (2.900,2.635) (2.925,2.683) (2.950,2.733) (2.975,2.786) (3.000,2.841) (3.025,2.899) (3.046,2.950) (3.065,3.000) (3.083,3.050) (3.100,3.099) (3.117,3.150) (3.132,3.200) (3.146,3.250) (3.159,3.300) (3.171,3.350) (3.183,3.400) (3.193,3.450) (3.202,3.500) (3.211,3.550) (3.218,3.600) (3.224,3.650) (3.230,3.700) (3.235,3.750) (3.238,3.800) (3.241,3.850) (3.243,3.900) (3.244,3.950) (3.244,4.000) (3.243,4.050) (3.241,4.100) (3.238,4.150) (3.234,4.200) (3.229,4.250) (3.224,4.300) (3.217,4.350) (3.210,4.400) (3.201,4.450) (3.192,4.500) (3.182,4.550) (3.170,4.600) (3.158,4.650) (3.145,4.700) (3.130,4.750) (3.115,4.800) (3.099,4.850) (3.081,4.900) (3.063,4.950) (3.043,5.000) (3.022,5.050) (3.001,5.100) (2.978,5.150) (2.954,5.200) (2.928,5.250) (2.902,5.300) (2.875,5.348) (2.845,5.400) (2.814,5.450) (2.782,5.500) (2.750,5.549) (2.715,5.600) (2.678,5.650) (2.640,5.700) (2.601,5.750) (2.560,5.800) (2.525,5.840) (2.475,5.896) (2.425,5.949) (2.375,6.000) (2.325,6.049) (2.275,6.096) (2.225,6.140) (2.184,6.175) (2.125,6.224) (2.075,6.263) (2.026,6.300) (1.975,6.337) (1.925,6.372) (1.875,6.406) (1.825,6.438) (1.775,6.469) (1.725,6.499) (1.679,6.525) (1.634,6.550) (1.587,6.575) (1.538,6.600) (1.487,6.625) (1.433,6.650) (1.377,6.675) (1.325,6.697) (1.275,6.718) (1.225,6.737) (1.175,6.756) (1.125,6.773) (1.075,6.790) (1.025,6.806) (0.975,6.821) (0.925,6.835) (0.875,6.849) (0.825,6.861) (0.775,6.873) (0.725,6.884) (0.675,6.895) (0.625,6.904) (0.575,6.913) (0.525,6.921) (0.475,6.928) (0.425,6.935) (0.375,6.941) (0.325,6.946) (0.275,6.950) (0.225,6.954) (0.175,6.956) (0.125,6.959) (0.075,6.960) (0.025,6.961) (-0.050,6.960) (-0.100,6.959) (-0.150,6.958) (-0.200,6.955) (-0.250,6.952) (-0.300,6.948) (-0.350,6.943) (-0.400,6.938) (-0.450,6.932) (-0.500,6.925) (-0.550,6.917) (-0.600,6.909) (-0.650,6.899) (-0.700,6.890) (-0.750,6.879) (-0.800,6.867) (-0.850,6.855) (-0.900,6.842) (-0.950,6.828) (-1.000,6.814) (-1.050,6.798) (-1.100,6.782) (-1.150,6.765) (-1.200,6.747) (-1.250,6.728) (-1.300,6.708) (-1.350,6.687) (-1.400,6.665) (-1.450,6.642) (-1.500,6.619) (-1.550,6.594) (-1.600,6.568) (-1.650,6.541) (-1.700,6.513) (-1.750,6.484) (-1.800,6.454) (-1.845,6.425) (-1.900,6.389) (-1.950,6.355) (-1.992,6.325) (-2.050,6.282) (-2.092,6.250) (-2.150,6.204) (-2.200,6.162) (-2.242,6.125) (-2.297,6.075) (-2.350,6.025) (-2.400,5.975) (-2.448,5.925) (-2.494,5.875) (-2.538,5.825) (-2.575,5.782) (-2.621,5.725) (-2.660,5.675) (-2.697,5.625) (-2.732,5.575) (-2.766,5.525) (-2.799,5.475) (-2.830,5.425) (-2.859,5.375) (-2.888,5.325) (-2.915,5.275) (-2.941,5.225) (-2.966,5.175) (-2.989,5.125) (-3.012,5.075) (-3.033,5.025) (-3.053,4.975) (-3.072,4.925) (-3.090,4.875) (-3.107,4.825) (-3.123,4.775) (-3.138,4.725) (-3.151,4.675) (-3.164,4.625) (-3.176,4.575) (-3.187,4.525) (-3.197,4.475) (-3.206,4.425) (-3.214,4.375) (-3.221,4.325) (-3.227,4.275) (-3.232,4.225) (-3.236,4.175) (-3.239,4.125) (-3.242,4.075) (-3.243,4.025) (-3.244,3.975) (-3.243,3.925) (-3.242,3.875) (-3.240,3.825) (-3.237,3.775) (-3.232,3.725) (-3.227,3.675) (-3.221,3.625) (-3.214,3.575) (-3.207,3.525) (-3.198,3.475) (-3.188,3.425) (-3.177,3.375) (-3.165,3.325) (-3.153,3.275) (-3.139,3.225) (-3.124,3.175) (-3.109,3.125) (-3.092,3.075) (-3.075,3.027) (-3.056,2.975) (-3.036,2.925) (-3.015,2.875) (-2.993,2.825) (-2.970,2.775) (-2.946,2.725) (-2.921,2.675) (-2.895,2.625) (-2.867,2.575) (-2.839,2.525) (-2.809,2.475) (-2.778,2.425) (-2.750,2.381) (-2.713,2.325) (-2.678,2.275) (-2.643,2.225) (-2.606,2.175) (-2.575,2.135) (-2.528,2.075) (-2.488,2.025) (-2.450,1.980) (-2.402,1.925) (-2.358,1.875) (-2.312,1.825) (-2.275,1.785) (-2.225,1.733) (-2.175,1.683) (-2.125,1.633) (-2.075,1.585) (-2.038,1.550) (-2.000,1.515) (-1.957,1.475) (-1.901,1.425) (-1.850,1.381) (-1.800,1.337) (-1.756,1.300) (-1.700,1.253) (-1.650,1.212) (-1.605,1.175) (-1.550,1.131) (-1.500,1.091) (-1.450,1.052) (-1.400,1.013) (-1.351,0.975) (-1.300,0.936) (-1.252,0.900) (-1.200,0.861) (-1.152,0.825) (-1.100,0.787) (-1.051,0.750) (-1.000,0.713) (-0.950,0.677) (-0.900,0.640) (-0.850,0.604) (-0.800,0.568) (-0.750,0.532) (-0.705,0.500) (-0.650,0.461) (-0.600,0.425) (-0.550,0.390) (-0.500,0.354) (-0.450,0.319) (-0.400,0.284) (-0.354,0.250) (-0.300,0.213) (-0.250,0.177) (-0.220,0.155)}
\EQ{(-0.220,-0.155) (-0.275,-0.195) (-0.319,-0.225) (-0.375,-0.266) (-0.424,-0.300) (-0.475,-0.337) (-0.525,-0.372) (-0.575,-0.408) (-0.625,-0.443) (-0.670,-0.475) (-0.725,-0.515) (-0.774,-0.550) (-0.825,-0.587) (-0.875,-0.623) (-0.925,-0.659) (-0.975,-0.696) (-1.025,-0.733) (-1.075,-0.770) (-1.116,-0.800) (-1.175,-0.844) (-1.216,-0.875) (-1.275,-0.920) (-1.325,-0.959) (-1.375,-0.998) (-1.425,-1.037) (-1.472,-1.075) (-1.525,-1.118) (-1.565,-1.150) (-1.625,-1.200) (-1.675,-1.242) (-1.713,-1.275) (-1.771,-1.325) (-1.825,-1.374) (-1.875,-1.419) (-1.925,-1.466) (-1.962,-1.500) (-2.000,-1.537) (-2.038,-1.575) (-2.075,-1.612) (-2.112,-1.650) (-2.150,-1.690) (-2.200,-1.744) (-2.250,-1.799) (-2.294,-1.850) (-2.336,-1.900) (-2.375,-1.948) (-2.417,-2.000) (-2.450,-2.044) (-2.491,-2.100) (-2.525,-2.147) (-2.561,-2.200) (-2.594,-2.250) (-2.625,-2.299) (-2.656,-2.350) (-2.685,-2.400) (-2.713,-2.450) (-2.740,-2.500) (-2.765,-2.550) (-2.790,-2.600) (-2.813,-2.650) (-2.835,-2.700) (-2.856,-2.750) (-2.875,-2.799) (-2.894,-2.850) (-2.912,-2.900) (-2.928,-2.950) (-2.944,-3.000) (-2.958,-3.050) (-2.972,-3.100) (-2.984,-3.150) (-2.995,-3.200) (-3.005,-3.250) (-3.015,-3.300) (-3.023,-3.350) (-3.030,-3.400) (-3.036,-3.450) (-3.041,-3.500) (-3.045,-3.550) (-3.049,-3.600) (-3.051,-3.650) (-3.052,-3.700) (-3.052,-3.750) (-3.052,-3.800) (-3.050,-3.850) (-3.047,-3.900) (-3.044,-3.950) (-3.039,-4.000) (-3.033,-4.050) (-3.027,-4.100) (-3.019,-4.150) (-3.010,-4.200) (-3.001,-4.250) (-2.990,-4.300) (-2.978,-4.350) (-2.965,-4.400) (-2.951,-4.450) (-2.937,-4.500) (-2.921,-4.550) (-2.903,-4.600) (-2.885,-4.650) (-2.866,-4.700) (-2.845,-4.750) (-2.825,-4.796) (-2.800,-4.850) (-2.776,-4.900) (-2.751,-4.950) (-2.725,-4.998) (-2.700,-5.042) (-2.666,-5.100) (-2.635,-5.150) (-2.602,-5.200) (-2.568,-5.250) (-2.533,-5.300) (-2.500,-5.344) (-2.456,-5.400) (-2.415,-5.450) (-2.375,-5.497) (-2.328,-5.550) (-2.281,-5.600) (-2.231,-5.650) (-2.180,-5.700) (-2.125,-5.750) (-2.075,-5.794) (-2.025,-5.836) (-1.977,-5.875) (-1.925,-5.915) (-1.878,-5.950) (-1.825,-5.987) (-1.775,-6.021) (-1.731,-6.050) (-1.675,-6.085) (-1.625,-6.115) (-1.575,-6.143) (-1.525,-6.170) (-1.475,-6.197) (-1.425,-6.222) (-1.375,-6.246) (-1.325,-6.269) (-1.275,-6.291) (-1.225,-6.311) (-1.175,-6.331) (-1.126,-6.350) (-1.075,-6.368) (-1.025,-6.385) (-0.975,-6.401) (-0.925,-6.417) (-0.875,-6.431) (-0.825,-6.445) (-0.775,-6.457) (-0.725,-6.469) (-0.675,-6.480) (-0.625,-6.490) (-0.575,-6.499) (-0.525,-6.508) (-0.475,-6.516) (-0.425,-6.523) (-0.375,-6.529) (-0.325,-6.534) (-0.275,-6.539) (-0.225,-6.543) (-0.175,-6.546) (-0.125,-6.548) (-0.075,-6.549) (-0.025,-6.550) (0.047,-6.550) (0.100,-6.549) (0.150,-6.547) (0.200,-6.544) (0.250,-6.541) (0.300,-6.537) (0.350,-6.532) (0.400,-6.526) (0.450,-6.519) (0.500,-6.512) (0.550,-6.504) (0.600,-6.495) (0.650,-6.485) (0.700,-6.475) (0.750,-6.463) (0.800,-6.451) (0.850,-6.438) (0.900,-6.424) (0.950,-6.409) (1.000,-6.393) (1.050,-6.377) (1.100,-6.359) (1.150,-6.341) (1.200,-6.322) (1.250,-6.301) (1.300,-6.280) (1.350,-6.257) (1.400,-6.234) (1.450,-6.209) (1.500,-6.184) (1.550,-6.157) (1.600,-6.129) (1.650,-6.100) (1.700,-6.069) (1.750,-6.038) (1.800,-6.004) (1.842,-5.975) (1.900,-5.934) (1.945,-5.900) (2.000,-5.857) (2.050,-5.816) (2.097,-5.775) (2.150,-5.728) (2.200,-5.681) (2.250,-5.632) (2.300,-5.580) (2.350,-5.525) (2.394,-5.475) (2.436,-5.425) (2.475,-5.376) (2.514,-5.325) (2.550,-5.276) (2.586,-5.225) (2.619,-5.175) (2.650,-5.126) (2.681,-5.075) (2.710,-5.025) (2.737,-4.975) (2.763,-4.925) (2.788,-4.875) (2.812,-4.825) (2.834,-4.775) (2.856,-4.725) (2.875,-4.676) (2.894,-4.625) (2.912,-4.575) (2.929,-4.525) (2.944,-4.475) (2.959,-4.425) (2.972,-4.375) (2.984,-4.325) (2.995,-4.275) (3.006,-4.225) (3.015,-4.175) (3.023,-4.125) (3.030,-4.075) (3.036,-4.025) (3.041,-3.975) (3.046,-3.925) (3.049,-3.875) (3.051,-3.825) (3.052,-3.775) (3.052,-3.725) (3.052,-3.675) (3.050,-3.625) (3.047,-3.575) (3.044,-3.525) (3.039,-3.475) (3.033,-3.425) (3.026,-3.375) (3.019,-3.325) (3.010,-3.275) (3.000,-3.225) (2.990,-3.175) (2.978,-3.125) (2.965,-3.075) (2.951,-3.025) (2.936,-2.975) (2.920,-2.925) (2.903,-2.875) (2.885,-2.825) (2.866,-2.775) (2.845,-2.725) (2.825,-2.677) (2.801,-2.625) (2.778,-2.575) (2.753,-2.525) (2.726,-2.475) (2.700,-2.427) (2.675,-2.383) (2.641,-2.325) (2.610,-2.275) (2.578,-2.225) (2.544,-2.175) (2.509,-2.125) (2.475,-2.078) (2.436,-2.025) (2.400,-1.979) (2.357,-1.925) (2.315,-1.875) (2.275,-1.828) (2.228,-1.775) (2.183,-1.725) (2.136,-1.675) (2.100,-1.638) (2.063,-1.600) (2.025,-1.562) (1.987,-1.525) (1.950,-1.489) (1.909,-1.450) (1.854,-1.400) (1.800,-1.351) (1.750,-1.307) (1.700,-1.264) (1.655,-1.225) (1.600,-1.179) (1.550,-1.138) (1.504,-1.100) (1.450,-1.057) (1.410,-1.025) (1.350,-0.978) (1.300,-0.940) (1.250,-0.901) (1.200,-0.863) (1.150,-0.826) (1.100,-0.788) (1.050,-0.751) (1.000,-0.714) (0.950,-0.678) (0.900,-0.641) (0.850,-0.605) (0.809,-0.575) (0.750,-0.533) (0.705,-0.500) (0.650,-0.461) (0.600,-0.425) (0.550,-0.390) (0.500,-0.354) (0.450,-0.319) (0.400,-0.284) (0.354,-0.250) (0.300,-0.213) (0.250,-0.177) (0.206,-0.144)}
\EQ{(0.206,0.144) (0.248,0.175) (0.300,0.213) (0.350,0.248) (0.400,0.284) (0.450,0.319) (0.495,0.350) (0.550,0.390) (0.600,0.425) (0.650,0.461) (0.700,0.497) (0.750,0.533) (0.800,0.569) (0.844,0.600) (0.900,0.641) (0.947,0.675) (1.000,0.714) (1.049,0.750) (1.100,0.788) (1.149,0.825) (1.200,0.863) (1.249,0.900) (1.300,0.940) (1.346,0.975) (1.400,1.018) (1.441,1.050) (1.500,1.097) (1.550,1.138) (1.595,1.175) (1.650,1.221) (1.700,1.264) (1.742,1.300) (1.799,1.350) (1.850,1.396) (1.900,1.442) (1.950,1.489) (1.987,1.525) (2.025,1.562) (2.063,1.600) (2.100,1.638) (2.136,1.675) (2.175,1.717) (2.225,1.771) (2.273,1.825) (2.315,1.875) (2.350,1.917) (2.397,1.975) (2.436,2.025) (2.473,2.075) (2.509,2.125) (2.544,2.175) (2.575,2.221) (2.610,2.275) (2.641,2.325) (2.671,2.375) (2.699,2.425) (2.725,2.472) (2.750,2.520) (2.775,2.570) (2.800,2.622) (2.824,2.675) (2.845,2.725) (2.866,2.775) (2.885,2.825) (2.903,2.875) (2.920,2.925) (2.936,2.975) (2.951,3.025) (2.965,3.075) (2.978,3.125) (2.990,3.175) (3.000,3.225) (3.010,3.275) (3.019,3.325) (3.026,3.375) (3.033,3.425) (3.039,3.475) (3.044,3.525) (3.047,3.575) (3.050,3.625) (3.052,3.675) (3.052,3.725) (3.052,3.775) (3.051,3.825) (3.049,3.875) (3.046,3.925) (3.041,3.975) (3.036,4.025) (3.030,4.075) (3.023,4.125) (3.015,4.175) (3.006,4.225) (2.995,4.275) (2.984,4.325) (2.972,4.375) (2.959,4.425) (2.944,4.475) (2.929,4.525) (2.912,4.575) (2.894,4.625) (2.876,4.675) (2.856,4.725) (2.834,4.775) (2.812,4.825) (2.788,4.875) (2.763,4.925) (2.737,4.975) (2.710,5.025) (2.681,5.075) (2.651,5.125) (2.619,5.175) (2.586,5.225) (2.551,5.275) (2.514,5.325) (2.476,5.375) (2.436,5.425) (2.400,5.468) (2.350,5.525) (2.305,5.575) (2.256,5.625) (2.206,5.675) (2.153,5.725) (2.100,5.773) (2.050,5.816) (2.008,5.850) (1.950,5.896) (1.900,5.934) (1.850,5.970) (1.806,6.000) (1.750,6.038) (1.700,6.069) (1.650,6.100) (1.600,6.129) (1.550,6.157) (1.500,6.184) (1.450,6.209) (1.400,6.234) (1.350,6.257) (1.300,6.280) (1.253,6.300) (1.200,6.322) (1.150,6.341) (1.100,6.359) (1.050,6.377) (1.000,6.393) (0.950,6.409) (0.900,6.424) (0.850,6.438) (0.800,6.451) (0.750,6.463) (0.700,6.475) (0.650,6.485) (0.600,6.495) (0.550,6.504) (0.500,6.512) (0.450,6.519) (0.400,6.526) (0.350,6.532) (0.300,6.537) (0.250,6.541) (0.200,6.544) (0.150,6.547) (0.100,6.549) (0.050,6.550) (0.000,6.550) (-0.050,6.550) (-0.100,6.549) (-0.150,6.547) (-0.200,6.544) (-0.250,6.541) (-0.300,6.537) (-0.350,6.532) (-0.400,6.526) (-0.450,6.519) (-0.500,6.512) (-0.550,6.504) (-0.600,6.495) (-0.650,6.485) (-0.700,6.475) (-0.750,6.463) (-0.800,6.451) (-0.850,6.438) (-0.900,6.424) (-0.950,6.409) (-1.000,6.393) (-1.050,6.377) (-1.100,6.359) (-1.150,6.341) (-1.200,6.322) (-1.250,6.301) (-1.300,6.280) (-1.350,6.257) (-1.400,6.234) (-1.450,6.209) (-1.500,6.184) (-1.550,6.157) (-1.600,6.129) (-1.650,6.100) (-1.700,6.069) (-1.750,6.038) (-1.800,6.004) (-1.842,5.975) (-1.900,5.934) (-1.945,5.900) (-2.000,5.857) (-2.050,5.816) (-2.097,5.775) (-2.150,5.728) (-2.200,5.681) (-2.250,5.632) (-2.300,5.580) (-2.350,5.525) (-2.394,5.475) (-2.436,5.425) (-2.475,5.376) (-2.514,5.325) (-2.550,5.276) (-2.586,5.225) (-2.619,5.175) (-2.650,5.126) (-2.681,5.075) (-2.710,5.025) (-2.737,4.975) (-2.763,4.925) (-2.788,4.875) (-2.812,4.825) (-2.834,4.775) (-2.856,4.725) (-2.875,4.676) (-2.894,4.625) (-2.912,4.575) (-2.929,4.525) (-2.944,4.475) (-2.959,4.425) (-2.972,4.375) (-2.984,4.325) (-2.995,4.275) (-3.006,4.225) (-3.015,4.175) (-3.023,4.125) (-3.030,4.075) (-3.036,4.025) (-3.041,3.975) (-3.046,3.925) (-3.049,3.875) (-3.051,3.825) (-3.052,3.775) (-3.052,3.725) (-3.052,3.675) (-3.050,3.625) (-3.047,3.575) (-3.044,3.525) (-3.039,3.475) (-3.033,3.425) (-3.026,3.375) (-3.019,3.325) (-3.010,3.275) (-3.000,3.225) (-2.990,3.175) (-2.978,3.125) (-2.965,3.075) (-2.951,3.025) (-2.936,2.975) (-2.920,2.925) (-2.903,2.875) (-2.885,2.825) (-2.866,2.775) (-2.845,2.725) (-2.825,2.677) (-2.801,2.625) (-2.778,2.575) (-2.753,2.525) (-2.726,2.475) (-2.700,2.427) (-2.675,2.383) (-2.641,2.325) (-2.610,2.275) (-2.578,2.225) (-2.544,2.175) (-2.509,2.125) (-2.475,2.078) (-2.436,2.025) (-2.400,1.979) (-2.357,1.925) (-2.315,1.875) (-2.275,1.828) (-2.228,1.775) (-2.183,1.725) (-2.136,1.675) (-2.100,1.638) (-2.063,1.600) (-2.025,1.562) (-1.987,1.525) (-1.950,1.489) (-1.909,1.450) (-1.854,1.400) (-1.800,1.351) (-1.750,1.307) (-1.700,1.264) (-1.655,1.225) (-1.600,1.179) (-1.550,1.138) (-1.504,1.100) (-1.450,1.057) (-1.410,1.025) (-1.350,0.978) (-1.300,0.940) (-1.250,0.901) (-1.200,0.863) (-1.150,0.826) (-1.100,0.788) (-1.050,0.751) (-1.000,0.714) (-0.950,0.678) (-0.900,0.641) (-0.850,0.605) (-0.809,0.575) (-0.750,0.533) (-0.705,0.500) (-0.650,0.461) (-0.600,0.425) (-0.550,0.390) (-0.500,0.354) (-0.450,0.319) (-0.400,0.284) (-0.354,0.250) (-0.300,0.213) (-0.250,0.177) (-0.220,0.155)}
\EQ{(-0.220,-0.155) (-0.275,-0.195) (-0.319,-0.225) (-0.375,-0.266) (-0.424,-0.300) (-0.475,-0.337) (-0.525,-0.372) (-0.575,-0.408) (-0.625,-0.444) (-0.670,-0.475) (-0.725,-0.515) (-0.774,-0.550) (-0.825,-0.587) (-0.875,-0.624) (-0.925,-0.660) (-0.975,-0.697) (-1.025,-0.734) (-1.075,-0.771) (-1.125,-0.809) (-1.175,-0.847) (-1.225,-0.885) (-1.275,-0.924) (-1.325,-0.963) (-1.372,-1.000) (-1.425,-1.043) (-1.465,-1.075) (-1.525,-1.124) (-1.575,-1.166) (-1.615,-1.200) (-1.672,-1.250) (-1.725,-1.297) (-1.775,-1.342) (-1.825,-1.389) (-1.863,-1.425) (-1.900,-1.461) (-1.940,-1.500) (-1.990,-1.550) (-2.025,-1.587) (-2.061,-1.625) (-2.100,-1.668) (-2.150,-1.724) (-2.194,-1.775) (-2.236,-1.825) (-2.275,-1.874) (-2.315,-1.925) (-2.350,-1.972) (-2.388,-2.025) (-2.423,-2.075) (-2.456,-2.125) (-2.488,-2.175) (-2.519,-2.225) (-2.548,-2.275) (-2.575,-2.324) (-2.600,-2.370) (-2.625,-2.419) (-2.650,-2.470) (-2.675,-2.524) (-2.697,-2.575) (-2.718,-2.625) (-2.737,-2.675) (-2.756,-2.725) (-2.773,-2.775) (-2.789,-2.825) (-2.804,-2.875) (-2.818,-2.925) (-2.830,-2.975) (-2.842,-3.025) (-2.852,-3.075) (-2.862,-3.125) (-2.870,-3.175) (-2.877,-3.225) (-2.884,-3.275) (-2.889,-3.325) (-2.893,-3.375) (-2.896,-3.425) (-2.898,-3.475) (-2.899,-3.525) (-2.899,-3.575) (-2.898,-3.625) (-2.896,-3.675) (-2.893,-3.725) (-2.889,-3.775) (-2.884,-3.825) (-2.878,-3.875) (-2.871,-3.925) (-2.862,-3.975) (-2.853,-4.025) (-2.842,-4.075) (-2.831,-4.125) (-2.818,-4.175) (-2.804,-4.225) (-2.789,-4.275) (-2.773,-4.325) (-2.756,-4.375) (-2.738,-4.425) (-2.718,-4.475) (-2.697,-4.525) (-2.675,-4.575) (-2.652,-4.625) (-2.627,-4.675) (-2.601,-4.725) (-2.575,-4.772) (-2.550,-4.815) (-2.514,-4.875) (-2.482,-4.925) (-2.450,-4.972) (-2.413,-5.025) (-2.376,-5.075) (-2.337,-5.125) (-2.300,-5.170) (-2.253,-5.225) (-2.208,-5.275) (-2.161,-5.325) (-2.125,-5.362) (-2.086,-5.400) (-2.032,-5.450) (-1.976,-5.500) (-1.925,-5.543) (-1.885,-5.575) (-1.825,-5.621) (-1.775,-5.658) (-1.725,-5.693) (-1.678,-5.725) (-1.625,-5.759) (-1.575,-5.790) (-1.525,-5.819) (-1.475,-5.847) (-1.425,-5.874) (-1.375,-5.900) (-1.325,-5.924) (-1.275,-5.947) (-1.225,-5.970) (-1.175,-5.991) (-1.125,-6.011) (-1.075,-6.030) (-1.025,-6.048) (-0.975,-6.065) (-0.925,-6.081) (-0.875,-6.096) (-0.825,-6.111) (-0.775,-6.124) (-0.725,-6.137) (-0.675,-6.148) (-0.625,-6.159) (-0.575,-6.169) (-0.525,-6.178) (-0.475,-6.186) (-0.425,-6.193) (-0.375,-6.200) (-0.325,-6.205) (-0.275,-6.210) (-0.225,-6.214) (-0.175,-6.217) (-0.125,-6.220) (-0.075,-6.221) (-0.025,-6.222) (0.050,-6.222) (0.100,-6.221) (0.150,-6.219) (0.200,-6.216) (0.250,-6.212) (0.300,-6.208) (0.350,-6.203) (0.400,-6.196) (0.450,-6.190) (0.500,-6.182) (0.550,-6.173) (0.600,-6.164) (0.650,-6.154) (0.700,-6.142) (0.750,-6.130) (0.800,-6.117) (0.850,-6.104) (0.900,-6.089) (0.950,-6.073) (1.000,-6.057) (1.050,-6.039) (1.100,-6.020) (1.150,-6.001) (1.200,-5.980) (1.250,-5.959) (1.300,-5.936) (1.350,-5.912) (1.400,-5.887) (1.450,-5.861) (1.500,-5.833) (1.550,-5.805) (1.599,-5.775) (1.650,-5.743) (1.700,-5.710) (1.750,-5.676) (1.800,-5.640) (1.850,-5.602) (1.900,-5.563) (1.946,-5.525) (2.000,-5.479) (2.050,-5.434) (2.100,-5.386) (2.137,-5.350) (2.175,-5.311) (2.225,-5.257) (2.275,-5.200) (2.317,-5.150) (2.350,-5.108) (2.394,-5.050) (2.425,-5.008) (2.465,-4.950) (2.498,-4.900) (2.525,-4.857) (2.559,-4.800) (2.587,-4.750) (2.614,-4.700) (2.639,-4.650) (2.663,-4.600) (2.686,-4.550) (2.708,-4.500) (2.728,-4.450) (2.747,-4.400) (2.765,-4.350) (2.781,-4.300) (2.797,-4.250) (2.811,-4.200) (2.825,-4.150) (2.837,-4.100) (2.848,-4.050) (2.858,-4.000) (2.867,-3.950) (2.874,-3.900) (2.881,-3.850) (2.887,-3.800) (2.891,-3.750) (2.895,-3.700) (2.898,-3.650) (2.899,-3.600) (2.900,-3.550) (2.899,-3.500) (2.897,-3.450) (2.895,-3.400) (2.891,-3.350) (2.887,-3.300) (2.881,-3.250) (2.874,-3.200) (2.866,-3.150) (2.857,-3.100) (2.847,-3.050) (2.836,-3.000) (2.824,-2.950) (2.811,-2.900) (2.796,-2.850) (2.781,-2.800) (2.764,-2.750) (2.747,-2.700) (2.728,-2.650) (2.708,-2.600) (2.687,-2.550) (2.664,-2.500) (2.640,-2.450) (2.616,-2.400) (2.589,-2.350) (2.562,-2.300) (2.533,-2.250) (2.503,-2.200) (2.475,-2.155) (2.440,-2.100) (2.406,-2.050) (2.375,-2.007) (2.334,-1.950) (2.300,-1.906) (2.256,-1.850) (2.215,-1.800) (2.175,-1.752) (2.129,-1.700) (2.084,-1.650) (2.050,-1.613) (2.014,-1.575) (1.975,-1.535) (1.925,-1.485) (1.889,-1.450) (1.850,-1.413) (1.811,-1.375) (1.756,-1.325) (1.700,-1.275) (1.650,-1.231) (1.600,-1.188) (1.556,-1.150) (1.500,-1.104) (1.450,-1.063) (1.404,-1.025) (1.350,-0.983) (1.309,-0.950) (1.250,-0.904) (1.200,-0.866) (1.150,-0.828) (1.100,-0.790) (1.050,-0.753) (1.000,-0.716) (0.950,-0.679) (0.900,-0.642) (0.850,-0.606) (0.809,-0.575) (0.750,-0.533) (0.705,-0.500) (0.650,-0.461) (0.600,-0.425) (0.550,-0.390) (0.500,-0.354) (0.450,-0.319) (0.400,-0.284) (0.354,-0.250) (0.300,-0.213) (0.250,-0.177) (0.206,-0.144)}
\EQ{(0.206,0.144) (0.248,0.175) (0.300,0.213) (0.350,0.248) (0.400,0.284) (0.450,0.319) (0.495,0.350) (0.550,0.390) (0.600,0.425) (0.650,0.461) (0.700,0.497) (0.750,0.533) (0.800,0.569) (0.843,0.600) (0.900,0.642) (0.946,0.675) (1.000,0.716) (1.047,0.750) (1.100,0.790) (1.147,0.825) (1.200,0.866) (1.245,0.900) (1.300,0.943) (1.341,0.975) (1.400,1.022) (1.450,1.063) (1.496,1.100) (1.550,1.145) (1.600,1.188) (1.644,1.225) (1.700,1.275) (1.750,1.320) (1.800,1.366) (1.837,1.400) (1.875,1.436) (1.915,1.475) (1.965,1.525) (2.000,1.561) (2.038,1.600) (2.075,1.640) (2.125,1.695) (2.173,1.750) (2.215,1.800) (2.250,1.843) (2.296,1.900) (2.334,1.950) (2.370,2.000) (2.400,2.042) (2.440,2.100) (2.472,2.150) (2.500,2.195) (2.533,2.250) (2.562,2.300) (2.589,2.350) (2.616,2.400) (2.640,2.450) (2.664,2.500) (2.687,2.550) (2.708,2.600) (2.728,2.650) (2.747,2.700) (2.764,2.750) (2.781,2.800) (2.796,2.850) (2.811,2.900) (2.824,2.950) (2.836,3.000) (2.847,3.050) (2.857,3.100) (2.866,3.150) (2.874,3.200) (2.881,3.250) (2.887,3.300) (2.891,3.350) (2.895,3.400) (2.897,3.450) (2.899,3.500) (2.900,3.550) (2.899,3.600) (2.898,3.650) (2.895,3.700) (2.891,3.750) (2.887,3.800) (2.881,3.850) (2.874,3.900) (2.867,3.950) (2.858,4.000) (2.848,4.050) (2.837,4.100) (2.825,4.150) (2.811,4.200) (2.797,4.250) (2.781,4.300) (2.765,4.350) (2.747,4.400) (2.728,4.450) (2.708,4.500) (2.686,4.550) (2.663,4.600) (2.639,4.650) (2.614,4.700) (2.587,4.750) (2.559,4.800) (2.529,4.850) (2.500,4.897) (2.465,4.950) (2.431,5.000) (2.400,5.043) (2.357,5.100) (2.325,5.140) (2.275,5.200) (2.231,5.250) (2.185,5.300) (2.150,5.337) (2.112,5.375) (2.075,5.410) (2.032,5.450) (1.976,5.500) (1.925,5.543) (1.885,5.575) (1.825,5.621) (1.775,5.658) (1.725,5.693) (1.678,5.725) (1.625,5.759) (1.575,5.790) (1.525,5.819) (1.475,5.847) (1.425,5.874) (1.375,5.900) (1.325,5.924) (1.275,5.947) (1.225,5.970) (1.175,5.991) (1.125,6.011) (1.075,6.030) (1.025,6.048) (0.975,6.065) (0.925,6.081) (0.875,6.096) (0.825,6.111) (0.775,6.124) (0.725,6.137) (0.675,6.148) (0.625,6.159) (0.575,6.169) (0.525,6.178) (0.475,6.186) (0.425,6.193) (0.375,6.200) (0.325,6.205) (0.275,6.210) (0.225,6.214) (0.175,6.217) (0.125,6.220) (0.075,6.221) (0.025,6.222) (-0.050,6.222) (-0.100,6.221) (-0.150,6.219) (-0.200,6.216) (-0.250,6.212) (-0.300,6.208) (-0.350,6.203) (-0.400,6.196) (-0.450,6.190) (-0.500,6.182) (-0.550,6.173) (-0.600,6.164) (-0.650,6.154) (-0.700,6.142) (-0.750,6.130) (-0.800,6.117) (-0.850,6.104) (-0.900,6.089) (-0.950,6.073) (-1.000,6.057) (-1.050,6.039) (-1.100,6.020) (-1.150,6.001) (-1.200,5.980) (-1.250,5.959) (-1.300,5.936) (-1.350,5.912) (-1.400,5.887) (-1.450,5.861) (-1.500,5.833) (-1.550,5.805) (-1.599,5.775) (-1.650,5.743) (-1.700,5.710) (-1.750,5.676) (-1.800,5.640) (-1.850,5.602) (-1.900,5.563) (-1.946,5.525) (-2.000,5.479) (-2.050,5.434) (-2.100,5.386) (-2.137,5.350) (-2.175,5.311) (-2.225,5.257) (-2.275,5.200) (-2.317,5.150) (-2.350,5.108) (-2.394,5.050) (-2.425,5.008) (-2.465,4.950) (-2.498,4.900) (-2.525,4.857) (-2.559,4.800) (-2.587,4.750) (-2.614,4.700) (-2.639,4.650) (-2.663,4.600) (-2.686,4.550) (-2.708,4.500) (-2.728,4.450) (-2.747,4.400) (-2.765,4.350) (-2.781,4.300) (-2.797,4.250) (-2.811,4.200) (-2.825,4.150) (-2.837,4.100) (-2.848,4.050) (-2.858,4.000) (-2.867,3.950) (-2.874,3.900) (-2.881,3.850) (-2.887,3.800) (-2.891,3.750) (-2.895,3.700) (-2.898,3.650) (-2.899,3.600) (-2.900,3.550) (-2.899,3.500) (-2.897,3.450) (-2.895,3.400) (-2.891,3.350) (-2.887,3.300) (-2.881,3.250) (-2.874,3.200) (-2.866,3.150) (-2.857,3.100) (-2.847,3.050) (-2.836,3.000) (-2.824,2.950) (-2.811,2.900) (-2.796,2.850) (-2.781,2.800) (-2.764,2.750) (-2.747,2.700) (-2.728,2.650) (-2.708,2.600) (-2.687,2.550) (-2.664,2.500) (-2.640,2.450) (-2.616,2.400) (-2.589,2.350) (-2.562,2.300) (-2.533,2.250) (-2.503,2.200) (-2.475,2.155) (-2.440,2.100) (-2.406,2.050) (-2.375,2.007) (-2.334,1.950) (-2.300,1.906) (-2.256,1.850) (-2.215,1.800) (-2.175,1.752) (-2.129,1.700) (-2.084,1.650) (-2.050,1.613) (-2.014,1.575) (-1.975,1.535) (-1.925,1.485) (-1.889,1.450) (-1.850,1.413) (-1.811,1.375) (-1.756,1.325) (-1.700,1.275) (-1.650,1.231) (-1.600,1.188) (-1.556,1.150) (-1.500,1.104) (-1.450,1.063) (-1.404,1.025) (-1.350,0.983) (-1.309,0.950) (-1.250,0.904) (-1.200,0.866) (-1.150,0.828) (-1.100,0.790) (-1.050,0.753) (-1.000,0.716) (-0.950,0.679) (-0.900,0.642) (-0.850,0.606) (-0.809,0.575) (-0.750,0.533) (-0.705,0.500) (-0.650,0.461) (-0.600,0.425) (-0.550,0.390) (-0.500,0.354) (-0.450,0.319) (-0.400,0.284) (-0.354,0.250) (-0.300,0.213) (-0.250,0.177) (-0.220,0.155)}
\EQ{(-0.220,-0.155) (-0.275,-0.195) (-0.319,-0.225) (-0.375,-0.266) (-0.424,-0.300) (-0.475,-0.337) (-0.525,-0.372) (-0.575,-0.408) (-0.625,-0.444) (-0.670,-0.475) (-0.725,-0.516) (-0.773,-0.550) (-0.825,-0.588) (-0.875,-0.624) (-0.925,-0.661) (-0.975,-0.698) (-1.025,-0.735) (-1.075,-0.773) (-1.125,-0.811) (-1.175,-0.849) (-1.225,-0.888) (-1.273,-0.925) (-1.325,-0.967) (-1.367,-1.000) (-1.425,-1.048) (-1.475,-1.089) (-1.518,-1.125) (-1.575,-1.174) (-1.625,-1.218) (-1.675,-1.263) (-1.715,-1.300) (-1.769,-1.350) (-1.821,-1.400) (-1.871,-1.450) (-1.920,-1.500) (-1.967,-1.550) (-2.013,-1.600) (-2.050,-1.641) (-2.100,-1.699) (-2.142,-1.750) (-2.175,-1.791) (-2.221,-1.850) (-2.258,-1.900) (-2.293,-1.950) (-2.325,-1.996) (-2.361,-2.050) (-2.392,-2.100) (-2.422,-2.150) (-2.450,-2.199) (-2.475,-2.245) (-2.500,-2.292) (-2.525,-2.343) (-2.550,-2.396) (-2.574,-2.450) (-2.595,-2.500) (-2.615,-2.550) (-2.633,-2.600) (-2.650,-2.648) (-2.667,-2.700) (-2.682,-2.750) (-2.696,-2.800) (-2.708,-2.850) (-2.720,-2.900) (-2.730,-2.950) (-2.739,-3.000) (-2.747,-3.050) (-2.754,-3.100) (-2.760,-3.150) (-2.765,-3.200) (-2.769,-3.250) (-2.771,-3.300) (-2.773,-3.350) (-2.773,-3.400) (-2.773,-3.450) (-2.771,-3.500) (-2.768,-3.550) (-2.764,-3.600) (-2.759,-3.650) (-2.753,-3.700) (-2.746,-3.750) (-2.738,-3.800) (-2.729,-3.850) (-2.718,-3.900) (-2.706,-3.950) (-2.694,-4.000) (-2.680,-4.050) (-2.664,-4.100) (-2.648,-4.150) (-2.630,-4.200) (-2.612,-4.250) (-2.592,-4.300) (-2.570,-4.350) (-2.547,-4.400) (-2.525,-4.446) (-2.500,-4.496) (-2.475,-4.542) (-2.450,-4.587) (-2.412,-4.650) (-2.381,-4.700) (-2.350,-4.746) (-2.312,-4.800) (-2.275,-4.850) (-2.237,-4.900) (-2.200,-4.945) (-2.153,-5.000) (-2.108,-5.050) (-2.061,-5.100) (-2.025,-5.136) (-1.984,-5.175) (-1.930,-5.225) (-1.875,-5.273) (-1.825,-5.315) (-1.781,-5.350) (-1.725,-5.392) (-1.679,-5.425) (-1.625,-5.462) (-1.575,-5.495) (-1.527,-5.525) (-1.475,-5.556) (-1.425,-5.585) (-1.375,-5.612) (-1.325,-5.638) (-1.275,-5.662) (-1.225,-5.686) (-1.175,-5.708) (-1.125,-5.729) (-1.075,-5.749) (-1.025,-5.769) (-0.975,-5.787) (-0.925,-5.804) (-0.875,-5.819) (-0.825,-5.834) (-0.775,-5.848) (-0.725,-5.862) (-0.675,-5.874) (-0.625,-5.885) (-0.575,-5.895) (-0.525,-5.905) (-0.475,-5.913) (-0.425,-5.921) (-0.375,-5.928) (-0.325,-5.934) (-0.275,-5.939) (-0.225,-5.943) (-0.175,-5.946) (-0.125,-5.949) (-0.075,-5.950) (-0.025,-5.951) (0.050,-5.951) (0.100,-5.950) (0.150,-5.948) (0.200,-5.945) (0.250,-5.941) (0.300,-5.936) (0.350,-5.931) (0.400,-5.924) (0.450,-5.917) (0.500,-5.909) (0.550,-5.900) (0.600,-5.890) (0.650,-5.879) (0.700,-5.868) (0.750,-5.855) (0.800,-5.842) (0.850,-5.827) (0.900,-5.812) (0.950,-5.795) (1.000,-5.778) (1.050,-5.759) (1.100,-5.740) (1.150,-5.719) (1.200,-5.697) (1.249,-5.675) (1.300,-5.650) (1.350,-5.625) (1.397,-5.600) (1.442,-5.575) (1.486,-5.550) (1.550,-5.511) (1.600,-5.479) (1.643,-5.450) (1.700,-5.410) (1.748,-5.375) (1.800,-5.335) (1.843,-5.300) (1.900,-5.252) (1.950,-5.207) (2.000,-5.160) (2.036,-5.125) (2.075,-5.085) (2.125,-5.031) (2.175,-4.975) (2.217,-4.925) (2.250,-4.883) (2.294,-4.825) (2.325,-4.782) (2.364,-4.725) (2.397,-4.675) (2.425,-4.629) (2.457,-4.575) (2.484,-4.525) (2.511,-4.475) (2.535,-4.425) (2.559,-4.375) (2.581,-4.325) (2.602,-4.275) (2.621,-4.225) (2.639,-4.175) (2.656,-4.125) (2.672,-4.075) (2.687,-4.025) (2.700,-3.975) (2.712,-3.925) (2.723,-3.875) (2.733,-3.825) (2.742,-3.775) (2.750,-3.725) (2.757,-3.675) (2.762,-3.625) (2.766,-3.575) (2.770,-3.525) (2.772,-3.475) (2.773,-3.425) (2.773,-3.375) (2.772,-3.325) (2.770,-3.275) (2.767,-3.225) (2.763,-3.175) (2.757,-3.125) (2.751,-3.075) (2.743,-3.025) (2.735,-2.975) (2.725,-2.926) (2.714,-2.875) (2.702,-2.825) (2.689,-2.775) (2.674,-2.725) (2.659,-2.675) (2.642,-2.625) (2.625,-2.577) (2.605,-2.525) (2.585,-2.475) (2.563,-2.425) (2.541,-2.375) (2.517,-2.325) (2.491,-2.275) (2.464,-2.225) (2.436,-2.175) (2.407,-2.125) (2.376,-2.075) (2.344,-2.025) (2.311,-1.975) (2.276,-1.925) (2.240,-1.875) (2.202,-1.825) (2.162,-1.775) (2.125,-1.729) (2.079,-1.675) (2.036,-1.625) (2.000,-1.586) (1.950,-1.532) (1.900,-1.480) (1.850,-1.429) (1.800,-1.380) (1.750,-1.333) (1.700,-1.286) (1.661,-1.250) (1.604,-1.200) (1.550,-1.153) (1.500,-1.110) (1.458,-1.075) (1.400,-1.027) (1.350,-0.987) (1.304,-0.950) (1.250,-0.907) (1.209,-0.875) (1.150,-0.830) (1.100,-0.792) (1.050,-0.754) (1.000,-0.717) (0.950,-0.680) (0.900,-0.643) (0.850,-0.606) (0.808,-0.575) (0.750,-0.534) (0.704,-0.500) (0.650,-0.462) (0.600,-0.425) (0.550,-0.390) (0.500,-0.354) (0.450,-0.319) (0.400,-0.284) (0.354,-0.250) (0.300,-0.213) (0.250,-0.177) (0.206,-0.144)}
\EQ{(0.206,0.144) (0.248,0.175) (0.300,0.213) (0.350,0.248) (0.400,0.284) (0.450,0.319) (0.495,0.350) (0.550,0.390) (0.599,0.425) (0.650,0.462) (0.700,0.497) (0.750,0.534) (0.800,0.570) (0.842,0.600) (0.900,0.643) (0.944,0.675) (1.000,0.717) (1.045,0.750) (1.100,0.792) (1.144,0.825) (1.200,0.868) (1.241,0.900) (1.300,0.947) (1.350,0.987) (1.398,1.025) (1.450,1.068) (1.500,1.110) (1.547,1.150) (1.600,1.196) (1.650,1.241) (1.688,1.275) (1.725,1.309) (1.769,1.350) (1.821,1.400) (1.871,1.450) (1.920,1.500) (1.967,1.550) (2.013,1.600) (2.050,1.641) (2.100,1.699) (2.142,1.750) (2.175,1.791) (2.221,1.850) (2.258,1.900) (2.293,1.950) (2.325,1.996) (2.361,2.050) (2.392,2.100) (2.422,2.150) (2.450,2.199) (2.475,2.245) (2.500,2.292) (2.525,2.343) (2.550,2.396) (2.574,2.450) (2.595,2.500) (2.615,2.550) (2.633,2.600) (2.650,2.648) (2.667,2.700) (2.682,2.750) (2.696,2.800) (2.708,2.850) (2.720,2.900) (2.730,2.950) (2.739,3.000) (2.747,3.050) (2.754,3.100) (2.760,3.150) (2.765,3.200) (2.769,3.250) (2.771,3.300) (2.773,3.350) (2.773,3.400) (2.773,3.450) (2.771,3.500) (2.768,3.550) (2.764,3.600) (2.759,3.650) (2.753,3.700) (2.746,3.750) (2.738,3.800) (2.729,3.850) (2.718,3.900) (2.706,3.950) (2.694,4.000) (2.680,4.050) (2.664,4.100) (2.648,4.150) (2.630,4.200) (2.612,4.250) (2.592,4.300) (2.570,4.350) (2.547,4.400) (2.525,4.446) (2.500,4.496) (2.475,4.542) (2.450,4.587) (2.412,4.650) (2.381,4.700) (2.350,4.746) (2.312,4.800) (2.275,4.850) (2.237,4.900) (2.200,4.945) (2.153,5.000) (2.108,5.050) (2.061,5.100) (2.025,5.136) (1.984,5.175) (1.930,5.225) (1.875,5.273) (1.825,5.315) (1.781,5.350) (1.725,5.392) (1.679,5.425) (1.625,5.462) (1.575,5.495) (1.527,5.525) (1.475,5.556) (1.425,5.585) (1.375,5.612) (1.325,5.638) (1.275,5.662) (1.225,5.686) (1.175,5.708) (1.125,5.729) (1.075,5.749) (1.025,5.769) (0.975,5.787) (0.925,5.804) (0.875,5.819) (0.825,5.834) (0.775,5.848) (0.725,5.862) (0.675,5.874) (0.625,5.885) (0.575,5.895) (0.525,5.905) (0.475,5.913) (0.425,5.921) (0.375,5.928) (0.325,5.934) (0.275,5.939) (0.225,5.943) (0.175,5.946) (0.125,5.949) (0.075,5.950) (0.025,5.951) (-0.050,5.951) (-0.100,5.950) (-0.150,5.948) (-0.200,5.945) (-0.250,5.941) (-0.300,5.936) (-0.350,5.931) (-0.400,5.924) (-0.450,5.917) (-0.500,5.909) (-0.550,5.900) (-0.600,5.890) (-0.650,5.879) (-0.700,5.868) (-0.750,5.855) (-0.800,5.842) (-0.850,5.827) (-0.900,5.812) (-0.950,5.795) (-1.000,5.778) (-1.050,5.759) (-1.100,5.740) (-1.150,5.719) (-1.200,5.697) (-1.249,5.675) (-1.300,5.650) (-1.350,5.625) (-1.397,5.600) (-1.442,5.575) (-1.486,5.550) (-1.550,5.511) (-1.600,5.479) (-1.643,5.450) (-1.700,5.410) (-1.748,5.375) (-1.800,5.335) (-1.843,5.300) (-1.900,5.252) (-1.950,5.207) (-2.000,5.160) (-2.036,5.125) (-2.075,5.085) (-2.125,5.031) (-2.175,4.975) (-2.217,4.925) (-2.250,4.883) (-2.294,4.825) (-2.325,4.782) (-2.364,4.725) (-2.397,4.675) (-2.425,4.629) (-2.457,4.575) (-2.484,4.525) (-2.511,4.475) (-2.535,4.425) (-2.559,4.375) (-2.581,4.325) (-2.602,4.275) (-2.621,4.225) (-2.639,4.175) (-2.656,4.125) (-2.672,4.075) (-2.687,4.025) (-2.700,3.975) (-2.712,3.925) (-2.723,3.875) (-2.733,3.825) (-2.742,3.775) (-2.750,3.725) (-2.757,3.675) (-2.762,3.625) (-2.766,3.575) (-2.770,3.525) (-2.772,3.475) (-2.773,3.425) (-2.773,3.375) (-2.772,3.325) (-2.770,3.275) (-2.767,3.225) (-2.763,3.175) (-2.757,3.125) (-2.751,3.075) (-2.743,3.025) (-2.735,2.975) (-2.725,2.926) (-2.714,2.875) (-2.702,2.825) (-2.689,2.775) (-2.674,2.725) (-2.659,2.675) (-2.642,2.625) (-2.625,2.577) (-2.605,2.525) (-2.585,2.475) (-2.563,2.425) (-2.541,2.375) (-2.517,2.325) (-2.491,2.275) (-2.464,2.225) (-2.436,2.175) (-2.407,2.125) (-2.376,2.075) (-2.344,2.025) (-2.311,1.975) (-2.276,1.925) (-2.240,1.875) (-2.202,1.825) (-2.162,1.775) (-2.125,1.729) (-2.079,1.675) (-2.036,1.625) (-2.000,1.586) (-1.950,1.532) (-1.900,1.480) (-1.850,1.429) (-1.800,1.380) (-1.750,1.333) (-1.700,1.286) (-1.661,1.250) (-1.604,1.200) (-1.550,1.153) (-1.500,1.110) (-1.458,1.075) (-1.400,1.027) (-1.350,0.987) (-1.304,0.950) (-1.250,0.907) (-1.209,0.875) (-1.150,0.830) (-1.100,0.792) (-1.050,0.754) (-1.000,0.717) (-0.950,0.680) (-0.900,0.643) (-0.850,0.606) (-0.808,0.575) (-0.750,0.534) (-0.704,0.500) (-0.650,0.462) (-0.600,0.425) (-0.550,0.390) (-0.500,0.354) (-0.450,0.319) (-0.400,0.284) (-0.354,0.250) (-0.300,0.213) (-0.250,0.177) (-0.220,0.155)}
\EQ{(-0.220,-0.155) (-0.275,-0.195) (-0.319,-0.225) (-0.375,-0.266) (-0.424,-0.300) (-0.475,-0.337) (-0.525,-0.372) (-0.575,-0.408) (-0.625,-0.444) (-0.669,-0.475) (-0.725,-0.516) (-0.773,-0.550) (-0.825,-0.589) (-0.875,-0.625) (-0.925,-0.662) (-0.975,-0.699) (-1.025,-0.737) (-1.075,-0.774) (-1.125,-0.813) (-1.173,-0.850) (-1.225,-0.891) (-1.269,-0.925) (-1.325,-0.970) (-1.375,-1.011) (-1.422,-1.050) (-1.475,-1.095) (-1.525,-1.139) (-1.567,-1.175) (-1.622,-1.225) (-1.675,-1.274) (-1.725,-1.322) (-1.775,-1.370) (-1.825,-1.421) (-1.875,-1.473) (-1.924,-1.525) (-1.969,-1.575) (-2.012,-1.625) (-2.050,-1.671) (-2.094,-1.725) (-2.125,-1.766) (-2.170,-1.825) (-2.200,-1.868) (-2.239,-1.925) (-2.272,-1.975) (-2.300,-2.020) (-2.333,-2.075) (-2.362,-2.125) (-2.389,-2.175) (-2.414,-2.225) (-2.439,-2.275) (-2.462,-2.325) (-2.483,-2.375) (-2.504,-2.425) (-2.523,-2.475) (-2.541,-2.525) (-2.558,-2.575) (-2.573,-2.625) (-2.587,-2.675) (-2.600,-2.725) (-2.612,-2.775) (-2.623,-2.825) (-2.632,-2.875) (-2.640,-2.925) (-2.647,-2.975) (-2.654,-3.025) (-2.658,-3.075) (-2.662,-3.125) (-2.665,-3.175) (-2.666,-3.225) (-2.667,-3.275) (-2.666,-3.325) (-2.664,-3.375) (-2.661,-3.425) (-2.657,-3.475) (-2.651,-3.525) (-2.645,-3.575) (-2.637,-3.625) (-2.629,-3.675) (-2.619,-3.725) (-2.608,-3.775) (-2.595,-3.825) (-2.582,-3.875) (-2.567,-3.925) (-2.551,-3.975) (-2.534,-4.025) (-2.515,-4.075) (-2.496,-4.125) (-2.475,-4.174) (-2.452,-4.225) (-2.428,-4.275) (-2.402,-4.325) (-2.375,-4.375) (-2.350,-4.420) (-2.317,-4.475) (-2.285,-4.525) (-2.251,-4.575) (-2.216,-4.625) (-2.179,-4.675) (-2.140,-4.725) (-2.100,-4.773) (-2.055,-4.825) (-2.009,-4.875) (-1.960,-4.925) (-1.909,-4.975) (-1.855,-5.025) (-1.800,-5.073) (-1.750,-5.114) (-1.704,-5.150) (-1.650,-5.191) (-1.602,-5.225) (-1.550,-5.260) (-1.500,-5.293) (-1.450,-5.323) (-1.405,-5.350) (-1.361,-5.375) (-1.314,-5.400) (-1.265,-5.425) (-1.214,-5.450) (-1.160,-5.475) (-1.103,-5.500) (-1.050,-5.522) (-1.000,-5.541) (-0.950,-5.559) (-0.900,-5.577) (-0.850,-5.593) (-0.800,-5.608) (-0.750,-5.622) (-0.700,-5.635) (-0.650,-5.647) (-0.600,-5.659) (-0.550,-5.669) (-0.500,-5.678) (-0.450,-5.687) (-0.400,-5.694) (-0.350,-5.701) (-0.300,-5.707) (-0.250,-5.711) (-0.200,-5.715) (-0.150,-5.718) (-0.100,-5.721) (-0.050,-5.722) (0.000,-5.722) (0.050,-5.722) (0.100,-5.721) (0.150,-5.718) (0.200,-5.715) (0.250,-5.711) (0.300,-5.707) (0.350,-5.701) (0.400,-5.694) (0.450,-5.687) (0.500,-5.678) (0.550,-5.669) (0.600,-5.659) (0.650,-5.647) (0.700,-5.635) (0.750,-5.622) (0.800,-5.608) (0.850,-5.593) (0.900,-5.577) (0.950,-5.559) (1.000,-5.541) (1.050,-5.522) (1.100,-5.501) (1.150,-5.480) (1.200,-5.457) (1.250,-5.433) (1.300,-5.407) (1.350,-5.381) (1.400,-5.353) (1.447,-5.325) (1.500,-5.293) (1.550,-5.260) (1.600,-5.226) (1.650,-5.191) (1.700,-5.153) (1.750,-5.114) (1.797,-5.075) (1.850,-5.029) (1.900,-4.983) (1.950,-4.935) (2.000,-4.884) (2.050,-4.830) (2.098,-4.775) (2.140,-4.725) (2.175,-4.680) (2.216,-4.625) (2.250,-4.577) (2.285,-4.525) (2.317,-4.475) (2.347,-4.425) (2.375,-4.376) (2.400,-4.330) (2.425,-4.281) (2.450,-4.229) (2.474,-4.175) (2.496,-4.125) (2.515,-4.075) (2.534,-4.025) (2.551,-3.975) (2.567,-3.925) (2.582,-3.875) (2.595,-3.825) (2.608,-3.775) (2.619,-3.725) (2.629,-3.675) (2.637,-3.625) (2.645,-3.575) (2.651,-3.525) (2.657,-3.475) (2.661,-3.425) (2.664,-3.375) (2.666,-3.325) (2.667,-3.275) (2.666,-3.225) (2.665,-3.175) (2.662,-3.125) (2.658,-3.075) (2.654,-3.025) (2.647,-2.975) (2.640,-2.925) (2.632,-2.875) (2.623,-2.825) (2.612,-2.775) (2.600,-2.725) (2.587,-2.675) (2.573,-2.625) (2.558,-2.575) (2.541,-2.525) (2.523,-2.475) (2.504,-2.425) (2.483,-2.375) (2.462,-2.325) (2.439,-2.275) (2.414,-2.225) (2.389,-2.175) (2.362,-2.125) (2.333,-2.075) (2.303,-2.025) (2.275,-1.980) (2.239,-1.925) (2.205,-1.875) (2.175,-1.833) (2.132,-1.775) (2.100,-1.733) (2.053,-1.675) (2.012,-1.625) (1.975,-1.583) (1.925,-1.527) (1.877,-1.475) (1.829,-1.425) (1.780,-1.375) (1.729,-1.325) (1.676,-1.275) (1.625,-1.228) (1.575,-1.183) (1.525,-1.139) (1.481,-1.100) (1.425,-1.053) (1.375,-1.011) (1.331,-0.975) (1.275,-0.930) (1.225,-0.891) (1.175,-0.851) (1.125,-0.813) (1.076,-0.775) (1.025,-0.737) (0.976,-0.700) (0.925,-0.662) (0.875,-0.625) (0.825,-0.589) (0.775,-0.552) (0.725,-0.516) (0.675,-0.480) (0.635,-0.450) (0.575,-0.408) (0.530,-0.375) (0.475,-0.337) (0.425,-0.301) (0.375,-0.266) (0.325,-0.230) (0.284,-0.200) (0.225,-0.160) (0.206,-0.144)}
\EQ{(0.206,0.144) (0.248,0.175) (0.300,0.213) (0.350,0.248) (0.400,0.284) (0.450,0.319) (0.495,0.350) (0.550,0.390) (0.599,0.425) (0.650,0.462) (0.700,0.498) (0.750,0.534) (0.800,0.570) (0.842,0.600) (0.900,0.644) (0.943,0.675) (1.000,0.718) (1.043,0.750) (1.100,0.794) (1.141,0.825) (1.200,0.871) (1.250,0.910) (1.300,0.950) (1.350,0.991) (1.392,1.025) (1.450,1.074) (1.500,1.117) (1.538,1.150) (1.595,1.200) (1.649,1.250) (1.700,1.298) (1.750,1.346) (1.800,1.395) (1.850,1.447) (1.900,1.499) (1.946,1.550) (1.990,1.600) (2.025,1.641) (2.074,1.700) (2.113,1.750) (2.150,1.799) (2.188,1.850) (2.222,1.900) (2.256,1.950) (2.288,2.000) (2.318,2.050) (2.348,2.100) (2.375,2.149) (2.400,2.197) (2.425,2.247) (2.450,2.299) (2.473,2.350) (2.494,2.400) (2.514,2.450) (2.532,2.500) (2.549,2.550) (2.565,2.600) (2.580,2.650) (2.594,2.700) (2.606,2.750) (2.617,2.800) (2.627,2.850) (2.636,2.900) (2.644,2.950) (2.651,3.000) (2.656,3.050) (2.660,3.100) (2.664,3.150) (2.666,3.200) (2.667,3.250) (2.666,3.300) (2.665,3.350) (2.663,3.400) (2.659,3.450) (2.654,3.500) (2.648,3.550) (2.641,3.600) (2.633,3.650) (2.624,3.700) (2.613,3.750) (2.602,3.800) (2.589,3.850) (2.575,3.899) (2.559,3.950) (2.543,4.000) (2.525,4.050) (2.506,4.100) (2.485,4.150) (2.463,4.200) (2.440,4.250) (2.415,4.300) (2.389,4.350) (2.361,4.400) (2.332,4.450) (2.301,4.500) (2.268,4.550) (2.234,4.600) (2.200,4.647) (2.160,4.700) (2.125,4.743) (2.077,4.800) (2.032,4.850) (1.985,4.900) (1.935,4.950) (1.882,5.000) (1.826,5.050) (1.775,5.094) (1.725,5.134) (1.675,5.172) (1.625,5.209) (1.575,5.244) (1.528,5.275) (1.475,5.308) (1.425,5.338) (1.375,5.367) (1.325,5.394) (1.275,5.420) (1.225,5.445) (1.175,5.468) (1.125,5.491) (1.075,5.512) (1.025,5.532) (0.976,5.550) (0.925,5.568) (0.875,5.585) (0.826,5.600) (0.775,5.615) (0.725,5.629) (0.675,5.641) (0.625,5.653) (0.575,5.664) (0.525,5.674) (0.475,5.683) (0.425,5.691) (0.375,5.698) (0.325,5.704) (0.275,5.709) (0.225,5.713) (0.175,5.717) (0.125,5.720) (0.075,5.721) (0.025,5.722) (-0.050,5.722) (-0.100,5.721) (-0.150,5.718) (-0.200,5.715) (-0.250,5.711) (-0.300,5.707) (-0.350,5.701) (-0.400,5.694) (-0.450,5.687) (-0.500,5.678) (-0.550,5.669) (-0.600,5.659) (-0.650,5.647) (-0.700,5.635) (-0.750,5.622) (-0.800,5.608) (-0.850,5.593) (-0.900,5.577) (-0.950,5.559) (-1.000,5.541) (-1.050,5.522) (-1.100,5.501) (-1.150,5.480) (-1.200,5.457) (-1.250,5.433) (-1.300,5.407) (-1.350,5.381) (-1.400,5.353) (-1.447,5.325) (-1.500,5.293) (-1.550,5.260) (-1.600,5.226) (-1.650,5.191) (-1.700,5.153) (-1.750,5.114) (-1.797,5.075) (-1.850,5.029) (-1.900,4.983) (-1.950,4.935) (-2.000,4.884) (-2.050,4.830) (-2.098,4.775) (-2.140,4.725) (-2.175,4.680) (-2.216,4.625) (-2.250,4.577) (-2.285,4.525) (-2.317,4.475) (-2.347,4.425) (-2.375,4.376) (-2.400,4.330) (-2.425,4.281) (-2.450,4.229) (-2.474,4.175) (-2.496,4.125) (-2.515,4.075) (-2.534,4.025) (-2.551,3.975) (-2.567,3.925) (-2.582,3.875) (-2.595,3.825) (-2.608,3.775) (-2.619,3.725) (-2.629,3.675) (-2.637,3.625) (-2.645,3.575) (-2.651,3.525) (-2.657,3.475) (-2.661,3.425) (-2.664,3.375) (-2.666,3.325) (-2.667,3.275) (-2.666,3.225) (-2.665,3.175) (-2.662,3.125) (-2.658,3.075) (-2.654,3.025) (-2.647,2.975) (-2.640,2.925) (-2.632,2.875) (-2.623,2.825) (-2.612,2.775) (-2.600,2.725) (-2.587,2.675) (-2.573,2.625) (-2.558,2.575) (-2.541,2.525) (-2.523,2.475) (-2.504,2.425) (-2.483,2.375) (-2.462,2.325) (-2.439,2.275) (-2.414,2.225) (-2.389,2.175) (-2.362,2.125) (-2.333,2.075) (-2.303,2.025) (-2.275,1.980) (-2.239,1.925) (-2.205,1.875) (-2.175,1.833) (-2.132,1.775) (-2.100,1.733) (-2.053,1.675) (-2.012,1.625) (-1.975,1.583) (-1.925,1.527) (-1.877,1.475) (-1.829,1.425) (-1.780,1.375) (-1.729,1.325) (-1.676,1.275) (-1.625,1.228) (-1.575,1.183) (-1.525,1.139) (-1.481,1.100) (-1.425,1.053) (-1.375,1.011) (-1.331,0.975) (-1.275,0.930) (-1.225,0.891) (-1.175,0.851) (-1.125,0.813) (-1.076,0.775) (-1.025,0.737) (-0.976,0.700) (-0.925,0.662) (-0.875,0.625) (-0.825,0.589) (-0.775,0.552) (-0.725,0.516) (-0.675,0.480) (-0.635,0.450) (-0.575,0.408) (-0.530,0.375) (-0.475,0.337) (-0.425,0.301) (-0.375,0.266) (-0.325,0.230) (-0.284,0.200) (-0.225,0.160) (-0.220,0.155)}
\EQ{(-0.220,-0.155) (-0.275,-0.195) (-0.319,-0.225) (-0.375,-0.266) (-0.424,-0.300) (-0.475,-0.337) (-0.525,-0.372) (-0.575,-0.408) (-0.625,-0.444) (-0.669,-0.475) (-0.725,-0.516) (-0.772,-0.550) (-0.825,-0.589) (-0.874,-0.625) (-0.925,-0.663) (-0.975,-0.700) (-1.025,-0.738) (-1.074,-0.775) (-1.125,-0.815) (-1.171,-0.850) (-1.225,-0.893) (-1.265,-0.925) (-1.325,-0.974) (-1.375,-1.016) (-1.416,-1.050) (-1.473,-1.100) (-1.525,-1.146) (-1.575,-1.191) (-1.612,-1.225) (-1.650,-1.261) (-1.691,-1.300) (-1.741,-1.350) (-1.790,-1.400) (-1.825,-1.438) (-1.860,-1.475) (-1.900,-1.520) (-1.947,-1.575) (-1.989,-1.625) (-2.025,-1.670) (-2.067,-1.725) (-2.100,-1.770) (-2.139,-1.825) (-2.173,-1.875) (-2.205,-1.925) (-2.236,-1.975) (-2.265,-2.025) (-2.293,-2.075) (-2.319,-2.125) (-2.344,-2.175) (-2.368,-2.225) (-2.390,-2.275) (-2.411,-2.325) (-2.431,-2.375) (-2.449,-2.425) (-2.466,-2.475) (-2.481,-2.525) (-2.496,-2.575) (-2.509,-2.625) (-2.521,-2.675) (-2.532,-2.725) (-2.541,-2.775) (-2.549,-2.825) (-2.557,-2.875) (-2.562,-2.925) (-2.567,-2.975) (-2.571,-3.025) (-2.573,-3.075) (-2.574,-3.125) (-2.574,-3.175) (-2.573,-3.225) (-2.571,-3.275) (-2.568,-3.325) (-2.563,-3.375) (-2.557,-3.425) (-2.550,-3.475) (-2.542,-3.525) (-2.533,-3.575) (-2.522,-3.625) (-2.510,-3.675) (-2.497,-3.725) (-2.483,-3.775) (-2.467,-3.825) (-2.450,-3.875) (-2.432,-3.925) (-2.413,-3.975) (-2.392,-4.025) (-2.369,-4.075) (-2.345,-4.125) (-2.320,-4.175) (-2.293,-4.225) (-2.264,-4.275) (-2.234,-4.325) (-2.202,-4.375) (-2.168,-4.425) (-2.133,-4.475) (-2.100,-4.518) (-2.055,-4.575) (-2.013,-4.625) (-1.975,-4.668) (-1.925,-4.722) (-1.875,-4.772) (-1.825,-4.820) (-1.775,-4.866) (-1.735,-4.900) (-1.675,-4.949) (-1.625,-4.988) (-1.575,-5.025) (-1.525,-5.060) (-1.475,-5.093) (-1.425,-5.124) (-1.375,-5.155) (-1.325,-5.183) (-1.275,-5.210) (-1.225,-5.236) (-1.175,-5.261) (-1.125,-5.284) (-1.075,-5.306) (-1.025,-5.327) (-0.975,-5.346) (-0.925,-5.365) (-0.875,-5.382) (-0.825,-5.398) (-0.775,-5.414) (-0.725,-5.428) (-0.675,-5.441) (-0.625,-5.453) (-0.575,-5.464) (-0.525,-5.474) (-0.475,-5.484) (-0.425,-5.492) (-0.375,-5.499) (-0.325,-5.506) (-0.275,-5.511) (-0.225,-5.516) (-0.175,-5.519) (-0.125,-5.522) (-0.075,-5.524) (-0.025,-5.525) (0.050,-5.524) (0.100,-5.523) (0.150,-5.521) (0.200,-5.518) (0.250,-5.513) (0.300,-5.508) (0.350,-5.503) (0.400,-5.496) (0.450,-5.488) (0.500,-5.479) (0.550,-5.469) (0.600,-5.459) (0.650,-5.447) (0.700,-5.434) (0.750,-5.421) (0.800,-5.406) (0.850,-5.390) (0.900,-5.374) (0.950,-5.356) (1.000,-5.337) (1.050,-5.316) (1.100,-5.295) (1.150,-5.272) (1.197,-5.250) (1.247,-5.225) (1.294,-5.200) (1.340,-5.175) (1.400,-5.140) (1.450,-5.109) (1.500,-5.076) (1.550,-5.042) (1.600,-5.006) (1.642,-4.975) (1.700,-4.929) (1.750,-4.887) (1.792,-4.850) (1.846,-4.800) (1.898,-4.750) (1.946,-4.700) (1.991,-4.650) (2.025,-4.611) (2.075,-4.550) (2.114,-4.500) (2.150,-4.451) (2.185,-4.400) (2.218,-4.350) (2.249,-4.300) (2.275,-4.257) (2.300,-4.212) (2.325,-4.165) (2.350,-4.115) (2.375,-4.062) (2.400,-4.005) (2.423,-3.950) (2.442,-3.900) (2.459,-3.850) (2.475,-3.801) (2.490,-3.750) (2.504,-3.700) (2.516,-3.650) (2.528,-3.600) (2.538,-3.550) (2.546,-3.500) (2.554,-3.450) (2.560,-3.400) (2.566,-3.350) (2.570,-3.300) (2.572,-3.250) (2.574,-3.200) (2.575,-3.150) (2.574,-3.100) (2.572,-3.050) (2.569,-3.000) (2.565,-2.950) (2.560,-2.900) (2.553,-2.850) (2.545,-2.800) (2.537,-2.750) (2.526,-2.700) (2.515,-2.650) (2.502,-2.600) (2.489,-2.550) (2.474,-2.500) (2.457,-2.450) (2.440,-2.400) (2.421,-2.350) (2.401,-2.300) (2.379,-2.250) (2.356,-2.200) (2.332,-2.150) (2.306,-2.100) (2.279,-2.050) (2.251,-2.000) (2.221,-1.950) (2.189,-1.900) (2.156,-1.850) (2.125,-1.805) (2.086,-1.750) (2.050,-1.702) (2.009,-1.650) (1.975,-1.608) (1.926,-1.550) (1.882,-1.500) (1.837,-1.450) (1.800,-1.411) (1.765,-1.375) (1.725,-1.334) (1.675,-1.285) (1.639,-1.250) (1.600,-1.214) (1.558,-1.175) (1.502,-1.125) (1.450,-1.080) (1.400,-1.037) (1.356,-1.000) (1.300,-0.954) (1.250,-0.913) (1.202,-0.875) (1.150,-0.834) (1.107,-0.800) (1.050,-0.757) (1.008,-0.725) (0.950,-0.682) (0.908,-0.650) (0.850,-0.607) (0.807,-0.575) (0.750,-0.534) (0.703,-0.500) (0.650,-0.462) (0.600,-0.426) (0.550,-0.390) (0.500,-0.355) (0.450,-0.319) (0.400,-0.284) (0.354,-0.250) (0.300,-0.213) (0.250,-0.177) (0.206,-0.144)}
\EQ{(0.206,0.144) (0.248,0.175) (0.300,0.213) (0.350,0.248) (0.400,0.284) (0.450,0.319) (0.495,0.350) (0.550,0.390) (0.599,0.425) (0.650,0.462) (0.700,0.498) (0.750,0.534) (0.800,0.571) (0.841,0.600) (0.900,0.644) (0.942,0.675) (1.000,0.719) (1.041,0.750) (1.100,0.795) (1.150,0.834) (1.200,0.873) (1.250,0.913) (1.296,0.950) (1.350,0.995) (1.400,1.037) (1.445,1.075) (1.500,1.123) (1.550,1.168) (1.600,1.214) (1.639,1.250) (1.675,1.285) (1.716,1.325) (1.765,1.375) (1.800,1.411) (1.837,1.450) (1.875,1.492) (1.925,1.549) (1.968,1.600) (2.000,1.639) (2.048,1.700) (2.086,1.750) (2.122,1.800) (2.156,1.850) (2.189,1.900) (2.221,1.950) (2.250,1.999) (2.275,2.043) (2.300,2.088) (2.325,2.136) (2.350,2.187) (2.375,2.241) (2.400,2.299) (2.421,2.350) (2.440,2.400) (2.457,2.450) (2.474,2.500) (2.489,2.550) (2.502,2.600) (2.515,2.650) (2.526,2.700) (2.537,2.750) (2.545,2.800) (2.553,2.850) (2.560,2.900) (2.565,2.950) (2.569,3.000) (2.572,3.050) (2.574,3.100) (2.575,3.150) (2.574,3.200) (2.572,3.250) (2.570,3.300) (2.566,3.350) (2.560,3.400) (2.554,3.450) (2.546,3.500) (2.538,3.550) (2.528,3.600) (2.516,3.650) (2.504,3.700) (2.490,3.750) (2.475,3.800) (2.459,3.850) (2.442,3.900) (2.423,3.950) (2.402,4.000) (2.381,4.050) (2.357,4.100) (2.333,4.150) (2.307,4.200) (2.279,4.250) (2.250,4.299) (2.218,4.350) (2.185,4.400) (2.151,4.450) (2.114,4.500) (2.075,4.550) (2.034,4.600) (2.000,4.640) (1.950,4.695) (1.900,4.747) (1.850,4.797) (1.800,4.843) (1.750,4.887) (1.705,4.925) (1.650,4.969) (1.609,5.000) (1.550,5.042) (1.502,5.075) (1.450,5.109) (1.400,5.140) (1.350,5.169) (1.300,5.197) (1.250,5.223) (1.200,5.249) (1.150,5.272) (1.100,5.295) (1.050,5.316) (1.000,5.337) (0.950,5.356) (0.900,5.374) (0.850,5.390) (0.800,5.406) (0.750,5.421) (0.700,5.434) (0.650,5.447) (0.600,5.459) (0.550,5.469) (0.500,5.479) (0.450,5.488) (0.400,5.496) (0.350,5.503) (0.300,5.508) (0.250,5.513) (0.200,5.518) (0.150,5.521) (0.100,5.523) (0.050,5.524) (0.000,5.525) (-0.050,5.524) (-0.100,5.523) (-0.150,5.521) (-0.200,5.518) (-0.250,5.513) (-0.300,5.508) (-0.350,5.503) (-0.400,5.496) (-0.450,5.488) (-0.500,5.479) (-0.550,5.469) (-0.600,5.459) (-0.650,5.447) (-0.700,5.434) (-0.750,5.421) (-0.800,5.406) (-0.850,5.390) (-0.900,5.374) (-0.950,5.356) (-1.000,5.337) (-1.050,5.316) (-1.100,5.295) (-1.150,5.272) (-1.197,5.250) (-1.247,5.225) (-1.294,5.200) (-1.340,5.175) (-1.400,5.140) (-1.450,5.109) (-1.500,5.076) (-1.550,5.042) (-1.600,5.006) (-1.642,4.975) (-1.700,4.929) (-1.750,4.887) (-1.792,4.850) (-1.846,4.800) (-1.898,4.750) (-1.946,4.700) (-1.991,4.650) (-2.025,4.611) (-2.075,4.550) (-2.114,4.500) (-2.150,4.451) (-2.185,4.400) (-2.218,4.350) (-2.249,4.300) (-2.275,4.257) (-2.300,4.212) (-2.325,4.165) (-2.350,4.115) (-2.375,4.062) (-2.400,4.005) (-2.423,3.950) (-2.442,3.900) (-2.459,3.850) (-2.475,3.801) (-2.490,3.750) (-2.504,3.700) (-2.516,3.650) (-2.528,3.600) (-2.538,3.550) (-2.546,3.500) (-2.554,3.450) (-2.560,3.400) (-2.566,3.350) (-2.570,3.300) (-2.572,3.250) (-2.574,3.200) (-2.575,3.150) (-2.574,3.100) (-2.572,3.050) (-2.569,3.000) (-2.565,2.950) (-2.560,2.900) (-2.553,2.850) (-2.545,2.800) (-2.537,2.750) (-2.526,2.700) (-2.515,2.650) (-2.502,2.600) (-2.489,2.550) (-2.474,2.500) (-2.457,2.450) (-2.440,2.400) (-2.421,2.350) (-2.401,2.300) (-2.379,2.250) (-2.356,2.200) (-2.332,2.150) (-2.306,2.100) (-2.279,2.050) (-2.251,2.000) (-2.221,1.950) (-2.189,1.900) (-2.156,1.850) (-2.125,1.805) (-2.086,1.750) (-2.050,1.702) (-2.009,1.650) (-1.975,1.608) (-1.926,1.550) (-1.882,1.500) (-1.837,1.450) (-1.800,1.411) (-1.765,1.375) (-1.725,1.334) (-1.675,1.285) (-1.639,1.250) (-1.600,1.214) (-1.558,1.175) (-1.502,1.125) (-1.450,1.080) (-1.400,1.037) (-1.356,1.000) (-1.300,0.954) (-1.250,0.913) (-1.202,0.875) (-1.150,0.834) (-1.107,0.800) (-1.050,0.757) (-1.008,0.725) (-0.950,0.682) (-0.908,0.650) (-0.850,0.607) (-0.807,0.575) (-0.750,0.534) (-0.703,0.500) (-0.650,0.462) (-0.600,0.426) (-0.550,0.390) (-0.500,0.355) (-0.450,0.319) (-0.400,0.284) (-0.354,0.250) (-0.300,0.213) (-0.250,0.177) (-0.220,0.155)}
}
  \icono
  \marcoMMA{10}{5}{$x$ (u.\,arb.)}{$z$ (u.\,arb.)}
\end{tikzpicture}\\[1pt]{\small\scshape (a)}
\end{minipage}\hfill
\begin{minipage}[t]{0.325\textwidth}\centering
{\footnotesize Equipotenciales y líneas de campo}\\[3pt]
\begin{tikzpicture}[x=0.198cm, y=0.198cm]
  \panelMMA{10}{mmafondo}{1.2}{%
    % Cuadrupolo puntual - equipotenciales
% Generado por scripts/generar_datos.py (no editar a mano).
\EQ{(-0.220,-0.155) (-0.275,-0.195) (-0.319,-0.225) (-0.375,-0.266) (-0.425,-0.300) (-0.475,-0.336) (-0.525,-0.371) (-0.568,-0.400) (-0.625,-0.441) (-0.675,-0.475) (-0.725,-0.510) (-0.775,-0.544) (-0.821,-0.575) (-0.875,-0.612) (-0.925,-0.646) (-0.969,-0.675) (-1.025,-0.713) (-1.075,-0.745) (-1.121,-0.775) (-1.175,-0.810) (-1.225,-0.841) (-1.275,-0.872) (-1.321,-0.900) (-1.375,-0.933) (-1.425,-0.963) (-1.475,-0.991) (-1.525,-1.020) (-1.575,-1.047) (-1.625,-1.074) (-1.674,-1.100) (-1.723,-1.125) (-1.774,-1.150) (-1.825,-1.175) (-1.875,-1.198) (-1.925,-1.220) (-1.975,-1.241) (-2.025,-1.262) (-2.075,-1.281) (-2.125,-1.300) (-2.175,-1.317) (-2.225,-1.334) (-2.275,-1.349) (-2.325,-1.364) (-2.375,-1.377) (-2.425,-1.389) (-2.475,-1.400) (-2.525,-1.410) (-2.575,-1.419) (-2.625,-1.426) (-2.675,-1.433) (-2.725,-1.438) (-2.775,-1.441) (-2.825,-1.444) (-2.875,-1.445) (-2.925,-1.445) (-2.975,-1.443) (-3.025,-1.440) (-3.075,-1.435) (-3.125,-1.429) (-3.175,-1.421) (-3.225,-1.412) (-3.275,-1.401) (-3.325,-1.388) (-3.375,-1.373) (-3.425,-1.357) (-3.475,-1.338) (-3.525,-1.318) (-3.575,-1.295) (-3.625,-1.270) (-3.675,-1.242) (-3.725,-1.212) (-3.775,-1.179) (-3.825,-1.142) (-3.875,-1.102) (-3.925,-1.058) (-3.975,-1.010) (-4.025,-0.956) (-4.072,-0.900) (-4.110,-0.850) (-4.145,-0.800) (-4.175,-0.753) (-4.206,-0.700) (-4.232,-0.650) (-4.256,-0.600) (-4.277,-0.550) (-4.297,-0.500) (-4.314,-0.450) (-4.329,-0.400) (-4.343,-0.350) (-4.354,-0.300) (-4.364,-0.250) (-4.371,-0.200) (-4.377,-0.150) (-4.382,-0.100) (-4.384,-0.050) (-4.385,0.000) (-4.384,0.050) (-4.382,0.100) (-4.377,0.150) (-4.371,0.200) (-4.364,0.250) (-4.354,0.300) (-4.343,0.350) (-4.329,0.400) (-4.314,0.450) (-4.297,0.500) (-4.277,0.550) (-4.256,0.600) (-4.232,0.650) (-4.206,0.700) (-4.177,0.750) (-4.150,0.792) (-4.110,0.850) (-4.075,0.896) (-4.030,0.950) (-3.985,1.000) (-3.934,1.050) (-3.878,1.100) (-3.825,1.142) (-3.780,1.175) (-3.725,1.212) (-3.675,1.242) (-3.625,1.270) (-3.575,1.295) (-3.525,1.318) (-3.475,1.338) (-3.425,1.357) (-3.375,1.373) (-3.325,1.388) (-3.275,1.401) (-3.225,1.412) (-3.175,1.421) (-3.125,1.429) (-3.075,1.435) (-3.025,1.440) (-2.975,1.443) (-2.925,1.445) (-2.875,1.445) (-2.825,1.444) (-2.775,1.441) (-2.725,1.438) (-2.675,1.433) (-2.625,1.426) (-2.575,1.419) (-2.525,1.410) (-2.475,1.400) (-2.425,1.389) (-2.375,1.377) (-2.325,1.364) (-2.275,1.349) (-2.225,1.334) (-2.175,1.317) (-2.126,1.300) (-2.075,1.281) (-2.025,1.262) (-1.975,1.241) (-1.925,1.220) (-1.875,1.198) (-1.826,1.175) (-1.775,1.151) (-1.725,1.126) (-1.675,1.100) (-1.627,1.075) (-1.580,1.050) (-1.535,1.025) (-1.491,1.000) (-1.447,0.975) (-1.400,0.948) (-1.350,0.918) (-1.300,0.888) (-1.250,0.857) (-1.200,0.825) (-1.150,0.794) (-1.100,0.762) (-1.050,0.729) (-1.007,0.700) (-0.950,0.663) (-0.900,0.629) (-0.858,0.600) (-0.800,0.561) (-0.750,0.527) (-0.700,0.493) (-0.650,0.458) (-0.604,0.425) (-0.550,0.388) (-0.500,0.353) (-0.450,0.318) (-0.400,0.283) (-0.354,0.250) (-0.300,0.213) (-0.250,0.177) (-0.220,0.155)}
\EQ{(0.206,0.144) (0.248,0.175) (0.300,0.213) (0.350,0.248) (0.400,0.283) (0.450,0.318) (0.496,0.350) (0.550,0.388) (0.600,0.423) (0.650,0.458) (0.700,0.493) (0.748,0.525) (0.800,0.561) (0.850,0.595) (0.895,0.625) (0.950,0.663) (1.000,0.696) (1.045,0.725) (1.100,0.762) (1.150,0.794) (1.200,0.825) (1.250,0.857) (1.300,0.888) (1.350,0.918) (1.400,0.948) (1.447,0.975) (1.491,1.000) (1.535,1.025) (1.580,1.050) (1.625,1.074) (1.674,1.100) (1.723,1.125) (1.774,1.150) (1.825,1.175) (1.875,1.198) (1.925,1.220) (1.975,1.241) (2.025,1.262) (2.075,1.281) (2.125,1.300) (2.175,1.317) (2.225,1.334) (2.275,1.349) (2.325,1.364) (2.375,1.377) (2.425,1.389) (2.475,1.400) (2.525,1.410) (2.575,1.419) (2.625,1.426) (2.675,1.433) (2.725,1.438) (2.775,1.441) (2.825,1.444) (2.875,1.445) (2.925,1.445) (2.975,1.443) (3.025,1.440) (3.075,1.435) (3.125,1.429) (3.175,1.421) (3.225,1.412) (3.275,1.401) (3.325,1.388) (3.375,1.373) (3.425,1.357) (3.475,1.338) (3.525,1.318) (3.575,1.295) (3.625,1.270) (3.675,1.242) (3.725,1.212) (3.775,1.179) (3.825,1.142) (3.875,1.102) (3.925,1.058) (3.975,1.010) (4.025,0.956) (4.072,0.900) (4.110,0.850) (4.145,0.800) (4.175,0.753) (4.206,0.700) (4.232,0.650) (4.256,0.600) (4.277,0.550) (4.297,0.500) (4.314,0.450) (4.329,0.400) (4.343,0.350) (4.354,0.300) (4.364,0.250) (4.371,0.200) (4.377,0.150) (4.382,0.100) (4.384,0.050) (4.385,0.000) (4.384,-0.050) (4.382,-0.100) (4.377,-0.150) (4.371,-0.200) (4.364,-0.250) (4.354,-0.300) (4.343,-0.350) (4.329,-0.400) (4.314,-0.450) (4.297,-0.500) (4.277,-0.550) (4.256,-0.600) (4.232,-0.650) (4.206,-0.700) (4.177,-0.750) (4.150,-0.792) (4.110,-0.850) (4.075,-0.896) (4.030,-0.950) (3.985,-1.000) (3.934,-1.050) (3.878,-1.100) (3.825,-1.142) (3.780,-1.175) (3.725,-1.212) (3.675,-1.242) (3.625,-1.270) (3.575,-1.295) (3.525,-1.318) (3.475,-1.338) (3.425,-1.357) (3.375,-1.373) (3.325,-1.388) (3.275,-1.401) (3.225,-1.412) (3.175,-1.421) (3.125,-1.429) (3.075,-1.435) (3.025,-1.440) (2.975,-1.443) (2.925,-1.445) (2.875,-1.445) (2.825,-1.444) (2.775,-1.441) (2.725,-1.438) (2.675,-1.433) (2.625,-1.426) (2.575,-1.419) (2.525,-1.410) (2.475,-1.400) (2.425,-1.389) (2.375,-1.377) (2.325,-1.364) (2.275,-1.349) (2.225,-1.334) (2.175,-1.317) (2.126,-1.300) (2.075,-1.281) (2.025,-1.262) (1.975,-1.241) (1.925,-1.220) (1.875,-1.198) (1.826,-1.175) (1.775,-1.151) (1.725,-1.126) (1.675,-1.100) (1.627,-1.075) (1.580,-1.050) (1.535,-1.025) (1.491,-1.000) (1.447,-0.975) (1.400,-0.948) (1.350,-0.918) (1.300,-0.888) (1.250,-0.857) (1.200,-0.825) (1.150,-0.794) (1.100,-0.762) (1.050,-0.729) (1.007,-0.700) (0.950,-0.663) (0.900,-0.629) (0.858,-0.600) (0.800,-0.561) (0.750,-0.527) (0.700,-0.493) (0.650,-0.458) (0.604,-0.425) (0.550,-0.388) (0.500,-0.353) (0.450,-0.318) (0.400,-0.283) (0.354,-0.250) (0.300,-0.213) (0.250,-0.177) (0.206,-0.144)}
\EQ{(-0.220,-0.155) (-0.275,-0.195) (-0.319,-0.225) (-0.375,-0.266) (-0.425,-0.300) (-0.475,-0.336) (-0.525,-0.371) (-0.568,-0.400) (-0.625,-0.441) (-0.675,-0.475) (-0.725,-0.510) (-0.775,-0.545) (-0.820,-0.575) (-0.875,-0.613) (-0.925,-0.647) (-0.968,-0.675) (-1.025,-0.714) (-1.075,-0.747) (-1.119,-0.775) (-1.175,-0.812) (-1.225,-0.844) (-1.275,-0.875) (-1.325,-0.906) (-1.375,-0.937) (-1.425,-0.967) (-1.475,-0.996) (-1.525,-1.025) (-1.569,-1.050) (-1.614,-1.075) (-1.660,-1.100) (-1.707,-1.125) (-1.755,-1.150) (-1.800,-1.173) (-1.850,-1.197) (-1.900,-1.221) (-1.950,-1.244) (-2.000,-1.266) (-2.050,-1.288) (-2.100,-1.308) (-2.150,-1.327) (-2.200,-1.346) (-2.250,-1.363) (-2.300,-1.380) (-2.350,-1.396) (-2.400,-1.410) (-2.450,-1.424) (-2.500,-1.436) (-2.550,-1.447) (-2.600,-1.458) (-2.650,-1.467) (-2.700,-1.475) (-2.750,-1.481) (-2.800,-1.487) (-2.850,-1.491) (-2.900,-1.494) (-2.950,-1.496) (-3.000,-1.497) (-3.050,-1.496) (-3.100,-1.493) (-3.150,-1.490) (-3.200,-1.485) (-3.250,-1.478) (-3.300,-1.470) (-3.350,-1.460) (-3.400,-1.449) (-3.450,-1.436) (-3.500,-1.421) (-3.550,-1.405) (-3.600,-1.386) (-3.650,-1.365) (-3.700,-1.343) (-3.750,-1.318) (-3.800,-1.290) (-3.850,-1.260) (-3.900,-1.228) (-3.950,-1.192) (-4.000,-1.153) (-4.050,-1.110) (-4.088,-1.075) (-4.125,-1.038) (-4.160,-1.000) (-4.200,-0.954) (-4.243,-0.900) (-4.275,-0.855) (-4.312,-0.800) (-4.342,-0.750) (-4.369,-0.700) (-4.395,-0.650) (-4.417,-0.600) (-4.438,-0.550) (-4.457,-0.500) (-4.473,-0.450) (-4.488,-0.400) (-4.501,-0.350) (-4.512,-0.300) (-4.521,-0.250) (-4.529,-0.200) (-4.534,-0.150) (-4.539,-0.100) (-4.541,-0.050) (-4.542,0.000) (-4.541,0.050) (-4.539,0.100) (-4.534,0.150) (-4.529,0.200) (-4.521,0.250) (-4.512,0.300) (-4.501,0.350) (-4.488,0.400) (-4.473,0.450) (-4.457,0.500) (-4.438,0.550) (-4.417,0.600) (-4.395,0.650) (-4.369,0.700) (-4.342,0.750) (-4.312,0.800) (-4.279,0.850) (-4.243,0.900) (-4.203,0.950) (-4.160,1.000) (-4.125,1.038) (-4.088,1.075) (-4.050,1.110) (-4.003,1.150) (-3.950,1.192) (-3.904,1.225) (-3.850,1.260) (-3.800,1.290) (-3.750,1.318) (-3.700,1.343) (-3.650,1.365) (-3.600,1.386) (-3.550,1.405) (-3.500,1.421) (-3.450,1.436) (-3.400,1.449) (-3.350,1.460) (-3.300,1.470) (-3.250,1.478) (-3.200,1.485) (-3.150,1.490) (-3.100,1.493) (-3.050,1.496) (-3.000,1.497) (-2.950,1.496) (-2.900,1.494) (-2.850,1.491) (-2.800,1.487) (-2.750,1.481) (-2.700,1.475) (-2.650,1.467) (-2.600,1.458) (-2.550,1.447) (-2.500,1.436) (-2.450,1.424) (-2.400,1.410) (-2.350,1.396) (-2.300,1.380) (-2.250,1.363) (-2.200,1.346) (-2.150,1.327) (-2.100,1.308) (-2.050,1.288) (-2.000,1.266) (-1.950,1.244) (-1.900,1.221) (-1.850,1.197) (-1.804,1.175) (-1.755,1.150) (-1.707,1.125) (-1.660,1.100) (-1.614,1.075) (-1.569,1.050) (-1.525,1.025) (-1.482,1.000) (-1.425,0.967) (-1.375,0.937) (-1.325,0.906) (-1.275,0.875) (-1.225,0.844) (-1.175,0.812) (-1.125,0.779) (-1.081,0.750) (-1.025,0.714) (-0.975,0.680) (-0.930,0.650) (-0.875,0.613) (-0.825,0.579) (-0.784,0.550) (-0.725,0.510) (-0.675,0.475) (-0.625,0.441) (-0.575,0.406) (-0.532,0.375) (-0.475,0.336) (-0.425,0.300) (-0.375,0.266) (-0.325,0.230) (-0.284,0.200) (-0.225,0.160) (-0.220,0.155)}
\EQ{(0.206,0.144) (0.248,0.175) (0.300,0.213) (0.350,0.248) (0.400,0.283) (0.450,0.318) (0.496,0.350) (0.550,0.388) (0.600,0.423) (0.650,0.458) (0.700,0.493) (0.747,0.525) (0.800,0.562) (0.850,0.596) (0.894,0.625) (0.950,0.664) (1.000,0.697) (1.043,0.725) (1.100,0.763) (1.150,0.796) (1.196,0.825) (1.250,0.859) (1.300,0.891) (1.350,0.921) (1.397,0.950) (1.450,0.982) (1.500,1.011) (1.550,1.040) (1.600,1.068) (1.650,1.095) (1.700,1.122) (1.750,1.148) (1.800,1.173) (1.850,1.197) (1.900,1.221) (1.950,1.244) (2.000,1.266) (2.050,1.288) (2.100,1.308) (2.150,1.327) (2.200,1.346) (2.250,1.363) (2.300,1.380) (2.350,1.396) (2.400,1.410) (2.450,1.424) (2.500,1.436) (2.550,1.447) (2.600,1.458) (2.650,1.467) (2.700,1.475) (2.750,1.481) (2.800,1.487) (2.850,1.491) (2.900,1.494) (2.950,1.496) (3.000,1.497) (3.050,1.496) (3.100,1.493) (3.150,1.490) (3.200,1.485) (3.250,1.478) (3.300,1.470) (3.350,1.460) (3.400,1.449) (3.450,1.436) (3.500,1.421) (3.550,1.405) (3.600,1.386) (3.650,1.365) (3.700,1.343) (3.750,1.318) (3.800,1.290) (3.850,1.260) (3.900,1.228) (3.950,1.192) (4.000,1.153) (4.050,1.110) (4.088,1.075) (4.125,1.038) (4.160,1.000) (4.200,0.954) (4.243,0.900) (4.275,0.855) (4.312,0.800) (4.342,0.750) (4.369,0.700) (4.395,0.650) (4.417,0.600) (4.438,0.550) (4.457,0.500) (4.473,0.450) (4.488,0.400) (4.501,0.350) (4.512,0.300) (4.521,0.250) (4.529,0.200) (4.534,0.150) (4.539,0.100) (4.541,0.050) (4.542,0.000) (4.541,-0.050) (4.539,-0.100) (4.534,-0.150) (4.529,-0.200) (4.521,-0.250) (4.512,-0.300) (4.501,-0.350) (4.488,-0.400) (4.473,-0.450) (4.457,-0.500) (4.438,-0.550) (4.417,-0.600) (4.395,-0.650) (4.369,-0.700) (4.342,-0.750) (4.312,-0.800) (4.279,-0.850) (4.243,-0.900) (4.203,-0.950) (4.160,-1.000) (4.125,-1.038) (4.088,-1.075) (4.050,-1.110) (4.003,-1.150) (3.950,-1.192) (3.904,-1.225) (3.850,-1.260) (3.800,-1.290) (3.750,-1.318) (3.700,-1.343) (3.650,-1.365) (3.600,-1.386) (3.550,-1.405) (3.500,-1.421) (3.450,-1.436) (3.400,-1.449) (3.350,-1.460) (3.300,-1.470) (3.250,-1.478) (3.200,-1.485) (3.150,-1.490) (3.100,-1.493) (3.050,-1.496) (3.000,-1.497) (2.950,-1.496) (2.900,-1.494) (2.850,-1.491) (2.800,-1.487) (2.750,-1.481) (2.700,-1.475) (2.650,-1.467) (2.600,-1.458) (2.550,-1.447) (2.500,-1.436) (2.450,-1.424) (2.400,-1.410) (2.350,-1.396) (2.300,-1.380) (2.250,-1.363) (2.200,-1.346) (2.150,-1.327) (2.100,-1.308) (2.050,-1.288) (2.000,-1.266) (1.950,-1.244) (1.900,-1.221) (1.850,-1.197) (1.804,-1.175) (1.755,-1.150) (1.707,-1.125) (1.660,-1.100) (1.614,-1.075) (1.569,-1.050) (1.525,-1.025) (1.482,-1.000) (1.425,-0.967) (1.375,-0.937) (1.325,-0.906) (1.275,-0.875) (1.225,-0.844) (1.175,-0.812) (1.125,-0.779) (1.081,-0.750) (1.025,-0.714) (0.975,-0.680) (0.930,-0.650) (0.875,-0.613) (0.825,-0.579) (0.784,-0.550) (0.725,-0.510) (0.675,-0.475) (0.625,-0.441) (0.575,-0.406) (0.532,-0.375) (0.475,-0.336) (0.425,-0.300) (0.375,-0.266) (0.325,-0.230) (0.284,-0.200) (0.225,-0.160) (0.206,-0.144)}
\EQ{(-0.220,-0.155) (-0.275,-0.195) (-0.319,-0.225) (-0.375,-0.266) (-0.425,-0.300) (-0.475,-0.336) (-0.525,-0.371) (-0.568,-0.400) (-0.625,-0.441) (-0.674,-0.475) (-0.725,-0.510) (-0.775,-0.545) (-0.819,-0.575) (-0.875,-0.614) (-0.925,-0.647) (-0.975,-0.681) (-1.025,-0.715) (-1.075,-0.748) (-1.125,-0.781) (-1.175,-0.814) (-1.225,-0.846) (-1.271,-0.875) (-1.325,-0.910) (-1.375,-0.941) (-1.425,-0.971) (-1.473,-1.000) (-1.525,-1.031) (-1.575,-1.060) (-1.625,-1.089) (-1.675,-1.116) (-1.725,-1.144) (-1.775,-1.170) (-1.825,-1.196) (-1.875,-1.221) (-1.925,-1.246) (-1.975,-1.270) (-2.025,-1.293) (-2.075,-1.315) (-2.125,-1.336) (-2.175,-1.357) (-2.222,-1.375) (-2.275,-1.395) (-2.325,-1.413) (-2.375,-1.430) (-2.425,-1.446) (-2.475,-1.461) (-2.525,-1.474) (-2.575,-1.487) (-2.625,-1.499) (-2.675,-1.510) (-2.725,-1.520) (-2.775,-1.529) (-2.825,-1.536) (-2.875,-1.542) (-2.925,-1.548) (-2.975,-1.552) (-3.025,-1.555) (-3.075,-1.556) (-3.125,-1.556) (-3.175,-1.556) (-3.225,-1.553) (-3.275,-1.550) (-3.325,-1.545) (-3.375,-1.538) (-3.425,-1.530) (-3.475,-1.521) (-3.525,-1.510) (-3.575,-1.497) (-3.625,-1.483) (-3.675,-1.467) (-3.725,-1.449) (-3.775,-1.429) (-3.825,-1.407) (-3.875,-1.383) (-3.925,-1.357) (-3.975,-1.329) (-4.021,-1.300) (-4.075,-1.264) (-4.125,-1.227) (-4.175,-1.187) (-4.217,-1.150) (-4.270,-1.100) (-4.318,-1.050) (-4.362,-1.000) (-4.400,-0.952) (-4.439,-0.900) (-4.473,-0.850) (-4.500,-0.807) (-4.533,-0.750) (-4.559,-0.700) (-4.583,-0.650) (-4.605,-0.600) (-4.624,-0.550) (-4.642,-0.500) (-4.658,-0.450) (-4.672,-0.400) (-4.684,-0.350) (-4.695,-0.300) (-4.704,-0.250) (-4.711,-0.200) (-4.717,-0.150) (-4.720,-0.100) (-4.723,-0.050) (-4.724,0.000) (-4.723,0.050) (-4.720,0.100) (-4.717,0.150) (-4.711,0.200) (-4.704,0.250) (-4.695,0.300) (-4.684,0.350) (-4.672,0.400) (-4.658,0.450) (-4.642,0.500) (-4.625,0.548) (-4.605,0.600) (-4.583,0.650) (-4.559,0.700) (-4.533,0.750) (-4.504,0.800) (-4.475,0.847) (-4.439,0.900) (-4.402,0.950) (-4.362,1.000) (-4.325,1.042) (-4.275,1.095) (-4.225,1.143) (-4.175,1.187) (-4.127,1.225) (-4.075,1.264) (-4.025,1.298) (-3.981,1.325) (-3.925,1.357) (-3.875,1.383) (-3.825,1.407) (-3.775,1.429) (-3.725,1.449) (-3.675,1.467) (-3.625,1.483) (-3.575,1.497) (-3.525,1.510) (-3.475,1.521) (-3.425,1.530) (-3.375,1.538) (-3.325,1.545) (-3.275,1.550) (-3.225,1.553) (-3.175,1.556) (-3.125,1.556) (-3.075,1.556) (-3.025,1.555) (-2.975,1.552) (-2.925,1.548) (-2.875,1.542) (-2.825,1.536) (-2.775,1.529) (-2.725,1.520) (-2.675,1.510) (-2.625,1.499) (-2.575,1.487) (-2.525,1.474) (-2.475,1.461) (-2.425,1.446) (-2.375,1.430) (-2.325,1.413) (-2.275,1.395) (-2.225,1.376) (-2.175,1.357) (-2.125,1.336) (-2.075,1.315) (-2.025,1.293) (-1.975,1.270) (-1.925,1.246) (-1.875,1.221) (-1.825,1.196) (-1.775,1.170) (-1.725,1.144) (-1.675,1.116) (-1.625,1.089) (-1.575,1.060) (-1.525,1.031) (-1.475,1.001) (-1.432,0.975) (-1.375,0.941) (-1.325,0.910) (-1.275,0.878) (-1.232,0.850) (-1.175,0.814) (-1.125,0.781) (-1.078,0.750) (-1.025,0.715) (-0.975,0.681) (-0.929,0.650) (-0.875,0.614) (-0.825,0.579) (-0.783,0.550) (-0.725,0.510) (-0.675,0.475) (-0.625,0.441) (-0.575,0.406) (-0.532,0.375) (-0.475,0.336) (-0.425,0.300) (-0.375,0.266) (-0.325,0.230) (-0.284,0.200) (-0.225,0.160) (-0.220,0.155)}
\EQ{(0.206,0.144) (0.248,0.175) (0.300,0.213) (0.350,0.248) (0.400,0.283) (0.450,0.318) (0.496,0.350) (0.550,0.389) (0.600,0.423) (0.650,0.458) (0.700,0.493) (0.747,0.525) (0.800,0.562) (0.850,0.596) (0.893,0.625) (0.950,0.665) (1.000,0.698) (1.050,0.732) (1.100,0.765) (1.150,0.797) (1.193,0.825) (1.250,0.862) (1.300,0.894) (1.350,0.925) (1.400,0.956) (1.450,0.986) (1.500,1.016) (1.550,1.046) (1.600,1.074) (1.646,1.100) (1.691,1.125) (1.737,1.150) (1.785,1.175) (1.833,1.200) (1.883,1.225) (1.934,1.250) (1.987,1.275) (2.042,1.300) (2.099,1.325) (2.150,1.346) (2.200,1.367) (2.250,1.386) (2.300,1.404) (2.350,1.421) (2.400,1.438) (2.450,1.453) (2.500,1.468) (2.550,1.481) (2.600,1.493) (2.650,1.505) (2.700,1.515) (2.750,1.524) (2.800,1.532) (2.850,1.539) (2.900,1.545) (2.950,1.550) (3.000,1.553) (3.050,1.556) (3.100,1.556) (3.150,1.556) (3.200,1.555) (3.250,1.552) (3.300,1.547) (3.350,1.542) (3.400,1.534) (3.450,1.526) (3.500,1.515) (3.550,1.504) (3.600,1.490) (3.650,1.475) (3.700,1.458) (3.750,1.439) (3.800,1.418) (3.850,1.396) (3.900,1.371) (3.950,1.343) (4.000,1.313) (4.050,1.281) (4.094,1.250) (4.150,1.207) (4.189,1.175) (4.244,1.125) (4.294,1.075) (4.340,1.025) (4.375,0.984) (4.421,0.925) (4.450,0.884) (4.489,0.825) (4.519,0.775) (4.546,0.725) (4.571,0.675) (4.594,0.625) (4.615,0.575) (4.633,0.525) (4.650,0.475) (4.665,0.425) (4.678,0.375) (4.690,0.325) (4.699,0.275) (4.708,0.225) (4.714,0.175) (4.719,0.125) (4.722,0.075) (4.723,0.025) (4.723,-0.050) (4.720,-0.100) (4.717,-0.150) (4.711,-0.200) (4.704,-0.250) (4.695,-0.300) (4.684,-0.350) (4.672,-0.400) (4.658,-0.450) (4.642,-0.500) (4.625,-0.548) (4.605,-0.600) (4.583,-0.650) (4.559,-0.700) (4.533,-0.750) (4.504,-0.800) (4.475,-0.847) (4.439,-0.900) (4.402,-0.950) (4.362,-1.000) (4.325,-1.042) (4.275,-1.095) (4.225,-1.143) (4.175,-1.187) (4.127,-1.225) (4.075,-1.264) (4.025,-1.298) (3.981,-1.325) (3.925,-1.357) (3.875,-1.383) (3.825,-1.407) (3.775,-1.429) (3.725,-1.449) (3.675,-1.467) (3.625,-1.483) (3.575,-1.497) (3.525,-1.510) (3.475,-1.521) (3.425,-1.530) (3.375,-1.538) (3.325,-1.545) (3.275,-1.550) (3.225,-1.553) (3.175,-1.556) (3.125,-1.556) (3.075,-1.556) (3.025,-1.555) (2.975,-1.552) (2.925,-1.548) (2.875,-1.542) (2.825,-1.536) (2.775,-1.529) (2.725,-1.520) (2.675,-1.510) (2.625,-1.499) (2.575,-1.487) (2.525,-1.474) (2.475,-1.461) (2.425,-1.446) (2.375,-1.430) (2.325,-1.413) (2.275,-1.395) (2.225,-1.376) (2.175,-1.357) (2.125,-1.336) (2.075,-1.315) (2.025,-1.293) (1.975,-1.270) (1.925,-1.246) (1.875,-1.221) (1.825,-1.196) (1.775,-1.170) (1.725,-1.144) (1.675,-1.116) (1.625,-1.089) (1.575,-1.060) (1.525,-1.031) (1.475,-1.001) (1.432,-0.975) (1.375,-0.941) (1.325,-0.910) (1.275,-0.878) (1.232,-0.850) (1.175,-0.814) (1.125,-0.781) (1.078,-0.750) (1.025,-0.715) (0.975,-0.681) (0.929,-0.650) (0.875,-0.614) (0.825,-0.579) (0.783,-0.550) (0.725,-0.510) (0.675,-0.475) (0.625,-0.441) (0.575,-0.406) (0.532,-0.375) (0.475,-0.336) (0.425,-0.300) (0.375,-0.266) (0.325,-0.230) (0.284,-0.200) (0.225,-0.160) (0.206,-0.144)}
\EQ{(-0.220,-0.155) (-0.275,-0.195) (-0.319,-0.225) (-0.375,-0.266) (-0.425,-0.300) (-0.475,-0.336) (-0.525,-0.371) (-0.568,-0.400) (-0.625,-0.441) (-0.674,-0.475) (-0.725,-0.511) (-0.775,-0.545) (-0.819,-0.575) (-0.875,-0.614) (-0.925,-0.648) (-0.975,-0.682) (-1.025,-0.716) (-1.075,-0.750) (-1.125,-0.783) (-1.175,-0.816) (-1.225,-0.848) (-1.275,-0.881) (-1.325,-0.913) (-1.375,-0.944) (-1.424,-0.975) (-1.475,-1.006) (-1.525,-1.037) (-1.575,-1.066) (-1.625,-1.096) (-1.675,-1.124) (-1.721,-1.150) (-1.766,-1.175) (-1.812,-1.200) (-1.859,-1.225) (-1.907,-1.250) (-1.957,-1.275) (-2.008,-1.300) (-2.060,-1.325) (-2.114,-1.350) (-2.171,-1.375) (-2.225,-1.398) (-2.275,-1.419) (-2.325,-1.438) (-2.375,-1.457) (-2.425,-1.475) (-2.475,-1.492) (-2.525,-1.508) (-2.575,-1.523) (-2.625,-1.538) (-2.675,-1.551) (-2.725,-1.563) (-2.775,-1.574) (-2.825,-1.585) (-2.875,-1.594) (-2.925,-1.602) (-2.975,-1.609) (-3.025,-1.615) (-3.075,-1.620) (-3.125,-1.623) (-3.175,-1.626) (-3.225,-1.627) (-3.275,-1.627) (-3.325,-1.626) (-3.375,-1.624) (-3.425,-1.620) (-3.475,-1.615) (-3.525,-1.609) (-3.575,-1.601) (-3.625,-1.592) (-3.675,-1.581) (-3.725,-1.569) (-3.775,-1.555) (-3.825,-1.539) (-3.875,-1.522) (-3.925,-1.503) (-3.975,-1.482) (-4.025,-1.459) (-4.075,-1.434) (-4.125,-1.407) (-4.175,-1.378) (-4.218,-1.350) (-4.275,-1.311) (-4.323,-1.275) (-4.375,-1.232) (-4.425,-1.188) (-4.464,-1.150) (-4.500,-1.113) (-4.535,-1.075) (-4.575,-1.028) (-4.616,-0.975) (-4.650,-0.928) (-4.685,-0.875) (-4.716,-0.825) (-4.744,-0.775) (-4.770,-0.725) (-4.793,-0.675) (-4.815,-0.625) (-4.835,-0.575) (-4.853,-0.525) (-4.869,-0.475) (-4.883,-0.425) (-4.895,-0.375) (-4.906,-0.325) (-4.916,-0.275) (-4.923,-0.225) (-4.929,-0.175) (-4.934,-0.125) (-4.937,-0.075) (-4.938,-0.025) (-4.938,0.050) (-4.936,0.100) (-4.932,0.150) (-4.926,0.200) (-4.920,0.250) (-4.911,0.300) (-4.901,0.350) (-4.889,0.400) (-4.876,0.450) (-4.861,0.500) (-4.844,0.550) (-4.825,0.600) (-4.804,0.650) (-4.782,0.700) (-4.757,0.750) (-4.730,0.800) (-4.701,0.850) (-4.669,0.900) (-4.635,0.950) (-4.600,0.996) (-4.556,1.050) (-4.512,1.100) (-4.475,1.139) (-4.438,1.175) (-4.400,1.210) (-4.354,1.250) (-4.300,1.292) (-4.255,1.325) (-4.200,1.362) (-4.150,1.393) (-4.100,1.421) (-4.050,1.447) (-4.000,1.471) (-3.950,1.493) (-3.900,1.513) (-3.850,1.531) (-3.800,1.547) (-3.750,1.562) (-3.700,1.575) (-3.650,1.586) (-3.600,1.596) (-3.550,1.605) (-3.500,1.612) (-3.450,1.618) (-3.400,1.622) (-3.350,1.625) (-3.300,1.627) (-3.250,1.627) (-3.200,1.627) (-3.150,1.625) (-3.100,1.622) (-3.050,1.617) (-3.000,1.612) (-2.950,1.605) (-2.900,1.598) (-2.850,1.589) (-2.800,1.580) (-2.750,1.569) (-2.700,1.557) (-2.650,1.544) (-2.600,1.531) (-2.550,1.516) (-2.500,1.500) (-2.450,1.484) (-2.400,1.466) (-2.350,1.448) (-2.300,1.429) (-2.250,1.409) (-2.200,1.388) (-2.150,1.366) (-2.100,1.344) (-2.050,1.320) (-2.000,1.296) (-1.950,1.272) (-1.900,1.246) (-1.850,1.220) (-1.800,1.194) (-1.750,1.167) (-1.700,1.139) (-1.650,1.110) (-1.600,1.081) (-1.550,1.051) (-1.506,1.025) (-1.450,0.991) (-1.400,0.960) (-1.350,0.929) (-1.305,0.900) (-1.250,0.865) (-1.200,0.832) (-1.151,0.800) (-1.100,0.766) (-1.050,0.733) (-1.001,0.700) (-0.950,0.665) (-0.900,0.631) (-0.855,0.600) (-0.800,0.563) (-0.750,0.528) (-0.700,0.493) (-0.650,0.459) (-0.603,0.425) (-0.550,0.389) (-0.500,0.353) (-0.450,0.318) (-0.400,0.283) (-0.354,0.250) (-0.300,0.213) (-0.250,0.177) (-0.220,0.155)}
\EQ{(0.206,0.144) (0.248,0.175) (0.300,0.213) (0.350,0.248) (0.400,0.283) (0.450,0.318) (0.496,0.350) (0.550,0.389) (0.600,0.423) (0.650,0.459) (0.700,0.493) (0.746,0.525) (0.800,0.563) (0.850,0.597) (0.891,0.625) (0.950,0.665) (1.000,0.699) (1.050,0.733) (1.100,0.766) (1.150,0.799) (1.200,0.832) (1.250,0.865) (1.300,0.897) (1.345,0.925) (1.400,0.960) (1.450,0.991) (1.500,1.021) (1.548,1.050) (1.600,1.081) (1.650,1.110) (1.700,1.139) (1.750,1.167) (1.800,1.194) (1.850,1.220) (1.900,1.246) (1.950,1.272) (2.000,1.296) (2.050,1.320) (2.100,1.344) (2.150,1.366) (2.200,1.388) (2.250,1.409) (2.300,1.429) (2.350,1.448) (2.400,1.466) (2.450,1.484) (2.500,1.500) (2.550,1.516) (2.600,1.531) (2.650,1.544) (2.700,1.557) (2.750,1.569) (2.800,1.580) (2.850,1.589) (2.900,1.598) (2.950,1.605) (3.000,1.612) (3.050,1.617) (3.100,1.622) (3.150,1.625) (3.200,1.627) (3.250,1.627) (3.300,1.627) (3.350,1.625) (3.400,1.622) (3.450,1.618) (3.500,1.612) (3.550,1.605) (3.600,1.596) (3.650,1.586) (3.700,1.575) (3.750,1.562) (3.800,1.547) (3.850,1.531) (3.900,1.513) (3.950,1.493) (4.000,1.471) (4.050,1.447) (4.100,1.421) (4.150,1.393) (4.200,1.362) (4.250,1.328) (4.300,1.292) (4.350,1.253) (4.400,1.210) (4.438,1.175) (4.475,1.139) (4.512,1.100) (4.550,1.058) (4.597,1.000) (4.635,0.950) (4.669,0.900) (4.700,0.851) (4.730,0.800) (4.757,0.750) (4.782,0.700) (4.804,0.650) (4.825,0.600) (4.844,0.550) (4.861,0.500) (4.876,0.450) (4.889,0.400) (4.901,0.350) (4.911,0.300) (4.920,0.250) (4.926,0.200) (4.932,0.150) (4.936,0.100) (4.938,0.050) (4.939,0.000) (4.938,-0.050) (4.936,-0.100) (4.932,-0.150) (4.926,-0.200) (4.920,-0.250) (4.911,-0.300) (4.901,-0.350) (4.889,-0.400) (4.876,-0.450) (4.861,-0.500) (4.844,-0.550) (4.825,-0.600) (4.804,-0.650) (4.782,-0.700) (4.757,-0.750) (4.730,-0.800) (4.701,-0.850) (4.669,-0.900) (4.635,-0.950) (4.600,-0.996) (4.556,-1.050) (4.512,-1.100) (4.475,-1.139) (4.438,-1.175) (4.400,-1.210) (4.354,-1.250) (4.300,-1.292) (4.255,-1.325) (4.200,-1.362) (4.150,-1.393) (4.100,-1.421) (4.050,-1.447) (4.000,-1.471) (3.950,-1.493) (3.900,-1.513) (3.850,-1.531) (3.800,-1.547) (3.750,-1.562) (3.700,-1.575) (3.650,-1.586) (3.600,-1.596) (3.550,-1.605) (3.500,-1.612) (3.450,-1.618) (3.400,-1.622) (3.350,-1.625) (3.300,-1.627) (3.250,-1.627) (3.200,-1.627) (3.150,-1.625) (3.100,-1.622) (3.050,-1.617) (3.000,-1.612) (2.950,-1.605) (2.900,-1.598) (2.850,-1.589) (2.800,-1.580) (2.750,-1.569) (2.700,-1.557) (2.650,-1.544) (2.600,-1.531) (2.550,-1.516) (2.500,-1.500) (2.450,-1.484) (2.400,-1.466) (2.350,-1.448) (2.300,-1.429) (2.250,-1.409) (2.200,-1.388) (2.150,-1.366) (2.100,-1.344) (2.050,-1.320) (2.000,-1.296) (1.950,-1.272) (1.900,-1.246) (1.850,-1.220) (1.800,-1.194) (1.750,-1.167) (1.700,-1.139) (1.650,-1.110) (1.600,-1.081) (1.550,-1.051) (1.506,-1.025) (1.450,-0.991) (1.400,-0.960) (1.350,-0.929) (1.305,-0.900) (1.250,-0.865) (1.200,-0.832) (1.151,-0.800) (1.100,-0.766) (1.050,-0.733) (1.001,-0.700) (0.950,-0.665) (0.900,-0.631) (0.855,-0.600) (0.800,-0.563) (0.750,-0.528) (0.700,-0.493) (0.650,-0.459) (0.603,-0.425) (0.550,-0.389) (0.500,-0.353) (0.450,-0.318) (0.400,-0.283) (0.354,-0.250) (0.300,-0.213) (0.250,-0.177) (0.206,-0.144)}
\EQ{(-0.220,-0.155) (-0.275,-0.195) (-0.319,-0.225) (-0.375,-0.266) (-0.425,-0.300) (-0.475,-0.336) (-0.525,-0.371) (-0.567,-0.400) (-0.625,-0.441) (-0.674,-0.475) (-0.725,-0.511) (-0.775,-0.546) (-0.818,-0.575) (-0.875,-0.615) (-0.925,-0.649) (-0.975,-0.683) (-1.025,-0.717) (-1.074,-0.750) (-1.125,-0.785) (-1.175,-0.818) (-1.224,-0.850) (-1.275,-0.884) (-1.325,-0.916) (-1.375,-0.948) (-1.417,-0.975) (-1.475,-1.011) (-1.525,-1.042) (-1.575,-1.073) (-1.620,-1.100) (-1.675,-1.133) (-1.725,-1.162) (-1.775,-1.190) (-1.825,-1.218) (-1.875,-1.246) (-1.925,-1.273) (-1.975,-1.299) (-2.025,-1.325) (-2.075,-1.350) (-2.125,-1.374) (-2.175,-1.398) (-2.225,-1.421) (-2.275,-1.443) (-2.325,-1.465) (-2.375,-1.486) (-2.425,-1.505) (-2.475,-1.525) (-2.525,-1.543) (-2.575,-1.560) (-2.625,-1.577) (-2.675,-1.593) (-2.725,-1.608) (-2.775,-1.622) (-2.825,-1.635) (-2.875,-1.647) (-2.925,-1.658) (-2.975,-1.668) (-3.025,-1.677) (-3.075,-1.686) (-3.125,-1.693) (-3.175,-1.699) (-3.225,-1.704) (-3.275,-1.708) (-3.325,-1.711) (-3.375,-1.713) (-3.425,-1.713) (-3.475,-1.713) (-3.525,-1.711) (-3.575,-1.708) (-3.625,-1.704) (-3.675,-1.698) (-3.725,-1.691) (-3.775,-1.683) (-3.825,-1.674) (-3.875,-1.663) (-3.925,-1.650) (-3.975,-1.636) (-4.025,-1.621) (-4.075,-1.604) (-4.125,-1.585) (-4.175,-1.564) (-4.225,-1.542) (-4.275,-1.518) (-4.325,-1.491) (-4.375,-1.462) (-4.425,-1.431) (-4.472,-1.400) (-4.525,-1.362) (-4.572,-1.325) (-4.625,-1.280) (-4.675,-1.234) (-4.725,-1.185) (-4.775,-1.130) (-4.822,-1.075) (-4.861,-1.025) (-4.897,-0.975) (-4.925,-0.932) (-4.961,-0.875) (-4.989,-0.825) (-5.015,-0.775) (-5.040,-0.725) (-5.062,-0.675) (-5.082,-0.625) (-5.100,-0.577) (-5.117,-0.525) (-5.133,-0.475) (-5.146,-0.425) (-5.158,-0.375) (-5.168,-0.325) (-5.177,-0.275) (-5.184,-0.225) (-5.190,-0.175) (-5.195,-0.125) (-5.197,-0.075) (-5.199,-0.025) (-5.198,0.050) (-5.196,0.100) (-5.193,0.150) (-5.187,0.200) (-5.181,0.250) (-5.173,0.300) (-5.163,0.350) (-5.152,0.400) (-5.140,0.450) (-5.125,0.500) (-5.109,0.550) (-5.092,0.600) (-5.072,0.650) (-5.051,0.700) (-5.028,0.750) (-5.002,0.800) (-4.975,0.850) (-4.946,0.900) (-4.914,0.950) (-4.879,1.000) (-4.842,1.050) (-4.801,1.100) (-4.758,1.150) (-4.710,1.200) (-4.659,1.250) (-4.602,1.300) (-4.550,1.343) (-4.507,1.375) (-4.450,1.415) (-4.400,1.447) (-4.353,1.475) (-4.309,1.500) (-4.260,1.525) (-4.208,1.550) (-4.150,1.575) (-4.100,1.595) (-4.050,1.613) (-4.000,1.629) (-3.950,1.644) (-3.900,1.657) (-3.850,1.668) (-3.800,1.679) (-3.750,1.687) (-3.700,1.695) (-3.650,1.701) (-3.600,1.706) (-3.550,1.710) (-3.500,1.712) (-3.450,1.713) (-3.400,1.713) (-3.350,1.712) (-3.300,1.710) (-3.250,1.706) (-3.200,1.702) (-3.150,1.696) (-3.100,1.689) (-3.050,1.682) (-3.000,1.673) (-2.950,1.663) (-2.900,1.653) (-2.850,1.641) (-2.800,1.628) (-2.750,1.615) (-2.700,1.600) (-2.650,1.585) (-2.600,1.569) (-2.550,1.552) (-2.500,1.534) (-2.450,1.515) (-2.400,1.496) (-2.350,1.475) (-2.300,1.454) (-2.250,1.432) (-2.200,1.410) (-2.150,1.386) (-2.100,1.362) (-2.050,1.338) (-2.000,1.312) (-1.950,1.286) (-1.900,1.259) (-1.850,1.232) (-1.800,1.204) (-1.750,1.176) (-1.705,1.150) (-1.650,1.118) (-1.600,1.088) (-1.550,1.058) (-1.500,1.027) (-1.457,1.000) (-1.400,0.964) (-1.350,0.932) (-1.300,0.900) (-1.250,0.867) (-1.200,0.835) (-1.150,0.801) (-1.100,0.768) (-1.050,0.734) (-1.000,0.700) (-0.950,0.666) (-0.900,0.632) (-0.854,0.600) (-0.800,0.563) (-0.750,0.528) (-0.700,0.494) (-0.650,0.459) (-0.603,0.425) (-0.550,0.389) (-0.500,0.353) (-0.450,0.318) (-0.400,0.283) (-0.354,0.250) (-0.300,0.213) (-0.250,0.177) (-0.220,0.155)}
\EQ{(0.206,0.144) (0.248,0.175) (0.300,0.213) (0.350,0.248) (0.400,0.283) (0.450,0.318) (0.496,0.350) (0.550,0.389) (0.600,0.423) (0.650,0.459) (0.700,0.494) (0.746,0.525) (0.800,0.563) (0.850,0.598) (0.900,0.632) (0.950,0.666) (1.000,0.700) (1.050,0.734) (1.100,0.768) (1.148,0.800) (1.200,0.835) (1.250,0.867) (1.300,0.900) (1.350,0.932) (1.400,0.964) (1.450,0.996) (1.497,1.025) (1.550,1.058) (1.600,1.088) (1.650,1.118) (1.700,1.147) (1.748,1.175) (1.792,1.200) (1.837,1.225) (1.883,1.250) (1.929,1.275) (1.975,1.299) (2.025,1.325) (2.075,1.350) (2.125,1.374) (2.175,1.398) (2.225,1.421) (2.275,1.443) (2.325,1.465) (2.375,1.486) (2.425,1.505) (2.475,1.525) (2.525,1.543) (2.575,1.560) (2.625,1.577) (2.675,1.593) (2.725,1.608) (2.775,1.622) (2.825,1.635) (2.875,1.647) (2.925,1.658) (2.975,1.668) (3.025,1.677) (3.075,1.686) (3.125,1.693) (3.175,1.699) (3.225,1.704) (3.275,1.708) (3.325,1.711) (3.375,1.713) (3.425,1.713) (3.475,1.713) (3.525,1.711) (3.575,1.708) (3.625,1.704) (3.675,1.698) (3.725,1.691) (3.775,1.683) (3.825,1.674) (3.875,1.663) (3.925,1.650) (3.975,1.636) (4.025,1.621) (4.075,1.604) (4.125,1.585) (4.175,1.564) (4.225,1.542) (4.275,1.518) (4.325,1.491) (4.375,1.462) (4.425,1.431) (4.472,1.400) (4.525,1.362) (4.572,1.325) (4.625,1.280) (4.675,1.234) (4.725,1.185) (4.775,1.130) (4.822,1.075) (4.861,1.025) (4.897,0.975) (4.925,0.932) (4.961,0.875) (4.989,0.825) (5.015,0.775) (5.040,0.725) (5.062,0.675) (5.082,0.625) (5.100,0.577) (5.117,0.525) (5.133,0.475) (5.146,0.425) (5.158,0.375) (5.168,0.325) (5.177,0.275) (5.184,0.225) (5.190,0.175) (5.195,0.125) (5.197,0.075) (5.199,0.025) (5.198,-0.050) (5.196,-0.100) (5.193,-0.150) (5.187,-0.200) (5.181,-0.250) (5.173,-0.300) (5.163,-0.350) (5.152,-0.400) (5.140,-0.450) (5.125,-0.500) (5.109,-0.550) (5.092,-0.600) (5.072,-0.650) (5.051,-0.700) (5.028,-0.750) (5.002,-0.800) (4.975,-0.850) (4.946,-0.900) (4.914,-0.950) (4.879,-1.000) (4.842,-1.050) (4.801,-1.100) (4.758,-1.150) (4.710,-1.200) (4.659,-1.250) (4.602,-1.300) (4.550,-1.343) (4.507,-1.375) (4.450,-1.415) (4.400,-1.447) (4.353,-1.475) (4.309,-1.500) (4.260,-1.525) (4.208,-1.550) (4.150,-1.575) (4.100,-1.595) (4.050,-1.613) (4.000,-1.629) (3.950,-1.644) (3.900,-1.657) (3.850,-1.668) (3.800,-1.679) (3.750,-1.687) (3.700,-1.695) (3.650,-1.701) (3.600,-1.706) (3.550,-1.710) (3.500,-1.712) (3.450,-1.713) (3.400,-1.713) (3.350,-1.712) (3.300,-1.710) (3.250,-1.706) (3.200,-1.702) (3.150,-1.696) (3.100,-1.689) (3.050,-1.682) (3.000,-1.673) (2.950,-1.663) (2.900,-1.653) (2.850,-1.641) (2.800,-1.628) (2.750,-1.615) (2.700,-1.600) (2.650,-1.585) (2.600,-1.569) (2.550,-1.552) (2.500,-1.534) (2.450,-1.515) (2.400,-1.496) (2.350,-1.475) (2.300,-1.454) (2.250,-1.432) (2.200,-1.410) (2.150,-1.386) (2.100,-1.362) (2.050,-1.338) (2.000,-1.312) (1.950,-1.286) (1.900,-1.259) (1.850,-1.232) (1.800,-1.204) (1.750,-1.176) (1.705,-1.150) (1.650,-1.118) (1.600,-1.088) (1.550,-1.058) (1.500,-1.027) (1.457,-1.000) (1.400,-0.964) (1.350,-0.932) (1.300,-0.900) (1.250,-0.867) (1.200,-0.835) (1.150,-0.801) (1.100,-0.768) (1.050,-0.734) (1.000,-0.700) (0.950,-0.666) (0.900,-0.632) (0.854,-0.600) (0.800,-0.563) (0.750,-0.528) (0.700,-0.494) (0.650,-0.459) (0.603,-0.425) (0.550,-0.389) (0.500,-0.353) (0.450,-0.318) (0.400,-0.283) (0.354,-0.250) (0.300,-0.213) (0.250,-0.177) (0.206,-0.144)}
\EQ{(-0.220,-0.155) (-0.275,-0.195) (-0.319,-0.225) (-0.375,-0.266) (-0.425,-0.300) (-0.475,-0.336) (-0.525,-0.371) (-0.567,-0.400) (-0.625,-0.441) (-0.674,-0.475) (-0.725,-0.511) (-0.775,-0.546) (-0.817,-0.575) (-0.875,-0.616) (-0.925,-0.650) (-0.975,-0.684) (-1.025,-0.719) (-1.072,-0.750) (-1.125,-0.787) (-1.175,-0.820) (-1.220,-0.850) (-1.275,-0.887) (-1.325,-0.920) (-1.372,-0.950) (-1.425,-0.985) (-1.475,-1.017) (-1.525,-1.048) (-1.568,-1.075) (-1.625,-1.110) (-1.675,-1.141) (-1.725,-1.171) (-1.774,-1.200) (-1.825,-1.230) (-1.875,-1.259) (-1.925,-1.287) (-1.975,-1.314) (-2.025,-1.342) (-2.075,-1.368) (-2.125,-1.394) (-2.175,-1.420) (-2.225,-1.444) (-2.275,-1.468) (-2.325,-1.492) (-2.375,-1.515) (-2.425,-1.537) (-2.475,-1.558) (-2.525,-1.579) (-2.575,-1.599) (-2.625,-1.618) (-2.675,-1.637) (-2.725,-1.654) (-2.775,-1.671) (-2.825,-1.687) (-2.875,-1.702) (-2.925,-1.717) (-2.975,-1.730) (-3.025,-1.743) (-3.075,-1.755) (-3.125,-1.765) (-3.175,-1.775) (-3.225,-1.784) (-3.275,-1.792) (-3.325,-1.799) (-3.375,-1.805) (-3.425,-1.811) (-3.475,-1.815) (-3.525,-1.818) (-3.575,-1.820) (-3.625,-1.820) (-3.675,-1.820) (-3.725,-1.819) (-3.775,-1.816) (-3.825,-1.813) (-3.875,-1.808) (-3.925,-1.802) (-3.975,-1.795) (-4.025,-1.786) (-4.075,-1.776) (-4.125,-1.765) (-4.175,-1.753) (-4.225,-1.739) (-4.275,-1.723) (-4.325,-1.706) (-4.375,-1.688) (-4.425,-1.667) (-4.475,-1.645) (-4.525,-1.622) (-4.575,-1.596) (-4.625,-1.568) (-4.675,-1.538) (-4.725,-1.506) (-4.770,-1.475) (-4.825,-1.435) (-4.868,-1.400) (-4.925,-1.352) (-4.975,-1.305) (-5.025,-1.255) (-5.075,-1.200) (-5.118,-1.150) (-5.150,-1.109) (-5.194,-1.050) (-5.225,-1.004) (-5.259,-0.950) (-5.289,-0.900) (-5.316,-0.850) (-5.341,-0.800) (-5.365,-0.750) (-5.386,-0.700) (-5.406,-0.650) (-5.424,-0.600) (-5.441,-0.550) (-5.456,-0.500) (-5.469,-0.450) (-5.481,-0.400) (-5.491,-0.350) (-5.500,-0.300) (-5.508,-0.250) (-5.514,-0.200) (-5.519,-0.150) (-5.522,-0.100) (-5.524,-0.050) (-5.525,0.000) (-5.524,0.050) (-5.522,0.100) (-5.519,0.150) (-5.514,0.200) (-5.508,0.250) (-5.500,0.300) (-5.491,0.350) (-5.481,0.400) (-5.469,0.450) (-5.456,0.500) (-5.441,0.550) (-5.424,0.600) (-5.406,0.650) (-5.386,0.700) (-5.365,0.750) (-5.341,0.800) (-5.316,0.850) (-5.289,0.900) (-5.259,0.950) (-5.228,1.000) (-5.194,1.050) (-5.157,1.100) (-5.125,1.141) (-5.075,1.200) (-5.030,1.250) (-4.980,1.300) (-4.927,1.350) (-4.875,1.395) (-4.825,1.435) (-4.775,1.472) (-4.725,1.506) (-4.675,1.538) (-4.625,1.568) (-4.575,1.596) (-4.525,1.622) (-4.475,1.645) (-4.425,1.667) (-4.375,1.688) (-4.325,1.706) (-4.275,1.723) (-4.225,1.739) (-4.175,1.753) (-4.125,1.765) (-4.075,1.776) (-4.025,1.786) (-3.975,1.795) (-3.925,1.802) (-3.875,1.808) (-3.825,1.813) (-3.775,1.816) (-3.725,1.819) (-3.675,1.820) (-3.625,1.820) (-3.575,1.820) (-3.525,1.818) (-3.475,1.815) (-3.425,1.811) (-3.375,1.805) (-3.325,1.799) (-3.275,1.792) (-3.225,1.784) (-3.175,1.775) (-3.125,1.765) (-3.075,1.755) (-3.025,1.743) (-2.975,1.730) (-2.925,1.717) (-2.875,1.702) (-2.825,1.687) (-2.775,1.671) (-2.725,1.654) (-2.675,1.637) (-2.625,1.618) (-2.578,1.600) (-2.525,1.579) (-2.475,1.558) (-2.425,1.537) (-2.375,1.515) (-2.325,1.492) (-2.275,1.468) (-2.225,1.444) (-2.175,1.420) (-2.125,1.394) (-2.075,1.368) (-2.025,1.342) (-1.975,1.314) (-1.925,1.287) (-1.875,1.259) (-1.825,1.230) (-1.775,1.201) (-1.725,1.171) (-1.675,1.141) (-1.625,1.110) (-1.575,1.079) (-1.528,1.050) (-1.475,1.017) (-1.425,0.985) (-1.375,0.952) (-1.325,0.920) (-1.275,0.887) (-1.225,0.854) (-1.183,0.825) (-1.125,0.787) (-1.075,0.753) (-1.025,0.719) (-0.975,0.684) (-0.925,0.650) (-0.875,0.616) (-0.825,0.581) (-0.781,0.550) (-0.725,0.511) (-0.675,0.476) (-0.625,0.441) (-0.575,0.406) (-0.532,0.375) (-0.475,0.336) (-0.425,0.300) (-0.375,0.266) (-0.325,0.230) (-0.284,0.200) (-0.225,0.160) (-0.220,0.155)}
\EQ{(0.206,0.144) (0.248,0.175) (0.300,0.213) (0.350,0.248) (0.400,0.283) (0.450,0.318) (0.496,0.350) (0.550,0.389) (0.600,0.424) (0.650,0.459) (0.700,0.494) (0.745,0.525) (0.800,0.564) (0.850,0.598) (0.900,0.633) (0.950,0.667) (0.998,0.700) (1.050,0.736) (1.100,0.770) (1.145,0.800) (1.200,0.837) (1.250,0.870) (1.296,0.900) (1.350,0.936) (1.400,0.968) (1.449,1.000) (1.500,1.032) (1.550,1.064) (1.600,1.095) (1.649,1.125) (1.700,1.156) (1.750,1.186) (1.800,1.215) (1.850,1.244) (1.900,1.273) (1.949,1.300) (1.995,1.325) (2.041,1.350) (2.088,1.375) (2.137,1.400) (2.186,1.425) (2.237,1.450) (2.289,1.475) (2.343,1.500) (2.398,1.525) (2.450,1.548) (2.500,1.569) (2.550,1.589) (2.600,1.609) (2.650,1.628) (2.700,1.646) (2.750,1.663) (2.800,1.679) (2.850,1.695) (2.900,1.710) (2.950,1.724) (3.000,1.737) (3.050,1.749) (3.100,1.760) (3.150,1.771) (3.200,1.780) (3.250,1.788) (3.300,1.796) (3.350,1.803) (3.400,1.808) (3.450,1.813) (3.500,1.816) (3.550,1.819) (3.600,1.820) (3.650,1.820) (3.700,1.820) (3.750,1.818) (3.800,1.815) (3.850,1.811) (3.900,1.805) (3.950,1.799) (4.000,1.791) (4.050,1.782) (4.100,1.771) (4.150,1.759) (4.200,1.746) (4.250,1.731) (4.300,1.715) (4.350,1.697) (4.400,1.678) (4.450,1.657) (4.500,1.634) (4.550,1.609) (4.600,1.582) (4.650,1.553) (4.696,1.525) (4.750,1.489) (4.800,1.453) (4.850,1.415) (4.898,1.375) (4.950,1.329) (5.000,1.280) (5.050,1.228) (5.097,1.175) (5.138,1.125) (5.175,1.076) (5.211,1.025) (5.244,0.975) (5.274,0.925) (5.300,0.880) (5.325,0.832) (5.350,0.782) (5.375,0.726) (5.396,0.675) (5.415,0.625) (5.433,0.575) (5.448,0.525) (5.462,0.475) (5.475,0.425) (5.486,0.375) (5.496,0.325) (5.504,0.275) (5.511,0.225) (5.516,0.175) (5.521,0.125) (5.523,0.075) (5.525,0.025) (5.524,-0.050) (5.522,-0.100) (5.519,-0.150) (5.514,-0.200) (5.508,-0.250) (5.500,-0.300) (5.491,-0.350) (5.481,-0.400) (5.469,-0.450) (5.456,-0.500) (5.441,-0.550) (5.424,-0.600) (5.406,-0.650) (5.386,-0.700) (5.365,-0.750) (5.341,-0.800) (5.316,-0.850) (5.289,-0.900) (5.259,-0.950) (5.228,-1.000) (5.194,-1.050) (5.157,-1.100) (5.125,-1.141) (5.075,-1.200) (5.030,-1.250) (4.980,-1.300) (4.927,-1.350) (4.875,-1.395) (4.825,-1.435) (4.775,-1.472) (4.725,-1.506) (4.675,-1.538) (4.625,-1.568) (4.575,-1.596) (4.525,-1.622) (4.475,-1.645) (4.425,-1.667) (4.375,-1.688) (4.325,-1.706) (4.275,-1.723) (4.225,-1.739) (4.175,-1.753) (4.125,-1.765) (4.075,-1.776) (4.025,-1.786) (3.975,-1.795) (3.925,-1.802) (3.875,-1.808) (3.825,-1.813) (3.775,-1.816) (3.725,-1.819) (3.675,-1.820) (3.625,-1.820) (3.575,-1.820) (3.525,-1.818) (3.475,-1.815) (3.425,-1.811) (3.375,-1.805) (3.325,-1.799) (3.275,-1.792) (3.225,-1.784) (3.175,-1.775) (3.125,-1.765) (3.075,-1.755) (3.025,-1.743) (2.975,-1.730) (2.925,-1.717) (2.875,-1.702) (2.825,-1.687) (2.775,-1.671) (2.725,-1.654) (2.675,-1.637) (2.625,-1.618) (2.578,-1.600) (2.525,-1.579) (2.475,-1.558) (2.425,-1.537) (2.375,-1.515) (2.325,-1.492) (2.275,-1.468) (2.225,-1.444) (2.175,-1.420) (2.125,-1.394) (2.075,-1.368) (2.025,-1.342) (1.975,-1.314) (1.925,-1.287) (1.875,-1.259) (1.825,-1.230) (1.775,-1.201) (1.725,-1.171) (1.675,-1.141) (1.625,-1.110) (1.575,-1.079) (1.528,-1.050) (1.475,-1.017) (1.425,-0.985) (1.375,-0.952) (1.325,-0.920) (1.275,-0.887) (1.225,-0.854) (1.183,-0.825) (1.125,-0.787) (1.075,-0.753) (1.025,-0.719) (0.975,-0.684) (0.925,-0.650) (0.875,-0.616) (0.825,-0.581) (0.781,-0.550) (0.725,-0.511) (0.675,-0.476) (0.625,-0.441) (0.575,-0.406) (0.532,-0.375) (0.475,-0.336) (0.425,-0.300) (0.375,-0.266) (0.325,-0.230) (0.284,-0.200) (0.225,-0.160) (0.206,-0.144)}
\EQ{(-0.220,-0.155) (-0.275,-0.195) (-0.319,-0.225) (-0.375,-0.266) (-0.425,-0.300) (-0.475,-0.336) (-0.525,-0.371) (-0.567,-0.400) (-0.625,-0.442) (-0.673,-0.475) (-0.725,-0.512) (-0.775,-0.547) (-0.817,-0.575) (-0.875,-0.616) (-0.924,-0.650) (-0.975,-0.686) (-1.025,-0.720) (-1.069,-0.750) (-1.125,-0.788) (-1.175,-0.822) (-1.225,-0.856) (-1.275,-0.890) (-1.325,-0.923) (-1.375,-0.956) (-1.425,-0.989) (-1.475,-1.022) (-1.519,-1.050) (-1.575,-1.086) (-1.625,-1.118) (-1.675,-1.149) (-1.725,-1.180) (-1.775,-1.211) (-1.825,-1.242) (-1.875,-1.271) (-1.924,-1.300) (-1.975,-1.330) (-2.025,-1.359) (-2.075,-1.387) (-2.125,-1.415) (-2.175,-1.442) (-2.225,-1.468) (-2.275,-1.495) (-2.325,-1.520) (-2.375,-1.545) (-2.425,-1.570) (-2.475,-1.593) (-2.525,-1.617) (-2.575,-1.639) (-2.625,-1.661) (-2.675,-1.682) (-2.725,-1.703) (-2.775,-1.723) (-2.825,-1.742) (-2.875,-1.761) (-2.925,-1.778) (-2.975,-1.796) (-3.025,-1.812) (-3.075,-1.827) (-3.125,-1.842) (-3.175,-1.856) (-3.225,-1.869) (-3.275,-1.881) (-3.325,-1.893) (-3.375,-1.904) (-3.425,-1.913) (-3.475,-1.922) (-3.525,-1.930) (-3.575,-1.937) (-3.625,-1.944) (-3.675,-1.949) (-3.725,-1.953) (-3.775,-1.957) (-3.825,-1.959) (-3.875,-1.961) (-3.925,-1.961) (-3.975,-1.960) (-4.025,-1.959) (-4.075,-1.956) (-4.125,-1.952) (-4.175,-1.948) (-4.225,-1.942) (-4.275,-1.934) (-4.325,-1.926) (-4.375,-1.917) (-4.425,-1.906) (-4.475,-1.894) (-4.525,-1.880) (-4.575,-1.866) (-4.624,-1.850) (-4.675,-1.832) (-4.725,-1.813) (-4.775,-1.793) (-4.825,-1.770) (-4.875,-1.746) (-4.925,-1.721) (-4.975,-1.693) (-5.025,-1.664) (-5.075,-1.632) (-5.123,-1.600) (-5.175,-1.562) (-5.223,-1.525) (-5.275,-1.482) (-5.325,-1.438) (-5.365,-1.400) (-5.415,-1.350) (-5.450,-1.312) (-5.500,-1.254) (-5.544,-1.200) (-5.575,-1.158) (-5.616,-1.100) (-5.649,-1.050) (-5.675,-1.007) (-5.708,-0.950) (-5.735,-0.900) (-5.759,-0.850) (-5.782,-0.800) (-5.804,-0.750) (-5.823,-0.700) (-5.842,-0.650) (-5.858,-0.600) (-5.874,-0.550) (-5.887,-0.500) (-5.900,-0.450) (-5.911,-0.400) (-5.920,-0.350) (-5.929,-0.300) (-5.936,-0.250) (-5.941,-0.200) (-5.946,-0.150) (-5.949,-0.100) (-5.951,-0.050) (-5.951,0.000) (-5.951,0.050) (-5.949,0.100) (-5.946,0.150) (-5.941,0.200) (-5.936,0.250) (-5.929,0.300) (-5.920,0.350) (-5.911,0.400) (-5.900,0.450) (-5.887,0.500) (-5.874,0.550) (-5.858,0.600) (-5.842,0.650) (-5.823,0.700) (-5.804,0.750) (-5.782,0.800) (-5.759,0.850) (-5.735,0.900) (-5.708,0.950) (-5.679,1.000) (-5.650,1.048) (-5.616,1.100) (-5.581,1.150) (-5.550,1.192) (-5.504,1.250) (-5.461,1.300) (-5.425,1.339) (-5.375,1.390) (-5.325,1.438) (-5.284,1.475) (-5.225,1.524) (-5.175,1.562) (-5.125,1.599) (-5.075,1.632) (-5.025,1.664) (-4.975,1.693) (-4.925,1.721) (-4.875,1.746) (-4.825,1.770) (-4.775,1.793) (-4.725,1.813) (-4.675,1.832) (-4.625,1.850) (-4.575,1.866) (-4.525,1.880) (-4.475,1.894) (-4.425,1.906) (-4.375,1.917) (-4.325,1.926) (-4.275,1.934) (-4.225,1.942) (-4.175,1.948) (-4.125,1.952) (-4.075,1.956) (-4.025,1.959) (-3.975,1.960) (-3.925,1.961) (-3.875,1.961) (-3.825,1.959) (-3.775,1.957) (-3.725,1.953) (-3.675,1.949) (-3.625,1.944) (-3.575,1.937) (-3.525,1.930) (-3.475,1.922) (-3.425,1.913) (-3.375,1.904) (-3.325,1.893) (-3.275,1.881) (-3.225,1.869) (-3.175,1.856) (-3.125,1.842) (-3.075,1.827) (-3.025,1.812) (-2.975,1.796) (-2.925,1.778) (-2.875,1.761) (-2.825,1.742) (-2.775,1.723) (-2.725,1.703) (-2.675,1.682) (-2.625,1.661) (-2.575,1.639) (-2.525,1.617) (-2.475,1.593) (-2.425,1.570) (-2.375,1.545) (-2.325,1.520) (-2.275,1.495) (-2.225,1.468) (-2.175,1.442) (-2.125,1.415) (-2.075,1.387) (-2.025,1.359) (-1.975,1.330) (-1.925,1.301) (-1.881,1.275) (-1.825,1.242) (-1.775,1.211) (-1.725,1.180) (-1.676,1.150) (-1.625,1.118) (-1.575,1.086) (-1.525,1.054) (-1.480,1.025) (-1.425,0.989) (-1.375,0.956) (-1.328,0.925) (-1.275,0.890) (-1.225,0.856) (-1.179,0.825) (-1.125,0.788) (-1.075,0.754) (-1.033,0.725) (-0.975,0.686) (-0.925,0.651) (-0.875,0.616) (-0.825,0.581) (-0.781,0.550) (-0.725,0.512) (-0.675,0.476) (-0.625,0.442) (-0.575,0.406) (-0.531,0.375) (-0.475,0.336) (-0.425,0.300) (-0.375,0.266) (-0.325,0.230) (-0.284,0.200) (-0.225,0.160) (-0.220,0.155)}
\EQ{(0.206,0.144) (0.248,0.175) (0.300,0.213) (0.350,0.248) (0.400,0.283) (0.450,0.318) (0.496,0.350) (0.550,0.389) (0.600,0.424) (0.650,0.459) (0.700,0.494) (0.745,0.525) (0.800,0.564) (0.850,0.599) (0.900,0.634) (0.950,0.668) (0.997,0.700) (1.050,0.737) (1.100,0.771) (1.143,0.800) (1.200,0.839) (1.250,0.873) (1.300,0.906) (1.350,0.940) (1.400,0.973) (1.442,1.000) (1.500,1.038) (1.550,1.070) (1.597,1.100) (1.650,1.134) (1.700,1.165) (1.750,1.196) (1.798,1.225) (1.850,1.257) (1.900,1.286) (1.950,1.316) (2.000,1.344) (2.050,1.373) (2.099,1.400) (2.144,1.425) (2.191,1.450) (2.238,1.475) (2.286,1.500) (2.335,1.525) (2.385,1.550) (2.437,1.575) (2.489,1.600) (2.544,1.625) (2.600,1.650) (2.650,1.672) (2.700,1.693) (2.750,1.713) (2.800,1.733) (2.850,1.751) (2.900,1.770) (2.950,1.787) (3.000,1.804) (3.050,1.820) (3.100,1.835) (3.150,1.849) (3.200,1.863) (3.249,1.875) (3.300,1.887) (3.350,1.898) (3.400,1.909) (3.450,1.918) (3.500,1.926) (3.550,1.934) (3.600,1.941) (3.650,1.946) (3.700,1.951) (3.750,1.955) (3.800,1.958) (3.850,1.960) (3.900,1.961) (3.950,1.961) (4.000,1.960) (4.050,1.958) (4.100,1.954) (4.150,1.950) (4.200,1.945) (4.250,1.938) (4.300,1.930) (4.350,1.922) (4.400,1.911) (4.450,1.900) (4.500,1.887) (4.550,1.873) (4.600,1.858) (4.650,1.841) (4.700,1.823) (4.750,1.803) (4.800,1.782) (4.850,1.759) (4.900,1.734) (4.950,1.707) (5.000,1.679) (5.047,1.650) (5.100,1.616) (5.150,1.581) (5.191,1.550) (5.250,1.503) (5.300,1.461) (5.339,1.425) (5.390,1.375) (5.438,1.325) (5.475,1.284) (5.524,1.225) (5.563,1.175) (5.599,1.125) (5.633,1.075) (5.664,1.025) (5.694,0.975) (5.721,0.925) (5.747,0.875) (5.771,0.825) (5.793,0.775) (5.814,0.725) (5.833,0.675) (5.850,0.625) (5.866,0.575) (5.881,0.525) (5.894,0.475) (5.905,0.425) (5.916,0.375) (5.925,0.325) (5.932,0.275) (5.939,0.225) (5.944,0.175) (5.947,0.125) (5.950,0.075) (5.951,0.025) (5.951,-0.050) (5.949,-0.100) (5.946,-0.150) (5.941,-0.200) (5.936,-0.250) (5.929,-0.300) (5.920,-0.350) (5.911,-0.400) (5.900,-0.450) (5.887,-0.500) (5.874,-0.550) (5.858,-0.600) (5.842,-0.650) (5.823,-0.700) (5.804,-0.750) (5.782,-0.800) (5.759,-0.850) (5.735,-0.900) (5.708,-0.950) (5.679,-1.000) (5.650,-1.048) (5.616,-1.100) (5.581,-1.150) (5.550,-1.192) (5.504,-1.250) (5.461,-1.300) (5.425,-1.339) (5.375,-1.390) (5.325,-1.438) (5.284,-1.475) (5.225,-1.524) (5.175,-1.562) (5.125,-1.599) (5.075,-1.632) (5.025,-1.664) (4.975,-1.693) (4.925,-1.721) (4.875,-1.746) (4.825,-1.770) (4.775,-1.793) (4.725,-1.813) (4.675,-1.832) (4.625,-1.850) (4.575,-1.866) (4.525,-1.880) (4.475,-1.894) (4.425,-1.906) (4.375,-1.917) (4.325,-1.926) (4.275,-1.934) (4.225,-1.942) (4.175,-1.948) (4.125,-1.952) (4.075,-1.956) (4.025,-1.959) (3.975,-1.960) (3.925,-1.961) (3.875,-1.961) (3.825,-1.959) (3.775,-1.957) (3.725,-1.953) (3.675,-1.949) (3.625,-1.944) (3.575,-1.937) (3.525,-1.930) (3.475,-1.922) (3.425,-1.913) (3.375,-1.904) (3.325,-1.893) (3.275,-1.881) (3.225,-1.869) (3.175,-1.856) (3.125,-1.842) (3.075,-1.827) (3.025,-1.812) (2.975,-1.796) (2.925,-1.778) (2.875,-1.761) (2.825,-1.742) (2.775,-1.723) (2.725,-1.703) (2.675,-1.682) (2.625,-1.661) (2.575,-1.639) (2.525,-1.617) (2.475,-1.593) (2.425,-1.570) (2.375,-1.545) (2.325,-1.520) (2.275,-1.495) (2.225,-1.468) (2.175,-1.442) (2.125,-1.415) (2.075,-1.387) (2.025,-1.359) (1.975,-1.330) (1.925,-1.301) (1.881,-1.275) (1.825,-1.242) (1.775,-1.211) (1.725,-1.180) (1.676,-1.150) (1.625,-1.118) (1.575,-1.086) (1.525,-1.054) (1.480,-1.025) (1.425,-0.989) (1.375,-0.956) (1.328,-0.925) (1.275,-0.890) (1.225,-0.856) (1.179,-0.825) (1.125,-0.788) (1.075,-0.754) (1.033,-0.725) (0.975,-0.686) (0.925,-0.651) (0.875,-0.616) (0.825,-0.581) (0.781,-0.550) (0.725,-0.512) (0.675,-0.476) (0.625,-0.442) (0.575,-0.406) (0.531,-0.375) (0.475,-0.336) (0.425,-0.300) (0.375,-0.266) (0.325,-0.230) (0.284,-0.200) (0.225,-0.160) (0.206,-0.144)}
\EQ{(-0.220,-0.155) (-0.275,-0.195) (-0.319,-0.225) (-0.375,-0.266) (-0.425,-0.300) (-0.475,-0.336) (-0.525,-0.371) (-0.567,-0.400) (-0.625,-0.442) (-0.673,-0.475) (-0.725,-0.512) (-0.775,-0.547) (-0.825,-0.582) (-0.875,-0.617) (-0.923,-0.650) (-0.975,-0.687) (-1.025,-0.721) (-1.067,-0.750) (-1.125,-0.790) (-1.175,-0.824) (-1.225,-0.859) (-1.275,-0.893) (-1.323,-0.925) (-1.375,-0.960) (-1.425,-0.994) (-1.472,-1.025) (-1.525,-1.060) (-1.575,-1.093) (-1.624,-1.125) (-1.675,-1.158) (-1.725,-1.190) (-1.775,-1.222) (-1.820,-1.250) (-1.875,-1.285) (-1.925,-1.316) (-1.975,-1.346) (-2.023,-1.375) (-2.075,-1.406) (-2.125,-1.436) (-2.175,-1.465) (-2.225,-1.493) (-2.275,-1.521) (-2.325,-1.549) (-2.372,-1.575) (-2.419,-1.600) (-2.466,-1.625) (-2.514,-1.650) (-2.563,-1.675) (-2.613,-1.700) (-2.664,-1.725) (-2.716,-1.750) (-2.770,-1.775) (-2.825,-1.800) (-2.875,-1.822) (-2.925,-1.844) (-2.975,-1.865) (-3.025,-1.885) (-3.075,-1.904) (-3.125,-1.923) (-3.175,-1.942) (-3.225,-1.959) (-3.275,-1.976) (-3.325,-1.993) (-3.375,-2.008) (-3.425,-2.023) (-3.475,-2.037) (-3.525,-2.051) (-3.575,-2.064) (-3.625,-2.076) (-3.675,-2.087) (-3.725,-2.097) (-3.775,-2.107) (-3.825,-2.116) (-3.875,-2.124) (-3.925,-2.131) (-3.975,-2.138) (-4.025,-2.143) (-4.075,-2.148) (-4.125,-2.152) (-4.175,-2.155) (-4.225,-2.157) (-4.275,-2.158) (-4.325,-2.158) (-4.375,-2.158) (-4.425,-2.156) (-4.475,-2.154) (-4.525,-2.150) (-4.575,-2.146) (-4.625,-2.140) (-4.675,-2.134) (-4.725,-2.126) (-4.775,-2.117) (-4.825,-2.108) (-4.875,-2.097) (-4.925,-2.085) (-4.975,-2.071) (-5.025,-2.057) (-5.075,-2.041) (-5.125,-2.024) (-5.175,-2.006) (-5.225,-1.986) (-5.275,-1.965) (-5.325,-1.942) (-5.375,-1.918) (-5.425,-1.892) (-5.475,-1.864) (-5.525,-1.835) (-5.575,-1.804) (-5.618,-1.775) (-5.675,-1.735) (-5.722,-1.700) (-5.775,-1.658) (-5.814,-1.625) (-5.869,-1.575) (-5.921,-1.525) (-5.970,-1.475) (-6.016,-1.425) (-6.050,-1.385) (-6.098,-1.325) (-6.136,-1.275) (-6.172,-1.225) (-6.200,-1.183) (-6.236,-1.125) (-6.266,-1.075) (-6.294,-1.025) (-6.320,-0.975) (-6.344,-0.925) (-6.367,-0.875) (-6.388,-0.825) (-6.408,-0.775) (-6.426,-0.725) (-6.443,-0.675) (-6.459,-0.625) (-6.473,-0.575) (-6.486,-0.525) (-6.498,-0.475) (-6.509,-0.425) (-6.518,-0.375) (-6.526,-0.325) (-6.533,-0.275) (-6.539,-0.225) (-6.543,-0.175) (-6.547,-0.125) (-6.549,-0.075) (-6.550,-0.025) (-6.550,0.037) (-6.548,0.100) (-6.545,0.150) (-6.541,0.200) (-6.536,0.250) (-6.530,0.300) (-6.522,0.350) (-6.513,0.400) (-6.503,0.450) (-6.492,0.500) (-6.480,0.550) (-6.466,0.600) (-6.451,0.650) (-6.435,0.700) (-6.417,0.750) (-6.398,0.800) (-6.378,0.850) (-6.356,0.900) (-6.332,0.950) (-6.307,1.000) (-6.280,1.050) (-6.251,1.100) (-6.225,1.143) (-6.189,1.200) (-6.154,1.250) (-6.118,1.300) (-6.079,1.350) (-6.037,1.400) (-6.000,1.442) (-5.950,1.496) (-5.900,1.546) (-5.850,1.593) (-5.800,1.637) (-5.754,1.675) (-5.700,1.717) (-5.654,1.750) (-5.600,1.787) (-5.550,1.819) (-5.500,1.850) (-5.456,1.875) (-5.409,1.900) (-5.360,1.925) (-5.307,1.950) (-5.251,1.975) (-5.200,1.996) (-5.150,2.015) (-5.100,2.033) (-5.050,2.049) (-5.000,2.064) (-4.950,2.078) (-4.900,2.091) (-4.850,2.102) (-4.800,2.113) (-4.750,2.122) (-4.700,2.130) (-4.650,2.137) (-4.600,2.143) (-4.550,2.148) (-4.500,2.152) (-4.450,2.155) (-4.400,2.157) (-4.350,2.158) (-4.300,2.158) (-4.250,2.158) (-4.200,2.156) (-4.150,2.153) (-4.100,2.150) (-4.050,2.146) (-4.000,2.140) (-3.950,2.134) (-3.900,2.128) (-3.850,2.120) (-3.800,2.111) (-3.750,2.102) (-3.700,2.092) (-3.650,2.081) (-3.600,2.070) (-3.550,2.057) (-3.500,2.044) (-3.450,2.030) (-3.400,2.016) (-3.350,2.001) (-3.300,1.985) (-3.250,1.968) (-3.200,1.951) (-3.150,1.933) (-3.100,1.914) (-3.050,1.895) (-3.001,1.875) (-2.950,1.854) (-2.900,1.833) (-2.850,1.811) (-2.800,1.789) (-2.750,1.766) (-2.700,1.742) (-2.650,1.718) (-2.600,1.694) (-2.550,1.669) (-2.500,1.643) (-2.450,1.617) (-2.400,1.590) (-2.350,1.563) (-2.300,1.535) (-2.250,1.507) (-2.200,1.479) (-2.150,1.450) (-2.100,1.421) (-2.050,1.391) (-2.000,1.361) (-1.950,1.331) (-1.900,1.300) (-1.850,1.269) (-1.800,1.238) (-1.750,1.206) (-1.702,1.175) (-1.650,1.142) (-1.600,1.109) (-1.550,1.076) (-1.500,1.044) (-1.450,1.010) (-1.400,0.977) (-1.350,0.943) (-1.300,0.910) (-1.250,0.876) (-1.200,0.842) (-1.150,0.807) (-1.103,0.775) (-1.050,0.739) (-1.000,0.704) (-0.950,0.669) (-0.900,0.634) (-0.851,0.600) (-0.800,0.565) (-0.750,0.529) (-0.709,0.500) (-0.650,0.459) (-0.602,0.425) (-0.550,0.389) (-0.500,0.354) (-0.450,0.319) (-0.400,0.283) (-0.354,0.250) (-0.300,0.213) (-0.250,0.177) (-0.220,0.155)}
\EQ{(0.206,0.144) (0.248,0.175) (0.300,0.213) (0.350,0.248) (0.400,0.283) (0.450,0.319) (0.496,0.350) (0.550,0.389) (0.600,0.424) (0.650,0.459) (0.700,0.494) (0.744,0.525) (0.800,0.565) (0.850,0.599) (0.900,0.634) (0.950,0.669) (0.995,0.700) (1.050,0.739) (1.100,0.773) (1.150,0.807) (1.200,0.842) (1.249,0.875) (1.300,0.910) (1.350,0.943) (1.397,0.975) (1.450,1.010) (1.500,1.044) (1.548,1.075) (1.600,1.109) (1.650,1.142) (1.700,1.174) (1.750,1.206) (1.800,1.238) (1.850,1.269) (1.900,1.300) (1.950,1.331) (2.000,1.361) (2.050,1.391) (2.100,1.421) (2.150,1.450) (2.193,1.475) (2.237,1.500) (2.282,1.525) (2.327,1.550) (2.372,1.575) (2.419,1.600) (2.466,1.625) (2.514,1.650) (2.563,1.675) (2.613,1.700) (2.664,1.725) (2.716,1.750) (2.770,1.775) (2.825,1.800) (2.875,1.822) (2.925,1.844) (2.975,1.865) (3.025,1.885) (3.075,1.904) (3.125,1.923) (3.175,1.942) (3.225,1.959) (3.275,1.976) (3.325,1.993) (3.375,2.008) (3.425,2.023) (3.475,2.037) (3.525,2.051) (3.575,2.064) (3.625,2.076) (3.675,2.087) (3.725,2.097) (3.775,2.107) (3.825,2.116) (3.875,2.124) (3.925,2.131) (3.975,2.138) (4.025,2.143) (4.075,2.148) (4.125,2.152) (4.175,2.155) (4.225,2.157) (4.275,2.158) (4.325,2.158) (4.375,2.158) (4.425,2.156) (4.475,2.154) (4.525,2.150) (4.575,2.146) (4.625,2.140) (4.675,2.134) (4.725,2.126) (4.775,2.117) (4.825,2.108) (4.875,2.097) (4.925,2.085) (4.975,2.071) (5.025,2.057) (5.075,2.041) (5.125,2.024) (5.175,2.006) (5.225,1.986) (5.275,1.965) (5.325,1.942) (5.375,1.918) (5.425,1.892) (5.475,1.864) (5.525,1.835) (5.575,1.804) (5.618,1.775) (5.675,1.735) (5.722,1.700) (5.775,1.658) (5.814,1.625) (5.869,1.575) (5.921,1.525) (5.970,1.475) (6.016,1.425) (6.050,1.385) (6.098,1.325) (6.136,1.275) (6.172,1.225) (6.200,1.183) (6.236,1.125) (6.266,1.075) (6.294,1.025) (6.320,0.975) (6.344,0.925) (6.367,0.875) (6.388,0.825) (6.408,0.775) (6.426,0.725) (6.443,0.675) (6.459,0.625) (6.473,0.575) (6.486,0.525) (6.498,0.475) (6.509,0.425) (6.518,0.375) (6.526,0.325) (6.533,0.275) (6.539,0.225) (6.543,0.175) (6.547,0.125) (6.549,0.075) (6.550,0.025) (6.550,-0.037) (6.548,-0.100) (6.545,-0.150) (6.541,-0.200) (6.536,-0.250) (6.530,-0.300) (6.522,-0.350) (6.513,-0.400) (6.503,-0.450) (6.492,-0.500) (6.480,-0.550) (6.466,-0.600) (6.451,-0.650) (6.435,-0.700) (6.417,-0.750) (6.398,-0.800) (6.378,-0.850) (6.356,-0.900) (6.332,-0.950) (6.307,-1.000) (6.280,-1.050) (6.251,-1.100) (6.225,-1.143) (6.189,-1.200) (6.154,-1.250) (6.118,-1.300) (6.079,-1.350) (6.037,-1.400) (6.000,-1.442) (5.950,-1.496) (5.900,-1.546) (5.850,-1.593) (5.800,-1.637) (5.754,-1.675) (5.700,-1.717) (5.654,-1.750) (5.600,-1.787) (5.550,-1.819) (5.500,-1.850) (5.456,-1.875) (5.409,-1.900) (5.360,-1.925) (5.307,-1.950) (5.251,-1.975) (5.200,-1.996) (5.150,-2.015) (5.100,-2.033) (5.050,-2.049) (5.000,-2.064) (4.950,-2.078) (4.900,-2.091) (4.850,-2.102) (4.800,-2.113) (4.750,-2.122) (4.700,-2.130) (4.650,-2.137) (4.600,-2.143) (4.550,-2.148) (4.500,-2.152) (4.450,-2.155) (4.400,-2.157) (4.350,-2.158) (4.300,-2.158) (4.250,-2.158) (4.200,-2.156) (4.150,-2.153) (4.100,-2.150) (4.050,-2.146) (4.000,-2.140) (3.950,-2.134) (3.900,-2.128) (3.850,-2.120) (3.800,-2.111) (3.750,-2.102) (3.700,-2.092) (3.650,-2.081) (3.600,-2.070) (3.550,-2.057) (3.500,-2.044) (3.450,-2.030) (3.400,-2.016) (3.350,-2.001) (3.300,-1.985) (3.250,-1.968) (3.200,-1.951) (3.150,-1.933) (3.100,-1.914) (3.050,-1.895) (3.001,-1.875) (2.950,-1.854) (2.900,-1.833) (2.850,-1.811) (2.800,-1.789) (2.750,-1.766) (2.700,-1.742) (2.650,-1.718) (2.600,-1.694) (2.550,-1.669) (2.500,-1.643) (2.450,-1.617) (2.400,-1.590) (2.350,-1.563) (2.300,-1.535) (2.250,-1.507) (2.200,-1.479) (2.150,-1.450) (2.100,-1.421) (2.050,-1.391) (2.000,-1.361) (1.950,-1.331) (1.900,-1.300) (1.850,-1.269) (1.800,-1.238) (1.750,-1.206) (1.702,-1.175) (1.650,-1.142) (1.600,-1.109) (1.550,-1.076) (1.500,-1.044) (1.450,-1.010) (1.400,-0.977) (1.350,-0.943) (1.300,-0.910) (1.250,-0.876) (1.200,-0.842) (1.150,-0.807) (1.103,-0.775) (1.050,-0.739) (1.000,-0.704) (0.950,-0.669) (0.900,-0.634) (0.851,-0.600) (0.800,-0.565) (0.750,-0.529) (0.709,-0.500) (0.650,-0.459) (0.602,-0.425) (0.550,-0.389) (0.500,-0.354) (0.450,-0.319) (0.400,-0.283) (0.354,-0.250) (0.300,-0.213) (0.250,-0.177) (0.206,-0.144)}
\EQ{(-0.220,-0.155) (-0.275,-0.195) (-0.319,-0.225) (-0.375,-0.266) (-0.424,-0.300) (-0.475,-0.336) (-0.525,-0.371) (-0.567,-0.400) (-0.625,-0.442) (-0.673,-0.475) (-0.725,-0.512) (-0.775,-0.547) (-0.825,-0.583) (-0.875,-0.618) (-0.922,-0.650) (-0.975,-0.688) (-1.025,-0.722) (-1.075,-0.757) (-1.125,-0.792) (-1.173,-0.825) (-1.225,-0.861) (-1.275,-0.896) (-1.318,-0.925) (-1.375,-0.964) (-1.425,-0.998) (-1.475,-1.032) (-1.525,-1.066) (-1.575,-1.100) (-1.625,-1.133) (-1.675,-1.167) (-1.725,-1.200) (-1.775,-1.233) (-1.825,-1.266) (-1.875,-1.298) (-1.925,-1.330) (-1.975,-1.362) (-2.025,-1.394) (-2.074,-1.425) (-2.125,-1.457) (-2.175,-1.488) (-2.225,-1.519) (-2.275,-1.549) (-2.325,-1.580) (-2.375,-1.609) (-2.425,-1.639) (-2.475,-1.668) (-2.525,-1.697) (-2.575,-1.725) (-2.619,-1.750) (-2.665,-1.775) (-2.710,-1.800) (-2.757,-1.825) (-2.803,-1.850) (-2.850,-1.874) (-2.900,-1.900) (-2.949,-1.925) (-2.999,-1.950) (-3.050,-1.975) (-3.100,-1.999) (-3.150,-2.022) (-3.200,-2.045) (-3.250,-2.068) (-3.300,-2.090) (-3.350,-2.111) (-3.400,-2.132) (-3.450,-2.152) (-3.500,-2.172) (-3.550,-2.191) (-3.600,-2.210) (-3.650,-2.228) (-3.700,-2.246) (-3.750,-2.263) (-3.800,-2.279) (-3.850,-2.295) (-3.900,-2.310) (-3.950,-2.324) (-4.000,-2.338) (-4.050,-2.352) (-4.100,-2.364) (-4.150,-2.376) (-4.200,-2.387) (-4.250,-2.398) (-4.300,-2.408) (-4.350,-2.417) (-4.400,-2.425) (-4.450,-2.433) (-4.500,-2.440) (-4.550,-2.447) (-4.600,-2.452) (-4.650,-2.457) (-4.700,-2.462) (-4.750,-2.465) (-4.800,-2.468) (-4.850,-2.469) (-4.900,-2.470) (-4.950,-2.471) (-5.000,-2.470) (-5.050,-2.469) (-5.100,-2.467) (-5.150,-2.464) (-5.200,-2.460) (-5.250,-2.455) (-5.300,-2.449) (-5.350,-2.443) (-5.400,-2.435) (-5.450,-2.427) (-5.500,-2.417) (-5.550,-2.407) (-5.600,-2.395) (-5.650,-2.383) (-5.700,-2.370) (-5.750,-2.355) (-5.800,-2.339) (-5.850,-2.323) (-5.900,-2.305) (-5.950,-2.286) (-6.000,-2.265) (-6.050,-2.244) (-6.100,-2.221) (-6.150,-2.196) (-6.200,-2.171) (-6.250,-2.144) (-6.300,-2.115) (-6.350,-2.085) (-6.400,-2.053) (-6.450,-2.019) (-6.500,-1.983) (-6.545,-1.950) (-6.600,-1.906) (-6.650,-1.864) (-6.695,-1.825) (-6.748,-1.775) (-6.799,-1.725) (-6.846,-1.675) (-6.892,-1.625) (-6.925,-1.586) (-6.974,-1.525) (-7.013,-1.475) (-7.049,-1.425) (-7.083,-1.375) (-7.116,-1.325) (-7.147,-1.275) (-7.175,-1.226) (-7.200,-1.181) (-7.225,-1.134) (-7.250,-1.084) (-7.275,-1.031) (-7.300,-0.975) (-7.320,-0.925) (-7.340,-0.875) (-7.358,-0.825) (-7.375,-0.775) (-7.391,-0.725) (-7.405,-0.675) (-7.419,-0.625) (-7.431,-0.575) (-7.443,-0.525) (-7.453,-0.475) (-7.462,-0.425) (-7.470,-0.375) (-7.477,-0.325) (-7.483,-0.275) (-7.488,-0.225) (-7.492,-0.175) (-7.495,-0.125) (-7.497,-0.075) (-7.498,-0.025) (-7.498,0.050) (-7.496,0.100) (-7.494,0.150) (-7.490,0.200) (-7.486,0.250) (-7.480,0.300) (-7.474,0.350) (-7.466,0.400) (-7.457,0.450) (-7.448,0.500) (-7.437,0.550) (-7.425,0.600) (-7.412,0.650) (-7.398,0.700) (-7.383,0.750) (-7.367,0.800) (-7.349,0.850) (-7.330,0.900) (-7.310,0.950) (-7.289,1.000) (-7.266,1.050) (-7.242,1.100) (-7.217,1.150) (-7.190,1.200) (-7.161,1.250) (-7.131,1.300) (-7.100,1.350) (-7.066,1.400) (-7.031,1.450) (-7.000,1.492) (-6.955,1.550) (-6.913,1.600) (-6.875,1.644) (-6.825,1.698) (-6.775,1.749) (-6.725,1.797) (-6.675,1.843) (-6.625,1.886) (-6.577,1.925) (-6.525,1.965) (-6.477,2.000) (-6.425,2.036) (-6.375,2.069) (-6.325,2.100) (-6.275,2.129) (-6.225,2.157) (-6.175,2.184) (-6.125,2.209) (-6.075,2.232) (-6.025,2.255) (-5.977,2.275) (-5.925,2.295) (-5.875,2.314) (-5.825,2.331) (-5.775,2.347) (-5.725,2.362) (-5.675,2.376) (-5.625,2.389) (-5.575,2.401) (-5.525,2.412) (-5.475,2.422) (-5.425,2.431) (-5.375,2.439) (-5.325,2.446) (-5.275,2.452) (-5.225,2.457) (-5.175,2.462) (-5.125,2.465) (-5.075,2.468) (-5.025,2.470) (-4.975,2.471) (-4.925,2.471) (-4.875,2.470) (-4.825,2.469) (-4.775,2.466) (-4.725,2.463) (-4.675,2.460) (-4.625,2.455) (-4.575,2.450) (-4.525,2.444) (-4.475,2.437) (-4.425,2.430) (-4.375,2.421) (-4.325,2.412) (-4.275,2.403) (-4.225,2.393) (-4.175,2.382) (-4.125,2.370) (-4.075,2.358) (-4.025,2.345) (-3.975,2.331) (-3.925,2.317) (-3.875,2.302) (-3.825,2.287) (-3.775,2.271) (-3.725,2.254) (-3.675,2.237) (-3.625,2.219) (-3.575,2.201) (-3.525,2.182) (-3.475,2.162) (-3.425,2.142) (-3.375,2.122) (-3.325,2.100) (-3.275,2.079) (-3.225,2.057) (-3.175,2.034) (-3.125,2.011) (-3.075,1.987) (-3.025,1.963) (-2.975,1.938) (-2.925,1.913) (-2.875,1.887) (-2.825,1.861) (-2.775,1.835) (-2.725,1.808) (-2.675,1.781) (-2.625,1.753) (-2.575,1.725) (-2.531,1.700) (-2.487,1.675) (-2.425,1.639) (-2.375,1.609) (-2.325,1.580) (-2.276,1.550) (-2.225,1.519) (-2.175,1.488) (-2.125,1.457) (-2.075,1.426) (-2.025,1.394) (-1.975,1.362) (-1.925,1.330) (-1.878,1.300) (-1.825,1.266) (-1.775,1.233) (-1.725,1.200) (-1.675,1.167) (-1.625,1.133) (-1.576,1.100) (-1.525,1.066) (-1.475,1.032) (-1.428,1.000) (-1.375,0.964) (-1.325,0.930) (-1.282,0.900) (-1.225,0.861) (-1.175,0.827) (-1.125,0.792) (-1.075,0.757) (-1.029,0.725) (-0.975,0.688) (-0.925,0.652) (-0.875,0.618) (-0.825,0.583) (-0.779,0.550) (-0.725,0.512) (-0.675,0.477) (-0.625,0.442) (-0.575,0.407) (-0.531,0.375) (-0.475,0.336) (-0.425,0.300) (-0.375,0.266) (-0.325,0.230) (-0.284,0.200) (-0.225,0.160) (-0.220,0.155)}
\EQ{(0.206,0.144) (0.248,0.175) (0.300,0.213) (0.350,0.248) (0.400,0.283) (0.450,0.319) (0.496,0.350) (0.550,0.389) (0.600,0.424) (0.650,0.460) (0.700,0.495) (0.744,0.525) (0.800,0.565) (0.850,0.600) (0.900,0.635) (0.950,0.670) (0.993,0.700) (1.050,0.740) (1.100,0.774) (1.150,0.809) (1.200,0.844) (1.245,0.875) (1.300,0.913) (1.350,0.947) (1.400,0.981) (1.450,1.015) (1.500,1.049) (1.550,1.083) (1.600,1.117) (1.650,1.150) (1.700,1.183) (1.750,1.216) (1.800,1.249) (1.850,1.282) (1.900,1.314) (1.950,1.346) (1.995,1.375) (2.050,1.410) (2.100,1.442) (2.150,1.473) (2.194,1.500) (2.250,1.534) (2.300,1.565) (2.350,1.594) (2.400,1.624) (2.444,1.650) (2.500,1.682) (2.550,1.711) (2.600,1.739) (2.650,1.767) (2.700,1.795) (2.750,1.822) (2.800,1.848) (2.850,1.874) (2.900,1.900) (2.949,1.925) (2.999,1.950) (3.050,1.975) (3.100,1.999) (3.150,2.022) (3.200,2.045) (3.250,2.068) (3.300,2.090) (3.350,2.111) (3.400,2.132) (3.450,2.152) (3.500,2.172) (3.550,2.191) (3.600,2.210) (3.650,2.228) (3.700,2.246) (3.750,2.263) (3.800,2.279) (3.850,2.295) (3.900,2.310) (3.950,2.324) (4.000,2.338) (4.050,2.352) (4.100,2.364) (4.150,2.376) (4.200,2.387) (4.250,2.398) (4.300,2.408) (4.350,2.417) (4.400,2.425) (4.450,2.433) (4.500,2.440) (4.550,2.447) (4.600,2.452) (4.650,2.457) (4.700,2.462) (4.750,2.465) (4.800,2.468) (4.850,2.469) (4.900,2.470) (4.950,2.471) (5.000,2.470) (5.050,2.469) (5.100,2.467) (5.150,2.464) (5.200,2.460) (5.250,2.455) (5.300,2.449) (5.350,2.443) (5.400,2.435) (5.450,2.427) (5.500,2.417) (5.550,2.407) (5.600,2.395) (5.650,2.383) (5.700,2.370) (5.750,2.355) (5.800,2.339) (5.850,2.323) (5.900,2.305) (5.950,2.286) (6.000,2.265) (6.050,2.244) (6.100,2.221) (6.150,2.196) (6.200,2.171) (6.250,2.144) (6.300,2.115) (6.350,2.085) (6.400,2.053) (6.450,2.019) (6.500,1.983) (6.545,1.950) (6.600,1.906) (6.650,1.864) (6.695,1.825) (6.748,1.775) (6.799,1.725) (6.846,1.675) (6.892,1.625) (6.925,1.586) (6.974,1.525) (7.013,1.475) (7.049,1.425) (7.083,1.375) (7.116,1.325) (7.147,1.275) (7.175,1.226) (7.200,1.181) (7.225,1.134) (7.250,1.084) (7.275,1.031) (7.300,0.975) (7.320,0.925) (7.340,0.875) (7.358,0.825) (7.375,0.775) (7.391,0.725) (7.405,0.675) (7.419,0.625) (7.431,0.575) (7.443,0.525) (7.453,0.475) (7.462,0.425) (7.470,0.375) (7.477,0.325) (7.483,0.275) (7.488,0.225) (7.492,0.175) (7.495,0.125) (7.497,0.075) (7.498,0.025) (7.498,-0.050) (7.496,-0.100) (7.494,-0.150) (7.490,-0.200) (7.486,-0.250) (7.480,-0.300) (7.474,-0.350) (7.466,-0.400) (7.457,-0.450) (7.448,-0.500) (7.437,-0.550) (7.425,-0.600) (7.412,-0.650) (7.398,-0.700) (7.383,-0.750) (7.367,-0.800) (7.349,-0.850) (7.330,-0.900) (7.310,-0.950) (7.289,-1.000) (7.266,-1.050) (7.242,-1.100) (7.217,-1.150) (7.190,-1.200) (7.161,-1.250) (7.131,-1.300) (7.100,-1.350) (7.066,-1.400) (7.031,-1.450) (7.000,-1.492) (6.955,-1.550) (6.913,-1.600) (6.875,-1.644) (6.825,-1.698) (6.775,-1.749) (6.725,-1.797) (6.675,-1.843) (6.625,-1.886) (6.577,-1.925) (6.525,-1.965) (6.477,-2.000) (6.425,-2.036) (6.375,-2.069) (6.325,-2.100) (6.275,-2.129) (6.225,-2.157) (6.175,-2.184) (6.125,-2.209) (6.075,-2.232) (6.025,-2.255) (5.977,-2.275) (5.925,-2.295) (5.875,-2.314) (5.825,-2.331) (5.775,-2.347) (5.725,-2.362) (5.675,-2.376) (5.625,-2.389) (5.575,-2.401) (5.525,-2.412) (5.475,-2.422) (5.425,-2.431) (5.375,-2.439) (5.325,-2.446) (5.275,-2.452) (5.225,-2.457) (5.175,-2.462) (5.125,-2.465) (5.075,-2.468) (5.025,-2.470) (4.975,-2.471) (4.925,-2.471) (4.875,-2.470) (4.825,-2.469) (4.775,-2.466) (4.725,-2.463) (4.675,-2.460) (4.625,-2.455) (4.575,-2.450) (4.525,-2.444) (4.475,-2.437) (4.425,-2.430) (4.375,-2.421) (4.325,-2.412) (4.275,-2.403) (4.225,-2.393) (4.175,-2.382) (4.125,-2.370) (4.075,-2.358) (4.025,-2.345) (3.975,-2.331) (3.925,-2.317) (3.875,-2.302) (3.825,-2.287) (3.775,-2.271) (3.725,-2.254) (3.675,-2.237) (3.625,-2.219) (3.575,-2.201) (3.525,-2.182) (3.475,-2.162) (3.425,-2.142) (3.375,-2.122) (3.325,-2.100) (3.275,-2.079) (3.225,-2.057) (3.175,-2.034) (3.125,-2.011) (3.075,-1.987) (3.025,-1.963) (2.975,-1.938) (2.925,-1.913) (2.875,-1.887) (2.825,-1.861) (2.775,-1.835) (2.725,-1.808) (2.675,-1.781) (2.625,-1.753) (2.575,-1.725) (2.531,-1.700) (2.487,-1.675) (2.425,-1.639) (2.375,-1.609) (2.325,-1.580) (2.276,-1.550) (2.225,-1.519) (2.175,-1.488) (2.125,-1.457) (2.075,-1.426) (2.025,-1.394) (1.975,-1.362) (1.925,-1.330) (1.878,-1.300) (1.825,-1.266) (1.775,-1.233) (1.725,-1.200) (1.675,-1.167) (1.625,-1.133) (1.576,-1.100) (1.525,-1.066) (1.475,-1.032) (1.428,-1.000) (1.375,-0.964) (1.325,-0.930) (1.282,-0.900) (1.225,-0.861) (1.175,-0.827) (1.125,-0.792) (1.075,-0.757) (1.029,-0.725) (0.975,-0.688) (0.925,-0.652) (0.875,-0.618) (0.825,-0.583) (0.779,-0.550) (0.725,-0.512) (0.675,-0.477) (0.625,-0.442) (0.575,-0.407) (0.531,-0.375) (0.475,-0.336) (0.425,-0.300) (0.375,-0.266) (0.325,-0.230) (0.284,-0.200) (0.225,-0.160) (0.206,-0.144)}
\EQ{(-0.220,-0.155) (-0.275,-0.195) (-0.319,-0.225) (-0.375,-0.266) (-0.424,-0.300) (-0.475,-0.336) (-0.525,-0.371) (-0.567,-0.400) (-0.625,-0.442) (-0.672,-0.475) (-0.725,-0.513) (-0.775,-0.548) (-0.825,-0.583) (-0.875,-0.618) (-0.921,-0.650) (-0.975,-0.689) (-1.025,-0.724) (-1.075,-0.759) (-1.125,-0.794) (-1.170,-0.825) (-1.225,-0.864) (-1.275,-0.899) (-1.325,-0.934) (-1.375,-0.968) (-1.421,-1.000) (-1.475,-1.038) (-1.525,-1.072) (-1.575,-1.107) (-1.625,-1.141) (-1.674,-1.175) (-1.725,-1.210) (-1.775,-1.244) (-1.821,-1.275) (-1.875,-1.312) (-1.925,-1.346) (-1.969,-1.375) (-2.025,-1.413) (-2.075,-1.446) (-2.118,-1.475) (-2.175,-1.513) (-2.225,-1.546) (-2.270,-1.575) (-2.325,-1.611) (-2.375,-1.644) (-2.424,-1.675) (-2.475,-1.708) (-2.525,-1.740) (-2.575,-1.772) (-2.620,-1.800) (-2.675,-1.834) (-2.725,-1.866) (-2.775,-1.896) (-2.822,-1.925) (-2.875,-1.957) (-2.925,-1.988) (-2.975,-2.017) (-3.025,-2.047) (-3.073,-2.075) (-3.125,-2.105) (-3.175,-2.134) (-3.225,-2.163) (-3.275,-2.191) (-3.325,-2.219) (-3.375,-2.246) (-3.425,-2.274) (-3.474,-2.300) (-3.521,-2.325) (-3.569,-2.350) (-3.617,-2.375) (-3.665,-2.400) (-3.715,-2.425) (-3.765,-2.450) (-3.816,-2.475) (-3.867,-2.500) (-3.920,-2.525) (-3.973,-2.550) (-4.025,-2.574) (-4.075,-2.596) (-4.125,-2.619) (-4.175,-2.640) (-4.225,-2.662) (-4.275,-2.683) (-4.325,-2.703) (-4.375,-2.723) (-4.425,-2.743) (-4.475,-2.762) (-4.525,-2.781) (-4.575,-2.799) (-4.625,-2.817) (-4.675,-2.834) (-4.725,-2.851) (-4.775,-2.867) (-4.825,-2.883) (-4.875,-2.899) (-4.925,-2.914) (-4.975,-2.928) (-5.025,-2.942) (-5.075,-2.956) (-5.125,-2.968) (-5.175,-2.981) (-5.225,-2.993) (-5.275,-3.004) (-5.325,-3.015) (-5.375,-3.025) (-5.425,-3.035) (-5.475,-3.044) (-5.525,-3.053) (-5.575,-3.061) (-5.625,-3.069) (-5.675,-3.075) (-5.725,-3.082) (-5.775,-3.088) (-5.825,-3.093) (-5.875,-3.098) (-5.925,-3.102) (-5.975,-3.105) (-6.025,-3.108) (-6.075,-3.110) (-6.125,-3.112) (-6.175,-3.113) (-6.225,-3.113) (-6.275,-3.113) (-6.325,-3.112) (-6.375,-3.110) (-6.425,-3.108) (-6.475,-3.105) (-6.525,-3.101) (-6.575,-3.097) (-6.625,-3.092) (-6.675,-3.086) (-6.725,-3.080) (-6.775,-3.072) (-6.825,-3.065) (-6.875,-3.056) (-6.925,-3.046) (-6.975,-3.036) (-7.025,-3.025) (-7.075,-3.013) (-7.125,-3.001) (-7.175,-2.987) (-7.225,-2.973) (-7.275,-2.958) (-7.325,-2.942) (-7.375,-2.925) (-7.425,-2.907) (-7.475,-2.888) (-7.525,-2.868) (-7.575,-2.848) (-7.625,-2.826) (-7.675,-2.803) (-7.725,-2.779) (-7.775,-2.754) (-7.825,-2.728) (-7.875,-2.701) (-7.920,-2.675) (-7.975,-2.642) (-8.025,-2.611) (-8.075,-2.579) (-8.117,-2.550) (-8.175,-2.509) (-8.222,-2.475) (-8.275,-2.434) (-8.318,-2.400) (-8.375,-2.352) (-8.425,-2.308) (-8.475,-2.262) (-8.514,-2.225) (-8.550,-2.189) (-8.589,-2.150) (-8.625,-2.111) (-8.675,-2.056) (-8.723,-2.000) (-8.765,-1.950) (-8.800,-1.905) (-8.842,-1.850) (-8.875,-1.804) (-8.912,-1.750) (-8.946,-1.700) (-8.975,-1.654) (-9.008,-1.600) (-9.037,-1.550) (-9.065,-1.500) (-9.092,-1.450) (-9.118,-1.400) (-9.142,-1.350) (-9.165,-1.300) (-9.188,-1.250) (-9.209,-1.200) (-9.229,-1.150) (-9.248,-1.100) (-9.267,-1.050) (-9.284,-1.000) (-9.300,-0.951) (-9.316,-0.900) (-9.330,-0.850) (-9.344,-0.800) (-9.357,-0.750) (-9.368,-0.700) (-9.379,-0.650) (-9.390,-0.600) (-9.399,-0.550) (-9.407,-0.500) (-9.415,-0.450) (-9.422,-0.400) (-9.428,-0.350) (-9.433,-0.300) (-9.437,-0.250) (-9.441,-0.200) (-9.444,-0.150) (-9.446,-0.100) (-9.447,-0.050) (-9.447,0.000) (-9.447,0.050) (-9.446,0.100) (-9.444,0.150) (-9.441,0.200) (-9.437,0.250) (-9.433,0.300) (-9.428,0.350) (-9.422,0.400) (-9.415,0.450) (-9.407,0.500) (-9.399,0.550) (-9.390,0.600) (-9.379,0.650) (-9.368,0.700) (-9.357,0.750) (-9.344,0.800) (-9.330,0.850) (-9.316,0.900) (-9.300,0.950) (-9.284,1.000) (-9.267,1.050) (-9.248,1.100) (-9.229,1.150) (-9.209,1.200) (-9.188,1.250) (-9.165,1.300) (-9.142,1.350) (-9.118,1.400) (-9.092,1.450) (-9.065,1.500) (-9.037,1.550) (-9.008,1.600) (-8.978,1.650) (-8.950,1.693) (-8.912,1.750) (-8.878,1.800) (-8.842,1.850) (-8.804,1.900) (-8.765,1.950) (-8.725,1.998) (-8.680,2.050) (-8.636,2.100) (-8.600,2.138) (-8.564,2.175) (-8.525,2.214) (-8.488,2.250) (-8.450,2.285) (-8.406,2.325) (-8.350,2.373) (-8.300,2.414) (-8.255,2.450) (-8.200,2.491) (-8.153,2.525) (-8.100,2.562) (-8.050,2.595) (-8.003,2.625) (-7.950,2.657) (-7.900,2.686) (-7.850,2.714) (-7.800,2.741) (-7.750,2.767) (-7.700,2.791) (-7.650,2.814) (-7.600,2.837) (-7.550,2.858) (-7.500,2.878) (-7.450,2.898) (-7.400,2.916) (-7.350,2.933) (-7.300,2.950) (-7.250,2.966) (-7.200,2.980) (-7.150,2.994) (-7.100,3.007) (-7.050,3.019) (-7.000,3.031) (-6.950,3.041) (-6.900,3.051) (-6.850,3.060) (-6.800,3.069) (-6.750,3.076) (-6.700,3.083) (-6.650,3.089) (-6.600,3.094) (-6.550,3.099) (-6.500,3.103) (-6.450,3.106) (-6.400,3.109) (-6.350,3.111) (-6.300,3.112) (-6.250,3.113) (-6.200,3.113) (-6.150,3.112) (-6.100,3.111) (-6.050,3.109) (-6.000,3.107) (-5.950,3.103) (-5.900,3.100) (-5.850,3.095) (-5.800,3.090) (-5.750,3.085) (-5.700,3.079) (-5.650,3.072) (-5.600,3.065) (-5.550,3.057) (-5.500,3.049) (-5.450,3.040) (-5.400,3.030) (-5.350,3.020) (-5.300,3.010) (-5.250,2.998) (-5.200,2.987) (-5.151,2.975) (-5.100,2.962) (-5.050,2.949) (-5.000,2.935) (-4.950,2.921) (-4.900,2.906) (-4.850,2.891) (-4.800,2.875) (-4.750,2.859) (-4.700,2.843) (-4.650,2.825) (-4.600,2.808) (-4.550,2.790) (-4.500,2.771) (-4.450,2.752) (-4.400,2.733) (-4.350,2.713) (-4.300,2.693) (-4.250,2.672) (-4.200,2.651) (-4.150,2.629) (-4.100,2.608) (-4.050,2.585) (-4.000,2.562) (-3.950,2.539) (-3.900,2.516) (-3.850,2.492) (-3.800,2.467) (-3.750,2.443) (-3.700,2.418) (-3.650,2.392) (-3.600,2.366) (-3.550,2.340) (-3.500,2.314) (-3.450,2.287) (-3.400,2.260) (-3.350,2.233) (-3.300,2.205) (-3.250,2.177) (-3.203,2.150) (-3.159,2.125) (-3.116,2.100) (-3.050,2.062) (-3.000,2.032) (-2.950,2.002) (-2.904,1.975) (-2.850,1.942) (-2.800,1.912) (-2.750,1.881) (-2.700,1.850) (-2.650,1.819) (-2.600,1.787) (-2.550,1.756) (-2.502,1.725) (-2.450,1.692) (-2.400,1.660) (-2.350,1.627) (-2.300,1.595) (-2.250,1.562) (-2.200,1.529) (-2.156,1.500) (-2.100,1.463) (-2.050,1.430) (-2.006,1.400) (-1.950,1.363) (-1.900,1.329) (-1.858,1.300) (-1.800,1.261) (-1.750,1.227) (-1.700,1.193) (-1.650,1.158) (-1.602,1.125) (-1.550,1.090) (-1.500,1.055) (-1.457,1.025) (-1.400,0.986) (-1.350,0.951) (-1.300,0.916) (-1.250,0.881) (-1.206,0.850) (-1.150,0.811) (-1.100,0.776) (-1.050,0.741) (-1.000,0.706) (-0.956,0.675) (-0.900,0.636) (-0.850,0.600) (-0.800,0.566) (-0.750,0.530) (-0.708,0.500) (-0.650,0.460) (-0.601,0.425) (-0.550,0.389) (-0.500,0.354) (-0.450,0.319) (-0.400,0.283) (-0.354,0.250) (-0.300,0.213) (-0.250,0.177) (-0.220,0.155)}
\EQ{(0.206,0.144) (0.248,0.175) (0.300,0.213) (0.350,0.248) (0.400,0.283) (0.450,0.319) (0.496,0.350) (0.550,0.389) (0.600,0.424) (0.650,0.460) (0.700,0.495) (0.743,0.525) (0.800,0.566) (0.849,0.600) (0.900,0.636) (0.950,0.671) (0.992,0.700) (1.050,0.741) (1.098,0.775) (1.150,0.811) (1.200,0.846) (1.242,0.875) (1.300,0.916) (1.349,0.950) (1.400,0.986) (1.450,1.020) (1.493,1.050) (1.550,1.090) (1.600,1.124) (1.650,1.158) (1.700,1.193) (1.748,1.225) (1.800,1.261) (1.850,1.295) (1.895,1.325) (1.950,1.363) (2.000,1.396) (2.043,1.425) (2.100,1.463) (2.150,1.496) (2.194,1.525) (2.250,1.562) (2.300,1.595) (2.347,1.625) (2.400,1.660) (2.450,1.692) (2.500,1.724) (2.550,1.756) (2.600,1.787) (2.650,1.819) (2.700,1.850) (2.750,1.881) (2.800,1.912) (2.850,1.942) (2.900,1.972) (2.946,2.000) (3.000,2.032) (3.050,2.062) (3.100,2.091) (3.150,2.120) (3.200,2.148) (3.247,2.175) (3.292,2.200) (3.337,2.225) (3.382,2.250) (3.428,2.275) (3.474,2.300) (3.521,2.325) (3.569,2.350) (3.617,2.375) (3.665,2.400) (3.715,2.425) (3.765,2.450) (3.816,2.475) (3.867,2.500) (3.920,2.525) (3.973,2.550) (4.025,2.574) (4.075,2.596) (4.125,2.619) (4.175,2.640) (4.225,2.662) (4.275,2.683) (4.325,2.703) (4.375,2.723) (4.425,2.743) (4.475,2.762) (4.525,2.781) (4.575,2.799) (4.625,2.817) (4.675,2.834) (4.725,2.851) (4.775,2.867) (4.825,2.883) (4.875,2.899) (4.925,2.914) (4.975,2.928) (5.025,2.942) (5.075,2.956) (5.125,2.968) (5.175,2.981) (5.225,2.993) (5.275,3.004) (5.325,3.015) (5.375,3.025) (5.425,3.035) (5.475,3.044) (5.525,3.053) (5.575,3.061) (5.625,3.069) (5.675,3.075) (5.725,3.082) (5.775,3.088) (5.825,3.093) (5.875,3.098) (5.925,3.102) (5.975,3.105) (6.025,3.108) (6.075,3.110) (6.125,3.112) (6.175,3.113) (6.225,3.113) (6.275,3.113) (6.325,3.112) (6.375,3.110) (6.425,3.108) (6.475,3.105) (6.525,3.101) (6.575,3.097) (6.625,3.092) (6.675,3.086) (6.725,3.080) (6.775,3.072) (6.825,3.065) (6.875,3.056) (6.925,3.046) (6.975,3.036) (7.025,3.025) (7.075,3.013) (7.125,3.001) (7.175,2.987) (7.225,2.973) (7.275,2.958) (7.325,2.942) (7.375,2.925) (7.425,2.907) (7.475,2.888) (7.525,2.868) (7.575,2.848) (7.625,2.826) (7.675,2.803) (7.725,2.779) (7.775,2.754) (7.825,2.728) (7.875,2.701) (7.920,2.675) (7.975,2.642) (8.025,2.611) (8.075,2.579) (8.117,2.550) (8.175,2.509) (8.222,2.475) (8.275,2.434) (8.318,2.400) (8.375,2.352) (8.425,2.308) (8.475,2.262) (8.514,2.225) (8.550,2.189) (8.589,2.150) (8.625,2.111) (8.675,2.056) (8.723,2.000) (8.765,1.950) (8.800,1.905) (8.842,1.850) (8.875,1.804) (8.912,1.750) (8.946,1.700) (8.975,1.654) (9.008,1.600) (9.037,1.550) (9.065,1.500) (9.092,1.450) (9.118,1.400) (9.142,1.350) (9.165,1.300) (9.188,1.250) (9.209,1.200) (9.229,1.150) (9.248,1.100) (9.267,1.050) (9.284,1.000) (9.300,0.951) (9.316,0.900) (9.330,0.850) (9.344,0.800) (9.357,0.750) (9.368,0.700) (9.379,0.650) (9.390,0.600) (9.399,0.550) (9.407,0.500) (9.415,0.450) (9.422,0.400) (9.428,0.350) (9.433,0.300) (9.437,0.250) (9.441,0.200) (9.444,0.150) (9.446,0.100) (9.447,0.050) (9.447,0.000) (9.447,-0.050) (9.446,-0.100) (9.444,-0.150) (9.441,-0.200) (9.437,-0.250) (9.433,-0.300) (9.428,-0.350) (9.422,-0.400) (9.415,-0.450) (9.407,-0.500) (9.399,-0.550) (9.390,-0.600) (9.379,-0.650) (9.368,-0.700) (9.357,-0.750) (9.344,-0.800) (9.330,-0.850) (9.316,-0.900) (9.300,-0.950) (9.284,-1.000) (9.267,-1.050) (9.248,-1.100) (9.229,-1.150) (9.209,-1.200) (9.188,-1.250) (9.165,-1.300) (9.142,-1.350) (9.118,-1.400) (9.092,-1.450) (9.065,-1.500) (9.037,-1.550) (9.008,-1.600) (8.978,-1.650) (8.950,-1.693) (8.912,-1.750) (8.878,-1.800) (8.842,-1.850) (8.804,-1.900) (8.765,-1.950) (8.725,-1.998) (8.680,-2.050) (8.636,-2.100) (8.600,-2.138) (8.564,-2.175) (8.525,-2.214) (8.488,-2.250) (8.450,-2.285) (8.406,-2.325) (8.350,-2.373) (8.300,-2.414) (8.255,-2.450) (8.200,-2.491) (8.153,-2.525) (8.100,-2.562) (8.050,-2.595) (8.003,-2.625) (7.950,-2.657) (7.900,-2.686) (7.850,-2.714) (7.800,-2.741) (7.750,-2.767) (7.700,-2.791) (7.650,-2.814) (7.600,-2.837) (7.550,-2.858) (7.500,-2.878) (7.450,-2.898) (7.400,-2.916) (7.350,-2.933) (7.300,-2.950) (7.250,-2.966) (7.200,-2.980) (7.150,-2.994) (7.100,-3.007) (7.050,-3.019) (7.000,-3.031) (6.950,-3.041) (6.900,-3.051) (6.850,-3.060) (6.800,-3.069) (6.750,-3.076) (6.700,-3.083) (6.650,-3.089) (6.600,-3.094) (6.550,-3.099) (6.500,-3.103) (6.450,-3.106) (6.400,-3.109) (6.350,-3.111) (6.300,-3.112) (6.250,-3.113) (6.200,-3.113) (6.150,-3.112) (6.100,-3.111) (6.050,-3.109) (6.000,-3.107) (5.950,-3.103) (5.900,-3.100) (5.850,-3.095) (5.800,-3.090) (5.750,-3.085) (5.700,-3.079) (5.650,-3.072) (5.600,-3.065) (5.550,-3.057) (5.500,-3.049) (5.450,-3.040) (5.400,-3.030) (5.350,-3.020) (5.300,-3.010) (5.250,-2.998) (5.200,-2.987) (5.151,-2.975) (5.100,-2.962) (5.050,-2.949) (5.000,-2.935) (4.950,-2.921) (4.900,-2.906) (4.850,-2.891) (4.800,-2.875) (4.750,-2.859) (4.700,-2.843) (4.650,-2.825) (4.600,-2.808) (4.550,-2.790) (4.500,-2.771) (4.450,-2.752) (4.400,-2.733) (4.350,-2.713) (4.300,-2.693) (4.250,-2.672) (4.200,-2.651) (4.150,-2.629) (4.100,-2.608) (4.050,-2.585) (4.000,-2.562) (3.950,-2.539) (3.900,-2.516) (3.850,-2.492) (3.800,-2.467) (3.750,-2.443) (3.700,-2.418) (3.650,-2.392) (3.600,-2.366) (3.550,-2.340) (3.500,-2.314) (3.450,-2.287) (3.400,-2.260) (3.350,-2.233) (3.300,-2.205) (3.250,-2.177) (3.203,-2.150) (3.159,-2.125) (3.116,-2.100) (3.050,-2.062) (3.000,-2.032) (2.950,-2.002) (2.904,-1.975) (2.850,-1.942) (2.800,-1.912) (2.750,-1.881) (2.700,-1.850) (2.650,-1.819) (2.600,-1.787) (2.550,-1.756) (2.502,-1.725) (2.450,-1.692) (2.400,-1.660) (2.350,-1.627) (2.300,-1.595) (2.250,-1.562) (2.200,-1.529) (2.156,-1.500) (2.100,-1.463) (2.050,-1.430) (2.006,-1.400) (1.950,-1.363) (1.900,-1.329) (1.858,-1.300) (1.800,-1.261) (1.750,-1.227) (1.700,-1.193) (1.650,-1.158) (1.602,-1.125) (1.550,-1.090) (1.500,-1.055) (1.457,-1.025) (1.400,-0.986) (1.350,-0.951) (1.300,-0.916) (1.250,-0.881) (1.206,-0.850) (1.150,-0.811) (1.100,-0.776) (1.050,-0.741) (1.000,-0.706) (0.956,-0.675) (0.900,-0.636) (0.850,-0.600) (0.800,-0.566) (0.750,-0.530) (0.708,-0.500) (0.650,-0.460) (0.601,-0.425) (0.550,-0.389) (0.500,-0.354) (0.450,-0.319) (0.400,-0.283) (0.354,-0.250) (0.300,-0.213) (0.250,-0.177) (0.206,-0.144)}
\EQ{(10.000,-7.071) (9.950,-7.036) (9.900,-7.000) (9.850,-6.965) (9.800,-6.930) (9.758,-6.900) (9.700,-6.859) (9.652,-6.825) (9.600,-6.788) (9.550,-6.753) (9.500,-6.718) (9.450,-6.682) (9.405,-6.650) (9.350,-6.611) (9.300,-6.576) (9.250,-6.541) (9.200,-6.505) (9.157,-6.475) (9.100,-6.435) (9.051,-6.400) (9.000,-6.364) (8.950,-6.329) (8.900,-6.293) (8.850,-6.258) (8.804,-6.225) (8.750,-6.187) (8.700,-6.152) (8.650,-6.117) (8.600,-6.081) (8.556,-6.050) (8.500,-6.010) (8.450,-5.975) (8.400,-5.940) (8.350,-5.904) (8.309,-5.875) (8.250,-5.834) (8.202,-5.800) (8.150,-5.763) (8.100,-5.728) (8.050,-5.692) (8.000,-5.657) (7.955,-5.625) (7.900,-5.586) (7.850,-5.551) (7.800,-5.515) (7.750,-5.480) (7.708,-5.450) (7.650,-5.409) (7.601,-5.375) (7.550,-5.339) (7.500,-5.303) (7.450,-5.268) (7.400,-5.233) (7.354,-5.200) (7.300,-5.162) (7.250,-5.127) (7.200,-5.091) (7.150,-5.056) (7.106,-5.025) (7.050,-4.985) (7.000,-4.950) (6.950,-4.914) (6.900,-4.879) (6.859,-4.850) (6.800,-4.808) (6.753,-4.775) (6.700,-4.738) (6.650,-4.702) (6.600,-4.667) (6.550,-4.632) (6.505,-4.600) (6.450,-4.561) (6.400,-4.525) (6.350,-4.490) (6.300,-4.455) (6.258,-4.425) (6.200,-4.384) (6.152,-4.350) (6.100,-4.313) (6.050,-4.278) (6.000,-4.243) (5.950,-4.207) (5.904,-4.175) (5.850,-4.137) (5.800,-4.101) (5.750,-4.066) (5.700,-4.031) (5.657,-4.000) (5.600,-3.960) (5.551,-3.925) (5.500,-3.889) (5.450,-3.854) (5.400,-3.818) (5.350,-3.783) (5.303,-3.750) (5.250,-3.712) (5.200,-3.677) (5.150,-3.642) (5.100,-3.606) (5.056,-3.575) (5.000,-3.536) (4.950,-3.500) (4.900,-3.465) (4.850,-3.429) (4.808,-3.400) (4.750,-3.359) (4.702,-3.325) (4.650,-3.288) (4.600,-3.253) (4.550,-3.217) (4.500,-3.182) (4.455,-3.150) (4.400,-3.111) (4.350,-3.076) (4.300,-3.041) (4.250,-3.005) (4.207,-2.975) (4.150,-2.935) (4.101,-2.900) (4.050,-2.864) (4.000,-2.828) (3.950,-2.793) (3.900,-2.758) (3.854,-2.725) (3.800,-2.687) (3.750,-2.652) (3.700,-2.616) (3.650,-2.581) (3.606,-2.550) (3.550,-2.510) (3.500,-2.475) (3.450,-2.440) (3.400,-2.404) (3.359,-2.375) (3.300,-2.334) (3.253,-2.300) (3.200,-2.263) (3.150,-2.227) (3.100,-2.192) (3.050,-2.157) (3.005,-2.125) (2.950,-2.086) (2.900,-2.051) (2.850,-2.015) (2.800,-1.980) (2.758,-1.950) (2.700,-1.909) (2.652,-1.875) (2.600,-1.839) (2.550,-1.803) (2.500,-1.768) (2.450,-1.732) (2.404,-1.700) (2.350,-1.662) (2.300,-1.626) (2.250,-1.591) (2.200,-1.556) (2.157,-1.525) (2.100,-1.485) (2.051,-1.450) (2.000,-1.414) (1.950,-1.379) (1.900,-1.344) (1.850,-1.308) (1.803,-1.275) (1.750,-1.238) (1.700,-1.202) (1.650,-1.167) (1.600,-1.131) (1.556,-1.100) (1.500,-1.061) (1.450,-1.025) (1.400,-0.990) (1.350,-0.955) (1.308,-0.925) (1.250,-0.884) (1.202,-0.850) (1.150,-0.813) (1.100,-0.778) (1.050,-0.743) (1.000,-0.707) (0.955,-0.675) (0.900,-0.637) (0.850,-0.601) (0.800,-0.566) (0.750,-0.531) (0.708,-0.500) (0.650,-0.460) (0.601,-0.425) (0.550,-0.389) (0.500,-0.354) (0.450,-0.319) (0.400,-0.283) (0.354,-0.250) (0.300,-0.213) (0.250,-0.177) (0.206,-0.144)}
\EQ{(-0.220,-0.155) (-0.275,-0.195) (-0.319,-0.225) (-0.375,-0.266) (-0.424,-0.300) (-0.475,-0.336) (-0.525,-0.371) (-0.566,-0.400) (-0.625,-0.442) (-0.672,-0.475) (-0.725,-0.513) (-0.775,-0.548) (-0.825,-0.584) (-0.875,-0.619) (-0.920,-0.650) (-0.975,-0.690) (-1.025,-0.725) (-1.075,-0.760) (-1.125,-0.796) (-1.167,-0.825) (-1.225,-0.866) (-1.273,-0.900) (-1.325,-0.937) (-1.375,-0.972) (-1.425,-1.008) (-1.475,-1.043) (-1.520,-1.075) (-1.575,-1.114) (-1.625,-1.149) (-1.675,-1.185) (-1.725,-1.220) (-1.768,-1.250) (-1.825,-1.291) (-1.874,-1.325) (-1.925,-1.361) (-1.975,-1.397) (-2.025,-1.432) (-2.075,-1.467) (-2.121,-1.500) (-2.175,-1.538) (-2.225,-1.573) (-2.275,-1.609) (-2.325,-1.644) (-2.369,-1.675) (-2.425,-1.715) (-2.475,-1.750) (-2.525,-1.786) (-2.575,-1.821) (-2.616,-1.850) (-2.675,-1.892) (-2.722,-1.925) (-2.775,-1.962) (-2.825,-1.998) (-2.875,-2.033) (-2.925,-2.068) (-2.970,-2.100) (-3.025,-2.139) (-3.075,-2.174) (-3.125,-2.210) (-3.175,-2.245) (-3.217,-2.275) (-3.275,-2.316) (-3.323,-2.350) (-3.375,-2.387) (-3.425,-2.422) (-3.475,-2.457) (-3.525,-2.493) (-3.571,-2.525) (-3.625,-2.563) (-3.675,-2.599) (-3.725,-2.634) (-3.775,-2.669) (-3.818,-2.700) (-3.875,-2.740) (-3.924,-2.775) (-3.975,-2.811) (-4.025,-2.846) (-4.066,-2.875) (-4.125,-2.917) (-4.172,-2.950) (-4.225,-2.988) (-4.275,-3.023) (-4.325,-3.058) (-4.375,-3.094) (-4.419,-3.125) (-4.475,-3.164) (-4.525,-3.200) (-4.575,-3.235) (-4.625,-3.270) (-4.667,-3.300) (-4.725,-3.341) (-4.773,-3.375) (-4.825,-3.412) (-4.875,-3.447) (-4.925,-3.483) (-4.975,-3.518) (-5.021,-3.550) (-5.075,-3.589) (-5.125,-3.624) (-5.175,-3.659) (-5.225,-3.695) (-5.268,-3.725) (-5.325,-3.765) (-5.374,-3.800) (-5.425,-3.836) (-5.475,-3.871) (-5.525,-3.907) (-5.575,-3.942) (-5.622,-3.975) (-5.675,-4.013) (-5.725,-4.048) (-5.775,-4.084) (-5.825,-4.119) (-5.869,-4.150) (-5.925,-4.190) (-5.975,-4.225) (-6.025,-4.260) (-6.075,-4.296) (-6.117,-4.325) (-6.175,-4.366) (-6.223,-4.400) (-6.275,-4.437) (-6.325,-4.472) (-6.375,-4.508) (-6.425,-4.543) (-6.470,-4.575) (-6.525,-4.614) (-6.575,-4.649) (-6.625,-4.685) (-6.675,-4.720) (-6.718,-4.750) (-6.775,-4.791) (-6.824,-4.825) (-6.875,-4.861) (-6.925,-4.897) (-6.975,-4.932) (-7.025,-4.967) (-7.071,-5.000) (-7.125,-5.038) (-7.175,-5.073) (-7.225,-5.109) (-7.275,-5.144) (-7.319,-5.175) (-7.375,-5.215) (-7.425,-5.250) (-7.475,-5.286) (-7.525,-5.321) (-7.566,-5.350) (-7.625,-5.392) (-7.672,-5.425) (-7.725,-5.462) (-7.775,-5.498) (-7.825,-5.533) (-7.875,-5.568) (-7.920,-5.600) (-7.975,-5.639) (-8.025,-5.675) (-8.075,-5.710) (-8.125,-5.745) (-8.167,-5.775) (-8.225,-5.816) (-8.273,-5.850) (-8.325,-5.887) (-8.375,-5.922) (-8.425,-5.957) (-8.475,-5.993) (-8.521,-6.025) (-8.575,-6.063) (-8.625,-6.099) (-8.675,-6.134) (-8.725,-6.170) (-8.768,-6.200) (-8.825,-6.240) (-8.874,-6.275) (-8.925,-6.311) (-8.975,-6.346) (-9.025,-6.382) (-9.075,-6.417) (-9.122,-6.450) (-9.175,-6.488) (-9.225,-6.523) (-9.275,-6.558) (-9.325,-6.594) (-9.369,-6.625) (-9.425,-6.665) (-9.475,-6.700) (-9.525,-6.735) (-9.575,-6.771) (-9.617,-6.800) (-9.675,-6.841) (-9.723,-6.875) (-9.775,-6.912) (-9.825,-6.947) (-9.875,-6.983) (-9.925,-7.018) (-9.970,-7.050) (-10.000,-7.071)}
\EQ{(0.206,0.144) (0.248,0.175) (0.300,0.213) (0.350,0.248) (0.400,0.283) (0.450,0.319) (0.496,0.350) (0.550,0.389) (0.600,0.424) (0.650,0.460) (0.700,0.495) (0.743,0.525) (0.800,0.566) (0.849,0.600) (0.900,0.637) (0.950,0.672) (1.000,0.707) (1.050,0.743) (1.096,0.775) (1.150,0.813) (1.200,0.849) (1.250,0.884) (1.300,0.919) (1.344,0.950) (1.400,0.990) (1.450,1.025) (1.500,1.061) (1.550,1.096) (1.591,1.125) (1.650,1.167) (1.697,1.200) (1.750,1.238) (1.800,1.273) (1.850,1.308) (1.900,1.344) (1.945,1.375) (2.000,1.414) (2.050,1.450) (2.100,1.485) (2.150,1.520) (2.192,1.550) (2.250,1.591) (2.298,1.625) (2.350,1.662) (2.400,1.697) (2.450,1.732) (2.500,1.768) (2.546,1.800) (2.600,1.839) (2.650,1.874) (2.700,1.909) (2.750,1.945) (2.793,1.975) (2.850,2.015) (2.899,2.050) (2.950,2.086) (3.000,2.121) (3.050,2.157) (3.100,2.192) (3.147,2.225) (3.200,2.263) (3.250,2.298) (3.300,2.334) (3.350,2.369) (3.394,2.400) (3.450,2.440) (3.500,2.475) (3.550,2.510) (3.600,2.546) (3.642,2.575) (3.700,2.616) (3.748,2.650) (3.800,2.687) (3.850,2.722) (3.900,2.758) (3.950,2.793) (3.995,2.825) (4.050,2.864) (4.100,2.899) (4.150,2.935) (4.200,2.970) (4.243,3.000) (4.300,3.041) (4.349,3.075) (4.400,3.111) (4.450,3.147) (4.500,3.182) (4.550,3.217) (4.596,3.250) (4.650,3.288) (4.700,3.323) (4.750,3.359) (4.800,3.394) (4.844,3.425) (4.900,3.465) (4.950,3.500) (5.000,3.536) (5.050,3.571) (5.091,3.600) (5.150,3.642) (5.197,3.675) (5.250,3.712) (5.300,3.748) (5.350,3.783) (5.400,3.818) (5.445,3.850) (5.500,3.889) (5.550,3.924) (5.600,3.960) (5.650,3.995) (5.692,4.025) (5.750,4.066) (5.798,4.100) (5.850,4.137) (5.900,4.172) (5.950,4.207) (6.000,4.243) (6.046,4.275) (6.100,4.313) (6.150,4.349) (6.200,4.384) (6.250,4.419) (6.293,4.450) (6.350,4.490) (6.399,4.525) (6.450,4.561) (6.500,4.596) (6.550,4.632) (6.600,4.667) (6.647,4.700) (6.700,4.738) (6.750,4.773) (6.800,4.808) (6.850,4.844) (6.894,4.875) (6.950,4.914) (7.000,4.950) (7.050,4.985) (7.100,5.020) (7.142,5.050) (7.200,5.091) (7.248,5.125) (7.300,5.162) (7.350,5.197) (7.400,5.233) (7.450,5.268) (7.495,5.300) (7.550,5.339) (7.600,5.374) (7.650,5.409) (7.700,5.445) (7.743,5.475) (7.800,5.515) (7.849,5.550) (7.900,5.586) (7.950,5.622) (8.000,5.657) (8.050,5.692) (8.096,5.725) (8.150,5.763) (8.200,5.798) (8.250,5.834) (8.300,5.869) (8.344,5.900) (8.400,5.940) (8.450,5.975) (8.500,6.010) (8.550,6.046) (8.591,6.075) (8.650,6.117) (8.697,6.150) (8.750,6.187) (8.800,6.223) (8.850,6.258) (8.900,6.293) (8.945,6.325) (9.000,6.364) (9.050,6.399) (9.100,6.435) (9.150,6.470) (9.192,6.500) (9.250,6.541) (9.298,6.575) (9.350,6.611) (9.400,6.647) (9.450,6.682) (9.500,6.718) (9.546,6.750) (9.600,6.788) (9.650,6.824) (9.700,6.859) (9.750,6.894) (9.793,6.925) (9.850,6.965) (9.899,7.000) (9.950,7.036) (10.000,7.071)}
\EQ{(-10.000,7.071) (-9.950,7.036) (-9.900,7.000) (-9.850,6.965) (-9.800,6.930) (-9.758,6.900) (-9.700,6.859) (-9.652,6.825) (-9.600,6.788) (-9.550,6.753) (-9.500,6.718) (-9.450,6.682) (-9.405,6.650) (-9.350,6.611) (-9.300,6.576) (-9.250,6.541) (-9.200,6.505) (-9.157,6.475) (-9.100,6.435) (-9.051,6.400) (-9.000,6.364) (-8.950,6.329) (-8.900,6.293) (-8.850,6.258) (-8.804,6.225) (-8.750,6.187) (-8.700,6.152) (-8.650,6.117) (-8.600,6.081) (-8.556,6.050) (-8.500,6.010) (-8.450,5.975) (-8.400,5.940) (-8.350,5.904) (-8.309,5.875) (-8.250,5.834) (-8.202,5.800) (-8.150,5.763) (-8.100,5.728) (-8.050,5.692) (-8.000,5.657) (-7.955,5.625) (-7.900,5.586) (-7.850,5.551) (-7.800,5.515) (-7.750,5.480) (-7.708,5.450) (-7.650,5.409) (-7.601,5.375) (-7.550,5.339) (-7.500,5.303) (-7.450,5.268) (-7.400,5.233) (-7.354,5.200) (-7.300,5.162) (-7.250,5.127) (-7.200,5.091) (-7.150,5.056) (-7.106,5.025) (-7.050,4.985) (-7.000,4.950) (-6.950,4.914) (-6.900,4.879) (-6.859,4.850) (-6.800,4.808) (-6.753,4.775) (-6.700,4.738) (-6.650,4.702) (-6.600,4.667) (-6.550,4.632) (-6.505,4.600) (-6.450,4.561) (-6.400,4.525) (-6.350,4.490) (-6.300,4.455) (-6.258,4.425) (-6.200,4.384) (-6.152,4.350) (-6.100,4.313) (-6.050,4.278) (-6.000,4.243) (-5.950,4.207) (-5.904,4.175) (-5.850,4.137) (-5.800,4.101) (-5.750,4.066) (-5.700,4.031) (-5.657,4.000) (-5.600,3.960) (-5.551,3.925) (-5.500,3.889) (-5.450,3.854) (-5.400,3.818) (-5.350,3.783) (-5.303,3.750) (-5.250,3.712) (-5.200,3.677) (-5.150,3.642) (-5.100,3.606) (-5.056,3.575) (-5.000,3.536) (-4.950,3.500) (-4.900,3.465) (-4.850,3.429) (-4.808,3.400) (-4.750,3.359) (-4.702,3.325) (-4.650,3.288) (-4.600,3.253) (-4.550,3.217) (-4.500,3.182) (-4.455,3.150) (-4.400,3.111) (-4.350,3.076) (-4.300,3.041) (-4.250,3.005) (-4.207,2.975) (-4.150,2.935) (-4.101,2.900) (-4.050,2.864) (-4.000,2.828) (-3.950,2.793) (-3.900,2.758) (-3.854,2.725) (-3.800,2.687) (-3.750,2.652) (-3.700,2.616) (-3.650,2.581) (-3.606,2.550) (-3.550,2.510) (-3.500,2.475) (-3.450,2.440) (-3.400,2.404) (-3.359,2.375) (-3.300,2.334) (-3.253,2.300) (-3.200,2.263) (-3.150,2.227) (-3.100,2.192) (-3.050,2.157) (-3.005,2.125) (-2.950,2.086) (-2.900,2.051) (-2.850,2.015) (-2.800,1.980) (-2.758,1.950) (-2.700,1.909) (-2.652,1.875) (-2.600,1.839) (-2.550,1.803) (-2.500,1.768) (-2.450,1.732) (-2.404,1.700) (-2.350,1.662) (-2.300,1.626) (-2.250,1.591) (-2.200,1.556) (-2.157,1.525) (-2.100,1.485) (-2.051,1.450) (-2.000,1.414) (-1.950,1.379) (-1.900,1.344) (-1.850,1.308) (-1.803,1.275) (-1.750,1.238) (-1.700,1.202) (-1.650,1.167) (-1.600,1.131) (-1.556,1.100) (-1.500,1.061) (-1.450,1.025) (-1.400,0.990) (-1.350,0.955) (-1.308,0.925) (-1.250,0.884) (-1.202,0.850) (-1.150,0.813) (-1.100,0.778) (-1.050,0.743) (-1.000,0.707) (-0.955,0.675) (-0.900,0.637) (-0.850,0.601) (-0.800,0.566) (-0.750,0.531) (-0.708,0.500) (-0.650,0.460) (-0.601,0.425) (-0.550,0.389) (-0.500,0.354) (-0.450,0.319) (-0.400,0.283) (-0.354,0.250) (-0.300,0.213) (-0.250,0.177) (-0.220,0.155)}
\EQ{(4.306,-10.000) (4.350,-9.950) (4.392,-9.900) (4.425,-9.860) (4.473,-9.800) (4.513,-9.750) (4.550,-9.701) (4.588,-9.650) (4.624,-9.600) (4.660,-9.550) (4.695,-9.500) (4.725,-9.455) (4.761,-9.400) (4.793,-9.350) (4.824,-9.300) (4.855,-9.250) (4.884,-9.200) (4.913,-9.150) (4.941,-9.100) (4.969,-9.050) (4.995,-9.000) (5.021,-8.950) (5.046,-8.900) (5.071,-8.850) (5.095,-8.800) (5.118,-8.750) (5.140,-8.700) (5.162,-8.650) (5.183,-8.600) (5.203,-8.550) (5.223,-8.500) (5.242,-8.450) (5.261,-8.400) (5.279,-8.350) (5.296,-8.300) (5.313,-8.250) (5.329,-8.200) (5.344,-8.150) (5.359,-8.100) (5.373,-8.050) (5.387,-8.000) (5.400,-7.951) (5.413,-7.900) (5.425,-7.850) (5.436,-7.800) (5.447,-7.750) (5.457,-7.700) (5.467,-7.650) (5.476,-7.600) (5.484,-7.550) (5.492,-7.500) (5.500,-7.450) (5.507,-7.400) (5.513,-7.350) (5.519,-7.300) (5.524,-7.250) (5.529,-7.200) (5.533,-7.150) (5.537,-7.100) (5.540,-7.050) (5.542,-7.000) (5.544,-6.950) (5.545,-6.900) (5.546,-6.850) (5.547,-6.800) (5.546,-6.750) (5.546,-6.700) (5.544,-6.650) (5.543,-6.600) (5.540,-6.550) (5.537,-6.500) (5.534,-6.450) (5.530,-6.400) (5.525,-6.350) (5.520,-6.300) (5.515,-6.250) (5.509,-6.200) (5.502,-6.150) (5.495,-6.100) (5.487,-6.050) (5.478,-6.000) (5.469,-5.950) (5.460,-5.900) (5.450,-5.850) (5.439,-5.800) (5.428,-5.750) (5.416,-5.700) (5.404,-5.650) (5.391,-5.600) (5.378,-5.550) (5.363,-5.500) (5.349,-5.450) (5.334,-5.400) (5.318,-5.350) (5.301,-5.300) (5.284,-5.250) (5.267,-5.200) (5.249,-5.150) (5.230,-5.100) (5.210,-5.050) (5.190,-5.000) (5.170,-4.950) (5.150,-4.903) (5.127,-4.850) (5.104,-4.800) (5.081,-4.750) (5.057,-4.700) (5.033,-4.650) (5.008,-4.600) (4.982,-4.550) (4.956,-4.500) (4.929,-4.450) (4.901,-4.400) (4.875,-4.354) (4.844,-4.300) (4.814,-4.250) (4.784,-4.200) (4.753,-4.150) (4.725,-4.107) (4.688,-4.050) (4.655,-4.000) (4.625,-3.955) (4.587,-3.900) (4.551,-3.850) (4.516,-3.800) (4.479,-3.750) (4.441,-3.700) (4.403,-3.650) (4.364,-3.600) (4.325,-3.551) (4.284,-3.500) (4.250,-3.458) (4.201,-3.400) (4.159,-3.350) (4.125,-3.311) (4.075,-3.254) (4.026,-3.200) (3.981,-3.150) (3.935,-3.100) (3.900,-3.063) (3.864,-3.025) (3.825,-2.985) (3.775,-2.934) (3.725,-2.883) (3.675,-2.834) (3.625,-2.785) (3.575,-2.736) (3.537,-2.700) (3.500,-2.665) (3.457,-2.625) (3.403,-2.575) (3.350,-2.526) (3.300,-2.481) (3.250,-2.437) (3.209,-2.400) (3.151,-2.350) (3.100,-2.306) (3.050,-2.263) (3.006,-2.225) (2.950,-2.178) (2.900,-2.137) (2.856,-2.100) (2.800,-2.054) (2.750,-2.014) (2.702,-1.975) (2.650,-1.933) (2.609,-1.900) (2.550,-1.854) (2.500,-1.814) (2.450,-1.775) (2.400,-1.737) (2.353,-1.700) (2.300,-1.660) (2.255,-1.625) (2.200,-1.583) (2.156,-1.550) (2.100,-1.508) (2.056,-1.475) (2.000,-1.433) (1.956,-1.400) (1.900,-1.359) (1.854,-1.325) (1.800,-1.285) (1.752,-1.250) (1.700,-1.212) (1.650,-1.175) (1.600,-1.139) (1.550,-1.103) (1.500,-1.067) (1.450,-1.031) (1.408,-1.000) (1.350,-0.959) (1.304,-0.925) (1.250,-0.887) (1.200,-0.851) (1.150,-0.815) (1.100,-0.780) (1.059,-0.750) (1.000,-0.708) (0.953,-0.675) (0.900,-0.637) (0.850,-0.602) (0.800,-0.566) (0.750,-0.531) (0.707,-0.500) (0.650,-0.460) (0.601,-0.425) (0.550,-0.389) (0.500,-0.354) (0.450,-0.319) (0.400,-0.283) (0.354,-0.250) (0.300,-0.213) (0.250,-0.177) (0.206,-0.144)}
\EQ{(-0.220,-0.155) (-0.275,-0.195) (-0.319,-0.225) (-0.375,-0.266) (-0.424,-0.300) (-0.475,-0.336) (-0.525,-0.372) (-0.566,-0.400) (-0.625,-0.442) (-0.672,-0.475) (-0.725,-0.513) (-0.775,-0.549) (-0.825,-0.584) (-0.875,-0.620) (-0.918,-0.650) (-0.975,-0.691) (-1.024,-0.725) (-1.075,-0.762) (-1.125,-0.797) (-1.175,-0.833) (-1.225,-0.869) (-1.269,-0.900) (-1.325,-0.941) (-1.373,-0.975) (-1.425,-1.013) (-1.475,-1.049) (-1.525,-1.085) (-1.575,-1.121) (-1.625,-1.157) (-1.675,-1.194) (-1.718,-1.225) (-1.775,-1.267) (-1.820,-1.300) (-1.875,-1.340) (-1.922,-1.375) (-1.975,-1.415) (-2.023,-1.450) (-2.075,-1.489) (-2.123,-1.525) (-2.175,-1.564) (-2.222,-1.600) (-2.275,-1.640) (-2.320,-1.675) (-2.375,-1.717) (-2.418,-1.750) (-2.475,-1.795) (-2.525,-1.834) (-2.575,-1.873) (-2.625,-1.913) (-2.671,-1.950) (-2.725,-1.993) (-2.764,-2.025) (-2.825,-2.075) (-2.875,-2.116) (-2.916,-2.150) (-2.975,-2.199) (-3.025,-2.241) (-3.065,-2.275) (-3.123,-2.325) (-3.175,-2.371) (-3.225,-2.415) (-3.265,-2.450) (-3.321,-2.500) (-3.375,-2.549) (-3.425,-2.595) (-3.475,-2.642) (-3.525,-2.689) (-3.563,-2.725) (-3.600,-2.760) (-3.641,-2.800) (-3.692,-2.850) (-3.742,-2.900) (-3.791,-2.950) (-3.840,-3.000) (-3.875,-3.037) (-3.911,-3.075) (-3.950,-3.117) (-4.000,-3.171) (-4.049,-3.225) (-4.093,-3.275) (-4.137,-3.325) (-4.175,-3.369) (-4.222,-3.425) (-4.264,-3.475) (-4.300,-3.519) (-4.345,-3.575) (-4.384,-3.625) (-4.422,-3.675) (-4.460,-3.725) (-4.497,-3.775) (-4.534,-3.825) (-4.569,-3.875) (-4.600,-3.919) (-4.638,-3.975) (-4.672,-4.025) (-4.700,-4.068) (-4.737,-4.125) (-4.768,-4.175) (-4.799,-4.225) (-4.825,-4.269) (-4.858,-4.325) (-4.887,-4.375) (-4.915,-4.425) (-4.942,-4.475) (-4.969,-4.525) (-4.995,-4.575) (-5.020,-4.625) (-5.045,-4.675) (-5.069,-4.725) (-5.093,-4.775) (-5.116,-4.825) (-5.138,-4.875) (-5.159,-4.925) (-5.180,-4.975) (-5.200,-5.024) (-5.220,-5.075) (-5.239,-5.125) (-5.258,-5.175) (-5.275,-5.223) (-5.293,-5.275) (-5.310,-5.325) (-5.325,-5.373) (-5.341,-5.425) (-5.356,-5.475) (-5.371,-5.525) (-5.384,-5.575) (-5.398,-5.625) (-5.410,-5.675) (-5.422,-5.725) (-5.434,-5.775) (-5.445,-5.825) (-5.455,-5.875) (-5.465,-5.925) (-5.474,-5.975) (-5.483,-6.025) (-5.491,-6.075) (-5.498,-6.125) (-5.505,-6.175) (-5.512,-6.225) (-5.518,-6.275) (-5.523,-6.325) (-5.528,-6.375) (-5.532,-6.425) (-5.536,-6.475) (-5.539,-6.525) (-5.542,-6.575) (-5.544,-6.625) (-5.545,-6.675) (-5.546,-6.725) (-5.547,-6.775) (-5.547,-6.825) (-5.546,-6.875) (-5.545,-6.925) (-5.543,-6.975) (-5.541,-7.025) (-5.538,-7.075) (-5.535,-7.125) (-5.531,-7.175) (-5.526,-7.225) (-5.522,-7.275) (-5.516,-7.325) (-5.510,-7.375) (-5.503,-7.425) (-5.496,-7.475) (-5.489,-7.525) (-5.480,-7.575) (-5.471,-7.625) (-5.462,-7.675) (-5.452,-7.725) (-5.442,-7.775) (-5.430,-7.825) (-5.419,-7.875) (-5.407,-7.925) (-5.394,-7.975) (-5.380,-8.025) (-5.366,-8.075) (-5.352,-8.125) (-5.337,-8.175) (-5.321,-8.225) (-5.305,-8.275) (-5.288,-8.325) (-5.270,-8.375) (-5.252,-8.425) (-5.233,-8.475) (-5.213,-8.525) (-5.193,-8.575) (-5.173,-8.625) (-5.151,-8.675) (-5.129,-8.725) (-5.106,-8.775) (-5.083,-8.825) (-5.059,-8.875) (-5.034,-8.925) (-5.008,-8.975) (-4.982,-9.025) (-4.955,-9.075) (-4.927,-9.125) (-4.900,-9.173) (-4.870,-9.225) (-4.840,-9.275) (-4.809,-9.325) (-4.777,-9.375) (-4.750,-9.417) (-4.711,-9.475) (-4.677,-9.525) (-4.642,-9.575) (-4.606,-9.625) (-4.575,-9.668) (-4.532,-9.725) (-4.500,-9.766) (-4.453,-9.825) (-4.413,-9.875) (-4.375,-9.920) (-4.328,-9.975) (-4.306,-10.000)}
\EQ{(0.206,0.144) (0.248,0.175) (0.300,0.213) (0.350,0.248) (0.400,0.283) (0.450,0.319) (0.495,0.350) (0.550,0.389) (0.600,0.424) (0.650,0.460) (0.700,0.495) (0.742,0.525) (0.800,0.566) (0.848,0.600) (0.900,0.637) (0.950,0.673) (1.000,0.708) (1.050,0.744) (1.094,0.775) (1.150,0.815) (1.199,0.850) (1.250,0.887) (1.300,0.923) (1.350,0.959) (1.400,0.995) (1.443,1.025) (1.500,1.067) (1.546,1.100) (1.600,1.139) (1.650,1.175) (1.700,1.212) (1.750,1.248) (1.800,1.285) (1.850,1.322) (1.900,1.359) (1.950,1.396) (2.000,1.433) (2.050,1.470) (2.100,1.508) (2.150,1.546) (2.200,1.583) (2.250,1.621) (2.300,1.660) (2.350,1.698) (2.400,1.737) (2.450,1.775) (2.500,1.814) (2.546,1.850) (2.600,1.893) (2.640,1.925) (2.700,1.973) (2.750,2.014) (2.795,2.050) (2.850,2.095) (2.900,2.137) (2.946,2.175) (3.000,2.220) (3.050,2.263) (3.094,2.300) (3.150,2.349) (3.200,2.393) (3.250,2.437) (3.293,2.475) (3.349,2.525) (3.400,2.572) (3.450,2.618) (3.500,2.665) (3.537,2.700) (3.575,2.736) (3.615,2.775) (3.666,2.825) (3.717,2.875) (3.767,2.925) (3.815,2.975) (3.864,3.025) (3.900,3.063) (3.935,3.100) (3.975,3.144) (4.025,3.198) (4.071,3.250) (4.115,3.300) (4.150,3.340) (4.200,3.399) (4.243,3.450) (4.275,3.489) (4.325,3.550) (4.364,3.600) (4.400,3.646) (4.441,3.700) (4.475,3.745) (4.516,3.800) (4.550,3.848) (4.587,3.900) (4.621,3.950) (4.650,3.992) (4.688,4.050) (4.721,4.100) (4.750,4.146) (4.784,4.200) (4.814,4.250) (4.844,4.300) (4.873,4.350) (4.900,4.398) (4.925,4.443) (4.950,4.489) (4.975,4.536) (5.000,4.585) (5.025,4.634) (5.050,4.685) (5.075,4.737) (5.100,4.791) (5.125,4.846) (5.149,4.900) (5.170,4.950) (5.190,5.000) (5.210,5.050) (5.230,5.100) (5.249,5.150) (5.267,5.200) (5.284,5.250) (5.301,5.300) (5.318,5.350) (5.334,5.400) (5.349,5.450) (5.363,5.500) (5.378,5.550) (5.391,5.600) (5.404,5.650) (5.416,5.700) (5.428,5.750) (5.439,5.800) (5.450,5.850) (5.460,5.900) (5.469,5.950) (5.478,6.000) (5.487,6.050) (5.495,6.100) (5.502,6.150) (5.509,6.200) (5.515,6.250) (5.520,6.300) (5.525,6.350) (5.530,6.400) (5.534,6.450) (5.537,6.500) (5.540,6.550) (5.543,6.600) (5.544,6.650) (5.546,6.700) (5.546,6.750) (5.547,6.800) (5.546,6.850) (5.545,6.900) (5.544,6.950) (5.542,7.000) (5.540,7.050) (5.537,7.100) (5.533,7.150) (5.529,7.200) (5.524,7.250) (5.519,7.300) (5.513,7.350) (5.507,7.400) (5.500,7.450) (5.492,7.500) (5.484,7.550) (5.476,7.600) (5.467,7.650) (5.457,7.700) (5.447,7.750) (5.436,7.800) (5.425,7.850) (5.413,7.900) (5.400,7.950) (5.387,8.000) (5.373,8.050) (5.359,8.100) (5.344,8.150) (5.329,8.200) (5.313,8.250) (5.296,8.300) (5.279,8.350) (5.261,8.400) (5.242,8.450) (5.223,8.500) (5.203,8.550) (5.183,8.600) (5.162,8.650) (5.140,8.700) (5.118,8.750) (5.095,8.800) (5.071,8.850) (5.046,8.900) (5.021,8.950) (4.995,9.000) (4.969,9.050) (4.941,9.100) (4.913,9.150) (4.884,9.200) (4.855,9.250) (4.825,9.299) (4.793,9.350) (4.761,9.400) (4.728,9.450) (4.700,9.492) (4.660,9.550) (4.625,9.599) (4.588,9.650) (4.551,9.700) (4.513,9.750) (4.475,9.798) (4.433,9.850) (4.400,9.890) (4.350,9.949) (4.306,10.000)}
\EQ{(-4.306,10.000) (-4.350,9.950) (-4.392,9.900) (-4.425,9.860) (-4.473,9.800) (-4.513,9.750) (-4.550,9.701) (-4.588,9.650) (-4.624,9.600) (-4.660,9.550) (-4.695,9.500) (-4.725,9.455) (-4.761,9.400) (-4.793,9.350) (-4.824,9.300) (-4.855,9.250) (-4.884,9.200) (-4.913,9.150) (-4.941,9.100) (-4.969,9.050) (-4.995,9.000) (-5.021,8.950) (-5.046,8.900) (-5.071,8.850) (-5.095,8.800) (-5.118,8.750) (-5.140,8.700) (-5.162,8.650) (-5.183,8.600) (-5.203,8.550) (-5.223,8.500) (-5.242,8.450) (-5.261,8.400) (-5.279,8.350) (-5.296,8.300) (-5.313,8.250) (-5.329,8.200) (-5.344,8.150) (-5.359,8.100) (-5.373,8.050) (-5.387,8.000) (-5.400,7.951) (-5.413,7.900) (-5.425,7.850) (-5.436,7.800) (-5.447,7.750) (-5.457,7.700) (-5.467,7.650) (-5.476,7.600) (-5.484,7.550) (-5.492,7.500) (-5.500,7.450) (-5.507,7.400) (-5.513,7.350) (-5.519,7.300) (-5.524,7.250) (-5.529,7.200) (-5.533,7.150) (-5.537,7.100) (-5.540,7.050) (-5.542,7.000) (-5.544,6.950) (-5.545,6.900) (-5.546,6.850) (-5.547,6.800) (-5.546,6.750) (-5.546,6.700) (-5.544,6.650) (-5.543,6.600) (-5.540,6.550) (-5.537,6.500) (-5.534,6.450) (-5.530,6.400) (-5.525,6.350) (-5.520,6.300) (-5.515,6.250) (-5.509,6.200) (-5.502,6.150) (-5.495,6.100) (-5.487,6.050) (-5.478,6.000) (-5.469,5.950) (-5.460,5.900) (-5.450,5.850) (-5.439,5.800) (-5.428,5.750) (-5.416,5.700) (-5.404,5.650) (-5.391,5.600) (-5.378,5.550) (-5.363,5.500) (-5.349,5.450) (-5.334,5.400) (-5.318,5.350) (-5.301,5.300) (-5.284,5.250) (-5.267,5.200) (-5.249,5.150) (-5.230,5.100) (-5.210,5.050) (-5.190,5.000) (-5.170,4.950) (-5.150,4.903) (-5.127,4.850) (-5.104,4.800) (-5.081,4.750) (-5.057,4.700) (-5.033,4.650) (-5.008,4.600) (-4.982,4.550) (-4.956,4.500) (-4.929,4.450) (-4.901,4.400) (-4.875,4.354) (-4.844,4.300) (-4.814,4.250) (-4.784,4.200) (-4.753,4.150) (-4.725,4.107) (-4.688,4.050) (-4.655,4.000) (-4.625,3.955) (-4.587,3.900) (-4.551,3.850) (-4.516,3.800) (-4.479,3.750) (-4.441,3.700) (-4.403,3.650) (-4.364,3.600) (-4.325,3.551) (-4.284,3.500) (-4.250,3.458) (-4.201,3.400) (-4.159,3.350) (-4.125,3.311) (-4.075,3.254) (-4.026,3.200) (-3.981,3.150) (-3.935,3.100) (-3.900,3.063) (-3.864,3.025) (-3.825,2.985) (-3.775,2.934) (-3.725,2.883) (-3.675,2.834) (-3.625,2.785) (-3.575,2.736) (-3.537,2.700) (-3.500,2.665) (-3.457,2.625) (-3.403,2.575) (-3.350,2.526) (-3.300,2.481) (-3.250,2.437) (-3.209,2.400) (-3.151,2.350) (-3.100,2.306) (-3.050,2.263) (-3.006,2.225) (-2.950,2.178) (-2.900,2.137) (-2.856,2.100) (-2.800,2.054) (-2.750,2.014) (-2.702,1.975) (-2.650,1.933) (-2.609,1.900) (-2.550,1.854) (-2.500,1.814) (-2.450,1.775) (-2.400,1.737) (-2.353,1.700) (-2.300,1.660) (-2.255,1.625) (-2.200,1.583) (-2.156,1.550) (-2.100,1.508) (-2.056,1.475) (-2.000,1.433) (-1.956,1.400) (-1.900,1.359) (-1.854,1.325) (-1.800,1.285) (-1.752,1.250) (-1.700,1.212) (-1.650,1.175) (-1.600,1.139) (-1.550,1.103) (-1.500,1.067) (-1.450,1.031) (-1.408,1.000) (-1.350,0.959) (-1.304,0.925) (-1.250,0.887) (-1.200,0.851) (-1.150,0.815) (-1.100,0.780) (-1.059,0.750) (-1.000,0.708) (-0.953,0.675) (-0.900,0.637) (-0.850,0.602) (-0.800,0.566) (-0.750,0.531) (-0.707,0.500) (-0.650,0.460) (-0.601,0.425) (-0.550,0.389) (-0.500,0.354) (-0.450,0.319) (-0.400,0.283) (-0.354,0.250) (-0.300,0.213) (-0.250,0.177) (-0.220,0.155)}
\EQ{(-0.220,-0.155) (-0.275,-0.195) (-0.319,-0.225) (-0.375,-0.266) (-0.424,-0.300) (-0.475,-0.337) (-0.525,-0.372) (-0.575,-0.407) (-0.625,-0.443) (-0.671,-0.475) (-0.725,-0.514) (-0.775,-0.549) (-0.825,-0.585) (-0.875,-0.620) (-0.917,-0.650) (-0.975,-0.692) (-1.022,-0.725) (-1.075,-0.763) (-1.125,-0.799) (-1.175,-0.835) (-1.225,-0.872) (-1.275,-0.908) (-1.325,-0.944) (-1.368,-0.975) (-1.425,-1.017) (-1.470,-1.050) (-1.525,-1.091) (-1.571,-1.125) (-1.625,-1.166) (-1.671,-1.200) (-1.725,-1.241) (-1.770,-1.275) (-1.825,-1.317) (-1.869,-1.350) (-1.925,-1.394) (-1.965,-1.425) (-2.025,-1.472) (-2.075,-1.512) (-2.123,-1.550) (-2.175,-1.592) (-2.216,-1.625) (-2.275,-1.674) (-2.325,-1.715) (-2.367,-1.750) (-2.425,-1.799) (-2.475,-1.843) (-2.525,-1.886) (-2.570,-1.925) (-2.625,-1.975) (-2.675,-2.020) (-2.725,-2.066) (-2.762,-2.100) (-2.800,-2.136) (-2.841,-2.175) (-2.893,-2.225) (-2.943,-2.275) (-2.993,-2.325) (-3.042,-2.375) (-3.090,-2.425) (-3.125,-2.463) (-3.160,-2.500) (-3.200,-2.545) (-3.249,-2.600) (-3.293,-2.650) (-3.336,-2.700) (-3.375,-2.747) (-3.418,-2.800) (-3.450,-2.840) (-3.497,-2.900) (-3.535,-2.950) (-3.572,-3.000) (-3.608,-3.050) (-3.644,-3.100) (-3.675,-3.146) (-3.712,-3.200) (-3.744,-3.250) (-3.775,-3.299) (-3.807,-3.350) (-3.837,-3.400) (-3.866,-3.450) (-3.894,-3.500) (-3.922,-3.550) (-3.948,-3.600) (-3.974,-3.650) (-3.999,-3.700) (-4.023,-3.750) (-4.047,-3.800) (-4.069,-3.850) (-4.091,-3.900) (-4.112,-3.950) (-4.132,-4.000) (-4.152,-4.050) (-4.170,-4.100) (-4.188,-4.150) (-4.206,-4.200) (-4.222,-4.250) (-4.238,-4.300) (-4.253,-4.350) (-4.267,-4.400) (-4.280,-4.450) (-4.293,-4.500) (-4.305,-4.550) (-4.316,-4.600) (-4.327,-4.650) (-4.337,-4.700) (-4.346,-4.750) (-4.354,-4.800) (-4.362,-4.850) (-4.369,-4.900) (-4.376,-4.950) (-4.381,-5.000) (-4.386,-5.050) (-4.391,-5.100) (-4.394,-5.150) (-4.397,-5.200) (-4.400,-5.250) (-4.401,-5.300) (-4.402,-5.350) (-4.402,-5.400) (-4.402,-5.450) (-4.401,-5.500) (-4.399,-5.550) (-4.396,-5.600) (-4.393,-5.650) (-4.389,-5.700) (-4.385,-5.750) (-4.380,-5.800) (-4.374,-5.850) (-4.367,-5.900) (-4.360,-5.950) (-4.352,-6.000) (-4.343,-6.050) (-4.334,-6.100) (-4.323,-6.150) (-4.313,-6.200) (-4.301,-6.250) (-4.289,-6.300) (-4.276,-6.350) (-4.262,-6.400) (-4.247,-6.450) (-4.232,-6.500) (-4.216,-6.550) (-4.200,-6.598) (-4.182,-6.650) (-4.163,-6.700) (-4.144,-6.750) (-4.125,-6.798) (-4.104,-6.850) (-4.082,-6.900) (-4.060,-6.950) (-4.036,-7.000) (-4.012,-7.050) (-3.987,-7.100) (-3.961,-7.150) (-3.934,-7.200) (-3.907,-7.250) (-3.878,-7.300) (-3.850,-7.347) (-3.818,-7.400) (-3.786,-7.450) (-3.753,-7.500) (-3.725,-7.542) (-3.684,-7.600) (-3.650,-7.648) (-3.611,-7.700) (-3.575,-7.747) (-3.533,-7.800) (-3.500,-7.841) (-3.450,-7.900) (-3.406,-7.950) (-3.361,-8.000) (-3.325,-8.039) (-3.275,-8.092) (-3.225,-8.142) (-3.175,-8.191) (-3.125,-8.238) (-3.085,-8.275) (-3.029,-8.325) (-2.975,-8.372) (-2.925,-8.414) (-2.880,-8.450) (-2.825,-8.494) (-2.784,-8.525) (-2.725,-8.569) (-2.682,-8.600) (-2.625,-8.640) (-2.575,-8.674) (-2.525,-8.707) (-2.475,-8.739) (-2.425,-8.770) (-2.375,-8.800) (-2.325,-8.829) (-2.275,-8.858) (-2.225,-8.886) (-2.175,-8.912) (-2.125,-8.938) (-2.075,-8.964) (-2.025,-8.988) (-1.975,-9.012) (-1.925,-9.035) (-1.875,-9.057) (-1.825,-9.078) (-1.775,-9.099) (-1.725,-9.119) (-1.675,-9.139) (-1.625,-9.158) (-1.577,-9.175) (-1.525,-9.193) (-1.475,-9.210) (-1.425,-9.226) (-1.375,-9.242) (-1.325,-9.257) (-1.275,-9.272) (-1.225,-9.285) (-1.175,-9.298) (-1.125,-9.311) (-1.075,-9.323) (-1.025,-9.335) (-0.975,-9.345) (-0.925,-9.356) (-0.875,-9.365) (-0.825,-9.375) (-0.775,-9.383) (-0.725,-9.391) (-0.675,-9.399) (-0.625,-9.406) (-0.575,-9.412) (-0.525,-9.418) (-0.475,-9.423) (-0.425,-9.428) (-0.375,-9.432) (-0.325,-9.436) (-0.275,-9.439) (-0.225,-9.442) (-0.175,-9.444) (-0.125,-9.446) (-0.075,-9.447) (-0.025,-9.447) (0.050,-9.447) (0.100,-9.446) (0.150,-9.445) (0.200,-9.443) (0.250,-9.441) (0.300,-9.438) (0.350,-9.434) (0.400,-9.430) (0.450,-9.426) (0.500,-9.421) (0.550,-9.415) (0.600,-9.409) (0.650,-9.402) (0.700,-9.395) (0.750,-9.387) (0.800,-9.379) (0.850,-9.370) (0.900,-9.361) (0.950,-9.351) (1.000,-9.340) (1.050,-9.329) (1.100,-9.317) (1.150,-9.305) (1.200,-9.292) (1.250,-9.278) (1.300,-9.264) (1.349,-9.250) (1.400,-9.234) (1.450,-9.218) (1.500,-9.202) (1.550,-9.185) (1.600,-9.167) (1.650,-9.148) (1.700,-9.129) (1.750,-9.109) (1.800,-9.089) (1.850,-9.068) (1.900,-9.046) (1.946,-9.025) (2.000,-9.000) (2.050,-8.976) (2.100,-8.951) (2.150,-8.925) (2.198,-8.900) (2.244,-8.875) (2.289,-8.850) (2.333,-8.825) (2.400,-8.785) (2.450,-8.755) (2.497,-8.725) (2.550,-8.691) (2.600,-8.657) (2.647,-8.625) (2.700,-8.587) (2.750,-8.551) (2.800,-8.513) (2.849,-8.475) (2.900,-8.434) (2.942,-8.400) (3.000,-8.351) (3.050,-8.307) (3.100,-8.262) (3.139,-8.225) (3.192,-8.175) (3.242,-8.125) (3.291,-8.075) (3.338,-8.025) (3.375,-7.985) (3.425,-7.929) (3.471,-7.875) (3.513,-7.825) (3.550,-7.779) (3.592,-7.725) (3.625,-7.682) (3.667,-7.625) (3.700,-7.578) (3.736,-7.525) (3.770,-7.475) (3.800,-7.428) (3.833,-7.375) (3.863,-7.325) (3.892,-7.275) (3.921,-7.225) (3.948,-7.175) (3.974,-7.125) (4.000,-7.075) (4.024,-7.025) (4.048,-6.975) (4.071,-6.925) (4.093,-6.875) (4.114,-6.825) (4.134,-6.775) (4.154,-6.725) (4.173,-6.675) (4.191,-6.625) (4.208,-6.575) (4.224,-6.525) (4.240,-6.475) (4.255,-6.425) (4.269,-6.375) (4.282,-6.325) (4.295,-6.275) (4.307,-6.225) (4.318,-6.175) (4.329,-6.125) (4.338,-6.075) (4.347,-6.025) (4.356,-5.975) (4.364,-5.925) (4.370,-5.875) (4.377,-5.825) (4.382,-5.775) (4.387,-5.725) (4.391,-5.675) (4.395,-5.625) (4.398,-5.575) (4.400,-5.525) (4.401,-5.475) (4.402,-5.425) (4.402,-5.375) (4.402,-5.325) (4.401,-5.275) (4.399,-5.225) (4.396,-5.175) (4.393,-5.125) (4.389,-5.075) (4.384,-5.025) (4.379,-4.975) (4.373,-4.925) (4.366,-4.875) (4.358,-4.825) (4.350,-4.775) (4.342,-4.725) (4.332,-4.675) (4.322,-4.625) (4.311,-4.575) (4.299,-4.525) (4.287,-4.475) (4.274,-4.425) (4.260,-4.375) (4.245,-4.325) (4.230,-4.275) (4.214,-4.225) (4.197,-4.175) (4.180,-4.125) (4.161,-4.075) (4.142,-4.025) (4.122,-3.975) (4.102,-3.925) (4.080,-3.875) (4.058,-3.825) (4.035,-3.775) (4.011,-3.725) (3.987,-3.675) (3.961,-3.625) (3.935,-3.575) (3.908,-3.525) (3.880,-3.475) (3.852,-3.425) (3.825,-3.380) (3.792,-3.325) (3.760,-3.275) (3.728,-3.225) (3.700,-3.183) (3.661,-3.125) (3.626,-3.075) (3.590,-3.025) (3.554,-2.975) (3.516,-2.925) (3.478,-2.875) (3.438,-2.825) (3.400,-2.778) (3.357,-2.725) (3.315,-2.675) (3.275,-2.629) (3.227,-2.575) (3.182,-2.525) (3.136,-2.475) (3.100,-2.436) (3.050,-2.384) (3.000,-2.332) (2.950,-2.282) (2.900,-2.232) (2.850,-2.184) (2.800,-2.136) (2.762,-2.100) (2.725,-2.066) (2.681,-2.025) (2.626,-1.975) (2.575,-1.930) (2.525,-1.886) (2.484,-1.850) (2.426,-1.800) (2.375,-1.757) (2.325,-1.715) (2.277,-1.675) (2.225,-1.633) (2.185,-1.600) (2.125,-1.552) (2.075,-1.512) (2.029,-1.475) (1.975,-1.433) (1.933,-1.400) (1.875,-1.355) (1.825,-1.317) (1.775,-1.279) (1.725,-1.241) (1.675,-1.203) (1.625,-1.166) (1.575,-1.128) (1.525,-1.091) (1.475,-1.054) (1.425,-1.017) (1.375,-0.981) (1.333,-0.950) (1.275,-0.908) (1.230,-0.875) (1.175,-0.835) (1.126,-0.800) (1.075,-0.763) (1.025,-0.727) (0.975,-0.692) (0.925,-0.656) (0.882,-0.625) (0.825,-0.585) (0.777,-0.550) (0.725,-0.514) (0.675,-0.478) (0.625,-0.443) (0.575,-0.407) (0.531,-0.375) (0.475,-0.337) (0.425,-0.301) (0.375,-0.266) (0.325,-0.230) (0.284,-0.200) (0.225,-0.160) (0.206,-0.144)}
\EQ{(0.206,0.144) (0.248,0.175) (0.300,0.213) (0.350,0.248) (0.400,0.283) (0.450,0.319) (0.495,0.350) (0.550,0.390) (0.600,0.425) (0.650,0.460) (0.700,0.496) (0.742,0.525) (0.800,0.567) (0.847,0.600) (0.900,0.638) (0.950,0.674) (1.000,0.710) (1.050,0.745) (1.092,0.775) (1.150,0.817) (1.196,0.850) (1.250,0.890) (1.299,0.925) (1.350,0.963) (1.400,0.999) (1.450,1.036) (1.500,1.073) (1.550,1.110) (1.600,1.147) (1.650,1.184) (1.700,1.222) (1.750,1.260) (1.800,1.298) (1.850,1.336) (1.900,1.374) (1.950,1.413) (1.997,1.450) (2.050,1.492) (2.092,1.525) (2.150,1.572) (2.200,1.612) (2.246,1.650) (2.300,1.694) (2.350,1.736) (2.396,1.775) (2.450,1.821) (2.500,1.864) (2.541,1.900) (2.598,1.950) (2.650,1.997) (2.700,2.043) (2.750,2.089) (2.789,2.125) (2.825,2.160) (2.867,2.200) (2.918,2.250) (2.968,2.300) (3.017,2.350) (3.066,2.400) (3.113,2.450) (3.150,2.490) (3.200,2.545) (3.249,2.600) (3.293,2.650) (3.336,2.700) (3.375,2.747) (3.418,2.800) (3.450,2.840) (3.497,2.900) (3.535,2.950) (3.572,3.000) (3.608,3.050) (3.644,3.100) (3.675,3.146) (3.712,3.200) (3.744,3.250) (3.775,3.299) (3.807,3.350) (3.837,3.400) (3.866,3.450) (3.894,3.500) (3.922,3.550) (3.948,3.600) (3.974,3.650) (3.999,3.700) (4.023,3.750) (4.047,3.800) (4.069,3.850) (4.091,3.900) (4.112,3.950) (4.132,4.000) (4.152,4.050) (4.170,4.100) (4.188,4.150) (4.206,4.200) (4.222,4.250) (4.238,4.300) (4.253,4.350) (4.267,4.400) (4.280,4.450) (4.293,4.500) (4.305,4.550) (4.316,4.600) (4.327,4.650) (4.337,4.700) (4.346,4.750) (4.354,4.800) (4.362,4.850) (4.369,4.900) (4.376,4.950) (4.381,5.000) (4.386,5.050) (4.391,5.100) (4.394,5.150) (4.397,5.200) (4.400,5.250) (4.401,5.300) (4.402,5.350) (4.402,5.400) (4.402,5.450) (4.401,5.500) (4.399,5.550) (4.396,5.600) (4.393,5.650) (4.389,5.700) (4.385,5.750) (4.380,5.800) (4.374,5.850) (4.367,5.900) (4.360,5.950) (4.352,6.000) (4.343,6.050) (4.334,6.100) (4.323,6.150) (4.313,6.200) (4.301,6.250) (4.289,6.300) (4.276,6.350) (4.262,6.400) (4.247,6.450) (4.232,6.500) (4.216,6.550) (4.200,6.598) (4.182,6.650) (4.163,6.700) (4.144,6.750) (4.125,6.798) (4.104,6.850) (4.082,6.900) (4.060,6.950) (4.036,7.000) (4.012,7.050) (3.987,7.100) (3.961,7.150) (3.934,7.200) (3.907,7.250) (3.878,7.300) (3.850,7.347) (3.818,7.400) (3.786,7.450) (3.753,7.500) (3.725,7.542) (3.684,7.600) (3.650,7.648) (3.611,7.700) (3.575,7.747) (3.533,7.800) (3.500,7.841) (3.450,7.900) (3.406,7.950) (3.361,8.000) (3.325,8.039) (3.275,8.092) (3.225,8.142) (3.175,8.191) (3.125,8.238) (3.085,8.275) (3.029,8.325) (2.975,8.372) (2.925,8.414) (2.880,8.450) (2.825,8.494) (2.784,8.525) (2.725,8.569) (2.682,8.600) (2.625,8.640) (2.575,8.674) (2.525,8.707) (2.475,8.739) (2.425,8.770) (2.375,8.800) (2.325,8.829) (2.275,8.858) (2.225,8.886) (2.175,8.912) (2.125,8.938) (2.075,8.964) (2.025,8.988) (1.975,9.012) (1.925,9.035) (1.875,9.057) (1.825,9.078) (1.775,9.099) (1.725,9.119) (1.675,9.139) (1.625,9.158) (1.577,9.175) (1.525,9.193) (1.475,9.210) (1.425,9.226) (1.375,9.242) (1.325,9.257) (1.275,9.272) (1.225,9.285) (1.175,9.298) (1.125,9.311) (1.075,9.323) (1.025,9.335) (0.975,9.345) (0.925,9.356) (0.875,9.365) (0.825,9.375) (0.775,9.383) (0.725,9.391) (0.675,9.399) (0.625,9.406) (0.575,9.412) (0.525,9.418) (0.475,9.423) (0.425,9.428) (0.375,9.432) (0.325,9.436) (0.275,9.439) (0.225,9.442) (0.175,9.444) (0.125,9.446) (0.075,9.447) (0.025,9.447) (-0.050,9.447) (-0.100,9.446) (-0.150,9.445) (-0.200,9.443) (-0.250,9.441) (-0.300,9.438) (-0.350,9.434) (-0.400,9.430) (-0.450,9.426) (-0.500,9.421) (-0.550,9.415) (-0.600,9.409) (-0.650,9.402) (-0.700,9.395) (-0.750,9.387) (-0.800,9.379) (-0.850,9.370) (-0.900,9.361) (-0.950,9.351) (-1.000,9.340) (-1.050,9.329) (-1.100,9.317) (-1.150,9.305) (-1.200,9.292) (-1.250,9.278) (-1.300,9.264) (-1.349,9.250) (-1.400,9.234) (-1.450,9.218) (-1.500,9.202) (-1.550,9.185) (-1.600,9.167) (-1.650,9.148) (-1.700,9.129) (-1.750,9.109) (-1.800,9.089) (-1.850,9.068) (-1.900,9.046) (-1.946,9.025) (-2.000,9.000) (-2.050,8.976) (-2.100,8.951) (-2.150,8.925) (-2.198,8.900) (-2.244,8.875) (-2.289,8.850) (-2.333,8.825) (-2.400,8.785) (-2.450,8.755) (-2.497,8.725) (-2.550,8.691) (-2.600,8.657) (-2.647,8.625) (-2.700,8.587) (-2.750,8.551) (-2.800,8.513) (-2.849,8.475) (-2.900,8.434) (-2.942,8.400) (-3.000,8.351) (-3.050,8.307) (-3.100,8.262) (-3.139,8.225) (-3.192,8.175) (-3.242,8.125) (-3.291,8.075) (-3.338,8.025) (-3.375,7.985) (-3.425,7.929) (-3.471,7.875) (-3.513,7.825) (-3.550,7.779) (-3.592,7.725) (-3.625,7.682) (-3.667,7.625) (-3.700,7.578) (-3.736,7.525) (-3.770,7.475) (-3.800,7.428) (-3.833,7.375) (-3.863,7.325) (-3.892,7.275) (-3.921,7.225) (-3.948,7.175) (-3.974,7.125) (-4.000,7.075) (-4.024,7.025) (-4.048,6.975) (-4.071,6.925) (-4.093,6.875) (-4.114,6.825) (-4.134,6.775) (-4.154,6.725) (-4.173,6.675) (-4.191,6.625) (-4.208,6.575) (-4.224,6.525) (-4.240,6.475) (-4.255,6.425) (-4.269,6.375) (-4.282,6.325) (-4.295,6.275) (-4.307,6.225) (-4.318,6.175) (-4.329,6.125) (-4.338,6.075) (-4.347,6.025) (-4.356,5.975) (-4.364,5.925) (-4.370,5.875) (-4.377,5.825) (-4.382,5.775) (-4.387,5.725) (-4.391,5.675) (-4.395,5.625) (-4.398,5.575) (-4.400,5.525) (-4.401,5.475) (-4.402,5.425) (-4.402,5.375) (-4.402,5.325) (-4.401,5.275) (-4.399,5.225) (-4.396,5.175) (-4.393,5.125) (-4.389,5.075) (-4.384,5.025) (-4.379,4.975) (-4.373,4.925) (-4.366,4.875) (-4.358,4.825) (-4.350,4.775) (-4.342,4.725) (-4.332,4.675) (-4.322,4.625) (-4.311,4.575) (-4.299,4.525) (-4.287,4.475) (-4.274,4.425) (-4.260,4.375) (-4.245,4.325) (-4.230,4.275) (-4.214,4.225) (-4.197,4.175) (-4.180,4.125) (-4.161,4.075) (-4.142,4.025) (-4.122,3.975) (-4.102,3.925) (-4.080,3.875) (-4.058,3.825) (-4.035,3.775) (-4.011,3.725) (-3.987,3.675) (-3.961,3.625) (-3.935,3.575) (-3.908,3.525) (-3.880,3.475) (-3.852,3.425) (-3.825,3.380) (-3.792,3.325) (-3.760,3.275) (-3.728,3.225) (-3.700,3.183) (-3.661,3.125) (-3.626,3.075) (-3.590,3.025) (-3.554,2.975) (-3.516,2.925) (-3.478,2.875) (-3.438,2.825) (-3.400,2.778) (-3.357,2.725) (-3.315,2.675) (-3.275,2.629) (-3.227,2.575) (-3.182,2.525) (-3.136,2.475) (-3.100,2.436) (-3.050,2.384) (-3.000,2.332) (-2.950,2.282) (-2.900,2.232) (-2.850,2.184) (-2.800,2.136) (-2.762,2.100) (-2.725,2.066) (-2.681,2.025) (-2.626,1.975) (-2.575,1.930) (-2.525,1.886) (-2.484,1.850) (-2.426,1.800) (-2.375,1.757) (-2.325,1.715) (-2.277,1.675) (-2.225,1.633) (-2.185,1.600) (-2.125,1.552) (-2.075,1.512) (-2.029,1.475) (-1.975,1.433) (-1.933,1.400) (-1.875,1.355) (-1.825,1.317) (-1.775,1.279) (-1.725,1.241) (-1.675,1.203) (-1.625,1.166) (-1.575,1.128) (-1.525,1.091) (-1.475,1.054) (-1.425,1.017) (-1.375,0.981) (-1.333,0.950) (-1.275,0.908) (-1.230,0.875) (-1.175,0.835) (-1.126,0.800) (-1.075,0.763) (-1.025,0.727) (-0.975,0.692) (-0.925,0.656) (-0.882,0.625) (-0.825,0.585) (-0.777,0.550) (-0.725,0.514) (-0.675,0.478) (-0.625,0.443) (-0.575,0.407) (-0.531,0.375) (-0.475,0.337) (-0.425,0.301) (-0.375,0.266) (-0.325,0.230) (-0.284,0.200) (-0.225,0.160) (-0.220,0.155)}
\EQ{(-0.220,-0.155) (-0.275,-0.195) (-0.319,-0.225) (-0.375,-0.266) (-0.424,-0.300) (-0.475,-0.337) (-0.525,-0.372) (-0.575,-0.407) (-0.625,-0.443) (-0.671,-0.475) (-0.725,-0.514) (-0.775,-0.549) (-0.825,-0.585) (-0.875,-0.621) (-0.916,-0.650) (-0.975,-0.693) (-1.020,-0.725) (-1.075,-0.765) (-1.124,-0.800) (-1.175,-0.838) (-1.225,-0.874) (-1.275,-0.911) (-1.325,-0.948) (-1.375,-0.985) (-1.425,-1.022) (-1.475,-1.060) (-1.525,-1.098) (-1.575,-1.136) (-1.625,-1.174) (-1.675,-1.213) (-1.723,-1.250) (-1.775,-1.291) (-1.818,-1.325) (-1.875,-1.370) (-1.925,-1.411) (-1.973,-1.450) (-2.025,-1.493) (-2.075,-1.535) (-2.122,-1.575) (-2.175,-1.621) (-2.225,-1.664) (-2.265,-1.700) (-2.321,-1.750) (-2.375,-1.800) (-2.425,-1.846) (-2.475,-1.894) (-2.525,-1.942) (-2.575,-1.991) (-2.625,-2.042) (-2.675,-2.093) (-2.725,-2.146) (-2.775,-2.200) (-2.820,-2.250) (-2.864,-2.300) (-2.900,-2.342) (-2.949,-2.400) (-2.989,-2.450) (-3.025,-2.495) (-3.068,-2.550) (-3.100,-2.593) (-3.142,-2.650) (-3.175,-2.697) (-3.212,-2.750) (-3.245,-2.800) (-3.275,-2.846) (-3.309,-2.900) (-3.340,-2.950) (-3.369,-3.000) (-3.398,-3.050) (-3.425,-3.099) (-3.450,-3.146) (-3.475,-3.195) (-3.500,-3.245) (-3.525,-3.298) (-3.549,-3.350) (-3.571,-3.400) (-3.592,-3.450) (-3.612,-3.500) (-3.631,-3.550) (-3.650,-3.600) (-3.667,-3.650) (-3.684,-3.700) (-3.700,-3.750) (-3.715,-3.800) (-3.729,-3.850) (-3.742,-3.900) (-3.755,-3.950) (-3.767,-4.000) (-3.777,-4.050) (-3.787,-4.100) (-3.797,-4.150) (-3.805,-4.200) (-3.813,-4.250) (-3.820,-4.300) (-3.826,-4.350) (-3.831,-4.400) (-3.835,-4.450) (-3.839,-4.500) (-3.842,-4.550) (-3.844,-4.600) (-3.845,-4.650) (-3.846,-4.700) (-3.846,-4.750) (-3.845,-4.800) (-3.843,-4.850) (-3.840,-4.900) (-3.837,-4.950) (-3.833,-5.000) (-3.828,-5.050) (-3.822,-5.100) (-3.816,-5.150) (-3.808,-5.200) (-3.800,-5.250) (-3.791,-5.300) (-3.782,-5.350) (-3.771,-5.400) (-3.760,-5.450) (-3.747,-5.500) (-3.734,-5.550) (-3.721,-5.600) (-3.706,-5.650) (-3.690,-5.700) (-3.674,-5.750) (-3.656,-5.800) (-3.638,-5.850) (-3.619,-5.900) (-3.600,-5.948) (-3.578,-6.000) (-3.556,-6.050) (-3.533,-6.100) (-3.509,-6.150) (-3.484,-6.200) (-3.458,-6.250) (-3.431,-6.300) (-3.403,-6.350) (-3.375,-6.399) (-3.344,-6.450) (-3.313,-6.500) (-3.280,-6.550) (-3.250,-6.594) (-3.211,-6.650) (-3.175,-6.699) (-3.137,-6.750) (-3.100,-6.797) (-3.057,-6.850) (-3.015,-6.900) (-2.975,-6.946) (-2.926,-7.000) (-2.879,-7.050) (-2.830,-7.100) (-2.779,-7.150) (-2.727,-7.200) (-2.675,-7.247) (-2.625,-7.291) (-2.585,-7.325) (-2.525,-7.375) (-2.475,-7.414) (-2.429,-7.450) (-2.375,-7.490) (-2.326,-7.525) (-2.275,-7.560) (-2.225,-7.594) (-2.176,-7.625) (-2.125,-7.657) (-2.075,-7.688) (-2.025,-7.717) (-1.975,-7.745) (-1.925,-7.772) (-1.875,-7.799) (-1.825,-7.824) (-1.775,-7.849) (-1.725,-7.873) (-1.675,-7.896) (-1.625,-7.918) (-1.575,-7.939) (-1.525,-7.959) (-1.475,-7.979) (-1.425,-7.998) (-1.375,-8.016) (-1.325,-8.034) (-1.276,-8.050) (-1.225,-8.066) (-1.175,-8.082) (-1.125,-8.096) (-1.075,-8.110) (-1.025,-8.123) (-0.975,-8.136) (-0.925,-8.148) (-0.875,-8.159) (-0.825,-8.170) (-0.775,-8.179) (-0.725,-8.189) (-0.675,-8.197) (-0.625,-8.205) (-0.575,-8.213) (-0.525,-8.219) (-0.475,-8.225) (-0.425,-8.231) (-0.375,-8.236) (-0.325,-8.240) (-0.275,-8.244) (-0.225,-8.247) (-0.175,-8.249) (-0.125,-8.251) (-0.075,-8.252) (-0.025,-8.253) (0.050,-8.253) (0.100,-8.252) (0.150,-8.250) (0.200,-8.248) (0.250,-8.245) (0.300,-8.242) (0.350,-8.238) (0.400,-8.234) (0.450,-8.228) (0.500,-8.223) (0.550,-8.216) (0.600,-8.209) (0.650,-8.201) (0.700,-8.193) (0.750,-8.184) (0.800,-8.175) (0.850,-8.164) (0.900,-8.153) (0.950,-8.142) (1.000,-8.130) (1.050,-8.117) (1.100,-8.103) (1.150,-8.089) (1.200,-8.074) (1.250,-8.058) (1.300,-8.042) (1.350,-8.025) (1.400,-8.007) (1.450,-7.989) (1.500,-7.969) (1.548,-7.950) (1.600,-7.928) (1.650,-7.907) (1.700,-7.884) (1.750,-7.861) (1.800,-7.837) (1.850,-7.812) (1.900,-7.786) (1.950,-7.759) (2.000,-7.731) (2.050,-7.702) (2.096,-7.675) (2.150,-7.642) (2.200,-7.610) (2.250,-7.577) (2.300,-7.543) (2.350,-7.508) (2.395,-7.475) (2.450,-7.434) (2.493,-7.400) (2.550,-7.354) (2.600,-7.313) (2.644,-7.275) (2.700,-7.225) (2.750,-7.178) (2.800,-7.130) (2.850,-7.080) (2.900,-7.028) (2.949,-6.975) (2.993,-6.925) (3.036,-6.875) (3.075,-6.828) (3.117,-6.775) (3.150,-6.733) (3.193,-6.675) (3.225,-6.630) (3.263,-6.575) (3.296,-6.525) (3.325,-6.480) (3.359,-6.425) (3.389,-6.375) (3.418,-6.325) (3.445,-6.275) (3.472,-6.225) (3.497,-6.175) (3.521,-6.125) (3.545,-6.075) (3.567,-6.025) (3.589,-5.975) (3.609,-5.925) (3.629,-5.875) (3.647,-5.825) (3.665,-5.775) (3.682,-5.725) (3.698,-5.675) (3.713,-5.625) (3.728,-5.575) (3.741,-5.525) (3.754,-5.475) (3.765,-5.425) (3.776,-5.375) (3.787,-5.325) (3.796,-5.275) (3.804,-5.225) (3.812,-5.175) (3.819,-5.125) (3.825,-5.075) (3.830,-5.025) (3.835,-4.975) (3.839,-4.925) (3.842,-4.875) (3.844,-4.825) (3.845,-4.775) (3.846,-4.725) (3.846,-4.675) (3.845,-4.625) (3.843,-4.575) (3.841,-4.525) (3.837,-4.475) (3.833,-4.425) (3.828,-4.375) (3.823,-4.325) (3.816,-4.275) (3.809,-4.225) (3.801,-4.175) (3.792,-4.125) (3.783,-4.075) (3.772,-4.025) (3.761,-3.975) (3.749,-3.925) (3.736,-3.875) (3.722,-3.825) (3.708,-3.775) (3.692,-3.725) (3.676,-3.675) (3.659,-3.625) (3.641,-3.575) (3.622,-3.525) (3.602,-3.475) (3.582,-3.425) (3.560,-3.375) (3.538,-3.325) (3.514,-3.275) (3.490,-3.225) (3.465,-3.175) (3.439,-3.125) (3.412,-3.075) (3.384,-3.025) (3.355,-2.975) (3.325,-2.926) (3.294,-2.875) (3.262,-2.825) (3.229,-2.775) (3.200,-2.733) (3.160,-2.675) (3.125,-2.627) (3.087,-2.575) (3.050,-2.527) (3.009,-2.475) (2.975,-2.432) (2.928,-2.375) (2.886,-2.325) (2.850,-2.284) (2.800,-2.228) (2.752,-2.175) (2.705,-2.125) (2.657,-2.075) (2.609,-2.025) (2.559,-1.975) (2.508,-1.925) (2.456,-1.875) (2.403,-1.825) (2.350,-1.777) (2.300,-1.731) (2.250,-1.687) (2.209,-1.650) (2.151,-1.600) (2.100,-1.556) (2.050,-1.514) (2.003,-1.475) (1.950,-1.431) (1.900,-1.391) (1.850,-1.350) (1.800,-1.311) (1.755,-1.275) (1.700,-1.232) (1.659,-1.200) (1.600,-1.155) (1.550,-1.117) (1.500,-1.079) (1.450,-1.041) (1.400,-1.004) (1.350,-0.966) (1.300,-0.929) (1.250,-0.893) (1.200,-0.856) (1.158,-0.825) (1.100,-0.783) (1.055,-0.750) (1.000,-0.711) (0.951,-0.675) (0.900,-0.639) (0.850,-0.603) (0.800,-0.567) (0.750,-0.532) (0.706,-0.500) (0.650,-0.461) (0.600,-0.425) (0.550,-0.390) (0.500,-0.354) (0.450,-0.319) (0.400,-0.284) (0.354,-0.250) (0.300,-0.213) (0.250,-0.177) (0.206,-0.144)}
\EQ{(0.206,0.144) (0.248,0.175) (0.300,0.213) (0.350,0.248) (0.400,0.284) (0.450,0.319) (0.495,0.350) (0.550,0.390) (0.600,0.425) (0.650,0.461) (0.700,0.496) (0.742,0.525) (0.800,0.567) (0.846,0.600) (0.900,0.639) (0.950,0.675) (1.000,0.711) (1.050,0.747) (1.100,0.783) (1.150,0.819) (1.192,0.850) (1.250,0.893) (1.294,0.925) (1.350,0.966) (1.396,1.000) (1.450,1.041) (1.495,1.075) (1.550,1.117) (1.594,1.150) (1.650,1.193) (1.691,1.225) (1.750,1.271) (1.800,1.311) (1.850,1.350) (1.900,1.391) (1.942,1.425) (2.000,1.472) (2.050,1.514) (2.093,1.550) (2.150,1.599) (2.200,1.643) (2.250,1.687) (2.293,1.725) (2.348,1.775) (2.400,1.823) (2.450,1.870) (2.500,1.918) (2.550,1.967) (2.600,2.017) (2.650,2.068) (2.700,2.120) (2.750,2.173) (2.797,2.225) (2.842,2.275) (2.875,2.313) (2.925,2.372) (2.969,2.425) (3.009,2.475) (3.049,2.525) (3.087,2.575) (3.124,2.625) (3.160,2.675) (3.195,2.725) (3.225,2.770) (3.262,2.825) (3.294,2.875) (3.325,2.925) (3.355,2.975) (3.384,3.025) (3.412,3.075) (3.439,3.125) (3.465,3.175) (3.490,3.225) (3.514,3.275) (3.538,3.325) (3.560,3.375) (3.582,3.425) (3.602,3.475) (3.622,3.525) (3.641,3.575) (3.659,3.625) (3.675,3.673) (3.692,3.725) (3.708,3.775) (3.722,3.825) (3.736,3.875) (3.749,3.925) (3.761,3.975) (3.772,4.025) (3.783,4.075) (3.792,4.125) (3.801,4.175) (3.809,4.225) (3.816,4.275) (3.823,4.325) (3.828,4.375) (3.833,4.425) (3.837,4.475) (3.841,4.525) (3.843,4.575) (3.845,4.625) (3.846,4.675) (3.846,4.725) (3.845,4.775) (3.844,4.825) (3.842,4.875) (3.839,4.925) (3.835,4.975) (3.830,5.025) (3.825,5.075) (3.819,5.125) (3.812,5.175) (3.804,5.225) (3.796,5.275) (3.787,5.325) (3.776,5.375) (3.765,5.425) (3.754,5.475) (3.741,5.525) (3.728,5.575) (3.713,5.625) (3.698,5.675) (3.682,5.725) (3.665,5.775) (3.647,5.825) (3.629,5.875) (3.609,5.925) (3.589,5.975) (3.567,6.025) (3.545,6.075) (3.521,6.125) (3.497,6.175) (3.472,6.225) (3.445,6.275) (3.418,6.325) (3.389,6.375) (3.359,6.425) (3.328,6.475) (3.300,6.519) (3.263,6.575) (3.229,6.625) (3.193,6.675) (3.156,6.725) (3.125,6.765) (3.077,6.825) (3.036,6.875) (3.000,6.917) (2.950,6.974) (2.903,7.025) (2.855,7.075) (2.805,7.125) (2.753,7.175) (2.700,7.225) (2.650,7.269) (2.600,7.313) (2.555,7.350) (2.500,7.395) (2.450,7.434) (2.400,7.471) (2.350,7.508) (2.300,7.543) (2.253,7.575) (2.200,7.610) (2.150,7.642) (2.100,7.673) (2.054,7.700) (2.000,7.731) (1.950,7.759) (1.900,7.786) (1.850,7.812) (1.800,7.837) (1.750,7.861) (1.700,7.884) (1.650,7.907) (1.600,7.928) (1.550,7.949) (1.500,7.969) (1.450,7.989) (1.400,8.007) (1.350,8.025) (1.300,8.042) (1.250,8.058) (1.200,8.074) (1.150,8.089) (1.100,8.103) (1.050,8.117) (1.000,8.130) (0.950,8.142) (0.900,8.153) (0.850,8.164) (0.800,8.175) (0.750,8.184) (0.700,8.193) (0.650,8.201) (0.600,8.209) (0.550,8.216) (0.500,8.223) (0.450,8.228) (0.400,8.234) (0.350,8.238) (0.300,8.242) (0.250,8.245) (0.200,8.248) (0.150,8.250) (0.100,8.252) (0.050,8.253) (0.000,8.253) (-0.050,8.253) (-0.100,8.252) (-0.150,8.250) (-0.200,8.248) (-0.250,8.245) (-0.300,8.242) (-0.350,8.238) (-0.400,8.234) (-0.450,8.228) (-0.500,8.223) (-0.550,8.216) (-0.600,8.209) (-0.650,8.201) (-0.700,8.193) (-0.750,8.184) (-0.800,8.175) (-0.850,8.164) (-0.900,8.153) (-0.950,8.142) (-1.000,8.130) (-1.050,8.117) (-1.100,8.103) (-1.150,8.089) (-1.200,8.074) (-1.250,8.058) (-1.300,8.042) (-1.350,8.025) (-1.400,8.007) (-1.450,7.989) (-1.500,7.969) (-1.548,7.950) (-1.600,7.928) (-1.650,7.907) (-1.700,7.884) (-1.750,7.861) (-1.800,7.837) (-1.850,7.812) (-1.900,7.786) (-1.950,7.759) (-2.000,7.731) (-2.050,7.702) (-2.096,7.675) (-2.150,7.642) (-2.200,7.610) (-2.250,7.577) (-2.300,7.543) (-2.350,7.508) (-2.395,7.475) (-2.450,7.434) (-2.493,7.400) (-2.550,7.354) (-2.600,7.313) (-2.644,7.275) (-2.700,7.225) (-2.750,7.178) (-2.800,7.130) (-2.850,7.080) (-2.900,7.028) (-2.949,6.975) (-2.993,6.925) (-3.036,6.875) (-3.075,6.828) (-3.117,6.775) (-3.150,6.733) (-3.193,6.675) (-3.225,6.630) (-3.263,6.575) (-3.296,6.525) (-3.325,6.480) (-3.359,6.425) (-3.389,6.375) (-3.418,6.325) (-3.445,6.275) (-3.472,6.225) (-3.497,6.175) (-3.521,6.125) (-3.545,6.075) (-3.567,6.025) (-3.589,5.975) (-3.609,5.925) (-3.629,5.875) (-3.647,5.825) (-3.665,5.775) (-3.682,5.725) (-3.698,5.675) (-3.713,5.625) (-3.728,5.575) (-3.741,5.525) (-3.754,5.475) (-3.765,5.425) (-3.776,5.375) (-3.787,5.325) (-3.796,5.275) (-3.804,5.225) (-3.812,5.175) (-3.819,5.125) (-3.825,5.075) (-3.830,5.025) (-3.835,4.975) (-3.839,4.925) (-3.842,4.875) (-3.844,4.825) (-3.845,4.775) (-3.846,4.725) (-3.846,4.675) (-3.845,4.625) (-3.843,4.575) (-3.841,4.525) (-3.837,4.475) (-3.833,4.425) (-3.828,4.375) (-3.823,4.325) (-3.816,4.275) (-3.809,4.225) (-3.801,4.175) (-3.792,4.125) (-3.783,4.075) (-3.772,4.025) (-3.761,3.975) (-3.749,3.925) (-3.736,3.875) (-3.722,3.825) (-3.708,3.775) (-3.692,3.725) (-3.676,3.675) (-3.659,3.625) (-3.641,3.575) (-3.622,3.525) (-3.602,3.475) (-3.582,3.425) (-3.560,3.375) (-3.538,3.325) (-3.514,3.275) (-3.490,3.225) (-3.465,3.175) (-3.439,3.125) (-3.412,3.075) (-3.384,3.025) (-3.355,2.975) (-3.325,2.926) (-3.294,2.875) (-3.262,2.825) (-3.229,2.775) (-3.200,2.733) (-3.160,2.675) (-3.125,2.627) (-3.087,2.575) (-3.050,2.527) (-3.009,2.475) (-2.975,2.432) (-2.928,2.375) (-2.886,2.325) (-2.850,2.284) (-2.800,2.228) (-2.752,2.175) (-2.705,2.125) (-2.657,2.075) (-2.609,2.025) (-2.559,1.975) (-2.508,1.925) (-2.456,1.875) (-2.403,1.825) (-2.350,1.777) (-2.300,1.731) (-2.250,1.687) (-2.209,1.650) (-2.151,1.600) (-2.100,1.556) (-2.050,1.514) (-2.003,1.475) (-1.950,1.431) (-1.900,1.391) (-1.850,1.350) (-1.800,1.311) (-1.755,1.275) (-1.700,1.232) (-1.659,1.200) (-1.600,1.155) (-1.550,1.117) (-1.500,1.079) (-1.450,1.041) (-1.400,1.004) (-1.350,0.966) (-1.300,0.929) (-1.250,0.893) (-1.200,0.856) (-1.158,0.825) (-1.100,0.783) (-1.055,0.750) (-1.000,0.711) (-0.951,0.675) (-0.900,0.639) (-0.850,0.603) (-0.800,0.567) (-0.750,0.532) (-0.706,0.500) (-0.650,0.461) (-0.600,0.425) (-0.550,0.390) (-0.500,0.354) (-0.450,0.319) (-0.400,0.284) (-0.354,0.250) (-0.300,0.213) (-0.250,0.177) (-0.220,0.155)}
\EQ{(-0.220,-0.155) (-0.275,-0.195) (-0.319,-0.225) (-0.375,-0.266) (-0.424,-0.300) (-0.475,-0.337) (-0.525,-0.372) (-0.575,-0.407) (-0.625,-0.443) (-0.671,-0.475) (-0.725,-0.514) (-0.775,-0.550) (-0.825,-0.586) (-0.875,-0.622) (-0.925,-0.658) (-0.975,-0.694) (-1.019,-0.725) (-1.075,-0.767) (-1.121,-0.800) (-1.175,-0.840) (-1.223,-0.875) (-1.275,-0.914) (-1.323,-0.950) (-1.375,-0.989) (-1.422,-1.025) (-1.475,-1.066) (-1.520,-1.100) (-1.575,-1.143) (-1.616,-1.175) (-1.675,-1.222) (-1.725,-1.263) (-1.771,-1.300) (-1.825,-1.344) (-1.875,-1.386) (-1.921,-1.425) (-1.975,-1.472) (-2.025,-1.515) (-2.064,-1.550) (-2.120,-1.600) (-2.174,-1.650) (-2.225,-1.698) (-2.275,-1.746) (-2.325,-1.795) (-2.375,-1.845) (-2.425,-1.896) (-2.475,-1.949) (-2.522,-2.000) (-2.567,-2.050) (-2.611,-2.100) (-2.650,-2.146) (-2.695,-2.200) (-2.735,-2.250) (-2.774,-2.300) (-2.812,-2.350) (-2.849,-2.400) (-2.885,-2.450) (-2.919,-2.500) (-2.950,-2.546) (-2.985,-2.600) (-3.016,-2.650) (-3.046,-2.700) (-3.075,-2.750) (-3.100,-2.794) (-3.125,-2.841) (-3.150,-2.889) (-3.175,-2.939) (-3.200,-2.991) (-3.225,-3.045) (-3.249,-3.100) (-3.270,-3.150) (-3.290,-3.200) (-3.309,-3.250) (-3.326,-3.300) (-3.343,-3.350) (-3.359,-3.400) (-3.374,-3.450) (-3.389,-3.500) (-3.402,-3.550) (-3.414,-3.600) (-3.426,-3.650) (-3.436,-3.700) (-3.446,-3.750) (-3.455,-3.800) (-3.462,-3.850) (-3.469,-3.900) (-3.476,-3.950) (-3.481,-4.000) (-3.485,-4.050) (-3.489,-4.100) (-3.491,-4.150) (-3.493,-4.200) (-3.494,-4.250) (-3.494,-4.300) (-3.493,-4.350) (-3.492,-4.400) (-3.489,-4.450) (-3.486,-4.500) (-3.482,-4.550) (-3.477,-4.600) (-3.471,-4.650) (-3.464,-4.700) (-3.456,-4.750) (-3.447,-4.800) (-3.438,-4.850) (-3.428,-4.900) (-3.416,-4.950) (-3.404,-5.000) (-3.391,-5.050) (-3.377,-5.100) (-3.362,-5.150) (-3.346,-5.200) (-3.329,-5.250) (-3.311,-5.300) (-3.292,-5.350) (-3.272,-5.400) (-3.251,-5.450) (-3.229,-5.500) (-3.206,-5.550) (-3.182,-5.600) (-3.157,-5.650) (-3.131,-5.700) (-3.103,-5.750) (-3.075,-5.799) (-3.044,-5.850) (-3.013,-5.900) (-2.981,-5.950) (-2.950,-5.995) (-2.912,-6.050) (-2.875,-6.100) (-2.837,-6.150) (-2.800,-6.196) (-2.755,-6.250) (-2.712,-6.300) (-2.675,-6.342) (-2.625,-6.396) (-2.575,-6.447) (-2.525,-6.496) (-2.475,-6.544) (-2.425,-6.589) (-2.384,-6.625) (-2.325,-6.675) (-2.275,-6.716) (-2.231,-6.750) (-2.175,-6.792) (-2.130,-6.825) (-2.075,-6.863) (-2.025,-6.897) (-1.982,-6.925) (-1.925,-6.961) (-1.875,-6.991) (-1.825,-7.020) (-1.775,-7.048) (-1.725,-7.074) (-1.676,-7.100) (-1.625,-7.125) (-1.575,-7.149) (-1.525,-7.172) (-1.475,-7.194) (-1.425,-7.215) (-1.375,-7.236) (-1.325,-7.255) (-1.275,-7.274) (-1.225,-7.292) (-1.175,-7.309) (-1.125,-7.325) (-1.075,-7.340) (-1.025,-7.355) (-0.975,-7.369) (-0.925,-7.382) (-0.875,-7.395) (-0.825,-7.406) (-0.775,-7.417) (-0.725,-7.427) (-0.675,-7.437) (-0.625,-7.446) (-0.575,-7.454) (-0.525,-7.461) (-0.475,-7.468) (-0.425,-7.474) (-0.375,-7.479) (-0.325,-7.484) (-0.275,-7.488) (-0.225,-7.492) (-0.175,-7.494) (-0.125,-7.496) (-0.075,-7.498) (-0.025,-7.498) (0.050,-7.498) (0.100,-7.497) (0.150,-7.495) (0.200,-7.493) (0.250,-7.490) (0.300,-7.486) (0.350,-7.482) (0.400,-7.477) (0.450,-7.471) (0.500,-7.465) (0.550,-7.458) (0.599,-7.450) (0.650,-7.441) (0.700,-7.432) (0.750,-7.422) (0.800,-7.412) (0.850,-7.400) (0.900,-7.388) (0.950,-7.376) (1.000,-7.362) (1.050,-7.348) (1.100,-7.333) (1.150,-7.317) (1.200,-7.300) (1.250,-7.283) (1.300,-7.265) (1.350,-7.246) (1.400,-7.226) (1.450,-7.205) (1.500,-7.183) (1.550,-7.161) (1.600,-7.137) (1.650,-7.113) (1.700,-7.088) (1.750,-7.061) (1.800,-7.034) (1.850,-7.005) (1.900,-6.976) (1.950,-6.945) (2.000,-6.913) (2.050,-6.880) (2.094,-6.850) (2.150,-6.811) (2.198,-6.775) (2.250,-6.735) (2.295,-6.700) (2.350,-6.654) (2.400,-6.611) (2.441,-6.575) (2.495,-6.525) (2.547,-6.475) (2.597,-6.425) (2.644,-6.375) (2.690,-6.325) (2.725,-6.286) (2.775,-6.227) (2.817,-6.175) (2.850,-6.133) (2.893,-6.075) (2.925,-6.031) (2.964,-5.975) (2.997,-5.925) (3.025,-5.881) (3.060,-5.825) (3.089,-5.775) (3.117,-5.725) (3.144,-5.675) (3.170,-5.625) (3.195,-5.575) (3.218,-5.525) (3.241,-5.475) (3.262,-5.425) (3.282,-5.375) (3.302,-5.325) (3.320,-5.275) (3.338,-5.225) (3.354,-5.175) (3.370,-5.125) (3.384,-5.075) (3.398,-5.025) (3.410,-4.975) (3.422,-4.925) (3.433,-4.875) (3.443,-4.825) (3.452,-4.775) (3.460,-4.725) (3.467,-4.675) (3.474,-4.625) (3.479,-4.575) (3.484,-4.525) (3.488,-4.475) (3.491,-4.425) (3.493,-4.375) (3.494,-4.325) (3.494,-4.275) (3.494,-4.225) (3.492,-4.175) (3.490,-4.125) (3.487,-4.075) (3.483,-4.025) (3.478,-3.975) (3.473,-3.925) (3.466,-3.875) (3.459,-3.825) (3.450,-3.775) (3.441,-3.725) (3.431,-3.675) (3.420,-3.625) (3.408,-3.575) (3.395,-3.525) (3.382,-3.475) (3.367,-3.425) (3.352,-3.375) (3.335,-3.325) (3.318,-3.275) (3.300,-3.227) (3.280,-3.175) (3.260,-3.125) (3.238,-3.075) (3.216,-3.025) (3.193,-2.975) (3.168,-2.925) (3.143,-2.875) (3.117,-2.825) (3.089,-2.775) (3.061,-2.725) (3.031,-2.675) (3.001,-2.625) (2.969,-2.575) (2.936,-2.525) (2.902,-2.475) (2.867,-2.425) (2.831,-2.375) (2.800,-2.334) (2.755,-2.275) (2.715,-2.225) (2.675,-2.176) (2.632,-2.125) (2.589,-2.075) (2.550,-2.031) (2.500,-1.976) (2.452,-1.925) (2.404,-1.875) (2.355,-1.825) (2.305,-1.775) (2.254,-1.725) (2.201,-1.675) (2.150,-1.628) (2.100,-1.582) (2.050,-1.537) (2.008,-1.500) (1.950,-1.450) (1.900,-1.407) (1.850,-1.365) (1.802,-1.325) (1.750,-1.283) (1.710,-1.250) (1.650,-1.202) (1.600,-1.163) (1.552,-1.125) (1.500,-1.085) (1.455,-1.050) (1.400,-1.008) (1.357,-0.975) (1.300,-0.933) (1.257,-0.900) (1.200,-0.858) (1.155,-0.825) (1.100,-0.785) (1.053,-0.750) (1.000,-0.712) (0.950,-0.676) (0.900,-0.640) (0.850,-0.604) (0.800,-0.568) (0.750,-0.532) (0.706,-0.500) (0.650,-0.461) (0.600,-0.425) (0.550,-0.390) (0.500,-0.354) (0.450,-0.319) (0.400,-0.284) (0.354,-0.250) (0.300,-0.213) (0.250,-0.177) (0.206,-0.144)}
\EQ{(0.206,0.144) (0.248,0.175) (0.300,0.213) (0.350,0.248) (0.400,0.284) (0.450,0.319) (0.495,0.350) (0.550,0.390) (0.600,0.425) (0.650,0.461) (0.700,0.496) (0.741,0.525) (0.800,0.568) (0.845,0.600) (0.900,0.640) (0.949,0.675) (1.000,0.712) (1.050,0.748) (1.100,0.785) (1.150,0.821) (1.200,0.858) (1.250,0.895) (1.300,0.933) (1.350,0.970) (1.400,1.008) (1.450,1.046) (1.500,1.085) (1.550,1.124) (1.600,1.163) (1.647,1.200) (1.700,1.242) (1.741,1.275) (1.800,1.324) (1.850,1.365) (1.892,1.400) (1.950,1.450) (2.000,1.493) (2.050,1.537) (2.092,1.575) (2.147,1.625) (2.200,1.674) (2.250,1.722) (2.300,1.770) (2.350,1.820) (2.400,1.871) (2.450,1.923) (2.499,1.975) (2.545,2.025) (2.589,2.075) (2.625,2.117) (2.674,2.175) (2.715,2.225) (2.750,2.269) (2.793,2.325) (2.825,2.367) (2.867,2.425) (2.900,2.472) (2.936,2.525) (2.969,2.575) (3.000,2.624) (3.031,2.675) (3.061,2.725) (3.089,2.775) (3.117,2.825) (3.143,2.875) (3.168,2.925) (3.193,2.975) (3.216,3.025) (3.238,3.075) (3.260,3.125) (3.280,3.175) (3.299,3.225) (3.318,3.275) (3.335,3.325) (3.352,3.375) (3.367,3.425) (3.382,3.475) (3.395,3.525) (3.408,3.575) (3.420,3.625) (3.431,3.675) (3.441,3.725) (3.450,3.775) (3.459,3.825) (3.466,3.875) (3.473,3.925) (3.478,3.975) (3.483,4.025) (3.487,4.075) (3.490,4.125) (3.492,4.175) (3.494,4.225) (3.494,4.275) (3.494,4.325) (3.493,4.375) (3.491,4.425) (3.488,4.475) (3.484,4.525) (3.479,4.575) (3.474,4.625) (3.467,4.675) (3.460,4.725) (3.452,4.775) (3.443,4.825) (3.433,4.875) (3.422,4.925) (3.410,4.975) (3.398,5.025) (3.384,5.075) (3.370,5.125) (3.354,5.175) (3.338,5.225) (3.320,5.275) (3.302,5.325) (3.282,5.375) (3.262,5.425) (3.241,5.475) (3.218,5.525) (3.195,5.575) (3.170,5.625) (3.144,5.675) (3.117,5.725) (3.089,5.775) (3.060,5.825) (3.029,5.875) (3.000,5.921) (2.964,5.975) (2.929,6.025) (2.900,6.066) (2.856,6.125) (2.825,6.165) (2.776,6.225) (2.734,6.275) (2.700,6.314) (2.650,6.369) (2.600,6.422) (2.550,6.472) (2.500,6.520) (2.450,6.567) (2.400,6.611) (2.355,6.650) (2.300,6.696) (2.250,6.735) (2.200,6.774) (2.150,6.811) (2.100,6.846) (2.058,6.875) (2.000,6.913) (1.950,6.945) (1.901,6.975) (1.850,7.005) (1.800,7.034) (1.750,7.061) (1.700,7.088) (1.650,7.113) (1.600,7.137) (1.550,7.161) (1.500,7.183) (1.450,7.205) (1.402,7.225) (1.350,7.246) (1.300,7.265) (1.250,7.283) (1.201,7.300) (1.150,7.317) (1.100,7.333) (1.050,7.348) (1.000,7.362) (0.950,7.376) (0.900,7.388) (0.850,7.400) (0.800,7.412) (0.750,7.422) (0.700,7.432) (0.650,7.441) (0.600,7.450) (0.550,7.458) (0.500,7.465) (0.450,7.471) (0.400,7.477) (0.350,7.482) (0.300,7.486) (0.250,7.490) (0.200,7.493) (0.150,7.495) (0.100,7.497) (0.050,7.498) (0.000,7.498) (-0.050,7.498) (-0.100,7.497) (-0.150,7.495) (-0.200,7.493) (-0.250,7.490) (-0.300,7.486) (-0.350,7.482) (-0.400,7.477) (-0.450,7.471) (-0.500,7.465) (-0.550,7.458) (-0.599,7.450) (-0.650,7.441) (-0.700,7.432) (-0.750,7.422) (-0.800,7.412) (-0.850,7.400) (-0.900,7.388) (-0.950,7.376) (-1.000,7.362) (-1.050,7.348) (-1.100,7.333) (-1.150,7.317) (-1.200,7.300) (-1.250,7.283) (-1.300,7.265) (-1.350,7.246) (-1.400,7.226) (-1.450,7.205) (-1.500,7.183) (-1.550,7.161) (-1.600,7.137) (-1.650,7.113) (-1.700,7.088) (-1.750,7.061) (-1.800,7.034) (-1.850,7.005) (-1.900,6.976) (-1.950,6.945) (-2.000,6.913) (-2.050,6.880) (-2.094,6.850) (-2.150,6.811) (-2.198,6.775) (-2.250,6.735) (-2.295,6.700) (-2.350,6.654) (-2.400,6.611) (-2.441,6.575) (-2.495,6.525) (-2.547,6.475) (-2.597,6.425) (-2.644,6.375) (-2.690,6.325) (-2.725,6.286) (-2.775,6.227) (-2.817,6.175) (-2.850,6.133) (-2.893,6.075) (-2.925,6.031) (-2.964,5.975) (-2.997,5.925) (-3.025,5.881) (-3.060,5.825) (-3.089,5.775) (-3.117,5.725) (-3.144,5.675) (-3.170,5.625) (-3.195,5.575) (-3.218,5.525) (-3.241,5.475) (-3.262,5.425) (-3.282,5.375) (-3.302,5.325) (-3.320,5.275) (-3.338,5.225) (-3.354,5.175) (-3.370,5.125) (-3.384,5.075) (-3.398,5.025) (-3.410,4.975) (-3.422,4.925) (-3.433,4.875) (-3.443,4.825) (-3.452,4.775) (-3.460,4.725) (-3.467,4.675) (-3.474,4.625) (-3.479,4.575) (-3.484,4.525) (-3.488,4.475) (-3.491,4.425) (-3.493,4.375) (-3.494,4.325) (-3.494,4.275) (-3.494,4.225) (-3.492,4.175) (-3.490,4.125) (-3.487,4.075) (-3.483,4.025) (-3.478,3.975) (-3.473,3.925) (-3.466,3.875) (-3.459,3.825) (-3.450,3.775) (-3.441,3.725) (-3.431,3.675) (-3.420,3.625) (-3.408,3.575) (-3.395,3.525) (-3.382,3.475) (-3.367,3.425) (-3.352,3.375) (-3.335,3.325) (-3.318,3.275) (-3.300,3.227) (-3.280,3.175) (-3.260,3.125) (-3.238,3.075) (-3.216,3.025) (-3.193,2.975) (-3.168,2.925) (-3.143,2.875) (-3.117,2.825) (-3.089,2.775) (-3.061,2.725) (-3.031,2.675) (-3.001,2.625) (-2.969,2.575) (-2.936,2.525) (-2.902,2.475) (-2.867,2.425) (-2.831,2.375) (-2.800,2.334) (-2.755,2.275) (-2.715,2.225) (-2.675,2.176) (-2.632,2.125) (-2.589,2.075) (-2.550,2.031) (-2.500,1.976) (-2.452,1.925) (-2.404,1.875) (-2.355,1.825) (-2.305,1.775) (-2.254,1.725) (-2.201,1.675) (-2.150,1.628) (-2.100,1.582) (-2.050,1.537) (-2.008,1.500) (-1.950,1.450) (-1.900,1.407) (-1.850,1.365) (-1.802,1.325) (-1.750,1.283) (-1.710,1.250) (-1.650,1.202) (-1.600,1.163) (-1.552,1.125) (-1.500,1.085) (-1.455,1.050) (-1.400,1.008) (-1.357,0.975) (-1.300,0.933) (-1.257,0.900) (-1.200,0.858) (-1.155,0.825) (-1.100,0.785) (-1.053,0.750) (-1.000,0.712) (-0.950,0.676) (-0.900,0.640) (-0.850,0.604) (-0.800,0.568) (-0.750,0.532) (-0.706,0.500) (-0.650,0.461) (-0.600,0.425) (-0.550,0.390) (-0.500,0.354) (-0.450,0.319) (-0.400,0.284) (-0.354,0.250) (-0.300,0.213) (-0.250,0.177) (-0.220,0.155)}
\EQ{(-0.220,-0.155) (-0.275,-0.195) (-0.319,-0.225) (-0.375,-0.266) (-0.424,-0.300) (-0.475,-0.337) (-0.525,-0.372) (-0.575,-0.408) (-0.625,-0.443) (-0.671,-0.475) (-0.725,-0.515) (-0.775,-0.550) (-0.825,-0.586) (-0.875,-0.622) (-0.925,-0.659) (-0.975,-0.695) (-1.017,-0.725) (-1.075,-0.768) (-1.119,-0.800) (-1.175,-0.842) (-1.219,-0.875) (-1.275,-0.917) (-1.319,-0.950) (-1.375,-0.994) (-1.416,-1.025) (-1.475,-1.071) (-1.525,-1.111) (-1.574,-1.150) (-1.625,-1.191) (-1.667,-1.225) (-1.725,-1.274) (-1.775,-1.316) (-1.815,-1.350) (-1.872,-1.400) (-1.925,-1.447) (-1.975,-1.492) (-2.025,-1.538) (-2.064,-1.575) (-2.117,-1.625) (-2.167,-1.675) (-2.217,-1.725) (-2.265,-1.775) (-2.300,-1.812) (-2.335,-1.850) (-2.375,-1.894) (-2.424,-1.950) (-2.467,-2.000) (-2.500,-2.040) (-2.548,-2.100) (-2.587,-2.150) (-2.624,-2.200) (-2.661,-2.250) (-2.696,-2.300) (-2.725,-2.343) (-2.762,-2.400) (-2.794,-2.450) (-2.824,-2.500) (-2.850,-2.545) (-2.875,-2.589) (-2.900,-2.635) (-2.925,-2.683) (-2.950,-2.733) (-2.975,-2.786) (-3.000,-2.841) (-3.025,-2.899) (-3.046,-2.950) (-3.065,-3.000) (-3.083,-3.050) (-3.100,-3.099) (-3.117,-3.150) (-3.132,-3.200) (-3.146,-3.250) (-3.159,-3.300) (-3.171,-3.350) (-3.183,-3.400) (-3.193,-3.450) (-3.202,-3.500) (-3.211,-3.550) (-3.218,-3.600) (-3.224,-3.650) (-3.230,-3.700) (-3.235,-3.750) (-3.238,-3.800) (-3.241,-3.850) (-3.243,-3.900) (-3.244,-3.950) (-3.244,-4.000) (-3.243,-4.050) (-3.241,-4.100) (-3.238,-4.150) (-3.234,-4.200) (-3.229,-4.250) (-3.224,-4.300) (-3.217,-4.350) (-3.210,-4.400) (-3.201,-4.450) (-3.192,-4.500) (-3.182,-4.550) (-3.170,-4.600) (-3.158,-4.650) (-3.145,-4.700) (-3.130,-4.750) (-3.115,-4.800) (-3.099,-4.850) (-3.081,-4.900) (-3.063,-4.950) (-3.043,-5.000) (-3.022,-5.050) (-3.001,-5.100) (-2.978,-5.150) (-2.954,-5.200) (-2.928,-5.250) (-2.902,-5.300) (-2.875,-5.348) (-2.845,-5.400) (-2.814,-5.450) (-2.782,-5.500) (-2.750,-5.549) (-2.715,-5.600) (-2.678,-5.650) (-2.640,-5.700) (-2.601,-5.750) (-2.560,-5.800) (-2.525,-5.840) (-2.475,-5.896) (-2.425,-5.949) (-2.375,-6.000) (-2.325,-6.049) (-2.275,-6.096) (-2.225,-6.140) (-2.184,-6.175) (-2.125,-6.224) (-2.075,-6.263) (-2.026,-6.300) (-1.975,-6.337) (-1.925,-6.372) (-1.875,-6.406) (-1.825,-6.438) (-1.775,-6.469) (-1.725,-6.499) (-1.679,-6.525) (-1.634,-6.550) (-1.587,-6.575) (-1.538,-6.600) (-1.487,-6.625) (-1.433,-6.650) (-1.377,-6.675) (-1.325,-6.697) (-1.275,-6.718) (-1.225,-6.737) (-1.175,-6.756) (-1.125,-6.773) (-1.075,-6.790) (-1.025,-6.806) (-0.975,-6.821) (-0.925,-6.835) (-0.875,-6.849) (-0.825,-6.861) (-0.775,-6.873) (-0.725,-6.884) (-0.675,-6.895) (-0.625,-6.904) (-0.575,-6.913) (-0.525,-6.921) (-0.475,-6.928) (-0.425,-6.935) (-0.375,-6.941) (-0.325,-6.946) (-0.275,-6.950) (-0.225,-6.954) (-0.175,-6.956) (-0.125,-6.959) (-0.075,-6.960) (-0.025,-6.961) (0.050,-6.960) (0.100,-6.959) (0.150,-6.958) (0.200,-6.955) (0.250,-6.952) (0.300,-6.948) (0.350,-6.943) (0.400,-6.938) (0.450,-6.932) (0.500,-6.925) (0.550,-6.917) (0.600,-6.909) (0.650,-6.899) (0.700,-6.890) (0.750,-6.879) (0.800,-6.867) (0.850,-6.855) (0.900,-6.842) (0.950,-6.828) (1.000,-6.814) (1.050,-6.798) (1.100,-6.782) (1.150,-6.765) (1.200,-6.747) (1.250,-6.728) (1.300,-6.708) (1.350,-6.687) (1.400,-6.665) (1.450,-6.642) (1.500,-6.619) (1.550,-6.594) (1.600,-6.568) (1.650,-6.541) (1.700,-6.513) (1.750,-6.484) (1.800,-6.454) (1.845,-6.425) (1.900,-6.389) (1.950,-6.355) (1.992,-6.325) (2.050,-6.282) (2.092,-6.250) (2.150,-6.204) (2.200,-6.162) (2.242,-6.125) (2.297,-6.075) (2.350,-6.025) (2.400,-5.975) (2.448,-5.925) (2.494,-5.875) (2.538,-5.825) (2.575,-5.782) (2.621,-5.725) (2.660,-5.675) (2.697,-5.625) (2.732,-5.575) (2.766,-5.525) (2.799,-5.475) (2.830,-5.425) (2.859,-5.375) (2.888,-5.325) (2.915,-5.275) (2.941,-5.225) (2.966,-5.175) (2.989,-5.125) (3.012,-5.075) (3.033,-5.025) (3.053,-4.975) (3.072,-4.925) (3.090,-4.875) (3.107,-4.825) (3.123,-4.775) (3.138,-4.725) (3.151,-4.675) (3.164,-4.625) (3.176,-4.575) (3.187,-4.525) (3.197,-4.475) (3.206,-4.425) (3.214,-4.375) (3.221,-4.325) (3.227,-4.275) (3.232,-4.225) (3.236,-4.175) (3.239,-4.125) (3.242,-4.075) (3.243,-4.025) (3.244,-3.975) (3.243,-3.925) (3.242,-3.875) (3.240,-3.825) (3.237,-3.775) (3.232,-3.725) (3.227,-3.675) (3.221,-3.625) (3.214,-3.575) (3.207,-3.525) (3.198,-3.475) (3.188,-3.425) (3.177,-3.375) (3.165,-3.325) (3.153,-3.275) (3.139,-3.225) (3.124,-3.175) (3.109,-3.125) (3.092,-3.075) (3.075,-3.027) (3.056,-2.975) (3.036,-2.925) (3.015,-2.875) (2.993,-2.825) (2.970,-2.775) (2.946,-2.725) (2.921,-2.675) (2.895,-2.625) (2.867,-2.575) (2.839,-2.525) (2.809,-2.475) (2.778,-2.425) (2.750,-2.381) (2.713,-2.325) (2.678,-2.275) (2.643,-2.225) (2.606,-2.175) (2.575,-2.135) (2.528,-2.075) (2.488,-2.025) (2.450,-1.980) (2.402,-1.925) (2.358,-1.875) (2.312,-1.825) (2.275,-1.785) (2.225,-1.733) (2.175,-1.683) (2.125,-1.633) (2.075,-1.585) (2.038,-1.550) (2.000,-1.515) (1.957,-1.475) (1.901,-1.425) (1.850,-1.381) (1.800,-1.337) (1.756,-1.300) (1.700,-1.253) (1.650,-1.212) (1.605,-1.175) (1.550,-1.131) (1.500,-1.091) (1.450,-1.052) (1.400,-1.013) (1.351,-0.975) (1.300,-0.936) (1.252,-0.900) (1.200,-0.861) (1.152,-0.825) (1.100,-0.787) (1.051,-0.750) (1.000,-0.713) (0.950,-0.677) (0.900,-0.640) (0.850,-0.604) (0.800,-0.568) (0.750,-0.532) (0.705,-0.500) (0.650,-0.461) (0.600,-0.425) (0.550,-0.390) (0.500,-0.354) (0.450,-0.319) (0.400,-0.284) (0.354,-0.250) (0.300,-0.213) (0.250,-0.177) (0.206,-0.144)}
\EQ{(0.206,0.144) (0.248,0.175) (0.300,0.213) (0.350,0.248) (0.400,0.284) (0.450,0.319) (0.495,0.350) (0.550,0.390) (0.600,0.425) (0.650,0.461) (0.700,0.497) (0.750,0.532) (0.800,0.568) (0.845,0.600) (0.900,0.640) (0.948,0.675) (1.000,0.713) (1.050,0.750) (1.100,0.787) (1.150,0.823) (1.200,0.861) (1.250,0.898) (1.300,0.936) (1.350,0.974) (1.400,1.013) (1.448,1.050) (1.500,1.091) (1.543,1.125) (1.600,1.171) (1.650,1.212) (1.697,1.250) (1.750,1.295) (1.800,1.337) (1.844,1.375) (1.900,1.424) (1.950,1.469) (2.000,1.515) (2.038,1.550) (2.075,1.585) (2.117,1.625) (2.167,1.675) (2.217,1.725) (2.265,1.775) (2.300,1.812) (2.335,1.850) (2.375,1.894) (2.424,1.950) (2.467,2.000) (2.500,2.040) (2.548,2.100) (2.587,2.150) (2.624,2.200) (2.661,2.250) (2.696,2.300) (2.725,2.343) (2.762,2.400) (2.794,2.450) (2.824,2.500) (2.850,2.545) (2.875,2.589) (2.900,2.635) (2.925,2.683) (2.950,2.733) (2.975,2.786) (3.000,2.841) (3.025,2.899) (3.046,2.950) (3.065,3.000) (3.083,3.050) (3.100,3.099) (3.117,3.150) (3.132,3.200) (3.146,3.250) (3.159,3.300) (3.171,3.350) (3.183,3.400) (3.193,3.450) (3.202,3.500) (3.211,3.550) (3.218,3.600) (3.224,3.650) (3.230,3.700) (3.235,3.750) (3.238,3.800) (3.241,3.850) (3.243,3.900) (3.244,3.950) (3.244,4.000) (3.243,4.050) (3.241,4.100) (3.238,4.150) (3.234,4.200) (3.229,4.250) (3.224,4.300) (3.217,4.350) (3.210,4.400) (3.201,4.450) (3.192,4.500) (3.182,4.550) (3.170,4.600) (3.158,4.650) (3.145,4.700) (3.130,4.750) (3.115,4.800) (3.099,4.850) (3.081,4.900) (3.063,4.950) (3.043,5.000) (3.022,5.050) (3.001,5.100) (2.978,5.150) (2.954,5.200) (2.928,5.250) (2.902,5.300) (2.875,5.348) (2.845,5.400) (2.814,5.450) (2.782,5.500) (2.750,5.549) (2.715,5.600) (2.678,5.650) (2.640,5.700) (2.601,5.750) (2.560,5.800) (2.525,5.840) (2.475,5.896) (2.425,5.949) (2.375,6.000) (2.325,6.049) (2.275,6.096) (2.225,6.140) (2.184,6.175) (2.125,6.224) (2.075,6.263) (2.026,6.300) (1.975,6.337) (1.925,6.372) (1.875,6.406) (1.825,6.438) (1.775,6.469) (1.725,6.499) (1.679,6.525) (1.634,6.550) (1.587,6.575) (1.538,6.600) (1.487,6.625) (1.433,6.650) (1.377,6.675) (1.325,6.697) (1.275,6.718) (1.225,6.737) (1.175,6.756) (1.125,6.773) (1.075,6.790) (1.025,6.806) (0.975,6.821) (0.925,6.835) (0.875,6.849) (0.825,6.861) (0.775,6.873) (0.725,6.884) (0.675,6.895) (0.625,6.904) (0.575,6.913) (0.525,6.921) (0.475,6.928) (0.425,6.935) (0.375,6.941) (0.325,6.946) (0.275,6.950) (0.225,6.954) (0.175,6.956) (0.125,6.959) (0.075,6.960) (0.025,6.961) (-0.050,6.960) (-0.100,6.959) (-0.150,6.958) (-0.200,6.955) (-0.250,6.952) (-0.300,6.948) (-0.350,6.943) (-0.400,6.938) (-0.450,6.932) (-0.500,6.925) (-0.550,6.917) (-0.600,6.909) (-0.650,6.899) (-0.700,6.890) (-0.750,6.879) (-0.800,6.867) (-0.850,6.855) (-0.900,6.842) (-0.950,6.828) (-1.000,6.814) (-1.050,6.798) (-1.100,6.782) (-1.150,6.765) (-1.200,6.747) (-1.250,6.728) (-1.300,6.708) (-1.350,6.687) (-1.400,6.665) (-1.450,6.642) (-1.500,6.619) (-1.550,6.594) (-1.600,6.568) (-1.650,6.541) (-1.700,6.513) (-1.750,6.484) (-1.800,6.454) (-1.845,6.425) (-1.900,6.389) (-1.950,6.355) (-1.992,6.325) (-2.050,6.282) (-2.092,6.250) (-2.150,6.204) (-2.200,6.162) (-2.242,6.125) (-2.297,6.075) (-2.350,6.025) (-2.400,5.975) (-2.448,5.925) (-2.494,5.875) (-2.538,5.825) (-2.575,5.782) (-2.621,5.725) (-2.660,5.675) (-2.697,5.625) (-2.732,5.575) (-2.766,5.525) (-2.799,5.475) (-2.830,5.425) (-2.859,5.375) (-2.888,5.325) (-2.915,5.275) (-2.941,5.225) (-2.966,5.175) (-2.989,5.125) (-3.012,5.075) (-3.033,5.025) (-3.053,4.975) (-3.072,4.925) (-3.090,4.875) (-3.107,4.825) (-3.123,4.775) (-3.138,4.725) (-3.151,4.675) (-3.164,4.625) (-3.176,4.575) (-3.187,4.525) (-3.197,4.475) (-3.206,4.425) (-3.214,4.375) (-3.221,4.325) (-3.227,4.275) (-3.232,4.225) (-3.236,4.175) (-3.239,4.125) (-3.242,4.075) (-3.243,4.025) (-3.244,3.975) (-3.243,3.925) (-3.242,3.875) (-3.240,3.825) (-3.237,3.775) (-3.232,3.725) (-3.227,3.675) (-3.221,3.625) (-3.214,3.575) (-3.207,3.525) (-3.198,3.475) (-3.188,3.425) (-3.177,3.375) (-3.165,3.325) (-3.153,3.275) (-3.139,3.225) (-3.124,3.175) (-3.109,3.125) (-3.092,3.075) (-3.075,3.027) (-3.056,2.975) (-3.036,2.925) (-3.015,2.875) (-2.993,2.825) (-2.970,2.775) (-2.946,2.725) (-2.921,2.675) (-2.895,2.625) (-2.867,2.575) (-2.839,2.525) (-2.809,2.475) (-2.778,2.425) (-2.750,2.381) (-2.713,2.325) (-2.678,2.275) (-2.643,2.225) (-2.606,2.175) (-2.575,2.135) (-2.528,2.075) (-2.488,2.025) (-2.450,1.980) (-2.402,1.925) (-2.358,1.875) (-2.312,1.825) (-2.275,1.785) (-2.225,1.733) (-2.175,1.683) (-2.125,1.633) (-2.075,1.585) (-2.038,1.550) (-2.000,1.515) (-1.957,1.475) (-1.901,1.425) (-1.850,1.381) (-1.800,1.337) (-1.756,1.300) (-1.700,1.253) (-1.650,1.212) (-1.605,1.175) (-1.550,1.131) (-1.500,1.091) (-1.450,1.052) (-1.400,1.013) (-1.351,0.975) (-1.300,0.936) (-1.252,0.900) (-1.200,0.861) (-1.152,0.825) (-1.100,0.787) (-1.051,0.750) (-1.000,0.713) (-0.950,0.677) (-0.900,0.640) (-0.850,0.604) (-0.800,0.568) (-0.750,0.532) (-0.705,0.500) (-0.650,0.461) (-0.600,0.425) (-0.550,0.390) (-0.500,0.354) (-0.450,0.319) (-0.400,0.284) (-0.354,0.250) (-0.300,0.213) (-0.250,0.177) (-0.220,0.155)}
\EQ{(-0.220,-0.155) (-0.275,-0.195) (-0.319,-0.225) (-0.375,-0.266) (-0.424,-0.300) (-0.475,-0.337) (-0.525,-0.372) (-0.575,-0.408) (-0.625,-0.443) (-0.670,-0.475) (-0.725,-0.515) (-0.774,-0.550) (-0.825,-0.587) (-0.875,-0.623) (-0.925,-0.659) (-0.975,-0.696) (-1.025,-0.733) (-1.075,-0.770) (-1.116,-0.800) (-1.175,-0.844) (-1.216,-0.875) (-1.275,-0.920) (-1.325,-0.959) (-1.375,-0.998) (-1.425,-1.037) (-1.472,-1.075) (-1.525,-1.118) (-1.565,-1.150) (-1.625,-1.200) (-1.675,-1.242) (-1.713,-1.275) (-1.771,-1.325) (-1.825,-1.374) (-1.875,-1.419) (-1.925,-1.466) (-1.962,-1.500) (-2.000,-1.537) (-2.038,-1.575) (-2.075,-1.612) (-2.112,-1.650) (-2.150,-1.690) (-2.200,-1.744) (-2.250,-1.799) (-2.294,-1.850) (-2.336,-1.900) (-2.375,-1.948) (-2.417,-2.000) (-2.450,-2.044) (-2.491,-2.100) (-2.525,-2.147) (-2.561,-2.200) (-2.594,-2.250) (-2.625,-2.299) (-2.656,-2.350) (-2.685,-2.400) (-2.713,-2.450) (-2.740,-2.500) (-2.765,-2.550) (-2.790,-2.600) (-2.813,-2.650) (-2.835,-2.700) (-2.856,-2.750) (-2.875,-2.799) (-2.894,-2.850) (-2.912,-2.900) (-2.928,-2.950) (-2.944,-3.000) (-2.958,-3.050) (-2.972,-3.100) (-2.984,-3.150) (-2.995,-3.200) (-3.005,-3.250) (-3.015,-3.300) (-3.023,-3.350) (-3.030,-3.400) (-3.036,-3.450) (-3.041,-3.500) (-3.045,-3.550) (-3.049,-3.600) (-3.051,-3.650) (-3.052,-3.700) (-3.052,-3.750) (-3.052,-3.800) (-3.050,-3.850) (-3.047,-3.900) (-3.044,-3.950) (-3.039,-4.000) (-3.033,-4.050) (-3.027,-4.100) (-3.019,-4.150) (-3.010,-4.200) (-3.001,-4.250) (-2.990,-4.300) (-2.978,-4.350) (-2.965,-4.400) (-2.951,-4.450) (-2.937,-4.500) (-2.921,-4.550) (-2.903,-4.600) (-2.885,-4.650) (-2.866,-4.700) (-2.845,-4.750) (-2.825,-4.796) (-2.800,-4.850) (-2.776,-4.900) (-2.751,-4.950) (-2.725,-4.998) (-2.700,-5.042) (-2.666,-5.100) (-2.635,-5.150) (-2.602,-5.200) (-2.568,-5.250) (-2.533,-5.300) (-2.500,-5.344) (-2.456,-5.400) (-2.415,-5.450) (-2.375,-5.497) (-2.328,-5.550) (-2.281,-5.600) (-2.231,-5.650) (-2.180,-5.700) (-2.125,-5.750) (-2.075,-5.794) (-2.025,-5.836) (-1.977,-5.875) (-1.925,-5.915) (-1.878,-5.950) (-1.825,-5.987) (-1.775,-6.021) (-1.731,-6.050) (-1.675,-6.085) (-1.625,-6.115) (-1.575,-6.143) (-1.525,-6.170) (-1.475,-6.197) (-1.425,-6.222) (-1.375,-6.246) (-1.325,-6.269) (-1.275,-6.291) (-1.225,-6.311) (-1.175,-6.331) (-1.126,-6.350) (-1.075,-6.368) (-1.025,-6.385) (-0.975,-6.401) (-0.925,-6.417) (-0.875,-6.431) (-0.825,-6.445) (-0.775,-6.457) (-0.725,-6.469) (-0.675,-6.480) (-0.625,-6.490) (-0.575,-6.499) (-0.525,-6.508) (-0.475,-6.516) (-0.425,-6.523) (-0.375,-6.529) (-0.325,-6.534) (-0.275,-6.539) (-0.225,-6.543) (-0.175,-6.546) (-0.125,-6.548) (-0.075,-6.549) (-0.025,-6.550) (0.047,-6.550) (0.100,-6.549) (0.150,-6.547) (0.200,-6.544) (0.250,-6.541) (0.300,-6.537) (0.350,-6.532) (0.400,-6.526) (0.450,-6.519) (0.500,-6.512) (0.550,-6.504) (0.600,-6.495) (0.650,-6.485) (0.700,-6.475) (0.750,-6.463) (0.800,-6.451) (0.850,-6.438) (0.900,-6.424) (0.950,-6.409) (1.000,-6.393) (1.050,-6.377) (1.100,-6.359) (1.150,-6.341) (1.200,-6.322) (1.250,-6.301) (1.300,-6.280) (1.350,-6.257) (1.400,-6.234) (1.450,-6.209) (1.500,-6.184) (1.550,-6.157) (1.600,-6.129) (1.650,-6.100) (1.700,-6.069) (1.750,-6.038) (1.800,-6.004) (1.842,-5.975) (1.900,-5.934) (1.945,-5.900) (2.000,-5.857) (2.050,-5.816) (2.097,-5.775) (2.150,-5.728) (2.200,-5.681) (2.250,-5.632) (2.300,-5.580) (2.350,-5.525) (2.394,-5.475) (2.436,-5.425) (2.475,-5.376) (2.514,-5.325) (2.550,-5.276) (2.586,-5.225) (2.619,-5.175) (2.650,-5.126) (2.681,-5.075) (2.710,-5.025) (2.737,-4.975) (2.763,-4.925) (2.788,-4.875) (2.812,-4.825) (2.834,-4.775) (2.856,-4.725) (2.875,-4.676) (2.894,-4.625) (2.912,-4.575) (2.929,-4.525) (2.944,-4.475) (2.959,-4.425) (2.972,-4.375) (2.984,-4.325) (2.995,-4.275) (3.006,-4.225) (3.015,-4.175) (3.023,-4.125) (3.030,-4.075) (3.036,-4.025) (3.041,-3.975) (3.046,-3.925) (3.049,-3.875) (3.051,-3.825) (3.052,-3.775) (3.052,-3.725) (3.052,-3.675) (3.050,-3.625) (3.047,-3.575) (3.044,-3.525) (3.039,-3.475) (3.033,-3.425) (3.026,-3.375) (3.019,-3.325) (3.010,-3.275) (3.000,-3.225) (2.990,-3.175) (2.978,-3.125) (2.965,-3.075) (2.951,-3.025) (2.936,-2.975) (2.920,-2.925) (2.903,-2.875) (2.885,-2.825) (2.866,-2.775) (2.845,-2.725) (2.825,-2.677) (2.801,-2.625) (2.778,-2.575) (2.753,-2.525) (2.726,-2.475) (2.700,-2.427) (2.675,-2.383) (2.641,-2.325) (2.610,-2.275) (2.578,-2.225) (2.544,-2.175) (2.509,-2.125) (2.475,-2.078) (2.436,-2.025) (2.400,-1.979) (2.357,-1.925) (2.315,-1.875) (2.275,-1.828) (2.228,-1.775) (2.183,-1.725) (2.136,-1.675) (2.100,-1.638) (2.063,-1.600) (2.025,-1.562) (1.987,-1.525) (1.950,-1.489) (1.909,-1.450) (1.854,-1.400) (1.800,-1.351) (1.750,-1.307) (1.700,-1.264) (1.655,-1.225) (1.600,-1.179) (1.550,-1.138) (1.504,-1.100) (1.450,-1.057) (1.410,-1.025) (1.350,-0.978) (1.300,-0.940) (1.250,-0.901) (1.200,-0.863) (1.150,-0.826) (1.100,-0.788) (1.050,-0.751) (1.000,-0.714) (0.950,-0.678) (0.900,-0.641) (0.850,-0.605) (0.809,-0.575) (0.750,-0.533) (0.705,-0.500) (0.650,-0.461) (0.600,-0.425) (0.550,-0.390) (0.500,-0.354) (0.450,-0.319) (0.400,-0.284) (0.354,-0.250) (0.300,-0.213) (0.250,-0.177) (0.206,-0.144)}
\EQ{(0.206,0.144) (0.248,0.175) (0.300,0.213) (0.350,0.248) (0.400,0.284) (0.450,0.319) (0.495,0.350) (0.550,0.390) (0.600,0.425) (0.650,0.461) (0.700,0.497) (0.750,0.533) (0.800,0.569) (0.844,0.600) (0.900,0.641) (0.947,0.675) (1.000,0.714) (1.049,0.750) (1.100,0.788) (1.149,0.825) (1.200,0.863) (1.249,0.900) (1.300,0.940) (1.346,0.975) (1.400,1.018) (1.441,1.050) (1.500,1.097) (1.550,1.138) (1.595,1.175) (1.650,1.221) (1.700,1.264) (1.742,1.300) (1.799,1.350) (1.850,1.396) (1.900,1.442) (1.950,1.489) (1.987,1.525) (2.025,1.562) (2.063,1.600) (2.100,1.638) (2.136,1.675) (2.175,1.717) (2.225,1.771) (2.273,1.825) (2.315,1.875) (2.350,1.917) (2.397,1.975) (2.436,2.025) (2.473,2.075) (2.509,2.125) (2.544,2.175) (2.575,2.221) (2.610,2.275) (2.641,2.325) (2.671,2.375) (2.699,2.425) (2.725,2.472) (2.750,2.520) (2.775,2.570) (2.800,2.622) (2.824,2.675) (2.845,2.725) (2.866,2.775) (2.885,2.825) (2.903,2.875) (2.920,2.925) (2.936,2.975) (2.951,3.025) (2.965,3.075) (2.978,3.125) (2.990,3.175) (3.000,3.225) (3.010,3.275) (3.019,3.325) (3.026,3.375) (3.033,3.425) (3.039,3.475) (3.044,3.525) (3.047,3.575) (3.050,3.625) (3.052,3.675) (3.052,3.725) (3.052,3.775) (3.051,3.825) (3.049,3.875) (3.046,3.925) (3.041,3.975) (3.036,4.025) (3.030,4.075) (3.023,4.125) (3.015,4.175) (3.006,4.225) (2.995,4.275) (2.984,4.325) (2.972,4.375) (2.959,4.425) (2.944,4.475) (2.929,4.525) (2.912,4.575) (2.894,4.625) (2.876,4.675) (2.856,4.725) (2.834,4.775) (2.812,4.825) (2.788,4.875) (2.763,4.925) (2.737,4.975) (2.710,5.025) (2.681,5.075) (2.651,5.125) (2.619,5.175) (2.586,5.225) (2.551,5.275) (2.514,5.325) (2.476,5.375) (2.436,5.425) (2.400,5.468) (2.350,5.525) (2.305,5.575) (2.256,5.625) (2.206,5.675) (2.153,5.725) (2.100,5.773) (2.050,5.816) (2.008,5.850) (1.950,5.896) (1.900,5.934) (1.850,5.970) (1.806,6.000) (1.750,6.038) (1.700,6.069) (1.650,6.100) (1.600,6.129) (1.550,6.157) (1.500,6.184) (1.450,6.209) (1.400,6.234) (1.350,6.257) (1.300,6.280) (1.253,6.300) (1.200,6.322) (1.150,6.341) (1.100,6.359) (1.050,6.377) (1.000,6.393) (0.950,6.409) (0.900,6.424) (0.850,6.438) (0.800,6.451) (0.750,6.463) (0.700,6.475) (0.650,6.485) (0.600,6.495) (0.550,6.504) (0.500,6.512) (0.450,6.519) (0.400,6.526) (0.350,6.532) (0.300,6.537) (0.250,6.541) (0.200,6.544) (0.150,6.547) (0.100,6.549) (0.050,6.550) (0.000,6.550) (-0.050,6.550) (-0.100,6.549) (-0.150,6.547) (-0.200,6.544) (-0.250,6.541) (-0.300,6.537) (-0.350,6.532) (-0.400,6.526) (-0.450,6.519) (-0.500,6.512) (-0.550,6.504) (-0.600,6.495) (-0.650,6.485) (-0.700,6.475) (-0.750,6.463) (-0.800,6.451) (-0.850,6.438) (-0.900,6.424) (-0.950,6.409) (-1.000,6.393) (-1.050,6.377) (-1.100,6.359) (-1.150,6.341) (-1.200,6.322) (-1.250,6.301) (-1.300,6.280) (-1.350,6.257) (-1.400,6.234) (-1.450,6.209) (-1.500,6.184) (-1.550,6.157) (-1.600,6.129) (-1.650,6.100) (-1.700,6.069) (-1.750,6.038) (-1.800,6.004) (-1.842,5.975) (-1.900,5.934) (-1.945,5.900) (-2.000,5.857) (-2.050,5.816) (-2.097,5.775) (-2.150,5.728) (-2.200,5.681) (-2.250,5.632) (-2.300,5.580) (-2.350,5.525) (-2.394,5.475) (-2.436,5.425) (-2.475,5.376) (-2.514,5.325) (-2.550,5.276) (-2.586,5.225) (-2.619,5.175) (-2.650,5.126) (-2.681,5.075) (-2.710,5.025) (-2.737,4.975) (-2.763,4.925) (-2.788,4.875) (-2.812,4.825) (-2.834,4.775) (-2.856,4.725) (-2.875,4.676) (-2.894,4.625) (-2.912,4.575) (-2.929,4.525) (-2.944,4.475) (-2.959,4.425) (-2.972,4.375) (-2.984,4.325) (-2.995,4.275) (-3.006,4.225) (-3.015,4.175) (-3.023,4.125) (-3.030,4.075) (-3.036,4.025) (-3.041,3.975) (-3.046,3.925) (-3.049,3.875) (-3.051,3.825) (-3.052,3.775) (-3.052,3.725) (-3.052,3.675) (-3.050,3.625) (-3.047,3.575) (-3.044,3.525) (-3.039,3.475) (-3.033,3.425) (-3.026,3.375) (-3.019,3.325) (-3.010,3.275) (-3.000,3.225) (-2.990,3.175) (-2.978,3.125) (-2.965,3.075) (-2.951,3.025) (-2.936,2.975) (-2.920,2.925) (-2.903,2.875) (-2.885,2.825) (-2.866,2.775) (-2.845,2.725) (-2.825,2.677) (-2.801,2.625) (-2.778,2.575) (-2.753,2.525) (-2.726,2.475) (-2.700,2.427) (-2.675,2.383) (-2.641,2.325) (-2.610,2.275) (-2.578,2.225) (-2.544,2.175) (-2.509,2.125) (-2.475,2.078) (-2.436,2.025) (-2.400,1.979) (-2.357,1.925) (-2.315,1.875) (-2.275,1.828) (-2.228,1.775) (-2.183,1.725) (-2.136,1.675) (-2.100,1.638) (-2.063,1.600) (-2.025,1.562) (-1.987,1.525) (-1.950,1.489) (-1.909,1.450) (-1.854,1.400) (-1.800,1.351) (-1.750,1.307) (-1.700,1.264) (-1.655,1.225) (-1.600,1.179) (-1.550,1.138) (-1.504,1.100) (-1.450,1.057) (-1.410,1.025) (-1.350,0.978) (-1.300,0.940) (-1.250,0.901) (-1.200,0.863) (-1.150,0.826) (-1.100,0.788) (-1.050,0.751) (-1.000,0.714) (-0.950,0.678) (-0.900,0.641) (-0.850,0.605) (-0.809,0.575) (-0.750,0.533) (-0.705,0.500) (-0.650,0.461) (-0.600,0.425) (-0.550,0.390) (-0.500,0.354) (-0.450,0.319) (-0.400,0.284) (-0.354,0.250) (-0.300,0.213) (-0.250,0.177) (-0.220,0.155)}
\EQ{(-0.220,-0.155) (-0.275,-0.195) (-0.319,-0.225) (-0.375,-0.266) (-0.424,-0.300) (-0.475,-0.337) (-0.525,-0.372) (-0.575,-0.408) (-0.625,-0.444) (-0.670,-0.475) (-0.725,-0.515) (-0.774,-0.550) (-0.825,-0.587) (-0.875,-0.624) (-0.925,-0.660) (-0.975,-0.697) (-1.025,-0.734) (-1.075,-0.771) (-1.125,-0.809) (-1.175,-0.847) (-1.225,-0.885) (-1.275,-0.924) (-1.325,-0.963) (-1.372,-1.000) (-1.425,-1.043) (-1.465,-1.075) (-1.525,-1.124) (-1.575,-1.166) (-1.615,-1.200) (-1.672,-1.250) (-1.725,-1.297) (-1.775,-1.342) (-1.825,-1.389) (-1.863,-1.425) (-1.900,-1.461) (-1.940,-1.500) (-1.990,-1.550) (-2.025,-1.587) (-2.061,-1.625) (-2.100,-1.668) (-2.150,-1.724) (-2.194,-1.775) (-2.236,-1.825) (-2.275,-1.874) (-2.315,-1.925) (-2.350,-1.972) (-2.388,-2.025) (-2.423,-2.075) (-2.456,-2.125) (-2.488,-2.175) (-2.519,-2.225) (-2.548,-2.275) (-2.575,-2.324) (-2.600,-2.370) (-2.625,-2.419) (-2.650,-2.470) (-2.675,-2.524) (-2.697,-2.575) (-2.718,-2.625) (-2.737,-2.675) (-2.756,-2.725) (-2.773,-2.775) (-2.789,-2.825) (-2.804,-2.875) (-2.818,-2.925) (-2.830,-2.975) (-2.842,-3.025) (-2.852,-3.075) (-2.862,-3.125) (-2.870,-3.175) (-2.877,-3.225) (-2.884,-3.275) (-2.889,-3.325) (-2.893,-3.375) (-2.896,-3.425) (-2.898,-3.475) (-2.899,-3.525) (-2.899,-3.575) (-2.898,-3.625) (-2.896,-3.675) (-2.893,-3.725) (-2.889,-3.775) (-2.884,-3.825) (-2.878,-3.875) (-2.871,-3.925) (-2.862,-3.975) (-2.853,-4.025) (-2.842,-4.075) (-2.831,-4.125) (-2.818,-4.175) (-2.804,-4.225) (-2.789,-4.275) (-2.773,-4.325) (-2.756,-4.375) (-2.738,-4.425) (-2.718,-4.475) (-2.697,-4.525) (-2.675,-4.575) (-2.652,-4.625) (-2.627,-4.675) (-2.601,-4.725) (-2.575,-4.772) (-2.550,-4.815) (-2.514,-4.875) (-2.482,-4.925) (-2.450,-4.972) (-2.413,-5.025) (-2.376,-5.075) (-2.337,-5.125) (-2.300,-5.170) (-2.253,-5.225) (-2.208,-5.275) (-2.161,-5.325) (-2.125,-5.362) (-2.086,-5.400) (-2.032,-5.450) (-1.976,-5.500) (-1.925,-5.543) (-1.885,-5.575) (-1.825,-5.621) (-1.775,-5.658) (-1.725,-5.693) (-1.678,-5.725) (-1.625,-5.759) (-1.575,-5.790) (-1.525,-5.819) (-1.475,-5.847) (-1.425,-5.874) (-1.375,-5.900) (-1.325,-5.924) (-1.275,-5.947) (-1.225,-5.970) (-1.175,-5.991) (-1.125,-6.011) (-1.075,-6.030) (-1.025,-6.048) (-0.975,-6.065) (-0.925,-6.081) (-0.875,-6.096) (-0.825,-6.111) (-0.775,-6.124) (-0.725,-6.137) (-0.675,-6.148) (-0.625,-6.159) (-0.575,-6.169) (-0.525,-6.178) (-0.475,-6.186) (-0.425,-6.193) (-0.375,-6.200) (-0.325,-6.205) (-0.275,-6.210) (-0.225,-6.214) (-0.175,-6.217) (-0.125,-6.220) (-0.075,-6.221) (-0.025,-6.222) (0.050,-6.222) (0.100,-6.221) (0.150,-6.219) (0.200,-6.216) (0.250,-6.212) (0.300,-6.208) (0.350,-6.203) (0.400,-6.196) (0.450,-6.190) (0.500,-6.182) (0.550,-6.173) (0.600,-6.164) (0.650,-6.154) (0.700,-6.142) (0.750,-6.130) (0.800,-6.117) (0.850,-6.104) (0.900,-6.089) (0.950,-6.073) (1.000,-6.057) (1.050,-6.039) (1.100,-6.020) (1.150,-6.001) (1.200,-5.980) (1.250,-5.959) (1.300,-5.936) (1.350,-5.912) (1.400,-5.887) (1.450,-5.861) (1.500,-5.833) (1.550,-5.805) (1.599,-5.775) (1.650,-5.743) (1.700,-5.710) (1.750,-5.676) (1.800,-5.640) (1.850,-5.602) (1.900,-5.563) (1.946,-5.525) (2.000,-5.479) (2.050,-5.434) (2.100,-5.386) (2.137,-5.350) (2.175,-5.311) (2.225,-5.257) (2.275,-5.200) (2.317,-5.150) (2.350,-5.108) (2.394,-5.050) (2.425,-5.008) (2.465,-4.950) (2.498,-4.900) (2.525,-4.857) (2.559,-4.800) (2.587,-4.750) (2.614,-4.700) (2.639,-4.650) (2.663,-4.600) (2.686,-4.550) (2.708,-4.500) (2.728,-4.450) (2.747,-4.400) (2.765,-4.350) (2.781,-4.300) (2.797,-4.250) (2.811,-4.200) (2.825,-4.150) (2.837,-4.100) (2.848,-4.050) (2.858,-4.000) (2.867,-3.950) (2.874,-3.900) (2.881,-3.850) (2.887,-3.800) (2.891,-3.750) (2.895,-3.700) (2.898,-3.650) (2.899,-3.600) (2.900,-3.550) (2.899,-3.500) (2.897,-3.450) (2.895,-3.400) (2.891,-3.350) (2.887,-3.300) (2.881,-3.250) (2.874,-3.200) (2.866,-3.150) (2.857,-3.100) (2.847,-3.050) (2.836,-3.000) (2.824,-2.950) (2.811,-2.900) (2.796,-2.850) (2.781,-2.800) (2.764,-2.750) (2.747,-2.700) (2.728,-2.650) (2.708,-2.600) (2.687,-2.550) (2.664,-2.500) (2.640,-2.450) (2.616,-2.400) (2.589,-2.350) (2.562,-2.300) (2.533,-2.250) (2.503,-2.200) (2.475,-2.155) (2.440,-2.100) (2.406,-2.050) (2.375,-2.007) (2.334,-1.950) (2.300,-1.906) (2.256,-1.850) (2.215,-1.800) (2.175,-1.752) (2.129,-1.700) (2.084,-1.650) (2.050,-1.613) (2.014,-1.575) (1.975,-1.535) (1.925,-1.485) (1.889,-1.450) (1.850,-1.413) (1.811,-1.375) (1.756,-1.325) (1.700,-1.275) (1.650,-1.231) (1.600,-1.188) (1.556,-1.150) (1.500,-1.104) (1.450,-1.063) (1.404,-1.025) (1.350,-0.983) (1.309,-0.950) (1.250,-0.904) (1.200,-0.866) (1.150,-0.828) (1.100,-0.790) (1.050,-0.753) (1.000,-0.716) (0.950,-0.679) (0.900,-0.642) (0.850,-0.606) (0.809,-0.575) (0.750,-0.533) (0.705,-0.500) (0.650,-0.461) (0.600,-0.425) (0.550,-0.390) (0.500,-0.354) (0.450,-0.319) (0.400,-0.284) (0.354,-0.250) (0.300,-0.213) (0.250,-0.177) (0.206,-0.144)}
\EQ{(0.206,0.144) (0.248,0.175) (0.300,0.213) (0.350,0.248) (0.400,0.284) (0.450,0.319) (0.495,0.350) (0.550,0.390) (0.600,0.425) (0.650,0.461) (0.700,0.497) (0.750,0.533) (0.800,0.569) (0.843,0.600) (0.900,0.642) (0.946,0.675) (1.000,0.716) (1.047,0.750) (1.100,0.790) (1.147,0.825) (1.200,0.866) (1.245,0.900) (1.300,0.943) (1.341,0.975) (1.400,1.022) (1.450,1.063) (1.496,1.100) (1.550,1.145) (1.600,1.188) (1.644,1.225) (1.700,1.275) (1.750,1.320) (1.800,1.366) (1.837,1.400) (1.875,1.436) (1.915,1.475) (1.965,1.525) (2.000,1.561) (2.038,1.600) (2.075,1.640) (2.125,1.695) (2.173,1.750) (2.215,1.800) (2.250,1.843) (2.296,1.900) (2.334,1.950) (2.370,2.000) (2.400,2.042) (2.440,2.100) (2.472,2.150) (2.500,2.195) (2.533,2.250) (2.562,2.300) (2.589,2.350) (2.616,2.400) (2.640,2.450) (2.664,2.500) (2.687,2.550) (2.708,2.600) (2.728,2.650) (2.747,2.700) (2.764,2.750) (2.781,2.800) (2.796,2.850) (2.811,2.900) (2.824,2.950) (2.836,3.000) (2.847,3.050) (2.857,3.100) (2.866,3.150) (2.874,3.200) (2.881,3.250) (2.887,3.300) (2.891,3.350) (2.895,3.400) (2.897,3.450) (2.899,3.500) (2.900,3.550) (2.899,3.600) (2.898,3.650) (2.895,3.700) (2.891,3.750) (2.887,3.800) (2.881,3.850) (2.874,3.900) (2.867,3.950) (2.858,4.000) (2.848,4.050) (2.837,4.100) (2.825,4.150) (2.811,4.200) (2.797,4.250) (2.781,4.300) (2.765,4.350) (2.747,4.400) (2.728,4.450) (2.708,4.500) (2.686,4.550) (2.663,4.600) (2.639,4.650) (2.614,4.700) (2.587,4.750) (2.559,4.800) (2.529,4.850) (2.500,4.897) (2.465,4.950) (2.431,5.000) (2.400,5.043) (2.357,5.100) (2.325,5.140) (2.275,5.200) (2.231,5.250) (2.185,5.300) (2.150,5.337) (2.112,5.375) (2.075,5.410) (2.032,5.450) (1.976,5.500) (1.925,5.543) (1.885,5.575) (1.825,5.621) (1.775,5.658) (1.725,5.693) (1.678,5.725) (1.625,5.759) (1.575,5.790) (1.525,5.819) (1.475,5.847) (1.425,5.874) (1.375,5.900) (1.325,5.924) (1.275,5.947) (1.225,5.970) (1.175,5.991) (1.125,6.011) (1.075,6.030) (1.025,6.048) (0.975,6.065) (0.925,6.081) (0.875,6.096) (0.825,6.111) (0.775,6.124) (0.725,6.137) (0.675,6.148) (0.625,6.159) (0.575,6.169) (0.525,6.178) (0.475,6.186) (0.425,6.193) (0.375,6.200) (0.325,6.205) (0.275,6.210) (0.225,6.214) (0.175,6.217) (0.125,6.220) (0.075,6.221) (0.025,6.222) (-0.050,6.222) (-0.100,6.221) (-0.150,6.219) (-0.200,6.216) (-0.250,6.212) (-0.300,6.208) (-0.350,6.203) (-0.400,6.196) (-0.450,6.190) (-0.500,6.182) (-0.550,6.173) (-0.600,6.164) (-0.650,6.154) (-0.700,6.142) (-0.750,6.130) (-0.800,6.117) (-0.850,6.104) (-0.900,6.089) (-0.950,6.073) (-1.000,6.057) (-1.050,6.039) (-1.100,6.020) (-1.150,6.001) (-1.200,5.980) (-1.250,5.959) (-1.300,5.936) (-1.350,5.912) (-1.400,5.887) (-1.450,5.861) (-1.500,5.833) (-1.550,5.805) (-1.599,5.775) (-1.650,5.743) (-1.700,5.710) (-1.750,5.676) (-1.800,5.640) (-1.850,5.602) (-1.900,5.563) (-1.946,5.525) (-2.000,5.479) (-2.050,5.434) (-2.100,5.386) (-2.137,5.350) (-2.175,5.311) (-2.225,5.257) (-2.275,5.200) (-2.317,5.150) (-2.350,5.108) (-2.394,5.050) (-2.425,5.008) (-2.465,4.950) (-2.498,4.900) (-2.525,4.857) (-2.559,4.800) (-2.587,4.750) (-2.614,4.700) (-2.639,4.650) (-2.663,4.600) (-2.686,4.550) (-2.708,4.500) (-2.728,4.450) (-2.747,4.400) (-2.765,4.350) (-2.781,4.300) (-2.797,4.250) (-2.811,4.200) (-2.825,4.150) (-2.837,4.100) (-2.848,4.050) (-2.858,4.000) (-2.867,3.950) (-2.874,3.900) (-2.881,3.850) (-2.887,3.800) (-2.891,3.750) (-2.895,3.700) (-2.898,3.650) (-2.899,3.600) (-2.900,3.550) (-2.899,3.500) (-2.897,3.450) (-2.895,3.400) (-2.891,3.350) (-2.887,3.300) (-2.881,3.250) (-2.874,3.200) (-2.866,3.150) (-2.857,3.100) (-2.847,3.050) (-2.836,3.000) (-2.824,2.950) (-2.811,2.900) (-2.796,2.850) (-2.781,2.800) (-2.764,2.750) (-2.747,2.700) (-2.728,2.650) (-2.708,2.600) (-2.687,2.550) (-2.664,2.500) (-2.640,2.450) (-2.616,2.400) (-2.589,2.350) (-2.562,2.300) (-2.533,2.250) (-2.503,2.200) (-2.475,2.155) (-2.440,2.100) (-2.406,2.050) (-2.375,2.007) (-2.334,1.950) (-2.300,1.906) (-2.256,1.850) (-2.215,1.800) (-2.175,1.752) (-2.129,1.700) (-2.084,1.650) (-2.050,1.613) (-2.014,1.575) (-1.975,1.535) (-1.925,1.485) (-1.889,1.450) (-1.850,1.413) (-1.811,1.375) (-1.756,1.325) (-1.700,1.275) (-1.650,1.231) (-1.600,1.188) (-1.556,1.150) (-1.500,1.104) (-1.450,1.063) (-1.404,1.025) (-1.350,0.983) (-1.309,0.950) (-1.250,0.904) (-1.200,0.866) (-1.150,0.828) (-1.100,0.790) (-1.050,0.753) (-1.000,0.716) (-0.950,0.679) (-0.900,0.642) (-0.850,0.606) (-0.809,0.575) (-0.750,0.533) (-0.705,0.500) (-0.650,0.461) (-0.600,0.425) (-0.550,0.390) (-0.500,0.354) (-0.450,0.319) (-0.400,0.284) (-0.354,0.250) (-0.300,0.213) (-0.250,0.177) (-0.220,0.155)}
\EQ{(-0.220,-0.155) (-0.275,-0.195) (-0.319,-0.225) (-0.375,-0.266) (-0.424,-0.300) (-0.475,-0.337) (-0.525,-0.372) (-0.575,-0.408) (-0.625,-0.444) (-0.670,-0.475) (-0.725,-0.516) (-0.773,-0.550) (-0.825,-0.588) (-0.875,-0.624) (-0.925,-0.661) (-0.975,-0.698) (-1.025,-0.735) (-1.075,-0.773) (-1.125,-0.811) (-1.175,-0.849) (-1.225,-0.888) (-1.273,-0.925) (-1.325,-0.967) (-1.367,-1.000) (-1.425,-1.048) (-1.475,-1.089) (-1.518,-1.125) (-1.575,-1.174) (-1.625,-1.218) (-1.675,-1.263) (-1.715,-1.300) (-1.769,-1.350) (-1.821,-1.400) (-1.871,-1.450) (-1.920,-1.500) (-1.967,-1.550) (-2.013,-1.600) (-2.050,-1.641) (-2.100,-1.699) (-2.142,-1.750) (-2.175,-1.791) (-2.221,-1.850) (-2.258,-1.900) (-2.293,-1.950) (-2.325,-1.996) (-2.361,-2.050) (-2.392,-2.100) (-2.422,-2.150) (-2.450,-2.199) (-2.475,-2.245) (-2.500,-2.292) (-2.525,-2.343) (-2.550,-2.396) (-2.574,-2.450) (-2.595,-2.500) (-2.615,-2.550) (-2.633,-2.600) (-2.650,-2.648) (-2.667,-2.700) (-2.682,-2.750) (-2.696,-2.800) (-2.708,-2.850) (-2.720,-2.900) (-2.730,-2.950) (-2.739,-3.000) (-2.747,-3.050) (-2.754,-3.100) (-2.760,-3.150) (-2.765,-3.200) (-2.769,-3.250) (-2.771,-3.300) (-2.773,-3.350) (-2.773,-3.400) (-2.773,-3.450) (-2.771,-3.500) (-2.768,-3.550) (-2.764,-3.600) (-2.759,-3.650) (-2.753,-3.700) (-2.746,-3.750) (-2.738,-3.800) (-2.729,-3.850) (-2.718,-3.900) (-2.706,-3.950) (-2.694,-4.000) (-2.680,-4.050) (-2.664,-4.100) (-2.648,-4.150) (-2.630,-4.200) (-2.612,-4.250) (-2.592,-4.300) (-2.570,-4.350) (-2.547,-4.400) (-2.525,-4.446) (-2.500,-4.496) (-2.475,-4.542) (-2.450,-4.587) (-2.412,-4.650) (-2.381,-4.700) (-2.350,-4.746) (-2.312,-4.800) (-2.275,-4.850) (-2.237,-4.900) (-2.200,-4.945) (-2.153,-5.000) (-2.108,-5.050) (-2.061,-5.100) (-2.025,-5.136) (-1.984,-5.175) (-1.930,-5.225) (-1.875,-5.273) (-1.825,-5.315) (-1.781,-5.350) (-1.725,-5.392) (-1.679,-5.425) (-1.625,-5.462) (-1.575,-5.495) (-1.527,-5.525) (-1.475,-5.556) (-1.425,-5.585) (-1.375,-5.612) (-1.325,-5.638) (-1.275,-5.662) (-1.225,-5.686) (-1.175,-5.708) (-1.125,-5.729) (-1.075,-5.749) (-1.025,-5.769) (-0.975,-5.787) (-0.925,-5.804) (-0.875,-5.819) (-0.825,-5.834) (-0.775,-5.848) (-0.725,-5.862) (-0.675,-5.874) (-0.625,-5.885) (-0.575,-5.895) (-0.525,-5.905) (-0.475,-5.913) (-0.425,-5.921) (-0.375,-5.928) (-0.325,-5.934) (-0.275,-5.939) (-0.225,-5.943) (-0.175,-5.946) (-0.125,-5.949) (-0.075,-5.950) (-0.025,-5.951) (0.050,-5.951) (0.100,-5.950) (0.150,-5.948) (0.200,-5.945) (0.250,-5.941) (0.300,-5.936) (0.350,-5.931) (0.400,-5.924) (0.450,-5.917) (0.500,-5.909) (0.550,-5.900) (0.600,-5.890) (0.650,-5.879) (0.700,-5.868) (0.750,-5.855) (0.800,-5.842) (0.850,-5.827) (0.900,-5.812) (0.950,-5.795) (1.000,-5.778) (1.050,-5.759) (1.100,-5.740) (1.150,-5.719) (1.200,-5.697) (1.249,-5.675) (1.300,-5.650) (1.350,-5.625) (1.397,-5.600) (1.442,-5.575) (1.486,-5.550) (1.550,-5.511) (1.600,-5.479) (1.643,-5.450) (1.700,-5.410) (1.748,-5.375) (1.800,-5.335) (1.843,-5.300) (1.900,-5.252) (1.950,-5.207) (2.000,-5.160) (2.036,-5.125) (2.075,-5.085) (2.125,-5.031) (2.175,-4.975) (2.217,-4.925) (2.250,-4.883) (2.294,-4.825) (2.325,-4.782) (2.364,-4.725) (2.397,-4.675) (2.425,-4.629) (2.457,-4.575) (2.484,-4.525) (2.511,-4.475) (2.535,-4.425) (2.559,-4.375) (2.581,-4.325) (2.602,-4.275) (2.621,-4.225) (2.639,-4.175) (2.656,-4.125) (2.672,-4.075) (2.687,-4.025) (2.700,-3.975) (2.712,-3.925) (2.723,-3.875) (2.733,-3.825) (2.742,-3.775) (2.750,-3.725) (2.757,-3.675) (2.762,-3.625) (2.766,-3.575) (2.770,-3.525) (2.772,-3.475) (2.773,-3.425) (2.773,-3.375) (2.772,-3.325) (2.770,-3.275) (2.767,-3.225) (2.763,-3.175) (2.757,-3.125) (2.751,-3.075) (2.743,-3.025) (2.735,-2.975) (2.725,-2.926) (2.714,-2.875) (2.702,-2.825) (2.689,-2.775) (2.674,-2.725) (2.659,-2.675) (2.642,-2.625) (2.625,-2.577) (2.605,-2.525) (2.585,-2.475) (2.563,-2.425) (2.541,-2.375) (2.517,-2.325) (2.491,-2.275) (2.464,-2.225) (2.436,-2.175) (2.407,-2.125) (2.376,-2.075) (2.344,-2.025) (2.311,-1.975) (2.276,-1.925) (2.240,-1.875) (2.202,-1.825) (2.162,-1.775) (2.125,-1.729) (2.079,-1.675) (2.036,-1.625) (2.000,-1.586) (1.950,-1.532) (1.900,-1.480) (1.850,-1.429) (1.800,-1.380) (1.750,-1.333) (1.700,-1.286) (1.661,-1.250) (1.604,-1.200) (1.550,-1.153) (1.500,-1.110) (1.458,-1.075) (1.400,-1.027) (1.350,-0.987) (1.304,-0.950) (1.250,-0.907) (1.209,-0.875) (1.150,-0.830) (1.100,-0.792) (1.050,-0.754) (1.000,-0.717) (0.950,-0.680) (0.900,-0.643) (0.850,-0.606) (0.808,-0.575) (0.750,-0.534) (0.704,-0.500) (0.650,-0.462) (0.600,-0.425) (0.550,-0.390) (0.500,-0.354) (0.450,-0.319) (0.400,-0.284) (0.354,-0.250) (0.300,-0.213) (0.250,-0.177) (0.206,-0.144)}
\EQ{(0.206,0.144) (0.248,0.175) (0.300,0.213) (0.350,0.248) (0.400,0.284) (0.450,0.319) (0.495,0.350) (0.550,0.390) (0.599,0.425) (0.650,0.462) (0.700,0.497) (0.750,0.534) (0.800,0.570) (0.842,0.600) (0.900,0.643) (0.944,0.675) (1.000,0.717) (1.045,0.750) (1.100,0.792) (1.144,0.825) (1.200,0.868) (1.241,0.900) (1.300,0.947) (1.350,0.987) (1.398,1.025) (1.450,1.068) (1.500,1.110) (1.547,1.150) (1.600,1.196) (1.650,1.241) (1.688,1.275) (1.725,1.309) (1.769,1.350) (1.821,1.400) (1.871,1.450) (1.920,1.500) (1.967,1.550) (2.013,1.600) (2.050,1.641) (2.100,1.699) (2.142,1.750) (2.175,1.791) (2.221,1.850) (2.258,1.900) (2.293,1.950) (2.325,1.996) (2.361,2.050) (2.392,2.100) (2.422,2.150) (2.450,2.199) (2.475,2.245) (2.500,2.292) (2.525,2.343) (2.550,2.396) (2.574,2.450) (2.595,2.500) (2.615,2.550) (2.633,2.600) (2.650,2.648) (2.667,2.700) (2.682,2.750) (2.696,2.800) (2.708,2.850) (2.720,2.900) (2.730,2.950) (2.739,3.000) (2.747,3.050) (2.754,3.100) (2.760,3.150) (2.765,3.200) (2.769,3.250) (2.771,3.300) (2.773,3.350) (2.773,3.400) (2.773,3.450) (2.771,3.500) (2.768,3.550) (2.764,3.600) (2.759,3.650) (2.753,3.700) (2.746,3.750) (2.738,3.800) (2.729,3.850) (2.718,3.900) (2.706,3.950) (2.694,4.000) (2.680,4.050) (2.664,4.100) (2.648,4.150) (2.630,4.200) (2.612,4.250) (2.592,4.300) (2.570,4.350) (2.547,4.400) (2.525,4.446) (2.500,4.496) (2.475,4.542) (2.450,4.587) (2.412,4.650) (2.381,4.700) (2.350,4.746) (2.312,4.800) (2.275,4.850) (2.237,4.900) (2.200,4.945) (2.153,5.000) (2.108,5.050) (2.061,5.100) (2.025,5.136) (1.984,5.175) (1.930,5.225) (1.875,5.273) (1.825,5.315) (1.781,5.350) (1.725,5.392) (1.679,5.425) (1.625,5.462) (1.575,5.495) (1.527,5.525) (1.475,5.556) (1.425,5.585) (1.375,5.612) (1.325,5.638) (1.275,5.662) (1.225,5.686) (1.175,5.708) (1.125,5.729) (1.075,5.749) (1.025,5.769) (0.975,5.787) (0.925,5.804) (0.875,5.819) (0.825,5.834) (0.775,5.848) (0.725,5.862) (0.675,5.874) (0.625,5.885) (0.575,5.895) (0.525,5.905) (0.475,5.913) (0.425,5.921) (0.375,5.928) (0.325,5.934) (0.275,5.939) (0.225,5.943) (0.175,5.946) (0.125,5.949) (0.075,5.950) (0.025,5.951) (-0.050,5.951) (-0.100,5.950) (-0.150,5.948) (-0.200,5.945) (-0.250,5.941) (-0.300,5.936) (-0.350,5.931) (-0.400,5.924) (-0.450,5.917) (-0.500,5.909) (-0.550,5.900) (-0.600,5.890) (-0.650,5.879) (-0.700,5.868) (-0.750,5.855) (-0.800,5.842) (-0.850,5.827) (-0.900,5.812) (-0.950,5.795) (-1.000,5.778) (-1.050,5.759) (-1.100,5.740) (-1.150,5.719) (-1.200,5.697) (-1.249,5.675) (-1.300,5.650) (-1.350,5.625) (-1.397,5.600) (-1.442,5.575) (-1.486,5.550) (-1.550,5.511) (-1.600,5.479) (-1.643,5.450) (-1.700,5.410) (-1.748,5.375) (-1.800,5.335) (-1.843,5.300) (-1.900,5.252) (-1.950,5.207) (-2.000,5.160) (-2.036,5.125) (-2.075,5.085) (-2.125,5.031) (-2.175,4.975) (-2.217,4.925) (-2.250,4.883) (-2.294,4.825) (-2.325,4.782) (-2.364,4.725) (-2.397,4.675) (-2.425,4.629) (-2.457,4.575) (-2.484,4.525) (-2.511,4.475) (-2.535,4.425) (-2.559,4.375) (-2.581,4.325) (-2.602,4.275) (-2.621,4.225) (-2.639,4.175) (-2.656,4.125) (-2.672,4.075) (-2.687,4.025) (-2.700,3.975) (-2.712,3.925) (-2.723,3.875) (-2.733,3.825) (-2.742,3.775) (-2.750,3.725) (-2.757,3.675) (-2.762,3.625) (-2.766,3.575) (-2.770,3.525) (-2.772,3.475) (-2.773,3.425) (-2.773,3.375) (-2.772,3.325) (-2.770,3.275) (-2.767,3.225) (-2.763,3.175) (-2.757,3.125) (-2.751,3.075) (-2.743,3.025) (-2.735,2.975) (-2.725,2.926) (-2.714,2.875) (-2.702,2.825) (-2.689,2.775) (-2.674,2.725) (-2.659,2.675) (-2.642,2.625) (-2.625,2.577) (-2.605,2.525) (-2.585,2.475) (-2.563,2.425) (-2.541,2.375) (-2.517,2.325) (-2.491,2.275) (-2.464,2.225) (-2.436,2.175) (-2.407,2.125) (-2.376,2.075) (-2.344,2.025) (-2.311,1.975) (-2.276,1.925) (-2.240,1.875) (-2.202,1.825) (-2.162,1.775) (-2.125,1.729) (-2.079,1.675) (-2.036,1.625) (-2.000,1.586) (-1.950,1.532) (-1.900,1.480) (-1.850,1.429) (-1.800,1.380) (-1.750,1.333) (-1.700,1.286) (-1.661,1.250) (-1.604,1.200) (-1.550,1.153) (-1.500,1.110) (-1.458,1.075) (-1.400,1.027) (-1.350,0.987) (-1.304,0.950) (-1.250,0.907) (-1.209,0.875) (-1.150,0.830) (-1.100,0.792) (-1.050,0.754) (-1.000,0.717) (-0.950,0.680) (-0.900,0.643) (-0.850,0.606) (-0.808,0.575) (-0.750,0.534) (-0.704,0.500) (-0.650,0.462) (-0.600,0.425) (-0.550,0.390) (-0.500,0.354) (-0.450,0.319) (-0.400,0.284) (-0.354,0.250) (-0.300,0.213) (-0.250,0.177) (-0.220,0.155)}
\EQ{(-0.220,-0.155) (-0.275,-0.195) (-0.319,-0.225) (-0.375,-0.266) (-0.424,-0.300) (-0.475,-0.337) (-0.525,-0.372) (-0.575,-0.408) (-0.625,-0.444) (-0.669,-0.475) (-0.725,-0.516) (-0.773,-0.550) (-0.825,-0.589) (-0.875,-0.625) (-0.925,-0.662) (-0.975,-0.699) (-1.025,-0.737) (-1.075,-0.774) (-1.125,-0.813) (-1.173,-0.850) (-1.225,-0.891) (-1.269,-0.925) (-1.325,-0.970) (-1.375,-1.011) (-1.422,-1.050) (-1.475,-1.095) (-1.525,-1.139) (-1.567,-1.175) (-1.622,-1.225) (-1.675,-1.274) (-1.725,-1.322) (-1.775,-1.370) (-1.825,-1.421) (-1.875,-1.473) (-1.924,-1.525) (-1.969,-1.575) (-2.012,-1.625) (-2.050,-1.671) (-2.094,-1.725) (-2.125,-1.766) (-2.170,-1.825) (-2.200,-1.868) (-2.239,-1.925) (-2.272,-1.975) (-2.300,-2.020) (-2.333,-2.075) (-2.362,-2.125) (-2.389,-2.175) (-2.414,-2.225) (-2.439,-2.275) (-2.462,-2.325) (-2.483,-2.375) (-2.504,-2.425) (-2.523,-2.475) (-2.541,-2.525) (-2.558,-2.575) (-2.573,-2.625) (-2.587,-2.675) (-2.600,-2.725) (-2.612,-2.775) (-2.623,-2.825) (-2.632,-2.875) (-2.640,-2.925) (-2.647,-2.975) (-2.654,-3.025) (-2.658,-3.075) (-2.662,-3.125) (-2.665,-3.175) (-2.666,-3.225) (-2.667,-3.275) (-2.666,-3.325) (-2.664,-3.375) (-2.661,-3.425) (-2.657,-3.475) (-2.651,-3.525) (-2.645,-3.575) (-2.637,-3.625) (-2.629,-3.675) (-2.619,-3.725) (-2.608,-3.775) (-2.595,-3.825) (-2.582,-3.875) (-2.567,-3.925) (-2.551,-3.975) (-2.534,-4.025) (-2.515,-4.075) (-2.496,-4.125) (-2.475,-4.174) (-2.452,-4.225) (-2.428,-4.275) (-2.402,-4.325) (-2.375,-4.375) (-2.350,-4.420) (-2.317,-4.475) (-2.285,-4.525) (-2.251,-4.575) (-2.216,-4.625) (-2.179,-4.675) (-2.140,-4.725) (-2.100,-4.773) (-2.055,-4.825) (-2.009,-4.875) (-1.960,-4.925) (-1.909,-4.975) (-1.855,-5.025) (-1.800,-5.073) (-1.750,-5.114) (-1.704,-5.150) (-1.650,-5.191) (-1.602,-5.225) (-1.550,-5.260) (-1.500,-5.293) (-1.450,-5.323) (-1.405,-5.350) (-1.361,-5.375) (-1.314,-5.400) (-1.265,-5.425) (-1.214,-5.450) (-1.160,-5.475) (-1.103,-5.500) (-1.050,-5.522) (-1.000,-5.541) (-0.950,-5.559) (-0.900,-5.577) (-0.850,-5.593) (-0.800,-5.608) (-0.750,-5.622) (-0.700,-5.635) (-0.650,-5.647) (-0.600,-5.659) (-0.550,-5.669) (-0.500,-5.678) (-0.450,-5.687) (-0.400,-5.694) (-0.350,-5.701) (-0.300,-5.707) (-0.250,-5.711) (-0.200,-5.715) (-0.150,-5.718) (-0.100,-5.721) (-0.050,-5.722) (0.000,-5.722) (0.050,-5.722) (0.100,-5.721) (0.150,-5.718) (0.200,-5.715) (0.250,-5.711) (0.300,-5.707) (0.350,-5.701) (0.400,-5.694) (0.450,-5.687) (0.500,-5.678) (0.550,-5.669) (0.600,-5.659) (0.650,-5.647) (0.700,-5.635) (0.750,-5.622) (0.800,-5.608) (0.850,-5.593) (0.900,-5.577) (0.950,-5.559) (1.000,-5.541) (1.050,-5.522) (1.100,-5.501) (1.150,-5.480) (1.200,-5.457) (1.250,-5.433) (1.300,-5.407) (1.350,-5.381) (1.400,-5.353) (1.447,-5.325) (1.500,-5.293) (1.550,-5.260) (1.600,-5.226) (1.650,-5.191) (1.700,-5.153) (1.750,-5.114) (1.797,-5.075) (1.850,-5.029) (1.900,-4.983) (1.950,-4.935) (2.000,-4.884) (2.050,-4.830) (2.098,-4.775) (2.140,-4.725) (2.175,-4.680) (2.216,-4.625) (2.250,-4.577) (2.285,-4.525) (2.317,-4.475) (2.347,-4.425) (2.375,-4.376) (2.400,-4.330) (2.425,-4.281) (2.450,-4.229) (2.474,-4.175) (2.496,-4.125) (2.515,-4.075) (2.534,-4.025) (2.551,-3.975) (2.567,-3.925) (2.582,-3.875) (2.595,-3.825) (2.608,-3.775) (2.619,-3.725) (2.629,-3.675) (2.637,-3.625) (2.645,-3.575) (2.651,-3.525) (2.657,-3.475) (2.661,-3.425) (2.664,-3.375) (2.666,-3.325) (2.667,-3.275) (2.666,-3.225) (2.665,-3.175) (2.662,-3.125) (2.658,-3.075) (2.654,-3.025) (2.647,-2.975) (2.640,-2.925) (2.632,-2.875) (2.623,-2.825) (2.612,-2.775) (2.600,-2.725) (2.587,-2.675) (2.573,-2.625) (2.558,-2.575) (2.541,-2.525) (2.523,-2.475) (2.504,-2.425) (2.483,-2.375) (2.462,-2.325) (2.439,-2.275) (2.414,-2.225) (2.389,-2.175) (2.362,-2.125) (2.333,-2.075) (2.303,-2.025) (2.275,-1.980) (2.239,-1.925) (2.205,-1.875) (2.175,-1.833) (2.132,-1.775) (2.100,-1.733) (2.053,-1.675) (2.012,-1.625) (1.975,-1.583) (1.925,-1.527) (1.877,-1.475) (1.829,-1.425) (1.780,-1.375) (1.729,-1.325) (1.676,-1.275) (1.625,-1.228) (1.575,-1.183) (1.525,-1.139) (1.481,-1.100) (1.425,-1.053) (1.375,-1.011) (1.331,-0.975) (1.275,-0.930) (1.225,-0.891) (1.175,-0.851) (1.125,-0.813) (1.076,-0.775) (1.025,-0.737) (0.976,-0.700) (0.925,-0.662) (0.875,-0.625) (0.825,-0.589) (0.775,-0.552) (0.725,-0.516) (0.675,-0.480) (0.635,-0.450) (0.575,-0.408) (0.530,-0.375) (0.475,-0.337) (0.425,-0.301) (0.375,-0.266) (0.325,-0.230) (0.284,-0.200) (0.225,-0.160) (0.206,-0.144)}
\EQ{(0.206,0.144) (0.248,0.175) (0.300,0.213) (0.350,0.248) (0.400,0.284) (0.450,0.319) (0.495,0.350) (0.550,0.390) (0.599,0.425) (0.650,0.462) (0.700,0.498) (0.750,0.534) (0.800,0.570) (0.842,0.600) (0.900,0.644) (0.943,0.675) (1.000,0.718) (1.043,0.750) (1.100,0.794) (1.141,0.825) (1.200,0.871) (1.250,0.910) (1.300,0.950) (1.350,0.991) (1.392,1.025) (1.450,1.074) (1.500,1.117) (1.538,1.150) (1.595,1.200) (1.649,1.250) (1.700,1.298) (1.750,1.346) (1.800,1.395) (1.850,1.447) (1.900,1.499) (1.946,1.550) (1.990,1.600) (2.025,1.641) (2.074,1.700) (2.113,1.750) (2.150,1.799) (2.188,1.850) (2.222,1.900) (2.256,1.950) (2.288,2.000) (2.318,2.050) (2.348,2.100) (2.375,2.149) (2.400,2.197) (2.425,2.247) (2.450,2.299) (2.473,2.350) (2.494,2.400) (2.514,2.450) (2.532,2.500) (2.549,2.550) (2.565,2.600) (2.580,2.650) (2.594,2.700) (2.606,2.750) (2.617,2.800) (2.627,2.850) (2.636,2.900) (2.644,2.950) (2.651,3.000) (2.656,3.050) (2.660,3.100) (2.664,3.150) (2.666,3.200) (2.667,3.250) (2.666,3.300) (2.665,3.350) (2.663,3.400) (2.659,3.450) (2.654,3.500) (2.648,3.550) (2.641,3.600) (2.633,3.650) (2.624,3.700) (2.613,3.750) (2.602,3.800) (2.589,3.850) (2.575,3.899) (2.559,3.950) (2.543,4.000) (2.525,4.050) (2.506,4.100) (2.485,4.150) (2.463,4.200) (2.440,4.250) (2.415,4.300) (2.389,4.350) (2.361,4.400) (2.332,4.450) (2.301,4.500) (2.268,4.550) (2.234,4.600) (2.200,4.647) (2.160,4.700) (2.125,4.743) (2.077,4.800) (2.032,4.850) (1.985,4.900) (1.935,4.950) (1.882,5.000) (1.826,5.050) (1.775,5.094) (1.725,5.134) (1.675,5.172) (1.625,5.209) (1.575,5.244) (1.528,5.275) (1.475,5.308) (1.425,5.338) (1.375,5.367) (1.325,5.394) (1.275,5.420) (1.225,5.445) (1.175,5.468) (1.125,5.491) (1.075,5.512) (1.025,5.532) (0.976,5.550) (0.925,5.568) (0.875,5.585) (0.826,5.600) (0.775,5.615) (0.725,5.629) (0.675,5.641) (0.625,5.653) (0.575,5.664) (0.525,5.674) (0.475,5.683) (0.425,5.691) (0.375,5.698) (0.325,5.704) (0.275,5.709) (0.225,5.713) (0.175,5.717) (0.125,5.720) (0.075,5.721) (0.025,5.722) (-0.050,5.722) (-0.100,5.721) (-0.150,5.718) (-0.200,5.715) (-0.250,5.711) (-0.300,5.707) (-0.350,5.701) (-0.400,5.694) (-0.450,5.687) (-0.500,5.678) (-0.550,5.669) (-0.600,5.659) (-0.650,5.647) (-0.700,5.635) (-0.750,5.622) (-0.800,5.608) (-0.850,5.593) (-0.900,5.577) (-0.950,5.559) (-1.000,5.541) (-1.050,5.522) (-1.100,5.501) (-1.150,5.480) (-1.200,5.457) (-1.250,5.433) (-1.300,5.407) (-1.350,5.381) (-1.400,5.353) (-1.447,5.325) (-1.500,5.293) (-1.550,5.260) (-1.600,5.226) (-1.650,5.191) (-1.700,5.153) (-1.750,5.114) (-1.797,5.075) (-1.850,5.029) (-1.900,4.983) (-1.950,4.935) (-2.000,4.884) (-2.050,4.830) (-2.098,4.775) (-2.140,4.725) (-2.175,4.680) (-2.216,4.625) (-2.250,4.577) (-2.285,4.525) (-2.317,4.475) (-2.347,4.425) (-2.375,4.376) (-2.400,4.330) (-2.425,4.281) (-2.450,4.229) (-2.474,4.175) (-2.496,4.125) (-2.515,4.075) (-2.534,4.025) (-2.551,3.975) (-2.567,3.925) (-2.582,3.875) (-2.595,3.825) (-2.608,3.775) (-2.619,3.725) (-2.629,3.675) (-2.637,3.625) (-2.645,3.575) (-2.651,3.525) (-2.657,3.475) (-2.661,3.425) (-2.664,3.375) (-2.666,3.325) (-2.667,3.275) (-2.666,3.225) (-2.665,3.175) (-2.662,3.125) (-2.658,3.075) (-2.654,3.025) (-2.647,2.975) (-2.640,2.925) (-2.632,2.875) (-2.623,2.825) (-2.612,2.775) (-2.600,2.725) (-2.587,2.675) (-2.573,2.625) (-2.558,2.575) (-2.541,2.525) (-2.523,2.475) (-2.504,2.425) (-2.483,2.375) (-2.462,2.325) (-2.439,2.275) (-2.414,2.225) (-2.389,2.175) (-2.362,2.125) (-2.333,2.075) (-2.303,2.025) (-2.275,1.980) (-2.239,1.925) (-2.205,1.875) (-2.175,1.833) (-2.132,1.775) (-2.100,1.733) (-2.053,1.675) (-2.012,1.625) (-1.975,1.583) (-1.925,1.527) (-1.877,1.475) (-1.829,1.425) (-1.780,1.375) (-1.729,1.325) (-1.676,1.275) (-1.625,1.228) (-1.575,1.183) (-1.525,1.139) (-1.481,1.100) (-1.425,1.053) (-1.375,1.011) (-1.331,0.975) (-1.275,0.930) (-1.225,0.891) (-1.175,0.851) (-1.125,0.813) (-1.076,0.775) (-1.025,0.737) (-0.976,0.700) (-0.925,0.662) (-0.875,0.625) (-0.825,0.589) (-0.775,0.552) (-0.725,0.516) (-0.675,0.480) (-0.635,0.450) (-0.575,0.408) (-0.530,0.375) (-0.475,0.337) (-0.425,0.301) (-0.375,0.266) (-0.325,0.230) (-0.284,0.200) (-0.225,0.160) (-0.220,0.155)}
\EQ{(-0.220,-0.155) (-0.275,-0.195) (-0.319,-0.225) (-0.375,-0.266) (-0.424,-0.300) (-0.475,-0.337) (-0.525,-0.372) (-0.575,-0.408) (-0.625,-0.444) (-0.669,-0.475) (-0.725,-0.516) (-0.772,-0.550) (-0.825,-0.589) (-0.874,-0.625) (-0.925,-0.663) (-0.975,-0.700) (-1.025,-0.738) (-1.074,-0.775) (-1.125,-0.815) (-1.171,-0.850) (-1.225,-0.893) (-1.265,-0.925) (-1.325,-0.974) (-1.375,-1.016) (-1.416,-1.050) (-1.473,-1.100) (-1.525,-1.146) (-1.575,-1.191) (-1.612,-1.225) (-1.650,-1.261) (-1.691,-1.300) (-1.741,-1.350) (-1.790,-1.400) (-1.825,-1.438) (-1.860,-1.475) (-1.900,-1.520) (-1.947,-1.575) (-1.989,-1.625) (-2.025,-1.670) (-2.067,-1.725) (-2.100,-1.770) (-2.139,-1.825) (-2.173,-1.875) (-2.205,-1.925) (-2.236,-1.975) (-2.265,-2.025) (-2.293,-2.075) (-2.319,-2.125) (-2.344,-2.175) (-2.368,-2.225) (-2.390,-2.275) (-2.411,-2.325) (-2.431,-2.375) (-2.449,-2.425) (-2.466,-2.475) (-2.481,-2.525) (-2.496,-2.575) (-2.509,-2.625) (-2.521,-2.675) (-2.532,-2.725) (-2.541,-2.775) (-2.549,-2.825) (-2.557,-2.875) (-2.562,-2.925) (-2.567,-2.975) (-2.571,-3.025) (-2.573,-3.075) (-2.574,-3.125) (-2.574,-3.175) (-2.573,-3.225) (-2.571,-3.275) (-2.568,-3.325) (-2.563,-3.375) (-2.557,-3.425) (-2.550,-3.475) (-2.542,-3.525) (-2.533,-3.575) (-2.522,-3.625) (-2.510,-3.675) (-2.497,-3.725) (-2.483,-3.775) (-2.467,-3.825) (-2.450,-3.875) (-2.432,-3.925) (-2.413,-3.975) (-2.392,-4.025) (-2.369,-4.075) (-2.345,-4.125) (-2.320,-4.175) (-2.293,-4.225) (-2.264,-4.275) (-2.234,-4.325) (-2.202,-4.375) (-2.168,-4.425) (-2.133,-4.475) (-2.100,-4.518) (-2.055,-4.575) (-2.013,-4.625) (-1.975,-4.668) (-1.925,-4.722) (-1.875,-4.772) (-1.825,-4.820) (-1.775,-4.866) (-1.735,-4.900) (-1.675,-4.949) (-1.625,-4.988) (-1.575,-5.025) (-1.525,-5.060) (-1.475,-5.093) (-1.425,-5.124) (-1.375,-5.155) (-1.325,-5.183) (-1.275,-5.210) (-1.225,-5.236) (-1.175,-5.261) (-1.125,-5.284) (-1.075,-5.306) (-1.025,-5.327) (-0.975,-5.346) (-0.925,-5.365) (-0.875,-5.382) (-0.825,-5.398) (-0.775,-5.414) (-0.725,-5.428) (-0.675,-5.441) (-0.625,-5.453) (-0.575,-5.464) (-0.525,-5.474) (-0.475,-5.484) (-0.425,-5.492) (-0.375,-5.499) (-0.325,-5.506) (-0.275,-5.511) (-0.225,-5.516) (-0.175,-5.519) (-0.125,-5.522) (-0.075,-5.524) (-0.025,-5.525) (0.050,-5.524) (0.100,-5.523) (0.150,-5.521) (0.200,-5.518) (0.250,-5.513) (0.300,-5.508) (0.350,-5.503) (0.400,-5.496) (0.450,-5.488) (0.500,-5.479) (0.550,-5.469) (0.600,-5.459) (0.650,-5.447) (0.700,-5.434) (0.750,-5.421) (0.800,-5.406) (0.850,-5.390) (0.900,-5.374) (0.950,-5.356) (1.000,-5.337) (1.050,-5.316) (1.100,-5.295) (1.150,-5.272) (1.197,-5.250) (1.247,-5.225) (1.294,-5.200) (1.340,-5.175) (1.400,-5.140) (1.450,-5.109) (1.500,-5.076) (1.550,-5.042) (1.600,-5.006) (1.642,-4.975) (1.700,-4.929) (1.750,-4.887) (1.792,-4.850) (1.846,-4.800) (1.898,-4.750) (1.946,-4.700) (1.991,-4.650) (2.025,-4.611) (2.075,-4.550) (2.114,-4.500) (2.150,-4.451) (2.185,-4.400) (2.218,-4.350) (2.249,-4.300) (2.275,-4.257) (2.300,-4.212) (2.325,-4.165) (2.350,-4.115) (2.375,-4.062) (2.400,-4.005) (2.423,-3.950) (2.442,-3.900) (2.459,-3.850) (2.475,-3.801) (2.490,-3.750) (2.504,-3.700) (2.516,-3.650) (2.528,-3.600) (2.538,-3.550) (2.546,-3.500) (2.554,-3.450) (2.560,-3.400) (2.566,-3.350) (2.570,-3.300) (2.572,-3.250) (2.574,-3.200) (2.575,-3.150) (2.574,-3.100) (2.572,-3.050) (2.569,-3.000) (2.565,-2.950) (2.560,-2.900) (2.553,-2.850) (2.545,-2.800) (2.537,-2.750) (2.526,-2.700) (2.515,-2.650) (2.502,-2.600) (2.489,-2.550) (2.474,-2.500) (2.457,-2.450) (2.440,-2.400) (2.421,-2.350) (2.401,-2.300) (2.379,-2.250) (2.356,-2.200) (2.332,-2.150) (2.306,-2.100) (2.279,-2.050) (2.251,-2.000) (2.221,-1.950) (2.189,-1.900) (2.156,-1.850) (2.125,-1.805) (2.086,-1.750) (2.050,-1.702) (2.009,-1.650) (1.975,-1.608) (1.926,-1.550) (1.882,-1.500) (1.837,-1.450) (1.800,-1.411) (1.765,-1.375) (1.725,-1.334) (1.675,-1.285) (1.639,-1.250) (1.600,-1.214) (1.558,-1.175) (1.502,-1.125) (1.450,-1.080) (1.400,-1.037) (1.356,-1.000) (1.300,-0.954) (1.250,-0.913) (1.202,-0.875) (1.150,-0.834) (1.107,-0.800) (1.050,-0.757) (1.008,-0.725) (0.950,-0.682) (0.908,-0.650) (0.850,-0.607) (0.807,-0.575) (0.750,-0.534) (0.703,-0.500) (0.650,-0.462) (0.600,-0.426) (0.550,-0.390) (0.500,-0.355) (0.450,-0.319) (0.400,-0.284) (0.354,-0.250) (0.300,-0.213) (0.250,-0.177) (0.206,-0.144)}
\EQ{(0.206,0.144) (0.248,0.175) (0.300,0.213) (0.350,0.248) (0.400,0.284) (0.450,0.319) (0.495,0.350) (0.550,0.390) (0.599,0.425) (0.650,0.462) (0.700,0.498) (0.750,0.534) (0.800,0.571) (0.841,0.600) (0.900,0.644) (0.942,0.675) (1.000,0.719) (1.041,0.750) (1.100,0.795) (1.150,0.834) (1.200,0.873) (1.250,0.913) (1.296,0.950) (1.350,0.995) (1.400,1.037) (1.445,1.075) (1.500,1.123) (1.550,1.168) (1.600,1.214) (1.639,1.250) (1.675,1.285) (1.716,1.325) (1.765,1.375) (1.800,1.411) (1.837,1.450) (1.875,1.492) (1.925,1.549) (1.968,1.600) (2.000,1.639) (2.048,1.700) (2.086,1.750) (2.122,1.800) (2.156,1.850) (2.189,1.900) (2.221,1.950) (2.250,1.999) (2.275,2.043) (2.300,2.088) (2.325,2.136) (2.350,2.187) (2.375,2.241) (2.400,2.299) (2.421,2.350) (2.440,2.400) (2.457,2.450) (2.474,2.500) (2.489,2.550) (2.502,2.600) (2.515,2.650) (2.526,2.700) (2.537,2.750) (2.545,2.800) (2.553,2.850) (2.560,2.900) (2.565,2.950) (2.569,3.000) (2.572,3.050) (2.574,3.100) (2.575,3.150) (2.574,3.200) (2.572,3.250) (2.570,3.300) (2.566,3.350) (2.560,3.400) (2.554,3.450) (2.546,3.500) (2.538,3.550) (2.528,3.600) (2.516,3.650) (2.504,3.700) (2.490,3.750) (2.475,3.800) (2.459,3.850) (2.442,3.900) (2.423,3.950) (2.402,4.000) (2.381,4.050) (2.357,4.100) (2.333,4.150) (2.307,4.200) (2.279,4.250) (2.250,4.299) (2.218,4.350) (2.185,4.400) (2.151,4.450) (2.114,4.500) (2.075,4.550) (2.034,4.600) (2.000,4.640) (1.950,4.695) (1.900,4.747) (1.850,4.797) (1.800,4.843) (1.750,4.887) (1.705,4.925) (1.650,4.969) (1.609,5.000) (1.550,5.042) (1.502,5.075) (1.450,5.109) (1.400,5.140) (1.350,5.169) (1.300,5.197) (1.250,5.223) (1.200,5.249) (1.150,5.272) (1.100,5.295) (1.050,5.316) (1.000,5.337) (0.950,5.356) (0.900,5.374) (0.850,5.390) (0.800,5.406) (0.750,5.421) (0.700,5.434) (0.650,5.447) (0.600,5.459) (0.550,5.469) (0.500,5.479) (0.450,5.488) (0.400,5.496) (0.350,5.503) (0.300,5.508) (0.250,5.513) (0.200,5.518) (0.150,5.521) (0.100,5.523) (0.050,5.524) (0.000,5.525) (-0.050,5.524) (-0.100,5.523) (-0.150,5.521) (-0.200,5.518) (-0.250,5.513) (-0.300,5.508) (-0.350,5.503) (-0.400,5.496) (-0.450,5.488) (-0.500,5.479) (-0.550,5.469) (-0.600,5.459) (-0.650,5.447) (-0.700,5.434) (-0.750,5.421) (-0.800,5.406) (-0.850,5.390) (-0.900,5.374) (-0.950,5.356) (-1.000,5.337) (-1.050,5.316) (-1.100,5.295) (-1.150,5.272) (-1.197,5.250) (-1.247,5.225) (-1.294,5.200) (-1.340,5.175) (-1.400,5.140) (-1.450,5.109) (-1.500,5.076) (-1.550,5.042) (-1.600,5.006) (-1.642,4.975) (-1.700,4.929) (-1.750,4.887) (-1.792,4.850) (-1.846,4.800) (-1.898,4.750) (-1.946,4.700) (-1.991,4.650) (-2.025,4.611) (-2.075,4.550) (-2.114,4.500) (-2.150,4.451) (-2.185,4.400) (-2.218,4.350) (-2.249,4.300) (-2.275,4.257) (-2.300,4.212) (-2.325,4.165) (-2.350,4.115) (-2.375,4.062) (-2.400,4.005) (-2.423,3.950) (-2.442,3.900) (-2.459,3.850) (-2.475,3.801) (-2.490,3.750) (-2.504,3.700) (-2.516,3.650) (-2.528,3.600) (-2.538,3.550) (-2.546,3.500) (-2.554,3.450) (-2.560,3.400) (-2.566,3.350) (-2.570,3.300) (-2.572,3.250) (-2.574,3.200) (-2.575,3.150) (-2.574,3.100) (-2.572,3.050) (-2.569,3.000) (-2.565,2.950) (-2.560,2.900) (-2.553,2.850) (-2.545,2.800) (-2.537,2.750) (-2.526,2.700) (-2.515,2.650) (-2.502,2.600) (-2.489,2.550) (-2.474,2.500) (-2.457,2.450) (-2.440,2.400) (-2.421,2.350) (-2.401,2.300) (-2.379,2.250) (-2.356,2.200) (-2.332,2.150) (-2.306,2.100) (-2.279,2.050) (-2.251,2.000) (-2.221,1.950) (-2.189,1.900) (-2.156,1.850) (-2.125,1.805) (-2.086,1.750) (-2.050,1.702) (-2.009,1.650) (-1.975,1.608) (-1.926,1.550) (-1.882,1.500) (-1.837,1.450) (-1.800,1.411) (-1.765,1.375) (-1.725,1.334) (-1.675,1.285) (-1.639,1.250) (-1.600,1.214) (-1.558,1.175) (-1.502,1.125) (-1.450,1.080) (-1.400,1.037) (-1.356,1.000) (-1.300,0.954) (-1.250,0.913) (-1.202,0.875) (-1.150,0.834) (-1.107,0.800) (-1.050,0.757) (-1.008,0.725) (-0.950,0.682) (-0.908,0.650) (-0.850,0.607) (-0.807,0.575) (-0.750,0.534) (-0.703,0.500) (-0.650,0.462) (-0.600,0.426) (-0.550,0.390) (-0.500,0.355) (-0.450,0.319) (-0.400,0.284) (-0.354,0.250) (-0.300,0.213) (-0.250,0.177) (-0.220,0.155)}

    % Cuadrupolo puntual - líneas de campo
% Generado por scripts/generar_datos.py (no editar a mano).
\LC{(0.435,1.114) (0.489,1.167) (0.546,1.216) (0.605,1.261) (0.667,1.304) (0.732,1.342) (0.799,1.375) (0.868,1.405) (0.939,1.429) (1.011,1.448) (1.085,1.461) (1.160,1.468) (1.235,1.468) (1.310,1.461) (1.383,1.446) (1.455,1.424) (1.523,1.394) (1.588,1.355) (1.646,1.309) (1.698,1.255) (1.742,1.194) (1.777,1.128) (1.803,1.057) (1.817,0.984) (1.822,0.909) (1.817,0.834) (1.803,0.761) (1.780,0.689) (1.749,0.621) (1.712,0.556) (1.669,0.495) (1.621,0.437) (1.568,0.384) (1.511,0.335) (1.451,0.290) (1.388,0.249) (1.323,0.213) (1.255,0.179) (1.186,0.150)}
\LC{(-0.449,1.110) (-0.504,1.161) (-0.562,1.208) (-0.623,1.252) (-0.687,1.292) (-0.753,1.327) (-0.821,1.358) (-0.891,1.384) (-0.963,1.405) (-1.037,1.420) (-1.111,1.428) (-1.186,1.430) (-1.261,1.424) (-1.335,1.411) (-1.407,1.390) (-1.476,1.361) (-1.541,1.323) (-1.600,1.278) (-1.652,1.224) (-1.696,1.163) (-1.731,1.097) (-1.756,1.026) (-1.770,0.953) (-1.774,0.878) (-1.768,0.803) (-1.753,0.730) (-1.728,0.659) (-1.697,0.591) (-1.658,0.527) (-1.613,0.467) (-1.564,0.411) (-1.510,0.359) (-1.452,0.311) (-1.391,0.268) (-1.327,0.228) (-1.261,0.193) (-1.193,0.161) (-1.170,0.152)}
\LC{(-0.693,-0.961) (-0.768,-0.970) (-0.843,-0.968) (-0.916,-0.955) (-0.987,-0.930) (-1.052,-0.893) (-1.109,-0.844) (-1.154,-0.784) (-1.185,-0.716) (-1.201,-0.643) (-1.202,-0.568) (-1.189,-0.495) (-1.163,-0.424) (-1.126,-0.359)}
\LC{(0.153,-1.185) (0.173,-1.258) (0.193,-1.330) (0.215,-1.401) (0.239,-1.473) (0.263,-1.544) (0.289,-1.614) (0.316,-1.684) (0.343,-1.754) (0.373,-1.823) (0.403,-1.892) (0.434,-1.960) (0.467,-2.027) (0.501,-2.094) (0.536,-2.160) (0.572,-2.226) (0.609,-2.291) (0.648,-2.356) (0.687,-2.419) (0.728,-2.482) (0.770,-2.545) (0.813,-2.606) (0.857,-2.667) (0.902,-2.727) (0.948,-2.786) (0.995,-2.844) (1.044,-2.901) (1.094,-2.957) (1.144,-3.013) (1.196,-3.067) (1.249,-3.120) (1.303,-3.172) (1.358,-3.223) (1.413,-3.273) (1.470,-3.322) (1.528,-3.370) (1.587,-3.416) (1.647,-3.461) (1.708,-3.505) (1.770,-3.547) (1.833,-3.588) (1.897,-3.628) (1.961,-3.666) (2.027,-3.703) (2.093,-3.738) (2.160,-3.771) (2.228,-3.803) (2.297,-3.833) (2.366,-3.861) (2.436,-3.888) (2.507,-3.913) (2.578,-3.936) (2.651,-3.956) (2.723,-3.975) (2.796,-3.992) (2.870,-4.007) (2.944,-4.019) (3.018,-4.030) (3.093,-4.038) (3.167,-4.043) (3.242,-4.046) (3.317,-4.047) (3.392,-4.046) (3.467,-4.041) (3.542,-4.035) (3.616,-4.025) (3.690,-4.013) (3.764,-3.998) (3.836,-3.980) (3.909,-3.959) (3.980,-3.935) (4.050,-3.909) (4.119,-3.879) (4.186,-3.847) (4.253,-3.812) (4.317,-3.773) (4.380,-3.732) (4.440,-3.688) (4.499,-3.641) (4.555,-3.591) (4.608,-3.539) (4.659,-3.484) (4.707,-3.426) (4.752,-3.366) (4.794,-3.304) (4.833,-3.240) (4.868,-3.173) (4.900,-3.105) (4.928,-3.036) (4.953,-2.965) (4.973,-2.893) (4.991,-2.820) (5.004,-2.746) (5.014,-2.672) (5.020,-2.597) (5.022,-2.522) (5.021,-2.447) (5.016,-2.372) (5.008,-2.298) (4.996,-2.224) (4.981,-2.150) (4.963,-2.077) (4.942,-2.005) (4.918,-1.934) (4.891,-1.864) (4.862,-1.795) (4.830,-1.728) (4.795,-1.661) (4.758,-1.596) (4.719,-1.532) (4.677,-1.469) (4.634,-1.408) (4.588,-1.349) (4.541,-1.291) (4.492,-1.234) (4.441,-1.179) (4.389,-1.125) (4.335,-1.073) (4.280,-1.022) (4.223,-0.973) (4.165,-0.925) (4.106,-0.879) (4.046,-0.834) (3.985,-0.791) (3.923,-0.749) (3.859,-0.709) (3.795,-0.670) (3.730,-0.632) (3.665,-0.596) (3.598,-0.561) (3.531,-0.528) (3.463,-0.496) (3.395,-0.465) (3.326,-0.436) (3.256,-0.407) (3.186,-0.380) (3.116,-0.355) (3.045,-0.330) (2.974,-0.306) (2.902,-0.284) (2.830,-0.263) (2.758,-0.243) (2.686,-0.224) (2.613,-0.206) (2.540,-0.189) (2.467,-0.172) (2.393,-0.157) (2.319,-0.143) (2.246,-0.130) (2.172,-0.117) (2.098,-0.105) (2.023,-0.095) (1.949,-0.084) (1.875,-0.075) (1.800,-0.066) (1.751,-0.061)}
\LC{(0.283,-2.625) (0.299,-2.698) (0.316,-2.771) (0.333,-2.844) (0.350,-2.917) (0.368,-2.990) (0.387,-3.063) (0.406,-3.135) (0.425,-3.208) (0.445,-3.280) (0.466,-3.352) (0.486,-3.424) (0.507,-3.496) (0.529,-3.568) (0.551,-3.640) (0.574,-3.711) (0.597,-3.783) (0.620,-3.854) (0.644,-3.925) (0.669,-3.996) (0.693,-4.067) (0.719,-4.137) (0.744,-4.208) (0.770,-4.278) (0.797,-4.348) (0.824,-4.418) (0.852,-4.488) (0.879,-4.557) (0.908,-4.627) (0.937,-4.696) (0.966,-4.765) (0.996,-4.834) (1.026,-4.903) (1.056,-4.971) (1.087,-5.039) (1.119,-5.108) (1.151,-5.175) (1.183,-5.243) (1.216,-5.310) (1.249,-5.378) (1.283,-5.445) (1.317,-5.511) (1.352,-5.578) (1.387,-5.644) (1.422,-5.710) (1.458,-5.776) (1.494,-5.842) (1.531,-5.907) (1.568,-5.972) (1.606,-6.037) (1.644,-6.102) (1.682,-6.166) (1.721,-6.230) (1.761,-6.294) (1.801,-6.358) (1.841,-6.421) (1.882,-6.484) (1.923,-6.547) (1.964,-6.609) (2.006,-6.671) (2.048,-6.733) (2.091,-6.795) (2.135,-6.856) (2.178,-6.917) (2.222,-6.978) (2.267,-7.038) (2.312,-7.098) (2.357,-7.158) (2.403,-7.217) (2.449,-7.276) (2.496,-7.335) (2.543,-7.394) (2.590,-7.452) (2.638,-7.509) (2.686,-7.567) (2.735,-7.624) (2.784,-7.680) (2.834,-7.737) (2.884,-7.793) (2.934,-7.848) (2.985,-7.903) (3.036,-7.958) (3.088,-8.013) (3.140,-8.067) (3.192,-8.120) (3.245,-8.174) (3.298,-8.227) (3.352,-8.279) (3.406,-8.331) (3.460,-8.383) (3.515,-8.434) (3.570,-8.485) (3.626,-8.535) (3.682,-8.585) (3.738,-8.634) (3.795,-8.683) (3.852,-8.732) (3.909,-8.780) (3.967,-8.828) (4.026,-8.875) (4.084,-8.922) (4.143,-8.968) (4.203,-9.014) (4.262,-9.059) (4.323,-9.104) (4.383,-9.148) (4.444,-9.192) (4.505,-9.235) (4.567,-9.278) (4.629,-9.320) (4.691,-9.362) (4.754,-9.403) (4.817,-9.444) (4.880,-9.484) (4.944,-9.524) (5.008,-9.563) (5.072,-9.601) (5.137,-9.639) (5.202,-9.676) (5.267,-9.713) (5.333,-9.749) (5.399,-9.785) (5.466,-9.820) (5.532,-9.854) (5.599,-9.888) (5.666,-9.921) (5.734,-9.954) (5.802,-9.986) (5.847,-10.000)}
\LC{(-0.162,-1.698) (-0.176,-1.772) (-0.191,-1.845) (-0.207,-1.919) (-0.223,-1.992) (-0.240,-2.065) (-0.258,-2.138) (-0.276,-2.210) (-0.295,-2.283) (-0.314,-2.356) (-0.334,-2.428) (-0.355,-2.500) (-0.376,-2.572) (-0.398,-2.644) (-0.421,-2.715) (-0.444,-2.786) (-0.468,-2.858) (-0.492,-2.928) (-0.517,-2.999) (-0.543,-3.070) (-0.569,-3.140) (-0.596,-3.210) (-0.623,-3.280) (-0.651,-3.349) (-0.680,-3.419) (-0.709,-3.488) (-0.739,-3.557) (-0.769,-3.625) (-0.800,-3.693) (-0.832,-3.761) (-0.864,-3.829) (-0.897,-3.896) (-0.931,-3.963) (-0.965,-4.030) (-1.000,-4.097) (-1.035,-4.163) (-1.071,-4.229) (-1.107,-4.294) (-1.144,-4.359) (-1.182,-4.424) (-1.220,-4.489) (-1.259,-4.553) (-1.299,-4.617) (-1.339,-4.680) (-1.380,-4.743) (-1.421,-4.806) (-1.463,-4.868) (-1.505,-4.930) (-1.548,-4.991) (-1.592,-5.052) (-1.636,-5.113) (-1.681,-5.173) (-1.726,-5.233) (-1.772,-5.292) (-1.819,-5.351) (-1.866,-5.409) (-1.913,-5.467) (-1.961,-5.525) (-2.010,-5.582) (-2.060,-5.638) (-2.110,-5.694) (-2.160,-5.749) (-2.211,-5.804) (-2.263,-5.859) (-2.315,-5.913) (-2.368,-5.966) (-2.421,-6.019) (-2.475,-6.071) (-2.529,-6.123) (-2.584,-6.174) (-2.639,-6.224) (-2.695,-6.274) (-2.752,-6.324) (-2.809,-6.372) (-2.866,-6.420) (-2.924,-6.468) (-2.983,-6.515) (-3.042,-6.561) (-3.102,-6.606) (-3.162,-6.651) (-3.222,-6.695) (-3.284,-6.739) (-3.345,-6.782) (-3.407,-6.824) (-3.470,-6.865) (-3.533,-6.906) (-3.597,-6.945) (-3.661,-6.984) (-3.725,-7.023) (-3.790,-7.060) (-3.855,-7.097) (-3.921,-7.133) (-3.987,-7.168) (-4.054,-7.202) (-4.121,-7.236) (-4.189,-7.268) (-4.257,-7.300) (-4.325,-7.331) (-4.394,-7.361) (-4.463,-7.390) (-4.533,-7.418) (-4.602,-7.445) (-4.673,-7.472) (-4.743,-7.497) (-4.814,-7.521) (-4.886,-7.545) (-4.957,-7.567) (-5.029,-7.588) (-5.101,-7.609) (-5.174,-7.628) (-5.246,-7.646) (-5.319,-7.663) (-5.393,-7.679) (-5.466,-7.694) (-5.540,-7.708) (-5.614,-7.721) (-5.688,-7.732) (-5.762,-7.743) (-5.837,-7.752) (-5.911,-7.760) (-5.986,-7.767) (-6.061,-7.773) (-6.136,-7.777) (-6.210,-7.780) (-6.285,-7.782) (-6.360,-7.782) (-6.435,-7.782) (-6.510,-7.779) (-6.585,-7.776) (-6.660,-7.771) (-6.735,-7.765) (-6.810,-7.758) (-6.884,-7.749) (-6.958,-7.739) (-7.032,-7.727) (-7.106,-7.714) (-7.180,-7.699) (-7.253,-7.683) (-7.326,-7.665) (-7.399,-7.646) (-7.471,-7.626) (-7.542,-7.604) (-7.614,-7.581) (-7.684,-7.556) (-7.754,-7.529) (-7.824,-7.501) (-7.893,-7.471) (-7.961,-7.440) (-8.029,-7.408) (-8.095,-7.373) (-8.161,-7.338) (-8.227,-7.300) (-8.291,-7.262) (-8.354,-7.221) (-8.416,-7.180) (-8.477,-7.136) (-8.538,-7.091) (-8.597,-7.045) (-8.654,-6.997) (-8.711,-6.948) (-8.766,-6.898) (-8.820,-6.845) (-8.873,-6.792) (-8.924,-6.737) (-8.974,-6.681) (-9.022,-6.624) (-9.069,-6.565) (-9.114,-6.505) (-9.158,-6.444) (-9.200,-6.382) (-9.240,-6.319) (-9.278,-6.254) (-9.315,-6.189) (-9.350,-6.123) (-9.383,-6.055) (-9.415,-5.987) (-9.444,-5.918) (-9.472,-5.849) (-9.497,-5.778) (-9.521,-5.707) (-9.543,-5.635) (-9.563,-5.563) (-9.581,-5.490) (-9.597,-5.417) (-9.611,-5.343) (-9.623,-5.269) (-9.634,-5.195) (-9.642,-5.120) (-9.648,-5.046) (-9.653,-4.971) (-9.656,-4.896) (-9.656,-4.821) (-9.655,-4.746) (-9.652,-4.671) (-9.648,-4.596) (-9.641,-4.521) (-9.633,-4.447) (-9.622,-4.373) (-9.611,-4.298) (-9.597,-4.225) (-9.582,-4.151) (-9.565,-4.078) (-9.546,-4.006) (-9.526,-3.933) (-9.505,-3.861) (-9.482,-3.790) (-9.457,-3.719) (-9.431,-3.649) (-9.403,-3.579) (-9.374,-3.510) (-9.344,-3.442) (-9.312,-3.374) (-9.279,-3.306) (-9.245,-3.240) (-9.209,-3.174) (-9.172,-3.108) (-9.134,-3.044) (-9.095,-2.980) (-9.055,-2.916) (-9.013,-2.854) (-8.971,-2.792) (-8.927,-2.731) (-8.882,-2.671) (-8.837,-2.611) (-8.790,-2.553) (-8.743,-2.495) (-8.694,-2.437) (-8.645,-2.381) (-8.595,-2.325) (-8.544,-2.270) (-8.492,-2.216) (-8.439,-2.163) (-8.385,-2.110) (-8.331,-2.058) (-8.276,-2.007) (-8.221,-1.957) (-8.164,-1.908) (-8.107,-1.859) (-8.049,-1.811) (-7.991,-1.764) (-7.932,-1.718) (-7.872,-1.672) (-7.812,-1.627) (-7.751,-1.583) (-7.690,-1.540) (-7.628,-1.498) (-7.566,-1.456) (-7.503,-1.415) (-7.440,-1.375) (-7.376,-1.335) (-7.312,-1.297) (-7.247,-1.259) (-7.182,-1.221) (-7.117,-1.185) (-7.051,-1.149) (-6.985,-1.114) (-6.918,-1.079) (-6.851,-1.046) (-6.784,-1.012) (-6.716,-0.980) (-6.648,-0.948) (-6.580,-0.917) (-6.511,-0.887) (-6.442,-0.858) (-6.373,-0.829) (-6.304,-0.800) (-6.234,-0.773) (-6.164,-0.745) (-6.094,-0.719) (-6.023,-0.693) (-5.953,-0.668) (-5.882,-0.643) (-5.811,-0.620) (-5.740,-0.596) (-5.668,-0.573) (-5.596,-0.551) (-5.525,-0.530) (-5.453,-0.508) (-5.380,-0.488) (-5.308,-0.468) (-5.236,-0.449) (-5.163,-0.430) (-5.090,-0.412) (-5.018,-0.394) (-4.945,-0.376) (-4.871,-0.360) (-4.798,-0.343) (-4.725,-0.328) (-4.652,-0.312) (-4.578,-0.298) (-4.504,-0.283) (-4.431,-0.269) (-4.357,-0.256) (-4.283,-0.243) (-4.209,-0.231) (-4.135,-0.219) (-4.061,-0.207) (-3.987,-0.196) (-3.913,-0.185) (-3.838,-0.174) (-3.764,-0.164) (-3.690,-0.155) (-3.615,-0.146) (-3.541,-0.137) (-3.466,-0.128) (-3.392,-0.120) (-3.317,-0.112) (-3.242,-0.105) (-3.168,-0.098) (-3.093,-0.091) (-3.018,-0.085) (-2.944,-0.078) (-2.869,-0.073) (-2.794,-0.067) (-2.719,-0.062) (-2.644,-0.057) (-2.569,-0.052) (-2.495,-0.048) (-2.420,-0.044) (-2.345,-0.040) (-2.270,-0.036) (-2.195,-0.032) (-2.120,-0.029) (-2.045,-0.026) (-1.970,-0.023) (-1.895,-0.021) (-1.820,-0.019) (-1.770,-0.017)}
\LC{(-0.196,1.492) (-0.216,1.565) (-0.237,1.637) (-0.259,1.708) (-0.282,1.780) (-0.305,1.851) (-0.330,1.922) (-0.356,1.992) (-0.382,2.062) (-0.410,2.132) (-0.439,2.201) (-0.468,2.270) (-0.499,2.339) (-0.530,2.407) (-0.562,2.475) (-0.595,2.542) (-0.630,2.609) (-0.665,2.675) (-0.701,2.741) (-0.738,2.806) (-0.776,2.871) (-0.815,2.935) (-0.854,2.998) (-0.895,3.061) (-0.937,3.124) (-0.979,3.185) (-1.023,3.247) (-1.067,3.307) (-1.112,3.367) (-1.158,3.426) (-1.206,3.484) (-1.253,3.542) (-1.302,3.599) (-1.352,3.655) (-1.403,3.711) (-1.454,3.765) (-1.506,3.819) (-1.559,3.872) (-1.614,3.924) (-1.668,3.975) (-1.724,4.025) (-1.781,4.074) (-1.838,4.123) (-1.896,4.170) (-1.955,4.217) (-2.015,4.262) (-2.075,4.306) (-2.137,4.349) (-2.199,4.391) (-2.262,4.432) (-2.325,4.472) (-2.389,4.511) (-2.454,4.548) (-2.520,4.584) (-2.587,4.619) (-2.654,4.653) (-2.721,4.685) (-2.790,4.716) (-2.859,4.745) (-2.928,4.773) (-2.998,4.800) (-3.069,4.825) (-3.140,4.849) (-3.212,4.871) (-3.284,4.891) (-3.357,4.910) (-3.430,4.928) (-3.503,4.943) (-3.577,4.957) (-3.651,4.969) (-3.725,4.979) (-3.800,4.987) (-3.874,4.994) (-3.949,4.998) (-4.024,5.001) (-4.099,5.002) (-4.174,5.000) (-4.249,4.997) (-4.324,4.991) (-4.398,4.983) (-4.473,4.973) (-4.547,4.961) (-4.620,4.947) (-4.693,4.930) (-4.766,4.911) (-4.838,4.890) (-4.909,4.866) (-4.979,4.840) (-5.049,4.812) (-5.117,4.781) (-5.185,4.748) (-5.251,4.713) (-5.315,4.675) (-5.379,4.635) (-5.440,4.592) (-5.500,4.547) (-5.559,4.500) (-5.615,4.450) (-5.669,4.399) (-5.722,4.345) (-5.772,4.289) (-5.819,4.231) (-5.865,4.171) (-5.907,4.110) (-5.947,4.046) (-5.985,3.981) (-6.019,3.915) (-6.051,3.847) (-6.080,3.778) (-6.106,3.707) (-6.129,3.636) (-6.149,3.564) (-6.166,3.491) (-6.180,3.417) (-6.191,3.343) (-6.199,3.268) (-6.204,3.193) (-6.206,3.118) (-6.205,3.043) (-6.201,2.968) (-6.195,2.894) (-6.186,2.819) (-6.174,2.745) (-6.159,2.672) (-6.142,2.599) (-6.122,2.526) (-6.100,2.455) (-6.075,2.384) (-6.049,2.314) (-6.020,2.245) (-5.989,2.176) (-5.955,2.109) (-5.920,2.043) (-5.883,1.978) (-5.844,1.913) (-5.804,1.850) (-5.761,1.789) (-5.717,1.728) (-5.672,1.668) (-5.624,1.610) (-5.576,1.553) (-5.526,1.497) (-5.475,1.442) (-5.422,1.389) (-5.368,1.336) (-5.314,1.285) (-5.258,1.235) (-5.200,1.187) (-5.142,1.139) (-5.083,1.093) (-5.023,1.048) (-4.963,1.004) (-4.901,0.961) (-4.838,0.920) (-4.775,0.880) (-4.711,0.840) (-4.647,0.802) (-4.581,0.765) (-4.515,0.730) (-4.449,0.695) (-4.382,0.661) (-4.314,0.629) (-4.246,0.597) (-4.178,0.567) (-4.109,0.538) (-4.039,0.509) (-3.969,0.482) (-3.899,0.456) (-3.828,0.430) (-3.758,0.406) (-3.686,0.382) (-3.615,0.360) (-3.543,0.338) (-3.495,0.324)}
\LC{(0.436,1.693) (0.474,1.758) (0.513,1.822) (0.553,1.885) (0.595,1.947) (0.639,2.009) (0.684,2.069) (0.730,2.127) (0.778,2.185) (0.827,2.242) (0.878,2.297) (0.930,2.351) (0.984,2.403) (1.039,2.454) (1.095,2.504) (1.153,2.552) (1.212,2.598) (1.272,2.643) (1.333,2.686) (1.396,2.727) (1.460,2.766) (1.525,2.803) (1.591,2.839) (1.659,2.872) (1.727,2.903) (1.796,2.931) (1.866,2.957) (1.938,2.981) (2.009,3.003) (2.082,3.021) (2.155,3.037) (2.229,3.050) (2.304,3.060) (2.378,3.067) (2.453,3.071) (2.528,3.072) (2.603,3.070) (2.678,3.064) (2.752,3.054) (2.826,3.041) (2.899,3.024) (2.971,3.004) (3.042,2.980) (3.112,2.951) (3.180,2.919) (3.245,2.883) (3.309,2.844) (3.370,2.800) (3.428,2.753) (3.483,2.702) (3.534,2.647) (3.582,2.589) (3.626,2.528) (3.665,2.465) (3.701,2.398) (3.731,2.330) (3.757,2.259) (3.777,2.187) (3.793,2.114) (3.804,2.040) (3.811,1.965) (3.812,1.890) (3.809,1.815) (3.801,1.741) (3.788,1.667) (3.772,1.594) (3.751,1.521) (3.727,1.451) (3.699,1.381) (3.667,1.313) (3.632,1.247) (3.594,1.182) (3.553,1.119) (3.509,1.058) (3.463,0.999) (3.415,0.942) (3.364,0.887) (3.311,0.834) (3.256,0.782) (3.200,0.733) (3.141,0.686) (3.082,0.641) (3.020,0.597) (2.958,0.556) (2.894,0.516) (2.829,0.479) (2.763,0.443) (2.697,0.409) (2.629,0.376) (2.560,0.346) (2.491,0.317) (2.445,0.299)}
\LC{(0.289,2.296) (0.308,2.368) (0.328,2.441) (0.348,2.513) (0.369,2.585) (0.390,2.657) (0.412,2.729) (0.434,2.800) (0.458,2.872) (0.481,2.943) (0.506,3.014) (0.531,3.084) (0.556,3.155) (0.582,3.225) (0.609,3.295) (0.636,3.365) (0.664,3.435) (0.693,3.504) (0.722,3.573) (0.752,3.642) (0.782,3.711) (0.813,3.779) (0.844,3.847) (0.876,3.915) (0.909,3.983) (0.942,4.050) (0.976,4.117) (1.010,4.183) (1.045,4.250) (1.081,4.316) (1.117,4.381) (1.154,4.447) (1.191,4.512) (1.229,4.577) (1.267,4.641) (1.306,4.705) (1.346,4.769) (1.386,4.832) (1.427,4.895) (1.468,4.958) (1.510,5.020) (1.553,5.082) (1.596,5.143) (1.639,5.204) (1.684,5.265) (1.728,5.325) (1.774,5.385) (1.819,5.444) (1.866,5.503) (1.913,5.561) (1.960,5.619) (2.009,5.677) (2.057,5.734) (2.106,5.791) (2.156,5.847) (2.206,5.902) (2.257,5.957) (2.309,6.012) (2.361,6.066) (2.413,6.120) (2.466,6.173) (2.520,6.225) (2.574,6.277) (2.628,6.329) (2.683,6.380) (2.739,6.430) (2.795,6.480) (2.852,6.529) (2.909,6.577) (2.967,6.625) (3.025,6.673) (3.083,6.719) (3.143,6.765) (3.202,6.811) (3.262,6.856) (3.323,6.900) (3.384,6.943) (3.446,6.986) (3.508,7.028) (3.571,7.069) (3.634,7.110) (3.697,7.150) (3.761,7.189) (3.825,7.228) (3.890,7.265) (3.956,7.302) (4.021,7.338) (4.087,7.374) (4.154,7.408) (4.221,7.442) (4.288,7.475) (4.356,7.507) (4.424,7.538) (4.493,7.568) (4.562,7.598) (4.631,7.627) (4.701,7.654) (4.771,7.681) (4.841,7.707) (4.912,7.732) (4.983,7.756) (5.054,7.779) (5.126,7.801) (5.198,7.822) (5.270,7.842) (5.343,7.861) (5.416,7.879) (5.489,7.896) (5.562,7.912) (5.636,7.927) (5.709,7.941) (5.783,7.954) (5.857,7.965) (5.932,7.976) (6.006,7.985) (6.081,7.993) (6.155,8.000) (6.230,8.006) (6.305,8.011) (6.380,8.014) (6.455,8.016) (6.530,8.017) (6.605,8.017) (6.680,8.015) (6.755,8.012) (6.830,8.008) (6.904,8.002) (6.979,7.995) (7.054,7.987) (7.128,7.978) (7.202,7.967) (7.276,7.954) (7.350,7.941) (7.423,7.926) (7.496,7.909) (7.569,7.891) (7.642,7.872) (7.714,7.851) (7.785,7.828) (7.856,7.805) (7.927,7.779) (7.997,7.753) (8.067,7.724) (8.136,7.695) (8.204,7.664) (8.271,7.631) (8.338,7.597) (8.404,7.561) (8.469,7.524) (8.533,7.485) (8.597,7.445) (8.659,7.404) (8.721,7.360) (8.781,7.316) (8.840,7.270) (8.898,7.223) (8.955,7.174) (9.011,7.124) (9.066,7.072) (9.119,7.019) (9.171,6.965) (9.221,6.910) (9.270,6.853) (9.318,6.795) (9.364,6.736) (9.408,6.675) (9.451,6.614) (9.492,6.551) (9.532,6.487) (9.570,6.423) (9.606,6.357) (9.641,6.290) (9.673,6.223) (9.704,6.155) (9.733,6.086) (9.761,6.016) (9.786,5.945) (9.810,5.874) (9.832,5.802) (9.852,5.730) (9.870,5.657) (9.886,5.584) (9.900,5.510) (9.912,5.436) (9.923,5.362) (9.932,5.288) (9.938,5.213) (9.943,5.138) (9.946,5.063) (9.948,4.988) (9.947,4.913) (9.945,4.838) (9.941,4.763) (9.935,4.689) (9.927,4.614) (9.918,4.539) (9.907,4.465) (9.894,4.391) (9.880,4.318) (9.864,4.244) (9.846,4.172) (9.827,4.099) (9.806,4.027) (9.784,3.955) (9.761,3.884) (9.736,3.813) (9.709,3.743) (9.681,3.674) (9.652,3.605) (9.621,3.536) (9.589,3.468) (9.556,3.401) (9.522,3.334) (9.486,3.268) (9.449,3.203) (9.411,3.138) (9.372,3.074) (9.332,3.011) (9.291,2.948) (9.248,2.887) (9.205,2.825) (9.161,2.765) (9.115,2.705) (9.069,2.646) (9.022,2.588) (8.973,2.531) (8.924,2.474) (8.874,2.418) (8.824,2.363) (8.772,2.308) (8.720,2.254) (8.667,2.201) (8.613,2.149) (8.558,2.098) (8.503,2.047) (8.447,1.997) (8.390,1.948) (8.333,1.900) (8.275,1.852) (8.217,1.805) (8.157,1.759) (8.098,1.714) (8.037,1.669) (7.977,1.625) (7.915,1.582) (7.853,1.540) (7.791,1.498) (7.728,1.457) (7.665,1.417) (7.601,1.377) (7.537,1.338) (7.472,1.300) (7.407,1.263) (7.342,1.226) (7.276,1.190) (7.210,1.155) (7.143,1.120) (7.077,1.086) (7.009,1.053) (6.942,1.021) (6.874,0.989) (6.806,0.958) (6.737,0.927) (6.668,0.897) (6.599,0.868) (6.530,0.839) (6.461,0.811) (6.391,0.784) (6.321,0.757) (6.250,0.731) (6.180,0.705) (6.109,0.680) (6.038,0.656) (5.967,0.632) (5.896,0.609) (5.824,0.586) (5.753,0.564) (5.681,0.542) (5.609,0.521) (5.537,0.501) (5.464,0.481) (5.392,0.462) (5.319,0.443) (5.247,0.425) (5.174,0.407) (5.101,0.389) (5.028,0.373) (4.954,0.356) (4.881,0.340) (4.808,0.325) (4.734,0.310) (4.661,0.296) (4.587,0.282) (4.513,0.268) (4.439,0.255) (4.365,0.242) (4.291,0.230) (4.217,0.218) (4.143,0.207) (4.069,0.196) (3.995,0.185) (3.920,0.175) (3.846,0.165) (3.772,0.156) (3.697,0.147) (3.623,0.138) (3.548,0.130) (3.474,0.122) (3.399,0.114) (3.324,0.107) (3.250,0.100) (3.200,0.095)}
\LC{(1.013,3.088) (1.059,3.147) (1.106,3.206) (1.154,3.264) (1.202,3.321) (1.252,3.377) (1.303,3.432) (1.355,3.486) (1.407,3.540) (1.461,3.592) (1.515,3.644) (1.570,3.695) (1.627,3.744) (1.684,3.793) (1.742,3.840) (1.801,3.887) (1.861,3.932) (1.921,3.976) (1.983,4.019) (2.045,4.061) (2.108,4.102) (2.172,4.141) (2.236,4.179) (2.302,4.216) (2.368,4.251) (2.435,4.285) (2.503,4.317) (2.571,4.348) (2.640,4.378) (2.709,4.406) (2.780,4.432) (2.850,4.457) (2.922,4.480) (2.994,4.501) (3.066,4.521) (3.139,4.539) (3.212,4.555) (3.286,4.569) (3.360,4.581) (3.434,4.591) (3.509,4.600) (3.583,4.606) (3.658,4.610) (3.733,4.612) (3.808,4.612) (3.883,4.610) (3.958,4.605) (4.033,4.598) (4.107,4.589) (4.181,4.577) (4.255,4.563) (4.328,4.547) (4.401,4.528) (4.472,4.506) (4.543,4.482) (4.614,4.455) (4.683,4.426) (4.751,4.394) (4.817,4.360) (4.882,4.323) (4.946,4.283) (5.008,4.241) (5.068,4.196) (5.127,4.149) (5.183,4.100) (5.237,4.048) (5.289,3.993) (5.338,3.937) (5.385,3.878) (5.429,3.818) (5.470,3.755) (5.508,3.690) (5.544,3.624) (5.576,3.557) (5.605,3.488) (5.631,3.417) (5.654,3.346) (5.674,3.273) (5.690,3.200) (5.703,3.126) (5.713,3.052) (5.719,2.977) (5.723,2.902) (5.723,2.827) (5.720,2.752) (5.714,2.678) (5.705,2.603) (5.693,2.529) (5.678,2.456) (5.660,2.383) (5.640,2.311) (5.617,2.239) (5.592,2.169) (5.564,2.099) (5.533,2.030) (5.501,1.963) (5.466,1.896) (5.430,1.831) (5.391,1.767) (5.350,1.704) (5.308,1.642) (5.263,1.581) (5.217,1.522) (5.170,1.464) (5.121,1.408) (5.070,1.352) (5.018,1.298) (4.965,1.246) (4.910,1.194) (4.854,1.144) (4.797,1.095) (4.739,1.048) (4.680,1.002) (4.620,0.957) (4.559,0.913) (4.497,0.871) (4.434,0.830) (4.370,0.790) (4.306,0.752) (4.241,0.715) (4.175,0.679) (4.153,0.667)}
\LC{(-10.000,1.452) (-9.943,1.420) (-9.875,1.389) (-9.807,1.358) (-9.738,1.328) (-9.669,1.298) (-9.600,1.269) (-9.531,1.240) (-9.461,1.211) (-9.392,1.183) (-9.322,1.156) (-9.252,1.129) (-9.182,1.102) (-9.112,1.076) (-9.042,1.050) (-8.971,1.024) (-8.900,0.999) (-8.829,0.974) (-8.759,0.950) (-8.687,0.926) (-8.616,0.903) (-8.545,0.880) (-8.473,0.857) (-8.402,0.835) (-8.330,0.813) (-8.258,0.791) (-8.186,0.770) (-8.114,0.749) (-8.042,0.729) (-7.970,0.709) (-7.897,0.689) (-7.825,0.670) (-7.752,0.651) (-7.680,0.633) (-7.607,0.614) (-7.534,0.597) (-7.461,0.579) (-7.388,0.562) (-7.315,0.545) (-7.242,0.529) (-7.169,0.512) (-7.095,0.497) (-7.022,0.481) (-6.948,0.466) (-6.875,0.451) (-6.801,0.437) (-6.728,0.422) (-6.654,0.408) (-6.580,0.395) (-6.506,0.382) (-6.433,0.369) (-6.359,0.356) (-6.285,0.343) (-6.211,0.331) (-6.137,0.320) (-6.062,0.308) (-5.988,0.297) (-5.914,0.286) (-5.840,0.275) (-5.766,0.265) (-5.691,0.254) (-5.617,0.245) (-5.543,0.235) (-5.468,0.225) (-5.394,0.216) (-5.319,0.207) (-5.245,0.199) (-5.170,0.190) (-5.096,0.182) (-5.021,0.174) (-4.946,0.167) (-4.872,0.159) (-4.797,0.152) (-4.723,0.145) (-4.698,0.143)}
\LC{(1.673,-5.217) (1.718,-5.277) (1.763,-5.337) (1.809,-5.396) (1.856,-5.455) (1.903,-5.514) (1.950,-5.571) (1.999,-5.629) (2.047,-5.686) (2.097,-5.742) (2.147,-5.798) (2.197,-5.854) (2.248,-5.909) (2.300,-5.963) (2.352,-6.017) (2.404,-6.071) (2.458,-6.123) (2.511,-6.176) (2.566,-6.228) (2.620,-6.279) (2.676,-6.329) (2.732,-6.380) (2.788,-6.429) (2.845,-6.478) (2.902,-6.526) (2.960,-6.574) (3.019,-6.621) (3.077,-6.667) (3.137,-6.713) (3.197,-6.758) (3.257,-6.803) (3.318,-6.846) (3.380,-6.889) (3.441,-6.932) (3.504,-6.973) (3.567,-7.014) (3.630,-7.055) (3.694,-7.094) (3.758,-7.133) (3.823,-7.171) (3.888,-7.208) (3.953,-7.244) (4.019,-7.280) (4.086,-7.315) (4.153,-7.349) (4.220,-7.382) (4.287,-7.414) (4.356,-7.446) (4.424,-7.477) (4.493,-7.506) (4.562,-7.535) (4.632,-7.563) (4.702,-7.590) (4.772,-7.616) (4.843,-7.641) (4.914,-7.666) (4.985,-7.689) (5.056,-7.711) (5.128,-7.732) (5.201,-7.753) (5.273,-7.772) (5.346,-7.790) (5.419,-7.807) (5.492,-7.823) (5.566,-7.838) (5.639,-7.852) (5.713,-7.865) (5.787,-7.877) (5.862,-7.887) (5.936,-7.897) (6.011,-7.905) (6.085,-7.912) (6.160,-7.918) (6.235,-7.922) (6.310,-7.926) (6.385,-7.928) (6.460,-7.929) (6.535,-7.928) (6.610,-7.927) (6.685,-7.924) (6.760,-7.919) (6.834,-7.914) (6.909,-7.907) (6.983,-7.898) (7.058,-7.889) (7.132,-7.877) (7.206,-7.865) (7.280,-7.851) (7.353,-7.836) (7.426,-7.819) (7.499,-7.801) (7.571,-7.781) (7.643,-7.760) (7.715,-7.737) (7.786,-7.713) (7.856,-7.687) (7.926,-7.660) (7.996,-7.632) (8.064,-7.602) (8.132,-7.570) (8.200,-7.537) (8.266,-7.502) (8.332,-7.466) (8.397,-7.429) (8.461,-7.389) (8.524,-7.349) (8.586,-7.307) (8.647,-7.263) (8.707,-7.218) (8.766,-7.172) (8.823,-7.124) (8.880,-7.074) (8.935,-7.024) (8.989,-6.971) (9.041,-6.918) (9.093,-6.863) (9.142,-6.807) (9.191,-6.750) (9.238,-6.691) (9.283,-6.631) (9.327,-6.570) (9.369,-6.508) (9.409,-6.445) (9.448,-6.381) (9.485,-6.316) (9.520,-6.249) (9.553,-6.182) (9.585,-6.114) (9.615,-6.045) (9.643,-5.976) (9.669,-5.906) (9.694,-5.835) (9.716,-5.763) (9.736,-5.691) (9.755,-5.618) (9.772,-5.545) (9.787,-5.472) (9.800,-5.398) (9.811,-5.324) (9.820,-5.249) (9.827,-5.175) (9.833,-5.100) (9.836,-5.025) (9.838,-4.950) (9.838,-4.875) (9.836,-4.800) (9.832,-4.725) (9.826,-4.650) (9.819,-4.576) (9.810,-4.501) (9.799,-4.427) (9.787,-4.353) (9.773,-4.279) (9.757,-4.206) (9.739,-4.133) (9.720,-4.060) (9.700,-3.988) (9.678,-3.917) (9.654,-3.845) (9.629,-3.775) (9.603,-3.705) (9.575,-3.635) (9.546,-3.566) (9.515,-3.497) (9.483,-3.429) (9.450,-3.362) (9.416,-3.296) (9.380,-3.230) (9.343,-3.164) (9.305,-3.100) (9.266,-3.036) (9.226,-2.973) (9.184,-2.910) (9.142,-2.848) (9.098,-2.787) (9.054,-2.727) (9.008,-2.667) (8.962,-2.608) (8.914,-2.550) (8.866,-2.493) (8.817,-2.436) (8.767,-2.380) (8.716,-2.325) (8.664,-2.271) (8.612,-2.217) (8.558,-2.165) (8.504,-2.113) (8.450,-2.061) (8.394,-2.011) (8.338,-1.961) (8.281,-1.912) (8.224,-1.864) (8.165,-1.817) (8.107,-1.770) (8.047,-1.724) (7.987,-1.679) (7.927,-1.635) (7.866,-1.591) (7.804,-1.548) (7.742,-1.506) (7.680,-1.465) (7.617,-1.424) (7.553,-1.384) (7.489,-1.345) (7.425,-1.306) (7.360,-1.269) (7.295,-1.232) (7.229,-1.195) (7.163,-1.160) (7.097,-1.125) (7.030,-1.091) (6.963,-1.057) (6.896,-1.024) (6.828,-0.992) (6.760,-0.960) (6.692,-0.929) (6.623,-0.899) (6.554,-0.870) (6.485,-0.841) (6.415,-0.813) (6.346,-0.785) (6.276,-0.758) (6.206,-0.731) (6.135,-0.706) (6.064,-0.681) (5.994,-0.656) (5.923,-0.632) (5.851,-0.609) (5.780,-0.586) (5.708,-0.564) (5.636,-0.542) (5.564,-0.521) (5.492,-0.500) (5.420,-0.480) (5.348,-0.461) (5.275,-0.442) (5.202,-0.423) (5.130,-0.405) (5.057,-0.388) (4.983,-0.371) (4.910,-0.355) (4.837,-0.339) (4.764,-0.323) (4.690,-0.308) (4.617,-0.294) (4.543,-0.280) (4.469,-0.266) (4.395,-0.253) (4.321,-0.240) (4.272,-0.232)}
\LC{(-0.492,3.275) (-0.514,3.347) (-0.537,3.418) (-0.561,3.490) (-0.585,3.561) (-0.609,3.632) (-0.634,3.702) (-0.659,3.773) (-0.685,3.843) (-0.711,3.913) (-0.738,3.984) (-0.766,4.053) (-0.793,4.123) (-0.822,4.192) (-0.850,4.262) (-0.880,4.331) (-0.910,4.400) (-0.940,4.468) (-0.971,4.537) (-1.002,4.605) (-1.034,4.673) (-1.066,4.740) (-1.099,4.808) (-1.132,4.875) (-1.166,4.942) (-1.200,5.009) (-1.234,5.075) (-1.270,5.142) (-1.305,5.208) (-1.341,5.273) (-1.378,5.339) (-1.415,5.404) (-1.452,5.469) (-1.491,5.534) (-1.529,5.598) (-1.568,5.662) (-1.608,5.726) (-1.647,5.789) (-1.688,5.852) (-1.729,5.915) (-1.770,5.978) (-1.812,6.040) (-1.854,6.102) (-1.897,6.164) (-1.940,6.225) (-1.984,6.286) (-2.028,6.346) (-2.073,6.407) (-2.118,6.466) (-2.164,6.526) (-2.210,6.585) (-2.256,6.644) (-2.303,6.702) (-2.351,6.761) (-2.399,6.818) (-2.447,6.876) (-2.496,6.932) (-2.545,6.989) (-2.595,7.045) (-2.645,7.101) (-2.696,7.156) (-2.747,7.211) (-2.799,7.266) (-2.851,7.320) (-2.903,7.373) (-2.956,7.426) (-3.009,7.479) (-3.063,7.531) (-3.117,7.583) (-3.172,7.635) (-3.227,7.686) (-3.282,7.736) (-3.338,7.786) (-3.395,7.835) (-3.452,7.884) (-3.509,7.933) (-3.566,7.981) (-3.624,8.028) (-3.683,8.075) (-3.742,8.122) (-3.801,8.168) (-3.861,8.213) (-3.921,8.258) (-3.982,8.302) (-4.042,8.346) (-4.104,8.389) (-4.166,8.432) (-4.228,8.474) (-4.290,8.515) (-4.353,8.556) (-4.416,8.596) (-4.480,8.636) (-4.544,8.675) (-4.608,8.713) (-4.673,8.751) (-4.738,8.788) (-4.804,8.825) (-4.870,8.861) (-4.936,8.896) (-5.003,8.930) (-5.069,8.964) (-5.137,8.997) (-5.204,9.030) (-5.272,9.062) (-5.340,9.093) (-5.409,9.123) (-5.478,9.153) (-5.547,9.182) (-5.617,9.210) (-5.686,9.237) (-5.757,9.264) (-5.827,9.290) (-5.898,9.315) (-5.968,9.340) (-6.040,9.363) (-6.111,9.386) (-6.183,9.408) (-6.255,9.429) (-6.327,9.450) (-6.399,9.469) (-6.472,9.488) (-6.545,9.505) (-6.618,9.522) (-6.691,9.538) (-6.765,9.554) (-6.838,9.568) (-6.912,9.581) (-6.986,9.594) (-7.060,9.605) (-7.134,9.616) (-7.209,9.625) (-7.283,9.634) (-7.358,9.642) (-7.433,9.649) (-7.507,9.654) (-7.582,9.659) (-7.657,9.663) (-7.732,9.666) (-7.807,9.667) (-7.882,9.668) (-7.957,9.668) (-8.032,9.666) (-8.107,9.664) (-8.182,9.660) (-8.257,9.656) (-8.331,9.650) (-8.406,9.643) (-8.481,9.635) (-8.555,9.626) (-8.630,9.616) (-8.704,9.605) (-8.778,9.592) (-8.851,9.578) (-8.925,9.564) (-8.998,9.548) (-9.071,9.530) (-9.144,9.512) (-9.216,9.493) (-9.288,9.472) (-9.360,9.450) (-9.431,9.427) (-9.502,9.402) (-9.573,9.377) (-9.643,9.350) (-9.712,9.322) (-9.781,9.292) (-9.850,9.262) (-9.918,9.230) (-9.985,9.197) (-10.000,9.186)}
\LC{(-10.000,-1.331) (-9.939,-1.302) (-9.870,-1.273) (-9.801,-1.245) (-9.731,-1.217) (-9.661,-1.190) (-9.591,-1.163) (-9.521,-1.136) (-9.451,-1.110) (-9.380,-1.084) (-9.310,-1.059) (-9.239,-1.034) (-9.168,-1.009) (-9.097,-0.985) (-9.026,-0.961) (-8.955,-0.938) (-8.883,-0.915) (-8.812,-0.892) (-8.740,-0.870) (-8.669,-0.848) (-8.597,-0.826) (-8.525,-0.805) (-8.453,-0.784) (-8.381,-0.764) (-8.308,-0.744) (-8.236,-0.724) (-8.164,-0.704) (-8.091,-0.685) (-8.019,-0.667) (-7.946,-0.648) (-7.873,-0.630) (-7.800,-0.612) (-7.727,-0.595) (-7.654,-0.578) (-7.581,-0.561) (-7.508,-0.545) (-7.435,-0.529) (-7.361,-0.513) (-7.288,-0.498) (-7.214,-0.483) (-7.141,-0.468) (-7.067,-0.453) (-6.994,-0.439) (-6.920,-0.425) (-6.846,-0.411) (-6.772,-0.398) (-6.699,-0.385) (-6.625,-0.372) (-6.551,-0.360) (-6.477,-0.348) (-6.403,-0.336) (-6.328,-0.324)}
\LC{(1.709,4.492) (1.761,4.547) (1.813,4.600) (1.866,4.653) (1.920,4.706) (1.975,4.757) (2.030,4.808) (2.086,4.858) (2.142,4.907) (2.199,4.956) (2.257,5.003) (2.316,5.050) (2.375,5.096) (2.435,5.141) (2.496,5.186) (2.557,5.229) (2.619,5.271) (2.681,5.313) (2.744,5.354) (2.808,5.393) (2.872,5.432) (2.937,5.469) (3.003,5.506) (3.069,5.542) (3.135,5.576) (3.202,5.609) (3.270,5.642) (3.338,5.673) (3.407,5.703) (3.476,5.732) (3.546,5.760) (3.616,5.786) (3.687,5.811) (3.758,5.835) (3.829,5.858) (3.901,5.879) (3.973,5.899) (4.046,5.918) (4.119,5.935) (4.192,5.951) (4.266,5.966) (4.340,5.979) (4.414,5.990) (4.488,6.000) (4.563,6.009) (4.637,6.016) (4.712,6.021) (4.787,6.025) (4.862,6.027) (4.937,6.028) (5.012,6.027) (5.087,6.024) (5.162,6.019) (5.236,6.013) (5.311,6.004) (5.385,5.995) (5.459,5.983) (5.533,5.969) (5.607,5.953) (5.679,5.936) (5.752,5.917) (5.824,5.895) (5.895,5.872) (5.966,5.847) (6.036,5.820) (6.105,5.791) (6.173,5.760) (6.240,5.727) (6.307,5.692) (6.372,5.655) (6.436,5.616) (6.499,5.575) (6.560,5.532) (6.620,5.487) (6.679,5.440) (6.736,5.391) (6.792,5.341) (6.845,5.289) (6.897,5.235) (6.947,5.179) (6.996,5.121) (7.042,5.062) (7.086,5.002) (7.128,4.940) (7.168,4.876) (7.205,4.811) (7.241,4.745) (7.274,4.678) (7.304,4.609) (7.333,4.540) (7.358,4.469) (7.382,4.398) (7.402,4.326) (7.421,4.253) (7.437,4.180) (7.450,4.106) (7.461,4.032) (7.469,3.957) (7.475,3.883) (7.478,3.808) (7.479,3.733) (7.478,3.658) (7.474,3.583) (7.468,3.508) (7.460,3.433) (7.449,3.359) (7.436,3.285) (7.421,3.212) (7.404,3.139) (7.384,3.066) (7.363,2.995) (7.340,2.923) (7.315,2.853) (7.287,2.783) (7.258,2.714) (7.228,2.645) (7.195,2.577) (7.161,2.511) (7.125,2.445) (7.088,2.380) (7.049,2.316) (7.009,2.252) (6.967,2.190) (6.924,2.129) (6.880,2.068) (6.834,2.009) (6.787,1.950) (6.739,1.893) (6.689,1.837) (6.639,1.781) (6.587,1.727) (6.534,1.673) (6.481,1.621) (6.426,1.569) (6.371,1.519) (6.314,1.470) (6.257,1.421) (6.199,1.374) (6.140,1.328) (6.080,1.282) (6.020,1.238) (5.958,1.194) (5.897,1.152) (5.834,1.111) (5.771,1.070) (5.707,1.031) (5.686,1.018)}
\LC{(0.123,-3.792) (0.128,-3.866) (0.133,-3.941) (0.138,-4.016) (0.144,-4.091) (0.149,-4.166) (0.154,-4.241) (0.160,-4.315) (0.165,-4.390) (0.171,-4.465) (0.177,-4.540) (0.183,-4.614) (0.189,-4.689) (0.195,-4.764) (0.201,-4.839) (0.207,-4.913) (0.214,-4.988) (0.220,-5.063) (0.227,-5.138) (0.233,-5.212) (0.240,-5.287) (0.247,-5.362) (0.254,-5.436) (0.261,-5.511) (0.268,-5.586) (0.275,-5.660) (0.283,-5.735) (0.290,-5.810) (0.298,-5.884) (0.305,-5.959) (0.313,-6.033) (0.321,-6.108) (0.329,-6.183) (0.337,-6.257) (0.345,-6.332) (0.353,-6.406) (0.361,-6.481) (0.370,-6.555) (0.378,-6.630) (0.387,-6.704) (0.395,-6.779) (0.404,-6.853) (0.413,-6.928) (0.422,-7.002) (0.431,-7.077) (0.440,-7.151) (0.449,-7.226) (0.459,-7.300) (0.468,-7.374) (0.478,-7.449) (0.487,-7.523) (0.497,-7.598) (0.507,-7.672) (0.517,-7.746) (0.527,-7.821) (0.537,-7.895) (0.547,-7.969) (0.558,-8.044) (0.568,-8.118) (0.579,-8.192) (0.589,-8.266) (0.600,-8.341) (0.611,-8.415) (0.621,-8.489) (0.632,-8.563) (0.644,-8.637) (0.655,-8.712) (0.666,-8.786) (0.677,-8.860) (0.689,-8.934) (0.700,-9.008) (0.712,-9.082) (0.724,-9.156) (0.736,-9.230) (0.748,-9.304) (0.760,-9.378) (0.772,-9.452) (0.784,-9.526) (0.796,-9.600) (0.809,-9.674) (0.821,-9.748) (0.834,-9.822) (0.847,-9.896) (0.859,-9.970) (0.868,-10.000)}
\LC{(-0.270,-4.625) (-0.278,-4.699) (-0.287,-4.774) (-0.296,-4.848) (-0.306,-4.923) (-0.315,-4.997) (-0.324,-5.071) (-0.334,-5.146) (-0.344,-5.220) (-0.354,-5.295) (-0.364,-5.369) (-0.374,-5.443) (-0.384,-5.517) (-0.395,-5.592) (-0.405,-5.666) (-0.416,-5.740) (-0.427,-5.814) (-0.438,-5.889) (-0.449,-5.963) (-0.461,-6.037) (-0.472,-6.111) (-0.484,-6.185) (-0.495,-6.259) (-0.507,-6.333) (-0.519,-6.407) (-0.532,-6.481) (-0.544,-6.555) (-0.556,-6.629) (-0.569,-6.703) (-0.582,-6.777) (-0.594,-6.851) (-0.608,-6.925) (-0.621,-6.999) (-0.634,-7.072) (-0.647,-7.146) (-0.661,-7.220) (-0.675,-7.294) (-0.689,-7.367) (-0.703,-7.441) (-0.717,-7.515) (-0.731,-7.588) (-0.745,-7.662) (-0.760,-7.736) (-0.775,-7.809) (-0.789,-7.883) (-0.804,-7.956) (-0.820,-8.030) (-0.835,-8.103) (-0.850,-8.176) (-0.866,-8.250) (-0.881,-8.323) (-0.897,-8.396) (-0.913,-8.470) (-0.929,-8.543) (-0.946,-8.616) (-0.962,-8.689) (-0.979,-8.763) (-0.995,-8.836) (-1.012,-8.909) (-1.029,-8.982) (-1.046,-9.055) (-1.063,-9.128) (-1.081,-9.201) (-1.098,-9.274) (-1.116,-9.347) (-1.134,-9.419) (-1.151,-9.492) (-1.170,-9.565) (-1.188,-9.638) (-1.206,-9.711) (-1.225,-9.783) (-1.243,-9.856) (-1.262,-9.929) (-1.281,-10.000)}
\LC{(1.386,-3.828) (1.435,-3.884) (1.486,-3.940) (1.537,-3.994) (1.589,-4.048) (1.642,-4.101) (1.696,-4.153) (1.751,-4.205) (1.806,-4.256) (1.862,-4.305) (1.919,-4.354) (1.977,-4.402) (2.035,-4.449) (2.094,-4.495) (2.154,-4.540) (2.215,-4.584) (2.276,-4.627) (2.339,-4.669) (2.401,-4.710) (2.465,-4.750) (2.529,-4.789) (2.594,-4.827) (2.659,-4.863) (2.726,-4.899) (2.792,-4.933) (2.860,-4.965) (2.928,-4.997) (2.997,-5.027) (3.066,-5.056) (3.135,-5.084) (3.206,-5.110) (3.276,-5.135) (3.348,-5.158) (3.419,-5.180) (3.492,-5.201) (3.564,-5.220) (3.637,-5.237) (3.710,-5.253) (3.784,-5.267) (3.858,-5.279) (3.932,-5.290) (4.007,-5.299) (4.081,-5.306) (4.156,-5.312) (4.231,-5.316) (4.306,-5.318) (4.381,-5.317) (4.456,-5.316) (4.531,-5.312) (4.606,-5.306) (4.680,-5.298) (4.755,-5.288) (4.829,-5.276) (4.902,-5.262) (4.975,-5.245) (5.048,-5.227) (5.120,-5.206) (5.192,-5.184) (5.262,-5.158) (5.332,-5.131) (5.401,-5.102) (5.469,-5.070) (5.536,-5.036) (5.602,-5.000) (5.666,-4.961) (5.729,-4.921) (5.791,-4.878) (5.851,-4.833) (5.909,-4.785) (5.965,-4.736) (6.020,-4.685) (6.072,-4.631) (6.123,-4.576) (6.171,-4.518) (6.217,-4.459) (6.261,-4.398) (6.302,-4.336) (6.341,-4.272) (6.378,-4.206) (6.411,-4.139) (6.442,-4.071) (6.470,-4.001) (6.496,-3.931) (6.518,-3.859) (6.538,-3.787) (6.555,-3.714) (6.569,-3.640) (6.581,-3.566) (6.589,-3.491) (6.595,-3.417) (6.598,-3.342) (6.598,-3.267) (6.596,-3.192) (6.590,-3.117) (6.582,-3.042) (6.572,-2.968) (6.559,-2.894) (6.544,-2.821) (6.526,-2.748) (6.506,-2.676) (6.484,-2.604) (6.459,-2.533) (6.432,-2.463) (6.403,-2.394) (6.373,-2.326) (6.340,-2.258) (6.305,-2.192) (6.269,-2.126) (6.230,-2.062) (6.191,-1.998) (6.149,-1.936) (6.106,-1.874) (6.061,-1.814) (6.015,-1.755) (5.968,-1.697) (5.919,-1.640) (5.869,-1.584) (5.817,-1.529) (5.765,-1.476) (5.711,-1.424) (5.656,-1.372) (5.600,-1.322) (5.543,-1.274) (5.486,-1.226) (5.427,-1.179) (5.367,-1.134) (5.307,-1.089) (5.245,-1.046) (5.183,-1.004) (5.121,-0.963) (5.057,-0.923) (4.993,-0.884) (4.928,-0.846) (4.863,-0.809)}
\LC{(-2.350,6.023) (-2.402,6.076) (-2.456,6.129) (-2.509,6.182) (-2.563,6.234) (-2.618,6.285) (-2.673,6.336) (-2.729,6.386) (-2.785,6.435) (-2.842,6.484) (-2.900,6.533) (-2.957,6.580) (-3.016,6.628) (-3.075,6.674) (-3.134,6.720) (-3.194,6.765) (-3.254,6.810) (-3.315,6.853) (-3.376,6.897) (-3.438,6.939) (-3.500,6.981) (-3.563,7.022) (-3.626,7.062) (-3.690,7.102) (-3.754,7.141) (-3.819,7.179) (-3.884,7.216) (-3.949,7.253) (-4.015,7.289) (-4.081,7.324) (-4.148,7.358) (-4.215,7.391) (-4.283,7.424) (-4.351,7.455) (-4.419,7.486) (-4.488,7.516) (-4.557,7.545) (-4.627,7.573) (-4.697,7.600) (-4.767,7.627) (-4.838,7.652) (-4.908,7.677) (-4.980,7.700) (-5.051,7.723) (-5.123,7.744) (-5.195,7.765) (-5.268,7.784) (-5.340,7.802) (-5.413,7.820) (-5.486,7.836) (-5.560,7.851) (-5.634,7.866) (-5.707,7.879) (-5.781,7.891) (-5.856,7.901) (-5.930,7.911) (-6.005,7.920) (-6.079,7.927) (-6.154,7.933) (-6.229,7.938) (-6.304,7.942) (-6.379,7.944) (-6.454,7.945) (-6.529,7.945) (-6.604,7.944) (-6.679,7.941) (-6.753,7.937) (-6.828,7.932) (-6.903,7.925) (-6.978,7.917) (-7.052,7.908) (-7.126,7.897) (-7.200,7.885) (-7.274,7.872) (-7.348,7.857) (-7.421,7.840) (-7.494,7.823) (-7.566,7.803) (-7.638,7.783) (-7.710,7.760) (-7.781,7.737) (-7.852,7.712) (-7.922,7.685) (-7.991,7.657) (-8.060,7.627) (-8.128,7.596) (-8.196,7.563) (-8.263,7.529) (-8.329,7.494) (-8.394,7.456) (-8.458,7.418) (-8.521,7.378) (-8.584,7.336) (-8.645,7.293) (-8.705,7.248) (-8.765,7.202) (-8.823,7.155) (-8.880,7.106) (-8.935,7.056) (-8.990,7.004) (-9.043,6.951) (-9.094,6.897) (-9.145,6.841) (-9.193,6.784) (-9.241,6.726) (-9.287,6.667) (-9.331,6.606) (-9.374,6.544) (-9.415,6.481) (-9.454,6.418) (-9.491,6.353) (-9.527,6.287) (-9.561,6.220) (-9.594,6.152) (-9.624,6.084) (-9.653,6.015) (-9.680,5.945) (-9.705,5.874) (-9.728,5.803) (-9.749,5.731) (-9.769,5.658) (-9.786,5.585) (-9.802,5.512) (-9.816,5.438) (-9.827,5.364) (-9.837,5.290) (-9.845,5.215) (-9.851,5.140) (-9.856,5.066) (-9.858,4.991) (-9.859,4.916) (-9.857,4.841) (-9.854,4.766) (-9.850,4.691) (-9.843,4.616) (-9.835,4.542) (-9.825,4.467) (-9.813,4.393) (-9.799,4.319) (-9.784,4.246) (-9.768,4.173) (-9.749,4.100) (-9.729,4.028) (-9.708,3.956) (-9.685,3.884) (-9.661,3.814) (-9.635,3.743) (-9.608,3.673) (-9.579,3.604) (-9.549,3.535) (-9.518,3.467) (-9.485,3.399) (-9.451,3.333) (-9.416,3.266) (-9.380,3.201) (-9.342,3.136) (-9.304,3.072) (-9.264,3.008) (-9.223,2.945) (-9.181,2.883) (-9.138,2.822) (-9.094,2.761) (-9.049,2.701) (-9.003,2.642) (-8.956,2.583) (-8.908,2.525) (-8.859,2.468) (-8.810,2.412) (-8.759,2.357) (-8.708,2.302) (-8.656,2.248) (-8.603,2.195) (-8.549,2.143) (-8.495,2.091) (-8.439,2.040) (-8.384,1.990) (-8.327,1.941) (-8.270,1.892) (-8.212,1.844) (-8.154,1.797) (-8.095,1.751) (-8.035,1.706) (-7.975,1.661) (-7.914,1.617) (-7.853,1.574) (-7.791,1.531) (-7.729,1.489) (-7.666,1.448) (-7.603,1.408) (-7.539,1.368) (-7.475,1.330) (-7.410,1.291) (-7.345,1.254) (-7.280,1.217) (-7.214,1.181) (-7.148,1.146) (-7.082,1.111) (-7.015,1.077) (-6.947,1.044) (-6.880,1.012) (-6.812,0.980) (-6.744,0.949) (-6.675,0.918) (-6.607,0.888) (-6.537,0.859) (-6.491,0.840)}
\LC{(10.000,-1.736) (9.939,-1.698) (9.874,-1.662) (9.808,-1.625) (9.742,-1.590) (9.676,-1.554) (9.609,-1.519) (9.543,-1.485) (9.476,-1.451) (9.408,-1.418) (9.341,-1.385) (9.273,-1.353) (9.205,-1.321) (9.137,-1.290) (9.069,-1.259) (9.000,-1.229) (8.931,-1.199) (8.862,-1.170) (8.793,-1.141) (8.724,-1.113) (8.654,-1.085) (8.584,-1.057) (8.514,-1.030) (8.444,-1.004) (8.374,-0.978) (8.303,-0.952) (8.233,-0.927) (8.162,-0.902) (8.091,-0.878) (8.020,-0.854) (7.949,-0.831) (7.877,-0.808) (7.806,-0.785) (7.734,-0.763) (7.662,-0.741) (7.590,-0.720) (7.518,-0.699) (7.446,-0.679) (7.374,-0.658) (7.302,-0.639) (7.229,-0.619) (7.156,-0.600) (7.084,-0.582) (7.011,-0.564) (6.938,-0.546)}
\LC{(4.710,-8.828) (4.775,-8.865) (4.841,-8.902) (4.907,-8.938) (4.973,-8.973) (5.039,-9.008) (5.106,-9.042) (5.173,-9.075) (5.241,-9.107) (5.309,-9.139) (5.377,-9.171) (5.445,-9.201) (5.514,-9.231) (5.583,-9.260) (5.653,-9.288) (5.723,-9.316) (5.793,-9.343) (5.863,-9.369) (5.933,-9.394) (6.004,-9.419) (6.075,-9.443) (6.147,-9.466) (6.218,-9.488) (6.290,-9.509) (6.362,-9.530) (6.435,-9.550) (6.507,-9.569) (6.580,-9.587) (6.653,-9.604) (6.726,-9.620) (6.800,-9.636) (6.873,-9.650) (6.947,-9.664) (7.021,-9.677) (7.095,-9.689) (7.169,-9.700) (7.243,-9.710) (7.318,-9.719) (7.392,-9.727) (7.467,-9.734) (7.542,-9.741) (7.617,-9.746) (7.692,-9.750) (7.766,-9.753) (7.841,-9.755) (7.916,-9.757) (7.991,-9.757) (8.066,-9.756) (8.141,-9.754) (8.216,-9.751) (8.291,-9.747) (8.366,-9.742) (8.441,-9.736) (8.515,-9.728) (8.590,-9.720) (8.664,-9.710) (8.739,-9.700) (8.813,-9.688) (8.887,-9.675) (8.960,-9.661) (9.034,-9.646) (9.107,-9.629) (9.180,-9.612) (9.252,-9.593) (9.325,-9.573) (9.397,-9.552) (9.468,-9.529) (9.539,-9.506) (9.610,-9.481) (9.681,-9.455) (9.750,-9.428) (9.820,-9.400) (9.889,-9.370) (9.957,-9.339) (10.000,-9.318)}
\LC{(2.598,-5.621) (2.657,-5.667) (2.718,-5.712) (2.778,-5.756) (2.840,-5.799) (2.901,-5.842) (2.964,-5.883) (3.027,-5.924) (3.090,-5.964) (3.154,-6.003) (3.219,-6.041) (3.284,-6.078) (3.350,-6.114) (3.416,-6.149) (3.483,-6.184) (3.550,-6.217) (3.618,-6.249) (3.686,-6.281) (3.754,-6.311) (3.823,-6.340) (3.893,-6.369) (3.963,-6.396) (4.033,-6.422) (4.104,-6.447) (4.175,-6.470) (4.247,-6.493) (4.318,-6.514) (4.391,-6.534) (4.463,-6.553) (4.536,-6.571) (4.609,-6.588) (4.683,-6.603) (4.756,-6.617) (4.830,-6.629) (4.904,-6.641) (4.979,-6.651) (5.053,-6.659) (5.128,-6.666) (5.203,-6.672) (5.278,-6.676) (5.353,-6.679) (5.428,-6.680) (5.503,-6.680) (5.578,-6.678) (5.652,-6.675) (5.727,-6.670) (5.802,-6.664) (5.877,-6.656) (5.951,-6.646) (6.025,-6.635) (6.099,-6.622) (6.172,-6.607) (6.246,-6.591) (6.318,-6.572) (6.391,-6.553) (6.463,-6.531) (6.534,-6.508) (6.605,-6.483) (6.675,-6.456) (6.744,-6.427) (6.812,-6.397) (6.880,-6.364) (6.947,-6.330) (7.013,-6.295) (7.078,-6.257) (7.142,-6.218) (7.204,-6.177) (7.266,-6.134) (7.326,-6.089) (7.385,-6.043) (7.443,-5.995) (7.499,-5.945) (7.554,-5.894) (7.607,-5.841) (7.658,-5.786) (7.708,-5.730) (7.756,-5.672) (7.802,-5.613) (7.846,-5.553) (7.889,-5.491) (7.929,-5.428) (7.968,-5.363) (8.004,-5.298) (8.038,-5.231) (8.071,-5.163) (8.101,-5.095) (8.129,-5.025) (8.154,-4.955) (8.178,-4.883) (8.199,-4.812) (8.218,-4.739) (8.235,-4.666) (8.249,-4.592) (8.261,-4.518) (8.271,-4.444) (8.279,-4.369) (8.285,-4.295) (8.288,-4.220) (8.289,-4.145) (8.288,-4.070) (8.285,-3.995) (8.279,-3.920) (8.272,-3.845) (8.262,-3.771) (8.251,-3.697) (8.237,-3.623) (8.222,-3.550) (8.205,-3.477) (8.185,-3.404) (8.164,-3.332) (8.141,-3.261) (8.117,-3.190) (8.090,-3.120) (8.062,-3.050) (8.033,-2.981) (8.002,-2.913) (7.969,-2.846) (7.935,-2.779) (7.899,-2.713) (7.862,-2.648) (7.823,-2.583) (7.783,-2.520) (7.742,-2.457) (7.700,-2.395) (7.656,-2.335) (7.611,-2.274) (7.565,-2.215) (7.518,-2.157) (7.469,-2.100) (7.420,-2.043) (7.370,-1.987) (7.319,-1.933) (7.266,-1.879) (7.213,-1.826) (7.159,-1.774) (7.104,-1.723) (7.048,-1.673) (6.992,-1.624) (6.934,-1.575) (6.876,-1.528) (6.857,-1.512)}
\LC{(-0.778,5.179) (-0.800,5.251) (-0.823,5.322) (-0.846,5.394) (-0.869,5.465) (-0.893,5.536) (-0.916,5.608) (-0.941,5.679) (-0.965,5.749) (-0.990,5.820) (-1.015,5.891) (-1.041,5.961) (-1.066,6.032) (-1.092,6.102) (-1.119,6.172) (-1.145,6.242) (-1.172,6.312) (-1.200,6.382) (-1.227,6.452) (-1.255,6.522) (-1.284,6.591) (-1.312,6.660) (-1.341,6.730) (-1.370,6.799) (-1.400,6.868) (-1.429,6.937) (-1.459,7.005) (-1.490,7.074) (-1.521,7.142) (-1.552,7.211) (-1.583,7.279) (-1.615,7.347) (-1.647,7.415) (-1.679,7.482) (-1.711,7.550) (-1.744,7.617) (-1.777,7.684) (-1.811,7.752) (-1.845,7.818) (-1.879,7.885) (-1.913,7.952) (-1.948,8.018) (-1.983,8.085) (-2.018,8.151) (-2.054,8.217) (-2.090,8.283) (-2.126,8.348) (-2.163,8.414) (-2.200,8.479) (-2.237,8.544) (-2.274,8.609) (-2.312,8.674) (-2.350,8.739) (-2.388,8.803) (-2.427,8.867) (-2.466,8.931) (-2.505,8.995) (-2.545,9.059) (-2.585,9.123) (-2.625,9.186) (-2.666,9.249) (-2.706,9.312) (-2.747,9.375) (-2.789,9.437) (-2.830,9.500) (-2.872,9.562) (-2.915,9.624) (-2.957,9.685) (-3.000,9.747) (-3.043,9.808) (-3.087,9.869) (-3.131,9.930) (-3.175,9.991) (-3.190,10.000)}
\LC{(-3.035,-7.115) (-3.091,-7.164) (-3.147,-7.213) (-3.204,-7.262) (-3.262,-7.310) (-3.320,-7.358) (-3.379,-7.405) (-3.437,-7.451) (-3.497,-7.497) (-3.557,-7.542) (-3.617,-7.587) (-3.678,-7.631) (-3.739,-7.674) (-3.800,-7.717) (-3.862,-7.759) (-3.925,-7.801) (-3.988,-7.841) (-4.051,-7.882) (-4.115,-7.921) (-4.179,-7.960) (-4.243,-7.999) (-4.308,-8.036) (-4.374,-8.073) (-4.439,-8.109) (-4.505,-8.145) (-4.572,-8.179) (-4.639,-8.213) (-4.706,-8.247) (-4.773,-8.279) (-4.841,-8.311) (-4.910,-8.342) (-4.978,-8.372) (-5.047,-8.401) (-5.117,-8.430) (-5.186,-8.458) (-5.256,-8.485) (-5.327,-8.511) (-5.397,-8.536) (-5.468,-8.561) (-5.539,-8.584) (-5.611,-8.607) (-5.683,-8.629) (-5.755,-8.650) (-5.827,-8.670) (-5.899,-8.689) (-5.972,-8.707) (-6.045,-8.724) (-6.118,-8.741) (-6.192,-8.756) (-6.265,-8.770) (-6.339,-8.784) (-6.413,-8.796) (-6.487,-8.808) (-6.562,-8.818) (-6.636,-8.827) (-6.710,-8.835) (-6.785,-8.843) (-6.860,-8.849) (-6.935,-8.854) (-7.010,-8.858) (-7.085,-8.861) (-7.160,-8.862) (-7.235,-8.863) (-7.310,-8.862) (-7.384,-8.861) (-7.459,-8.858) (-7.534,-8.854) (-7.609,-8.849) (-7.684,-8.842) (-7.758,-8.835) (-7.833,-8.826) (-7.907,-8.816) (-7.981,-8.804) (-8.055,-8.792) (-8.129,-8.778) (-8.203,-8.763) (-8.276,-8.746) (-8.349,-8.729) (-8.421,-8.710) (-8.493,-8.689) (-8.565,-8.668) (-8.637,-8.645) (-8.708,-8.621) (-8.778,-8.595) (-8.848,-8.568) (-8.917,-8.540) (-8.986,-8.510) (-9.055,-8.479) (-9.122,-8.447) (-9.189,-8.413) (-9.256,-8.378) (-9.321,-8.342) (-9.386,-8.304) (-9.450,-8.265) (-9.513,-8.225) (-9.576,-8.183) (-9.637,-8.140) (-9.698,-8.095) (-9.757,-8.050) (-9.816,-8.003) (-9.873,-7.955) (-9.929,-7.905) (-9.985,-7.854) (-10.000,-7.837)}
\LC{(-0.793,-5.153) (-0.816,-5.224) (-0.839,-5.296) (-0.863,-5.367) (-0.886,-5.438) (-0.910,-5.509) (-0.935,-5.580) (-0.960,-5.651) (-0.985,-5.721) (-1.010,-5.792) (-1.036,-5.862) (-1.062,-5.933) (-1.088,-6.003) (-1.115,-6.073) (-1.142,-6.143) (-1.169,-6.213) (-1.197,-6.283) (-1.225,-6.352) (-1.253,-6.422) (-1.282,-6.491) (-1.311,-6.560) (-1.340,-6.629) (-1.370,-6.698) (-1.400,-6.767) (-1.430,-6.835) (-1.460,-6.904) (-1.491,-6.972) (-1.522,-7.041) (-1.554,-7.109) (-1.586,-7.177) (-1.618,-7.244) (-1.650,-7.312) (-1.683,-7.379) (-1.716,-7.447) (-1.749,-7.514) (-1.783,-7.581) (-1.817,-7.648) (-1.851,-7.714) (-1.886,-7.781) (-1.921,-7.847) (-1.956,-7.914) (-1.992,-7.980) (-2.028,-8.045) (-2.064,-8.111) (-2.101,-8.177) (-2.137,-8.242) (-2.175,-8.307) (-2.212,-8.372) (-2.250,-8.437) (-2.288,-8.501) (-2.326,-8.566) (-2.365,-8.630) (-2.404,-8.694) (-2.443,-8.758) (-2.483,-8.822) (-2.523,-8.885) (-2.563,-8.948) (-2.604,-9.011) (-2.645,-9.074) (-2.686,-9.137) (-2.727,-9.199) (-2.769,-9.262) (-2.811,-9.324) (-2.854,-9.386) (-2.897,-9.447) (-2.940,-9.509) (-2.983,-9.570) (-3.027,-9.631) (-3.071,-9.692) (-3.115,-9.752) (-3.159,-9.812) (-3.204,-9.873) (-3.249,-9.932) (-3.295,-9.992) (-3.310,-10.000)}
\LC{(-1.651,4.550) (-1.701,4.606) (-1.751,4.662) (-1.802,4.717) (-1.854,4.771) (-1.906,4.824) (-1.960,4.877) (-2.013,4.930) (-2.068,4.981) (-2.123,5.032) (-2.179,5.082) (-2.235,5.131) (-2.293,5.180) (-2.350,5.228) (-2.409,5.275) (-2.468,5.321) (-2.527,5.366) (-2.588,5.411) (-2.649,5.455) (-2.710,5.498) (-2.772,5.540) (-2.835,5.581) (-2.898,5.621) (-2.962,5.660) (-3.027,5.698) (-3.092,5.736) (-3.157,5.772) (-3.224,5.808) (-3.290,5.842) (-3.357,5.875) (-3.425,5.907) (-3.493,5.939) (-3.562,5.969) (-3.631,5.998) (-3.701,6.026) (-3.771,6.052) (-3.842,6.078) (-3.912,6.102) (-3.984,6.125) (-4.056,6.147) (-4.128,6.168) (-4.200,6.187) (-4.273,6.205) (-4.346,6.222) (-4.419,6.237) (-4.493,6.251) (-4.567,6.263) (-4.641,6.275) (-4.716,6.284) (-4.790,6.293) (-4.865,6.299) (-4.940,6.305) (-5.015,6.308) (-5.090,6.310) (-5.165,6.311) (-5.240,6.310) (-5.315,6.307) (-5.389,6.303) (-5.464,6.297) (-5.539,6.289) (-5.613,6.280) (-5.687,6.268) (-5.761,6.255) (-5.835,6.241) (-5.908,6.224) (-5.981,6.206) (-6.053,6.186) (-6.125,6.164) (-6.196,6.140) (-6.266,6.114) (-6.336,6.086) (-6.405,6.057) (-6.473,6.026) (-6.540,5.992) (-6.606,5.957) (-6.671,5.920) (-6.736,5.881) (-6.798,5.840) (-6.860,5.797) (-6.920,5.753) (-6.979,5.706) (-7.037,5.658) (-7.093,5.608) (-7.147,5.557) (-7.200,5.503) (-7.251,5.448) (-7.300,5.392) (-7.347,5.333) (-7.392,5.273) (-7.436,5.212) (-7.477,5.150) (-7.516,5.086) (-7.553,5.020) (-7.588,4.954) (-7.620,4.886) (-7.651,4.818) (-7.679,4.748) (-7.704,4.678) (-7.728,4.607) (-7.749,4.535) (-7.767,4.462) (-7.783,4.389) (-7.797,4.315) (-7.809,4.241) (-7.818,4.166) (-7.824,4.092) (-7.829,4.017) (-7.830,3.942) (-7.830,3.867) (-7.828,3.792) (-7.823,3.717) (-7.816,3.642) (-7.806,3.568) (-7.795,3.494) (-7.781,3.420) (-7.766,3.347) (-7.748,3.274) (-7.729,3.201) (-7.707,3.130) (-7.684,3.058) (-7.658,2.988) (-7.631,2.918) (-7.603,2.848) (-7.572,2.780) (-7.540,2.712) (-7.506,2.645) (-7.471,2.579) (-7.434,2.514) (-7.396,2.449) (-7.356,2.386) (-7.315,2.323) (-7.272,2.261) (-7.229,2.200) (-7.184,2.140) (-7.137,2.081) (-7.090,2.023) (-7.041,1.966) (-6.992,1.910) (-6.941,1.854) (-6.889,1.800) (-6.836,1.747) (-6.783,1.695) (-6.728,1.643) (-6.673,1.593) (-6.616,1.543) (-6.559,1.495) (-6.501,1.447) (-6.442,1.401) (-6.383,1.355) (-6.322,1.310) (-6.261,1.267) (-6.200,1.224) (-6.138,1.182) (-6.075,1.141) (-6.011,1.101) (-5.947,1.062) (-5.883,1.024) (-5.818,0.987) (-5.752,0.950) (-5.686,0.915) (-5.620,0.880) (-5.553,0.846) (-5.485,0.813) (-5.417,0.781)}
\LC{(10.000,-0.013) (9.943,-0.013) (9.868,-0.013) (9.793,-0.012) (9.718,-0.012) (9.643,-0.012) (9.568,-0.012) (9.493,-0.011) (9.418,-0.011) (9.343,-0.011) (9.268,-0.011) (9.193,-0.010) (9.118,-0.010) (9.043,-0.010) (8.968,-0.010) (8.893,-0.009) (8.818,-0.009) (8.743,-0.009) (8.668,-0.009) (8.593,-0.008) (8.518,-0.008) (8.443,-0.008) (8.368,-0.008) (8.293,-0.008) (8.218,-0.007) (8.143,-0.007) (8.068,-0.007) (7.993,-0.007) (7.918,-0.007) (7.843,-0.006) (7.768,-0.006) (7.693,-0.006) (7.618,-0.006) (7.543,-0.006) (7.468,-0.006) (7.393,-0.005) (7.318,-0.005) (7.243,-0.005) (7.168,-0.005) (7.093,-0.005) (7.018,-0.005) (6.943,-0.004) (6.868,-0.004) (6.793,-0.004) (6.718,-0.004) (6.643,-0.004) (6.568,-0.004) (6.493,-0.004) (6.418,-0.004) (6.343,-0.003) (6.268,-0.003) (6.193,-0.003) (6.118,-0.003) (6.043,-0.003) (5.968,-0.003) (5.893,-0.003) (5.818,-0.003) (5.743,-0.003) (5.668,-0.002) (5.593,-0.002) (5.518,-0.002) (5.443,-0.002) (5.418,-0.002)}
\LC{(-0.001,2.928) (-0.001,3.003) (-0.001,3.078) (-0.001,3.153) (-0.001,3.228) (-0.002,3.303) (-0.002,3.378) (-0.002,3.453) (-0.002,3.528) (-0.002,3.603) (-0.002,3.678) (-0.002,3.753) (-0.002,3.828) (-0.002,3.903) (-0.002,3.978) (-0.002,4.053) (-0.002,4.128) (-0.002,4.203) (-0.003,4.278) (-0.003,4.353) (-0.003,4.428) (-0.003,4.503) (-0.003,4.578) (-0.003,4.653) (-0.003,4.728) (-0.003,4.803) (-0.003,4.878) (-0.003,4.953) (-0.004,5.028) (-0.004,5.103) (-0.004,5.178) (-0.004,5.253) (-0.004,5.328) (-0.004,5.403) (-0.004,5.478) (-0.004,5.553) (-0.004,5.628) (-0.005,5.703) (-0.005,5.778) (-0.005,5.853) (-0.005,5.928) (-0.005,6.003) (-0.005,6.078) (-0.005,6.153) (-0.005,6.228) (-0.006,6.303) (-0.006,6.378) (-0.006,6.453) (-0.006,6.528) (-0.006,6.603) (-0.006,6.678) (-0.006,6.753) (-0.007,6.828) (-0.007,6.903) (-0.007,6.978) (-0.007,7.053) (-0.007,7.128) (-0.007,7.203) (-0.007,7.278) (-0.008,7.353) (-0.008,7.428) (-0.008,7.503) (-0.008,7.578) (-0.008,7.653) (-0.008,7.728) (-0.009,7.803) (-0.009,7.878) (-0.009,7.953) (-0.009,8.028) (-0.009,8.103) (-0.009,8.178) (-0.010,8.253) (-0.010,8.328) (-0.010,8.403) (-0.010,8.478) (-0.010,8.553) (-0.010,8.628) (-0.011,8.703) (-0.011,8.778) (-0.011,8.853) (-0.011,8.928) (-0.011,9.003) (-0.012,9.078) (-0.012,9.153) (-0.012,9.228) (-0.012,9.303) (-0.012,9.378) (-0.013,9.453) (-0.013,9.528) (-0.013,9.603) (-0.013,9.678) (-0.013,9.753) (-0.014,9.828) (-0.014,9.903) (-0.014,9.978) (-0.014,10.000)}
\LC{(4.007,7.839) (4.071,7.879) (4.135,7.919) (4.199,7.957) (4.264,7.995) (4.329,8.032) (4.394,8.069) (4.460,8.105) (4.527,8.140) (4.593,8.174) (4.660,8.208) (4.728,8.241) (4.795,8.273) (4.864,8.304) (4.932,8.335) (5.001,8.365) (5.070,8.394) (5.140,8.422) (5.209,8.449) (5.280,8.476) (5.350,8.501) (5.421,8.526) (5.492,8.550) (5.563,8.573) (5.635,8.596) (5.707,8.617) (5.779,8.637) (5.851,8.657) (5.924,8.675) (5.997,8.693) (6.070,8.710) (6.143,8.726) (6.217,8.740) (6.291,8.754) (6.364,8.767) (6.439,8.779) (6.513,8.789) (6.587,8.799) (6.662,8.808) (6.736,8.816) (6.811,8.822) (6.886,8.828) (6.961,8.832) (7.036,8.835) (7.110,8.837) (7.185,8.838) (7.260,8.838) (7.335,8.837) (7.410,8.835) (7.485,8.831) (7.560,8.826) (7.635,8.820) (7.710,8.813) (7.784,8.805) (7.859,8.795) (7.933,8.784) (8.007,8.772) (8.081,8.759) (8.154,8.744) (8.227,8.728) (8.300,8.711) (8.373,8.692) (8.445,8.672) (8.517,8.651) (8.589,8.629) (8.660,8.605) (8.731,8.580) (8.801,8.553) (8.870,8.525) (8.939,8.496) (9.008,8.466) (9.076,8.434) (9.143,8.401) (9.210,8.366) (9.275,8.330) (9.341,8.293) (9.405,8.254) (9.468,8.214) (9.531,8.173) (9.593,8.130) (9.653,8.086) (9.713,8.041) (9.772,7.994) (9.830,7.946) (9.886,7.897) (9.942,7.847) (9.996,7.795) (10.000,7.778)}
\LC{(-0.891,-2.699) (-0.937,-2.758) (-0.984,-2.816) (-1.033,-2.873) (-1.082,-2.930) (-1.133,-2.985) (-1.184,-3.039) (-1.237,-3.093) (-1.291,-3.145) (-1.346,-3.196) (-1.402,-3.246) (-1.459,-3.295) (-1.517,-3.342) (-1.576,-3.388) (-1.636,-3.434) (-1.697,-3.477) (-1.759,-3.520) (-1.822,-3.560) (-1.885,-3.600) (-1.950,-3.638) (-2.016,-3.674) (-2.082,-3.709) (-2.149,-3.743) (-2.217,-3.774) (-2.286,-3.804) (-2.356,-3.832) (-2.426,-3.858) (-2.497,-3.883) (-2.568,-3.905) (-2.640,-3.926) (-2.713,-3.944) (-2.786,-3.961) (-2.860,-3.975) (-2.934,-3.987) (-3.008,-3.997) (-3.083,-4.004) (-3.158,-4.009) (-3.233,-4.012) (-3.308,-4.012) (-3.383,-4.010) (-3.457,-4.005) (-3.532,-3.997) (-3.606,-3.987) (-3.680,-3.974) (-3.753,-3.958) (-3.826,-3.939) (-3.898,-3.918) (-3.969,-3.893) (-4.038,-3.866) (-4.107,-3.835) (-4.174,-3.802) (-4.240,-3.765) (-4.303,-3.726) (-4.365,-3.684) (-4.425,-3.638) (-4.483,-3.590) (-4.538,-3.539) (-4.590,-3.486) (-4.640,-3.430) (-4.687,-3.371) (-4.730,-3.310) (-4.771,-3.247) (-4.808,-3.182) (-4.842,-3.115) (-4.872,-3.046) (-4.898,-2.976) (-4.921,-2.904) (-4.940,-2.832) (-4.955,-2.758) (-4.966,-2.684) (-4.974,-2.610) (-4.978,-2.535) (-4.979,-2.460) (-4.975,-2.385) (-4.969,-2.310) (-4.959,-2.236) (-4.945,-2.162) (-4.928,-2.089) (-4.909,-2.017) (-4.886,-1.945) (-4.860,-1.875) (-4.831,-1.805) (-4.800,-1.737) (-4.767,-1.670) (-4.730,-1.605) (-4.692,-1.540) (-4.651,-1.477) (-4.608,-1.416) (-4.563,-1.356) (-4.517,-1.297) (-4.468,-1.240) (-4.418,-1.184) (-4.366,-1.130) (-4.313,-1.077) (-4.258,-1.026) (-4.201,-0.977) (-4.144,-0.929) (-4.085,-0.882) (-4.025,-0.837) (-3.964,-0.793) (-3.902,-0.751) (-3.839,-0.711) (-3.775,-0.671) (-3.710,-0.634) (-3.667,-0.609)}
\LC{(1.347,2.119) (1.418,2.141) (1.491,2.160) (1.564,2.175) (1.638,2.187) (1.713,2.194) (1.788,2.196) (1.863,2.194) (1.938,2.187) (2.012,2.176) (2.085,2.159) (2.156,2.137) (2.226,2.109) (2.294,2.076) (2.358,2.038) (2.419,1.994) (2.476,1.945) (2.528,1.891) (2.575,1.832) (2.616,1.770) (2.651,1.703) (2.679,1.634) (2.701,1.562) (2.716,1.489) (2.723,1.414) (2.725,1.339) (2.719,1.264) (2.708,1.190) (2.691,1.117) (2.667,1.046) (2.639,0.977) (2.606,0.909) (2.582,0.866)}
\LC{(-1.963,-4.791) (-2.018,-4.843) (-2.073,-4.893) (-2.129,-4.943) (-2.186,-4.992) (-2.243,-5.041) (-2.302,-5.088) (-2.360,-5.135) (-2.420,-5.181) (-2.479,-5.226) (-2.540,-5.270) (-2.601,-5.313) (-2.663,-5.356) (-2.726,-5.397) (-2.789,-5.438) (-2.853,-5.477) (-2.917,-5.516) (-2.982,-5.553) (-3.047,-5.590) (-3.113,-5.626) (-3.180,-5.660) (-3.247,-5.694) (-3.315,-5.726) (-3.383,-5.757) (-3.451,-5.787) (-3.521,-5.816) (-3.590,-5.844) (-3.660,-5.871) (-3.731,-5.896) (-3.802,-5.921) (-3.873,-5.944) (-3.945,-5.965) (-4.017,-5.986) (-4.090,-6.005) (-4.163,-6.022) (-4.236,-6.039) (-4.309,-6.054) (-4.383,-6.067) (-4.457,-6.079) (-4.531,-6.090) (-4.606,-6.099) (-4.680,-6.106) (-4.755,-6.112) (-4.830,-6.117) (-4.905,-6.119) (-4.980,-6.121) (-5.055,-6.120) (-5.130,-6.118) (-5.205,-6.114) (-5.280,-6.109) (-5.354,-6.102) (-5.429,-6.092) (-5.503,-6.082) (-5.577,-6.069) (-5.650,-6.054) (-5.724,-6.038) (-5.796,-6.020) (-5.869,-6.000) (-5.940,-5.978) (-6.011,-5.954) (-6.082,-5.928) (-6.152,-5.900) (-6.220,-5.871) (-6.288,-5.839) (-6.355,-5.805) (-6.421,-5.770) (-6.486,-5.732) (-6.550,-5.693) (-6.613,-5.651) (-6.674,-5.608) (-6.734,-5.563) (-6.792,-5.516) (-6.849,-5.467) (-6.904,-5.416) (-6.958,-5.363) (-7.009,-5.309) (-7.059,-5.253) (-7.107,-5.196) (-7.153,-5.136) (-7.197,-5.076) (-7.239,-5.013) (-7.279,-4.950) (-7.317,-4.885) (-7.352,-4.819) (-7.385,-4.752) (-7.416,-4.683) (-7.444,-4.614) (-7.470,-4.543) (-7.493,-4.472) (-7.515,-4.400) (-7.533,-4.327) (-7.549,-4.254) (-7.563,-4.180) (-7.574,-4.106) (-7.583,-4.032) (-7.589,-3.957) (-7.593,-3.882) (-7.595,-3.807) (-7.594,-3.732) (-7.591,-3.657) (-7.585,-3.582) (-7.577,-3.508) (-7.567,-3.434) (-7.555,-3.360) (-7.541,-3.286) (-7.524,-3.213) (-7.505,-3.140) (-7.485,-3.068) (-7.462,-2.996) (-7.438,-2.926) (-7.412,-2.855) (-7.383,-2.786) (-7.353,-2.717) (-7.322,-2.649) (-7.288,-2.582) (-7.254,-2.516) (-7.217,-2.450) (-7.179,-2.385) (-7.139,-2.322) (-7.098,-2.259) (-7.056,-2.197) (-7.012,-2.136) (-6.967,-2.076) (-6.921,-2.017) (-6.874,-1.959) (-6.825,-1.902) (-6.775,-1.846) (-6.724,-1.791) (-6.672,-1.737) (-6.619,-1.683) (-6.566,-1.631) (-6.511,-1.580) (-6.455,-1.530) (-6.398,-1.481) (-6.341,-1.433) (-6.282,-1.386) (-6.223,-1.339) (-6.163,-1.294) (-6.103,-1.250) (-6.042,-1.207) (-5.980,-1.164) (-5.959,-1.151)}
\LC{(-3.628,-6.337) (-3.695,-6.369) (-3.764,-6.400) (-3.832,-6.430) (-3.901,-6.460) (-3.971,-6.488) (-4.041,-6.515) (-4.111,-6.541) (-4.182,-6.565) (-4.253,-6.589) (-4.325,-6.612) (-4.397,-6.633) (-4.469,-6.653) (-4.541,-6.672) (-4.614,-6.690) (-4.687,-6.707) (-4.761,-6.722) (-4.834,-6.736) (-4.908,-6.749) (-4.983,-6.760) (-5.057,-6.770) (-5.131,-6.779) (-5.206,-6.787) (-5.281,-6.793) (-5.356,-6.797) (-5.430,-6.800) (-5.505,-6.802) (-5.580,-6.802) (-5.655,-6.801) (-5.730,-6.798) (-5.805,-6.794) (-5.880,-6.788) (-5.955,-6.781) (-6.029,-6.771) (-6.103,-6.761) (-6.177,-6.749) (-6.251,-6.735) (-6.324,-6.719) (-6.397,-6.702) (-6.470,-6.683) (-6.542,-6.662) (-6.614,-6.640) (-6.685,-6.616) (-6.755,-6.590) (-6.825,-6.562) (-6.894,-6.533) (-6.962,-6.502) (-7.029,-6.469) (-7.096,-6.434) (-7.162,-6.398) (-7.226,-6.360) (-7.290,-6.320) (-7.352,-6.278) (-7.413,-6.235) (-7.473,-6.190) (-7.532,-6.143) (-7.589,-6.095) (-7.645,-6.045) (-7.700,-5.993) (-7.752,-5.940) (-7.804,-5.885) (-7.853,-5.829) (-7.901,-5.771) (-7.947,-5.712) (-7.991,-5.652) (-8.034,-5.590) (-8.074,-5.527) (-8.113,-5.462) (-8.149,-5.397) (-8.184,-5.330) (-8.216,-5.262) (-8.246,-5.194) (-8.274,-5.124) (-8.300,-5.054) (-8.324,-4.983) (-8.346,-4.911) (-8.365,-4.838) (-8.382,-4.765) (-8.397,-4.692) (-8.410,-4.618) (-8.420,-4.544) (-8.428,-4.469) (-8.435,-4.394) (-8.439,-4.320) (-8.440,-4.245) (-8.440,-4.170) (-8.437,-4.095) (-8.433,-4.020) (-8.426,-3.945) (-8.418,-3.871) (-8.407,-3.796) (-8.394,-3.722) (-8.380,-3.649) (-8.363,-3.576) (-8.345,-3.503) (-8.325,-3.431) (-8.303,-3.359) (-8.279,-3.288) (-8.254,-3.217) (-8.227,-3.147) (-8.198,-3.078) (-8.168,-3.009) (-8.136,-2.941) (-8.103,-2.874) (-8.068,-2.808) (-8.032,-2.742) (-7.994,-2.677) (-7.955,-2.613) (-7.915,-2.550) (-7.873,-2.487) (-7.831,-2.426) (-7.787,-2.365) (-7.742,-2.305) (-7.695,-2.246) (-7.648,-2.188) (-7.600,-2.130) (-7.550,-2.074) (-7.500,-2.019) (-7.448,-1.964) (-7.431,-1.946)}
\LC{(3.506,6.475) (3.571,6.511) (3.638,6.546) (3.704,6.580) (3.772,6.614) (3.839,6.646) (3.907,6.678) (3.976,6.708) (4.045,6.738) (4.114,6.767) (4.184,6.794) (4.254,6.821) (4.324,6.846) (4.395,6.871) (4.467,6.894) (4.538,6.916) (4.610,6.937) (4.682,6.957) (4.755,6.976) (4.828,6.994) (4.901,7.010) (4.974,7.026) (5.048,7.040) (5.122,7.053) (5.196,7.064) (5.270,7.075) (5.345,7.084) (5.419,7.092) (5.494,7.098) (5.569,7.103) (5.644,7.107) (5.719,7.109) (5.794,7.110) (5.869,7.110) (5.944,7.108) (6.019,7.104) (6.094,7.099) (6.168,7.093) (6.243,7.085) (6.317,7.076) (6.391,7.065) (6.465,7.052) (6.539,7.038) (6.612,7.023) (6.685,7.005) (6.758,6.986) (6.830,6.966) (6.902,6.944) (6.973,6.920) (7.043,6.894) (7.113,6.867) (7.183,6.839) (7.251,6.808) (7.319,6.776) (7.386,6.742) (7.452,6.706) (7.517,6.669) (7.581,6.630) (7.644,6.590) (7.706,6.548) (7.767,6.504) (7.827,6.458) (7.885,6.411) (7.942,6.363) (7.998,6.312) (8.052,6.261) (8.105,6.207) (8.156,6.152) (8.206,6.096) (8.254,6.039) (8.300,5.980) (8.344,5.919) (8.387,5.858) (8.428,5.795) (8.467,5.731) (8.504,5.666) (8.540,5.599) (8.573,5.532) (8.604,5.464) (8.633,5.395) (8.660,5.325) (8.685,5.254) (8.708,5.183) (8.729,5.111) (8.748,5.038) (8.765,4.965) (8.779,4.892) (8.792,4.818) (8.802,4.743) (8.810,4.669) (8.816,4.594) (8.820,4.519) (8.822,4.444) (8.822,4.369) (8.820,4.294) (8.816,4.219) (8.810,4.145) (8.801,4.070) (8.791,3.996) (8.780,3.922) (8.766,3.848) (8.750,3.775) (8.733,3.702) (8.714,3.629) (8.693,3.557) (8.670,3.486) (8.646,3.415) (8.620,3.344) (8.593,3.274) (8.564,3.205) (8.533,3.137) (8.501,3.069) (8.468,3.002) (8.433,2.935) (8.397,2.870) (8.359,2.805) (8.321,2.740) (8.281,2.677) (8.239,2.614) (8.197,2.552) (8.153,2.491) (8.109,2.431) (8.063,2.372) (8.016,2.313) (7.968,2.256) (7.919,2.199) (7.869,2.143) (7.818,2.088) (7.767,2.033)}
\LC{(-1.864,3.540) (-1.929,3.577) (-1.995,3.614) (-2.061,3.648) (-2.129,3.681) (-2.197,3.712) (-2.266,3.741) (-2.336,3.768) (-2.406,3.794) (-2.477,3.818) (-2.549,3.839) (-2.622,3.859) (-2.695,3.877) (-2.768,3.892) (-2.842,3.905) (-2.916,3.916) (-2.991,3.925) (-3.065,3.931) (-3.140,3.935) (-3.215,3.936) (-3.290,3.934) (-3.365,3.930) (-3.440,3.924) (-3.514,3.914) (-3.588,3.902) (-3.662,3.887) (-3.734,3.869) (-3.806,3.848) (-3.877,3.824) (-3.947,3.797) (-4.016,3.767) (-4.084,3.734) (-4.149,3.698) (-4.213,3.659) (-4.276,3.617) (-4.336,3.572) (-4.393,3.525) (-4.449,3.474) (-4.501,3.420) (-4.551,3.364) (-4.598,3.306) (-4.641,3.245) (-4.682,3.181) (-4.718,3.116) (-4.752,3.049) (-4.781,2.980) (-4.807,2.910) (-4.830,2.838) (-4.848,2.765) (-4.863,2.692) (-4.873,2.618) (-4.880,2.543) (-4.883,2.468) (-4.883,2.393) (-4.879,2.318) (-4.871,2.244) (-4.860,2.169) (-4.845,2.096) (-4.827,2.023) (-4.806,1.951) (-4.782,1.880) (-4.755,1.810) (-4.726,1.741) (-4.693,1.673) (-4.658,1.607) (-4.621,1.542) (-4.581,1.478) (-4.540,1.416) (-4.496,1.355) (-4.450,1.296) (-4.402,1.238) (-4.352,1.182) (-4.301,1.127)}
\LC{(2.522,3.823) (2.594,3.843) (2.667,3.861) (2.740,3.877) (2.814,3.891) (2.888,3.902) (2.962,3.911) (3.037,3.918) (3.112,3.923) (3.187,3.924) (3.262,3.924) (3.337,3.920) (3.412,3.914) (3.486,3.906) (3.560,3.894) (3.634,3.880) (3.707,3.862) (3.779,3.842) (3.850,3.819) (3.921,3.793) (3.990,3.763) (4.057,3.731) (4.124,3.696) (4.188,3.657) (4.250,3.616) (4.311,3.572) (4.369,3.524) (4.425,3.474) (4.478,3.421) (4.528,3.366) (4.576,3.308) (4.620,3.247) (4.661,3.184) (4.698,3.119) (4.732,3.052) (4.763,2.984) (4.789,2.914) (4.812,2.842) (4.831,2.770) (4.846,2.696) (4.858,2.622) (4.865,2.548) (4.869,2.473) (4.869,2.398) (4.865,2.323) (4.858,2.248) (4.847,2.174) (4.833,2.100) (4.816,2.027) (4.795,1.955) (4.772,1.884) (4.745,1.814) (4.716,1.745) (4.684,1.677) (4.649,1.610) (4.612,1.545) (4.599,1.524)}
\LC{(-2.405,-4.586) (-2.470,-4.624) (-2.535,-4.661) (-2.601,-4.697) (-2.668,-4.731) (-2.735,-4.764) (-2.803,-4.796) (-2.871,-4.827) (-2.940,-4.856) (-3.010,-4.883) (-3.080,-4.910) (-3.151,-4.934) (-3.223,-4.958) (-3.294,-4.979) (-3.367,-5.000) (-3.439,-5.018) (-3.512,-5.035) (-3.586,-5.050) (-3.659,-5.064) (-3.733,-5.076) (-3.808,-5.086) (-3.882,-5.094) (-3.957,-5.101) (-4.032,-5.105) (-4.107,-5.108) (-4.182,-5.109) (-4.257,-5.108) (-4.332,-5.104) (-4.407,-5.099) (-4.481,-5.091) (-4.556,-5.082) (-4.630,-5.070) (-4.703,-5.056) (-4.777,-5.040) (-4.849,-5.022) (-4.921,-5.001) (-4.993,-4.978) (-5.063,-4.953) (-5.133,-4.925) (-5.202,-4.895) (-5.270,-4.863) (-5.336,-4.829) (-5.402,-4.792) (-5.465,-4.753) (-5.528,-4.711) (-5.589,-4.667) (-5.648,-4.621) (-5.705,-4.573) (-5.761,-4.522) (-5.814,-4.470) (-5.865,-4.415) (-5.915,-4.358) (-5.961,-4.300) (-6.006,-4.239) (-6.047,-4.177) (-6.087,-4.113) (-6.123,-4.048) (-6.157,-3.981) (-6.188,-3.912) (-6.216,-3.843) (-6.242,-3.772) (-6.264,-3.701) (-6.284,-3.628) (-6.300,-3.555) (-6.314,-3.481) (-6.325,-3.407) (-6.332,-3.333) (-6.337,-3.258) (-6.339,-3.183) (-6.338,-3.108) (-6.334,-3.033) (-6.328,-2.958) (-6.319,-2.884) (-6.307,-2.810) (-6.292,-2.736) (-6.276,-2.663) (-6.256,-2.591) (-6.234,-2.519) (-6.210,-2.448) (-6.184,-2.378) (-6.155,-2.308) (-6.125,-2.240) (-6.092,-2.172) (-6.058,-2.106) (-6.021,-2.040) (-5.983,-1.976) (-5.943,-1.912) (-5.901,-1.850) (-5.857,-1.789) (-5.812,-1.729) (-5.766,-1.670) (-5.718,-1.612) (-5.669,-1.556) (-5.618,-1.500) (-5.566,-1.446) (-5.513,-1.393) (-5.459,-1.342) (-5.403,-1.291) (-5.347,-1.242) (-5.290,-1.193) (-5.231,-1.146) (-5.172,-1.101) (-5.111,-1.056) (-5.050,-1.013) (-4.988,-0.970) (-4.926,-0.929) (-4.862,-0.889) (-4.798,-0.850) (-4.733,-0.812) (-4.668,-0.776)}
\LC{(-3.747,6.524) (-3.815,6.556) (-3.883,6.587) (-3.952,6.617) (-4.021,6.646) (-4.091,6.674) (-4.161,6.701) (-4.231,6.727) (-4.302,6.751) (-4.373,6.775) (-4.445,6.798) (-4.517,6.819) (-4.589,6.840) (-4.661,6.859) (-4.734,6.877) (-4.807,6.893) (-4.880,6.909) (-4.954,6.923) (-5.028,6.936) (-5.102,6.948) (-5.176,6.959) (-5.251,6.968) (-5.325,6.976) (-5.400,6.982) (-5.475,6.987) (-5.550,6.991) (-5.625,6.994) (-5.700,6.995) (-5.775,6.994) (-5.850,6.992) (-5.925,6.989) (-5.999,6.984) (-6.074,6.977) (-6.149,6.969) (-6.223,6.960) (-6.297,6.948) (-6.371,6.936) (-6.445,6.921) (-6.518,6.905) (-6.591,6.888) (-6.663,6.869) (-6.736,6.848) (-6.807,6.825) (-6.878,6.801) (-6.948,6.775) (-7.018,6.747) (-7.087,6.718) (-7.155,6.687) (-7.223,6.654) (-7.290,6.620) (-7.355,6.584) (-7.420,6.546) (-7.484,6.506) (-7.546,6.465) (-7.608,6.422) (-7.668,6.378) (-7.727,6.331) (-7.785,6.283) (-7.841,6.234) (-7.896,6.183) (-7.950,6.130) (-8.002,6.076) (-8.052,6.021) (-8.101,5.964) (-8.148,5.905) (-8.193,5.845) (-8.236,5.784) (-8.278,5.722) (-8.317,5.658) (-8.355,5.593) (-8.391,5.527) (-8.425,5.460) (-8.457,5.392) (-8.486,5.324) (-8.514,5.254) (-8.539,5.183) (-8.563,5.112) (-8.584,5.040) (-8.603,4.968) (-8.620,4.895) (-8.635,4.821) (-8.648,4.747) (-8.658,4.673) (-8.666,4.598) (-8.673,4.524) (-8.677,4.449) (-8.679,4.374) (-8.679,4.299) (-8.676,4.224) (-8.672,4.149) (-8.666,4.074) (-8.658,4.000) (-8.648,3.925) (-8.636,3.851) (-8.622,3.778) (-8.606,3.704) (-8.588,3.631) (-8.569,3.559) (-8.548,3.487) (-8.525,3.416) (-8.500,3.345) (-8.474,3.274) (-8.446,3.205) (-8.417,3.136) (-8.386,3.067) (-8.354,3.000) (-8.320,2.933) (-8.285,2.867) (-8.248,2.801) (-8.210,2.736) (-8.171,2.672) (-8.130,2.609) (-8.089,2.547) (-8.046,2.486) (-8.002,2.425) (-7.956,2.365) (-7.910,2.306) (-7.863,2.248) (-7.814,2.191) (-7.765,2.134) (-7.714,2.079) (-7.663,2.024) (-7.646,2.006)}
\LC{(-1.421,-5.598) (-1.458,-5.663) (-1.495,-5.729) (-1.532,-5.794) (-1.570,-5.858) (-1.608,-5.923) (-1.647,-5.987) (-1.686,-6.051) (-1.726,-6.115) (-1.766,-6.178) (-1.807,-6.241) (-1.848,-6.304) (-1.889,-6.366) (-1.931,-6.429) (-1.973,-6.490) (-2.016,-6.552) (-2.059,-6.613) (-2.103,-6.674) (-2.147,-6.735) (-2.192,-6.795) (-2.237,-6.855) (-2.282,-6.915) (-2.328,-6.974) (-2.374,-7.033) (-2.421,-7.092) (-2.468,-7.150) (-2.516,-7.208) (-2.564,-7.266) (-2.613,-7.323) (-2.661,-7.380) (-2.711,-7.436) (-2.761,-7.492) (-2.811,-7.548) (-2.861,-7.603) (-2.912,-7.658) (-2.964,-7.713) (-3.016,-7.767) (-3.068,-7.821) (-3.121,-7.874) (-3.174,-7.927) (-3.228,-7.980) (-3.281,-8.032) (-3.336,-8.083) (-3.391,-8.135) (-3.446,-8.185) (-3.501,-8.236) (-3.557,-8.286) (-3.614,-8.335) (-3.671,-8.384) (-3.728,-8.432) (-3.786,-8.480) (-3.844,-8.528) (-3.902,-8.575) (-3.961,-8.622) (-4.020,-8.668) (-4.080,-8.713) (-4.139,-8.758) (-4.200,-8.803) (-4.261,-8.847) (-4.322,-8.890) (-4.383,-8.933) (-4.445,-8.976) (-4.507,-9.018) (-4.570,-9.059) (-4.633,-9.100) (-4.696,-9.140) (-4.760,-9.180) (-4.824,-9.219) (-4.888,-9.257) (-4.953,-9.295) (-5.018,-9.333) (-5.083,-9.369) (-5.149,-9.405) (-5.215,-9.441) (-5.281,-9.476) (-5.348,-9.510) (-5.415,-9.544) (-5.482,-9.577) (-5.550,-9.609) (-5.618,-9.641) (-5.686,-9.672) (-5.755,-9.702) (-5.824,-9.732) (-5.893,-9.761) (-5.962,-9.789) (-6.032,-9.817) (-6.102,-9.843) (-6.172,-9.870) (-6.243,-9.895) (-6.314,-9.920) (-6.385,-9.944) (-6.456,-9.967) (-6.528,-9.989) (-6.576,-10.000)}
\LC{(1.060,5.164) (1.090,5.232) (1.121,5.301) (1.151,5.369) (1.183,5.437) (1.214,5.505) (1.247,5.573) (1.279,5.641) (1.312,5.708) (1.345,5.775) (1.379,5.842) (1.413,5.909) (1.448,5.975) (1.483,6.042) (1.518,6.108) (1.554,6.174) (1.590,6.239) (1.627,6.305) (1.664,6.370) (1.701,6.435) (1.739,6.500) (1.778,6.564) (1.816,6.629) (1.855,6.693) (1.895,6.756) (1.935,6.820) (1.975,6.883) (2.016,6.946) (2.057,7.009) (2.098,7.071) (2.140,7.134) (2.182,7.196) (2.225,7.257) (2.268,7.319) (2.312,7.380) (2.355,7.441) (2.400,7.501) (2.444,7.561) (2.489,7.621) (2.535,7.681) (2.581,7.740) (2.627,7.799) (2.674,7.858) (2.721,7.916) (2.768,7.975) (2.816,8.032) (2.864,8.090) (2.913,8.147) (2.962,8.204) (3.011,8.260) (3.061,8.316) (3.111,8.372) (3.162,8.428) (3.212,8.483) (3.264,8.537) (3.315,8.592) (3.367,8.646) (3.420,8.699) (3.473,8.753) (3.526,8.806) (3.579,8.858) (3.633,8.910) (3.687,8.962) (3.742,9.013) (3.797,9.064) (3.853,9.115) (3.908,9.165) (3.964,9.215) (4.021,9.264) (4.078,9.313) (4.135,9.362) (4.192,9.410) (4.250,9.457) (4.309,9.505) (4.367,9.551) (4.426,9.598) (4.486,9.644) (4.545,9.689) (4.605,9.734) (4.666,9.779) (4.726,9.823) (4.787,9.866) (4.849,9.909) (4.910,9.952) (4.972,9.994) (4.993,10.000)}
\LC{(-4.211,7.912) (-4.276,7.950) (-4.342,7.986) (-4.408,8.022) (-4.474,8.057) (-4.540,8.092) (-4.607,8.126) (-4.675,8.159) (-4.743,8.191) (-4.811,8.222) (-4.879,8.253) (-4.948,8.283) (-5.017,8.312) (-5.086,8.340) (-5.156,8.367) (-5.226,8.394) (-5.297,8.420) (-5.368,8.445) (-5.439,8.469) (-5.510,8.492) (-5.582,8.514) (-5.654,8.535) (-5.726,8.556) (-5.798,8.575) (-5.871,8.594) (-5.944,8.611) (-6.017,8.628) (-6.090,8.644) (-6.164,8.659) (-6.237,8.672) (-6.311,8.685) (-6.385,8.697) (-6.460,8.707) (-6.534,8.717) (-6.609,8.725) (-6.683,8.733) (-6.758,8.739) (-6.833,8.745) (-6.908,8.749) (-6.983,8.752) (-7.058,8.754) (-7.133,8.755) (-7.208,8.754) (-7.283,8.753) (-7.357,8.750) (-7.432,8.746) (-7.507,8.741) (-7.582,8.735) (-7.657,8.727) (-7.731,8.719) (-7.805,8.709) (-7.880,8.697) (-7.953,8.685) (-8.027,8.671) (-8.101,8.656) (-8.174,8.640) (-8.247,8.622) (-8.319,8.603) (-8.391,8.583) (-8.463,8.561) (-8.535,8.538) (-8.606,8.514) (-8.676,8.488) (-8.746,8.461) (-8.815,8.432) (-8.884,8.403) (-8.952,8.371) (-9.020,8.339) (-9.087,8.305) (-9.153,8.270) (-9.219,8.233) (-9.283,8.195) (-9.347,8.156) (-9.410,8.115) (-9.472,8.073) (-9.534,8.030) (-9.594,7.985) (-9.653,7.939) (-9.711,7.892) (-9.768,7.844) (-9.824,7.794) (-9.879,7.743) (-9.933,7.690) (-9.985,7.637) (-10.000,7.618)}
\LC{(4.811,-7.150) (4.883,-7.169) (4.956,-7.187) (5.029,-7.203) (5.103,-7.219) (5.176,-7.233) (5.250,-7.246) (5.324,-7.258) (5.398,-7.269) (5.473,-7.278) (5.547,-7.286) (5.622,-7.293) (5.697,-7.299) (5.772,-7.303) (5.847,-7.306) (5.922,-7.307) (5.997,-7.307) (6.072,-7.306) (6.146,-7.303) (6.221,-7.299) (6.296,-7.294) (6.371,-7.287) (6.445,-7.278) (6.520,-7.268) (6.594,-7.257) (6.668,-7.244) (6.741,-7.229) (6.815,-7.213) (6.887,-7.195) (6.960,-7.176) (7.032,-7.155) (7.104,-7.133) (7.175,-7.109) (7.245,-7.083) (7.315,-7.056) (7.384,-7.027) (7.453,-6.997) (7.520,-6.965) (7.587,-6.931) (7.653,-6.895) (7.719,-6.858) (7.783,-6.820) (7.846,-6.779) (7.908,-6.737) (7.969,-6.694) (8.029,-6.649) (8.088,-6.602) (8.145,-6.554) (8.202,-6.504) (8.256,-6.453) (8.310,-6.400) (8.361,-6.346) (8.412,-6.290) (8.460,-6.233) (8.507,-6.175) (8.553,-6.115) (8.597,-6.054) (8.639,-5.992) (8.679,-5.929) (8.717,-5.864) (8.753,-5.799) (8.788,-5.732) (8.821,-5.664) (8.851,-5.596) (8.880,-5.527) (8.907,-5.457) (8.931,-5.386) (8.954,-5.314) (8.974,-5.242) (8.993,-5.169) (9.009,-5.096) (9.024,-5.023) (9.036,-4.949) (9.046,-4.874) (9.054,-4.800) (9.060,-4.725) (9.065,-4.650) (9.067,-4.575) (9.067,-4.500) (9.065,-4.425) (9.061,-4.350) (9.055,-4.276) (9.048,-4.201) (9.038,-4.127) (9.027,-4.052) (9.014,-3.979) (8.999,-3.905) (8.982,-3.832) (8.963,-3.759) (8.943,-3.687) (8.921,-3.615) (8.898,-3.544) (8.873,-3.473) (8.846,-3.403) (8.818,-3.334) (8.788,-3.265) (8.757,-3.197) (8.725,-3.129) (8.691,-3.062) (8.656,-2.996) (8.619,-2.930) (8.581,-2.866) (8.555,-2.823)}
\LC{(-0.995,2.331) (-1.051,2.381) (-1.109,2.428) (-1.169,2.474) (-1.230,2.518) (-1.292,2.560) (-1.355,2.600) (-1.420,2.638) (-1.486,2.674) (-1.553,2.707) (-1.621,2.739) (-1.690,2.767) (-1.760,2.794) (-1.831,2.818) (-1.903,2.839) (-1.976,2.857) (-2.049,2.873) (-2.123,2.885) (-2.198,2.895) (-2.272,2.901) (-2.347,2.904) (-2.422,2.903) (-2.497,2.899) (-2.572,2.891) (-2.646,2.879) (-2.719,2.864) (-2.792,2.844) (-2.863,2.821) (-2.933,2.793) (-3.001,2.762) (-3.067,2.726) (-3.130,2.686) (-3.191,2.642) (-3.249,2.594) (-3.303,2.542) (-3.354,2.487) (-3.400,2.428) (-3.442,2.366) (-3.480,2.302) (-3.513,2.234) (-3.541,2.165) (-3.564,2.093) (-3.581,2.020) (-3.594,1.946) (-3.601,1.872) (-3.603,1.797) (-3.600,1.722) (-3.593,1.647) (-3.580,1.573) (-3.564,1.500) (-3.543,1.428) (-3.517,1.358) (-3.488,1.288) (-3.456,1.221) (-3.420,1.155) (-3.381,1.091) (-3.338,1.029) (-3.293,0.969) (-3.246,0.911) (-3.196,0.855) (-3.143,0.802) (-3.089,0.750) (-3.033,0.700) (-2.975,0.653) (-2.915,0.607) (-2.854,0.564) (-2.791,0.523) (-2.728,0.483) (-2.662,0.446) (-2.596,0.410) (-2.574,0.399)}
\LC{(-10.000,-0.013) (-9.942,-0.013) (-9.867,-0.013) (-9.792,-0.012) (-9.717,-0.012) (-9.642,-0.012) (-9.567,-0.012) (-9.492,-0.011) (-9.417,-0.011) (-9.342,-0.011) (-9.267,-0.010) (-9.192,-0.010) (-9.117,-0.010) (-9.042,-0.010) (-8.967,-0.010) (-8.892,-0.009) (-8.817,-0.009) (-8.742,-0.009) (-8.667,-0.009) (-8.592,-0.008) (-8.517,-0.008) (-8.442,-0.008) (-8.367,-0.008) (-8.292,-0.008) (-8.217,-0.007) (-8.142,-0.007) (-8.067,-0.007) (-7.992,-0.007) (-7.917,-0.007) (-7.842,-0.006) (-7.767,-0.006) (-7.692,-0.006) (-7.617,-0.006) (-7.542,-0.006) (-7.467,-0.005) (-7.392,-0.005) (-7.317,-0.005) (-7.242,-0.005) (-7.167,-0.005) (-7.142,-0.005)}
\LC{(1.556,-6.813) (1.589,-6.880) (1.622,-6.947) (1.656,-7.014) (1.690,-7.081) (1.725,-7.147) (1.760,-7.214) (1.795,-7.280) (1.830,-7.346) (1.866,-7.412) (1.902,-7.478) (1.939,-7.543) (1.976,-7.608) (2.013,-7.673) (2.051,-7.738) (2.089,-7.803) (2.127,-7.868) (2.165,-7.932) (2.204,-7.996) (2.244,-8.060) (2.283,-8.124) (2.323,-8.187) (2.363,-8.250) (2.404,-8.314) (2.445,-8.376) (2.486,-8.439) (2.528,-8.501) (2.570,-8.564) (2.612,-8.626) (2.654,-8.687) (2.697,-8.749) (2.741,-8.810) (2.784,-8.871) (2.828,-8.932) (2.872,-8.992) (2.917,-9.053) (2.962,-9.113) (3.007,-9.173) (3.053,-9.232) (3.099,-9.292) (3.145,-9.351) (3.191,-9.409) (3.238,-9.468) (3.286,-9.526) (3.333,-9.584) (3.381,-9.642) (3.429,-9.699) (3.478,-9.757) (3.526,-9.814) (3.576,-9.870) (3.625,-9.927) (3.675,-9.983) (3.691,-10.000)}
\LC{(0.769,5.621) (0.789,5.693) (0.810,5.765) (0.831,5.837) (0.852,5.909) (0.873,5.981) (0.895,6.053) (0.917,6.125) (0.939,6.196) (0.961,6.268) (0.984,6.339) (1.007,6.411) (1.031,6.482) (1.054,6.553) (1.078,6.624) (1.102,6.695) (1.126,6.766) (1.151,6.837) (1.176,6.908) (1.201,6.979) (1.226,7.049) (1.252,7.120) (1.278,7.190) (1.304,7.260) (1.330,7.331) (1.357,7.401) (1.384,7.471) (1.411,7.541) (1.439,7.610) (1.466,7.680) (1.494,7.750) (1.523,7.819) (1.551,7.888) (1.580,7.958) (1.609,8.027) (1.638,8.096) (1.668,8.165) (1.698,8.234) (1.728,8.302) (1.758,8.371) (1.789,8.439) (1.820,8.508) (1.851,8.576) (1.882,8.644) (1.914,8.712) (1.946,8.780) (1.978,8.848) (2.010,8.915) (2.043,8.983) (2.076,9.050) (2.109,9.117) (2.143,9.184) (2.176,9.251) (2.210,9.318) (2.245,9.385) (2.279,9.452) (2.314,9.518) (2.349,9.584) (2.384,9.651) (2.420,9.717) (2.456,9.783) (2.492,9.848) (2.528,9.914) (2.565,9.979) (2.577,10.000)}
\LC{(3.537,-7.525) (3.597,-7.570) (3.658,-7.614) (3.719,-7.658) (3.780,-7.701) (3.842,-7.743) (3.904,-7.785) (3.967,-7.826) (4.030,-7.867) (4.094,-7.907) (4.158,-7.946) (4.222,-7.984) (4.287,-8.022) (4.352,-8.059) (4.418,-8.095) (4.484,-8.131) (4.550,-8.166) (4.617,-8.200) (4.684,-8.234) (4.751,-8.266) (4.819,-8.298) (4.887,-8.330) (4.956,-8.360) (5.025,-8.390) (5.094,-8.419) (5.164,-8.447) (5.234,-8.474) (5.304,-8.500) (5.374,-8.526) (5.445,-8.551) (5.516,-8.574) (5.588,-8.597) (5.659,-8.619) (5.731,-8.641) (5.803,-8.661) (5.876,-8.680) (5.948,-8.699) (6.021,-8.716) (6.095,-8.733) (6.168,-8.749) (6.241,-8.763) (6.315,-8.777) (6.389,-8.790) (6.463,-8.801) (6.537,-8.812) (6.612,-8.821) (6.686,-8.830) (6.761,-8.838) (6.836,-8.844) (6.911,-8.849) (6.985,-8.854) (7.060,-8.857) (7.135,-8.859) (7.210,-8.860) (7.285,-8.860) (7.360,-8.858) (7.435,-8.856) (7.510,-8.852) (7.585,-8.847) (7.660,-8.841) (7.734,-8.834) (7.809,-8.825) (7.883,-8.816) (7.957,-8.805) (8.031,-8.792) (8.105,-8.779) (8.179,-8.764) (8.252,-8.748) (8.325,-8.731) (8.398,-8.712) (8.470,-8.692) (8.542,-8.671) (8.613,-8.648) (8.684,-8.624) (8.755,-8.599) (8.825,-8.573) (8.895,-8.545) (8.964,-8.515) (9.032,-8.485) (9.100,-8.453) (9.167,-8.420) (9.234,-8.385) (9.300,-8.349) (9.365,-8.312) (9.429,-8.273) (9.493,-8.233) (9.555,-8.192) (9.617,-8.149) (9.678,-8.105) (9.737,-8.060) (9.796,-8.013) (9.854,-7.965) (9.910,-7.916) (9.966,-7.866) (10.000,-7.831)}
\LC{(0.740,-5.756) (0.760,-5.829) (0.779,-5.901) (0.799,-5.974) (0.819,-6.046) (0.839,-6.118) (0.859,-6.190) (0.880,-6.262) (0.901,-6.334) (0.922,-6.406) (0.943,-6.478) (0.965,-6.550) (0.987,-6.622) (1.009,-6.694) (1.031,-6.765) (1.054,-6.837) (1.076,-6.908) (1.099,-6.979) (1.123,-7.051) (1.146,-7.122) (1.170,-7.193) (1.194,-7.264) (1.218,-7.335) (1.243,-7.406) (1.268,-7.477) (1.293,-7.547) (1.318,-7.618) (1.344,-7.689) (1.369,-7.759) (1.395,-7.829) (1.422,-7.900) (1.448,-7.970) (1.475,-8.040) (1.502,-8.110) (1.529,-8.180) (1.556,-8.250) (1.584,-8.319) (1.612,-8.389) (1.640,-8.458) (1.668,-8.528) (1.697,-8.597) (1.726,-8.666) (1.755,-8.736) (1.784,-8.805) (1.814,-8.873) (1.844,-8.942) (1.874,-9.011) (1.904,-9.080) (1.935,-9.148) (1.966,-9.216) (1.997,-9.285) (2.028,-9.353) (2.060,-9.421) (2.091,-9.489) (2.123,-9.557) (2.156,-9.624) (2.188,-9.692) (2.221,-9.759) (2.254,-9.827) (2.287,-9.894) (2.320,-9.961) (2.343,-10.000)}
\LC{(0.443,6.926) (0.453,7.000) (0.463,7.075) (0.472,7.149) (0.482,7.223) (0.492,7.298) (0.503,7.372) (0.513,7.446) (0.523,7.521) (0.534,7.595) (0.544,7.669) (0.555,7.743) (0.566,7.818) (0.576,7.892) (0.587,7.966) (0.598,8.040) (0.610,8.114) (0.621,8.189) (0.632,8.263) (0.644,8.337) (0.655,8.411) (0.667,8.485) (0.679,8.559) (0.691,8.633) (0.703,8.707) (0.715,8.781) (0.727,8.855) (0.739,8.929) (0.752,9.003) (0.764,9.077) (0.777,9.151) (0.790,9.225) (0.802,9.299) (0.815,9.373) (0.828,9.447) (0.841,9.520) (0.855,9.594) (0.868,9.668) (0.881,9.742) (0.895,9.816) (0.909,9.889) (0.922,9.963) (0.932,10.000)}
\LC{(-0.456,6.301) (-0.467,6.376) (-0.478,6.450) (-0.489,6.524) (-0.500,6.598) (-0.511,6.672) (-0.523,6.746) (-0.535,6.820) (-0.546,6.894) (-0.558,6.969) (-0.570,7.043) (-0.582,7.117) (-0.595,7.191) (-0.607,7.265) (-0.620,7.338) (-0.632,7.412) (-0.645,7.486) (-0.658,7.560) (-0.671,7.634) (-0.684,7.708) (-0.697,7.782) (-0.711,7.855) (-0.724,7.929) (-0.738,8.003) (-0.752,8.077) (-0.766,8.150) (-0.780,8.224) (-0.794,8.298) (-0.808,8.371) (-0.823,8.445) (-0.837,8.518) (-0.852,8.592) (-0.867,8.666) (-0.882,8.739) (-0.897,8.812) (-0.912,8.886) (-0.928,8.959) (-0.943,9.033) (-0.959,9.106) (-0.974,9.179) (-0.990,9.253) (-1.006,9.326) (-1.022,9.399) (-1.039,9.472) (-1.055,9.546) (-1.071,9.619) (-1.088,9.692) (-1.105,9.765) (-1.122,9.838) (-1.139,9.911) (-1.156,9.984) (-1.162,10.000)}
\LC{(-5.902,-9.274) (-5.973,-9.298) (-6.045,-9.322) (-6.116,-9.344) (-6.188,-9.365) (-6.260,-9.386) (-6.333,-9.406) (-6.405,-9.425) (-6.478,-9.443) (-6.551,-9.460) (-6.624,-9.477) (-6.698,-9.492) (-6.771,-9.507) (-6.845,-9.521) (-6.919,-9.533) (-6.993,-9.545) (-7.067,-9.556) (-7.141,-9.566) (-7.216,-9.575) (-7.290,-9.583) (-7.365,-9.590) (-7.440,-9.596) (-7.515,-9.601) (-7.590,-9.605) (-7.664,-9.608) (-7.739,-9.611) (-7.814,-9.612) (-7.889,-9.611) (-7.964,-9.610) (-8.039,-9.608) (-8.114,-9.605) (-8.189,-9.601) (-8.264,-9.595) (-8.339,-9.589) (-8.413,-9.581) (-8.488,-9.572) (-8.562,-9.562) (-8.636,-9.551) (-8.710,-9.539) (-8.784,-9.525) (-8.858,-9.511) (-8.931,-9.495) (-9.004,-9.478) (-9.077,-9.460) (-9.149,-9.441) (-9.221,-9.420) (-9.293,-9.398) (-9.365,-9.375) (-9.436,-9.351) (-9.506,-9.326) (-9.576,-9.299) (-9.646,-9.271) (-9.715,-9.242) (-9.784,-9.212) (-9.852,-9.180) (-9.919,-9.147) (-9.986,-9.113) (-10.000,-9.102)}
\LC{(0.919,-1.901) (0.981,-1.943) (1.045,-1.983) (1.110,-2.020) (1.177,-2.054) (1.245,-2.085) (1.315,-2.113) (1.386,-2.138) (1.458,-2.159) (1.530,-2.176) (1.604,-2.189) (1.679,-2.199) (1.754,-2.204) (1.829,-2.204) (1.903,-2.200) (1.978,-2.191) (2.051,-2.177) (2.124,-2.158) (2.195,-2.133) (2.264,-2.103) (2.330,-2.068) (2.393,-2.028) (2.452,-1.982) (2.507,-1.931) (2.557,-1.875) (2.602,-1.814) (2.640,-1.750) (2.673,-1.683) (2.698,-1.612) (2.717,-1.540) (2.730,-1.466) (2.735,-1.391) (2.734,-1.316) (2.726,-1.241) (2.712,-1.168) (2.693,-1.095) (2.668,-1.025) (2.638,-0.956) (2.603,-0.889) (2.564,-0.825) (2.521,-0.764) (2.475,-0.705) (2.425,-0.649) (2.408,-0.630)}
\LC{(2.081,-3.206) (2.153,-3.228) (2.226,-3.247) (2.299,-3.263) (2.373,-3.277) (2.447,-3.288) (2.521,-3.297) (2.596,-3.302) (2.671,-3.305) (2.746,-3.305) (2.821,-3.301) (2.896,-3.294) (2.970,-3.284) (3.044,-3.271) (3.117,-3.254) (3.189,-3.234) (3.260,-3.210) (3.330,-3.183) (3.398,-3.152) (3.465,-3.117) (3.529,-3.079) (3.592,-3.037) (3.651,-2.992) (3.708,-2.943) (3.762,-2.891) (3.813,-2.836) (3.860,-2.777) (3.903,-2.716) (3.943,-2.652) (3.978,-2.586) (4.009,-2.518) (4.036,-2.448) (4.058,-2.376) (4.075,-2.303) (4.088,-2.229) (4.097,-2.155) (4.101,-2.080) (4.100,-2.005) (4.095,-1.930) (4.086,-1.856) (4.073,-1.782) (4.056,-1.709) (4.036,-1.637) (4.011,-1.566) (3.983,-1.496) (3.952,-1.428) (3.918,-1.361) (3.881,-1.296) (3.855,-1.253)}
\LC{(-0.657,-1.870) (-0.706,-1.926) (-0.757,-1.982) (-0.809,-2.036) (-0.863,-2.088) (-0.918,-2.138) (-0.975,-2.187) (-1.034,-2.234) (-1.094,-2.279) (-1.155,-2.322) (-1.218,-2.363) (-1.283,-2.401) (-1.348,-2.438) (-1.415,-2.471) (-1.483,-2.503) (-1.553,-2.532) (-1.623,-2.558) (-1.694,-2.581) (-1.766,-2.601) (-1.839,-2.618) (-1.913,-2.632) (-1.987,-2.642) (-2.062,-2.649) (-2.137,-2.652) (-2.212,-2.652) (-2.287,-2.648) (-2.361,-2.639) (-2.435,-2.627) (-2.508,-2.610) (-2.580,-2.589) (-2.651,-2.563) (-2.720,-2.533) (-2.786,-2.499) (-2.850,-2.460) (-2.911,-2.416) (-2.969,-2.369) (-3.023,-2.317) (-3.073,-2.261) (-3.119,-2.201) (-3.159,-2.138) (-3.195,-2.072) (-3.225,-2.003) (-3.249,-1.932) (-3.268,-1.860) (-3.282,-1.786) (-3.289,-1.712) (-3.291,-1.637) (-3.288,-1.562) (-3.279,-1.487) (-3.266,-1.413) (-3.247,-1.341) (-3.224,-1.269) (-3.197,-1.200) (-3.165,-1.132) (-3.130,-1.065) (-3.091,-1.001) (-3.049,-0.939) (-3.004,-0.879) (-2.956,-0.822) (-2.905,-0.766) (-2.852,-0.713) (-2.797,-0.663) (-2.740,-0.614) (-2.681,-0.568) (-2.620,-0.524) (-2.599,-0.510)}
\LC{(3.005,4.986) (3.075,5.014) (3.145,5.040) (3.216,5.066) (3.287,5.089) (3.358,5.112) (3.431,5.133) (3.503,5.152) (3.576,5.170) (3.649,5.186) (3.723,5.200) (3.797,5.213) (3.871,5.224) (3.945,5.234) (4.020,5.241) (4.095,5.247) (4.169,5.251) (4.244,5.253) (4.319,5.253) (4.394,5.252) (4.469,5.248) (4.544,5.242) (4.619,5.234) (4.693,5.225) (4.767,5.213) (4.841,5.199) (4.914,5.182) (4.987,5.164) (5.059,5.143) (5.130,5.121) (5.201,5.095) (5.271,5.068) (5.340,5.039) (5.408,5.007) (5.474,4.973) (5.540,4.936) (5.604,4.897) (5.667,4.857) (5.728,4.813) (5.788,4.768) (5.846,4.720) (5.902,4.671) (5.956,4.619) (6.009,4.565) (6.059,4.509) (6.107,4.452) (6.152,4.392) (6.195,4.331) (6.236,4.268) (6.274,4.203) (6.310,4.137) (6.343,4.070) (6.373,4.001) (6.400,3.931) (6.425,3.860) (6.447,3.789) (6.465,3.716) (6.481,3.643) (6.495,3.569) (6.505,3.495) (6.512,3.420) (6.517,3.345) (6.519,3.270) (6.518,3.195) (6.514,3.120) (6.508,3.046) (6.499,2.971) (6.487,2.897) (6.473,2.823) (6.456,2.750) (6.437,2.678) (6.416,2.606) (6.392,2.535) (6.367,2.464) (6.339,2.395) (6.309,2.326) (6.277,2.258) (6.243,2.191) (6.207,2.125) (6.169,2.060) (6.130,1.996) (6.089,1.934) (6.047,1.872) (6.002,1.811) (5.972,1.771)}
\LC{(10.000,1.463) (9.942,1.431) (9.874,1.399) (9.805,1.368) (9.737,1.338) (9.668,1.308) (9.599,1.278) (9.530,1.249) (9.461,1.221) (9.391,1.192) (9.322,1.165) (9.252,1.137) (9.182,1.110) (9.112,1.084) (9.041,1.058) (8.971,1.032) (8.900,1.007) (8.829,0.982) (8.758,0.957) (8.687,0.933) (8.616,0.910) (8.545,0.886) (8.474,0.864) (8.402,0.841) (8.330,0.819) (8.258,0.797) (8.187,0.776) (8.115,0.755) (8.042,0.735) (7.970,0.714) (7.898,0.695) (7.825,0.675) (7.753,0.656) (7.680,0.638) (7.608,0.619) (7.535,0.601) (7.462,0.584) (7.389,0.566) (7.316,0.549) (7.243,0.533) (7.169,0.516) (7.096,0.500) (7.023,0.485) (6.949,0.470) (6.876,0.455) (6.851,0.450)}
\LC{(3.473,-5.665) (3.543,-5.692) (3.614,-5.717) (3.685,-5.741) (3.756,-5.764) (3.828,-5.786) (3.900,-5.806) (3.973,-5.825) (4.045,-5.843) (4.119,-5.859) (4.192,-5.874) (4.266,-5.887) (4.340,-5.898) (4.414,-5.909) (4.489,-5.917) (4.564,-5.924) (4.638,-5.930) (4.713,-5.934) (4.788,-5.936) (4.863,-5.936) (4.938,-5.935) (5.013,-5.932) (5.088,-5.927) (5.163,-5.921) (5.237,-5.913) (5.312,-5.902) (5.386,-5.890) (5.459,-5.876) (5.533,-5.861) (5.606,-5.843) (5.678,-5.823) (5.750,-5.801) (5.821,-5.778) (5.891,-5.752) (5.961,-5.725) (6.030,-5.695) (6.098,-5.663) (6.165,-5.630) (6.231,-5.594) (6.296,-5.556) (6.360,-5.517) (6.422,-5.475) (6.483,-5.432) (6.542,-5.386) (6.601,-5.338) (6.657,-5.289) (6.712,-5.238) (6.765,-5.185) (6.816,-5.130) (6.865,-5.073) (6.912,-5.015) (6.957,-4.955) (7.001,-4.894) (7.041,-4.831) (7.080,-4.767) (7.116,-4.701) (7.150,-4.634) (7.182,-4.566) (7.211,-4.497) (7.238,-4.427) (7.262,-4.356) (7.284,-4.284) (7.303,-4.212) (7.319,-4.138) (7.333,-4.065) (7.345,-3.991) (7.354,-3.916) (7.360,-3.842) (7.364,-3.767) (7.366,-3.692) (7.365,-3.617) (7.362,-3.542) (7.356,-3.467) (7.348,-3.392) (7.337,-3.318) (7.325,-3.244) (7.310,-3.171) (7.293,-3.098) (7.274,-3.025) (7.253,-2.953) (7.229,-2.882) (7.204,-2.811) (7.177,-2.741) (7.148,-2.672) (7.117,-2.604) (7.085,-2.536) (7.051,-2.469) (7.015,-2.403) (6.978,-2.338) (6.939,-2.274) (6.898,-2.211) (6.856,-2.149) (6.813,-2.088) (6.768,-2.028)}
\LC{(-3.826,5.589) (-3.899,5.607) (-3.972,5.622) (-4.046,5.636) (-4.119,5.649) (-4.194,5.660) (-4.268,5.670) (-4.343,5.677) (-4.417,5.684) (-4.492,5.688) (-4.567,5.691) (-4.642,5.692) (-4.717,5.691) (-4.792,5.688) (-4.867,5.684) (-4.942,5.678) (-5.016,5.670) (-5.091,5.660) (-5.165,5.647) (-5.238,5.634) (-5.312,5.618) (-5.384,5.600) (-5.457,5.580) (-5.528,5.558) (-5.599,5.533) (-5.670,5.507) (-5.739,5.479) (-5.808,5.449) (-5.875,5.416) (-5.942,5.382) (-6.007,5.345) (-6.071,5.306) (-6.134,5.265) (-6.196,5.222) (-6.256,5.177) (-6.314,5.130) (-6.371,5.081) (-6.426,5.031) (-6.480,4.978) (-6.531,4.923) (-6.580,4.867) (-6.628,4.808) (-6.673,4.749) (-6.716,4.687) (-6.756,4.624) (-6.794,4.559) (-6.830,4.494) (-6.864,4.426) (-6.895,4.358) (-6.923,4.289) (-6.949,4.218) (-6.972,4.147) (-6.992,4.075) (-7.010,4.002) (-7.026,3.929) (-7.038,3.855) (-7.048,3.780) (-7.055,3.706) (-7.060,3.631) (-7.062,3.556) (-7.062,3.481) (-7.059,3.406) (-7.053,3.331) (-7.046,3.257) (-7.035,3.182) (-7.023,3.108) (-7.008,3.035) (-6.990,2.962) (-6.971,2.889) (-6.949,2.818) (-6.926,2.746) (-6.900,2.676) (-6.872,2.606) (-6.843,2.537) (-6.811,2.469) (-6.778,2.402) (-6.743,2.336) (-6.707,2.270)}
\LC{(0.590,-1.027) (0.659,-1.056) (0.731,-1.079) (0.804,-1.094) (0.879,-1.100) (0.954,-1.098) (1.028,-1.086) (1.099,-1.064) (1.167,-1.031) (1.227,-0.987) (1.279,-0.933) (1.320,-0.870) (1.349,-0.801) (1.363,-0.728) (1.365,-0.653) (1.353,-0.579) (1.330,-0.507) (1.298,-0.440) (1.256,-0.377) (1.208,-0.320) (1.190,-0.302)}
\LC{(-2.212,7.292) (-2.255,7.354) (-2.298,7.415) (-2.341,7.476) (-2.385,7.537) (-2.429,7.598) (-2.474,7.658) (-2.519,7.718) (-2.564,7.778) (-2.610,7.837) (-2.656,7.896) (-2.703,7.955) (-2.750,8.013) (-2.797,8.071) (-2.845,8.129) (-2.893,8.187) (-2.942,8.244) (-2.991,8.301) (-3.040,8.357) (-3.090,8.414) (-3.140,8.469) (-3.190,8.525) (-3.241,8.580) (-3.292,8.635) (-3.344,8.689) (-3.395,8.744) (-3.448,8.797) (-3.500,8.851) (-3.553,8.904) (-3.607,8.956) (-3.661,9.009) (-3.715,9.061) (-3.769,9.112) (-3.824,9.163) (-3.879,9.214) (-3.935,9.264) (-3.991,9.314) (-4.047,9.364) (-4.104,9.413) (-4.161,9.462) (-4.218,9.510) (-4.276,9.558) (-4.334,9.605) (-4.392,9.652) (-4.451,9.699) (-4.510,9.745) (-4.570,9.791) (-4.630,9.836) (-4.690,9.881) (-4.750,9.925) (-4.811,9.969) (-4.872,10.000)}
\LC{(3.158,7.740) (3.213,7.791) (3.267,7.843) (3.322,7.894) (3.378,7.944) (3.434,7.994) (3.491,8.043) (3.547,8.092) (3.605,8.140) (3.662,8.188) (3.721,8.236) (3.779,8.283) (3.838,8.329) (3.897,8.375) (3.957,8.420) (4.017,8.465) (4.078,8.510) (4.138,8.553) (4.200,8.597) (4.261,8.639) (4.323,8.681) (4.386,8.723) (4.449,8.764) (4.512,8.804) (4.575,8.844) (4.639,8.883) (4.704,8.922) (4.768,8.960) (4.833,8.997) (4.899,9.034) (4.964,9.070) (5.030,9.106) (5.097,9.141) (5.164,9.175) (5.231,9.208) (5.298,9.241) (5.366,9.274) (5.434,9.305) (5.502,9.336) (5.571,9.366) (5.640,9.395) (5.709,9.424) (5.779,9.452) (5.849,9.479) (5.919,9.506) (5.989,9.532) (6.060,9.557) (6.131,9.581) (6.202,9.605) (6.274,9.627) (6.345,9.649) (6.417,9.670) (6.489,9.691) (6.562,9.710) (6.635,9.729) (6.707,9.747) (6.780,9.764) (6.854,9.780) (6.927,9.795) (7.001,9.809) (7.074,9.823) (7.148,9.836) (7.222,9.847) (7.297,9.858) (7.371,9.868) (7.445,9.877) (7.520,9.885) (7.595,9.892) (7.669,9.898) (7.744,9.903) (7.819,9.908) (7.894,9.911) (7.969,9.913) (8.044,9.914) (8.119,9.914) (8.194,9.914) (8.269,9.912) (8.344,9.909) (8.419,9.905) (8.494,9.900) (8.568,9.894) (8.643,9.886) (8.718,9.878) (8.792,9.869) (8.866,9.858) (8.940,9.847) (9.014,9.834) (9.088,9.820) (9.162,9.805) (9.235,9.789) (9.308,9.772) (9.380,9.753) (9.453,9.734) (9.525,9.713) (9.597,9.691) (9.668,9.668) (9.739,9.644) (9.809,9.618) (9.880,9.591) (9.949,9.564) (10.000,9.534)}
\LC{(-1.193,-1.771) (-1.266,-1.788) (-1.340,-1.800) (-1.415,-1.807) (-1.490,-1.809) (-1.565,-1.805) (-1.639,-1.796) (-1.713,-1.780) (-1.784,-1.758) (-1.853,-1.729) (-1.919,-1.693) (-1.982,-1.651) (-2.039,-1.603) (-2.090,-1.549) (-2.136,-1.489) (-2.174,-1.424) (-2.204,-1.356) (-2.226,-1.284) (-2.240,-1.210) (-2.245,-1.135) (-2.242,-1.061) (-2.232,-0.986) (-2.215,-0.913) (-2.190,-0.842) (-2.160,-0.774) (-2.124,-0.708) (-2.083,-0.645) (-2.037,-0.586) (-1.988,-0.529) (-1.935,-0.477)}
\LC{(-2.191,-3.252) (-2.264,-3.270) (-2.337,-3.286) (-2.411,-3.299) (-2.485,-3.309) (-2.560,-3.317) (-2.635,-3.321) (-2.710,-3.323) (-2.785,-3.321) (-2.860,-3.317) (-2.934,-3.309) (-3.009,-3.298) (-3.082,-3.284) (-3.155,-3.266) (-3.227,-3.245) (-3.298,-3.220) (-3.367,-3.191) (-3.435,-3.159) (-3.501,-3.124) (-3.565,-3.085) (-3.626,-3.042) (-3.686,-2.996) (-3.742,-2.946) (-3.795,-2.893) (-3.845,-2.837) (-3.891,-2.778) (-3.933,-2.716) (-3.972,-2.652) (-4.006,-2.585) (-4.036,-2.517) (-4.062,-2.446) (-4.083,-2.374) (-4.100,-2.301) (-4.112,-2.227) (-4.120,-2.153) (-4.123,-2.078) (-4.122,-2.003) (-4.116,-1.928) (-4.106,-1.853) (-4.093,-1.780) (-4.075,-1.707) (-4.054,-1.635) (-4.029,-1.564) (-4.001,-1.495) (-3.969,-1.427) (-3.935,-1.360) (-3.923,-1.338)}
\LC{(-10.000,-2.495) (-9.955,-2.443) (-9.899,-2.393) (-9.844,-2.342) (-9.787,-2.293) (-9.730,-2.244) (-9.673,-2.196) (-9.615,-2.148) (-9.556,-2.102) (-9.497,-2.055) (-9.438,-2.010) (-9.378,-1.965) (-9.317,-1.921) (-9.256,-1.877) (-9.195,-1.834) (-9.133,-1.791) (-9.070,-1.750) (-9.008,-1.709) (-8.945,-1.668) (-8.881,-1.628) (-8.817,-1.589) (-8.753,-1.550) (-8.688,-1.512) (-8.623,-1.475) (-8.558,-1.438) (-8.492,-1.402) (-8.426,-1.366) (-8.360,-1.331) (-8.293,-1.297) (-8.226,-1.263) (-8.159,-1.230) (-8.092,-1.197) (-8.024,-1.165) (-7.956,-1.133) (-7.888,-1.102) (-7.865,-1.092)}
\LC{(-1.547,2.185) (-1.621,2.198) (-1.695,2.206) (-1.770,2.211) (-1.845,2.210) (-1.920,2.206) (-1.994,2.196) (-2.068,2.181) (-2.140,2.161) (-2.211,2.136) (-2.279,2.105) (-2.345,2.069) (-2.408,2.028) (-2.466,1.981) (-2.521,1.930) (-2.570,1.873) (-2.614,1.813) (-2.652,1.748) (-2.684,1.680) (-2.709,1.609) (-2.727,1.537) (-2.739,1.463) (-2.743,1.388) (-2.741,1.313) (-2.733,1.238) (-2.719,1.165) (-2.699,1.092) (-2.673,1.022) (-2.654,0.976)}

    % Cuadrupolo puntual - puntas de flecha
% Generado por scripts/generar_datos.py (no editar a mano).
\FL{1.698}{1.255}{-50.0}
\FL{-1.668}{1.205}{-127.5}
\FL{-1.140}{-0.805}{125.1}
\FL{0.742}{-2.503}{-56.3}
\FL{3.093}{-4.038}{-5.2}
\FL{5.016}{-2.372}{95.0}
\FL{3.140}{-0.363}{160.0}
\FL{0.727}{-4.161}{-70.0}
\FL{2.222}{-6.978}{-53.8}
\FL{4.485}{-9.221}{-35.2}
\FL{-0.525}{-3.023}{-109.9}
\FL{-1.850}{-5.390}{-129.0}
\FL{-3.921}{-7.133}{-151.7}
\FL{-6.535}{-7.778}{177.6}
\FL{-8.941}{-6.719}{131.8}
\FL{-9.587}{-4.176}{78.1}
\FL{-8.164}{-1.908}{40.8}
\FL{-5.787}{-0.612}{18.2}
\FL{-3.118}{-0.093}{5.1}
\FL{-0.841}{2.977}{122.2}
\FL{-3.357}{4.910}{166.1}
\FL{-6.080}{3.778}{-111.5}
\FL{-4.942}{0.990}{-34.9}
\FL{1.636}{2.861}{25.9}
\FL{3.688}{1.358}{-114.5}
\FL{0.772}{3.688}{66.1}
\FL{2.378}{6.084}{45.7}
\FL{4.794}{7.690}{20.3}
\FL{7.642}{7.872}{-15.6}
\FL{9.761}{6.016}{-69.4}
\FL{9.462}{3.225}{-119.6}
\FL{7.364}{1.238}{-150.8}
\FL{4.636}{0.291}{-169.2}
\FL{2.129}{4.115}{31.8}
\FL{4.967}{4.269}{-33.8}
\FL{5.293}{1.622}{-125.9}
\FL{-8.711}{0.934}{-18.5}
\FL{-6.038}{0.304}{-8.7}
\FL{2.566}{-6.228}{-43.4}
\FL{4.795}{-7.625}{-19.7}
\FL{7.426}{-7.819}{13.5}
\FL{9.508}{-6.272}{62.2}
\FL{9.603}{-3.705}{111.2}
\FL{7.947}{-1.649}{143.8}
\FL{5.564}{-0.521}{163.9}
\FL{-1.077}{4.763}{115.8}
\FL{-2.868}{7.337}{134.3}
\FL{-5.432}{9.133}{156.6}
\FL{-8.506}{9.632}{-173.2}
\FL{-8.188}{-0.711}{15.0}
\FL{2.765}{5.367}{32.0}
\FL{5.361}{5.998}{-7.9}
\FL{7.382}{4.398}{-72.9}
\FL{6.722}{1.874}{-131.0}
\FL{0.247}{-5.362}{-84.7}
\FL{0.618}{-8.464}{-81.6}
\FL{-0.453}{-5.987}{-98.7}
\FL{-0.962}{-8.689}{-102.7}
\FL{2.681}{-4.875}{-28.2}
\FL{5.831}{-4.848}{37.3}
\FL{6.106}{-1.874}{125.8}
\FL{-3.542}{7.008}{146.9}
\FL{-6.429}{7.945}{179.2}
\FL{-9.161}{6.822}{-130.9}
\FL{-9.701}{3.932}{-72.4}
\FL{-7.832}{1.559}{-34.6}
\FL{8.491}{-1.022}{159.2}
\FL{5.981}{-9.411}{-19.0}
\FL{8.689}{-9.707}{8.0}
\FL{3.731}{-6.301}{-23.7}
\FL{6.270}{-6.585}{13.7}
\FL{8.162}{-4.931}{71.5}
\FL{7.742}{-2.457}{123.9}
\FL{-1.237}{6.475}{111.8}
\FL{-2.440}{8.889}{121.3}
\FL{-4.616}{-8.202}{-153.1}
\FL{-8.349}{-8.729}{165.8}
\FL{-1.263}{-6.445}{-112.4}
\FL{-2.523}{-8.885}{-122.3}
\FL{-2.814}{5.567}{146.9}
\FL{-5.712}{6.264}{-170.3}
\FL{-7.797}{4.315}{-99.7}
\FL{-6.691}{1.609}{-42.1}
\FL{7.693}{-0.006}{179.9}
\FL{-0.003}{4.703}{90.1}
\FL{-0.010}{8.253}{90.1}
\FL{5.468}{8.542}{18.5}
\FL{8.660}{8.605}{-19.1}
\FL{-2.426}{-3.858}{-160.2}
\FL{-4.909}{-2.017}{73.5}
\FL{2.476}{1.945}{-43.4}
\FL{-3.025}{-5.578}{-150.9}
\FL{-5.552}{-6.073}{170.1}
\FL{-7.493}{-4.472}{107.3}
\FL{-6.921}{-2.017}{51.3}
\FL{-5.057}{-6.770}{-172.8}
\FL{-7.770}{-5.922}{133.4}
\FL{-8.236}{-3.171}{68.6}
\FL{4.999}{7.031}{11.0}
\FL{7.998}{6.312}{-42.9}
\FL{8.620}{3.344}{-110.8}
\FL{-3.140}{3.935}{178.1}
\FL{-4.876}{2.293}{-84.5}
\FL{4.461}{3.439}{-45.4}
\FL{-3.684}{-5.068}{-170.7}
\FL{-6.074}{-4.134}{121.2}
\FL{-5.718}{-1.612}{49.6}
\FL{-5.201}{6.962}{172.8}
\FL{-7.984}{6.094}{-133.5}
\FL{-8.465}{3.251}{-68.5}
\FL{-2.390}{-7.053}{-128.5}
\FL{-4.996}{-9.320}{-150.2}
\FL{1.790}{6.586}{59.0}
\FL{3.760}{9.030}{42.9}
\FL{-5.654}{8.535}{163.8}
\FL{-8.723}{8.470}{-158.7}
\FL{6.717}{-7.234}{11.4}
\FL{9.066}{-4.625}{88.1}
\FL{-3.337}{2.506}{-131.8}
\FL{-8.567}{-0.008}{0.2}
\FL{2.514}{-8.481}{-56.3}
\FL{1.542}{7.865}{67.6}
\FL{5.071}{-8.409}{-22.5}
\FL{8.494}{-8.685}{16.3}
\FL{1.430}{-7.923}{-69.4}
\FL{0.667}{8.485}{81.0}
\FL{-0.771}{8.175}{100.8}
\FL{-7.989}{-9.610}{178.5}
\FL{2.413}{-2.013}{37.0}
\FL{3.762}{-2.891}{45.7}
\FL{-2.808}{-2.486}{149.3}
\FL{4.544}{5.242}{-5.2}
\FL{6.518}{3.195}{-91.8}
\FL{8.450}{0.856}{-162.5}
\FL{5.188}{-5.918}{6.1}
\FL{7.365}{-3.617}{91.6}
\FL{-5.263}{5.628}{-168.0}
\FL{-7.060}{3.631}{-92.6}
\FL{1.263}{-0.952}{47.8}
\FL{-3.430}{8.779}{134.3}
\FL{4.639}{8.883}{31.3}
\FL{8.194}{9.914}{-1.0}
\FL{-2.162}{-1.446}{119.3}
\FL{-3.778}{-2.911}{134.6}
\FL{-8.987}{-1.695}{32.8}
\FL{-2.521}{1.930}{-133.9}
}
  \icono
  \marcoMMA{10}{5}{$x$ (u.\,arb.)}{$z$ (u.\,arb.)}
\end{tikzpicture}\\[1pt]{\small\scshape (b)}
\end{minipage}\hfill
\begin{minipage}[t]{0.325\textwidth}\centering
{\footnotesize Campo eléctrico}\\[3pt]
\begin{tikzpicture}[x=0.198cm, y=0.198cm]
  \panelMMA{10}{mmafondo}{1.2}{% Cuadrupolo puntual - campo vectorial
% Generado por scripts/generar_datos.py (no editar a mano).
\VF{-10.000}{-10.000}{-0.370}{0.123}{0.000}
\VF{-9.000}{-10.000}{-0.386}{0.058}{0.000}
\VF{-8.000}{-10.000}{-0.424}{-0.013}{0.037}
\VF{-7.000}{-10.000}{-0.449}{-0.097}{0.076}
\VF{-6.000}{-10.000}{-0.455}{-0.192}{0.114}
\VF{-5.000}{-10.000}{-0.438}{-0.292}{0.150}
\VF{-4.000}{-10.000}{-0.395}{-0.391}{0.182}
\VF{-3.000}{-10.000}{-0.325}{-0.480}{0.209}
\VF{-2.000}{-10.000}{-0.232}{-0.551}{0.229}
\VF{-1.000}{-10.000}{-0.121}{-0.598}{0.242}
\VF{0.000}{-10.000}{0.000}{-0.614}{0.246}
\VF{1.000}{-10.000}{0.121}{-0.598}{0.242}
\VF{2.000}{-10.000}{0.232}{-0.551}{0.229}
\VF{3.000}{-10.000}{0.325}{-0.480}{0.209}
\VF{4.000}{-10.000}{0.395}{-0.391}{0.182}
\VF{5.000}{-10.000}{0.438}{-0.292}{0.150}
\VF{6.000}{-10.000}{0.455}{-0.192}{0.114}
\VF{7.000}{-10.000}{0.449}{-0.097}{0.076}
\VF{8.000}{-10.000}{0.424}{-0.013}{0.037}
\VF{9.000}{-10.000}{0.386}{0.058}{0.000}
\VF{10.000}{-10.000}{0.370}{0.123}{0.000}
\VF{-10.000}{-9.000}{-0.341}{0.189}{0.000}
\VF{-9.000}{-9.000}{-0.388}{0.129}{0.021}
\VF{-8.000}{-9.000}{-0.444}{0.058}{0.064}
\VF{-7.000}{-9.000}{-0.486}{-0.034}{0.107}
\VF{-6.000}{-9.000}{-0.507}{-0.143}{0.150}
\VF{-5.000}{-9.000}{-0.500}{-0.262}{0.192}
\VF{-4.000}{-9.000}{-0.460}{-0.383}{0.229}
\VF{-3.000}{-9.000}{-0.386}{-0.496}{0.262}
\VF{-2.000}{-9.000}{-0.279}{-0.588}{0.286}
\VF{-1.000}{-9.000}{-0.146}{-0.648}{0.302}
\VF{0.000}{-9.000}{0.000}{-0.670}{0.307}
\VF{1.000}{-9.000}{0.146}{-0.648}{0.302}
\VF{2.000}{-9.000}{0.279}{-0.588}{0.286}
\VF{3.000}{-9.000}{0.386}{-0.496}{0.262}
\VF{4.000}{-9.000}{0.460}{-0.383}{0.229}
\VF{5.000}{-9.000}{0.500}{-0.262}{0.192}
\VF{6.000}{-9.000}{0.507}{-0.143}{0.150}
\VF{7.000}{-9.000}{0.486}{-0.034}{0.107}
\VF{8.000}{-9.000}{0.444}{0.058}{0.064}
\VF{9.000}{-9.000}{0.388}{0.129}{0.021}
\VF{10.000}{-9.000}{0.341}{0.189}{0.000}
\VF{-10.000}{-8.000}{-0.292}{0.258}{0.000}
\VF{-9.000}{-8.000}{-0.371}{0.217}{0.044}
\VF{-8.000}{-8.000}{-0.447}{0.149}{0.090}
\VF{-7.000}{-8.000}{-0.513}{0.054}{0.138}
\VF{-6.000}{-8.000}{-0.556}{-0.067}{0.187}
\VF{-5.000}{-8.000}{-0.567}{-0.208}{0.235}
\VF{-4.000}{-8.000}{-0.536}{-0.357}{0.280}
\VF{-3.000}{-8.000}{-0.460}{-0.501}{0.319}
\VF{-2.000}{-8.000}{-0.338}{-0.622}{0.349}
\VF{-1.000}{-8.000}{-0.179}{-0.703}{0.369}
\VF{0.000}{-8.000}{0.000}{-0.732}{0.376}
\VF{1.000}{-8.000}{0.179}{-0.703}{0.369}
\VF{2.000}{-8.000}{0.338}{-0.622}{0.349}
\VF{3.000}{-8.000}{0.460}{-0.501}{0.319}
\VF{4.000}{-8.000}{0.536}{-0.357}{0.280}
\VF{5.000}{-8.000}{0.567}{-0.208}{0.235}
\VF{6.000}{-8.000}{0.556}{-0.067}{0.187}
\VF{7.000}{-8.000}{0.513}{0.054}{0.138}
\VF{8.000}{-8.000}{0.447}{0.149}{0.090}
\VF{9.000}{-8.000}{0.371}{0.217}{0.044}
\VF{10.000}{-8.000}{0.292}{0.258}{0.000}
\VF{-10.000}{-7.000}{-0.229}{0.338}{0.020}
\VF{-9.000}{-7.000}{-0.321}{0.315}{0.066}
\VF{-8.000}{-7.000}{-0.420}{0.261}{0.115}
\VF{-7.000}{-7.000}{-0.514}{0.171}{0.167}
\VF{-6.000}{-7.000}{-0.591}{0.043}{0.222}
\VF{-5.000}{-7.000}{-0.632}{-0.119}{0.278}
\VF{-4.000}{-7.000}{-0.623}{-0.303}{0.332}
\VF{-3.000}{-7.000}{-0.551}{-0.488}{0.380}
\VF{-2.000}{-7.000}{-0.415}{-0.650}{0.419}
\VF{-1.000}{-7.000}{-0.224}{-0.762}{0.445}
\VF{0.000}{-7.000}{0.000}{-0.803}{0.453}
\VF{1.000}{-7.000}{0.224}{-0.762}{0.445}
\VF{2.000}{-7.000}{0.415}{-0.650}{0.419}
\VF{3.000}{-7.000}{0.551}{-0.488}{0.380}
\VF{4.000}{-7.000}{0.623}{-0.303}{0.332}
\VF{5.000}{-7.000}{0.632}{-0.119}{0.278}
\VF{6.000}{-7.000}{0.591}{0.043}{0.222}
\VF{7.000}{-7.000}{0.514}{0.171}{0.167}
\VF{8.000}{-7.000}{0.420}{0.261}{0.115}
\VF{9.000}{-7.000}{0.321}{0.315}{0.066}
\VF{10.000}{-7.000}{0.229}{0.338}{0.020}
\VF{-10.000}{-6.000}{-0.131}{0.407}{0.041}
\VF{-9.000}{-6.000}{-0.228}{0.412}{0.088}
\VF{-8.000}{-6.000}{-0.344}{0.387}{0.140}
\VF{-7.000}{-6.000}{-0.471}{0.319}{0.196}
\VF{-6.000}{-6.000}{-0.592}{0.197}{0.257}
\VF{-5.000}{-6.000}{-0.682}{0.021}{0.321}
\VF{-4.000}{-6.000}{-0.714}{-0.201}{0.386}
\VF{-3.000}{-6.000}{-0.663}{-0.442}{0.447}
\VF{-2.000}{-6.000}{-0.517}{-0.665}{0.497}
\VF{-1.000}{-6.000}{-0.285}{-0.825}{0.531}
\VF{0.000}{-6.000}{0.000}{-0.884}{0.543}
\VF{1.000}{-6.000}{0.285}{-0.825}{0.531}
\VF{2.000}{-6.000}{0.517}{-0.665}{0.497}
\VF{3.000}{-6.000}{0.663}{-0.442}{0.447}
\VF{4.000}{-6.000}{0.714}{-0.201}{0.386}
\VF{5.000}{-6.000}{0.682}{0.021}{0.321}
\VF{6.000}{-6.000}{0.592}{0.197}{0.257}
\VF{7.000}{-6.000}{0.471}{0.319}{0.196}
\VF{8.000}{-6.000}{0.344}{0.387}{0.140}
\VF{9.000}{-6.000}{0.228}{0.412}{0.088}
\VF{10.000}{-6.000}{0.131}{0.407}{0.041}
\VF{-10.000}{-5.000}{0.000}{0.448}{0.064}
\VF{-9.000}{-5.000}{-0.086}{0.485}{0.113}
\VF{-8.000}{-5.000}{-0.203}{0.501}{0.166}
\VF{-7.000}{-5.000}{-0.352}{0.479}{0.225}
\VF{-6.000}{-5.000}{-0.522}{0.394}{0.290}
\VF{-5.000}{-5.000}{-0.683}{0.228}{0.363}
\VF{-4.000}{-5.000}{-0.790}{-0.024}{0.440}
\VF{-3.000}{-5.000}{-0.793}{-0.334}{0.517}
\VF{-2.000}{-5.000}{-0.655}{-0.649}{0.585}
\VF{-1.000}{-5.000}{-0.375}{-0.889}{0.632}
\VF{0.000}{-5.000}{0.000}{-0.981}{0.649}
\VF{1.000}{-5.000}{0.375}{-0.889}{0.632}
\VF{2.000}{-5.000}{0.655}{-0.649}{0.585}
\VF{3.000}{-5.000}{0.793}{-0.334}{0.517}
\VF{4.000}{-5.000}{0.790}{-0.024}{0.440}
\VF{5.000}{-5.000}{0.683}{0.228}{0.363}
\VF{6.000}{-5.000}{0.522}{0.394}{0.290}
\VF{7.000}{-5.000}{0.352}{0.479}{0.225}
\VF{8.000}{-5.000}{0.203}{0.501}{0.166}
\VF{9.000}{-5.000}{0.086}{0.485}{0.113}
\VF{10.000}{-5.000}{-0.000}{0.448}{0.064}
\VF{-10.000}{-4.000}{0.150}{0.445}{0.088}
\VF{-9.000}{-4.000}{0.092}{0.508}{0.138}
\VF{-8.000}{-4.000}{0.000}{0.566}{0.194}
\VF{-7.000}{-4.000}{-0.138}{0.607}{0.255}
\VF{-6.000}{-4.000}{-0.331}{0.600}{0.324}
\VF{-5.000}{-4.000}{-0.568}{0.501}{0.403}
\VF{-4.000}{-4.000}{-0.795}{0.265}{0.493}
\VF{-3.000}{-4.000}{-0.920}{-0.112}{0.590}
\VF{-2.000}{-4.000}{-0.841}{-0.561}{0.683}
\VF{-1.000}{-4.000}{-0.513}{-0.944}{0.752}
\VF{0.000}{-4.000}{0.000}{-1.099}{0.779}
\VF{1.000}{-4.000}{0.513}{-0.944}{0.752}
\VF{2.000}{-4.000}{0.841}{-0.561}{0.683}
\VF{3.000}{-4.000}{0.920}{-0.112}{0.590}
\VF{4.000}{-4.000}{0.795}{0.265}{0.493}
\VF{5.000}{-4.000}{0.568}{0.501}{0.403}
\VF{6.000}{-4.000}{0.331}{0.600}{0.324}
\VF{7.000}{-4.000}{0.138}{0.607}{0.255}
\VF{8.000}{-4.000}{-0.000}{0.566}{0.194}
\VF{9.000}{-4.000}{-0.092}{0.508}{0.138}
\VF{10.000}{-4.000}{-0.150}{0.445}{0.088}
\VF{-10.000}{-3.000}{0.296}{0.391}{0.110}
\VF{-9.000}{-3.000}{0.277}{0.462}{0.164}
\VF{-8.000}{-3.000}{0.234}{0.545}{0.223}
\VF{-7.000}{-3.000}{0.149}{0.635}{0.288}
\VF{-6.000}{-3.000}{0.000}{0.719}{0.361}
\VF{-5.000}{-3.000}{-0.243}{0.756}{0.445}
\VF{-4.000}{-3.000}{-0.587}{0.661}{0.543}
\VF{-3.000}{-3.000}{-0.940}{0.313}{0.660}
\VF{-2.000}{-3.000}{-1.067}{-0.300}{0.789}
\VF{-1.000}{-3.000}{-0.742}{-0.955}{0.900}
\VF{0.000}{-3.000}{0.000}{-1.251}{0.946}
\VF{1.000}{-3.000}{0.742}{-0.955}{0.900}
\VF{2.000}{-3.000}{1.067}{-0.300}{0.789}
\VF{3.000}{-3.000}{0.940}{0.313}{0.660}
\VF{4.000}{-3.000}{0.587}{0.661}{0.543}
\VF{5.000}{-3.000}{0.243}{0.756}{0.445}
\VF{6.000}{-3.000}{-0.000}{0.719}{0.361}
\VF{7.000}{-3.000}{-0.149}{0.635}{0.288}
\VF{8.000}{-3.000}{-0.234}{0.545}{0.223}
\VF{9.000}{-3.000}{-0.277}{0.462}{0.164}
\VF{10.000}{-3.000}{-0.296}{0.391}{0.110}
\VF{-10.000}{-2.000}{0.416}{0.289}{0.129}
\VF{-9.000}{-2.000}{0.436}{0.350}{0.186}
\VF{-8.000}{-2.000}{0.446}{0.427}{0.250}
\VF{-7.000}{-2.000}{0.436}{0.524}{0.320}
\VF{-6.000}{-2.000}{0.388}{0.646}{0.400}
\VF{-5.000}{-2.000}{0.266}{0.793}{0.491}
\VF{-4.000}{-2.000}{0.000}{0.933}{0.597}
\VF{-3.000}{-2.000}{-0.509}{0.921}{0.727}
\VF{-2.000}{-2.000}{-1.143}{0.381}{0.896}
\VF{-1.000}{-2.000}{-1.082}{-0.721}{1.000}
\VF{0.000}{-2.000}{0.000}{-1.300}{1.000}
\VF{1.000}{-2.000}{1.082}{-0.721}{1.000}
\VF{2.000}{-2.000}{1.143}{0.381}{0.896}
\VF{3.000}{-2.000}{0.509}{0.921}{0.727}
\VF{4.000}{-2.000}{-0.000}{0.933}{0.597}
\VF{5.000}{-2.000}{-0.266}{0.793}{0.491}
\VF{6.000}{-2.000}{-0.388}{0.646}{0.400}
\VF{7.000}{-2.000}{-0.436}{0.524}{0.320}
\VF{8.000}{-2.000}{-0.446}{0.427}{0.250}
\VF{9.000}{-2.000}{-0.436}{0.350}{0.186}
\VF{10.000}{-2.000}{-0.416}{0.289}{0.129}
\VF{-10.000}{-1.000}{0.495}{0.154}{0.141}
\VF{-9.000}{-1.000}{0.541}{0.188}{0.201}
\VF{-8.000}{-1.000}{0.590}{0.233}{0.268}
\VF{-7.000}{-1.000}{0.639}{0.294}{0.344}
\VF{-6.000}{-1.000}{0.684}{0.378}{0.431}
\VF{-5.000}{-1.000}{0.717}{0.499}{0.532}
\VF{-4.000}{-1.000}{0.710}{0.681}{0.653}
\VF{-3.000}{-1.000}{0.577}{0.961}{0.803}
\VF{-2.000}{-1.000}{0.000}{1.300}{1.000}
\VF{-1.000}{-1.000}{-1.233}{0.411}{1.000}
\VF{1.000}{-1.000}{1.233}{0.411}{1.000}
\VF{2.000}{-1.000}{-0.000}{1.300}{1.000}
\VF{3.000}{-1.000}{-0.577}{0.961}{0.803}
\VF{4.000}{-1.000}{-0.710}{0.681}{0.653}
\VF{5.000}{-1.000}{-0.717}{0.499}{0.532}
\VF{6.000}{-1.000}{-0.684}{0.378}{0.431}
\VF{7.000}{-1.000}{-0.639}{0.294}{0.344}
\VF{8.000}{-1.000}{-0.590}{0.233}{0.268}
\VF{9.000}{-1.000}{-0.541}{0.188}{0.201}
\VF{10.000}{-1.000}{-0.495}{0.154}{0.141}
\VF{-10.000}{0.000}{0.522}{-0.000}{0.145}
\VF{-9.000}{0.000}{0.578}{-0.000}{0.206}
\VF{-8.000}{0.000}{0.640}{-0.000}{0.275}
\VF{-7.000}{0.000}{0.711}{-0.000}{0.353}
\VF{-6.000}{0.000}{0.792}{-0.000}{0.442}
\VF{-5.000}{0.000}{0.889}{-0.000}{0.548}
\VF{-4.000}{0.000}{1.007}{-0.000}{0.678}
\VF{-3.000}{0.000}{1.159}{-0.000}{0.845}
\VF{-2.000}{0.000}{1.300}{-0.000}{1.000}
\VF{2.000}{0.000}{-1.300}{-0.000}{1.000}
\VF{3.000}{0.000}{-1.159}{-0.000}{0.845}
\VF{4.000}{0.000}{-1.007}{-0.000}{0.678}
\VF{5.000}{0.000}{-0.889}{-0.000}{0.548}
\VF{6.000}{0.000}{-0.792}{-0.000}{0.442}
\VF{7.000}{0.000}{-0.711}{-0.000}{0.353}
\VF{8.000}{0.000}{-0.640}{-0.000}{0.275}
\VF{9.000}{0.000}{-0.578}{-0.000}{0.206}
\VF{10.000}{0.000}{-0.522}{-0.000}{0.145}
\VF{-10.000}{1.000}{0.495}{-0.154}{0.141}
\VF{-9.000}{1.000}{0.541}{-0.188}{0.201}
\VF{-8.000}{1.000}{0.590}{-0.233}{0.268}
\VF{-7.000}{1.000}{0.639}{-0.294}{0.344}
\VF{-6.000}{1.000}{0.684}{-0.378}{0.431}
\VF{-5.000}{1.000}{0.717}{-0.499}{0.532}
\VF{-4.000}{1.000}{0.710}{-0.681}{0.653}
\VF{-3.000}{1.000}{0.577}{-0.961}{0.803}
\VF{-2.000}{1.000}{0.000}{-1.300}{1.000}
\VF{-1.000}{1.000}{-1.233}{-0.411}{1.000}
\VF{1.000}{1.000}{1.233}{-0.411}{1.000}
\VF{2.000}{1.000}{-0.000}{-1.300}{1.000}
\VF{3.000}{1.000}{-0.577}{-0.961}{0.803}
\VF{4.000}{1.000}{-0.710}{-0.681}{0.653}
\VF{5.000}{1.000}{-0.717}{-0.499}{0.532}
\VF{6.000}{1.000}{-0.684}{-0.378}{0.431}
\VF{7.000}{1.000}{-0.639}{-0.294}{0.344}
\VF{8.000}{1.000}{-0.590}{-0.233}{0.268}
\VF{9.000}{1.000}{-0.541}{-0.188}{0.201}
\VF{10.000}{1.000}{-0.495}{-0.154}{0.141}
\VF{-10.000}{2.000}{0.416}{-0.289}{0.129}
\VF{-9.000}{2.000}{0.436}{-0.350}{0.186}
\VF{-8.000}{2.000}{0.446}{-0.427}{0.250}
\VF{-7.000}{2.000}{0.436}{-0.524}{0.320}
\VF{-6.000}{2.000}{0.388}{-0.646}{0.400}
\VF{-5.000}{2.000}{0.266}{-0.793}{0.491}
\VF{-4.000}{2.000}{0.000}{-0.933}{0.597}
\VF{-3.000}{2.000}{-0.509}{-0.921}{0.727}
\VF{-2.000}{2.000}{-1.143}{-0.381}{0.896}
\VF{-1.000}{2.000}{-1.082}{0.721}{1.000}
\VF{0.000}{2.000}{0.000}{1.300}{1.000}
\VF{1.000}{2.000}{1.082}{0.721}{1.000}
\VF{2.000}{2.000}{1.143}{-0.381}{0.896}
\VF{3.000}{2.000}{0.509}{-0.921}{0.727}
\VF{4.000}{2.000}{-0.000}{-0.933}{0.597}
\VF{5.000}{2.000}{-0.266}{-0.793}{0.491}
\VF{6.000}{2.000}{-0.388}{-0.646}{0.400}
\VF{7.000}{2.000}{-0.436}{-0.524}{0.320}
\VF{8.000}{2.000}{-0.446}{-0.427}{0.250}
\VF{9.000}{2.000}{-0.436}{-0.350}{0.186}
\VF{10.000}{2.000}{-0.416}{-0.289}{0.129}
\VF{-10.000}{3.000}{0.296}{-0.391}{0.110}
\VF{-9.000}{3.000}{0.277}{-0.462}{0.164}
\VF{-8.000}{3.000}{0.234}{-0.545}{0.223}
\VF{-7.000}{3.000}{0.149}{-0.635}{0.288}
\VF{-6.000}{3.000}{0.000}{-0.719}{0.361}
\VF{-5.000}{3.000}{-0.243}{-0.756}{0.445}
\VF{-4.000}{3.000}{-0.587}{-0.661}{0.543}
\VF{-3.000}{3.000}{-0.940}{-0.313}{0.660}
\VF{-2.000}{3.000}{-1.067}{0.300}{0.789}
\VF{-1.000}{3.000}{-0.742}{0.955}{0.900}
\VF{0.000}{3.000}{0.000}{1.251}{0.946}
\VF{1.000}{3.000}{0.742}{0.955}{0.900}
\VF{2.000}{3.000}{1.067}{0.300}{0.789}
\VF{3.000}{3.000}{0.940}{-0.313}{0.660}
\VF{4.000}{3.000}{0.587}{-0.661}{0.543}
\VF{5.000}{3.000}{0.243}{-0.756}{0.445}
\VF{6.000}{3.000}{-0.000}{-0.719}{0.361}
\VF{7.000}{3.000}{-0.149}{-0.635}{0.288}
\VF{8.000}{3.000}{-0.234}{-0.545}{0.223}
\VF{9.000}{3.000}{-0.277}{-0.462}{0.164}
\VF{10.000}{3.000}{-0.296}{-0.391}{0.110}
\VF{-10.000}{4.000}{0.150}{-0.445}{0.088}
\VF{-9.000}{4.000}{0.092}{-0.508}{0.138}
\VF{-8.000}{4.000}{0.000}{-0.566}{0.194}
\VF{-7.000}{4.000}{-0.138}{-0.607}{0.255}
\VF{-6.000}{4.000}{-0.331}{-0.600}{0.324}
\VF{-5.000}{4.000}{-0.568}{-0.501}{0.403}
\VF{-4.000}{4.000}{-0.795}{-0.265}{0.493}
\VF{-3.000}{4.000}{-0.920}{0.112}{0.590}
\VF{-2.000}{4.000}{-0.841}{0.561}{0.683}
\VF{-1.000}{4.000}{-0.513}{0.944}{0.752}
\VF{0.000}{4.000}{0.000}{1.099}{0.779}
\VF{1.000}{4.000}{0.513}{0.944}{0.752}
\VF{2.000}{4.000}{0.841}{0.561}{0.683}
\VF{3.000}{4.000}{0.920}{0.112}{0.590}
\VF{4.000}{4.000}{0.795}{-0.265}{0.493}
\VF{5.000}{4.000}{0.568}{-0.501}{0.403}
\VF{6.000}{4.000}{0.331}{-0.600}{0.324}
\VF{7.000}{4.000}{0.138}{-0.607}{0.255}
\VF{8.000}{4.000}{-0.000}{-0.566}{0.194}
\VF{9.000}{4.000}{-0.092}{-0.508}{0.138}
\VF{10.000}{4.000}{-0.150}{-0.445}{0.088}
\VF{-10.000}{5.000}{0.000}{-0.448}{0.064}
\VF{-9.000}{5.000}{-0.086}{-0.485}{0.113}
\VF{-8.000}{5.000}{-0.203}{-0.501}{0.166}
\VF{-7.000}{5.000}{-0.352}{-0.479}{0.225}
\VF{-6.000}{5.000}{-0.522}{-0.394}{0.290}
\VF{-5.000}{5.000}{-0.683}{-0.228}{0.363}
\VF{-4.000}{5.000}{-0.790}{0.024}{0.440}
\VF{-3.000}{5.000}{-0.793}{0.334}{0.517}
\VF{-2.000}{5.000}{-0.655}{0.649}{0.585}
\VF{-1.000}{5.000}{-0.375}{0.889}{0.632}
\VF{0.000}{5.000}{0.000}{0.981}{0.649}
\VF{1.000}{5.000}{0.375}{0.889}{0.632}
\VF{2.000}{5.000}{0.655}{0.649}{0.585}
\VF{3.000}{5.000}{0.793}{0.334}{0.517}
\VF{4.000}{5.000}{0.790}{0.024}{0.440}
\VF{5.000}{5.000}{0.683}{-0.228}{0.363}
\VF{6.000}{5.000}{0.522}{-0.394}{0.290}
\VF{7.000}{5.000}{0.352}{-0.479}{0.225}
\VF{8.000}{5.000}{0.203}{-0.501}{0.166}
\VF{9.000}{5.000}{0.086}{-0.485}{0.113}
\VF{10.000}{5.000}{-0.000}{-0.448}{0.064}
\VF{-10.000}{6.000}{-0.131}{-0.407}{0.041}
\VF{-9.000}{6.000}{-0.228}{-0.412}{0.088}
\VF{-8.000}{6.000}{-0.344}{-0.387}{0.140}
\VF{-7.000}{6.000}{-0.471}{-0.319}{0.196}
\VF{-6.000}{6.000}{-0.592}{-0.197}{0.257}
\VF{-5.000}{6.000}{-0.682}{-0.021}{0.321}
\VF{-4.000}{6.000}{-0.714}{0.201}{0.386}
\VF{-3.000}{6.000}{-0.663}{0.442}{0.447}
\VF{-2.000}{6.000}{-0.517}{0.665}{0.497}
\VF{-1.000}{6.000}{-0.285}{0.825}{0.531}
\VF{0.000}{6.000}{0.000}{0.884}{0.543}
\VF{1.000}{6.000}{0.285}{0.825}{0.531}
\VF{2.000}{6.000}{0.517}{0.665}{0.497}
\VF{3.000}{6.000}{0.663}{0.442}{0.447}
\VF{4.000}{6.000}{0.714}{0.201}{0.386}
\VF{5.000}{6.000}{0.682}{-0.021}{0.321}
\VF{6.000}{6.000}{0.592}{-0.197}{0.257}
\VF{7.000}{6.000}{0.471}{-0.319}{0.196}
\VF{8.000}{6.000}{0.344}{-0.387}{0.140}
\VF{9.000}{6.000}{0.228}{-0.412}{0.088}
\VF{10.000}{6.000}{0.131}{-0.407}{0.041}
\VF{-10.000}{7.000}{-0.229}{-0.338}{0.020}
\VF{-9.000}{7.000}{-0.321}{-0.315}{0.066}
\VF{-8.000}{7.000}{-0.420}{-0.261}{0.115}
\VF{-7.000}{7.000}{-0.514}{-0.171}{0.167}
\VF{-6.000}{7.000}{-0.591}{-0.043}{0.222}
\VF{-5.000}{7.000}{-0.632}{0.119}{0.278}
\VF{-4.000}{7.000}{-0.623}{0.303}{0.332}
\VF{-3.000}{7.000}{-0.551}{0.488}{0.380}
\VF{-2.000}{7.000}{-0.415}{0.650}{0.419}
\VF{-1.000}{7.000}{-0.224}{0.762}{0.445}
\VF{0.000}{7.000}{0.000}{0.803}{0.453}
\VF{1.000}{7.000}{0.224}{0.762}{0.445}
\VF{2.000}{7.000}{0.415}{0.650}{0.419}
\VF{3.000}{7.000}{0.551}{0.488}{0.380}
\VF{4.000}{7.000}{0.623}{0.303}{0.332}
\VF{5.000}{7.000}{0.632}{0.119}{0.278}
\VF{6.000}{7.000}{0.591}{-0.043}{0.222}
\VF{7.000}{7.000}{0.514}{-0.171}{0.167}
\VF{8.000}{7.000}{0.420}{-0.261}{0.115}
\VF{9.000}{7.000}{0.321}{-0.315}{0.066}
\VF{10.000}{7.000}{0.229}{-0.338}{0.020}
\VF{-10.000}{8.000}{-0.292}{-0.258}{0.000}
\VF{-9.000}{8.000}{-0.371}{-0.217}{0.044}
\VF{-8.000}{8.000}{-0.447}{-0.149}{0.090}
\VF{-7.000}{8.000}{-0.513}{-0.054}{0.138}
\VF{-6.000}{8.000}{-0.556}{0.067}{0.187}
\VF{-5.000}{8.000}{-0.567}{0.208}{0.235}
\VF{-4.000}{8.000}{-0.536}{0.357}{0.280}
\VF{-3.000}{8.000}{-0.460}{0.501}{0.319}
\VF{-2.000}{8.000}{-0.338}{0.622}{0.349}
\VF{-1.000}{8.000}{-0.179}{0.703}{0.369}
\VF{0.000}{8.000}{0.000}{0.732}{0.376}
\VF{1.000}{8.000}{0.179}{0.703}{0.369}
\VF{2.000}{8.000}{0.338}{0.622}{0.349}
\VF{3.000}{8.000}{0.460}{0.501}{0.319}
\VF{4.000}{8.000}{0.536}{0.357}{0.280}
\VF{5.000}{8.000}{0.567}{0.208}{0.235}
\VF{6.000}{8.000}{0.556}{0.067}{0.187}
\VF{7.000}{8.000}{0.513}{-0.054}{0.138}
\VF{8.000}{8.000}{0.447}{-0.149}{0.090}
\VF{9.000}{8.000}{0.371}{-0.217}{0.044}
\VF{10.000}{8.000}{0.292}{-0.258}{0.000}
\VF{-10.000}{9.000}{-0.341}{-0.189}{0.000}
\VF{-9.000}{9.000}{-0.388}{-0.129}{0.021}
\VF{-8.000}{9.000}{-0.444}{-0.058}{0.064}
\VF{-7.000}{9.000}{-0.486}{0.034}{0.107}
\VF{-6.000}{9.000}{-0.507}{0.143}{0.150}
\VF{-5.000}{9.000}{-0.500}{0.262}{0.192}
\VF{-4.000}{9.000}{-0.460}{0.383}{0.229}
\VF{-3.000}{9.000}{-0.386}{0.496}{0.262}
\VF{-2.000}{9.000}{-0.279}{0.588}{0.286}
\VF{-1.000}{9.000}{-0.146}{0.648}{0.302}
\VF{0.000}{9.000}{0.000}{0.670}{0.307}
\VF{1.000}{9.000}{0.146}{0.648}{0.302}
\VF{2.000}{9.000}{0.279}{0.588}{0.286}
\VF{3.000}{9.000}{0.386}{0.496}{0.262}
\VF{4.000}{9.000}{0.460}{0.383}{0.229}
\VF{5.000}{9.000}{0.500}{0.262}{0.192}
\VF{6.000}{9.000}{0.507}{0.143}{0.150}
\VF{7.000}{9.000}{0.486}{0.034}{0.107}
\VF{8.000}{9.000}{0.444}{-0.058}{0.064}
\VF{9.000}{9.000}{0.388}{-0.129}{0.021}
\VF{10.000}{9.000}{0.341}{-0.189}{0.000}
\VF{-10.000}{10.000}{-0.370}{-0.123}{0.000}
\VF{-9.000}{10.000}{-0.386}{-0.058}{0.000}
\VF{-8.000}{10.000}{-0.424}{0.013}{0.037}
\VF{-7.000}{10.000}{-0.449}{0.097}{0.076}
\VF{-6.000}{10.000}{-0.455}{0.192}{0.114}
\VF{-5.000}{10.000}{-0.438}{0.292}{0.150}
\VF{-4.000}{10.000}{-0.395}{0.391}{0.182}
\VF{-3.000}{10.000}{-0.325}{0.480}{0.209}
\VF{-2.000}{10.000}{-0.232}{0.551}{0.229}
\VF{-1.000}{10.000}{-0.121}{0.598}{0.242}
\VF{0.000}{10.000}{0.000}{0.614}{0.246}
\VF{1.000}{10.000}{0.121}{0.598}{0.242}
\VF{2.000}{10.000}{0.232}{0.551}{0.229}
\VF{3.000}{10.000}{0.325}{0.480}{0.209}
\VF{4.000}{10.000}{0.395}{0.391}{0.182}
\VF{5.000}{10.000}{0.438}{0.292}{0.150}
\VF{6.000}{10.000}{0.455}{0.192}{0.114}
\VF{7.000}{10.000}{0.449}{0.097}{0.076}
\VF{8.000}{10.000}{0.424}{0.013}{0.037}
\VF{9.000}{10.000}{0.386}{-0.058}{0.000}
\VF{10.000}{10.000}{0.370}{-0.123}{0.000}
}
  \marcoMMA{10}{5}{$x$ (u.\,arb.)}{$z$ (u.\,arb.)}
\end{tikzpicture}\\[1pt]{\small\scshape (c)}
\end{minipage}
\caption{Término cuadrupolar con simetría de revolución, $\phimul{2}\propto Q_{zz}\Pl{2}(\cos\theta)/r^{3}$ con $Q_{zz}>0$ (\cref{teo:revolucion}). (a)~Equipotenciales: los
lóbulos polares (potencial positivo) son más extensos que los ecuatoriales (negativo) porque $\Pl{2}(\pm1)=1$ y $\Pl{2}(0)=-\frac12$. (b)~Líneas de campo: salen a lo largo
del eje $z$, entran por el plano ecuatorial y cruzan el cono nodal. (c)~Campo eléctrico, con la misma escala de color que la \cref{fig:dipolo-mma}: cae como $1/r^{4}$,
aún más deprisa que el dipolar. El icono rojo representa dos dipolos opuestos. Recreación de~\cite{FIGURAS-T2}.}
\label{fig:cuadrupolo-mma}
\end{figure}
